%
% This file documents the biblatex-chicago package, which allows users
% of the biblatex package to format references according to the
% Chicago Manual of Style, 17th edition.
%
\documentclass[a4paper]{article}
% \usepackage[T1]{fontenc}
% \usepackage{textcomp}
% \usepackage[latin1]{inputenc}
\usepackage{fontspec}
\usepackage[american]{babel}
\usepackage[autostyle]{csquotes}
\usepackage{tabularx}
\usepackage{threeparttable}
\usepackage{booktabs}
\usepackage{afterpage}
%\usepackage[showframe=true]{geometry}
% \usepackage{vmargin}
% \setpapersize{A4}
% \setmarginsrb{1.65in}{.9in}{1.75in}{.6in}{0pt}{0pt}{12pt}{24pt}
\setlength{\marginparwidth}{1in}
\setlength{\belowcaptionskip}{1em}
\usepackage[colorlinks,urlcolor=blue,linkcolor=blue,%
filecolor=Teal]{hyperref}
% \usepackage{lmodern}
% \usepackage[scaled=0.9]{ClearSans}
% \usepackage[p]{zlmtt}
% \usepackage{gentium}
\setmainfont{GentiumPlus-Regular.ttf}[
ItalicFont = GentiumPlus-Italic.ttf,
BoldFont = GentiumPlus-Bold.ttf,
BoldItalicFont = GentiumPlus-BoldItalic.ttf]
\setsansfont{ClearSans-Regular.ttf}[
BoldFont = ClearSans-Bold.ttf,
ItalicFont = ClearSans-Italic.ttf,
BoldItalicFont = ClearSans-BoldItalic.ttf,
Scale = MatchLowercase]
\setmonofont{lmmonoprop10-regular.otf}[
BoldFont = lmmonoproplt10-bold.otf,
HyphenChar={-}]
%\usepackage[osf]{mathpazo}
%\usepackage[scaled]{helvet}
\usepackage[hyperref,svgnames]{xcolor}
%\usepackage[dvips]{xcolor}
\newcommand{\mycolor}[1]{\textcolor[HTML]{228B22}{#1}}
\usepackage{multicol}
% Some generic settings.
\newcommand{\cmd}[1]{\texttt{\textbackslash #1}}
\newcommand{\hc}{:\hspace{0pt}}
\hyphenation{text-cites mit-tel-pleis-to-zae-ne har-ley cal-li-ma-chus
  chi-ca-go au-to-cite case-changer beg-en-try note-ref-in-tro
  cite-in-cite ores-tes post-script car-to-gra-phy da-ta-bases
  ar-kan-sas pseu-do-nym head-less-cite da-ta-base bib-name-dash}
\setlength{\parindent}{0pt}
\usepackage{typearea}
\areaset[current]{360pt}{712pt}
%\special{dvipdfmx:config C 0x10}
\newcommand{\mymarginpar}[1]{\marginpar{\flushright#1}}
\newcommand{\colmarginpar}[1]{\mymarginpar{\mycolor{#1}}}
\newlength\myparskip
\newlength\myhalfskip
\myparskip=12pt plus 4pt minus 2pt
\myhalfskip=6pt plus 2pt minus 1pt
\newcommand{\mybigspace}{\vspace{\myparskip}}
\newcommand{\mylittlespace}{\vspace{\myhalfskip}}
\newcolumntype{L}{>{\raggedright}p}
\newcolumntype{H}{>{\sffamily\bfseries}l}
%\renewcommand{\textfraction}{0.0}
%\renewcommand{\topfraction}{1.0}
\setlength{\intextsep}{0pt}
\makeatletter
\renewcommand{\section}{\@startsection
  {section}%
  {1}%
  {0mm}%
  {\baselineskip}%
  {\baselineskip}%
  {\sffamily\normalsize\bfseries}}%
\renewcommand{\subsection}{\@startsection
  {subsection}%
  {1}%
  {0mm}%
  {\baselineskip}%
  {.5\baselineskip}%
  {\sffamily\normalsize\bfseries}}%
\renewcommand{\subsubsection}{\@startsection
  {subsubsection}%
  {1}%
  {0mm}%
  {\baselineskip}%
  {.5\baselineskip}%
  {\sffamily\normalsize\bfseries}}%
\renewcommand{\paragraph}{\@startsection
  {paragraph}%
  {1}%
  {\z@}%
  {\baselineskip}%3.25ex \@plus1ex \@minus.2ex}%
  {0mm}%
  {}}%
\long\def\@makecaption#1#2{%
  \vskip -3.5\abovecaptionskip
  \sbox\@tempboxa{\sffamily\bfseries #1: #2}%
  \ifdim \wd\@tempboxa >\hsize
    #1: #2\par
  \else
    \global \@minipagefalse
    \hb@xt@\hsize{\hfil\box\@tempboxa\hfil}%
  \fi
  \vskip\belowcaptionskip}
\def\ps@longtable{\setlength{\headsep}{36pt}\let\@mkboth\@gobbletwo
     \let\@oddhead\@empty\def\@oddfoot{\reset@font\hfil\thepage
     \hfil}\let\@evenhead\@empty\let\@evenfoot\@oddfoot}
\makeatother
\begin{document}
\begin{center}
  \sffamily\large\bfseries The biblatex-chicago package: \\
  Style files for biblatex

\vspace{.3\baselineskip}
\sffamily\normalsize\bfseries David Fussner\qquad Version 2.3b\\
\href{mailto:djf027@googlemail.com}{djf027@googlemail.com}\\ \today

\end{center} 
\setcounter{tocdepth}{3}
\begin{multicols}{2}
\footnotesize
\tableofcontents
\listoftables
\end{multicols}
\normalsize
\vspace{-.5\baselineskip}
\section{Notice}
\label{sec:Notice}

\textbf{This package is still under active development.  The
  \textsf{biblatex} package by Philipp Lehman, Philip Kime, Moritz
  Wemheuer, Audrey Boruvka, and Joseph Wright is now quite stable, but
  my task of incorporating the many enhancements it has accumulated in
  recent releases is ongoing.  The \textsf{biblatex-chicago} package
  itself now implements the 17th edition of the \emph{Chicago Manual
    of Style}, though I have made it possible to continue to use the
  16th edition files if that is imperative for you.  The package
  relies heavily, in all styles, on using \textsf{biber} as its
  backend; other backends will not work properly.}

\mylittlespace \textbf{I have tried to implement as much of the
  \emph{Manual's} specification as possible, though undoubtedly some
  gaps remain.  If it seems like this package could be of use to you,
  yet it doesn't do something you need/want it to do, please feel free
  to let me know, and of course any suggestions for solving problems
  more elegantly or accurately would be most welcome.}

\enlargethispage{\baselineskip}

\mylittlespace\textbf{Important Note:} If you have used
\textsf{biblatex-chicago} before, especially if you've been using
anything earlier than version 2.0, please make sure you have read the
RELEASE file that came with the package.  It details the changes
you'll need to make to your .bib database in order for it to work
properly with this release.  If you have continued to use the
16th-edition styles, I do strongly recommend that you switch to the
new edition, as I am deprecating the previous edition styles, and will
delete them in a future release.  Any required changes, as you can see
from the RELEASE file, in the main involve \emph{additions} to the
specification, with required alterations to your existing .bib
databases actually being rather rare, which should help ease the
transition.  If you are new to these styles or to \textsf{biblatex}
itself, please do continue reading at least the following section.

\section{Quickstart}
\label{quickstart}
\reversemarginpar

The \textsf{biblatex-chicago} package is designed for writers who wish
to use \LaTeX\ and \textsf{biblatex}, and who either want or need to
format their references according to one of the specifications defined
by the \emph{Chicago Manual of Style}.  This package includes two
versions of the \emph{Manual's} \enquote{author-date} system, favored
by many disciplines in the sciences and social sciences, and also its
\enquote{notes \&\ bibliography} style, generally favored in the
humanities.  The latter code produces a full reference in a first
footnote, shorter references in subsequent notes, and a full reference
in the bibliography. Some authors prefer to use the shorter note form
even for the first occurrence, relying on the bibliography to provide
the full information.  This, too, is supported by the code.  The
author-date styles produce a short, in-text citation inside
parentheses --- (Author Year) --- keyed to a list of references where
entries start with the same name and year.

\mylittlespace The documentation you are reading covers all three of
these Chicago styles and their variants.  I recommend that users new
to the package read this Quickstart section first, perhaps then
passing on to whichever of the two introductory files,
\href{file:cms-notes-intro.pdf}{\textsf{cms-notes-intro.pdf}} or
\href{file:cms-dates-intro.pdf}{\textsf{cms-dates-intro.pdf}}, is
relevant to their needs, returning here afterward for more details on
those parts of the functionality concerning which they still have
questions.  Much of what follows is relevant to all users, but I have
decided, after some experimentation, to keep the instructions for the
two author-date styles separate from those pertaining to the notes \&\
bibliography style, at least in sections~\ref{sec:Spec} and
\ref{sec:authdate}.  Information provided under one style will often
duplicate that found under the other, but efficiency's loss should, I
hope, be clarity's gain, and much of what you learn using one style
will be applicable without alteration to the other.  Within the
author-date section, the \textsf{authordate-trad} information really
only appears separately in section~\ref{sec:fields:authdate}, s.v.\
\enquote{title.} Throughout the documentation, any \mycolor{green}
text \colmarginpar{\textsf{New!}} indicates something \mycolor{new} in
this release, while \textcolor{Teal}{\textbf{blue-green}} text is a
clickable link to an external document.

\enlargethispage{\baselineskip}

\mylittlespace Here's a list of things you will need in order to use
\textsf{biblatex-chicago}:

\begin{itemize}{}{}
\item The \textsf{biblatex} package, of course!  The current version
  --- 3.20 at the time of writing --- has received extensive testing,
  and contains features and bug fixes upon which my code relies.
  Earlier versions may work, but I recommend the latest.
  \textsf{Biblatex} requires several packages, and it strongly
  recommends several more:
  \begin{itemize}{}{}
  \item \textsf{biber} --- the next-generation \textsc{Bib}\TeX\
    replacement by Philip Kime and Fran\c{c}ois Charette, available
    from SourceForge (required).  You should use the latest version,
    2.20, to work with \textsf{biblatex} 3.20 and
    \textsf{biblatex-chicago}; please note that any other backend will
    not produce accurate results.
  \item e-\TeX\ (required)
  \item \textsf{etoolbox} --- available from CTAN (required)
  \item \textsf{keyval} --- a standard package (required)
  \item \textsf{ifthen} --- a standard package (required)
  \item \textsf{url} --- a standard package (required).  In
    \textsf{biblatex-chicago.sty} you'll see that I set
    \cmd{urlstyle\{rm\}}.  If you happen to use a non-roman text font,
    you may want to set \cmd{urlstyle\{same\}} in your preamble
    instead.
  \item \textsf{babel} --- a standard package (\emph{strongly}
    recommended)
  \item \textsf{csquotes} --- available from CTAN (recommended).
    Please upgrade to the latest version of \textsf{csquotes} (5.2o).
  \end{itemize}
\item The standard \textsf{expl3} and \textsf{xparse} packages are
  loaded automatically for most users, and if they aren't
  \textsf{biblatex-chicago} does it for you.
\item With recent changes both to \textsf{biblatex} and to
  \textsf{biblatex-chicago}, \textsf{biblatex-chicago} itself now
  requires two packages, which are both loaded for you if you load
  \textsf{biblatex-chicago.sty}, but which you'll have to load
  manually if not.  They are:
  \begin{itemize}{}{}
  \item \textsf{nameref} --- a standard package, available
    in CTAN.
  \item \textsf{xstring} --- also standard and available in CTAN.
  \end{itemize}
\item The line:
  \begin{quote}
    \verb+\usepackage[notes,backend=biber]{biblatex-chicago}+
  \end{quote}
  in your document preamble to load the notes \&\ bibliography style,
  the line:
  \begin{quote}
    \verb+\usepackage[authordate,backend=biber]{biblatex-chicago}+
  \end{quote}
  to load the author-date style, or the line:
  \begin{quote}
    \verb+\usepackage[authordate-trad,backend=biber]{biblatex-chicago}+
  \end{quote}

  to load the traditional variant of the author-date style.  If you
  add \enquote{\texttt{16}} to any of the keys above, e.g.,
  \begin{quote}
    \verb+\usepackage[authordate16,backend=biber]{biblatex-chicago}+
  \end{quote}
  you can continue to use the 16th-edition styles, if that should
  prove necessary.  Any other options you usually pass to
  \textsf{biblatex} can be given to \textsf{biblatex-chicago} instead,
  but loading it this way sets up a large number of other parameters
  automatically, parameters whose absence may surprise you when
  processing your documents.  You can load the package via the usual
  \verb+\usepackage{biblatex}+, adding either
  \texttt{style=chicago-notes} or \texttt{style=chicago-authordate},
  but this is intended mainly for those, probably experienced users,
  who wish to set much of the low-level formatting of their documents
  themselves.  Please see sections~\ref{sec:loading} and
  \ref{sec:loading:auth} below for a fuller discussion of the issues
  involved here, and please also remember to load \textsf{xstring} and
  \textsf{nameref} manually if you use this latter method.
\item If you set
  \verb+\usepackage[notes,short,backend=biber]{biblatex-chicago}+
  you'll get the short note format even in the first reference of a
  notes \&\ bibliography document, letting the bibliography provide
  the full reference.
\item If you are accustomed to using the \textsf{natbib} or
  \textsf{mcite} compatibility options with \textsf{biblatex}, then
  you can continue to do so with \textsf{biblatex-chicago}.  If you
  are using \cmd{usepac\-kage\{biblatex-chicago\}} to load the style,
  then you can place (e.g.) \texttt{natbib} or \texttt{natbib=true}
  among its options to pass it through to \textsf{biblatex}.  Please
  see sections~\ref{sec:useropts} and \ref{sec:authuseropts}, below.
\item By far the simplest setup is to use \textsf{babel}, and to have
  \texttt{american} as the main text language.  (\textsf{Polyglossia}
  should work, too, and has received somewhat more testing for this
  release.)  As before, \textsf{babel}-less setups, and also those
  choosing \texttt{english} as the main text language, should work out
  of the box.  \textsf{Biblatex-chicago} also provides (at least
  partial) support for Brazilian Portuguese, British, Dutch, Finnish,
  French, German, Icelandic, Norwegian, Romanian, Spanish, and
  Swedish.  Please see below (section~\ref{sec:international}) for a
  fuller explanation of all the options.
\item \textsf{chicago-authordate.cbx, chicago-authordate-trad.cbx,
    chicago-dates-common.cbx},\break \textsf{chicago-authordate.bbx},
  \textsf{chicago-authordate-trad.bbx}, \textsf{chicago-notes.cbx},
  \textsf{chicago-notes.bbx}, \textsf{cms-american.lbx},
  \textsf{cms-brazilian.lbx}, \textsf{cms-british.lbx},
  \textsf{cms-dutch.lbx}, \textsf{cms-finnish.lbx},
  \textsf{cms-french.lbx}, \textsf{cms-german.lbx},
  \textsf{cms-icelandic.lbx}, \textsf{cms-ngerman.lbx},\break
  \textsf{cms-norsk.lbx}, \textsf{cms-norwegian.lbx},
  \textsf{cms-nynorsk.lbx}, \textsf{cms-romanian.lbx},
  \textsf{cms-spa-\break nish.lbx}, \textsf{cms-swedish.lbx},
  \textsf{biblatex-chicago.sty}, and \textsf{cmsdocs.sty}, all from
  \textsf{biblatex-chicago}, installed either in a system-wide \TeX\
  directory, or in the working directory where you keep your *.tex
  files.  The .zip file from CTAN contains subdirectories to help keep
  the growing number of files organized, so the files listed above can
  be found in the \texttt{latex/} subdirectory.  If you install in a
  system-wide directory, I suggest a standard layout like
  \texttt{<TEXMFLOCAL>\slash
    tex/latex/biblatex-contrib/biblatex-chicago}, where\
  \texttt{<TEXMFLOCAL>} is the root of your local \TeX\ installation
  --- for example, and depending on your operating system, something
  like \texttt{/usr/share/texmflocal},
  \texttt{/usr/local/share/texmf}, or \texttt{C:\textbackslash{}Local
    TeX Files\textbackslash}.  Then you can copy the contents of the
  \texttt{latex/} directory there.  (If you install into your working
  directory, then you'll need to copy the files directly there,
  without subdirectories.)  Of course, if you've placed them anywhere
  in the \texttt{texmf} tree, you'll need to update the file name
  database to make sure \TeX\ can find them.
\item The files \textsf{chicago-authordate16.cbx,
    chicago-authordate-trad16.cbx, chicago-dates-common16.cbx},
  \textsf{chicago-authordate16.bbx},
  \textsf{chicago-authordate-trad16.bbx},
  \textsf{chicago-notes16.cbx}, and \textsf{chicago-notes16.bbx},
  which, as their names suggest, allow you to continue using the
  16th-edition specifications alongside the most recent
  \textsf{biblatex}, if your project requires this.  They can be found
  in the same directory as the 17th-edition equivalents.
\item The dependent \LaTeX\ package \textsf{cmsendnotes.sty}, found
  with the previous files.  It offers additional functionality for
  those wishing to use the new \texttt{noteref} option with endnotes
  instead of footnotes.  See section~\ref{sec:noteref}, below, and
  also
  \href{file:cms-noteref-demo.pdf}{\textsf{cms-noteref-demo.pdf}}.
\item The file \mycolor{\textsf{cms.dbx}}, also to be found with the
  previous files, and a clone of the standard \textsf{biblatex} file
  \textsf{93-nameparts.dbx}.  It extends the default data model by
  adding new name parts, allowing the localized presentation of names
  from a number of linguistic and geographical contexts.  Please see
  section~\ref{sec:nameparts} below for details of how to use it.
\item The very clear and detailed documentation of the
  \textsf{biblatex} system, available in that package as
  \textsf{biblatex.pdf}.  Here the authors explain why you might want
  to use the system, the rules for constructing .bib files for it, and
  the (numerous) methods at your disposal for modifying the formatted
  output.
\item The files
  \href{file:cms-notes-intro.pdf}{\textsf{cms-notes-intro.pdf}},
  \href{file:cms-dates-intro.pdf}{\textsf{cms-dates-intro.pdf}},
  \href{file:cms-trad-appendix.pdf}{\textsf{cms-trad-appendix.pdf}},
  and\break
  \href{file:cms-noteref-demo.pdf}{\textsf{cms-noteref-demo.pdf}}, the
  first two of which contain introductions to some of the main
  features of the Chicago styles, while the third documents some of
  the alterations you might need to make to your .bib files to use the
  \texttt{trad} style.  The fourth gives a brief example of the usage
  of the \texttt{noteref} package option to the notes \&\ bibliography
  style.  All four are fully hyperlinked, the first three in
  particular allowing you easily to jump from notes or citations to an
  annotated bibliography or reference list, and thence to the .bib
  entries themselves.  If you ensure that all four are in the same
  directory as the document you are reading (the \TeX\ Live default),
  marginal links there will take you to further discussions here.  The
  file \textsf{cmsdocs.sty} contains code and kludges designed
  specifically for compiling \textsf{cms-dates-intro.tex},
  \textsf{cms-notes-intro.tex} and \textsf{cms-trad-appendix.tex}, so
  please \emph{do not} load it yourself anywhere else, as it redefines
  and interferes with some of the macros from the main package.
\item The annotated bibliography files \textsf{notes-test.bib} and
  \textsf{dates-test.bib}, and the not-yet-annotated
  \textsf{legal-test.bib}, all of which will acquaint you with many of
  the details on how to get started constructing your own .bib files
  for use with the three \textsf{biblatex-chicago} styles.
\item The files
  \href{file:cms-notes-sample.pdf}{\textsf{cms-notes-sample.pdf}},
  \href{file:cms-dates-sample.pdf}{\textsf{cms-dates-sample.pdf}},
  \href{file:cms-trad-sample.pdf}{\textsf{cms-trad-sample.pdf}}, and
  \href{file:cms-legal-sample.pdf}{\textsf{cms-legal-sample.pdf}}.
  The first shows how my system processes \textsf{notes-test.bib} and
  \textsf{cms-notes-sample.tex}, in both footnotes and bibliography,
  the second and third are the result of processing
  \textsf{dates-test.bib} with \textsf{cms-dates-sample.tex} or
  \textsf{cms-trad-sample.tex}, and the fourth processes
  \textsf{legal-test.bib} using \textsf{cms-legal-sample.tex}.  All of
  these files are in \texttt{doc/}, and the \textsf{sample} files,
  aside from the last named, are mainly included for testing purposes.
\item The file you are reading, \textsf{biblatex-chicago.pdf}, which
  aims to be as complete a description as possible of the rules for
  creating a .bib file that will, when processed by \LaTeX\ and
  \textsf{biber}, at least somewhat ease the burden when you try to
  implement the \emph{Chicago Manual of Style}'s specifications.
  These docs may seem frustratingly over-long, but remember that you
  only need to read the part(s) that apply to the style in which you
  are interested.  Much of the information in section~\ref{sec:Spec}
  is duplicated in section~\ref{sec:authdate}, so even if you have a
  need for multiple styles then using one will be excellent
  preparation for the others.  If you have used a previous version of
  this package, please pay particular attention to the sections on
  Obsolete and Deprecated Features, starting on
  page~\pageref{deprec:obsol}.  You will find the nineteen previous
  files in the \texttt{doc/} subdirectory once you've extracted
  \textsf{biblatex-chicago.zip}.  If you wish to place them in a
  system-wide directory, you can try:
  \texttt{<TEXMFLOCAL>/doc/latex/biblatex-contrib/biblatex-chicago},
  all the while remembering, of course, to update the file name
  database afterward.
\item Access to a copy of \emph{The Chicago Manual of Style} itself,
  which naturally contains incomparably more information than I can
  hope to present here.  It should always be your first port of call
  when any doubts arise as to exactly what the specification requires.
\end{itemize}

\subsection{License}
\label{sec:lppl}

Copyright \copyright\ 2008--2024 David Fussner.  This package is
author-maintained.  This work may be copied, distributed and/or
modified under the conditions of the \LaTeX\ Project Public License,
either version 1.3 of this license or (at your option) any later
version.  The latest version of this license is in
http://www.latex-project.org/lppl.txt and version 1.3 or later is part
of all distributions of \LaTeX\ version 2005/12/01 or later.  This
software is provided \enquote{as is,} without warranty of any kind,
either expressed or implied, including, but not limited to, the
implied warranties of merchantability and fitness for a particular
purpose.

\subsection{Acknowledgements}
\label{sec:acknowl}

Even a cursory glance at the cbx and bbx files in the package will
demonstrate how much of \textsf{biblatex's} code I've adapted and
re-used, and I've also followed some of the advice the authors have
given to others in the \texttt{comp.text.tex} newsgroup and on
\href{http://tex.stackexchange.com}{Stackexchange}.  In particular,
Philipp Lehman's advice on constructing \textsf{biblatex-chicago.sty}
was invaluable.  The code for formatting the footnote marks, and that
for printing the separating rule only after a run-on note, I've
adapted from the \textsf{footmisc} package by Robin Fairbairns, and
I've borrowed ideas for the \texttt{shorthandibid} option from Dominik
Wa{\ss}enhoven's \textsf{biblatex-dw} package.  I've adapted Audrey
Boruvka's \cmd{textcite} code from
\href{http://tex.stackexchange.com/questions/67837/citations-as-nouns-in-biblatex-chicago}{Stackexchange}
for the notes \&\ bibliography style, and her page-number-compression
code for both styles from the
\href{http://tex.stackexchange.com/questions/44492/biblatex-chicago-style-page-ranges}{same
  site}.  The dependent package \textsf{cmsendnotes.sty} contains code
by John Lavagnino and Ulrich Dirr.  I am very grateful to Mar\c{c}al
Orteu Punsola for the Spanish localization, to Patrick Danilevici for
the Romanian localization, to Wouter Lancee for the Dutch
localization, to Gustavo Barros for the Brazilian Portuguese
localization, to Stefan Bj\"{o}rk for the Swedish localization, to
Antti-Juhani Kaijahano for the Finnish localization, to Baldur
Kristinsson for providing the Icelandic localization, and to H{\aa}kon
Malmedal for the Norwegian localizations.  Kazuo Teramoto and Gildas
Hamel both sent patches to improve the package, and Arne Skj{\ae}rholt
provided some code to get me started on the \cmd{gentextcite}
commands.  If there's other \LaTeX\ code I've appropriated and
forgotten, please remind me.  Finally, Charles Schaum and Joseph
Reagle Jr.\ were both extremely generous with their help and advice
during the development of this package, and have both continued
indefatigably to test it and suggest needed improvements.  They were
particularly instrumental in encouraging the greatest possible degree
of compatibility with other \textsf{biblatex} styles.  Indeed, if the
task of adapting .bib files for use with the Chicago style seems
onerous now, you should have tried it before they got their hands on
it.

\section{Detailed Introduction}
\label{sec:Intro}

The \emph{Chicago Manual of Style}, implemented here in its 17th
edition, has long, in America at least, been one of the most
influential style guides for writers and publishers.  While one's
choices are now perhaps more extensive than ever, the \emph{Manual} at
least still provides a widely-recognized, and widely-utilized,
standard.  Indeed, when you add to this the sheer completeness of the
specification, its detailed instructions for referencing an enormous
number of different kinds of source material, then your choice (or
your publisher's choice) of the \emph{Manual} as a style guide seems
set to be a happy one.

\mylittlespace These very strengths, however, also make the style
difficult to use.  Admittedly, the \emph{Manual} emphasizes
consistency within a work, as opposed to rigid adherence to the
specification, at least when writer and publisher agree (14.4).
Sometimes a publisher demands such adherence, however, and anyone who
has attempted to produce it may well come away with the impression
that the specification itself is somewhat idiosyncratic in its
complexity, and I can't help but agree.  In the notes \&\ bibliography
style, the numerous differences in punctuation (and strings
identifying translators, editors, and the like) between footnotes and
bibliographies and the sometimes unusual location of page numbers; in
both styles the distinction between \enquote{journal} and
\enquote{magazine,} and the formatting differences between (e.g.)\ a
work from antiquity and one from the Renaissance, all of these tend to
overburden the writer who wants to comply with the standard.  Many of
these complexities, in truth, make the specification very nearly
impossible to implement straightforwardly in a system like
\textsf{biblatex} --- options multiply, each requiring a particular
sort of formatting, until one almost reaches the point of believing
that every individual book or article should have its own entry type.
Completeness and usability tend each to exclude the other, so the code
you have before you is a first attempt to achieve the former without
utterly sacrificing the latter.

\subsection*{What \textsf{biblatex-chicago} can and can't do}
\label{sec:bltries}

In short, the \textsf{biblatex} style files in this package try to
simplify the task of following the two Chicago specifications along
with their major variants.  In the notes \&\ bibliography style, the
two sorts of reference are treated separately (as are the two
different note forms, long and short), and you can choose always to
use the short note form, even at the first citation.  In the two
author-date styles, a series of options allows you to choose which
date (original printing, reprint, or both) appears in citations and at
the head of entries in the list of references.  In all styles,
punctuation is placed within quotation marks when needed, and as a
general rule as many parts of the style as possible are implemented as
transparently as possible.  Thanks to advice I received from Joseph
Reagle Jr.\ and Charles Schaum while these files were a work in
progress, I have attended as carefully as I can to backward
compatibility with the standard \textsf{biblatex} styles, and have
attempted to minimize both any changes you need to make to achieve
compliance with the Chicago specification, and indeed also any changes
necessary to switch between the two Chicago styles.  There is no doubt
room for improvement on this score, but even now, for a substantial
number of entries, any well-constructed .bib file that works for other
\textsf{biblatex} styles will \enquote{just work} under
\textsf{biblatex-chicago}.  By no means, however, will all entries in
such a .bib file produce equally satisfactory results.  Using this
documentation and the examples in \textsf{dates-test.bib} and/or
\textsf{notes-test.bib}, it should be possible to achieve compliance,
though the amount of revision necessary to do so will vary
significantly from .bib file to .bib file.  Conversely, once you have
created a database for \textsf{biblatex-chicago}, it won't necessarily
work well with other \textsf{biblatex} styles.  Indeed, most, quite
possibly all, users will find that they need to use special formatting
macros within the .bib file that would make such a file unusable in
any other context.  I strongly recommend, if you want to experiment
with this style, that you work on a copy of any .bib files that are
important to you, until you have determined that this package does
what you need/want it to do.

\mylittlespace When I first began working on this package, I made the
decision to alter as little as possible the main files from Lehman's
\textsf{biblatex}, so that my .bbx and .cbx files would use his
original \LaTeX\ .sty file and \textsc{Bib}\TeX\ .bst file.  As you
proceed, you will no doubt encounter some of the consequences of this
decision, with certain fields and entry types in the .bib file having
less-than-memorable names because I chose to use the supplementary
ones provided by \textsf{biblatex} rather than alter that package's
files.  With additions to the standard data model now possible, this
will be one of the directions for future development, particularly if
other styles are adopting certain broad conventions.  Needless to say,
I'm open to advice and suggestions on this score.

\section{The Specification:\ Notes\,\&\,Bibliography}
\label{sec:Spec}

In what follows, I attempt to explain all the parts of
\textsf{biblatex-chicago-notes} that might be considered somehow
\enquote{non standard,} at least with respect to the styles included
with \textsf{biblatex} itself, though in the section on entry fields I
have also duplicated a lot of the information in
\textsf{biblatex.pdf}, which I hope won't badly annoy expert users of
the system.  Headings in \mycolor{green} \colmarginpar{\textsf{New in
    this release}} indicate either material new to this release or old
material that has undergone significant revision.  Numbers in
parentheses refer to sections of the \emph{Chicago Manual of Style},
17th edition.  The file \textsf{notes-test.bib} contains many examples
from the \emph{Manual} which, when processed using
\textsf{biblatex-chicago-notes}, should produce the same output as you
see in the \emph{Manual} itself, or at least compliant output, where
the specifications are vague or open to interpretation, a state of
affairs which does sometimes occur.  I have provided
\href{file:cms-notes-sample.pdf}{\textsf{cms-notes-sample.pdf}}, which
shows how my system processes \textsf{notes-test.bib}, and I have also
included the reference keys from the latter file below in parentheses.

\subsection{Entry Types}
\label{sec:entrytypes}

The complete list of entry types currently available in
\textsf{biblatex-chicago-notes}, minus the odd \textsf{biblatex}
alias, is as follows: \textbf{article}, \textbf{artwork},
\textbf{audio}, \textbf{book}, \textbf{bookinbook}, \textbf{booklet},
\textbf{collection}, \textbf{customc}, \textbf{dataset},
\textbf{image}, \textbf{inbook}, \textbf{incollection},
\textbf{inproceedings}, \textbf{inreference}, \textbf{jurisdiction},
\textbf{legal}, \textbf{legislation}, \textbf{letter},
\textbf{manual}, \textbf{misc}, \textbf{music}, \textbf{mvbook},
\textbf{mvcollection}, \textbf{mvproceedings}, \textbf{mvreference},
\textbf{online} (with its alias \textbf{www}), \textbf{patent},
\textbf{performance}, \textbf{periodical}, \textbf{proceedings},
\textbf{reference}, \textbf{report} (with its alias
\textbf{techreport}), \textbf{review}, \textbf{standard},
\textbf{suppbook}, \textbf{suppcollection}, \textbf{suppperiodical},
\textbf{thesis} (with its aliases \textbf{mastersthesis} and
\textbf{phdthesis}), \textbf{unpublished}, and \textbf{video}.

\mylittlespace What follows is an attempt to specify all the
differences between these types and the standard provided by
\textsf{biblatex}.  If an entry type isn't discussed here, then it is
safe to assume that it works as it does in the standard styles.  In
general, I have attempted not to discuss specific entry fields here,
unless such a field is crucial to the overall operation of a given
entry type.  As a general and important rule, most entry types require
very few fields when you use \textsf{biblatex-chicago-notes}, so it
seemed to me better to gather information pertaining to fields in the
next section.

\mybigspace The \mymarginpar{\textbf{article}} \emph{Chicago Manual of
  Style} (14.164) recognizes three different sorts of periodical
publication, \enquote{journals,} \enquote{magazines,} and
\enquote{newspapers.} The first (14.166) is \enquote{a scholarly or
  professional periodical available mainly by subscription,} while the
second refers to \enquote{weekly or monthly (or sometimes daily)}
publications that are \enquote{available in individual issues at
  libraries or bookstores or newsstands or offered online, with or
  without a subscription.}  \enquote{Magazines} will tend to be
\enquote{more accessible to general readers,} and typically won't have
a volume number.  Indeed, by fiat I declare that should you need to
refer to a journal that identifies its issues mainly by year, month,
or week, then for the purposes of \textsf{biblatex-chicago-notes} such
a publication is a \enquote{magazine,} and not a \enquote{journal.}

\mylittlespace For articles in \enquote{journals} you can simply use
the traditional \textsc{Bib}\TeX\ --- and indeed \textsf{biblatex} ---
\textsf{article} entry type, which will work as expected and set off
the page numbers with a colon, as required by the \emph{Manual}.  If,
however, you need to refer to a \enquote{magazine} or a
\enquote{newspaper,} then you need to add an \textsf{entrysubtype}
field containing the exact string \texttt{magazine} or, now, its
synonym \texttt{newspaper}.  The main formatting differences between a
\texttt{magazine/newspaper} and a plain \textsf{article} are that the
year isn't placed within parentheses, and that page numbers are set
off by a comma rather than a colon.  Otherwise, the two sorts of
reference have much in common.  (For \textsf{article}, see
\emph{Manual} 14.168--87; batson, beattie:crime, friedman:learning,
garaud:gatine, garrett, hlatky:hrt, kern, lewis, loften:hamlet,
mcmillen:antebellum, rozner:liberation, saberhagen:beluga,
warr:ellison, white:callimachus. For \textsf{entrysubtype}
\texttt{magazine}, see 14.171, 14.188--200; assocpress:gun,
morgenson:market, reaves:rosen, stenger:privacy.)

\mylittlespace It gets worse.  The \emph{Manual} treats reviews (of
books, plays, performances, etc.) as a sort of recognizable subset of
\enquote{journals,} \enquote{magazines,} and \enquote{newspapers,}
distinguished mainly by the way one formats the title of the review
itself.  \textsf{Biblatex} provides a \textsf{review} entry type which
will handle a large subset of such citations, though not all.  The key
rule is this: if a review has a separate, non-generic title (gibbard;
osborne:poison) in addition to something that reads like
\enquote{review of \ldots,} then you need an \textsf{article} entry,
with or without the \texttt{magazine} \textsf{entrysubtype}, depending
on the sort of publication containing the review.  If the only title
is the generic \enquote{review of \ldots,} for example, then you'll
need the \textsf{review} entry type, with or without this same
\textsf{entrysubtype} toggle using \texttt{magazine}.  On
\textsf{review} entries, see below.

\mylittlespace In the case of a review with a specific as well as a
generic title, the former goes in the \textsf{title} field, and the
latter in the \textsf{titleaddon} field.  Standard \textsf{biblatex}
intends this field for use with additions to titles that may need to
be formatted differently from the titles themselves, and
\textsf{biblatex-chicago-notes} uses it in just this way, with the
additional wrinkle that it can, if needed, replace the \textsf{title}
entirely, and this in, effectively, any entry type, providing a fairly
powerful, if somewhat complicated, tool for getting \textsf{biblatex}
to do what you want.  Here, however, if all you need is a
\textsf{titleaddon}, then you want to switch to the \textsf{review}
type, where you can simply use the \textsf{title} field instead.

\mylittlespace \textsf{Biblatex-chicago} also, at the behest of
Bertold Schweitzer, supports the \textsf{relatedtype}
\texttt{reviewof}, which allows you to use the \textsf{related}
mechanism to provide information about the work being reviewed,
thereby simplifying how much information you need to provide in the
reviewing entry.  In particular, it relieves you of the need to
construct \textsf{titleaddon} or \textsf{title} fields like:
\texttt{review of \textbackslash mkbibemph\{Book Title\} by Author
  Name}, as the \textsf{related} entry's \textsf{title} automatically
provides the \textsf{titleaddon} in the \textsf{article} type and the
\textsf{title} in the \textsf{review} type, with the \textsf{related}
mechanism providing the connecting string.  This may be particularly
helpful if you need to cite multiple reviews of the same work; please
see section \ref{sec:related} for further details.

\mylittlespace No less than ten more things need explication here.
First, since the \emph{Manual} specifies that much of what goes into a
\textsf{titleaddon} field stays unformatted --- no italics, no
quotation marks --- this plain style is the default for such text,
which means that you'll have to format any titles within
\textsf{titleaddon} yourself, e.g., with \cmd{mkbibemph\{\}}.  (The
\textsf{related} mechanism just mentioned provides this
automatically.)  Second, the \emph{Manual} specifies a similar plain
style for the titles of other sorts of material found in
\enquote{magazines} and \enquote{newspapers,} e.g., obituaries,
letters to the editor, interviews, the names of regular columns, and
the like.  References may contain both the title of an individual
article and the name of the regular column, in which case the former
should go, as usual, in a \textsf{title} field, and the latter in
\textsf{titleaddon}.  As with reviews proper, if there is only the
generic title, then you want the \textsf{review} entry type.  (See
14.191, 14.195--96; morgenson:market, reaves:rosen.)

\mylittlespace Third, the \emph{Manual} has slightly complicated
instructions concerning \enquote{unsigned newspaper articles or
  features} (14.199).  First, it suggests that such pieces are
\enquote{best dealt with in text or notes.}  If, however, \enquote{a
  bibliography entry should be needed, the title of the newspaper
  stands in place of the author.}  The examples it provides,
therefore, suggest quite different treatments of the same material in
notes and bibliography, and they don't at any point that I can see
recommend a format for short notes.  I've implemented these
recommendations fairly literally, which means that in an
\textsf{article} entry, \textsf{entrysubtype} \texttt{magazine}, or in
a \textsf{review} entry, \textsf{entrysubtype} \texttt{magazine}, and
\emph{only} in such entries, a missing \textsf{author} field results
in the name of the periodical (in the \textsf{journaltitle} field)
being used as the missing author, but \emph{only} in the bibliography
and in short notes.  In long notes, the \textsf{title} will appear
first, before the \textsf{journaltitle}.  Note that the use of the
name of the newspaper as an author creates sorting issues in the
bibliography, issues that will mostly be solved for you if you use
\textsf{Biber} as the backend.  If you don't, or if the
\textsf{journaltitle} begins with a definite or indefinite article
with which you can't dispense, then you'll need a \textsf{sortkey}
field to ensure that the bibliography entry is alphabetized correctly.
(See lakeforester:pushcarts and, for the sorting issue,
\cmd{DeclareSortingTemplate} in section~\ref{sec:formatopts} below.)

\mylittlespace Fourth, Bertold Schweitzer has pointed out, following
the \emph{Manual} (14.183), that while an \textsf{issuetitle} often
has an \textsf{editor}, it is not too unusual for a \textsf{title} to
have, e.g., an \textsf{editor} and/or a \textsf{translator}.  In order
to allow as many permutations as possible on this theme, I have
brought the \textsf{article} entry type into line with most of the
other types in allowing the use of the \textsf{namea} and
\textsf{nameb} fields in order to associate an editor or a translator
specifically with the \textsf{title}.  The \textsf{editor} and
\textsf{translator} fields, in strict homology with other entry types,
are associated with the \textsf{issuetitle} if one is present, and
with the \textsf{title} otherwise.  The usual string concatenation
rules still apply --- cf.\ \textsf{editor} and \textsf{editortype} in
section~\ref{sec:entryfields}, below.

\mylittlespace Fifth, in certain fields just beginning your data with
a lowercase letter activates the mechanism for capitalizing that
letter depending on its context within a note or bibliography entry.
Please see \cmd{autocap} in section~\ref{sec:formatcommands} below for
the details, but both the \textsf{titleaddon} and \textsf{note} fields
are among those treating their data this way, and since both appear
regularly in \textsf{article} entries, I thought the problem merited a
mention here.

\mylittlespace Sixth, if you need to cite an entire issue of any sort
of periodical, rather than one article in an issue, then the
\textsf{periodical} entry type, once again with or without the
\texttt{magazine} toggle in \textsf{entrysubtype}, is what you'll
need.  (You can also use the \textsf{article} type, placing what would
normally be the \textsf{issuetitle} in the \textsf{title} field and
retaining the usual \textsf{journaltitle} field, but this arrangement
isn't compatible with standard \textsf{biblatex}.)  The \textsf{note}
field is where you place something like \enquote{special issue} or
\enquote{supplement} (with the small \enquote{s} enabling the
automatic capitalization routines), whether you are citing one article
or the whole issue (14.178--79; conley:fifthgrade, good:wholeissue).
Indeed, this is a somewhat specialized use of \textsf{note}, and if
you have other sorts of information you need to include in an
\textsf{article}, \textsf{periodical}, or \textsf{review} entry, then
you shouldn't put it in the \textsf{note} field, but rather in
\textsf{titleaddon} or perhaps \textsf{addendum} (brown:bremer).

\mylittlespace Seventh, if you wish to cite certain kinds of
television or radio broadcast, most notably interviews but perhaps
also news segments or other \enquote{journalistic} material, then the
\textsf{article} type, \textsf{entrysubtype} \texttt{magazine} is the
place for it.  The name of the program as a whole would go in
\textsf{journaltitle}, with the name of the episode or segment in
\textsf{title}, and the network's name in the \textsf{usera} field.
Of course, if the piece you are citing has only a generic name (an
interview, for example), then the \textsf{review} type would be the
best place for it (8.189, 14.213; see bundy:macneil for an example
of how this all might look in a .bib file.)  Other sorts of broadcast,
usually accessible through commercial recordings, would need one of
the audiovisual entry types, probably \textsf{audio}
(danforth:podcast) or \textsf{video} (friends:leia), while recordings
from archives fit best either into \textsf{online} or into
\textsf{misc} entries with an \textsf{entrysubtype} (coolidge:speech,
roosevelt:speech).

\mylittlespace Eighth, the \emph{Manual} (14.208) specifies that blogs
and other, similar online material should be presented like
\textsf{articles}, with \texttt{magazine} \textsf{entrysubtype}
(ellis:blog), and needn't appear in a bibliography.  The title of the
specific entry goes in \textsf{title}, the general title of the blog
goes in \textsf{journaltitle}, and the word \enquote{\texttt{blog}} in
the \textsf{location} field (though you could just use special
formatting in the \textsf{journaltitle} field itself, which may
sometimes be necessary).  The 17th edition specifies that
\enquote{blogs that are part of a larger publication should include
  the name of that publication.}  This usually involves a newspaper or
magazine which also publishes various blogs on its website, and it
means that such entries need a more general title than the
\textsf{journaltitle}.  It's not standard \textsf{biblatex} or
anything, but you can now put such information in \textsf{maintitle}
(with \textsf{mainsubtitle} and \textsf{maintitleaddon}, if needed),
but only in \textsf{article} and \textsf{review} entries with a
\texttt{magazine} \textsf{entrysubtype} (amlen:hoot).  To cite a whole
blog, you'll need the \textsf{periodical} entry type, with a
\textsf{title} instead of a \textsf{journaltitle}, along with a
(possible) \textsf{maintitle}.  Comments on blogs, with generic titles
like \enquote{comment on} or \enquote{reply to,} need a
\textsf{review} entry with the same \textsf{entrysubtype}.  Such
comments make particular use of the \textsf{eventdate} and of the
\textsf{nameaddon} fields; please see the documentation of
\textbf{review}, below, and also of the \textsf{relatedtype}
\texttt{commenton} in section~\ref{sec:related}.

\mylittlespace Ninth, the special \textsf{biblatex} field
\textsf{shortjournal} allows you to present shortened
\textsf{journaltitles} in \textsf{article}, \textsf{review}, and
\textsf{periodical} entries, as well as facilitating the creation of
lists of journal abbreviations in the manner of a \textsf{shorthand}
list.  Please see the documentation of \textbf{shortjournal} in
section~\ref{sec:entryfields} for all the details on how this works.

\mylittlespace Finally, the 17th edition (14.191) specifies that, for
news sites carrying \enquote{stories as they unfold, it may be
  appropriate to include a time stamp for an article that includes
  one.}  You can provide this by using the standard \textsf{biblatex}
time stamp format inside the \textsf{date} field, e.g.,
\texttt{2008-07-01T10:18:00}.  Since the \emph{Manual} prefers the
standard time zone initialisms, a separate \textsf{timezone} field
would be required if you want to provide one.

\mylittlespace If you're still with me, allow me to recommend that you
browse through \textsf{notes-test.bib} to get a feel for just how many
of the \emph{Manual}'s complexities the \textsf{article} and
\textsf{review} (and, indeed, \textsf{periodical}) types attempt to
address.  It may be that in future releases of
\textsf{biblatex-chicago-notes} I'll be able to simplify these
procedures somewhat, but in the meantime it might be of some comfort
that I have found in my own research that the unusual and/or limit
cases are really rather rare, and that the vast majority of sources
won't require any knowledge of these onerous details.

\mybigspace Arne \mymarginpar{\textbf{artwork}} Kjell Vikhagen
pointed out to me that none of the standard entry types were
straightforwardly adaptable when referring to visual artworks.  It's
unclear whether the \emph{Manual} (14.235) believes it necessary to
include them in the bibliographical apparatus at all, but it's easy
all the same to conceive of contexts in which a list of artworks
studied might be desirable, and \textsf{biblatex} includes entry types
for just this purpose, though the standard styles leave them
undefined.  \textsf{Biblatex-chicago} defines both \textsf{artwork}
and \textsf{image}, which are in fact now clones of each other, so you
can use either of them indifferently, the distinction existing only
for historical reasons.

\mylittlespace Constructing an entry is fairly straightforward.  As
one might expect, the artist goes in \textsf{author} and the name of
the work in \textsf{title}.  The \textsf{type} field is intended for
the medium --- e.g., oil on canvas, charcoal on paper, albumen print
--- and the \textsf{version} field might contain the state of an
etching.  You can place the dimensions of the work in \textsf{note},
and the current location in \textsf{organization},
\textsf{institution}, and/or \textsf{location}, in ascending order of
generality.  The \textsf{type} field, as in several other entry types,
uses \textsf{biblatex's} automatic capitalization routines, so if the
first word only needs a capital letter at the beginning of a sentence,
use lowercase in the .bib file and let \textsf{biblatex} handle it for
you.  (See \emph{Manual} 3.22, 8.198; leo:madonna, bedford:photo.)

\mylittlespace The 17th edition of the \emph{Manual} has included new
information in some of its examples, so I have added 4 new fields to
the drivers.  Alongside the usual \textsf{date} for the creation of a
work, you may also want to include the printing date of a particular
exemplar of a photograph or a print.  The system I have designed uses
the \emph{earlier} of the \textsf{date} and the \textsf{origdate} to
be the date of creation, and the \emph{later} to be the printing date.
The style will automatically prefix the printing date with the
localized \cmd{bibstring} \texttt{printed}, so if that's the wrong
string entirely then you can define \textsf{userd} any way you like to
change it.  If only \emph{one} of those two dates is available, it
will always serve as a creation date.

\mylittlespace One of the \emph{Manual's} examples is of a photograph
published in a periodical, and information about this publication
appears late in the entry, after the \textsf{type}.  I have included
the \textsf{howpublished} field so that you can give information about
the periodical (meaning that you'll have to format the title yourself
with \cmd{mkbibemph}), and the \textsf{eventdate} field for you to
provide the date of publication (mccurry:afghangirl).

\mylittlespace As a final complication, the \emph{Manual} (8.198) says
that \enquote{the names of works of antiquity \ldots\,are usually set
  in roman.}  If you should need to include such a work in the
reference apparatus, you can either define an \textsf{entrysubtype}
for an \textsf{artwork} entry --- anything will do --- or you could
use the \textsf{misc} entry type with an \textsf{entrysubtype}.
Assuming the complicated date handling I've just outlined isn't
required for such a work, in this instance the other fields in a
\textsf{misc} entry function pretty much as in \textsf{artwork}.

\mybigspace Following \mymarginpar{\textbf{audio}} the request of
Johan Nordstrom, I have included three entry types, all undefined by
the standard styles, designed to allow users to present audiovisual
sources in accordance with the Chicago specifications.  The
\emph{Manual's} presentation of such sources (14.261--68), though
admirably brief, seems to me somewhat inconsistent, though perhaps I'm
merely unable to spot the important regularities.  The proliferation
of online sources has made the task yet more complex, requiring the
inclusion of the \textbf{article}, the \textbf{online}, and even the
\textbf{misc} entry types, which see, under the audiovisual rubric.  I
shall attempt to delineate the main differences here, and though there
are likely to be occasions when your choice of entry type is not
obvious, at the very least \textsf{biblatex-chicago} should help you
maintain consistency.

\mylittlespace The \textbf{music} type is intended for all musical
recordings that do not have a video component.  This means, for
example, digital media (whether on CD or hard drive), vinyl records,
and tapes.  The \textbf{video} type includes most visual media,
whether it be films, TV shows, tapes and DVDs of the preceding or of
any sort of performance (including music), or online multimedia.  The
\emph{Manual's} treatment (14.267) of the latter suggests that online
video excerpts, short pieces, and interviews should generally use the
\textbf{online} type (horowitz:youtube, pollan:plant) or the
\textbf{article} type (harwood:biden, kessler:nyt), depending on
whether the pieces originate from an identifiably
\enquote{journalistic} outlet.  The \textbf{audio} type, our current
concern, fills gaps in the others, and presents its sources in a more
\enquote{book-like} manner.  Published musical scores need this type
--- unpublished ones would use \textsf{misc} with an
\textsf{entrysubtype} (shapey:partita) --- as do podcasts and such
favorite educational formats as the slideshow and the filmstrip
(danforth:podcast, greek:filmstrip, schubert:muellerin,
verdi:corsaro).  The \emph{Manual} (14.264) sometimes uses a similar
format for audio books (twain:audio), though, depending on the sorts
of publication facts you wish to present, this sort of material may
fall under \textsf{music} (auden:reading).  Dated audio recordings
that are part of an archive, online or no, may be presented either in
an \textbf{online} or in a \textbf{misc} entry with an
\textsf{entrysubtype}, the difference mainly being in just how closely
associated the \textsf{date} will be with the \textsf{title}
(coolidge:speech, roosevelt:speech).  Actual radio broadcasts (as
opposed to podcasts) pose something of a conundrum.  Interviews and
other sorts of \enquote{journalistic} material fit well into
\textsf{article} or \textsf{review} entries (14.213), but other sorts
of broadcast are not well represented in the \emph{Manual's} examples
(8.189), and what little there is suggests that, counter-intuitively,
the \textsf{video} type is the best fit, as it is well equipped to
present broadcasts of any sort.

\mylittlespace Once you've accepted the analogy of composer to
\textsf{author}, constructing an \textsf{audio} entry should be fairly
straightforward, since many of the fields function just as they do in
\textsf{book} or \textsf{inbook} entries.  Indeed, please note that I
compare it to both these other types as, in common with the other
audiovisual types, \textsf{audio} has to do double duty as an analogue
for both books and collections, so while there will normally be an
\textsf{author}, a \textsf{title}, a \textsf{publisher}, a
\textsf{date}, and a \textsf{location}, there may also be a
\textsf{booktitle} and/or a \textsf{maintitle} --- see
schubert:muellerin for an entry that uses all three in citing one song
from a cycle.  (As with the \textsf{music} and \textsf{video} types,
you can cite an individual piece separate from any large collection by
using the \textsf{title} field and by defining an
\textsf{entrysubtype}, which will stop \textsf{biblatex-chicago}
italicizing your \textsf{title} in the absence of a
\textsf{booktitle}.)  If the medium in question needs specifying, the
\textsf{type} field is the place for it.  Please note, also, that
while the \textsf{titleaddon} field can still specify creative or
editorial functions for which \textsf{biblatex-chicago} provides no
automated, localized handling, you can also now provide the string you
need in an \textsf{editor[abc]type} field, e.g.,
\enquote{\texttt{libretto by}} (verdi:corsaro).

\mylittlespace For podcasts, newly covered by the 17th edition
(14.267), the \textsf{audio} type provides the nearest analogue I
could find, and in general most of the data should fit comfortably
into the fields already discussed above, the episode name in
\textsf{title} and the name of the podcast in \textsf{booktitle}, for
starters.  Two details, however, need mentioning: first, the
\textsf{note} field as the place to specify that it is a podcast, and
the \textsf{eventdate} field for the date of publication of the
specific episode (\textsf{title}) cited, which appears in close
association with that \textsf{title}.  Indeed, the \textsf{eventdate}
field helps \textsf{biblatex-chicago} know that the \textsf{audio}
entry is a podcast episode, and helps it construct the entry
appropriately (danforth:podcast).

\mybigspace This \mymarginpar{\textbf{book}} is the standard
\textsf{biblatex} and \textsc{Bib}\TeX\ entry type, and the package
can automatically provide abbreviated references in notes and
bibliography when you use a \textsf{crossref} or an \textsf{xref}
field.  The functionality is not enabled by default, but you can
enable it in the preamble or in the \textsf{options} field using the
\texttt{booklongxref} option.  Please see \textbf{crossref} in
section~\ref{sec:entryfields} and \texttt{booklongxref} in
section~\ref{sec:chicpreset}, below.  Cf.\ harley:ancient:cart,
harley\hc cartography, and harley:hoc for how this might look.  The
\textsf{book} type is useful also to present multimedia app content,
the added fields \textsf{version} and \textsf{type} providing
information about the app's version and about the system on which it
runs (14.268; angry:birds).

%%\enlargethispage{-\baselineskip}

\mybigspace This \mymarginpar{\textbf{bookinbook}} type provides the
means of referring to parts of books that are considered, in other
contexts, themselves to be books, rather than chapters, essays, or
articles.  Such an entry can have a \textsf{title} and a
\textsf{booktitle}, but it can also contain a \textsf{maintitle}, all
three of which will be italicized when printed.  In general usage it
is, therefore, rather like the traditional \textsf{inbook} type, only
with its \textsf{title} in italics rather than in quotation marks.  As
with the \textsf{book} type, you can automatically enable abbreviated
references in notes and bibliography, though this isn't the default.
Please see \textbf{crossref} in section~\ref{sec:entryfields} and
\texttt{booklongxref} in section~\ref{sec:chicpreset}, below.  (Cf.\
\emph{Manual} 14.109, 14.122, 14.124; bernhard:boris, bernhard:ritter,
and bernhard:themacher for the abbreviating functionality; also
euripides:orestes [treated differently in 14.122 and 14.124],
plato:republic:gr.)

\mybigspace This \mymarginpar{\textbf{booklet}} is the first of two
entry types --- the other being \textsf{manual}, on which see below
--- which are traditional in \textsc{Bib}\TeX\ styles, but which the
\emph{Manual} (14.220) suggests may well be treated basically as
books.  In the interests of backward compatibility,
\textsf{biblatex-chica\-go-notes} will so format such an entry, which
uses the \textsf{howpublished} field instead of a standard
\textsf{publisher}, though of course if you do decide just to use a
\textsf{book} entry then any information you might have given in a
\textsf{howpublished} field should instead go in \textsf{publisher}.
(See clark:mesopot.)

\mybigspace This \mymarginpar{\textbf{collection}} is the standard
\textsf{biblatex} entry type, but the package can automatically
provide abbreviated references in notes and bibliography when you use
a \textsf{crossref} or an \textsf{xref} field.  The functionality is
not enabled by default, but you can enable it in the preamble or in
the \textsf{options} field using the new \texttt{booklongxref} option.
Please see \textbf{crossref} in section~\ref{sec:entryfields} and
\texttt{booklongxref} in section~\ref{sec:chicpreset}, below.  See
harley:ancient:cart, harley:cartography, and harley:hoc for how this
might look.

\mybigspace This \mymarginpar{\textbf{customc}} entry type allows you
to include alphabetized cross-references to other, separate entries in
the bibliography, particularly to other names or pseudonyms, as
recommen\-ded by the \emph{Manual}.  (This is different from the
\textsf{crossref}, \textsf{xref}, \textsf{userf} and \textsf{related}
mechanisms, all primarily designed to include cross-references to
other works.  Cf.\ 14.81--82).  The lecarre:cornwell entry, for
example, would allow your readers to find the more-commonly-used
pseudonym \enquote{John Le Carr\'e} even if they were, for some reason,
looking under his real name \enquote{David John Moore Cornwell.}\ As I
read the specification, these cross-references are particularly
encouraged, bordering on required, when \enquote{a bibliography
  includes two or more works published by the same author but under
  different pseudonyms.}\ The following entries in
\textsf{notes-test.bib} show one way of addressing this:
creasey:ashe:blast, creasey:york:death, creasey:morton:hide,
ashe:creasey, york:creasey and morton:creasey.

\mylittlespace In these latter cases, you would need merely to place
the pseudonym in the \textsf{author} field, and the author's real
name, under which their works are presented in the bibliography, in
the \textsf{title} field.  To make sure the cross-reference also
appears in the bibliography, you can either manually include the entry
key in a \cmd{nocite} command, or you can put that entry key in the
\textbf{userc} field in the work's main .bib entry, in which case
\textsf{biblatex-chicago} will print the cross-reference if and only
if you cite the main work.  (Cf.\ \textsf{userc}, below.)

\mylittlespace Under ordinary circumstances, \textsf{biblatex-chicago}
will connect the two parts of the cross-reference with the word
\enquote{\emph{See}} --- or its equivalent in the document's language
--- in italics.  If you wish to present it differently, you can put
the connecting word(s) into the \textsf{nameaddon} field, formatted as
you wish.

\mybigspace This \mymarginpar{\textbf{dataset}} entry type, new in
\textsf{biblatex} 3.13, allows you to cite scientific databases, for
which the \emph{Manual} (14.257) presents some rather specific, if
brief, instructions.  To construct your entry, you can put the name of
the database into \textsf{author}, a \enquote{descriptive phrase or
  record locator} in the \textsf{title} field, and if there's a
specific accession number needed beyond the record locator you can put
it into the \textsf{number} field, with the \textsf{type} field
reserved to help explain what sort of \textsf{number} is involved.
The \textsf{howpublished} field can also be used to provide extra
descriptive detail about the \textsf{number}, if needed.  More
generally, a \textsf{url} will locate the database as a whole and a
\textsf{urldate} will specify the date you accessed it.  If, for some
reason, an additional date is relevant, then the \textsf{date} field
is available, while the \textsf{pubstate} field will appear before the
\textsf{date} in case you need to modify the latter.  (See 14.257;
genbank:db, nasa:db.)

\mybigspace This \mymarginpar{\textbf{image}} entry type is now a
clone of the \textsf{artwork} type, which see.  I retain it here for
historical reasons (See 3.22, 8.198; bedford:photo.)

\mybigspace These \mymarginpar{\textbf{inbook}\\\textbf{incollection}}
two standard \textsf{biblatex} types have very nearly identical
formatting requirements as far as the Chicago specification is
concerned, but I have retained both of them for compatibility.
\textsf{Biblatex.pdf} (\S~2.1.1) intends the first for \enquote{a part
  of a book which forms a self-contained unit with its own title,}
while the second would hold \enquote{a contribution to a collection
  which forms a self-contained unit with a distinct author and its own
  title.}  The \textsf{title} of both sorts will be placed within
quotation marks, and in general you can use either type for most
material falling into these categories.  I have, in both types,
implemented the \emph{Manual's} recommendations for space-saving
abbreviations in notes and bibliography when you cite multiple pieces
from the same \textsf{collection}.  These abbreviations are activated
by default when you use the \textsf{crossref} or \textsf{xref} field
in \textsf{incollection} entries and in \textsf{inbook} entries,
because although the \emph{Manual} (14.108) here specifies a
\enquote{multiauthor book,} I believe the distinction between the two
is fine enough to encourage similar treatments.  (For more on this
mechanism see \textbf{crossref} in section~\ref{sec:entryfields},
below, and the new option \texttt{longcrossref} in
section~\ref{sec:chicpreset}.  Please note that it is also active by
default in \textsf{letter} and \textsf{inproceedings} entries.)  If
the part of a book to which you are referring has had a separate
publishing history as a book in its own right, then you may wish to
use the \textsf{bookinbook} type, instead, on which see above.  (See
\emph{Manual} 14.106--9; \textsf{inbook}: ashbrook:brain,
phibbs:diary, will:cohere; \textsf{incollection}: centinel:letters,
contrib:contrib, sirosh:visualcortex; ellet:galena, keating:dearborn,
and lippincott:chicago [and the \textsf{collection} entry
prairie:state] demonstrate the use of the \textsf{crossref} field with
its attendant abbreviations in notes and bibliography.)

\mylittlespace \textbf{NB}: The \emph{Manual} suggests that, when
referring to a chapter, one use either a chapter number or the
inclusive page numbers, not both.  If, however, you wish to refer in a
footnote to a specific page within the chapter,
\textsf{biblatex-chicago-notes} will always print the optional,
postnote argument of a \cmd{cite} command --- the page number, say ---
instead of any inclusive page numbers given in the .bib file
\textsf{incollection} entry.  This mechanism is quite general, that
is, any specific page reference given in any sort of \cmd{cite}
command overrides the contents of a \textsf{pages} field in a .bib
file entry.

\mybigspace This \mymarginpar{\textbf{inproceedings}} entry type works
pretty much as in standard \textsf{biblatex}, even more so now that,
after a request from Patrick Danilevici, I have included the
\textsf{eventdate}, \textsf{eventtitle}, \textsf{eventtitleaddon}, and
\textsf{venue} fields for specifying where and when the event occurred
that produced the proceedings.  These four fields are the main
difference between it and \textsf{incollection}, along with the lack
of an \textsf{edition} field and the possibility that an
\textsf{organization} may be cited alongside the \textsf{publisher},
even though the \emph{Manual} doesn't specify the use of any of these
supplementary fields (14.217).  Please note, also, that the
\textsf{crossref} and \textsf{xref} mechanism for shortening citations
of multiple pieces from the same \textsf{proceedings} is operative
here, just as it is in \textsf{incollection} and \textsf{inbook}
entries.  See \textbf{crossref} in section~\ref{sec:entryfields} and
the option \texttt{longcrossref} in section~\ref{sec:chicpreset} for
more details.

\mybigspace This \mymarginpar{\textbf{inreference}} entry type is
aliased to \textsf{incollection} in the standard styles, but the
\emph{Manual} has particular requirements, so if you are citing
\enquote{[w]ell-known reference books, such as major dictionaries and
  encyclopedias,} then this type should simplify the task of
conforming to the specifications (14.232--34).  The main thing to keep
in mind is that I have designed this entry type for
\enquote{alphabetically arranged} works, which you shouldn't cite by
page, but rather by the name(s) of the article(s).  Because of the
formatting required by the \emph{Manual}, we need one of
\textsf{biblatex's} list fields for this purpose, and in order to keep
all this out of the way of the standard styles, I have chosen the
\textsf{lista} field.  You should present these article names just as
they appear in the work, separated by the keyword
\enquote{\texttt{and}} if there is more than one, and
\textsf{biblatex-chicago-notes} will provide the appropriate prefatory
string (\texttt{s.v.}, plural \texttt{s.vv.}), and enclose each in its
own set of quotation marks (ency:britannica).  In a typical
\textsf{inreference} entry very few other fields are needed, but
\enquote{if a physical edition is cited, not only the edition number
  (if not the first) but also the date the volume or set was issued
  must be specified.}  In practice, this means a \textsf{title},
\textsf{date}, and possibly an \textsf{edition} field.

\mylittlespace There are quite a few other peculiarities to explain
here.  First of all, you should present any well-known works
\emph{only} in notes, not in a bibliography, as your readers are
assumed to know where to go for such a reference.  You can use the
\texttt{skipbib} option to achieve this.  For such works, and given
how little information will be present even in a full note, you may
wish to use \cmd{fullcite} or \cmd{footfullcite} in place of the short
form, especially if, for example, you are citing different versions of
an article appearing in different editions.

\mylittlespace If the work is slightly less well known, it may be that
full publication details are appropriate (times:guide).  Also, you can
put an article name in the \textsf{postnote} field of
\textsf{inreference} entries and have it properly formatted for you,
and this holds for both long and short notes, which could allow you to
refer separately to many different articles from the same reference
work using only one .bib entry.  (In a long note, any
\textsf{postnote} field stops the printing of the contents of
\textsf{lista}.)  The only limitation on this system is that the
\textsf{postnote} field, unlike \textsf{lista}, is not a
\textsf{biblatex} list, and therefore for the formatting to work
correctly you can only put one article name in it.  Despite this
limitation, I hope that the system might simplify things for users who
cite numerous works of reference.

\mylittlespace If it seems appropriate to include such a work in the
bibliography, be aware that the contents of the \textsf{lista} field
will also be presented there, which may not be what you want.  A
separate \textsf{reference} entry might well solve this problem.  (The
sorting issues with \textsf{inreference}, \textsf{mvreference}, and
\textsf{reference} entries should no longer exist, as they should now
always sort by \textsf{title} rather than by any \textsf{author},
\textsf{editor}, or \textsf{namec} that might also be present.  If the
\textsf{title} starts with a definite or indefinite article then a
\textsf{sortkey} may still be necessary.)

\mylittlespace Speaking of the \textsf{author}, this field holds the
author of the specific entry (in \textsf{lista}), not the author of
the \textsf{title} as a whole.  This name will be printed after the
entry's name (grove:sibelius).  If you wish to refer to a reference
work by author or indeed by editor, having either appear at the head
of the note (long or short) or bibliography entry, then you'll need to
use a \textsf{book} entry instead (cf.\ schellinger:novel), where the
\textsf{lista} mechanism will also work in the bibliography, but which
in every other way will be treated as a normal book, often a good
choice for unfamiliar or non-standard reference works.

\mylittlespace Finally, all of these rules apply to online reference
works, along with a few more.  The 17th edition of the \emph{Manual}
now allows, \enquote{subject to editorial discretion,} the alternative
treatment of an online reference work which \enquote{does not have
  (and never had) a printed counterpart} (14.206, 14.233).  In effect
this means that it can be treated more like an \textsf{online} entry
than a \textsf{book}, its \textsf{title} therefore in plain roman
rather than in italics.  You can achieve this in \textsf{inreference}
entries by providing an \textsf{entrysubtype} in the entry.  Online
reference works need not only a \textsf{url} but also, always, a
\textsf{urldate} (instead of a \textsf{date)}, as these sources are in
constant flux.  When that flux is of a particularly high frequency, as
with Wikipedia, then a time stamp may also be needed.  You can provide
this in the \textsf{urldate} field itself, using the standard
\textsf{biblatex} format, e.g., \texttt{2008-07-01T10:18:00}.  It is
\mymarginpar{\texttt{urlstamp=true}} possible to turn off the printing
of the \textsf{urltime} with the \texttt{urlstamp} option, which is
set to \texttt{true} by default, but which can be changed in your
preamble for the whole document, for specific entry types, or in the
\textsf{options} field of specific entries (wikiped:bibtex,
grove:sibelius).

\mybigspace I
\mymarginpar{\textbf{jurisdiction}\\\textbf{legal}\\\textbf{legislation}}
document these three types in section~\ref{sec:legal} below, both
because they all follow the specifications of the \emph{Bluebook}
instead of the \emph{Manual}, and also because they are the only entry
types treated identically by the notes \&\ bibliography style and the
author-date styles.

\mybigspace This \mymarginpar{\textbf{letter}} is the entry type to
use for citing letters, memoranda, or similar texts, but \emph{only}
when they appear in a published collection.  (Unpublished material of
this nature needs a \textsf{misc} entry, for which see below.)
Depending on what sort of information you need to present in a
citation, you may simply be able to get away with a standard
\textsf{book} entry, which may then be cited by page number (see
14.104; meredith:letters, adorno:benj).  If, however, for whatever
reason, you need to give full details of a specific letter, then
you'll need to use the \textsf{letter} entry type, which attempts to
simplify for you the \emph{Manual}'s rather complicated rules for
formatting such references.  (See 14.111; jackson:paulina:letter,
white:ross:memo, white:russ [a completely fictitious entry to show the
\textsf{crossref} mechanism], white:total [a \textsf{book} entry, for
the bibliography]).

\mylittlespace To start, the name of the letter writer goes in the
\textsf{author} field, while the \textsf{title} field contains both
the name of the writer and that of the recipient, in the form
\texttt{Author to Recipient}.  The \textsf{titleaddon} field contains,
optionally, the type of correspondence involved.  If it's a letter,
the type needn't be given, but if it's a memorandum or report or the
like, then this is the place to specify that fact.  Also, because the
\textsf{origdate} field only accepts numbers, if you want to use the
abbreviation \enquote{n.d.}  (or \verb+\bibstring{nodate}+) for
undated letters, then this is where you should put it.  If you need to
specify where a letter was written, then you can also use this field,
and, if both are present, remember to separate the location from the
type with a comma, like so: \texttt{memorandum, London}.
Alternatively, you can put the place of writing into the
\textsf{origlocation} field.  Most importantly, the date of the letter
itself goes in the \textsf{origdate} field (\texttt{year-month-day}),
which allows a full date specification, while the publishing date of
the whole collection goes in the \textsf{date} field.  As in other
entry types, then, the \textsf{date} field has its ordinary meaning of
\enquote{date of publication.}  (You may have noticed that the
presentation of the \textsf{origdate} in this sort of reference uses
the day-month-year format, unlike the month-day-year format seen
elsewhere.  This follows a suggestion that material with \enquote{many
  references to specific dates} may use this format [6.38, 9.35,
14.224].  I should, I guess, make this configurable.)  Another
difficulty arises when producing the short footnote form, which
requires you to provide a \textsf{shorttitle} field of the form
\enquote{\texttt{to Recipient},} the latter name as short as possible
while avoiding ambiguity.  The \cmd{letterdatelong} command can be
used in the \textsf{postnote} field of the citation to print the
\textsf{origdate}, a possible aid to disambiguation.  The remaining
fields are fairly self explanatory, but do remember that the title of
the published collection belongs in \textsf{booktitle} rather than in
\textsf{title}.

\mylittlespace Finally, the \emph{Manual} specifies that if you cite
more than one letter from a given published collection, then the
bibliography should contain only a reference to said collection,
rather than to each individual letter, while the form of footnotes
would remain the same.  This should be possible using
\textsc{Bib}\TeX's and \textsf{biblatex's} standard \textsf{crossref}
field, with each \textsf{letter} entry pointing to a
\textsf{collection} or \textsf{book} entry, for example.  (If you are
using \textsf{Biber}, then \textsf{letter} entries correctly inherit
fields from \textsf{book} and \textsf{collection} entries, and also
from the \textsf{mvbook} and \textsf{mvcollection} types ---
\textsf{titles} from the former provide a \textsf{booktitle} and from
the latter a \textsf{maintitle}.)  I shall discuss cross references at
length later (see esp.\ \textbf{crossref} in
section~\ref{sec:entryfields}, below), but I should mention here that
\textsf{letter} is one of the entry types in which a \textsf{crossref}
or an \textsf{xref} field automatically results in special shortened
forms in notes and bibliography if more than one piece from a single
collection is cited.  (The other entry types are \textsf{inbook},
\textsf{incollection}, and \textsf{inproceedings}; see 14.108 for the
\emph{Manual}'s specification.)  This ordinarily won't be an issue for
\textsf{letter} entries in the bibliography, as individual letters
aren't included there, but it is operative in notes, where you can
disable it by setting the \texttt{longcrossref=true} option, on which
see section~\ref{sec:chicpreset}, below.  To stop individual letters
turning up in the bibliography, you can use the \texttt{skipbib}
option in the \textsf{options} field.

\mybigspace This \mymarginpar{\textbf{manual}} is the second of two
traditional \textsc{Bib}\TeX\ entry types that the \emph{Manual}
suggests formatting as books, the other being \textsf{booklet}. As
with this latter, I have retained it in
\textsf{biblatex-chicago-notes} for backward compatibility, its main
peculiarity being that, in the absence of a named author, the
\textsf{organization} producing the manual will be printed both as
author and as publisher.  If you are using \textsf{Biber} you no
longer need a \textsf{sortkey} field to aid \textsf{biblatex's}
alphabetization routines, as the style takes care of this for you
(cf.\ section~\ref{sec:formatopts}, below).  You also don't need to
provide a \textsf{shortauthor} field, as the style will automatically
use \textsf{organization} in the absence of anything else.  Of course,
if you were to use a \textsf{book} entry for such a reference, then
you would need to define both \textsf{author} and \textsf{publisher}
using the name you here might have put in \textsf{organization}.  (See
14.84; chicago:manual, dyna:browser, natrecoff:camera.  Cp.\ also the
\textbf{standard} entry type.)

\mybigspace As \mymarginpar{\textbf{misc}} its name suggests, the
\textsf{misc} entry type was designed as a hold-all for citations that
didn't quite fit into other categories.  In
\textsf{biblatex-chicago-notes}, I have somewhat extended its
applicability, while retaining its traditional use.  Put simply, with
no \textsf{entrysubtype} field, a \textsf{misc} entry will retain
backward compatibility with the standard styles, so the usual
\textsf{howpublished}, \textsf{version}, and \textsf{type} fields are
all available for specifying an otherwise unclassifiable text, and the
\textsf{title} will be italicized.  (The \emph{Manual}, you may wish
to note, doesn't give specific instructions on how such citations
should be formatted, so when using the Chicago style I would recommend
you have recourse to this traditional entry type as sparingly as
possible.)

\mylittlespace If you do provide an \textsf{entrysubtype} field, the
\textsf{misc} type provides a means for citing unpublished letters,
memoranda, private contracts, wills, interviews, and the like, making
it something of an unpublished analogue to the \textsf{letter},
\textsf{article}, and \textsf{review} entry types (which see).  It
also works well for presenting online audio pieces, particularly dated
ones, like speeches.  Typically, such an entry will cite part of an
archive, and equally typically the text cited won't have a specific
title, but only a generic one, whereas an \textsf{unpublished} entry
will ordinarily have a specific author and title, and won't come from
a named archive.  The \textsf{misc} type with an \textsf{entrysubtype}
defined is the least formatted of all those specified by the
\emph{Manual}, so titles are in plain text, and any location details
take no parentheses in full footnotes.  (It is quite possible, though
somewhat unusual, for archival material to have a specific title,
rather than a generic one.  In these cases, you will need to enclose
the title inside a \cmd{mkbibquote} command manually.  Cf.\
roosevelt:speech, shapey:partita.)

\mylittlespace If you are presenting part of an unpublished archive,
then constructing most of your \,.bib entry is fairly straightforward.
\enquote{Letter-like} \textsf{misc} entries follow many of the same
conventions as \textsf{letter} entries presenting published material.
Titles are of the form \texttt{Author to Recipient}, and further
information can be given in the \textsf{titleaddon} field, including
the abbreviation \enquote{\texttt{n.d.}}\ (or
\verb+\bibstring{nodate}+) for undated examples.  The place where a
letter was written can go in \textsf{origlocation}, while the
\textsf{note}, \textsf{organization}, \textsf{institution}, and
\textsf{location} fields (in ascending order of generality) allow the
specification of which manuscript collection now holds the letter,
though the \emph{Manual} specifies (14.227) that well-known
depositories don't usually need a city, state or country specified.
(The traditional \textsf{misc} fields are all still available, also.)
Both the long and short note forms can use the same \textsf{title},
but in both cases you may need to use the \cmd{headlesscite} command
to avoid the awkward repetition of the author's name, though that name
will always appear in the bibliography (creel:house).  If the
\textsf{misc} entry isn't a letter, remember that, as in
\textsf{article} and \textsf{review} entries, words like
\texttt{interview} or \texttt{memorandum} needn't be capitalized
unless they follow a period --- the automatic capitalization routines
(with the \textsf{title} field starting with a lowercase letter [see
dinkel:agassiz, spock:interview, and \texttt{\textbackslash autocap}])
will ensure correctness.  Also, please note that you can give
additional information about the \textsf{author} in such entries by
using the \textsf{nameaddon} field, providing your own square brackets
if you're indicating that the \textsf{author} is pseudonymous, or
parentheses if it's another sort of information of interest to your
readers.  The package options \texttt{nameaddonformat} and
\texttt{nameaddonsep} can help here, as well.  See
sections~\ref{sec:useropts} and \ref{sec:chicpreset}, below.

\mylittlespace Now for the subtleties.  First, the \emph{Manual}
(14.224) allows in these entries, as it does in documentation
generally \enquote{if numerous dates occur} (9.35), for a more
streamlined presentation of dates using the day-month-year form,
different from the standard American month-day-year.  In
\textsf{letter} entries you use the \textsf{origdate} field to give
the date of individual letters, and it is always presented in the more
streamlined form.  Here, the same field will do exactly the same
thing, though with the added wrinkle that if you'd prefer to use the
standard day-month-year form you can, simply by putting the date into
the \textsf{date} field instead.  (Please choose one only in
\textsf{misc} entries with an \textsf{entrysubtype} --- in
\textsf{letter} entries the \textsf{date} refers to the published
collection.) Again just as in \textsf{letter} entries, if you want to
include the day-month-year in a short note, I have provided the
\cmd{letterdatelong} command for inclusion in the postnote field of
the citation command.  (The standard \textsf{biblatex} command
\cmd{printdate} will work if you prefer the standard date form.)

\enlargethispage{\baselineskip}

\mylittlespace Second, some material (roosevelt:speech) includes a
venue for the event recorded in the archive, so I have added the
\textsf{venue} field, which prints \emph{before} the date, with the
\textsf{origlocation} appearing after it.  Somewhat confusingly, in
published letters the \textsf{origlocation} itself prints before the
date, rather than after, so if the inconsistency between published and
unpublished letters bothers you then you could conceivably use
\textsf{venue} instead of \textsf{origlocation} for that purpose here.

\mylittlespace Finally, a few further notes.  First, please be aware
that defining an \textsf{entrysubtype} activates the automatic
capitalization mechanism in the \textsf{title} field of \textsf{misc}
entries, on which see\,\texttt{\textbackslash autocap} in
section~\ref{sec:formatcommands} below.  Second, and again as with
\textsf{letter} entries, the \emph{Manual} (14.222) suggests that
bibliography entries contain only the name of the manuscript
collection, unless only one item from that collection is cited.  The
\textsf{crossref} field can be used, as well as the \texttt{skipbib}
option, for preventing the individual items from turning up in the
bibliography.  Obviously, this is a matter for your discretion, and if
you're using only short notes (see the \texttt{short} option,
section~\ref{sec:useropts} below), you may feel the need to include
more information in the note if the bibliography doesn't contain a
full reference to an individual item.  Third, the \emph{Manual} offers
several examples of specific location information for pieces from an
archive, some of which appears \emph{before} the main archive name,
and some of which appears \emph{after} it.  I assume this may depend
on the exact nature of the archive itself, but in any case you can try
the \textsf{type} or \textsf{howpublished} fields for the first case
and the \textsf{number} field for the second.  Last, in all this class
of archived material, the \emph{Manual} (14.221) quite specifically
requires more consistency within your own work than conformity to some
external standard, so it is the former which you should pursue.  I
hope that \textsf{biblatex-chicago-notes} proves helpful in this
regard.  (See 14.211, 14.219, 14.221-231, 14.256, 14.264; creel:house,
dinkel:agassiz, roosevelt:speech, shapey:partita, spock:interview.)

\mybigspace This \mymarginpar{\textbf{music}} is one of three
audiovisual entry types, and is intended primarily to aid in the
presentation of musical recordings that do not have a video component,
though it can also include audio books (auden:reading).  A DVD or VHS
of an opera or other performance, by contrast, should use the
\textbf{video} type instead, while an online music video will probably
need an \textbf{online} entry.  (Cf.\ \textsf{online} and
\textsf{video}; handel:messiah, horowitz:youtube.)  Because
\textsf{biblatex} --- and \textsc{Bib}\TeX\ before it --- were
designed primarily for citing book-like objects, some choices needed
to be made in assigning the various roles found on the back of a CD to
the fields in a typical .bib entry.  I have also implemented several
bibstrings to help in identifying these roles within entries.  The
17th edition of the \emph{Manual} once again revised its
recommendations for this type, but fortunately the changes are
additive, i.e., you can re-use 16th-edition citations but are
encouraged to peruse the following guidelines to see if there's any
information you might think of adding to bring your citations more
into line with the spec.

\mylittlespace These guidelines, in summary form, are:

{\renewcommand{\descriptionlabel}[1]{\qquad\textsf{#1}}
\begin{description}
\item[author:] composer, songwriter, or performer(s), they will be
  closely associated with the \textsf{title}, either before it at the
  head of the entry or, at your discretion, just after it
  (holiday:fool).
\item[bookauthor:] Somewhat like an \textsf{author}, but it will hold
  the name associated with a whole album rather than an individual
  piece, should both be present, and will therefore appear in close
  association with the \textsf{booktitle}, rather than the
  \textsf{title} (rihanna:umbrella).
\item[editor, editora, editorb:] conductor, director or
  performer(s).  These will ordinarily follow the \textsf{title} of
  the work, though the usual \texttt{useauthor} and \texttt{useeditor}
  options can alter the presentation within an entry.  Because these
  are non-standard roles, you will need to identify them using the
  following:
\item[editortype, editoratype, editorbtype:] The most common roles,
  all associated with specific bibstrings (or their absence), will be
  \texttt{conductor}, \texttt{director}, \texttt{producer}, and,
  oddly, \texttt{none}.  The last is particularly useful when
  identifying the group performing a piece, as it usually doesn't need
  further specifying and this role prevents \textsf{biblatex} from
  falling back on the default \texttt{editor} bibstring.  The 17th
  edition (14.263) also seems to favor, in some circumstances, using
  strings to identify individual performers, e.g., \enquote{vocalist}
  or \enquote{pianist,} so even though there's no \cmd{bibstring}
  associated with these types you can now provide them, or anything
  else you need, in whatever form (\enquote{vocalist} or \enquote{sung
    by}) suits your citation.
\item[note:] This field can also hold contributors, perhaps
  collaborators or featured artists (holiday:fool, rihanna:umbrella).
\item[title, booktitle, maintitle:] As with the other audiovisual
  types, \textsf{music} serves as an analogue both to books and to
  collections, so the title will either be, e.g., the album title or a
  song title, in which latter case the album title would go into
  \textsf{booktitle}.  If you wish to cite a song that, as may happen,
  isn't part of any larger collection, your entry will in such a case
  have only a \textsf{title}, which \textsf{biblatex-chicago} would
  normally interpret as an album title.  You can now define an
  \textsf{entrysubtype} to let it know that the lone \textsf{title} is
  in fact a song (cf.\ naraya).  The \textsf{maintitle} might be
  necessary for something like a box set of \emph{Complete
    Symphonies}.
\item[chapter:] The 17th edition seems more keen on having track
  numbers for individual pieces, whether on a traditional format or on
  a streaming service.  The \textsf{chapter} field is the place for
  this information, and \textsf{biblatex-chicago} will automatically
  prepend the localized string \texttt{track} (cf.\ holiday:fool,
  rihanna:umbrella).
\item[publisher, series, number:] These three closely-associated
  fields are intended for presenting the catalog information provided
  by the music publisher.  The 17th edition generally only requires
  the \textsf{series} and \textsf{number} fields (nytrumpet:art),
  which hold the record label and catalog number, respectively.
  Alternatively, \textsf{publisher} would function as a synonym for
  \textsf{series} (holiday:fool), but there may be cases when you need
  or want to specify a publisher in addition to a label, as perhaps
  when a single publisher oversees more than one label.  You can
  certainly put all of this information into one of the above fields,
  but separating it may help make the .bib entry more readable.
\item[pubstate:] The \textsf{pubstate} field in \textsf{music} entries
  mainly has the usual meaning it has in other entry types, for which
  see the documentation of the field in section~\ref{sec:entryfields},
  below.  If the field contains \texttt{reprint}, however, this has a
  special meaning in \textsf{music} entries, where it will transform
  the \textsf{origdate} from a recording date for an entire album into
  an original release date for that album, notice of which will be
  printed towards the end of a note or bibliography entry.  No
  \texttt{reprint} \cmd{bibstring} will be printed, as only the syntax
  of the reference will have been altered.
\item[date, eventdate, origdate:] The 17th edition of the
  \emph{Manual}, like the 16th, considers \textsf{music} citations
  without a date to be \enquote{generally unacceptable} (14.263).
  Finding a date may take some research, but they will basically fall
  into two types, i.e., the date(s) of the recording or the copyright
  / publishing date(s).  Recording dates go either in
  \textsf{origdate} (for complete albums) or \textsf{eventdate} (for
  individual tracks).  The copyright or publishing dates go either in
  the \textsf{date} field (which applies to the current medium you are
  citing) or in the \textsf{origdate} field (which refers to the
  original release date).  You may have noticed that the
  \textsf{origdate} has two slightly different uses --- you can tell
  \textsf{biblatex-chicago} which sort you intend by using the string
  \texttt{reprint} in the \textsf{pubstate} field, which transforms
  the \textsf{origdate} from a recording date into an original release
  date.  The style will automatically prepend the bibstring
  \texttt{recorded} to the \textsf{eventdate} or, in the absence of
  this \textsf{pubstate} mechanism, to the \textsf{origdate}, or even
  to both, but you can modify what is printed there using the
  \textsf{userd} field, which acts as a sort of date type modifier.
  In \textsf{music} entries, \textsf{userd} will be prepended to an
  \textsf{eventdate} if there is one, barring that to the
  \textsf{origdate}, barring that to a \textsf{urldate}, and absent
  those three to the \textsf{date}.  (See floyd:atom, holiday:fool,
  nytrumpet:art.)
\item[type, howpublished:] As in all the audiovisual entry types, the
  \textsf{type} field holds the medium of the recording, e.g., vinyl,
  33 rpm, 8-track tape, cassette, compact disc, mp3, ogg vorbis.  The
  \textsf{howpublished} field, newly included for the 17th edition,
  can hold similar information \enquote{for streaming audio formats
    and downloads} (14.263). It can also, alternatively, hold the name
  of the streaming service, e.g., Spotify (cf.\ rihanna:umbrella).
\end{description}}

The entries in \textsf{notes-test.bib} should at least give you a good
idea of how this all works, and that file also contains an example of
an audio book presented in a \textsf{music} entry.  If you browse the
examples in the \emph{Manual} you will see the sheer variety of
possibilities for presenting these sources, my intention being that
judicious manipulation of\ .bib entries should allow you to make
\textsf{biblatex-chicago} do what you want.  Please let me know if
I've ignored something you need.  (Cf. 14.263--64; \textsf{eventdate},
\textsf{origdate}, \textsf{userd}; auden:reading, beethoven:sonata29,
bernstein:shostakovich, floyd:atom, holiday:fool, nytrumpet:art,
rubinstein:chopin.)

\mybigspace The \mymarginpar{\textbf{mvbook}\\\textbf{mvcollection}%
  \\\textbf{mvproceedings}\\\textbf{mvreference}} 17th edition of the
\emph{Manual} has deployed, in at least two contexts, a notable
syntactic change in the presentation of works that form part of other,
larger works.  Generally, the order of presentation, in
\textsf{biblatex} terms, has always been \textsf{title} --
\textsf{booktitle} -- \textsf{maintitle}, in increasing order of
generality.  In the vast majority of cases this order still holds, but
in TV episodes, for one example, the recommendation now is to present
the name of the series (\textsf{booktitle}) \emph{before} the name of
the episode (\textsf{title}).  The \textbf{video} type (14.265)
provides this by using an \textsf{entrysubtype}, \texttt{tvepisode},
which reverses the order for you in both long and short notes, and in
the bibliography.  The other context in which this reversal occurs is
multi-volume works (14.116--22).  Here, the preferred format, at least
for notes, appears to be \textsf{maintitle} -- \textsf{[book]title}
or, when all three titles are present, \textsf{title} --
\textsf{maintitle} -- \textsf{booktitle}.  The \emph{Manual} doesn't
carry this reordering through with absolute consistency, but I think
it important at least to offer it as a possibility to users of
\textsf{biblatex-chicago}.  Reluctant as I am simply to change the
data model and provide non-standard entry types, the least invasive
method seemed, and seems, to me to be to modify the \textbf{mv*} entry
types while maintaining backward compatibility with users' current
deployments of these types.

\mylittlespace So, while these types are no longer aliased to the
entry type that results from removing the \enquote{mv} from their
names, any\ .bib entries using them should, without modification,
continue to function as they always have.  Should you wish to ignore
the new syntax of presentation, and there are still examples in the
\emph{Manual} which do just that, then no changes are necessary.
These entries will still function, assuming you are using
\textsf{Biber}, as the target of cross-references from other entries,
the \textsf{title} of the \textbf{mv*} entry \emph{always} providing a
\textsf{maintitle} for the entry referencing it.  (If you want to
provide a \textsf{booktitle} for the referencing entry, please use
another entry type, e.g., \textbf{collection} for
\textbf{incollection} or \textbf{book} for \textbf{inbook}.  These
distinctions are particularly important to the correct functioning of
the abbreviated references that \textsf{biblatex-chicago}, in various
circumstances, provides.  Please see the documentation of the
\textbf{crossref} field in section~\ref{sec:entryfields}, below.)

\mylittlespace Also unchanged is the requirement, when multi-volume
works are presented in the reference apparatus, that any dates should
be appropriate to the specific nature of the citation.  This means
that a date range that is right for the presentation of a multi-volume
work in its entirety isn't right for citing, e.g., a single volume of
that work which appeared in one of the years contained in the date
range.  Because child entries will by default inherit all the date
fields from their parent (including the \textsf{endyear} of a date
range), I have turned off the inheritance of \textsf{date} and
\textsf{origdate} fields from all of the \textbf{mv*} entry types to
any other entry type.  When the dates of the parent and of the child
in such a situation are exactly the same, then this unfortunately
requires an extra field in the child's .bib entry.  When they're not
the same, as will, I believe, often be the case, this arrangement
saves a lot of annoying work in the child entry to suppress
wrongly-inherited fields.  Other sorts of parent entries aren't
affected by this, and of course you must be using \textsf{Biber} for
the settings to apply.

% %\enlargethispage{\baselineskip}

\mylittlespace Should you wish to employ the new,
\textsf{maintitle}-first syntax, then you'll
\mymarginpar{\texttt{maintitle}} need to use the \texttt{maintitle}
\textsf{relatedtype}.  In its simplest usage, to document one volume
of a multi-volume set, you would have, e.g., an \textsf{mvcollection}
entry with \textsf{relatedtype} \texttt{maintitle}, and a
\textsf{related} field pointing to a \textsf{collection} entry.  When
you cite the \textsf{mvcollection} entry itself, you'll get a long
note like \emph{MVCollTitle}, vol.\ 1, \emph{CollTitle}, and a short
note like \emph{MVCollTitle}, vol.\ 1., or, with a \textsf{postnote}
field, \emph{MVCollTitle}, 1:12, as the specification requires.  If
you wanted to cite one essay in the \textsf{collection}, then you
would, additionally, need an \textsf{incollection} entry with the
\texttt{maintitle} \textsf{relatedtype} and a \textsf{related} field
pointing to the \textsf{mvcollection} entry already mentioned, so
you're creating a chain of three different related entries but
presenting them in one reference.  It's important to keep in mind here
that, in effect, you're \emph{not} actually citing the
\textsf{mvcollection} entry, but the one volume of it represented by
the \textsf{collection} entry, or indeed an essay in that one volume.
Please consult the \emph{Manual} (14.116--22), and also see
harley:ancient:cart, harley:cartography, and harley:hoc for the
\enquote{old style} presentation with abbreviated cross references
using the \textsf{crossref} field, harleymt:ancient:cart,
harleymt:cartography and harleymt:hoc for the new presentation using
the \texttt{maintitle} \textsf{relatedtype}, and also plato:timaeus:gr
for an example of a three-work \texttt{maintitle} chain starting with
a \textsf{bookinbook} entry.

\mylittlespace The documentation of the \texttt{maintitle}
\textsf{relatedtype} in section~\ref{sec:related} contains all the
details, but there are several things I should like to mention here.
First, while you can happily mix these two methods of presentation in
your documents, please don't mix them within individual entries, which
means that if you are using a \textsf{crossref} field to an
\textsf{mvcollection} entry in a \textsf{collection} entry, say, and
the \textsf{collection} entry is itself the target of the
\textsf{mvcollection} entry's \textsf{related} field, please be
careful not to cite that \textsf{collection} entry independently, as
it can lead to unexpected results.  (If things don't look right to
you, try eliminating the use of \textsf{crossref} entirely from these
\textsf{related} chains and see if that helps, then send me a bug
report if it does.)  This restriction also means that, although the
\emph{Manual} prefers the \textsf{maintitle}-first format in notes and
allows either syntax in the bibliography, nonetheless with
\textsf{biblatex-chicago} whichever syntax you choose for the notes
will also appear in the bibliography.  Second, if you want to use a
three-work chain to cite one part of one volume, then this is possible
only by using the following entry types: \textsf{bookinbook},
\textsf{inbook}, \textsf{incollection}, \textsf{inproceedings}, and
\textsf{letter}.  All two-work chains must start with one of the
\textsf{mv*} types.  Third, as might be apparent from the previous
list, \textsf{mvreference} entries are special, in that their
\textsf{related} field should point to an \textsf{inreference} entry
if you want to cite an entry in an \enquote{alphabetically arranged
  work}, or to a \textsf{reference} entry otherwise.

\mylittlespace Fourth, please remember that, as these are citations
not of an \textsf{mv*} entry but rather of that entry's
\textsf{related} field, any \textsf{volumes} field in the former won't
be printed by default.  You can change this by setting the
\texttt{hidevolumes} option to \texttt{false} either in the preamble
or in the \textsf{options} field of the entry referenced by the
\textsf{mv*} entry's \textsf{related} field.  Finally, if you look at
the \enquote{mt} variants of the harley* entries mentioned above,
you'll see that harleymt:hoc has both subsidiary volumes included in
its \textsf{related} field.  You can create a separate \textsf{mv*}
entry pointing to each of it's subsidiary volumes, or you can list all
of those volumes in one \textsf{mv*} entry's \textsf{related} field
and \textsf{biblatex-chicago} will create separate clones for each
volume listed, clones with a standardized entry key looking like
\enquote{\texttt{mventrykey-singlevolumeentrykey},} which you should
then use for your citations.  The original \texttt{mventrykey}, in
this case, refers merely to the original \textsf{mv*} entry, as though
it had never had a \textsf{related} field, though you do need to cite
(or \cmd{nocite}) it somewhere in your document to make the
single-volume clones available in your reference apparatus.  The
mechanism's designed to save you some typing in common scenarios;
please see all of the (multifarious) details in
section~\ref{sec:related}, below.

\enlargethispage{\baselineskip}

\paragraph*{\protect\mymarginpar{\textbf{online}}}
\label{sec:online}

One of the features of the 17th edition of the \emph{Manual} is the
considerably extended, but still scattered, treatment of online
materials (8.189--92, 14.6--18, 14.159--63, 14.175--76, 14.187,
14.189, 14.205--10, 14.233).  The principles of that treatment have
changed somewhat, as the \emph{Manual} now places greater emphasis on
the \emph{location} of a source, which can in many cases outweigh, as
far as choosing an entry type goes, the \emph{nature} of the source.
Working out the correspondences between online sources and
\textsf{biblatex-chicago} entry types can, therefore, be tricky, so I
have included table~\ref{tab:online:types} summarizing the
increasingly detailed instructions in the \emph{Manual}, along with
some further annotations here that might help to clarify it.

\afterpage{\clearpage\thispagestyle{longtable}

\begin{table}[h!]
  \caption[\hspace{-1em}Online Entry Types - Notes \&\ Bibliography]%
  {Online materials and notes \&\ bibliography entry types}
  \label{tab:online:types}
  \centering\small\sffamily
  \hspace*{-6em}
  \begin{tabularx}{160mm}{@{}>{\raggedright}p{25mm}>{\raggedright}p{20mm}p{15mm}p{26mm}X@{}}
    \toprule
    Online Material & Entry Type & CMS Ref. & Sample Entry &
    Notes \\
    \cmidrule{1-1}\cmidrule(l){2-2}\cmidrule(l){3-3}\cmidrule(l){4-4}
    \cmidrule(l){5-5}
    Online edition of trad.\ publ.\ matter. &&&& Use the same
    entry type as you would choose were you citing it
    from a printed source.\\\addlinespace[.8mm]
    & @Book & 14.161-62 & james:ambassadors &
    CMS prefers (scanned) original page numbers
    to reflowable text. \\\addlinespace[.8mm]
    & @Article @Review & 14.175 & black:infectious & If
    no \enquote{suitable URL} is available, e.g., if it points
    to a generic portal page rather than to an abstract,
    use the name of the commercial database in an \textsf{addendum}
    field instead.
    \\\addlinespace[.8mm]
    \cmidrule(r{2em}){1-1}\addlinespace[.8mm]
    Blogs&& 14.208 &&\\\addlinespace[1.5mm]
    \hspace{.5em} Entire & @Periodical & & amlen:wordplay & The
    \textsf{maintitle} field holds the larger publication of which the
    blog is a part. \\
    \hspace{.5em} Single post & @Article & & amlen:hoot &
    \\\addlinespace[.8mm]
    \hspace{.5em} Comment & @Review & & viv:amlen & You can also
    use the \texttt{commenton} \textsf{relatedtype} for
    this. \\\addlinespace[.8mm]
    \cmidrule(r{2em}){1-1}\addlinespace[.8mm]
    Social Media & @Online &&& This includes anything --- posts, photos,
    videos --- on these and similar sites.  In other words, the
    \emph{location} of the material defines its treatment. \\\addlinespace[2mm]
    \hspace{.5em} Mailing list or \hspace*{.5em} forum post & & 14.210 &
    powell:email & Posts on private lists are to be treated as
    \enquote{personal communications,} using @Misc w/
    \textsf{entrysubtype}.\\\addlinespace[.8mm]
    \hspace{.5em} Facebook&& 14.209 & diaz:surprise &\\\addlinespace[.8mm]
    \hspace{.5em} Twitter&&& obrien:recycle &\\\addlinespace[.8mm]
    \hspace{.5em} Instagram&&& souza:obama &\\\addlinespace[.8mm]
    \hspace{.5em} Comments / \hspace*{.5em} replies & & 14.210
    & braun:reply &
    The \texttt{commenton} \textsf{relatedtype} is \emph{required}
    for this, and for the next entry, too.\\\addlinespace[.8mm]
    && 14.209 & licis:diazcomment &\\\addlinespace[.8mm]
    \cmidrule(r{2em}){1-1}\addlinespace[.8mm]
    Online Multimedia &&&&\\\addlinespace[.8mm]
    \hspace{.5em} Online video & @Online & 14.267 & pollan:plant &
    This category includes TED talks and most informal videos on
    YouTube and similar sites. \\\addlinespace[.8mm]
    \hspace{.5em} Online video, \hspace*{.5em} from a trad.\ \hspace*{.5em}
    journal & @Article &&
    kessler:nyt & You can use @Online, but this requires special
    formatting in the \textsf{note} or \textsf{titleaddon} field.
    \\\addlinespace[.8mm]
    \hspace{.5em} Published films in \hspace*{.5em} an archive & @Video &&
    weed:flatiron &\\\addlinespace[.8mm]
    \hspace{.5em} Podcasts & @Audio && danforth:podcast & Note the
    eventdate of the individual episode.\\\addlinespace[.8mm]
    \hspace{.5em} Archival audio & @Misc w/ \textsf{entrysubtype} & 14.264 &
    roosevelt:speech & Can have both a venue and an origlocation.
    \\\addlinespace[.8mm]
    \cmidrule(r{2em}){1-1}\addlinespace[.8mm]
    Streaming Media&&&& \\\addlinespace[.8mm]
    \hspace*{.5em} TV / Film & @Video & 14.265 & mayberry:brady &
    The streaming service is supplied by the URL.  The
    \texttt{tvepisode} entrysubtype is new in the 17th
    edition. \\\addlinespace[.8mm]
    \hspace*{.5em} Music & @Music & 14.263 & rihanna:umbrella &
    The streaming service is supplied by the howpublished field.
    \\\addlinespace[.8mm]
    \hspace*{.5em} News / Interviews & @Article @Review & 14.213 &
    bundy:macneil & Network information goes in the usera field.
    \\\addlinespace[.8mm]
    \cmidrule(r{2em}){1-1}\addlinespace[.8mm]
    Websites & @Online & 14.206-7 & evanston:library stenger:privacy &
    An online source \enquote{analogous to a traditionally printed
    work but [which] does not have (and never had) a printed counterpart}
    may now use an @Online entry, at your discretion.
    \\\addlinespace[.8mm]
    \hspace*{.5em} Reference works, \hspace*{.5em} cited by alpha-
    \hspace*{.5em} betized entry & @InReference w/ entrysub-\par
    type & 14.233 & wikiped:bibtex & As above, you can choose the
    @Online treatment of the title, but it's best achieved
    using an @InReference entry w/ entrysubtype. \\\addlinespace[.8mm]
    \hspace*{.5em} Scientific data- \hspace*{.5em} bases & @Dataset &
    14.257 & genbank:db & New in this release.\\\addlinespace[.8mm]
    \bottomrule
  \end{tabularx}
\end{table}}

\mylittlespace The basic principle, as I've cited in the penultimate
entry of table~\ref{tab:online:types}, is that \enquote{the title of a
  website that is analogous to a traditionally printed work but does
  not have (and never had) a printed counterpart can be treated like
  the titles of other websites, subject to editorial discretion}
(14.206).  This means that an intrinsically online entry like
stenger:privacy (citing CNN.com) need no longer be an \textsf{article}
but can be presented in an \textsf{online} entry.  (The same principle
applies to wikiped:bibtex, but because of the code facilitating
presentation of alphabetized entries in reference works, it's best in
this case to keep the \textsf{inreference} entry but add an
\textsf{entrysubtype} so that the \textsf{title} is presented as it
would be in an \textsf{online} entry.)  The corollary of the
principle, as the first entry in table~\ref{tab:online:types}
suggests, is that an online edition of a printed work will generally
require the same entry type as that printed work itself would.  Blogs
are, therefore, somewhat anomalous in requiring the various periodical
types, though the \emph{Manual} does specify that if you're not sure
whether a website is a blog, then it probably requires the
\textsf{online} type (14.206).  Social media, on the other hand, are
very much subject to the first principle, requiring \textsf{online}
entries no matter whether the citation is of text, a photo, or a
video.  Without pretending that all of the correspondences flow
deductively from the basic principles, I hope that the table might
simplify most of your choices.  If something remains unclear, please
let me know and I'll see if I can improve it.

%\enlargethispage{\baselineskip}

\mylittlespace A few more notes are in order. I designed the
\textsf{relatedtype} \texttt{commenton} to facilitate citation of
online comments, though it works slightly differently in the two entry
types in which it is available, \textsf{online} and \textsf{review}.
In both types it allows you to mimic thread structure by citing a
chain of replies to comments on posts, etc., all in a single entry,
while also simplifying your\ .bib entries.  This simplification works
differently depending on whether the comment itself has no specific
title, as always in \textsf{review} entries, or does have such a
title, as especially in \textsf{online} social media entries.  In the
former case, the \textsf{related} apparatus allows you not to provide
a \textsf{title} at all, but in the latter you still need a
\textsf{title}, which will be followed by the \textsf{relatedstring}.
In these latter entries, the \emph{only way to cite such comments} is
by using the \texttt{commenton} \textsf{relatedtype}
(licis:diazcomment).  If, in \textsf{online} entries, you decided
\emph{not} to use \textsf{commenton} in an entry like braun:reply, and
simply use a specially-crafted \textsf{titleaddon} field, you lose the
possibility of having two dates in the entry, one for the comment and
one for the original post, though to be fair it does end up looking
like the example in 14.210, where it is ambiguous to which part of the
citation the date applies.

\mylittlespace As for the thread structure, I've not tested how far
down the rabbit hole you can go, but a series of entries linked one to
the next by this \textsf{relatedtype} will all turn up if you cite the
first in the chain, though of course you can use the technique merely
as a convenient way to structure and simplify your\ .bib file, without
creating chains longer than 2 entries.  The default connecting string
is the localized \texttt{commenton}, but you can use
\textsf{relatedstring} to change it to \enquote{\texttt{reply to}} or
whatever else you need.  Please see the documentation of this
\textsf{relatedtype} in section~\ref{sec:related}, and also
diaz:surprise and licis:diazcomment.

\mylittlespace In general, constructing an \textsf{online}\ .bib file
entry is much the same as in \textsf{biblatex}.  The \textsf{title}
field would contain the title of the page, the \textsf{organization}
field could hold the title or owner of the whole site.  If there is no
specific title for a page, but only a generic one, then such a title
should go in \textsf{titleaddon}, not forgetting to begin that field
with a lowercase letter so that capitalization will work out
correctly.  It is worth remarking here, too, that the \emph{Manual}
(14.12--13) prefers, if they're available, revision dates to access
dates when documenting online material.  Indeed, given how rapidly
online sources may change (14.191, 14.209, 14.233), a time stamp may
often be necessary further to specify a revision date
(\textsf{urldate}) or the date of a comment or reply (\textsf{date}).
This time specification should be added to the date field using
\textsf{biblatex's} standard format, i.e.,
\texttt{2008-07-01T10:18:00}.  If a time zone is needed, then a
separate \textsf{timezone} or \textsf{urltimezone} field is the best
way, as it allows you to provide the initialisms that the
\emph{Manual} prefers (10.41, 14.191).  See \textsf{date},
\textsf{timezone}, \textsf{urldate}, and \textsf{userd} in
section~\ref{sec:entryfields}, below.

% %\enlargethispage{\baselineskip}

\mybigspace The \mymarginpar{\textbf{patent}} \emph{Manual} is very
brief on this subject (14.258), but very clear about which information
it wants you to present, so such entries may not work well with other
\textsf{biblatex} styles.  The important date, as far as Chicago is
concerned, is the filing date.  If a patent has been filed but not yet
granted, then you can place the filing date in either the
\textsf{date} field or the \textsf{origdate} field, and
\textsf{biblatex-chicago-notes} will automatically prepend the
bibstring \texttt{patentfiled} to it.  If the patent has been granted,
then you put the filing date in the \textsf{origdate} field, and you
put the date it was issued in the \textsf{date} field, to which the
bibstring \texttt{patentissued} will automatically be prepended.  You
can place additional information in the \textsf{addendum} field if
desired, and it will be printed in close association with the dates.
The patent number goes in the \textsf{number} field, and you should
use the standard \textsf{biblatex} bibstrings in the \textsf{type}
field.  Though it isn't mentioned by the \emph{Manual},
\textsf{biblatex-chicago-notes} will print the \textsf{holder} after
the \textsf{author}, if you provide one.  Finally, the style
automatically capitalizes \textsf{patent} titles sentence-style, so if
you need to keep a word capitalized then you should wrap it in curly
braces.  See petroff:impurity.

\mybigspace The \mymarginpar{\textbf{performance}} 17th edition of the
\emph{Manual} includes a new section (14.266) on citing live
performances, and even though such references can usually be limited
to the main text it may sometimes be useful to include them in notes.
Since \textsf{biblatex} provides the \textbf{performance} type, albeit
without using it in its standard styles, I though it might be useful
to define it for \textsf{biblatex-chicago}, particularly as the other
option for such material is the \textsf{misc} entry without any
\textsf{entrysubtype}, and that entry type is already somewhat
overloaded, though you can still use it if you wish.

\mylittlespace Such entries will generally have a \textsf{title}, a
\textsf{venue}, a \textsf{location} for the venue, and a \textsf{date}
for the performance, along with a possible plethora of authorial
and/or editorial roles depending on which sorts of contributor(s) you
wish to emphasize in the citation.  The \textsf{editor[abc]} and
\textsf{editor[abc]type} fields should be most helpful here.  I have
included strings for \texttt{choreographer} in all localization files,
but for others you may need to provide them in the
\textsf{editor[abc]type} fields as you wish them printed ---
\textsf{biblatex-chicago} will automatically capitalize any that start
with a lowercase letter.

\mybigspace This \mymarginpar{\textbf{periodical}} is the standard
\textsf{biblatex} entry type for presenting an entire issue of a
periodical, rather than one article within it.  It has the same
function in \textsf{biblatex-chicago-notes}, and in the main uses the
same fields, though in keeping with the system established in the
\textsf{article} entry type (which see) you'll need to provide
\textsf{entrysubtype} \texttt{magazine} if the periodical you are
citing is a \enquote{newspaper} or \enquote{magazine} instead of a
\enquote{journal.}  Also, remember that the \textsf{note} field is the
place for identifying strings like \enquote{special issue,} with its
initial lowercase letter to activate the automatic capitalization
routines.  (See \emph{Manual} 14.178; good:wholeissue.)

\mylittlespace It is worth noting a few things.  First, the
\textsf{titleaddon} field is now available in these entries, but as
the \textsf{title} here is analogous to the \textsf{journaltitle} in
\textsf{article} or \textsf{review} entries, the new
\mycolor{\texttt{jtitleaddon}} option (section~\ref{sec:chicpreset})
governs the punctuation separating the \textsf{titleaddon} from the
\textsf{title}.  Second, the special \textsf{biblatex} field
\textsf{shortjournal} allows you to present shortened
\textsf{journaltitles} in \textsf{article}, \textsf{review}, and
\textsf{periodical} entries, as well as facilitating the creation of
lists of journal abbreviations in the manner of a \textsf{shorthand}
list.  Because the \textsf{periodical} type uses the \textsf{title}
field instead of \textsf{journaltitle}, \textsf{biblatex-chicago}
automatically copies any \textsf{shorttitle} field, if one is present,
into \textsf{shortjournal}.  Please see the documentation of
\textbf{shortjournal} in section~\ref{sec:entryfields} for all the
details on how this works.  Finally, the \textsf{periodical} type is
the place for citing whole blogs, rather than individual blog posts,
which require either an \textsf{article} or a \textsf{review} entry.
In such citations the 17th edition (14.208) recommends that you
include the name of any larger (usually periodical) publication of
which the blog is a part.  The \textsf{maintitle} field (with
\textsf{mainsubtitle} and \textsf{maintitleaddon}, if needed) is the
place for it. Cf.\ amlen:wordplay.

\mybigspace This \mymarginpar{\textbf{proceedings}} is the standard
\textsf{biblatex} and \textsc{Bib}\TeX\ entry type, now also including
the \textsf{eventdate}, \textsf{eventtitle}, \textsf{eventtitleaddon},
and \textsf{venue} fields for identifying the event that produced the
\textsf{proceedings}.  The package can automatically provide
abbreviated references in notes and bibliography when you use a
\textsf{crossref} or an \textsf{xref} field.  The functionality is not
enabled by default, but you can enable it in the preamble or in the
\textsf{options} field using the \texttt{booklongxref} option.  Please
see \textbf{crossref} in section~\ref{sec:entryfields} and
\texttt{booklongxref} in section~\ref{sec:chicpreset}, below.

\mybigspace This \mymarginpar{\textbf{reference}} entry type is
aliased to \textsf{collection} by the standard \textsf{biblatex}
styles, but I intend it to be used in cases where you need to cite a
reference work but not an alphabetized entry or entries in that work.
This could be because it doesn't contain such entries, or perhaps
because you intend the citation to appear in a bibliography rather
than in notes.  Indeed, the only differences between it and
\textsf{inreference} are the lack of a \textsf{lista} field to present
an alphabetized entry, and the fact that any \textsf{postnote} field
will be printed verbatim, rather than formatted as an alphabetized
entry.  (See mla:style for an example of a reference work that uses
numbered sections rather than alphabetized entries, and that appears
in the bibliography as well.)

\mybigspace This \mymarginpar{\textbf{report}} entry type is a
\textsf{biblatex} generalization of the traditional \textsc{Bib}\TeX\
type \textsf{techreport}.  Instructions for such entries are rather
thin on the ground in the \emph{Manual} (8.186, 14.220), so I have
followed the generic advice about formatting it like a book, and hope
that the results conform to the specification.  At least one user has
indicated a need, now filled, for an \texttt{unpublished}
\textsf{entrysubtype}, which prints the \textsf{title} inside
quotation marks instead of in italics, but affects nothing else.  This
detail aside, the type's main peculiarities are the
\textsf{institution} field in place of a \textsf{publisher}, the
\textsf{type} field for identifying the kind of report in question,
the \textsf{number} field closely associated with the \textsf{type},
and the \textsf{isrn} field containing the International Standard
Technical Report Number of a technical report.  As in standard
\textsf{biblatex}, if you use a \textsf{techreport} entry, then the
\textsf{type} field automatically defaults to
\verb+\bibstring{techreport}+.  As with \textsf{booklet} and
\textsf{manual}, you can also use a \textsf{book} entry, putting the
report type in \textsf{note} and the \textsf{institution} in
\textsf{publisher}.  (See herwign:office.)

\mybigspace As \mymarginpar{\textbf{review}} its name suggests, the
\textsf{review} entry type was designed for reviews published in
periodicals, and if you've already read the \textsf{article}
instructions above --- if you haven't, I recommend doing so now ---
you'll know that \textsf{review} serves as well for citing other sorts
of material with generic titles, like letters to the editor,
obituaries, interviews, online comments and the like.  The primary
rule is that any piece that has only a generic title, like
\enquote{review of \ldots,} \enquote{interview with \ldots,} or
\enquote{obituary of \ldots,} calls for the \textsf{review} type.  Any
piece that also has a specific title, e.g., \enquote{\enquote{Lost in
    \textsc{Bib}\TeX,} an interview with \ldots,} requires an
\textsf{article} entry.  (This assumes the text is found in a
periodical of some sort.  Were it found in a book, then the
\textsf{incollection} type would serve your needs, and you could use
\textsf{title} and \textsf{titleaddon} there.  While we're on the
topic of exceptions, the \emph{Manual} includes an example (14.213)
where the \enquote{Interview} part of the title is considered a
subtitle rather than a titleaddon, said part therefore being included
inside the quotation marks and capitalized accordingly.  Not having
the journal in front of me I'm not sure what prompted that decision,
but \textsf{biblatex-chicago} would obviously have no difficulty
coping with such a situation.)

\mylittlespace Once you've decided to use \textsf{review}, then you
need to determine which sort of periodical you are citing, the rules
for which are the same as for an \textsf{article} entry.  If it is a
\enquote{magazine} or a \enquote{newspaper}, then you need an
\textsf{entrysubtype} \texttt{magazine}, or the synonymous
\textsf{entrysubtype} \texttt{newspaper}.  The generic title goes in
\textsf{title} and the other fields work just as as they do in an
\textsf{article} entry with the same \textsf{entrysubtype}, including
the substitution of the \textsf{journaltitle} for the \textsf{author}
if the latter is missing. (See 14.190--91, 14.195--96, 14.201--4,
14.213; barcott:review, bundy:macneil, Clemens:letter, gourmet:052006,
kozinn:review, nyt:obittrevor, nyt:trevorobit, unsigned:ranke,
wallraff:word.)  If, on the other hand, the piece comes from a
\enquote{journal,} then you don't need an \textsf{entrysubtype}.  The
generic title goes in \textsf{title}, and the remaining fields work
just as they do in a plain \textsf{article} entry.  (See 14.202;
ratliff:review.)

\mylittlespace \textsf{Biblatex-chicago} also, at the behest of
Bertold Schweitzer, supports the \textsf{relatedtype}
\texttt{reviewof}, which allows you to use the \textsf{related}
mechanism to provide information about the work being reviewed,
thereby simplifying how much information you need to provide in the
reviewing entry.  In particular, it relieves you of the need to
construct \textsf{title} or \textsf{titleaddon} fields like:
\verb+review of \mkbibemph{Book Title} by Author+, as the
\textsf{related} entry's \textsf{title} automatically provides the
\textsf{title} in the \textsf{review} type and the \textsf{titleaddon}
in the \textsf{article} type, with the \textsf{related} mechanism
providing the connecting string.  This may be particularly helpful if
you need to cite multiple reviews of the same work; please see section
\ref{sec:related} for further information.

\mylittlespace Most of the onerous details are the same as I described
them in the \textbf{article} section above, but I'll repeat some of
them briefly here.  If anything in the \textsf{title} needs
formatting, you need to provide those instructions yourself, as the
default is completely plain.  (As just mentioned, the \textsf{related}
mechanism provides this automatically.)  \textsf{Author}-less reviews
are treated just like similar newspaper articles --- in short notes
and in the bibliography the \textsf{journaltitle} replaces the author
and heads the entry, while in long notes the \textsf{title} comes
first.  The sorting of such entries is an issue, solved if you use
\textsf{Biber} as your backend, and otherwise requiring manual
intervention with, e.g., a \textsf{sortkey} (14.204; gourmet:052006,
nyt:trevorobit, unsigned:ranke, and see \cmd{DeclareSortingTemplate}
in section~\ref{sec:formatopts}, below.).  As in \textsf{misc} entries
with an \textsf{entrysubtype}, words like \enquote{interview,}
\enquote{review,} and \enquote{letter} only need capitalization after
a full stop, i.e., ordinarily in a bibliography and not a note, so
\textsf{biblatex-chicago-notes} automatically deals with this problem
itself if you start the \textsf{title} field with a lowercase letter.
The file \textsf{notes-test.bib} and the documentation of
\cmd{autocap} will provide guidance here.

\mylittlespace One detail of the \textsf{review} type is fairly new,
and in particular has changed between the 16th and 17th editions of
the \emph{Manual}.  As I mentioned above, blogs are best treated as
\textsf{articles} with \texttt{magazine} \textsf{entrysubtype},
whereas comments on those blogs --- or replies to those comments,
etc.\ --- need the \textsf{review} type with the same
\textsf{entrysubtype}.  (Neither need appear in the bibliography.)
What they also need is a date closely associated with the comment
(14.208; ac:comment), so I have included the \textsf{eventdate} in
\textsf{review} entries for just this purpose.  It will be printed
just after the \textsf{author} and before the \textsf{title}.  If you
need a time stamp in addition, as may frequently be the case with
multiple contributions by the same author to a single thread, then you
should now use the standard \textsf{biblatex} time-stamp format (e.g.,
\texttt{2008-07-01T10:18:00}) in the \textsf{eventdate} field itself,
which \textsf{biblatex-chicago} will format and print appropriately.
Please see the documentation concerning time stamps in
section~\ref{sec:entryfields}, s.v.\ \textsf{date}.  This change
allows the \textsf{nameaddon} field to revert to its primary use,
which is to provide extra information about the \textsf{author}.  In
blog comments, this could include the commenter's geographical
location, which you need to enclose in parentheses, as I've removed
the automatic square brackets from this field to allow it this more
general usefulness.  You can, of course, still provide your own square
brackets in \textsf{review} entries to indicate pseudonymous
authorship, which is the standard function of \textsf{nameaddon} in
most entry types.  The package options \texttt{nameaddonformat} and
\texttt{nameaddonsep} can help here, as well.  See
sections~\ref{sec:useropts} and \ref{sec:chicpreset}, below.

\mylittlespace In this context I should mention a small change to the
default behavior of \textsf{review} entries when they utilize a
\textsf{crossref} or \textsf{xref} field, as is really only useful
when the entry is a blog comment, as otherwise there won't generally
be any fields worth inheriting from the reviewed entry.  Assuming the
default values of the \textsf{biblatex-chicago} option
\texttt{longcrossref}, the driver now explicitly tests if the reviewed
entry has already been cited, and accordingly shortens the reviewing
citation, as the \emph{Manual} (14.208) suggests (viv:amlen).  (This
would be incorrect for, say, a book review, so you should either not
use a \textsf{crossref} or \textsf{xref} field there, or change the
state of the \texttt{longcrossref} option --- cf.\ the documentation
starting on page~\pageref{sec:crossref}, below.)  You
\mymarginpar{\texttt{blogurl}} can, if you wish, and while we're on
this subject, set the preamble option \texttt{blogurl} to allow your
child comments to inherit the URL from the parent blog.

\mylittlespace Another recent addition is the \textsf{relatedtype}
\texttt{commenton}, which allows you to simplify your\ .bib entries in
much the same way as the \texttt{reviewof} \textsf{relatedtype} does,
i.e., it constructs your \textsf{title} field for you (which the
\textsf{crossref} mechanism doesn't do).  It further allows you to
mimic thread structure by citing a chain of replies to comments on
blogs, etc., all in a single entry, while also simplifying your\ .bib
entries.  I've not tested how far down the rabbit hole you can go, but
a series of entries linked one to the next by this
\textsf{relatedtype} will all turn up if you cite the first in the
chain, though of course you can use the technique merely as a
convenient way to structure and simplify your\ .bib file, without
creating chains longer than 2 entries.  The default connecting string
is the localized \texttt{commenton}, but you can use
\textsf{relatedstring} to change it to \enquote{\texttt{reply to}} or
whatever else you need.  Please see the documentation of this
\textsf{relatedtype} in section~\ref{sec:related}, and also ellis:blog
and ac:comment.  Note also that this way of structuring your\ .bib
file is by no means required in \textsf{review} entries, though if you
want to cite replies and comments to social media threads, where you
need the \textsf{online} entry type, you will need to use this
\textsf{relatedtype}.

\mylittlespace Two more notes.  For the reasons I explained in the
\textsf{article} docs above, I have brought the \textsf{article} and
\textsf{review} entry types into line with most of the other types in
allowing the use of the \textsf{namea} and \textsf{nameb} fields in
order to associate an editor or a translator specifically with the
\textsf{title}.  The \textsf{editor} and \textsf{translator} fields,
in strict homology with other entry types, are associated with the
\textsf{issuetitle} if one is present, and with the \textsf{title}
otherwise.  The usual string concatenation rules still apply --- cf.\
\textsf{editor} and \textsf{editortype} in
section~\ref{sec:entryfields}, below.

\mylittlespace Finally, the special \textsf{biblatex} field
\textsf{shortjournal} allows you to present shortened
\textsf{journaltitles} in \textsf{review} entries, as well as in
\textsf{article} and \textsf{periodical} entries, and it facilitates
the creation of lists of journal abbreviations in the manner of a
\textsf{shorthand} list.  Please see the documentation of
\textbf{shortjournal} in section~\ref{sec:entryfields} for all the
details on how this works.

\mybigspace In \mymarginpar{\textbf{standard}} older releases it was
fairly straightforward to present published national or international
standards using a \textsf{book} entry, but with some additional
specifications now included in the 17th edition of the \emph{Manual}
(14.259) I think it might be helpful to provide a separate entry type.
The \textbf{standard} type has long existed in \textsf{biblatex},
though none of its included styles use it.  In
\textsf{biblatex-chicago} constructing such an entry is mostly
straightforward.  The organization responsible for the standard goes
in \textsf{organization}, the title in \textsf{title}, and the
\textsf{series} and \textsf{number} fields provide the ID of the
standard.  The \textsf{date} field generally provides the publication
date, though for some standards there may also be a later
reaffirmation date (or similar), for which you can use the
\textsf{eventdate}.

\mylittlespace Now, for the peculiarities.  In the bibliography, the
\textsf{organization} will appear at the head of the entry, and will
be reprinted as the publisher.  If you wish to provide a shortened
version for the second appearance, then the \textsf{publisher} field
is the place for it.  In long notes, the entry starts with the
\textsf{title}, so there the code prefers the \textsf{organization} as
publisher, because its shortened version may not be immediately
recognizable.  In short notes, only the \textsf{title} will appear
(along with any \textsf{pre} or \textsf{postnote} fields, obviously).
You can use the \textsf{author} field in addition to the
\textsf{organization}, but this is unnecessary.  If you absolutely
must have the \textsf{organization} or \textsf{author} appear at the
head of long and short notes, then providing any \textsf{entrysubtype}
whatsoever will accomplish this.  Any named \textsf{editor} or
\textsf{namec} will, as per the specification, \emph{not} appear at
the head of entries.  You can really only alter this by using a
\textsf{book} entry, instead.  (Cf.\ w3c:xml.)

\mylittlespace Finally, it is distinctly possible that an entry with
two dates will need somehow to specify just what sort of dates are
involved.  The usual \textsf{biblatex-chicago} method is the
\textsf{userd} field, and here that field will act as a date-type for
the \textsf{date} field itself.  For the \textsf{eventdate}, you'll
need to use \textsf{howpublished}, which I have commandeered for this
purpose in a few other entry types, as well.  (Cf.\ niso:bibref and
\textbf{howpublished} in section~\ref{sec:entryfields}, below.)

\mybigspace This \mymarginpar{\textbf{suppbook}} is the entry type to
use if the main focus of a reference is supplemental material in a
book or in a collection, e.g., an introduction, afterword, or forward,
either by the same or a different author.  In previous releases of
\textsf{biblatex-chicago} these three just-mentioned types of
material, and only these three types, could be referenced using the
\textsf{introduction}, \textsf{afterword}, or \textsf{foreword}
fields, a system that required you simply to define one of them in any
way and leave the others undefined.  The macros don't use the text
provided by such an entry, they merely check to see if one of them is
defined, in order to decide which sort of pre- or post-matter is at
stake, and to print the appropriate string before the \textsf{title}
in long notes, short notes, list of shorthands, and bibliography.  I
have retained this mechanism both for backward compatibility and
because it works without modification across multiple languages, but
have also added functionality which allows you to cite any sort of
supplemental material whatever, using the \textsf{type} field.  Under
this system, simply put the nature of the material, including the
relevant preposition, in that field, beginning with a lowercase letter
so \textsf{biblatex} can decide whether it needs capitalization
depending on the context.  Examples might be \enquote{\texttt{preface
    to}} or \enquote{\texttt{colophon of}.}  (Please note, however,
that unless you use a \cmd{bibstring} command in the \textsf{type}
field, the resultant entry will not be portable across languages.)

\mylittlespace There are a few other rules for constructing your .bib
entry.  The \textsf{author} field refers to the author of the
introduction or afterword, while \textsf{bookauthor} refers to the
author of the main text of the work, if the two differ.  The
\emph{Manual} requires the inclusion of the page range of the part in
question, though \emph{only} in the bibliography.  I have followed
this advice literally, so the \textsf{pages} field of a
\textsf{suppbook} entry won't automatically appear in a long note.  If
you wish to include those pages in a note, then you'll need to repeat
them in the \textsf{postnote} field of the citation command.

\mylittlespace Finally, if the focus of the reference is the main text
of the book, but you want to mention the name of the writer of an
introduction or afterword for bibliographical completeness, then the
normal \textsf{biblatex} rules apply, and you can just put their name
in the appropriate field of a \textsf{book} entry, that is, in the
\textsf{foreword}, \textsf{afterword}, or \textsf{introduction} field.
(See \emph{Manual} 14.110; polakow:afterw, prose:intro).

\mybigspace This \mymarginpar{\textbf{suppcollection}} fulfills a
function analogous to \textsf{suppbook}.  Indeed, I believe the
\textbf{suppbook} type can serve to present supplemental material in
both types of work, so this entry type is an alias to
\textsf{suppbook}, which see.

\mybigspace This \mymarginpar{\textbf{suppperiodical}} type is
intended to allow reference to generically-titled works in
periodicals, such as regular columns or letters to the editor.
\textsf{Biblatex-chicago-notes} provides the \textsf{review} type for
this purpose, and you can use either of these, as I've added
\textsf{suppperiodical} as an alias of \textsf{review}.  Please see
above under \textbf{review} for the full instructions on how to
construct a .bib entry for such a reference.

\mybigspace The \mymarginpar{\textbf{unpublished}}
\textsf{unpublished} entry type works largely as it does in standard
\textsf{biblatex}, though it's worth remembering that you should use a
lowercase letter at the start of your \textsf{note} field (or perhaps
an\ \cmd{autocap} command in the somewhat contradictory
\textsf{howpublished}, if you have one) for material that wouldn't
ordinarily be capitalized except at the beginning of a sentence.
Thanks to a bug report by Henry D. Hollithron, such entries will print
information about any \textsf{editor}, \textsf{translator},
\textsf{compiler}, etc., that you include in the .bib file.  Also,
conforming to the indications of the \emph{Manual}, and thanks to the
prompting of Jan David Hauck, you can use the \textsf{venue},
\textsf{eventdate}, \textsf{eventtitle}, and \textsf{eventtitleaddon}
fields further to specify unpublished conference papers and the like
(14.216--18; nass:address).

\mybigspace This \mymarginpar{\textbf{video}} is the last of the
three audiovisual entry types, and as its name suggests it is intended
for citing visual media, be it films of any sort or TV shows, whether
broadcast, on the Net, on VHS, DVD, or Blu-ray, though it will serve
as well, I think, for radio broadcasts of plays or drama serials.  As
with the \textsf{music} type discussed above, certain choices had to
be made when associating the production roles found, e.g., on a DVD,
to those bookish ones provided by \textsf{biblatex}.  The 17th edition
of the \emph{Manual} once again revised its recommendations for this
type, but fortunately the changes are additive, i.e., you can re-use
16th-edition citations but are encouraged to peruse the following
guidelines to see if there's any information you might think of adding
to bring your citations more into line with the spec.  Here are the
main guidelines:

{\renewcommand{\descriptionlabel}[1]{\qquad\textsf{#1}}
\begin{description}
\item[author:] This will not infrequently be left undefined, as the
  director of a film should be identified as such and therefore placed
  in the \textsf{editor} field with the appropriate
  \textsf{editortype} (see below).  You will need it, however, to
  identify the composer of, e.g., an oratorio on VHS (handel:messiah),
  or perhaps the provider of commentaries or other extras on a film
  DVD (cleese:holygrail).
\item[editor, editora, editorb =] director or producer, or possibly
  the performer or conductor in recorded musical performances.  These
  will ordinarily follow the \textsf{title} of the work, though the
  usual \texttt{useauthor} and \texttt{useeditor} options can alter
  the presentation within an entry.  Because these are non-standard
  roles, you will need to identify them using the following:
\item[editortype, editoratype, editorbtype:] The most common roles,
  all associated with specific bibstrings (or their absence), will
  likely be \texttt{director}, \texttt{producer}, and, oddly,
  \texttt{none}.  The last is particularly useful if you want to
  identify performers, as they usually don't need further specifying
  and this role prevents \textsf{biblatex} from falling back on the
  default \texttt{editor} bibstring.  Any other roles you want to
  emphasize, even if there is no pre-defined \cmd{bibstring}, can be
  provided here, and will be printed as-is, contextually capitalized.
  (Cf.\ hitchcock:nbynw.)
\item[title, titleaddon, booktitle, booktitleaddon, maintitle:] As
  with the other two audiovisual types, \textsf{video} serves as an
  analogue both to books and to collections, so the \textsf{title} may
  be of a whole film DVD or of a TV series, or it may identify one
  episode in a series or one scene in a film.  In the latter cases,
  the title of the whole would go in \textsf{booktitle}.  The
  \textsf{booktitleaddon} field is the place for specifying the season
  and/or episode number of a TV series, while the \textsf{titleaddon}
  is for any information that needs to come between the \textsf{title}
  and the \textsf{booktitle} (american:crime, cleese:holygrail,
  friends\hc leia, handel:messiah, hitchcock:nbynw,
  mayberry:brady).  As in the \textsf{music} type, a
  \textsf{maintitle} may be necessary for a boxed set or something
  similar.
\item[entrysubtype:] If, for some reason, you want to cite an
  individual episode or scene without reference to any larger unit,
  then your entry will contain only a \textsf{title}, which
  \textsf{biblatex-chicago} would normally interpret as the title of a
  complete film or TV series.  In such a case, you'll need to define
  an \textsf{entrysubtype} to let it know that the lone \textsf{title}
  is such a sub-unit.  In quite a different syntactic transformation,
  the 17th edition (14.265) now recommends that, when presenting
  episodes from a TV series, the name of the series
  (\textsf{booktitle}) comes before the episode name (\textsf{title}).
  The exact string \texttt{tvepisode} in the \textsf{entrysubtype}
  field achieves this reversal, which includes using the
  \textsf{booktitle} as a \textsf{sorttitle} in the bibliography and
  also as the \textsf{labeltitle} in short notes.
\item[usera:] When citing recordings of TV shows, the 17th edition now
  wants you to include the TV network for the original broadcast, and
  the \textsf{usera} field is the place for it.
  \textsf{Biblatex-chicago} has long used this field for this same
  purpose in \textsf{article, periodical}, and \textsf{review} entries
  with a \texttt{magazine} \textsf{entrysubtype}, so its inclusion
  here can at least hope to benefit from that prior acquaintance.  It
  will appear after the broadcast date, i.e., the \textsf{eventdate},
  and will be separated from it by the \cmd{bibstring}
  \enquote{\texttt{on.}}
\item[date, eventdate, origdate:] As with \textsf{music} entries, in
  order to follow the specifications of the \emph{Manual}, we need
  to provide three separate date fields for citing \textsf{video}
  sources, but their uses differ somewhat between the two types.  In
  both, the \textsf{date} will generally provide the publishing or
  copyright date of the medium you are referencing.  The
  \textsf{eventdate} will most commonly present either the broadcast
  date of a particular TV program, or the recording/performance date
  of, for example, an opera on DVD.  The style will automatically
  prepend the bibstring \texttt{broadcast} to such a date, though you
  can use the \textsf{userd} field to change the string printed there.
  (Absent an \textsf{eventdate}, the \textsf{userd} field in
  \textsf{video} entries will modify the \textsf{urldate}, and absent
  those two it will modify the \textsf{date}.)  The \textsf{origdate}
  has more or less the same function, and appears in the same places,
  as it does in standard book-like entries, providing the date of
  first release of a film, though there isn't any \texttt{reprint}
  string associated with it in this entry type.  Cf.\ friends:leia,
  handel:messiah, hitchcock:nbynw.
\item[type:] As in all the audiovisual entry types, the \textsf{type}
  field holds the medium of the \textsf{title}, e.g., 8 mm, VHS, DVD,
  Blu-ray, MPEG.
\end{description}}

As with the \textsf{music} type, entries in \textsf{notes-test.bib}
should at least give you a good idea of how all this works.  (Cf.\
14.265, 14.267; loc:city, weed:flatiron.)

\subsection{Entry Fields}
\label{sec:entryfields}

The following discussion presents, in alphabetical order, a complete
list of the entry fields you will need to use
\textsf{biblatex-chicago-notes}.  As in section \ref{sec:entrytypes},
I shall include references to the numbered paragraphs of the
\emph{Chicago Manual of Style}, and also to the entries in
\textsf{notes-test.bib}.  Many fields are most easily understood with
reference to other, related fields.  In such cases, cross references
should allow you to find the information you need.

\mybigspace As \mymarginpar{\textbf{addendum}} in standard
\textsf{biblatex}, this field allows you to add miscellaneous
information to the end of an entry, after publication data but before
any \textsf{url} or \textsf{doi} field.  In the \textsf{patent} entry
type (which see), it will be printed in close association with the
filing and issue dates.  In all other entry types this information
will come \emph{after} any \textsf{pages} or \textsf{postnote}
references present in long notes, allowing you in particular to use
the field to identify a particular type of book-like publication when
such data won't fit well in another part of an entry.  In any entry
type, if your data begins with a word that would ordinarily only be
capitalized at the beginning of a sentence, then simply ensure that
that word is in lowercase, and the style will take care of the rest.
Cf.\ \textsf{note}. (See \emph{Manual} 14.114, 14.159--63;
davenport:attention, natrecoff:camera.)

\mybigspace In most \mymarginpar{\textbf{afterword}} circumstances,
this field will function as it does in standard \textsf{biblatex},
i.e., you should include here the author(s) of an afterword to a given
work.  The \emph{Manual} suggests that, as a general rule, the
afterword would need to be of significant importance in its own right
to require mentioning in the reference apparatus, but this is clearly
a matter for the user's judgment.  As in \textsf{biblatex}, if the
name given here exactly matches that of an editor and/or a translator,
then \textsf{biblatex-chicago-notes} will concatenate these fields in
the formatted references.

%%\enlargethispage{-\baselineskip}

\mylittlespace As noted above, however, this field has a special
meaning in the \textsf{suppbook} entry type, used to make an
afterword, foreword, or introduction the main focus of a citation.  If
it's an afterword at issue, simply define \textsf{afterword} any way
you please, leave \textsf{foreword} and \textsf{introduction}
undefined, and \textsf{biblatex-chicago-notes} will do the rest. Cf.\
\textsf{foreword} and \textsf{introduction}. (See \emph{Manual}
14.105, 14.110; polakow:afterw.)

\paragraph*{\protect\mymarginpar{\textbf{annotation}}}
\label{sec:annote}

At the request of Emil Salim, \textsf{biblatex-chicago-notes} provides
a package option (see \texttt{annotation} below, section
\ref{sec:useropts}) to allow you to produce annotated bibliographies.
A recent feature request by Moritz Wemheuer referenced a
\href{https://tex.stackexchange.com/questions/528374/moving-addendum-field-to-the-end-in-biblatex-chicago/540755#540755}{StackExchange}
question which suggested that the possible uses for the
\textsf{annotation} field could well be more extensive, appearing as
it does at the very end of all entry types.  I have therefore modified
the \texttt{annotation} option so that you can print the field in the
bibliography (\texttt{=bib} or \texttt{=true}, the default), in long
notes (\texttt{=notes}), in both (\texttt{=all}), or in neither
(\texttt{=false}).  The two options \texttt{bibannotesep} and
\texttt{citeannotesep} allow you to choose the separator between the
rest of the entry and the \textsf{annotation}, and to choose a
different one in notes and bibliography.  The default formatting in
the bibliography (\texttt{vpar}) is to print the \textsf{annotation}
as a separate block using\ \verb+\par\nobreak\vskip\bibitemsep #1+,
while in long notes the default (\texttt{period}) is to print it
simply as an additional field, separated by a period.  The
\emph{Manual's} guidelines (14.64) allow for both these possibilities,
and I have provided a range of others, for which you should consult
the full documentation in section~\ref{sec:chicpreset}.  (Note also
that both options can be set globally or per-type in the preamble, or
per-entry in the \textsf{options} field of individual entries.  For
specialized needs, of course, you can re-declare the format
[\verb+\DeclareFieldFormat{annotation}+] in your preamble, or redefine
the \cmd{bibannotesep} and \cmd{citeannotesep} commands there.)  In
section~\ref{sec:useropts} you will find instructions for employing
the \texttt{formatbib} and \texttt{entrybreak} options to give you
fine-grained control over the formatting of the entire bibliography,
particularly with regard to \TeX's page-breaking algorithms.  The aim
is to remove, in most cases, any need for you to delve into the
low-level commands involved in these algorithms.

\mybigspace I \mymarginpar{\textbf{annotator}} have implemented this
\textsf{biblatex} field pretty much as that package's standard styles
do, even though the \emph{Manual} doesn't actually mention it.  It may
be useful for some purposes.  Cf.\ \textsf{commentator}.

\mybigspace For \mymarginpar{\textbf{author}} the most part, I have
implemented this field in a completely standard fashion.  Remember
that corporate or organizational authors need to have an extra set of
curly braces around them (e.g., \texttt{\{\{Associated Press\}\}}\,)
to prevent \textsf{biber} from treating one part of the name as a
surname (14.84, 14.200; assocpress:gun, chicago:manual).  If there is
no \textsf{author}, then \textsf{biblatex-chicago-notes} will, in the
bibliography and long notes, look in sequence, for a \textsf{namea},
an \textsf{editor}, a \textsf{nameb}, a \textsf{translator}, or a
\textsf{namec} (i.e., a compiler) and use that name (or those names)
instead, followed by the appropriate identifying string (esp.\ 14.103,
also 14.76, 14.121, 14.126, 14.180; boxer:china, brown:bremer,
harley:cartography, schellinger:novel, sechzer:women, silver:gawain,
soltes:georgia).  \textsf{Biblatex's} sorting algorithms will use the
first of those names found, which should ensure correct
alphabetization in the bibliography.  (See
\cmd{DeclareSortingTemplate} in section~\ref{sec:formatopts}, below.)
In short notes, where the \textsf{labelname} is used, the order
searched is somewhat augmented: \textsf{shortauthor, author,
  shorteditor, namea, editor, nameb, translator, namec}.  (See
\cmd{DeclareLabelname} in section~\ref{sec:formatopts}.)

\mylittlespace In the rare cases when this substitution mechanism
isn't appropriate, you have (at least) two options: either you can
(chaucer:liferecords) put all the information into a \textsf{note}
field rather than individual fields, or you can use the standard
\textsf{biblatex} options \texttt{useauthor=\hfill false},
\texttt{usenamea=false}, \texttt{useeditor=false},
\texttt{usenameb=false}, \texttt{usetranslator=false}, and
\texttt{usenamec=false} in the \textsf{options} field (chaucer:alt).
If you look at the chaucer:alt entry in \textsf{notes-test.bib},
you'll notice that you only need to turn off the fields that are
present in the entry, but please remember to use the new option
\texttt{usenamec} instead of the old \texttt{usecompiler}, as the
latter doesn't work as smoothly and completely as \textsf{biblatex's}
own name toggles.

% %\enlargethispage{\baselineskip}

\mylittlespace This system of options, then, can turn off
\textsf{biblatex-chicago-notes}'s mechanism for finding a name to
place at the head of an entry, but it also very usefully adds the
possibility of citing a work with an \textsf{author} by its editor,
compiler or translator instead (14.104; eliot:pound), something that
wasn't possible before.  For full details of how this works, see the
\textsf{editortype} documentation below.  (Of course, in
\textsf{collection}, \textsf{periodical} and \textsf{proceedings}
entries, an \textsf{author} isn't expected, so there the chain of
substitutions starts with \textsf{namea} and \textsf{editor}.  Also,
in \textsf{article} or \textsf{review} entries with
\textsf{entrysubtype} \texttt{magazine}, the absence of an
\textsf{author} triggers the use of the \textsf{journaltitle} in its
stead.  See those entry types for further details.)

\mylittlespace \textbf{NB}: The \emph{Manual} provides specific
instructions for formatting the names of both anonymous and
pseudonymous authors (14.79--82).  In the former case, if no author is
known or guessed at, then it may simply be omitted
(virginia:plantation).  The use of \enquote{Anonymous} as the name is
\enquote{generally to be avoided,} but may in some cases be useful
\enquote{in a bibliography in which several anonymous works need to be
  grouped.}  If, on the other hand, \enquote{the authorship is known
  or guessed at but was omitted on the title page,} then you need to
use the \textsf{authortype} field to let
\textsf{biblatex-chicago-notes} know this fact.  If the author is
known (horsley:prosodies), then put \texttt{anon} in the
\textsf{authortype} field, if guessed at (cook:sotweed) put
\texttt{anon?}\ there.  (In both cases,
\textsf{biblatex-chicago-notes} tests for these \emph{exact} strings,
so check your typing if it doesn't work.)  This will have the effect
of enclosing the name in square brackets, with or without the question
mark indicating doubt.  As long as you have the right string in the
\textsf{authortype} field, \textsf{biblatex-chicago-notes} will also
do the right thing automatically in the short note form.

\mylittlespace In most entry types (except \textsf{customc}), the
\textsf{nameaddon} field furnishes the means to cope with the case of
pseudonymous authorship.  If the author's real name isn't known,
simply put \texttt{pseud.}\,(or \verb+\bibstring{pseudonym}+) in that
field (centinel:letters).  If you wish to give a pseudonymous author's
real name, simply include it there, formatted as you wish it to
appear, as the contents of this field won't be manipulated as a name
by \textsf{biblatex} (lecarre:quest).  If you have given the author's
real name in the \textsf{author} field, then the pseudonym goes in
\textsf{nameaddon}, in the form \texttt{Firstname Lastname,\,pseud.}\
(creasey\hc ashe:blast, creasey:morton:hide, creasey:york:death).
This latter method will allow you to keep references to one author's
work under different pseudonyms grouped together in the bibliography,
as recommended by the \emph{Manual}, though it is now recommended
that, whichever system you employ, you include a cross-reference from
one name to the other in the bibliography.  You can do this using a
\textsf{customc} entry (ashe:creasey, morton:creasey, york:creasey).
Please see the entry on \textbf{nameaddon}, below, for circumstances
where you may need to provide your own square brackets when presenting
a pseudonym, and also the package options \texttt{nameaddonformat} and
\texttt{nameaddonsep} in sections~\ref{sec:useropts} and
\ref{sec:chicpreset}, below.

\mybigspace In \mymarginpar{\textbf{authortype}}
\textsf{biblatex-chicago}, this field serves a function very much in
keeping with the spirit of standard \textsf{biblatex}, if not with its
letter.  Instead of allowing you to change the string used to identify
an author, the field allows you to indicate when an author is
anonymous, that is, when their name doesn't appear on the title page
of the work you are citing.  As I've just detailed under
\textsf{author}, the \emph{Manual} generally discourages the use of
\enquote{Anonymous} as an author, preferring that you simply omit it.
If, however, the name of the author is known or guessed at, then
you're supposed to enclose that name within square brackets, which is
exactly what \textsf{biblatex-chicago} does for you when you put
either \texttt{anon} (author known) or \texttt{anon?} (author guessed
at) in the \textsf{authortype} field.  (Putting the square brackets in
yourself doesn't work right, hence this mechanism.)  The macros test
for these \emph{exact} strings, so check your typing if you don't see
the brackets.  Assuming the strings are correct,
\textsf{biblatex-chicago-notes} will also automatically do the right
thing in the short note form.  Cf.\ \textsf{author}.  (See 14.79--80;
cook:sotweed, horsley:prosodies.)

\mylittlespace The \emph{Manual} doesn't clarify how to treat multiple
works by the same \textsf{author}, in one or more of which their name
doesn't appear on the title page.  By default,
\textsf{biblatex-chicago} will, after the first appearance in the
bibliography, replace identical \textsf{authors} with the 3-em dash,
regardless of any \textsf{authortype} field that may be present.  If
you want to distinguish between works certainly written by and works
merely ascribed to a given author, then you can use the
\texttt{dashed} option in the \textsf{options} field of individual
entries, and possibly also a \textsf{sortname}, to get the results you
want.

\mybigspace For \mymarginpar{\textbf{bookauthor}} the most part, as in
\textsf{biblatex}, a \textsf{bookauthor} is the author of a
\textsf{booktitle}, so that, for example, if one chapter in a book has
different authorship from the book as a whole, you can include that
fact in a reference (will:cohere).  Keep in mind, however, that the
entry type for introductions, forewords and afterwords
(\textsf{suppbook}) uses \textsf{bookauthor} as the author of
\textsf{title} (polakow:afterw, prose:intro).

\enlargethispage{2\baselineskip}

\mybigspace This, \mymarginpar{\vspace{-8pt}\textbf{bookpagination}}
a standard \textsf{biblatex} field, allows you automatically to prefix
the appropriate string to information you provide in a \textsf{pages}
field.  If you leave it blank, the default is to print no identifying
string (the equivalent of setting it to \texttt{none}), as this is the
practice the \emph{Manual} recommends for nearly all page numbers.
Even if the numbers you cite aren't pages, but it is otherwise clear
from the context what they represent, you can still leave this blank.
If, however, you specifically need to identify what sort of unit the
\textsf{pages} field represents, then you can either hand-format that
field yourself, or use one of the provided bibstrings in the
\textsf{bookpagination} field.  These bibstrings currently are
\texttt{column,} \texttt{line,} \texttt{paragraph,} \texttt{page,}
\texttt{section,} and \texttt{verse}, all of which are used by
\textsf{biblatex's} standard styles.

\mylittlespace There are two points that may need explaining here.
First, all the bibstrings I have just listed follow the Chicago
specification, which may be confusing if they don't produce the
strings you expect.  Second, remember that \textsf{bookpagination}
applies only to the \textsf{pages} field --- if you need to format a
citation's \textsf{postnote} field, then you must use
\textsf{pagination}, which see (10.42--43, 14.147--56).

\mybigspace The \mymarginpar{\textbf{booksubtitle}} subtitle for a
\textsf{booktitle}.  See the next entry for further information.

\mybigspace In \mymarginpar{\textbf{booktitle}} the
\textsf{bookinbook}, \textsf{inbook}, \textsf{incollection},
\textsf{inproceedings}, and \textsf{letter} entry types, the
\textsf{booktitle} field holds the title of the larger volume in which
the \textsf{title} itself is contained as one part.  It is important
not to confuse this with the \textsf{maintitle}, which holds the more
general title of multiple volumes, e.g., \emph{Collected Works}.  It
is perfectly possible for one .bib file entry to contain all three
sorts of title (euripides:orestes, plato:republic:gr).  You may also
find a \textsf{booktitle} in other sorts of entries (e.g.,
\textsf{book} or \textsf{collection}), but there it will almost
invariably be providing information for the traditional
\textsc{Bib}\TeX\ cross-referencing apparatus, which I discuss below
(\textbf{crossref}).  This provision is now unnecessary, assuming you
are using \textsf{biber}.

\mybigspace An \mymarginpar{\textbf{booktitleaddon}} annex to the
\textsf{booktitle}.  It will be printed in the main text font, without
quotation marks.  If your data begins with a word that would
ordinarily only be capitalized at the beginning of a sentence, then
simply ensure that that word is in lowercase, and
\textsf{biblatex-chicago-notes} will automatically do the right thing.
The package and entry options \texttt{ptitleaddon} and
\texttt{ctitleaddon} (section~\ref{sec:chicpreset}) allow you to
customize the punctuation that appears before the
\textsf{booktitleaddon} field.

\mybigspace This \mymarginpar{\textbf{chapter}} field holds the
chapter number, mainly useful only in an \textsf{inbook} or an
\textsf{incollection} entry where you wish to cite a specific chapter
of a book (ashbrook:brain).  It now also holds the track number of
individual pieces of \textsf{music}, whether on a traditional format
or on a streaming service (holiday:fool, rihanna:umbrella).

\mybigspace I \mymarginpar{\textbf{commentator}} have implemented this
\textsf{biblatex} field pretty much as that package's standard styles
do, even though the \emph{Manual} doesn't actually mention it.  It may
be useful for some purposes.  Cf.\ \textsf{annotator}.

\paragraph*{\protect\mymarginpar{\textbf{crossref}}}
\label{sec:crossref}

This field is the standard \textsc{Bib}\TeX\ cross-referencing
mechanism, and \textsf{biblatex} has adopted it while also introducing
a modified one of its own (\textsf{xref}).  If you have used
\textsc{Bib}\TeX\ (or \textsf{bibtex8)} the \textsf{crossref} field
works much the same as it always has, while \textsf{xref} attempts to
remedy some of the deficiencies of the usual mechanism by ensuring
that child entries will inherit no data at all from their parents.
Section~2.4.1 of \textsf{biblatex.pdf} contains useful notes on
managing cross-referenced entries, and section~3.15 explains some of
the limitations of the traditional backends, which offer only a small
subset of \textsf{Biber's} features.  The functionality, discussed
below, for abbreviating references in \textsf{book},
\textsf{bookinbook}, \textsf{collection}, and \textsf{proceedings}
entries, and for using the \textsf{mv*} entry types to do so, will
prove extremely difficult to replicate with the older backends, so if
you plan on lots of cross-referencing in
\textsf{biblatex-chicago-notes} then I strongly recommend you use
\textsf{Biber}.

\mylittlespace (One reason for this is that when \textsf{Biber} is the
backend, \textsf{biblatex} defines a series of inheritance rules for
the \textsf{crossref} field which make it much more convenient to use.
Appendix B of \textsf{biblatex.pdf} explains the defaults, to which
\textsf{biblatex-chicago} has added several that I should mention
here: \textsf{incollection} entries can now inherit from \textsf{book}
and \textsf{mvbook} just as they do from \textsf{collection} and
\textsf{mvcollection} entries; \textsf{letter} entries now inherit
from \textsf{book}, \textsf{collection}, \textsf{mvbook}, and
\textsf{mvcollection} entries the same way an \textsf{inbook} or an
\textsf{incollection} entry would; the \textsf{namea}, \textsf{nameb},
\textsf{sortname}, \textsf{sorttitle}, and \textsf{sortyear} fields,
all highly single-entry specific, are no longer inheritable; and
\textsf{date} and \textsf{origdate} fields are not inheritable from
any of the new \textbf{mv*} entry types.)

\mylittlespace Turning now to the provision of abbreviated references
in \textsf{biblatex-chicago-notes}, the \emph{Manual} (14.108)
specifies that if you cite several contributions to the same
collection, all (including the collection itself) may be listed
separately in the bibliography, which the package does automatically,
using the default inclusion threshold of 2 in the case both of
\textsf{crossref}'ed and \textsf{xref}'ed entries.  (The familiar
\cmd{nocite} command may also help in some circumstances.)  In
footnotes the specification suggests that, after a citation of any one
contribution to the collection, all subsequent contributions may, even
in the first, long footnote, be cited using a slightly shortened form,
thus \enquote{avoiding clutter.}  In the bibliography the abbreviated
form is appropriate for all the child entries.  The
\textsf{biblatex-chicago-notes} package has always implemented these
instructions, but only if you use a \textsf{crossref} or an
\textsf{xref} field, and only in \textsf{incollection},
\textsf{inproceedings}, or \textsf{letter} entries (on the last named,
see just below).  Recent releases have considerably extended this
functionality.

\mylittlespace First, I added five entry types --- \textbf{book},
\textbf{bookinbook}, \textbf{collection}, \textbf{inbook},
\textbf{proceedings}, and \textbf{review} --- to the list of those
which use shortened cross references, and I added two options ---
\texttt{longcrossref} and \texttt{booklongxref}, on which more below
--- which you can use in the preamble or in the \textsf{options} field
of an entry to enable or disable the automatic provision of
abbreviated references.  (The \textsf{crossref} or \textsf{xref} field
are still necessary for this provision, but they are no longer
sufficient on their own.)  The \textsf{inbook} and \textsf{review}
types work exactly like \textsf{incollection} or
\textsf{inproceedings}; in previous releases, you could use
\textsf{inbook} instead of \textsf{incollection} to avoid the
automatic abbreviation, the two types being otherwise identical.  Now
that you can use an option to turn off abbreviated references even in
the presence of a \textsf{crossref} or \textsf{xref} field, I have
thought it sensible to include this entry type alongside the others.
(Cf.\ ellet:galena, keating:dearborn, lippincott:chicago, and
prairie:state to see this mechanism in action in both notes and
bibliography.)  In the \textsf{review} type the mechanism is aimed
primarily at blog comments, assuming you don't want to use the more
convenient \texttt{commenton} \textsf{relatedtype}, which absolves you
even of the need to provide a \textsf{title} field for such entries.
See the documentation of the \textsf{review} type above for the
details.

\mylittlespace The inclusion of \textbf{book}, \textbf{bookinbook},
\textbf{collection}, and \textbf{proceedings} entries fulfills a
request made by Kenneth L. Pearce, and allows you to obtain shortened
references to, for example, separate volumes within a multi-volume
work, or to different book-length works collected inside a single
volume.  Such references are not an explicit part of the
\emph{Manual's} specification, but they are a logical extension of it,
so the system of options for turning on this functionality behaves
differently for these four entry types than for the other 4 (see
below).  In \textsf{notes-test.bib} you can get a feel for how this
works by looking at bernhard:boris, bernhard:ritter,
bernhard:themacher, harley:ancient:cart, harley\hc cartography, and
harley:hoc.

\mylittlespace Before discussing the new package options, I should say
a little about some subtleties involved in this mechanism.  First, and
especially for \textsf{book}, \textsf{bookinbook},
\textsf{collection}, and \textsf{proceedings} entries, it is much
simpler if your backend is \textsf{Biber}, which allows you to provide
\textsf{maintitles} by cross-referencing an \textbf{mv*} entry, and
\textsf{booktitles} by cross-referencing \textsf{book} or
\textsf{collection} entries.  Second, where and when to print
\textsf{volume} information in these references is extremely complex,
and I confess that I designed the tests primarily with \textsf{Biber}
in mind.  Third, Andrew Goldstone long ago identified some other
difficulties in the package's treatment of abbreviated citations, both
in notes and bibliography, difficulties exacerbated now by the
extension of the mechanism to book-like entries.  If you refer
separately to chapters in a single-author \textsf{book}, then the
shortened part of the reference, to the whole book, won't repeat the
author's name before the title of the whole.  If, however, you refer
separately to parts of a \textsf{collection} or \textsf{proceedings},
even when the \textsf{editor} of the \textsf{collection} is the same
as the \textsf{author} of an essay in the collection, you will see the
name repeated before the abbreviated part referencing the whole parent
volume.

\mylittlespace Shortened references to book-like entries require, I
believe, a somewhat different treatment.  Here, repeated
\textsf{editors} are avoided if the abbreviated reference is to a
\textsf{collection} or \textsf{proceedings} entry, or to either of
their \textsf{mv*} versions, while for other entry types repeated
\textsf{authors} are avoided.  Because the code in these situations
tests for entry type, there may be corner cases where careful choice
of the parent entry type gets you what you want.  Likewise, judicious
use of the \textsf{editor} and \textsf{editortype} fields may also
help, in some circumstances, to clear names that are repeated
unnecessarily.  Also, because of the way dates are handled by the
\textsf{mv*} entry types, and by child entries cross-referenced to
such entry types, I thought it might help in these abbreviated
book-like entries to provide a date for the \textsf{title} when it's
part of a \textsf{maintitle}, though not when it's only part of a
\textsf{booktitle}.  If dates appear in shortened references where
you'd rather not have them, I have provided the \texttt{omitxrefdate}
option to turn them off, either in the preamble for the document as a
whole or in the \textsf{options} field of individual entries.  There
is also an \texttt{xrefurl} option available to control the printing
of \textsf{url}, \textsf{doi}, and \textsf{eprint} fields in
abbreviated references where such information might otherwise never
appear.  See \textbf{mvbook} in section~\ref{sec:entrytypes}, and both
\texttt{omitxrefdate} and \texttt{xrefurl} in
section~\ref{sec:useropts}.

\mylittlespace Finally, a published collection of letters also
requires different treatment (14.111).  If you cite more than one
letter from the same collection, then the \emph{Manual} specifies that
only the collection itself should appear in the bibliography.  In
footnotes, you can use the \textsf{letter} entry type, documented
above, for each individual letter, while the collection as a whole may
well require a \textsf{book} entry.  I have, after some consideration,
implemented the system of shortened references in \textsf{letter}
entries, even though the \emph{Manual} doesn't explicitly require it.
(See white:ross:memo, white:russ, and white:total, for examples of the
\textsf{crossref} field in action in this way, and please note that
the second of these entries is entirely fictitious, provided merely
for the sake of example.)  How then to keep the individual letters
from appearing in the bibliography?  The simplest mechanism is
probably just to use \enquote{\texttt{skipbib}} in the
\textsf{options} field.

\mylittlespace Returning, \mymarginpar{\texttt{longcrossref}} then, to
the package options which control whether and where the abbreviated
references appear, they function, by default, asymmetrically.  The
first, \texttt{longcrossref}, generally controls the settings for the
entry types more-or-less authorized by the \emph{Manual}:
\textsf{inbook}, \textsf{incollection}, \textsf{inproceedings},
\textsf{letter}, and \textsf{review}.

\begin{description}
\item[\qquad false:] This is the default.  If you use
  \textsf{crossref} or \textsf{xref} fields in the four mentioned
  entry types, you'll get the abbreviated references in both notes and
  bibliography.
\item[\qquad true:] You'll get no abbreviated references in these
  entry types, either in notes or in the bibliography.
\item[\qquad notes:] The abbreviated references will not appear in
  notes, but only in the bibliography.
\item[\qquad bib:] The abbreviated references will not appear in the
  bibliography, but only in notes.
\item[\qquad none:] This switch is special, allowing you with one
  setting to provide abbreviated references not just to the four entry
  types mentioned but also to \textsf{book}, \textsf{bookinbook},
  \textsf{collection}, and \textsf{proceedings} entries, both in notes
  and in the bibliography.
\end{description}

The \mymarginpar{\texttt{booklongxref}} second option,
\texttt{booklongxref}, controls the settings for \textsf{book},
\textsf{bookinbook}, \textsf{collection}, and \textsf{proceedings}
entries:

\begin{description}
\item[\qquad true:] This is the default.  If you use \textsf{crossref}
  or \textsf{xref} fields in these entry types, by default you will
  \emph{not} get any abbreviated references, either in notes or
  bibliography.
\item[\qquad false:] You'll get abbreviated references in these entry
  types both in notes and in the bibliography.
\item[\qquad notes:] The abbreviated references will not appear in
  notes, but only in the bibliography.
\item[\qquad bib:] The abbreviated references will not appear in the
  bibliography, but only in notes.
\end{description}

Please note that you can set both of these options either in the
preamble or in the \textsf{options} field of individual entries,
allowing you to change the settings on an entry-by-entry basis.

%\enlargethispage{2\baselineskip}

\mylittlespace Please further note that in earlier releases of
\textsf{biblatex-chicago} I recommended against using
\textsf{shorthand}, \textsf{reprinttitle} and/or \textsf{userf} fields
in combination with this abbreviated cross-referencing mechanism.  I
received, however, a request from Alexandre Roberts to allow the
shorthand to appear in the place of the abbreviated cross-reference as
an additional space-saving measure, and one from Kenneth Pearce to
permit the combination of the other two fields with \textsf{crossref},
as well.  The \textsf{userf} and \textsf{reprinttitle} fields should
just work automatically in such circumstances, but
\mymarginpar{\texttt{inheritshorthand}} the \textsf{shorthand} field
in parent entries needs to be enabled by setting the
\texttt{inheritshorthand} package option to \texttt{true}.  There are,
in addition, several other steps required to make this function
smoothly --- please see the documentation of the \textbf{shorthand}
field, below, for a full explanation.  (In case it isn't clear, the
combination of \textsf{userf}, \textsf{shorthand}, and
\textsf{crossref} functionality in a single entry is now possible.  If
you come across any problems or inaccuracies, please report them.)

\mybigspace I\mymarginpar{\textbf{date}} have now implemented all of
the applicable parts of \textsf{biblatex's} elegant, and long
standing, support for the \textsc{iso}8601-2 Extended Format
specification, which means the package now provides greatly enhanced
possibilities for presenting uncertain and unspecified dates and date
ranges, along with date eras, seasons, and time stamps.  I have also
implemented the \emph{Manual's} (9.64) guidelines for compressing year
ranges, as well as providing a few more extras to help with some of
the other tricky corners of the \emph{Manual's} instructions.  A
combination of \textsf{biblatex} and \textsf{biblatex-chicago} package
options allows you to define when, how, and where any of these
extended specifications will appear in your documents.  I have
attempted to provide as compliant a set of defaults as possible in
\textsf{biblatex-chicago.sty}, but you can alter any of them according
to your needs.  All are documented in section~\ref{sec:options},
below, but table~\ref{ad:date:extras}, located in the author-date
section, purports to serve as a convenient reference guide to how this
all works.

\mylittlespace There are several more general remarks about the
\textsf{date} field that may be helpful to users.  First, I highly
recommend familiarizing yourself with the extended date
specifications, as in many cases they will greatly simplify the
creation of your .bib databases.  The new \texttt{compressyears}
option (\texttt{true} by default), for example, takes a year range in
a date field and handles the somewhat tricky Chicago compression rules
for you, while also giving you a simple means of turning it off that
doesn't involve combing your .bib file for all the \textsf{year}
fields that contain your hand-formatted ranges.  Clearly, situations
may still arise when a specially-crafted \textsf{year} or
\textsf{origyear} field may be necessary, but if you can use the
enhanced specifications then I strongly advocate doing so.  Second,
the fine-grained specification of a time stamp is really only
necessary for news stories that are frequently updated \enquote{as
  they unfold} (14.191), for online sources that change rapidly enough
for a time stamp to be necessary (14.207, 14.233; wikiped:bibtex), or
for online posts, particularly comments, that may need a time stamp
for disambiguation (14.208--10).  If you wish to specify the time
zone, the \emph{Manual} (10.41) prefers initialisms like \enquote{EST}
or \enquote{PDT,} and these are most easily provided using the
\textsf{timezone} field, where you can include your own parentheses if
so desired (cp.\ 14.191).  For the \textsf{date} field itself, a time
stamp will only appear in \textsf{article}, \textsf{review},
\textsf{suppperiodical}, and \textsf{online} entries, the first three
only with a \texttt{magazine} \textsf{entrysubtype}.  All types can
print such a stamp from the \textsf{urldate} (controllable using the
new\mymarginpar{\texttt{urlstamp}} \texttt{urlstamp} option), while
only \textsf{review} and \textsf{suppperiodical} entries will print
this data from an \textsf{eventdate}.  If you find a context in which
a time stamp would be useful and which isn't included in this
discussion, please let me know.

\mylittlespace Third, an incomplete time specification will be ignored
by \textsf{biber}, so include the seconds in it, as in
table~\ref{ad:date:extras}, safe in the knowledge that they won't, by
default, ever appear in your documents.  Should you want that level of
discrimination, the \textsf{biblatex} option \texttt{seconds} set to
\texttt{true} provides it.  Fourth, in the \textsf{misc} entry type
the \textsf{date} field can help to distinguish between two classes of
archival material, letters and \enquote{letter-like} sources using
\textsf{origdate} while others (interviews, wills, contracts) use
\textsf{date}.  (See \textsf{misc} in section~\ref{sec:entrytypes} for
the details.)  Fifth, you can in most entry types qualify a
\textsf{date} with the \textsf{userd} field, assuming that the entry
contains no \textsf{urldate}.  For \textsf{music} and \textsf{video}
entries, there are several other requirements --- please see the
documentation of \textsf{userd}, below.

\enlargethispage{2\baselineskip}

\mylittlespace Sixth, and finally, please note that the
\textsf{nameaddon} field, which see, is no longer the place for time
stamps, as it was in the 16th-edition styles.  Any such data there
should be moved into the corresponding date field (either the
\textsf{date} or the \textsf{eventdate}, typically).  On all these
questions generally please cf.\ also \textsf{origdate},
\textsf{timezone}, and \textsf{year}, below; the \texttt{alldates},
\texttt{alltimes}, \texttt{alwaysrange}, \texttt{centuryrange},
\texttt{compressyears}, \texttt{datecirca}, \texttt{dateera},
\texttt{dateeraauto}, \texttt{dateuncertain}, \texttt{decaderange},
\texttt{nodatebrackets}, \texttt{nodates}, \texttt{noyearbrackets},
\texttt{timezones}, \texttt{urlstamp}, and \texttt{urltime} options in
sections~\ref{sec:presetopts}, \ref{sec:chicpreset}, and
\ref{sec:useropts}; and section~4.5.10 in \textsf{biblatex.pdf}

\mylittlespace (Users of the Chicago author-date style who wish to
minimize the labor needed to convert a .bib database for the notes \&\
bibliography style should be aware that the latter style includes
compatibility code for the \texttt{cmsdate} (silently ignored) and
\texttt{switchdates} options, along with the mechanism for reversing
\textsf{date} and \textsf{origdate}.  This means that you can, in
theory, leave all of this alone in your .bib file when making the
conversion, though I'm retaining the right to revoke this if the code
in question demonstrably interferes with the functioning of the notes
\&\ bibliography style.)

\mybigspace This \mymarginpar{\textbf{day}} field, as of
\textsf{biblatex} 0.9, is obsolete, and will be ignored if you use it
in your .bib files.  Use \textsf{date} instead.

\mybigspace Standard \mymarginpar{\textbf{doi}} \textsf{biblatex}
field, providing the Digital Object Identifier of the work.  The
\emph{Manual} specifies that, given their relative permanence compared
to URLs, \enquote{authors should prefer a DOI- or Handle-based URL
  whenever one is available} (14.8).  (14.175; friedman:learning).
Cf.\ \textsf{url}.

\mybigspace Standard \mymarginpar{\textbf{edition}} \textsf{biblatex}
field.  If you enter a plain cardinal number, \textsf{biblatex} will
convert it to an ordinal (chicago:manual), followed by the appropriate
string.  Any other sort of edition information will be printed as is,
though if your data begins with a word (or abbreviation) that would
ordinarily only be capitalized at the beginning of a sentence, then
simply ensure that that word (or abbreviation) is in lowercase, and
\textsf{biblatex-chicago-notes} will automatically do the right thing
(babb:peru, times:guide).  In most situations, the \emph{Manual}
generally recommends the use of abbreviations in both bibliography and
notes, but there is room for the user's discretion in specific
citations (emerson:nature).

\mybigspace As \mymarginpar{\textbf{editor}} far as possible, I have
implemented this field as \textsf{biblatex}'s standard styles do, but
the requirements specified by the \emph{Manual} present certain
complications that need explaining.  \textsf{Biblatex.pdf} points out
that the \textsf{editor} field will be associated with a
\textsf{title}, a \textsf{booktitle}, or a \textsf{maintitle},
depending on the sort of entry.  More specifically,
\textsf{biblatex-chicago} associates the \textsf{editor} with the most
comprehensive of those titles, that is, \textsf{maintitle} if there is
one, otherwise \textsf{booktitle}, otherwise \textsf{title}, if the
other two are lacking.  In a large number of cases, this is exactly
the correct behavior (adorno:benj, centinel:letters,
plato:republic:gr, among others).  Predictably, however, there are
numerous cases that require, for example, an additional editor for one
part of a collection or for one volume of a multi-volume work.  For
these cases I have provided the \textsf{namea} field.  You should
format names for this field as you would for \textsf{author} or
\textsf{editor}, and these names will always be associated with the
\textsf{title} (donne:var).

\mylittlespace As you will see below, I have also provided a
\textsf{nameb} field, which holds the translator of a given
\textsf{title} (euripides:orestes).  If \textsf{namea} and
\textsf{nameb} are the same, \textsf{biblatex-chicago} will
concatenate them, just as \textsf{biblatex} already does for
\textsf{editor}, \textsf{translator}, and \textsf{namec} (i.e., the
compiler).  Furthermore, it is conceivable that a given entry will
need separate editors for each of the three sorts of title.  For this,
and for various other tricky situations, there is the \cmd{partedit}
macro (and its siblings), designed to be used in a \textsf{note} field
or in one of the \textsf{titleaddon} fields (chaucer:liferecords).
(Because the strings identifying an editor differ in notes and
bibliography, one can't simply write them out in such a field, hence
the need for a macro, which I discuss further in the commands section
below [\ref{sec:formatcommands}].)  Please note that, when attempting
to find a name for the head of a note or a bibliography entry,
\textsf{namea} takes precedence over \textsf{editor}, and
\textsf{nameb} over \textsf{translator}.  Cf.\ \textsf{namea},
\textsf{nameb}, \textsf{namec}, and \textsf{translator}.

\mybigspace The \mymarginpar{\textbf{editora\\editorb\\editorc}} newer
releases of \textsf{biblatex} provide these fields as a means to
specify additional contributors to texts in a number of editorial
roles.  In the Chicago styles they seem most relevant for the
audiovisual types, especially \textsf{music} and \textsf{video}, and
now also the \textsf{performance} type, in all of which they can help
to identify conductors, directors, producers, and performers.  To
specify the role, use the fields \textsf{editoratype},
\textsf{editorbtype}, and \textsf{editorctype}, which see.  (Cf.\
bernstein:shostakovich, hamilton:miranda, handel:messiah.)

\mybigspace Normally, \mymarginpar{\textbf{editortype}} with the
exception of the \textsf{article} and \textsf{review} types,
\textsf{biblatex-chicago-notes} will automatically find a name to put
at the head of an entry, starting with an \textsf{author}, and
proceeding in order through \textsf{namea}, \textsf{editor},
\textsf{nameb}, \textsf{translator}, and \textsf{namec} (the
compiler).  If all six are missing, then the \textsf{title} will be
placed at the head.  (In \textsf{article} and \textsf{review} entries
with a \texttt{magazine} \textsf{entrysubtype}, a missing author
immediately prompts the use of \textsf{journaltitle} at the head of an
entry.  See above under \textsf{article} for details.)  The
\textsf{editortype} field provides even greater flexibility, giving
you the ability to indicate any number of roles at the head of an
entry.  You can do this even though an author is named (eliot:pound
shows this mechanism in action for a standard editor, rather than for
an alternative role).  Two things are necessary for this to happen.
First, in the \textsf{options} field you need to set
\texttt{useauthor=false}, then you need to put the name you wish to
see at the head of your entry into the \textsf{editor} or the
\textsf{namea} field.  If the \enquote{editor} is in fact a compiler,
then you need to put \texttt{compiler} into the \textsf{editortype}
field, and \textsf{biblatex} will print the correct string after the
name in both the bibliography and in the long note form.

\mylittlespace In previous releases of \textsf{biblatex-chicago} you
could only use defined \cmd{bibstrings} in this field, at least if you
wanted anything printed.  N.~Andrew Walsh pointed out that the
standard \textsf{biblatex} styles will just print the field as-is in
this case, allowing them to handle a great many unforeseen editorial
roles with comparative ease, so I've implemented this, too, making
sure to capitalize the string if the context demands it.  The string
you choose will differ depending on whether it will be printed after a
name at the head of an entry or before a name later on in the entry,
e.g., \enquote{cartographer} or \enquote{maps created by.}  A bit of
trial and error should see you through.

\mylittlespace There are a few more details of which you need to be
aware.  Because \textsf{biblatex-chicago} has added the \textsf{namea}
field, which gives you the ability to identify the editor specifically
of a \textsf{title} as opposed to a \textsf{maintitle} or a
\textsf{booktitle}, the name-finding algorithm checks first to see
whether a \textsf{namea} is defined.  If it is, that name will be used
at the head of the entry, if it isn't, or if you've set the option
\texttt{usenamea=false}, the algorithm will go ahead and look for an
\textsf{editor}.  The \textsf{editortype} field applies only to the
\textsf{editor}, but you can use \textsf{nameatype} to modify
\textsf{namea}.  Either of these names should be sorted properly in
the bibliography, but please be aware that if you want a shortened
form to appear in short notes then there's only the
\textsf{shorteditor}, which you should ensure presents whichever of
the two editors' names appears at the head of long notes or
bibliography entries.

\mylittlespace In \textsf{biblatex} 0.9 Lehman reworked the string
concatenation mechanism, for reasons he outlined in his RELEASE file,
and I have followed his lead.  In short, if you define the
\textsf{editortype} field, then concatenation is turned off, even if
the name of the \textsf{editor} matches, for example, that of the
\textsf{translator}.  In the absence of an \textsf{editortype} (or
\textsf{nameatype}), the usual mechanisms remain in place, that is, if
the \textsf{editor} exactly matches a \textsf{translator} and/or a
\textsf{namec}, or alternatively if \textsf{namea} exactly matches a
\textsf{nameb} and/or a \textsf{namec}, then \textsf{biblatex} will
print the appropriate strings.  The \emph{Manual} specifically (14.32)
recommends not using these identifying strings in the short note form,
and \textsf{biblatex-chicago-notes} follows their recommendation.  If
you nevertheless need to provide such a string, you'll have to do it
manually in the \textsf{shorteditor} field, or perhaps, in a different
sort of entry, in a \textsf{shortauthor} field.

\mylittlespace It may also be worth noting that because of certain
requirements in the specification -- absence of an \textsf{author},
for example -- the \texttt{useauthor=false} mechanism is either
unnecessary or won't work properly in the following entry types:
\textsf{collection}, \textsf{letter}, \textsf{patent},
\textsf{periodical}, \textsf{proceedings}, \textsf{review},
\textsf{suppbook}, \textsf{suppcollection}, and
\textsf{suppperiodical}.

\mybigspace These
\mymarginpar{\textbf{editoratype\\editorbtype\\editorctype}} fields
identify the exact role of the person named in the corresponding
\textsf{editor[a-c]} field, just as \textsf{editortype} (q.v.) does
for the \textsf{editor}.  Note that they are not part of the string
concatenation mechanism.  I have implemented them just as the standard
styles do, that is, if the field isn't a pre-defined \cmd{bibstring}
it will be printed as-is, contextually capitalized.  They have found a
use particularly in \textsf{music}, \textsf{performance}, and
\textsf{video} entries.  Cf.\ bernstein:shostakovich,
hamilton:miranda, handel:messiah.

\mybigspace Standard \mymarginpar{\textbf{eid}} \textsf{biblatex}
field, providing a string or number some journals use uniquely to
identify a particular article.  Only applicable to the
\textsf{article} entry type, and only to those without a
\texttt{magazine} \textsf{entrysubtype}.  The 17th edition of the
\emph{Manual} now specifies where to print this (14.174), and I have
moved it in accordance with its specifications.  It replaces the
\textsf{pages} field in long notes and bibliography, and appears after
any specific page cited in the \textsf{postnote} field of a long note.

\paragraph*{\protect\mymarginpar{\textbf{entrysubtype}}}
\label{sec:entrysub}

Standard and very powerful \textsf{biblatex} field, left undefined by
the standard styles.  In \textsf{bib\-latex-chicago-notes} it has eight
very specific uses, the first three of which I have designed in order
to maintain, as much as possible, backward compatibility with the
standard styles.  First, in \textsf{article}, \textsf{periodical}, and
\textsf{review} entries, the field allows you to differentiate between
scholarly \enquote{journals,} on the one hand, and \enquote{magazines}
and \enquote{newspapers} on the other.  Usage is fairly simple: you
need to put the exact string \texttt{magazine} into the
\textsf{entrysubtype} field if you are citing one of the latter two
types of source, whereas if your source is a \enquote{journal,} then
you need do nothing.

\mylittlespace The second use involves references to works from
classical antiquity and, according to the \emph{Manual}, from the
Middle Ages, as well.  When you cite such a work using the traditional
divisions into books, sections, lines, etc., divisions which are
presumed to be the same across all editions, then you need to put the
exact string \texttt{classical} into the \textsf{entrysubtype} field.
This has no effect in long notes or in the bibliography, but it does
affect the formatting of short notes, where it suppresses some of the
punctuation.  Ordinarily, you will use this toggle in a \textsf{book}
or a \textsf{bookinbook} entry, but it is possible that a journal
might well also present an edition of such a work.  Given the
tradition of using italics for the titles of such works, this may
require using a \textsf{titleaddon} field (with hand formatting)
instead of a \textsf{title}.  If you wish to reference a classical or
medieval work by the page numbers of a particular, non-standard
edition, then you shouldn't use the \textsf{entrysubtype} toggle.
Also, and the specification is reasonably clear about this, works from
the Renaissance and later, even if cited by the traditional divisions,
have short notes formatted normally, and therefore don't need an
\textsf{entrysubtype} field.  (See \emph{Manual} 14.242--54;
aristotle:metaphy:gr, herodotus:wilson, plato:republic:gr;
euripides:orestes is an example of a translation cited by page number
in a modern edition.  Cf.\ also the \mycolor{\texttt{notitle}} option
in section~\ref{sec:useropts}.)

\mylittlespace The third use occurs in \textsf{misc} entries.  If such
an entry contains no \textsf{entrysubtype} field, then the citation
will be treated just as the standard \textsf{biblatex} styles would,
including the use of italics for the \textsf{title}.  Any string at
all in \textsf{entrysubtype} tells \textsf{biblatex-chicago-notes} to
treat the source as part of an unpublished archive.  A \textsf{misc}
entry with an \textsf{entrysubtype} defined is the least formatted of
all those specified by the \emph{Manual} --- see
section~\ref{sec:entrytypes} above under \textbf{misc} for all the
details on how these citations work.

\mylittlespace Fourth, the field can be defined in the
\textsf{artwork} entry type in order to refer to a work from antiquity
whose title you do not wish to be italicized.  Please see the
documentation of \textbf{artwork} above for the details.  Fifth, you
can define it in a \textbf{standard} entry, q.v., to change the
appearance of both long and short notes.  Sixth, you can define it in
an \textbf{audio}, \textbf{music}, or \textbf{video} entry if such an
entry refers to an individual unit that isn't part of any larger
collection, the entry therefore having only a \textsf{title} and not a
\textsf{booktitle}, a \textsf{title} that \textsf{biblatex-chicago}
would normally interpret as the title of a larger unit (and therefore
italicize).  Seventh, and sticking with the \textbf{video} type,
though enacting quite a different syntactic transformation, the 17th
edition (14.265) now recommends that, when presenting episodes from a
TV series, the name of the series (\textsf{booktitle}) comes before
the episode name (\textsf{title}).  The exact string
\texttt{tvepisode} in the \textsf{entrysubtype} field achieves this
reversal, which includes using the \textsf{booktitle} as a
\textsf{sorttitle} in the bibliography and also as the
\textsf{labeltitle} in short notes.

\mylittlespace Eighth, and finally, you can use any
\textsf{entrysubtype} whatever in \textsf{inreference} entries in
order to treat them as inherently online works rather than standard
published works.  See the documentation of \textbf{online} and
\textbf{inreference} entries in section~\ref{sec:entrytypes}, above,
and also 14.233 and wikiped:bibtex.

\mybigspace Kazuo
\mymarginpar{\textbf{eprint}\\\textbf{eprintclass}\\\textbf{eprinttype}}
Teramoto suggested adding \textsf{biblatex's} excellent
\textsf{eprint} handling to \textsf{biblatex-chica\-go}, and he sent me
a patch implementing it.  I have applied it, with minor alterations,
so these three fields now work more or less as they do in standard
\textsf{biblatex}.  They may prove helpful in providing more
abbreviated references to online content than conventional URLs,
though I can find no specific reference to them in the \emph{Manual}.

% %\enlargethispage{-\baselineskip}

\mybigspace This \mymarginpar{\textbf{eventdate}} is a standard
\textsf{biblatex} field which has gradually accumulated functions in
\textsf{biblatex-chicago}.  It can now play a role in
\textsf{artwork}, \textsf{audio}, \textsf{image},
\textsf{inproceedings}, \textsf{music}, \textsf{proceedings},
\textsf{review}, \textsf{standard}, \textsf{suppperiodical},
\textsf{unpublished}, and \textsf{video} entries.  In \textsf{artwork}
and \textsf{image} entries it identifies the publication date of, most
frequently, a photograph, in association with the
\textsf{howpublished} field which identifies the periodical or other
medium in which it was published (mccurry:afghangirl).  In
\textsf{standard} entries it will also usually be associated with a
\textsf{howpublished} field, allowing you to specify a later renewal
or reaffirmation of a standard (niso\hc bibref).  In \textsf{audio}
entries, it specifies the release date of a single episode of a
podcast (danforth:podcast).  In \textsf{music} entries, it identifies
the recording or performance date of a particular song (rather than of
a whole disc, for which you would use \textsf{origdate}), whereas in
\textsf{video} entries it identifies either the original broadcast
date of a particular episode of a TV series or the date of a filmed
musical performance.  In both these cases \textsf{biblatex-chicago}
will automatically prepend a bibstring --- \texttt{recorded} and
\texttt{aired}, respectively --- to the date, but you can change this
string using the \textsf{userd} field, something you'll definitely
want to do for filmed musical performances (friends:leia,
handel:messiah, holiday:fool).

\mylittlespace In \textsf{inproceedings}, \textsf{proceedings}, and
\textsf{unpublished} entries it identifies the date of an event at
which a published or unpublished work was presented, though in truth
the \textsf{date} will do as well in \textsf{unpublished} entries
(nass:address).  The field's use in \textsf{review} entries, finally,
includes a possible time stamp.  In this context, an
\textsf{eventdate} helps to identify a particular comment on, or reply
to another comment on, a blog post.  Given that many such posts by a
single \textsf{author} could appear on the same day, you can
distinguish them by putting a time specification in the
\textsf{eventdate} field itself (ac:comment).  Please see the
\textbf{review} type, above, for the details of how to cite these
materials, possibly with the help of the \texttt{commenton}
\textsf{relatedtype}.  See also the \textsf{date} field docs, in
particular table~\ref{ad:date:extras} (located in the author-date
section), for details on how the \textsc{iso}8601-2 Extended Format
specifications offered by \textsf{biblatex}, including time stamps and
much else besides, have been implemented in \textsf{biblatex-chicago}.

\mybigspace This \mymarginpar{\textbf{eventtimezone}} field can, if
necessary, specify the time zone associated with a time stamp given as
part of an \textsf{eventdate}.  The \emph{Manual} prefers initialisms
like \enquote{EST} for this purpose, and you can provide parentheses
around it at your discretion (cp.\ 10.41 and 14.191).

\mybigspace A \mymarginpar{\textbf{eventtitle}} standard
\textsf{biblatex} field for identifying the name of the event that
produces either a published record (\textsf{inproceedings} and
\textsf{proceedings} entries) or an unpublished one
(\textsf{unpublished}).

\mybigspace Standard \mymarginpar{\textbf{eventtitleaddon}}
\textsf{biblatex} field for adding information about an
\textsf{eventtitle}, and available in the same entry types as that
field.

\mybigspace As \mymarginpar{\textbf{foreword}} with the
\textsf{afterword} field above, \textsf{foreword} will in general
function as it does in standard \textsf{biblatex}.  Like
\textsf{afterword} (and \textsf{introduction}), however, it has a
special meaning in a \textsf{suppbook} entry, where you simply need to
define it somehow (and leave \textsf{afterword} and
\textsf{introduction} undefined) to make a foreword the focus of a
citation.

\mybigspace A \mymarginpar{\textbf{holder}} standard \textsf{biblatex}
field for identifying a \textsf{patent}'s holder(s), if they differ
from the \textsf{author}.  The \emph{Manual} has nothing to say on the
subject, but \textsf{biblatex-chicago-notes} prints it (them), in
parentheses, just after the author(s).

\mybigspace Standard \mymarginpar{\textbf{howpublished}}
\textsf{biblatex} field which, like the \textsf{eventdate} field, is
gradually accumulating functions in \textsf{biblatex-chicago}.  In the
\textsf{booklet} type it retains something of its traditional usage,
replacing the \textsf{publisher}, and has a similar (somewhat
paradoxical) place in \textsf{unpublished} entries.  In the
\textsf{misc} and \textsf{performance} types it works almost as a
second \textsf{note} field, bringing in extra information about a work
in close association with the \textsf{type} and \textsf{version}
fields, while in \textsf{dataset} entries its information will be
associated with both those fields and also with the \textsf{number}
field.  17th-edition \textsf{music} entries require a field to provide
the medium of downloaded music and/or the name of the streaming
service, so \textsf{howpublished} works there as an online double of
\textsf{type} and of \textsf{publisher}.  Finally, in
\textsf{artwork}, \textsf{image}, and \textsf{standard} entries it
serves to qualify or modify an \textsf{eventdate}, almost as a
\textsf{userd} field modifies a \textsf{date} or \textsf{urldate}.
Please see the docs of those entry types for more information, and
also bedford:photo, clark:mesopot, mccurry:afghangirl, niso:bibref,
rihanna:umbrella.

\mybigspace Standard \mymarginpar{\textbf{institution}}
\textsf{biblatex} field.  In the \textsf{thesis} entry type, it will
usually identify the university for which the thesis was written,
while in a \textsf{report} entry it may identify any sort of
institution issuing the report.

% %\enlargethispage{\baselineskip}

\mybigspace As \mymarginpar{\textbf{introduction}} with the
\textsf{afterword} and \textsf{foreword} fields above,
\textsf{introduction} will in general function as it does in standard
\textsf{biblatex}.  Like those fields, however, it has a special
meaning in a \textsf{suppbook} entry, where you simply need to define
it somehow (and leave \textsf{afterword} and \textsf{foreword}
undefined) to make an introduction the focus of a citation.

\mybigspace Standard \mymarginpar{\textbf{isbn}} \textsf{biblatex}
field, for providing the International Standard Book Number of a
publication.  Not typically required by the \emph{Manual}.

\mybigspace Standard \mymarginpar{\textbf{isrn}} \textsf{biblatex}
field, for providing the International Standard Technical Report
Number of a report.  Only relevant to the \textsf{report} entry type,
and not typically required by the \emph{Manual}.

\mybigspace Standard \mymarginpar{\textbf{issn}} \textsf{biblatex}
field, for providing the International Standard Serial Number of a
periodical in an \textsf{article} or a \textsf{periodical} entry.  Not
typically required by the \emph{Manual}.

\mybigspace Standard \mymarginpar{\textbf{issue}} \textsf{biblatex}
field, designed for \textsf{article}, \textsf{periodical}, or
\textsf{review} entries identified by something like \enquote{Spring}
or \enquote{Summer} rather than by the usual \textsf{month} or
\textsf{number} fields (brown:bremer).  \textsf{Biblatex's} enhanced
date handling allows you to specify a season in the \textsf{date}
field, with the \enquote{months} 21--24 used for Spring, Summer,
Autumn, and Winter, respectively.  Cf.\ table~\ref{ad:date:extras},
below.

\mybigspace The \mymarginpar{\textbf{issuesubtitle}} subtitle for an
\textsf{issuetitle} --- see next entry.

\mybigspace Standard \mymarginpar{\textbf{issuetitle}}
\textsf{biblatex} field, intended to contain the title of a special
issue of any sort of periodical.  If the reference is to one article
within the special issue, then this field should be used in an
\textsf{article} entry (conley:fifthgrade), whereas if you are citing
the entire issue as a whole, then it would go in a \textsf{periodical}
entry, instead (good:wholeissue).  The \textsf{note} field is the
proper place to identify the type of issue, e.g.,\ \texttt{special
  issue}, with the initial letter lower-cased to enable automatic
contextual capitalization.

\mybigspace The \mymarginpar{\textbf{journalsubtitle}} subtitle for a
\textsf{journaltitle} --- see next entry.

\mybigspace Standard \mymarginpar{\textbf{journaltitle}}
\textsf{biblatex} field, replacing the standard \textsc{Bib}\TeX\
field \textsf{journal}, which, however, still works as an alias.  It
contains the name of any sort of periodical publication, and is found
in the \textsf{article} and \textsf{review} entry types.  In the case
where a piece in an \textsf{article} or \textsf{review}
(\textsf{entrysubtype} \texttt{magazine}) doesn't have an author,
\textsf{biblatex-chicago-notes} provides for this field to be used as
the author.  See above (section~\ref{sec:entrytypes}) under
\textbf{article} for details.  The lakeforester:pushcarts and
nyt:trevorobit entries in \textsf{notes-test.bib} will give you some
idea of how this works.  Please note there is a \textsf{shortjournal}
field which you can use to abbreviate the \textsf{journaltitle} in
notes and/or in the bibliography, and you can also use it to print a
list of journal abbreviations.  Cf.\ the \textsf{shortjournal}
documentation below.

\mybigspace An \colmarginpar{\textbf{journaltitleaddon}} annex to the
\textsf{journaltitle}, for which see previous entry.  Such an annex
would be printed in the main text font.  If your data begins with a
word that would ordinarily only be capitalized at the beginning of a
sentence, then simply ensure that that word is in lowercase, and
\textsf{biblatex-chicago-notes} will automatically do the right thing.
The package and entry option \mycolor{\texttt{jtitleaddon}}
(section~\ref{sec:chicpreset}) allows you to customize the punctuation
that appears before the \textsf{journaltitleaddon} field (hua:cms).
The default is a \texttt{space}.

\mybigspace This \mymarginpar{\textbf{keywords}} field is
\textsf{biblatex}'s powerful and flexible technique for filtering
bibliography entries, allowing you to subdivide a bibliography
according to just about any criteria you care to invent, or indeed to
prevent entries in notes from appearing in the bibliography, as the
\emph{Manual} sometimes recommends.  See \textsf{biblatex.pdf} (3.7)
for thorough documentation.

\mybigspace A \mymarginpar{\textbf{language}} standard
\textsf{biblatex} field, designed to allow you to specify the
language(s) in which a work is written.  As a general rule, the
Chicago style doesn't require you to provide this information, though
it may well be useful for clarifying the nature of certain works, such
as bilingual editions, for example.  There is at least one situation,
however, when the \emph{Manual} does specify this data, and that is
when the title of a work is given in translation, even though no
translation of the work has been published, something that might
happen when a title is in a language deemed to be unparseable by a
majority of your expected readership (14.99; pirumova,
rozner:liberation).  In such a case, you should provide the
language(s) involved using this field, connecting multiple languages
using the keyword \texttt{and}.  (I have retained \textsf{biblatex's}
\cmd{bibstring} mechanism here, which means that you can use the
standard bibstrings or, if one doesn't exist for the language you
need, just give the name of the language, capitalized as it should
appear in your text.  You can also mix these two modes inside one
entry without apparent harm.)

\mylittlespace An alternative arrangement suggested by the
\emph{Manual} is to retain the original title of a piece but then to
provide its translation, as well.  If you choose this option, you'll
need to make use of the \textbf{usere} field, on which see below.  In
effect, you'll probably only ever need to use one of these two fields
in any given entry, and in fact \textsf{biblatex-chicago-notes} will
only print one of them if both are present, preferring \textsf{usere}
over \textsf{language} for this purpose (see kern and weresz).  Note
also that both of these fields are universally associated with the
\textsf{title} of a work, rather than with a \textsf{booktitle} or a
\textsf{maintitle}.  If you need to attach a language or a translation
to either of the latter two, you could probably manage it with special
formatting inside those fields themselves.

\mybigspace I \mymarginpar{\textbf{lista}} intend this field
specifically for presenting citations from reference works that are
arranged alphabetically, where the name of the item rather than a page
or volume number should be given.  The field is a \textsf{biblatex}
list, which means you should separate multiple items with the keyword
\texttt{and}.  Each item receives its own set of quotation marks, and
the whole list will be prefixed by the appropriate string
(\enquote{s.v.,} \emph{sub verbo}, pl.\ \enquote{s.vv.}).
\textsf{Biblatex-chicago-notes} will only print such a field in a
\textsf{book} or an \textsf{inreference} entry, and you should look at
the documentation of these entry types for further details.  (See
\emph{Manual} 14.232--33; ency:britannica, grove:sibelius,
times:guide, wikiped:bibtex.)

\mybigspace This \mymarginpar{\textbf{location}} is
\textsf{biblatex}'s version of the usual \textsc{Bib}\TeX\ field
\textsf{address}, though the latter is accepted as an alias if that
simplifies the modification of older .bib files.  According to the
\emph{Manual} (14.129), a citation usually need only provide the first
city listed on any title page, though a list of cities separated by
the keyword \enquote{\texttt{and}} will be formatted appropriately.
If the place of publication is unknown, you can use
\verb+\autocap{n}.p.+\ instead (14.132).  For all cities, you should
use the common English version of the name, if such exists (14.131).

\mylittlespace Three more details need explanation here.  In
\textsf{article}, \textsf{periodical}, and \textsf{review} entries,
there is usually no need for a \textsf{location} field, but
\enquote{if a journal might be confused with another with a similar
  title, or if it might not be known to the users of a bibliography,}
then this field can present the place or institution where it is
published (14.182, 14.191, 14.193--94; lakeforester:pushcarts,
kimluu:diethyl, and garrett).  For blogs cited using \textsf{article}
entries, this is a good place to identify the nature of the source ---
i.e., the word \enquote{blog} --- letting the style automatically
provide the parentheses (14.208; ellis:blog).  Less predictably, it is
in the vicinity of the \textsf{location} that the \emph{Manual}
indicates that a particular book is a reprint edition (14.114), so in
such a case you can use the \textsf{biblatex-chicago} macro
\cmd{reprint}, followed by a comma, a space, and the location.
Somewhat more cleanly and simply, and more in keeping with standard
\textsf{biblatex} usage, you can just put the string \texttt{reprint}
into the \textsf{pubstate} field to achieve the same result.  See the
\textsf{pubstate} documentation below (aristotle:metaphy:gr,
schweitzer:bach).  The \textsf{origdate} field may be used to give the
original date of publication, and of course more complicated
situations should usually be amenable to inclusion in the
\textsf{note} field (emerson:nature).

\mybigspace The \mymarginpar{\textbf{mainsubtitle}} subtitle for a
\textsf{maintitle} --- see next entry.

\mybigspace The \mymarginpar{\textbf{maintitle}} main title for a
multi-volume work, e.g., \enquote{Opera} or \enquote{Collected Works.}
(See donne\hc var, euripides:orestes, harley:cartography, lach:asia,
pelikan:christian, and plato\hc republic:gr.)  When using a
\textsf{crossref} field and \textsf{Biber}, the \textsf{title} of
\textbf{mv*} entry types always becomes a \textsf{maintitle} in the
child entry.  See also the documentation of the \texttt{maintitle}
\textsf{relatedtype} in the \textbf{mvbook} docs in
section~\ref{sec:entrytypes}, above, and in section~\ref{sec:related},
below.

\mylittlespace Because the 17th edition of the \emph{Manual}
recommends that you present not only the names of blogs but also the
names of their parent (usually periodical) publications, I have added
this field to \textsf{article}, \textsf{periodical}, and
\textsf{review} entries for just this purpose.  See the documentation
of those entry types in section~\ref{sec:entrytypes}, above, and also
table~\ref{tab:online:types} (14.208; amlen:hoot).

% %\enlargethispage{\baselineskip}

\mybigspace An \mymarginpar{\textbf{maintitleaddon}} annex to the
\textsf{maintitle}, for which see previous entry.  Such an annex would
be printed in the main text font.  If your data begins with a word
that would ordinarily only be capitalized at the beginning of a
sentence, then simply ensure that that word is in lowercase, and
\textsf{biblatex-chicago-notes} will automatically do the right thing.
The package and entry options \texttt{ptitleaddon} and
\texttt{ctitleaddon} (section~\ref{sec:chicpreset}) allow you to
customize the punctuation that appears before the
\textsf{maintitleaddon} field (schubert:muellerin).

\mybigspace Standard \mymarginpar{\textbf{month}} \textsf{biblatex}
field, containing the month of publication.  This should be an
integer, i.e., \texttt{month=\{3\}} not \texttt{month=\{March\}}.  See
\textsf{date} for more information.

\mybigspace This \mymarginpar{\textbf{namea}} is one of the fields
\textsf{biblatex} provides for style writers to use, but which it
leaves undefined itself.  In \textsf{biblatex-chicago} it contains the
name(s) of the editor(s) of a \textsf{title}, if the entry has a
\textsf{booktitle} and/or a \textsf{maintitle}, in which situation the
\textsf{editor} would be associated with one of these latter fields
(donne:var).  (In \textsf{article} and \textsf{review} entries,
\textsf{namea} applies to the \textsf{title} instead of the
\textsf{issuetitle}, should the latter be present.)  You should
present names in the field exactly as you would those in an
\textsf{author} or \textsf{editor} field, and the package will
concatenate this field with \textsf{nameb} if they are identical.
When choosing a name to head a note or a bibliography entry,
\textsf{biblatex-chicago} gives precedence to \textsf{namea} over
\textsf{editor}.  See under \textbf{editor} above for the full
details.  Please note that, as the field is highly single-entry
specific, if you are using \textsf{Biber} \textsf{namea} isn't
inherited from a \textsf{crossref}'ed parent entry.  Please note,
also, that you can use the \textsf{nameatype} field to redefine this
role just as you can with \textsf{editortype}, which see.  Cf.\ also
\textsf{nameb}, \textsf{namec}, \textsf{translator}, and the macros
\cmd{partedit},\,\cmd{parttrans},\,\cmd{parteditandtrans},
\cmd{partcomp},\,\cmd{parteditandcomp}, \cmd{parttransandcomp}, and
\cmd{partedittransandcomp}, for which see
section~\ref{sec:formatcommands}.

\mybigspace This \mymarginpar{\textbf{nameaddon}} field is provided by
\textsf{biblatex}, though not used by the standard styles.  In
\textsf{biblatex-chicago} its primary use, in most entry types, has
always been to specify that an author's name is a pseudonym, or to
provide either the real name or the pseudonym itself, if the other is
being provided in the \textsf{author} field.  The abbreviation
\enquote{\texttt{pseud}.}\ (always lowercase in English) is specified,
either on its own or after the pseudonym (centinel:letters,
creasey:ashe:blast, creasey:morton:hide, creasey:york:death, and
lecarre:quest); remember that \verb+\bibstring{pseudonym}+ does the
work for you.  See under \textbf{author} above for the full details.

\mylittlespace The field has slowly accumulated other functions, so
when Philipp Immel made a feature request, and pointed to a discussion
on Stack Exchange which suggested a few more, I thought I might
generalize the field's functionality, providing three package options
to allow users to mould it to their needs.  Before discussing these,
allow me to emphasize that the package defaults remain exactly the
same as before, so that, absent any of the new options, the style
still provides square brackets around the \textsf{nameaddon} in most
entry types, no brackets of any sort in \textsf{online, review,} and
\textsf{suppperiodical} entries, as well as in \textsf{misc} entries
with an \textsf{entrysubtype}, and rather specialized handling in
\textsf{customc} entries (which ignore all but the first of the new
options --- see below).  If you're happy with the status quo, then no
changes to your documents or\ .bib databases are necessary.

\mylittlespace If you do need or want to put the field to a different
use, the following options may help.  All of them are available
globally, per type, and per entry.  The first
\mymarginpar{\texttt{nameaddon}} new option is simply called
\texttt{nameaddon}, and determines where and when the field will be
printed at all.  There are seven possible values:
\begin{description}
\setlength{\parskip}{-4pt}
\item[\qquad \texttt{all}:] This is the default; if an entry has a
  \textsf{nameaddon}, it will appear in both long notes and in the
  bibliography.
\item[\qquad \texttt{none}:] The field will appear neither in the
  bibliography nor in long notes.
\item[\qquad \texttt{bib}:] The field will appear only in the
  bibliography.
\item[\qquad \texttt{cite}:] The field will appear only in long notes.
\item[\qquad \texttt{first}:] (This key and the next two are only
  available as global options.)  Philipp Immel requested this as a way
  to provide an \textsf{author's} dates in the \textsf{nameaddon}
  field and only have them printed the first time that author appears
  in the bibliography.  A sequence of consecutive long notes citing
  works by the same author will be treated the same way.  The code
  tests for identical \textsf{nameaddon} fields in works by identical
  \textsf{authors}, so other sorts of \textsf{nameaddon} will be
  printed as usual.
\item[\qquad \texttt{bibfirst}:] Like \texttt{first}, but will not
  print the \textsf{nameaddon} field in long notes.
\item[\qquad \texttt{citefirst}:] Like \texttt{first}, but will not
  print the \textsf{nameaddon} field in the bibliography.
\end{description}

The \mymarginpar{\texttt{nameaddonsep}} \texttt{nameaddonsep} option
controls the punctuation that appears before the
\textsf{nameaddon}. It takes the following six keys:

\begin{description}
\setlength{\parskip}{-4pt}
\item[\qquad \texttt{space}] = \cmd{addspace}.  This is the default.
\item[\qquad \texttt{none}] = no separator at all.  It presumes that
  you will include one in the \textsf{nameaddon} field itself.
\item[\qquad \texttt{colon}] = \verb+\addcolon\addspace+.
\item[\qquad \texttt{comma}] = \verb+\addcomma\addspace+.
\item[\qquad \texttt{period}] = \verb+\addperiod\addspace+.
\item[\qquad \texttt{semicolon}] = \verb+\addsemicolon\addspace+.
\end{description}

The \mymarginpar{\texttt{nameaddon-\\format}} \texttt{nameaddonformat}
option allows you to change the format of the \textsf{nameaddon} field
on the fly, so its value should be a field format that
\textsf{biblatex} understands.  This includes standard formats like
\texttt{parens,\,brackets} or \texttt{emph}, and also custom formats
that you provide in your preamble using \cmd{DeclareFieldFormat}, in
case the standard ones aren't adequate.  If you don't define this
option, then the usual defaults apply, as delineated above, and you
can use your own parentheses in \textsf{online, review,} and
\textsf{suppperiodical} entries, as well as in \textsf{misc} entries
with an \textsf{entrysubtype}, to distinguish screen names or other
authorial information from traditional pseudonyms (in brackets).

\mylittlespace Finally, two more details.  If you are using the
17th-edition styles for the first time, please note that the 16th
edition of the \emph{Manual} recommended specifying comments to blogs
and other online content using a time stamp in parentheses after the
\textsf{author}, but the 17th edition handles time stamps both
differently and more widely, so in this case you would now put time
data into the \textsf{date} or \textsf{eventdate} field, particularly
when the date itself is too coarse a specification to identify a
comment unambiguously (cf.\ ac:comment, obrien:recycle).  In the
\textsf{customc} entry type, finally, which is used to create
alphabetized cross-references to other bibliography entries, the
\textsf{nameaddon} field allows you to change the default string
linking the two parts of the cross-reference.  The code automatically
tests for a known bibstring, which it will italicize.  Otherwise, it
prints the string as you've provided in the \textsf{nameaddon} field
itself.  The punctuation is fixed.

\mybigspace You \mymarginpar{\textbf{nameatype}} can use this field
to change the role of a \textsf{namea} just as you can use
\textsf{editortype} to change the role of an \textsf{editor}.  As with
the \textsf{editortype}, using this field prevents string
concatenation with identical \textsf{nameb} or \textsf{namec} fields.
Please see \textbf{editortype}, above, for the details.

\mybigspace Like \mymarginpar{\textbf{nameb}} \textsf{namea}, above,
this is a field left undefined by the standard \textsf{biblatex}
styles.  In \textsf{biblatex-chicago}, it contains the name(s) of the
translator(s) of a \textsf{title}, if the entry has a
\textsf{booktitle} or \textsf{maintitle}, or both, in which situation
the \textsf{translator} would be associated with one of these latter
fields (euripides:orestes).  (In \textsf{article} and \textsf{review}
entries, \textsf{nameb} applies to the \textsf{title} instead of the
\textsf{issuetitle}, should the latter be present.)  You should
present names in this field exactly as you would those in an
\textsf{author} or \textsf{translator} field, and the package will
concatenate this field with \textsf{namea} if they are identical.  See
under the \textbf{translator} field below for the full details.
Please note that, as the field is highly single-entry specific, if you
are using \textsf{Biber} \textsf{nameb} isn't inherited from a
\textsf{crossref}'ed parent entry.  Please note, also, that in
\textsf{biblatex-chicago's} name-finding algorithms \textsf{nameb}
takes precedence over \textsf{translator}.  Cf.\ \textsf{namea},
\textsf{namec}, \textsf{origlanguage} (section~\ref{sec:related}),
\textsf{translator}, \textsf{userf} (section~\ref{sec:related}), and
the macros \cmd{partedit}, \cmd{parttrans}, \cmd{parteditandtrans},
\cmd{partcomp}, \cmd{parteditandcomp}, \cmd{parttransandcomp}, and
\cmd{partedittransandcomp} in section~\ref{sec:formatcommands}.

\mybigspace The \mymarginpar{\textbf{namec}} \emph{Manual} (14.103)
specifies that works without an author may be listed under an editor,
translator, or compiler, assuming that one is available, and it also
specifies the strings to be used with the name(s) of compiler(s).  All
this suggests that the \emph{Manual} considers this to be standard
information that should be made available in a bibliographic
reference, so I have added that possibility to the many that
\textsf{biblatex} already provides, such as the \textsf{editor},
\textsf{translator}, \textsf{commentator}, \textsf{annotator}, and
\textsf{redactor}, along with writers of an \textsf{introduction},
\textsf{foreword}, or \textsf{afterword}.  Since \textsf{biblatex}
doesn't offer a \textsf{compiler} field, I have adopted for this
purpose the otherwise unused field \textsf{namec}.  It is important to
understand that, despite the analogous name, this field does not
function like \textsf{namea} or \textsf{nameb}, but rather like
\textsf{editor} or \textsf{translator}, and therefore if used will be
associated with whichever title field these latter two would be were
they present in the same entry.  Identical fields among these three
will be concatenated by the package, and concatenated too with the
(usually) unnecessary commentator, annotator and the rest.  Also
please note that I've arranged the concatenation algorithms to include
\textsf{namec} in the same test as \textsf{namea} and \textsf{nameb},
so in this particular circumstance you can, if needed, make
\textsf{namec} analogous to these two latter, \textsf{title}-only
fields.  (See above under \textbf{editortype} for details of how you
may, in certain circumstances, use that field, or the
\textsf{nameatype} field, to identify a compiler.)

\mylittlespace It might conceivably be necessary at some point to
identify the compiler(s) of a \textsf{title} separate from the
compiler(s) of a \textsf{booktitle} or \textsf{maintitle}, but for the
moment I've run out of available \textsf{name} fields, so you'll have
to fall back on the \cmd{partcomp} macro or the related
\cmd{parteditandcomp}, \cmd{parttransandcomp}, and
\cmd{partedittransandcomp}, on which see Commands
(section~\ref{sec:formatcommands}) below.  (Future releases may be
able to remedy this.)  It may be as well to mention here too that of
the names that can be substituted for the missing \textsf{author} at
the head of an entry, \textsf{biblatex-chicago-notes} will choose a
\textsf{namea} if present, then an \textsf{editor}, a \textsf{nameb},
or a \textsf{translator}, with \textsf{namec} coming last, assuming
that the fields aren't identical, and therefore to be concatenated.
The alphabetization routines should work properly for any of these
names, but do please remember that if you want the package to skip
over any names you can employ the \texttt{use<name>=false} options.
Indeed, \textsf{biblatex's} \texttt{usenamec} has replaced the old
Chicago-specific \texttt{usecompiler}, which is deprecated.

\mybigspace As \mymarginpar{\textbf{note}} in standard
\textsf{biblatex}, this field allows you to provide bibliographic data
that doesn't easily fit into any other field.  In this sense, it's
very like \textsf{addendum}, but the information provided here will be
printed just before the publication data.  (See chaucer:alt,
chaucer:liferecords, cook:sotweed, emerson:nature, and rodman:walk for
examples of this usage in action.)  It also has a specialized use in
all the periodical types (\textsf{article}, \textsf{periodical}, and
\textsf{review}), where it holds supplemental information about a
\textsf{journaltitle}, such as \enquote{special issue}
(conley:fifthgrade, good:wholeissue).  In all uses, if your data
begins with a word that would ordinarily only be capitalized at the
beginning of a sentence, then simply ensure that that word is in
lowercase, and \textsf{biblatex-chicago-notes} will automatically do
the right thing.  Cf.\ \textsf{addendum}.

\mybigspace This \mymarginpar{\textbf{number}} is a standard
\textsf{biblatex} field, steadily accumulating uses in
\textsf{biblatex-chicago}.  It may contain the number of a
\textsf{journaltitle} in an \textsf{article} or \textsf{review} entry,
the number of a \textsf{title} in a \textsf{periodical} entry, the
volume/number of a book (or musical recording) in a \textsf{series},
the (generally numerical) specifier of the \textsf{type} in a
\textsf{report} entry, the archive location (or database accession
number) of a \textsf{dataset} entry, and the number of a
national or international standard in a \textsf{standard}
entry.  Generally, in an \textsf{article}, \textsf{periodical}, or
\textsf{review} entry, this will be a plain cardinal number, but in
such entries \textsf{biblatex-chicago} now does the right thing if you
have a list or range of numbers (unsigned:ranke).  In any
\textsf{book}-like entry the field may well contain considerably more
information, including even a reference to \enquote{2nd ser.,} for
example, while the \textsf{series} field in such an entry will contain
the name of the series, rather than a number.  This field is also the
place for the patent number in a \textsf{patent} entry.  Cf.\
\textsf{issue} and \textsf{series}.  (Cf.\ 14.123--25 and boxer:china,
palmatary:pottery, wauchope:ceramics; 14.171 and beattie:crime,
conley:fifthgrade, friedman:learning, garrett, gibbard, hlatky:hrt,
mcmillen:antebellum, rozner:liberation, and warr:ellison; 14.257 and
genbank:db; 14.259 and niso:bibref; 14.263 and holiday:fool.)

\mylittlespace \textbf{NB}: This may be an opportune place to point
out that the \emph{Manual} (14.147) prefers arabic to roman numerals
in most circumstances (chapters, volumes, series numbers, etc.), even
when such numbers might be roman in the work cited.  The obvious
exception is page numbers, in which roman numerals indicate that the
citation came from the front matter, and should therefore be retained.

\mybigspace A \mymarginpar{\textbf{options}} standard
\textsf{biblatex} field, for setting certain options on a per-entry
basis rather than globally.  Information about some of the more common
options may be found above under \textsf{author} and below in
section~\ref{sec:options}.  See chaucer:alt, eliot:pound,
herwign:office, lecarre:quest, and mla:style for examples of the field
in use.

\mybigspace A \mymarginpar{\textbf{organization}} standard
\textsf{biblatex} field, retained mainly for use in the \textsf{misc},
\textsf{online}, and \textsf{manual} entry types, where it may be of
use to specify a publishing body that might not easily fit in other
categories.  In \textsf{biblatex}, it is also used to identify the
organization sponsoring a conference in a \textsf{proceedings} or
\textsf{inproceedings} entry, and I have retained this as a
possibility, though the \emph{Manual} is silent on the matter.

\mybigspace This \mymarginpar{\textbf{origdate}} \textsf{biblatex}
field allows you to provide more than one full date specification for
those references which need it.  As with the analogous \textsf{date}
field, you provide the date (or range of dates) in \textsc{iso}8601
format, i.e., \texttt{yyyy-mm-dd}.  (You can also provide a time stamp
in the field, after an uppercase \enquote{\texttt{T}}, but I foresee
this being very rarely needed in the notes \&\ bibliography style.
See table~\ref{ad:date:extras} for \textsf{biblatex-chicago's}
implementation of \textsf{biblatex's} enhanced date specifications.)
In most entry types, you would use \textsf{origdate} to provide the
date of first publication of a work, most usually needed only in the
case of reprint editions, but also recommended by the \emph{Manual}
for electronic editions of older works (14.114, 14.162;
aristotle:metaphy:gr, emerson:nature, james:ambassadors,
schweitzer:bach).  In the \textsf{letter} and \textsf{misc} (with
\textsf{entrysubtype}) entry types, the \textsf{origdate} identifies
when a letter (or similar) was written.  In such \textsf{misc}
entries, you can choose between an \textsf{origdate} and a
\textsf{date} field for this purpose, depending on how you want the
date formatted (day-month-year or month-day-year, respectively), while
in \textsf{letter} entries the \textsf{date} applies to the
publication of the whole collection.  If such a published collection
were itself a reprint, improvisation in the \textsf{location} field
might be able to rescue the situation.  (See jackson:paulina\hc
letter, white:ross:memo, white:russ, and white:total for how
\textsf{letter} entries usually work; creel:house shows the field in
action in a \textsf{misc} entry, while spock:interview uses
\textsf{date}.)

\mylittlespace In \textsf{music} entries, you can use the
\textsf{origdate} in two separate but related ways.  First, it can
identify the recording date of an entire disc, rather than of one
track on that disc, which would go in \textsf{eventdate}.  (Compare
holiday:fool with nytrumpet:art.)  The style will automatically
prepend the bibstring \texttt{recorded} to the date, but you can
change it with the \textsf{userd} field.  Be aware, however, that if
an entry also has an \textsf{eventdate}, then \textsf{userd} will
apply to that, instead, and you'll be forced to accept the default
string.  Second, the \textsf{origdate} can provide the original
release date of an album.  For this to happen, you need to put the
string \texttt{reprint} in the \textsf{pubstate} field, which is a
standard mechanism across many other entry types for identifying a
reprinted work.  (See floyd:atom.)

\mylittlespace A couple of further notes are in order.  First,
\textsf{artwork} and \textsf{image} entries (which see) have their own
scheme.  Here, the style uses the earlier of two dates as the creation
date of the work while the later is the printing date of, e.g., a
particular exemplar of a photograph or of an etching.  In such an
entry, the \textsf{origdate} may well be a creation date.  Second,
because the \textsf{origdate} field only accepts numbers, some
improvisation may be needed if you wish to include \enquote{n.d.}\
(\verb+\bibstring{nodate}+) in an entry.  In \textsf{letter} and
\textsf{misc}, this information can be placed in \textsf{titleaddon},
but in other entry types you may need to use the \textsf{location}
field.  (The \textsf{origyear} field usually works, too.)

\mybigspace See
\vspace{-14.2pt}
\mymarginpar{\textbf{origlanguage}\\
\textbf{origlocation}\\\textbf{origpublisher}}
section~\ref{sec:related}, below.
\vspace{18pt}

\mybigspace This \mymarginpar{\textbf{origtimezone}} field can, if
necessary, specify the time zone associated with a time stamp given as
part of an \textsf{origdate}.  The \emph{Manual} prefers initialisms
like \enquote{EST} for this purpose, and you can provide parentheses
around it at your discretion (cp.\ 10.41 and 14.191).

\mybigspace This \mymarginpar{\textbf{pages}} is the standard
\textsf{biblatex} field for providing page references.  In many
\textsf{article} and \textsf{review} entries you'll find this contains
something other than a page number, e.g. a section name or edition
specification (14.191; kozinn:review, nyt:obittrevor,
nyt:trevorobit).  Of course, the same may be true of almost any sort
of entry, though perhaps with less frequency.  Curious readers may
wish to look at brown:bremer (14.180) for an example of a
\textsf{pages} field used to facilitate reference to a two-part
journal article.  Cf.\ \textsf{number} for more information on the
\emph{Manual}'s preferences regarding the formatting of numerals;
\textsf{bookpagination} and \textsf{pagination} provide details about
\textsf{biblatex's} mechanisms for specifying what sort of division a
given \textsf{pages} field contains; and \textsf{usera} discusses a
different way to present the section information pertaining to a
newspaper article.

\mylittlespace David Gohlke brought to my attention a discussion that
took place a couple of years ago on
\href{http://tex.stackexchange.com/questions/44492/biblatex-chicago-style-page-ranges}{Stackexchange}
regarding the automatic compression of page ranges, e.g., 101-{-}109
in the .bib file or in the \textsf{postnote} field would become 101--9
in the document.  \textsf{Biblatex} has long had the facilities for
providing this, and though the \emph{Manual's} rules (9.61) are fairly
complicated, Audrey Boruvka fortunately provided in that discussion
code that implements the specifications.  As some users may well be
accustomed to compressing page ranges themselves in their .bib files,
and in their \textsf{postnote} fields, I have made the activation of
this code a package option, so setting \texttt{compresspages=true}
when loading \textsf{biblatex-chicago} should automatically give you
the Chicago-recommended page ranges.  \textbf{NB}: the code now
resides in \textsf{biblatex-chicago.sty}, so if you don't load that
package then you'll need to copy the code into your preamble for the
option to have the desired effect.

\mybigspace This, \mymarginpar{\textbf{pagination}} a standard
\textsf{biblatex} field, allows you automatically to prefix the
appropriate identifying string to information you provide in the
\textsf{postnote} field of a citation command, whereas
\textsf{bookpagination} allows you to prefix a string to the
\textsf{pages} field.  Please see \textbf{bookpagination} above for
all the details on this functionality, as aside from the difference
just mentioned the two fields are equivalent.

\mybigspace Standard \mymarginpar{\textbf{part}} \textsf{biblatex}
field, which identifies physical parts of a single logical volume in
\textsf{book}-like entries, not in periodicals.  It has the same
purpose in \textsf{biblatex-chicago-notes}, but because the
\emph{Manual} (14.121) calls such a thing a \enquote{book} and not a
\enquote{part,} the string printed in notes and bibliography will, at
least in English, be \enquote{\texttt{bk.}\hspace{-2pt}}\ instead of
the plain dot between volume number and part number
(harley:cartography, lach:asia).  If the field contains something
other than a number, \textsf{biblatex-chicago} will print it as is,
capitalizing it if necessary, rather than supplying the usual
bibstring, so this provides a mechanism for altering the string to
your liking.  The field will be printed in the same place in any entry
as would a \textsf{volume} number, and although it will most usually
be associated with such a number, it can also function independently,
allowing you to identify parts of works that don't fit into the
standard scheme.  If you need to identify \enquote{parts} or
\enquote{books} that are part of a published \textsf{series}, for
example, then you'll need to use a different field, (which in this
case would be \textsf{number} [palmatary:pottery]).  Cf.\
\textsf{volume}.

\mybigspace Standard \mymarginpar{\textbf{publisher}}
\textsf{biblatex} field.  Remember that \enquote{\texttt{and}} is a
keyword for connecting multiple publishers, so if a publisher's name
contains \enquote{and,} then you should either use the ampersand (\&)
or enclose the whole name in additional braces.  (See \emph{Manual}
14.133--41; aristotle:metaphy:gr, cohen:schiff, creasey:ashe:blast,
dunn:revolutions.)

%%\enlargethispage{.5\baselineskip}

\mylittlespace There are, as one might expect, a few further
subtleties involved here.  If you give two publishers in the field
they will both be printed, separated by a forward slash in both notes
and bibliography (14.90; sereny:cries).  The 17th edition generally is
rather keener than the 16th on using just one, particularly so in the
case when the parent company of an imprint is also listed on a title
page, in which case only the imprint need be included in your
apparatus (14.138).  If an academic publisher issues \enquote{certain
  books through a special publishing division or under a special
  imprint or as part of a publishing consortium (or joint imprint),}
this arrangement may be specified in the \textsf{publisher} field
(14.139; cohen:schiff).  If a book has two co-publishers \enquote{in
  different countries} (14.140), then the simplest thing to do is to
choose one, probably the nearest one geographically.  If you feel it
necessary to include both, then levistrauss:savage demonstrates one
way of doing so, using a combination of the \textsf{publisher} and
\textsf{location} fields.  If the work is self-published, you can
specify this in the \textsf{pubstate} field (see below), and any
commercial self-publishing platform would go in \textsf{publisher}
(14.137).  Books published before 1900 can, at your discretion,
include only the place (if known) and the date (14.128).  If for some
reason you need to indicate the absence of a publisher, the
abbreviation given by the \emph{Manual} is \texttt{n.p.}, though this
can also stand for \enquote{no place.}  The \emph{Manual} also
mentions {s.n.}\,(= \emph{sine nomine}) to specify the lack of a
publisher (10.42).

\mybigspace In \mymarginpar{\textbf{pubstate}} response to new
specifications in the 17th edition of the \emph{Manual} (esp.\
14.137), I have tried to generalize the functioning of the
\textsf{pubstate} field in all entry types.  The \texttt{reprint}
string still has a special status there, being ignored in
\textsf{video} entries and provoking a syntactic change in the
presentation of dates in \textsf{music} entries, while in other types
allowing the presentation of reprinted titles.  Other strings are
divided into two types: those which \textsf{biblatex-chicago} will
print as the \textsf{year}, which currently means \emph{only} those
for which \textsf{biblatex} contains bibstrings indicating works soon
to be published, i.e., \texttt{forthcoming}, \texttt{inpreparation},
\texttt{inpress}, and \texttt{submitted}; and those, i.e., everything
else, which will be printed before, and in close association with,
other information about the publisher of a work.  The four in the
first category will always be localized, as will \texttt{reprint} and
\texttt{selfpublished} (and anything else that \textsf{biblatex} finds
to be a \cmd{bibstring}) from the second category.  All other strings
will be printed as-is, capitalized if needed, just before the
publisher (author:forthcoming, contrib:contrib, schweitzer:bach).

\mybigspace I \mymarginpar{\textbf{redactor}} have implemented this
field just as \textsf{biblatex}'s standard styles do, even though the
\emph{Manual} doesn't actually mention it.  It may be useful for some
purposes.  Cf.\ \textsf{annotator} and \textsf{commentator}.

\mybigspace See \mymarginpar{\textbf{reprinttitle}}
section~\ref{sec:related}, below.

\mybigspace A \mymarginpar{\textbf{series}} standard \textsf{biblatex}
field, usually just a number in an \textsf{article},
\textsf{periodical}, or \textsf{review} entry, almost always the name
of a publication series in \textsf{book}-like entries, and providing
similar identifying information associated with a \textsf{number} in
\textsf{music} and \textsf{standard} entries.  If you need to attach
further information to the \textsf{series} name in a
\textsf{book}-like entry, then the \textsf{number} field is again the
place for it, whether it be a volume, a number, or even something like
\enquote{2nd ser.} or \enquote{\cmd{bibstring\{oldseries\}}.}  Of
course, you can also use \cmd{bibstring\{oldseries\}} or
\cmd{bibstring\{newseries\}} in an \textsf{article} entry, but there
you would place it in the \textsf{series} field itself.  (In fact, the
\textsf{series} field in \textsf{article}, \textsf{periodical}, and
\textsf{review} entries is one of the places where \textsf{biblatex}
allows you just to use the plain bibstring \texttt{oldseries}, for
example, rather than making you type \verb+\bibstring{oldseries}+.
The \textsf{type} field in \textsf{manual}, \textsf{patent},
\textsf{report}, and \textsf{thesis} entries also has this
auto-detection mechanism in place; see the discussion of
\cmd{bibstring} below for details.)  In whatever entry type, these
bibstrings produce the required abbreviation, which thankfully is the
same in both notes and bibliography.  (For books and similar entries,
see 14.123--26; boxer:china, browning:aurora, palmatary:pottery,
plato:republic:gr, wauchope:ceramics; for periodicals, see 14.184;
garaud:gatine, sewall:letter; also niso:bibref, nytrumpet:art) Cf.\
\textsf{number} for more information on the \emph{Manual}'s
preferences regarding the formatting of numerals.

\paragraph*{\protect\mymarginpar{\textbf{shortauthor}}}
\label{sec:shortauthor}

This is a standard \textsf{biblatex} field, but
\textsf{biblatex-chicago-notes} makes considerably greater use of it
than the standard styles.  For the purposes of the Chicago style, the
field provides the name to be used in the short form of a footnote.
In the vast majority of cases, you don't need to specify it, because
the \textsf{biblatex} system selects the author's last name from the
\textsf{author} field and uses it in such a reference, and if there is
no \textsf{author} it will search \textsf{namea}, \textsf{editor},
\textsf{nameb}, \textsf{translator}, and \textsf{namec}, in that
order.  In an author-less \textsf{article} or \textsf{review} entry
(\textsf{entrysubtype} \texttt{magazine}), where
\textsf{biblatex-chicago-notes} will use the \textsf{journaltitle} as
the author, you can use the \textsf{shortjournal} field instead, but
you'll need to set up the \texttt{journalabbrev} option to make sure
it's actually printed.  (See \textsf{shortjournal}, below.)  In
author-less \textsf{manual} entries, where the \textsf{organization}
will be so used, the style automatically uses any \textsf{shortauthor}
in the short note form, though it will sort by the
\textsf{organization} in the bibliography (dyna:browser,
gourmet:052006, lakeforester:pushcarts, nyt:trevorobit).

\mylittlespace As mentioned under \textsf{editortype}, the
\emph{Manual} (14.32) recommends against providing the identifying
string (e.g., ed.\ or trans.)\ in the short note form, and
\textsf{biblatex-chicago-notes} follows their recommendation.  If you
need to provide these strings in such a citation, then you'll have to
do so by hand in the \textsf{shortauthor} field, or in the
\textsf{shorteditor} field, whichever you are using.

\mybigspace Like \mymarginpar{\textbf{shorteditor}}
\textsf{shortauthor}, a field to provide a name for a short footnote,
in this case for, e.g., a \textsf{collection} entry that typically
lacks an author.  The \textsf{shortauthor} field works just as well in
most situations, but if you have set \texttt{useauthor=false} (and not
\texttt{useeditor=false}) in an entry's \textsf{options} field, then
only \textsf{shorteditor} will be recognized.  It may be worth
pointing out that, because \textsf{biblatex-chicago} also provides a
\textsf{namea} field for the editor of a \textsf{title} as opposed to
a \textsf{main-} or \textsf{booktitle}, and because in standard use
the \textsf{namea}, if present, will be chosen to head a bibliography
entry before the \textsf{editor}, you should present the shortened
\textsf{namea} here instead of a shortened \textsf{editor} in such
cases.  Cf.\ \textsf{editortype}, above.

\mybigspace This \mymarginpar{\textbf{shorthand}} is
\textsf{biblatex}'s mechanism for using abbreviations in place of the
usual short note form, and in previous releases I left it effectively
unmodified in \textsf{biblatex-chicago-notes}, apart from a few
formatting tweaks.  At the request of Kenneth Pearce, and following
some hints in the \emph{Manual}, I have made the system considerably
more flexible, which I hope might be useful for those with specialized
formatting needs.  In the default configuration, any entry which
contains a \textsf{shorthand} field will produce a normal first note,
either long or short according to your package options, informing the
reader that the work will hereafter be cited by this abbreviation.  As
in standard \textsf{biblatex}, the \cmd{printshorthands} command, now
an alias for \verb+\printbiblist{shorthand}+, will produce a formatted
list of abbreviations for reference purposes, a list which the
\emph{Manual} suggests should be placed either in the front matter
(when using footnotes) or before the endnotes, in case these are used.

\mylittlespace I have provided five options to alter these defaults.
First, there is a citation command, \cmd{shorthandcite}, which will
print the \textsf{shorthand} even at the first citation.  I have only
provided the most general form of this command, so you'll need to put
it inside parentheses or in a \cmd{footnote} command yourself.
Alternately, you can set the \texttt{shorthandfirst} option to
\texttt{true}, either in the preamble or in individual entries, to get
the same effect, somewhat more simply.  Third, I have included two
\texttt{bibenvironments} for use with the \texttt{env} option to the
\cmd{printshorthands} command: \texttt{losnotes} is designed to allow
a list of shorthands to appear inside footnotes, while
\texttt{losendnotes} does the same for endnotes.  Their main effect is
to change the font size, and in the latter case to clear up some
spurious punctuation and white space that I see on my system when
using endnotes.  (You'll probably also want to use the option
\texttt{heading=none} in order to get rid of the [oversized] default,
providing your own within the \cmd{footnote} command.)  Fourth, I have
provided a package option, \texttt{shorthandfull}, which prints
entries in the list of shorthands which contain full bibliographical
information, effectively allowing you to eschew the bibliography in
favor of a fortified shorthand list.  (See 13.67, 14.59--60, and also
\textsf{biblatex.pdf} for more information.)  Finally, the
\texttt{shorthandintro} option (section~\ref{sec:useropts}) allows
you to control whether the first, long citation will introduce the
\textsf{shorthand} at all, rather than just leaving it to the
shorthand list (or perhaps \enquote{common knowledge}) to clarify the
reference.

\mylittlespace Alexandre Roberts suggested a further refinement to
\textsf{shorthand} behavior, which allows for it to appear in the
place of the usual abbreviated citation of parent entries
cross-referenced by several different child entries.  In such a case,
instead of the usual \enquote{\ldots\,in Author, \emph{Title},
  24--38,} you would see instead \enquote{\ldots\,in \emph{ShrtHd},
  24--38.}  There are several steps required for enabling this
behavior.  First, you need to set the package option
\texttt{inheritshorthand} to \texttt{true}, which allows child entries
to inherit the necessary fields from their cross-referenced parents.
Second, you'll probably want to use the \textsf{shorthandintro} field
somehow to clarify that the \textsf{shorthand} applies to the
\emph{parent} rather than to the \emph{child}, as otherwise the
reference will be ambiguous.  Third, you'll need to put
\texttt{skipbiblist}, formerly \texttt{skiplos}, in the
\textsf{options} field of the child entries so that the
\textsf{shorthand} itself appears in the list of shorthands
\emph{only} next to the parent entry, and not also next to all of its
children.

\mylittlespace As I mentioned above under \textbf{crossref}, I
formerly recommended against using shorthands with cross-references,
but this extension of their use makes sense as an extra space-saving
measure.  I'm not certain that I've identified all the possible
drawbacks to enabling the \texttt{inheritshorthand} option, so care is
still needed, at least in the current state of
\textsf{biblatex-chicago-notes}.  Please report any problems you might
have with this functionality to the email address at the head of this
documentation.

\mybigspace When \mymarginpar{\textbf{shorthandintro}} you include a
\textsf{shorthand} in an entry, it will ordinarily appear the first
time you cite the work, at the end of a long note, surrounded by
parentheses and prefaced by the phrase \enquote{hereafter cited as.}
To modify this, you can either use the \texttt{shorthandintro} option
(section~\ref{sec:useropts}) or you can use this standard
\textsf{biblatex} field to change the formatting and the phrase to
suit your needs.  Please note, first, that you need to include the
shorthand in this field as you intend it to appear and, second, that
you still need the \textsf{shorthand} field present in order to ensure
the appropriate presentation of that shorthand in later citations and
in the list of shorthands.  Finally, I've tried to allow for as many
different styles of notification as possible, so by default the only
punctuation that will appear between the rest of the citation and the
\textsf{shorthandintro} is a space.  You can change this punctuation,
either in the preamble for the whole document or in individual
entries, using the \texttt{shorthandpunct} option, documented in
section~\ref{sec:chicpreset}.  If the available option keys aren't
adequate, you can use \texttt{none} and then provide custom
punctuation inside the \textsf{shorthandintro} field itself.

\mybigspace A \mymarginpar{\textbf{shortjournal}} special
\textsf{biblatex} field, used both to provide an abbreviated form of a
\textsf{journaltitle} in notes and/or bibliography and to facilitate
the creation of a list of journal abbreviations rather in the manner
of a \textsf{shorthand} list.  As requested by user BenVB, you can now
utilize this functionality in your documents, but there are several
steps to take in order to do so.  First, you'll need to provide both
\textsf{shortjournal} and \textsf{journaltitle} fields in the entry
types that use them, i.e., mainly \textsf{article} and \textsf{review}
entries.  In \textsf{periodical} entries the \textsf{title} field
presents what would be the \textsf{journaltitle} in the previous two,
so in such entries you can provide the standard \textsf{shorttitle}
field to accompany the \textsf{title}, and \textsf{biblatex-chicago}
will automatically copy the \textsf{shorttitle} into a
\textsf{shortjournal}.

\mylittlespace Having done this, you need to set the
\texttt{journalabbrev} option either when loading
\textsf{bibla\-tex-chicago} or in the \textsf{options} field of
individual .bib entries.  By default, this option is not set, so your
\textsf{shortjournal} fields will be silently ignored.  There are
three other settings:\ \texttt{true} prints the shortened fields both
in notes and bibliography, \texttt{notes} only in notes, and
\texttt{bib} only in the bibliography.  Should you wish to present a
list of these abbreviations with their expansions, then you need to
use the \verb+\printbiblist{shortjournal}+ command, perhaps with a
\texttt{title} option to differentiate the list from any
\textsf{shorthand} list.  As with \textsf{shorthand} lists, I have
provided two \texttt{bibenvironments} for printing this list in foot-
or endnotes (\texttt{sjnotes} and \texttt{sjendnotes}, respectively),
to be used with the \texttt{env} option to \cmd{printbiblist}.  Again
as with \textsf{shorthands}, you'll probably want to use the option
\texttt{heading=none} when using these environments, just to turn off
the (oversized) default, and perhaps provide your own title within the
\cmd{footnote} command.  Finally, if you don't like the default
formatting of the abbreviations in the list (bold italic), you can
roll your own using \verb+\DeclareFieldFormat{shortjournalwidth}+
--- you can see its default definition at the top of
\textsf{chicago-notes.bbx}.

\mybigspace A \mymarginpar{\textbf{shortseries}} special
\textsf{biblatex} field, used both to provide an abbreviated form of a
(book) \textsf{series} in notes and/or bibliography and to facilitate
the creation of a list of such abbreviations rather in the manner of a
\textsf{shorthand} list.  As with the \textsf{shortjournal} field, its
inclusion in \textsf{biblatex-chicago} was requested by user BenVB,
and it is now available in entry types which have book-like series
titles rather than journal-like numbers in the \textsf{series} field,
to wit: \textsf{audio, book, bookinbook, collection, inbook,
  incollection, inproceedings, inreference, letter, manual, music,
  mvbook, mvcollection, mvproceedings, mvreference, reference, report,
  standard, suppbook,} and \textsf{video}.  There are several steps to
take in order to use the field.  First, you'll need to provide both
\textsf{shortseries} and \textsf{series} fields in the entry, then
you'll need to set the \texttt{seriesabbrev} option either when
loading \textsf{biblatex-chicago}, for the whole document or for
specific entry types, or in the \textsf{options} field of individual
.bib entries.  By default, this option is not set, so your
\textsf{shortseries} fields will be silently ignored.  There are three
other settings:\ \texttt{true} prints the shortened fields both in
notes and bibliography, \texttt{notes} only in notes, and \texttt{bib}
only in the bibliography.  Should you wish to present a list of these
abbreviations with their expansions, then you need to use the
\verb+\printbiblist{shortseries}+ command, perhaps with a
\texttt{title} option to differentiate the list from any
\textsf{shorthand} list.  As with \textsf{shorthand} lists, I have
provided two \texttt{bibenvironments} for printing this list in foot-
or endnotes (\texttt{shsernotes} and \texttt{shserendnotes},
respectively), to be used with the \texttt{env} option to
\cmd{printbiblist}.  Again as with \textsf{shorthands}, you'll
probably want to use the option \texttt{heading=none} when using these
environments, just to turn off the (oversized) default, and perhaps
provide your own title within the \cmd{footnote} command.  Finally, if
you don't like the default formatting of the abbreviations in the list
(plain roman), you can roll your own using
\verb+\DeclareFieldFormat{shortserieswidth}+ --- you can see its
default definition at the top of \textsf{chicago-notes.bbx}.

\mybigspace A \mymarginpar{\textbf{shorttitle}} standard
\textsf{biblatex} field, primarily used to provide an abbreviated
title for short notes.  (It is also the way to hook
\textsf{periodical} entries into the \textsf{shortjournal} mechanism,
on which see the previous entry.)  In \textsf{biblatex-chicago-notes},
you need to take particular care with \textsf{letter} entries, where,
as explained above, the \emph{Manual} requires a special format
(\enquote{\texttt{to Recipient}}).  (See 14.111;
jackson:paulina:letter, white:ross:memo, white:russ.)  Some
\textsf{misc} entries (with an \textsf{entrysubtype}) also need
special attention.  (See creel:house, where the full \textsf{title} is
used as the \textsf{shortauthor} + \textsf{shorttitle} by using
\cmd{headlesscite} commands.)  Remember, also, that the generic titles
in \textsf{review} and \textsf{misc} entries may not want
capitalization in all contexts, so, as with the \textsf{title} field,
if you begin a \textsf{shorttitle} with a lowercase letter the style
will do the right thing (barcott:review, bundy:macneil,
Clemens:letter, kozinn:review, ratliff:review, unsigned:ranke).

\mybigspace Standard \mymarginpar{\textbf{sortkey} \\\textbf{sortname}
  \\\textbf{sorttitle} \\\textbf{sortyear}} \textsf{biblatex} fields,
designed to allow you to specify how you want an entry alphabetized in
a bibliography.  The \textsf{sortkey} field trumps all other sorting
information, while the others offer more fine-grained control.  In
general, if an entry doesn't turn up where you expect or want it,
these fields should provide the solution.  Entries with a corporate
author can now omit the definite or indefinite article, which should
help (14.70, 14.84; cotton:manufacture, nytrumpet:art).
\textsf{Biblatex-chicago} also includes the three supplemental name
fields (\textsf{name[a-c]}) in the sorting algorithm, so once again
you should find that a \textsf{sortkey} is needed less than before.
Still, some entries without a name field of any sort, particularly
those with a definite or indefinite article beginning the
\textsf{title}, may require assistance (greek:filmstrip,
grove:sibelius, nyt:obittrevor, virginia:plantation).  Please consult
\textsf{biblatex.pdf} and the remarks on \cmd{DeclareSortingTemplate}
in section~\ref{sec:formatopts}, below.

\mybigspace The \mymarginpar{\textbf{subtitle}} subtitle for a
\textsf{title} --- see next entry.

\mybigspace This \mymarginpar{\textbf{timezone}} field can, if
necessary, specify the time zone associated with a time stamp given as
part of an \textsf{date}.  The \emph{Manual} prefers initialisms like
\enquote{EST} for this purpose, and you can provide parentheses around
it at your discretion (cp.\ 10.41 and 14.191).

\mybigspace In \mymarginpar{\textbf{title}} the vast majority of
cases, this field works just as it always has in \textsc{Bib}\TeX, and
just as it does in \textsf{biblatex}.  Nearly every entry will have
one, the most likely exceptions being \textsf{incollection} or
\textsf{online} entries with a merely generic title, instead of a
specific one (centinel:letters, powell:email).  The main source of
difficulties flows from the \emph{Manual}'s rules for formatting
\textsf{titles}, rules which also hold for \textsf{booktitles} and
\textsf{maintitles}.  The whole point of using a
\textsf{biblatex}-based system is for it to do the formatting for you,
and in most cases \textsf{biblatex-chicago-notes} does just that,
surrounding titles with quotation marks, italicizing them, or
occasionally just leaving them alone.  When, however, a title is
quoted within a title, then you need to know some of the rules.  A
summary here should serve to clarify them, and help you to understand
when \textsf{biblatex-chicago-notes} might need your help in order to
comply with them.

\mylittlespace The internal rules of \textsf{biblatex-chicago-notes}
are as follows:

\begin{description}
\item[\qquad Italics:] \textsf{booktitle}, \textsf{maintitle}, and
  \textsf{journaltitle} in all entry types; \textsf{title} of
  \textsf{artwork}, \textsf{book}, \textsf{bookinbook},
  \textsf{booklet}, \textsf{collection}, \textsf{image},
  \textsf{manual}, \textsf{misc} (with no \textsf{entrysubtype}),
  \textsf{performance}, \textsf{periodical}, \textsf{proceedings},
  \textsf{report}, \textsf{standard}, \textsf{suppbook}, and
  \textsf{suppcollection} entry types.
\item[\qquad Quotation Marks:] \textsf{title} of \textsf{article},
  \textsf{inbook}, \textsf{incollection}, \textsf{inproceedings},
  \textsf{online}, \textsf{periodical}, \textsf{thesis}, and
  \textsf{unpublished} entry types, \textsf{issuetitle} in
  \textsf{article}, \textsf{periodical}, and \textsf{review} entry
  types.
\item[\qquad Sentence cased:] \textsf{title} in \textsf{patent}
  entries.
\item[\qquad Unformatted:] \textsf{booktitleaddon},
  \textsf{maintitleaddon}, and \textsf{titleaddon} in all entry types,
  \textsf{title} of \textsf{customc}, \textsf{letter}, \textsf{misc}
  (with an \textsf{entrysubtype}), \textsf{review}, and
  \textsf{suppperiodical} entry types.
\item[\qquad Italics or Quotation Marks:] All of the audiovisual entry
  types --- \textsf{audio}, \textsf{music}, and \textsf{video} ---
  have to serve as analogues both to \textsf{book} and to
  \textsf{inbook}.  Therefore, if there is both a \textsf{title} and a
  \textsf{booktitle}, then the \textsf{title} will be in quotation
  marks.  If there is no \textsf{booktitle}, then the \textsf{title}
  will be italicized, unless you provide an \textsf{entrysubtype}.
\end{description}

Now, the rules for which entry type to use for which sort of work tend
to be fairly straightforward, but in cases of doubt you can consult
section \ref{sec:entrytypes} above, the examples in
\textsf{notes-test.bib}, or go to the \emph{Manual} itself,
8.156--201.  Assuming, then, that you want to present a title within a
title, and you know what sort of formatting each of the two would, on
its own, require, then the following rules apply:

\begin{enumerate}
\item Inside an italicized title, all other titles are enclosed in
  quotation marks and italicized, so in such cases all you need to do
  is provide the quotation marks using \cmd{mkbibquote}, which will
  take care of any following punctuation that needs to be brought
  within the closing quotation mark(s) (14.94; donne:var,
  mchugh:wake).
\item Inside a quoted title, you should present another title as it
  would appear if it were on its own, so in such cases you'll need to
  do the formatting yourself.  Within the double quotes of the title
  another quoted title would take single quotes --- the
  \cmd{mkbibquote} command does this for you automatically, and also,
  I repeat, takes care of any following punctuation that needs to be
  brought within the closing quotation mark(s).  (See 14.94--95;
  garrett, loften:hamlet, murphy:silent, white:callimachus.)
\item Inside a plain title (most likely in a \textsf{review} entry or
  a \textsf{titleaddon} field), you should present another title as it
  would appear on its own, once again formatting it yourself using
  \cmd{mkbibemph} or \cmd{mkbibquote}.  (barcott:review, gibbard,
  osborne\hc poison, ratliff:review, unsigned:ranke).
\end{enumerate}

The \emph{Manual} provides a few more rules, as well.  A word normally
italicized in text should also be italicized in a quoted or plain-text
title, but should be in roman (\enquote{reverse italics}) in an
italicized title.  A quotation used as a (whole) title (with or
without a subtitle) retains, according to the 16th edition, its
quotation marks in an italicized title if it appears that way in the
source, but I can't find similar instructions in the 17th.  Such a
quotation always retains its quotation marks when the surrounding
title is quoted or plain (14.94; lewis).  A word or phrase in
quotation marks, but that isn't a quotation, retains those marks in
all title types (kimluu:diethyl).

\mylittlespace Finally, please note that in all \textsf{review} (and
\textsf{suppperiodical}) entries, and in \textsf{misc} entries with an
\textsf{entrysubtype}, and only in those entries,
\textsf{biblatex-chicago-notes} will automatically capitalize the
first word of the \textsf{title} after sentence-ending punctuation,
assuming that such a \textsf{title} begins with a lowercase letter in
your .bib database.  See\,\textbf{\textbackslash autocap} in
section~\ref{sec:formatcommands} below for more details.

% %\enlargethispage{-\baselineskip}

\mybigspace Standard \colmarginpar{\textbf{titleaddon}}
\textsf{biblatex} intends this field for use with additions to titles
that may need to be formatted differently from the titles themselves,
and \textsf{biblatex-chicago-notes} uses it in just this way, with the
additional wrinkle that it can, if needed, replace the \textsf{title}
entirely, and this in, effectively, any entry type, providing a fairly
powerful, if somewhat complicated, tool for getting \textsf{biblatex}
to do what you want (cf.\ centinel:letters, powell:email).  This field
will always be unformatted, that is, neither italicized nor placed
within quotation marks, so any formatting you may need within it
you'll need to provide manually yourself.  The single exception to
this rule is when your data begins with a word that would ordinarily
only be capitalized at the beginning of a sentence, in which case you
need then simply ensure that that word is in lowercase, and
\textsf{biblatex-chicago-notes} will automatically do the right thing.
See\,\textbf{\textbackslash autocap} in
section~\ref{sec:formatcommands}, below.  The package and entry
options \texttt{ptitleaddon} and \texttt{ctitleaddon}
(section~\ref{sec:chicpreset}) allow you to customize the punctuation
that appears before the \textsf{titleaddon} field.  Please note,
however, that I have added this field to the \textsf{periodical} entry
type, and that the punctuation there is governed by the
\mycolor{\texttt{jtitleaddon}} option, which defaults to a
\texttt{space}.  (Cf.\ brown:bremer, osborne:poison, reaves:rosen, and
white:ross:memo for examples where the field starts with a lowercase
letter; morgenson:market provides an example where the
\textsf{titleaddon} field, holding the name of a regular column in a
newspaper, is capitalized, a situation that is handled as you would
expect; coolidge:speech shows both entry options for controlling the
punctuation.)

\mybigspace As \mymarginpar{\textbf{translator}} far as possible, I
have implemented this field as \textsf{biblatex}'s standard styles do,
but the requirements specified by the \emph{Manual} present certain
complications that need explaining.  \textsf{Biblatex.pdf} points out
that the \textsf{translator} field will be associated with a
\textsf{title}, a \textsf{booktitle}, or a \textsf{maintitle},
depending on the sort of entry.  More specifically,
\textsf{biblatex-chicago} associates the \textsf{translator} with the
most comprehensive of those titles, that is, \textsf{maintitle} if
there is one, otherwise \textsf{booktitle}, otherwise \textsf{title},
if the other two are lacking.  In a large number of cases, this is
exactly the correct behavior (adorno:benj, centinel:letters,
plato:republic:gr, among others).  Predictably, however, there are
numerous cases that require, for example, an additional translator for
one part of a collection or for one volume of a multi-volume work.
For these cases I have provided the \textsf{nameb} field.  You should
format names for this field as you would for \textsf{author} or
\textsf{editor}, and these names will always be associated with the
\textsf{title} (euripides:orestes) In the algorithm for finding a name
for the head of notes and bibliography entries, \textsf{nameb} takes
precedence over \textsf{translator}.

\mylittlespace I have also provided a \textsf{namea} field, which
holds the editor of a given \textsf{title} (euripides\hc orestes).  If
\textsf{namea} and \textsf{nameb} are the same,
\textsf{biblatex-chicago} will concatenate them, just as
\textsf{biblatex} already does for \textsf{editor},
\textsf{translator}, and \textsf{namec} (i.e., the compiler).
Furthermore, it is conceivable that a given entry will need separate
translators for each of the three sorts of title.  For this, and for
various other tricky situations, there is the \cmd{parttrans} macro
(and its siblings), designed to be used in a \textsf{note} field or in
one of the \textsf{titleaddon} fields (ratliff:review).  (Because the
strings identifying a translator differ in notes and bibliography, one
can't simply write them out in such a field, hence the need for a
macro, which I discuss further in the commands section below
[\ref{sec:formatcommands}].)

\mylittlespace Finally, as I detailed above under \textbf{author}, in
the absence of an \textsf{author}, \textsf{namea}, \textsf{editor},
and \textsf{nameb}, the \textsf{translator} will be used at the head
of an entry (silver:gawain), and the bibliography entry alphabetized
by the translator's name, behavior that can be controlled with the
\texttt{use<name>} switches in the \textsf{options} field.  Cf.\
\textsf{author}, \textsf{editor}, \textsf{namea}, \textsf{nameb}, and
\textsf{namec}.

\mybigspace This \mymarginpar{\textbf{type}} is a standard
\textsf{biblatex} field, and in its normal usage serves to identify
the type of a \textsf{manual}, \textsf{patent}, \textsf{report}, or
\textsf{thesis} entry.  \textsf{Biblatex} 0.7 introduced the ability,
in some circumstances, to use a bibstring without inserting it in a
\cmd{bibstring} command, and in some entry types (\textsf{audio,
  manual, music, patent, report, suppbook, suppcollection, thesis,}
and \textsf{video}) the \textsf{type} field works this way, allowing
you simply to input, e.g., \texttt{patentus} rather than
\verb+\bibstring{patentus}+, though both will work.  (See
petroff:impurity; herwign:office, murphy:silent, and ross:thesis all
demonstrate how the \textsf{type} field may sometimes be automatically
set in such entries by using one of the standard entry-type aliases).
In other entry types (\textsf{artwork, image, book, online, article,
  review,} and \textsf{suppperiodical}) \textsf{biblatex-chicago} will
merely capitalize the contents according to context.

\mylittlespace In the \textsf{suppbook} entry type, and in its alias
\textsf{suppcollection}, you can use the \textsf{type} field to
specify what sort of supplemental material you are citing, e.g.,
\enquote{\texttt{preface to}} or \enquote{\texttt{postscript to}.}
Cf.\ \textsf{suppbook} above for the details.  (See \emph{Manual}
14.110; polakow:afterw, prose\hc intro).

\mylittlespace You can use the \textsf{type} field in
\textsf{artwork}, \textsf{audio}, \textsf{image}, \textsf{music}, and
\textsf{video} entries to identify the medium of the work, e.g.,
\texttt{oil on canvas}, \texttt{albumen print}, \texttt{compact disc}
or \texttt{MPEG}.  In \textsf{book} entries it will normally hold
system information about multimedia app content (14.268), while in
\textsf{online, article,} and \textsf{review} entries it will hold the
medium of online multimedia (14.267).  Cf.\ under these entry types in
section~\ref{sec:entrytypes}, above, for more details.  (See
auden:reading, bedford:photo, cleese:holygrail, leo:madonna,
nytrumpet:art.)

\mybigspace A standard \mymarginpar{\textbf{url}} \textsf{biblatex}
field, it holds the url of an online publication, though you can
provide one for all entry types.  The \emph{Manual} expresses a strong
preference for DOIs over URLs if the former is available --- cf.\
\textsf{doi} above, and also \textsf{urldate} just below.  The
required \LaTeX\ package \textsf{url} will ensure that your documents
format such references properly, in the text and in the reference
apparatus.  It may be worth noting that child entries generally won't
inherit \textsf{url} fields from their parents --- the information
seems entry-specific enough to warrant a little bit of extra typing if
you need to present the same locator in several entries.  You
\mymarginpar{\texttt{blogurl}} can, however, set the preamble option
\texttt{blogurl} to allow your child comments (\textsf{review}) to
inherit the URL from the parent blog (\textsf{article}).

\mybigspace A standard \mymarginpar{\textbf{urldate}}
\textsf{biblatex} field, it identifies exactly when you accessed a
given url, and is given in \textsc{iso}8601 format.  The \emph{Manual}
prefers DOIs to URLs; in the latter case it allows the use of access
dates, particularly in contexts that require it, but prefers that you
use revision dates, if these are available.  To enable you to specify
which date is at stake, I have provided the \textbf{userd} field,
documented below.  If an entry doesn't have a \textsf{userd}, then the
\textsf{urldate} will be treated as an access date (14.8, 14.12--13,
14.207; evanston:library, grove:sibelius, hlatky:hrt, osborne:poison,
sirosh:visualcortex, wikiped:bibtex).  You can also use the field to
specify a time stamp, should the date alone not be specific enough.
The time stamp follows the date, separated by an uppercase
\enquote{T}, like so: \texttt{yyyy-mm-dd\textbf{T}hh:mm:ss}.  If you
wish to specify the time zone, the \emph{Manual} (10.41) prefers
initialisms like \enquote{EST} or \enquote{PDT,} and these are most
easily provided using the \texttt{urltimezone} field, where you can
provide your own parentheses if so desired (cp.\ 14.191).  Following
the examples in the \emph{Manual}, any \textsf{urldate} will by
default be printed in 24-hour format, though other time stamps use
12-hour format.  The \textsf{biblatex} option \texttt{urltime},
discussed in section~\ref{sec:presetopts}, allows you to change this
in your preamble.

\mylittlespace A \textsf{urldate} time stamp (and
\texttt{urltimezone}) can appear in any entry whatsoever, if you judge
the online source to be the sort that changes rapidly enough for a
time stamp to be necessary (14.207, 14.233; wikiped:bibtex).  You can
stop it printing by setting the new \mymarginpar{\texttt{urlstamp}}
\texttt{urlstamp} option to \texttt{false} in your preamble for the
whole document or for specified entry types, or in the
\textsf{options} field of individual entries.  Please see the
documentation of \textbf{date}, above, and also
table~\ref{ad:date:extras}, below, for more details about time stamps
and other parts of \textsf{biblatex's} enhanced date specifications.
Table~\ref{tab:online:types} contains a summary of the current state
of \textsf{biblatex-chicago's} handling of online materials.

\mybigspace This \mymarginpar{\textbf{urltimezone}} field can, if
necessary, specify the time zone associated with a time stamp given as
part of an \textsf{urldate}.  The \emph{Manual} prefers initialisms
like \enquote{EST} for this purpose, and you can provide parentheses
around it at your discretion (cp.\ 10.41 and 14.191).

\mybigspace A \mymarginpar{\textbf{usera}} supplemental
\textsf{biblatex} field which in certain contexts in
\textsf{biblatex-chicago} will identify the broadcast network when you
cite a radio or television program.  In \textsf{article},
\textsf{periodical}, and \textsf{review} entries with
\textsf{entrysubtype} \texttt{magazine}, it acts almost as a
\enquote{\textsf{journaltitleaddon}} field, and its contents will be
placed, unformatted and between commas, after the
\textsf{journaltitle} and before the \textsf{date}.  In \textsf{video}
entries it comes after the \textsf{eventdate}, i.e., the date of first
broadcast, and is separated from that date by the \cmd{bibstring}
\enquote{\texttt{on}} (14.213, 14.265; american:crime, bundy:macneil,
friends:leia, mayberry:brady).

\mybigspace I \mymarginpar{\textbf{userc}} have implemented this
supplemental \textsf{biblatex} field as part of Chicago's name
cross-referencing system.  (The \enquote{c} part is meant as a sort of
mnemonic for this function, though it's perfectly possible to use the
field in other contexts.)  If you use the \textbf{customc} entry type
to include alphabetized cross-references to other, separate entries in
a bibliography, it is unlikely that you will cite the \textsf{customc}
entry in the body of your text.  Therefore, in order for it to appear
in the bibliography, you have two choices.  You can either include the
entry key of the \textsf{customc} entry in a \cmd{nocite} command
inside your document, or you can place that entry key in the
\textsf{userc} field of another .bib entry that you will be citing.
In the latter case, \textsf{biblatex-chicago} will call \cmd{nocite}
for you, and this method should ensure that there will be at least one
entry in the bibliography to which the cross-reference will point.
(See 14.81--82; creasey:ashe:blast, creasey:morton:hide,
creasey:york:death, lecarre:quest.)

%%\enlargethispage{2\baselineskip}

\mybigspace The \mymarginpar{\textbf{userd}} \textsf{userd} field acts
as a sort of \enquote{\textsf{datetype}} field, allowing you in most
entry types to identify whether a \textsf{urldate} is an access date
or a revision date.  The general usage is fairly simple.  If this
field is absent, then a \textsf{urldate} will be treated as an access
date, as has long been the default in \textsf{biblatex} and in
\textsf{biblatex-chicago}.  If you need to identify it in any other
way, what you include in \textsf{userd} will be printed \emph{before}
the \textsf{urldate}, so phrases like \enquote{\texttt{last modified}}
or \enquote{\texttt{last revised}} are what the field will typically
contain (14.12--13; wikiped:bibtex).  In the absence of a
\textsf{urldate} you can, in most entry types, include a
\textsf{userd} field to qualify a \textsf{date} in the same way it
would have modified a \textsf{urldate}.

\mylittlespace Because of the rather specialized needs of some
audio-visual references, this basic schema changes for \textsf{music}
and \textsf{video} entries.  In \textsf{music} entries where an
\textsf{eventdate} is present, \textsf{userd} will modify that date
instead of any \textsf{urldate} that may also be present, and it will
modify an \textsf{origdate} if it is present and there is no
\textsf{eventdate}.  It will modify a \textsf{date} only in the
absence of the other three.  In \textsf{video} entries it will modify
an \textsf{eventdate} if it is present, and in its absence the
\textsf{urldate}.  In the absence of those two, it can modify a
\textsf{date}.  Please see the documentation of the \textbf{music} and
\textbf{video} entry types, and especially of the \textsf{eventdate},
\textsf{origdate}, and \textsf{urldate} fields, above (14.263--65;
nytrumpet:art).

\mylittlespace In all cases, you can start the \textsf{userd} field
with a lowercase letter, and \textsf{biblatex} will take care of
automatic contextual capitalization for you.

\mybigspace Another \mymarginpar{\textbf{usere}} supplemental
\textsf{biblatex} field, which \textsf{biblatex-chicago} uses
specifically to provide a translated \textsf{title} of a work,
something that may be needed if you deem the original language
unparseable by a significant portion of your likely readership.  The
\emph{Manual} offers two alternatives in such a situation: either you
can translate the title and use that translation in your
\textsf{title} field, providing the original language in
\textsf{language}, or you can give the original title in
\textsf{title} and the translation in \textsf{usere}.  If you choose
the latter, you may need to provide a \textsf{shorttitle} so that the
short note form is also parseable.  Cf.\ \textbf{language}, above.
(See 14.99; kern, weresz.)

\mybigspace See \mymarginpar{\textbf{userf}}
section~\ref{sec:related}, below.

\mybigspace Standard \mymarginpar{\textbf{venue}} \textsf{biblatex}
offers this field for use in \textsf{proceedings} and
\textsf{inproceedings} entries, and after a request from Patrick
Danilevici I have followed suit.  I have also implemented the field in
the \textsf{misc} entry type, both with and without an
\textsf{entrysubtype}, in the \textsf{performance} type, and in the
\textsf{unpublished} type.  In all uses it will normally present the
actual venue of an event, as opposed, e.g., to the
\textsf{origlocation}, which might present where a letter was written
or where an earlier edition was printed.

\mybigspace Standard \mymarginpar{\textbf{version}} \textsf{biblatex}
field, formerly only available in \textsf{artwork}, \textsf{image},
\textsf{misc}, \textsf{music}, and \textsf{patent} entries in
\textsf{biblatex-chicago-notes}, but now also in \textsf{book} and
\textsf{performance} entries.  In most entry types it prints a
localized \enquote{\texttt{version}} string, but there may be
specialist needs in \textsf{artwork} and \textsf{image} entries, so
there you'll need to specify the type of data inside the field itself.
In the \textsf{book} type it is particularly needed for presenting
multimedia app content (14.268).

\mybigspace Standard \mymarginpar{\textbf{volume}} \textsf{biblatex}
field.  It holds the volume of a \textsf{journaltitle} in
\textsf{article} (and some \textsf{review}) entries, and also the
volume of a multi-volume work in many other sorts of entry.  The
treatment and placement of \textsf{volume} information in
\textsf{book}-like entries is rather complicated in the \emph{Manual}
(14.116--22).  In bibliography entries, the \textsf{volume} appears
either before the \textsf{maintitle} or before the publication
information.  In long notes, the same applies, but with the additional
possibility of this information appearing \emph{after} the publication
data, just before page numbers.  In the past, if you wanted the volume
information to appear here, you had to leave that information out of
your .bib entry and give it in the \textsf{pages} or \textsf{postnote}
field.  Now, you can use the \textsf{biblatex-chicago} option
\texttt{delayvolume} \mymarginpar{\texttt{delayvolume}} in your
preamble or in the \textsf{options} field of an entry to ensure that
any \textsf{volume} information that would normally have appeared just
before the publication data in a long note appears after it.

\mylittlespace The \textsf{volume} information in both books and
periodicals, and in both the bibliography and long notes, can appear
\emph{immediately before} the page number(s).  In such a case, the
\emph{Manual} prescribes the same treatment for both sorts of sources,
that is, that \enquote{a colon separates the volume number from the
  page number with no intervening space.}  I have implemented this,
but at the request of Clea~F.\ Rees I have made this punctuation
customizable, using the command \cmd{postvolpunct}.\
\mymarginpar{\cmd{postvolpunct}} By default it prints \cmd{addcolon},
but you can use \verb+\renewcommand{\postvolpunct}{...}+ in your
preamble to redefine it.  Cf.\ \textsf{part}, and the command
documentation in section~\ref{sec:formatcommands}.

\mybigspace Standard \mymarginpar{\textbf{volumes}} \textsf{biblatex}
field.  It holds the total number of volumes of a multi-volume work
(meredith:letters).  If both a \textsf{volume} and a \textsf{volumes}
field are present, as may occur particularly in cross-referenced
entries or in entries using the \texttt{maintitle}
\textsf{relatedtype}, then \textsf{biblatex-chicago} will ordinarily
suppress the \textsf{volumes} field in your references, except in some
instances when a \textsf{maintitle} is present.  If the
\textsf{volume} appears before the \textsf{maintitle}, the option
\texttt{hidevolumes}, \mymarginpar{\texttt{hidevolumes}} set to
\texttt{true} by default, controls whether to print the
\textsf{volumes} field after that title or not.  If it appears after
the \textsf{maintitle}, as with the \textsf{relatedtype} just
mentioned, the same option controls whether to print the
\textsf{volumes} field in close association with the \textsf{volume}.
Set the option to \texttt{false} either in the preamble or in the
\textsf{options} field of your entry to have the \textsf{volumes}
appear in these circumstances.  Cf.\ the option's documentation in
section~\ref{sec:chicpreset}, below.

\mybigspace A \mymarginpar{\textbf{xref}} modified \textsf{crossref}
field provided by \textsf{biblatex}, which prevents inheritance of any
data from the parent entry.  See \textbf{crossref}, above.

\mybigspace Standard \mymarginpar{\textbf{year}} \textsf{biblatex}
field.  It usually identifies the year of publication, though unlike
the \textsf{date} field it allows non-numeric input, so you can put
\enquote{n.d.}\ (or, to be language agnostic,
\verb+\bibstring{nodate}+) here if required, or indeed any other sort
of non-numerical date information.  For many kinds of uncertain and
unspecified dates it is now much simpler to make use of
\textsf{biblatex's} enhanced date specifications in the \textsf{date}
field, instead.  Please see table~\ref{ad:date:extras} for a summary
of how \textsf{biblatex-chicago} implements these enhancements.  Cf.\
bedford:photo, clark:mesopot, leo:madonna, ross:thesis.

\subsubsection{Fields for Related Entries}
\label{sec:related}

As \textsf{biblatex.pdf} puts it (\S~3.4), \enquote{Almost all
  bibliography styles require authors to specify certain types of
  relationship between entries such as \enquote{Reprint of},
  \enquote{Reprinted in,} etc. It is impossible to provide data fields
  to cover all of these relationships and so \textsf{biblatex}
  provides a general mechanism for this using the entry fields
  \textsf{related}, \textsf{relatedtype} and \textsf{relatedstring}.}
Before this mechanism was available \textsf{biblatex-chicago}
attempted to provide a similar but much more limited set of
inter-entry relationships using the \textsf{biblatex} fields
\textsf{origlanguage}, \textsf{origlocation}, \textsf{origpublisher},
\textsf{pubstate}, \textsf{reprinttitle}, and \textsf{userf}.  All of
these still work just as they always have or, I hope, somewhat better
than they always have after many recent bug fixes, but the more
general and more powerful \textsf{biblatex} \texttt{related} mechanism
is also available.  It can provide much of what the older system
provided and a great deal that it couldn't.  What follows is a
field-by-field discussion of the options now available.

\enlargethispage{\baselineskip}

\mybigspace In \mymarginpar{\textbf{origlanguage}} keeping with the
\emph{Manual}'s specifications, I have fairly thoroughly redefined
\textsf{biblatex}'s facilities for treating translations.  The
\textsf{origtitle} field isn't used, while the \textsf{language} and
\textsf{origdate} fields have been press-ganged for other duties.  The
\textsf{origlanguage} field, for its part, retains a dual role in
presenting translations in a bibliography.  The details of the
\emph{Manual}'s suggested treatment when both a translation and an
original are cited may be found below under \textbf{userf}.  Here,
however, I simply note that the introductory string used to connect
the translation's citation with the original's is \enquote{Originally
  published as,} which I suggest may well be inaccurate in a great
many cases, as for instance when citing a work from classical
antiquity, which will most certainly not \enquote{originally} have
been published in the Loeb Classical Library.  Although not, strictly
speaking, authorized by the \emph{Manual}, I have provided another way
to introduce the original text, using the \textsf{origlanguage} field,
which must be provided \emph{in the entry for the translation, not the
  original text} (aristotle:metaphy:trans).  If you put one of the
standard \textsf{biblatex} bibstrings there (enumerated below), then
the entry will work properly across multiple languages.  Otherwise,
just put the name of the language there, localized as necessary, and
\textsf{biblatex-chicago} will eschew \enquote{Originally published
  as} in favor of, e.g., \enquote{Greek edition:} or \enquote{French
  edition:}.  This has no effect in notes, where only the work cited
--- original or translation --- will be printed, but it may help to
make the \emph{Manual}'s suggestions for the bibliography more
palatable.  \textbf{NB:} You can use the \textsf{relatedtype}
\texttt{origpubas} with a customized \textsf{relatedstring} field to
achieve the same ends.

\mylittlespace That was the first usage, in keeping at least with the
spirit of the \emph{Manual}.  I have also, perhaps less in keeping
with that specification, retained some of \textsf{biblatex}'s
functionality for this field.  If an entry doesn't have a
\textsf{userf} field, and therefore won't be combining a text and its
translation in the bibliography, you can also use
\textsf{origlanguage} as the standard styles use it, so that instead
of saying, e.g., \enquote{translated by X,} the entry will read
\enquote{translated from the German by X.}  The \emph{Manual} doesn't
mention this, but it may conceivably help avoid certain ambiguities in
some citations.  As in \textsf{biblatex}, if you wish to use this
functionality, you have to provide \emph{not} the name of the
language, but rather a bibliography string, which may, at the time of
writing, be one of \texttt{american}, \texttt{brazilian},
\texttt{danish}, \texttt{dutch}, \texttt{english}, \texttt{french},
\texttt{german}, \texttt{greek}, \texttt{italian}, \texttt{latin},
\texttt{norwegian}, \texttt{portuguese}, \texttt{spanish}, or
\texttt{swedish}, to which I've added \texttt{russian}.

\mybigspace This \mymarginpar{\textbf{origlocation}} field mainly
serves to help document reprint editions and their corresponding
originals (14.114).  In \textsf{biblatex-chicago} you can provide both
an \textsf{origlocation} and an \textsf{origpublisher} to go along
with the \textsf{origdate}, should you so wish, and all of this
information will be printed in long notes and bibliography.  You can
also use this field in a \textsf{letter} or \textsf{misc} (with
\textsf{entrysubtype}) entry to give the place where a published or
unpublished letter was written (14.111, 14.229).  (Jonathan Robinson
has suggested that the \textsf{origlocation} may in some circumstances
actually be necessary for disambiguation, his example being early
printed editions of the same material printed in the same year but in
different cities.  The new functionality should make this simple to
achieve.  Cf.\ \textsf{origdate} (section~\ref{sec:entryfields}),
\textsf{origpublisher} and \textsf{pubstate}; schweitzer:bach.)
\textbf{NB:} It is impossible to present this same information, as
here, \emph{inside} a single entry using a \textsf{related} field,
though the \textsf{relatedtype} \texttt{origpubin} presents much the
same information \emph{after} the entry, using data extracted from a
separate entry.

\mybigspace As \mymarginpar{\textbf{origpublisher}} with the
\textsf{origlocation} field just above, this field mainly serves to
help document reprint editions and their corresponding originals
(14.114).  You can provide an \textsf{origpublisher} and/or an
\textsf{origlocation} in addition to the \textsf{origdate}, and all
will be presented in long notes and bibliography.  (Cf.\
\textsf{origdate} (section~\ref{sec:entryfields}),
\textsf{origlocation}, and \textsf{pubstate}; schweitzer:bach.)
\textbf{NB:} It is impossible to present this same information, as
here, \emph{inside} a single entry using a \textsf{related} field,
though the \textsf{relatedtype} \texttt{origpubin} presents much the
same information \emph{after} the entry, using data extracted from a
separate entry.

\mybigspace In \mymarginpar{\textbf{pubstate}} response to new
specifications in the 17th edition of the \emph{Manual} (esp.\
14.137), I have tried to generalize the functioning of the
\textsf{pubstate} field in all entry types.  The \texttt{reprint}
string still has a special status there, being ignored in
\textsf{video} entries and provoking a syntactic change in the
presentation of dates in \textsf{music} entries (14.263; floyd:atom),
while in other types allowing the presentation of reprinted titles.
Other strings are divided into two types: those which
\textsf{biblatex-chicago} will print as the \textsf{year}, which
currently means \emph{only} those for which \textsf{biblatex} contains
bibstrings indicating works soon to be published, i.e.,
\texttt{forthcoming}, \texttt{inpreparation}, \texttt{inpress}, and
\texttt{submitted}; and those, i.e., everything else, which will be
printed before, and in close association with, other information about
the publisher of a work.  The four in the first category will always
be localized, as will \texttt{reprint} and \texttt{selfpublished} (and
anything else that \textsf{biblatex} finds to be a \cmd{bibstring})
from the second category.  All other strings will be printed as-is,
capitalized if needed, just before the publisher (author:forthcoming,
contrib:contrib, schweitzer:bach).  \textbf{NB:} The \textsf{pubstate}
functionality currently has no equivalent using the \textsf{related}
field.

% %\enlargethispage{\baselineskip}

\mybigspace This \mymarginpar{\textbf{related}} field is required to
use \textsf{biblatex's} \textsf{related} functionality, and it should
contain the entry key or keys from which \textsf{biblatex} should
extract data for presentation not on its own, but rather in the
bibliography entry (or long note) which contains the \textsf{related}
field itself.  Indeed, unless you change the defaults using the
\textsf{relatedoptions} field this data will only appear in such
entries, never on its own.  Without a \textsf{relatedtype} field, this
will print the default type, equivalent to a long note citation
\emph{immediately after} the bibliography entry containing the
\textsf{related} field, with no intervening string.  You can specify a
string using the \textsf{relatedstring} field, so in effect this
presents a powerful mechanism for presenting full references to
related material of any sort whatsoever.

\mylittlespace By \mymarginpar{\texttt{related=bib}} default, the
package option \texttt{related} is set to print \textsf{related}
entries only in the bibliography.  If you would like them to appear
only in long notes, in both notes and bibliography, or indeed in
neither, you can set this option, either in your preamble (globally or
for specific entry types) or in the \textsf{options} field of the
relevant entry, to \texttt{notes}, \texttt{true}, or \texttt{false},
respectively (coolidge:speech and weed:flatiron).  For the three
\textsf{relatedtypes} that construct a single entry using data
extracted from related entries --- \texttt{commenton},
\texttt{maintitle}, and \texttt{reviewof} ---
\textsf{biblatex-chicago} will automatically set it to \texttt{true}
for you entry by entry, as this is required to get properly-formatted
citations in notes and bibliography.  See below for the details.

\mybigspace This \mymarginpar{\textbf{relatedoptions}} field will, I
should expect, only be needed very rarely.  If you want to set
entry-level options for a \textsf{related} entry this is where you can
do it, though please remember one important detail.  By default,
\textsf{Biber} sets this option to \texttt{dataonly}, which among
other things prevents the \textsf{related} entry from
appearing separately in the bibliography.  If you use the field
yourself, then you'll need to include \texttt{dataonly} as one of the
options therein to maintain this effect.  Of course, it may be you
don't want all the effects of \texttt{dataonly}, so you can tailor it
however you wish.  See \textsf{biblatex.pdf} \S~3.4.

\mybigspace The \mymarginpar{\textbf{relatedstring}} procedure for
choosing a string to connect the main entry with its related entry/ies
is straightforward, the default being a \texttt{bibstring}, if any,
with the same name as the \textsf{relatedtype}, or alternately a
string or strings defined within the driver for that
\textsf{relatedtype}, as happens with the types \texttt{origpubin} and
\texttt{bytranslator}.  Failing these, you can supply your own in the
\textsf{relatedstring} field, either in the form of the name of a
pre-defined \texttt{bibstring} or as any text you choose, and anything
in this field always takes precedence over the automatic choices.  If
your non-\texttt{bibstring} starts with a lowercase letter then
\textsf{biblatex-chicago} will capitalize it automatically for you
depending on context (weed:flatiron).  I have not altered the standard
\textsf{relatedtype} strings, and have indeed changed the
\textsf{reprinttitle} mechanism to use the \texttt{reprintfrom}
string, which works better syntactically in this context.

\mybigspace The \mymarginpar{\textbf{relatedtype}} standard
\textsf{biblatex} styles define six \textsf{relatedtypes}, and I have
either simply adopted them wholesale or adapted them to the needs of
the Chicago style, retaining the basic syntax as much as possible.  I
have also added four to these six (see below).  First, the standard
types:

\begin{description}
\item[\qquad bytranslator:] This prints a full citation of a
  translation, starting with the (localized) string
  \enquote{Translated by \textsf{translator} as \textsf{Title},
    \ldots} The reference is fuller in \textsf{biblatex-chicago} than
  in the standard styles, and for the first time allows users to
  choose the \emph{Manual's} alternate method for presenting original
  + translation (14.99; furet:related).  The old \textsf{userf}
  mechanism provides the other, as does the \texttt{origpubas}
  \textsf{relatedtype} (see below).
\item[\qquad default:] This is the macro used when no
  \textsf{relatedtype} is defined.  It prints, as in the standard
  styles, and with no intervening string, a full citation of
  \textsf{related} entries.  In \textsf{biblatex-chicago-notes}, the
  citation is in long note form, rather than bibliography form, as
  this is the usual practice in the \emph{Manual}.
\item[\qquad multivolume:] This briefly lists the individual volumes
  in a multi-volume work, and works much as in the standard styles.
  The \emph{Manual}, as far as I can see, has little to say on the
  matter.
\item[\qquad origpubas:] This type can, if you want, replace the old
  \textsf{userf} mechanism, described below, for presenting an
  original with its translation.  It's quite similar to the
  \texttt{default} type, but with a \texttt{bibstring} automatically
  connecting the entry with its \textsf{related} entries.  You can
  identify other sorts of relationships if you change the introductory
  string using \textsf{relatedstring}.
\item[\qquad origpubin:] I haven't altered this from the
  \textsf{biblatex} default at all, and it presents reprint
  information \emph{after} the main entry rather than within it.  The
  \emph{Manual} seems to prefer the latter for the notes \&
  bibliography style and, in some circumstances, the former for
  author-date.
\item[\qquad reprintfrom:] This type provides a replacement for the
  old \textsf{reprinttitle} mechanism described below.  As in the
  standard styles, it presents a fuller reference to the reprinted
  material than does \texttt{origpubin}, and is designed particularly
  for presenting pieces formerly printed in other collections or
  perhaps essays collected from various periodicals.  (In
  \textsf{biblatex-chicago} it contains some kludges to cope with
  possible \textsf{babel} language environments, so if you find it
  behaving oddly please let me know, including whether you are using
  \textsf{babel} [which I've tested] or \textsf{polyglossia} [which
  I've tested somewhat less].)
\end{description}
Now, the Chicago-specific types:
\begin{description}
\item[\qquad commenton:] I designed this \textsf{relatedtype} to
  facilitate citation of online comments, though it works slightly
  differently in the two entry types in which it is available,
  \textsf{online} and \textsf{review} (with its clone
  \textsf{suppperiodical}) (14.208--10).  In both types it allows you
  to mimic thread structure by citing a chain of replies to comments
  on posts, etc., all in a single entry, while also simplifying your\
  .bib entries.  This simplification works differently depending on
  whether the comment itself has no specific title, as always in
  \textsf{review} entries, or does have such a title, as especially in
  \textsf{online} social media entries.  In the former case, as you
  can see from ac:comment, the \textsf{related} apparatus allows you,
  and indeed encourages you, not to provide a \textsf{title} at all,
  as contrasted with the old system, still available of course, where
  your \textsf{title} field contained special formatting for the title
  of the blog on which this entry is a comment.  Note also here the
  \textsf{eventdate} field, a requirement, with its optional time
  stamp, which helps to differentiate multiple comments by the same
  author posted on the same day.  The \textsf{options} field can be
  used to prevent the entry appearing in the bibliography, and you can
  also provide a \textsf{url} specific to the comment, though this is
  by no means necessary.

  In social media threads comments and replies may well have their own
  title, so in such a case you still need a \textsf{title} field,
  which will be followed by the \textsf{relatedstring}.  In such
  \mymarginpar{NB!}  \textsf{online} entries, the \emph{only way to
    cite these comments} is by using the \texttt{commenton}
  \textsf{relatedtype} (licis:diazcomment).  Note that, unlike
  \textsf{review} entries, the date, and possible time stamp, of a
  comment should appear in the \textsf{date} field, not
  \textsf{eventdate}.  (Other fields, like \textsf{url} and
  \textsf{options}, have much the same uses as in \textsf{review}
  entries.)  If, in other \textsf{online} entries, you decided
  \emph{not} to use \textsf{commenton} in an entry like braun:reply,
  and simply use a specially-crafted \textsf{titleaddon} field, you
  lose the possibility of having two dates in the entry, one for the
  comment and one for the original post, though to be fair it does end
  up looking like the example in 14.210, where it is ambiguous to
  which part of the citation the date applies.

  As for the thread structure, I've not tested how far down the rabbit
  hole you can go, but a series of entries linked one to the next by
  this \textsf{relatedtype} will all turn up if you cite the first in
  the chain, though of course you can use the technique merely as a
  convenient way to structure and simplify your\ .bib file, without
  creating chains longer than 2 entries.  The default connecting
  string is the localized \texttt{commenton}, but you can use
  \textsf{relatedstring} to change it to \enquote{\texttt{reply to}}
  or whatever else you need.  I've tried to follow the rules for
  abbreviating parts of the various works included in the one
  reference, though in truth the \emph{Manual} provides no examples.
  Depending on whether the various parts have already been cited or
  not your references can take on quite varied appearances.  Let me
  know if something looks wrong to you.  Cf.\ ac:comment,
  diaz:surprise, ellis:blog, and licis:diazcomment for the use of the
  new \textsf{relatedtype}; amlen:hoot, amlen:wordplay, and viv:amlen
  for blogs and comments without the \textsf{related} mechanism.

  There are a few other things to remember.  As with the next two
  \textsf{relatedtypes}, \textsf{biblatex-chicago} will automatically
  set the \texttt{related} option to \texttt{true} entry-by-entry to
  ensure that the full data appears both in notes and in the
  bibliography.  If your parent entry has no \textsf{title} of its
  own, then, as with the \texttt{reviewof} \textsf{relatedtype}, it
  will use the \textsf{related} functionality also in short notes,
  which means that if you want to provide a \textsf{shorttitle} for
  them then it goes in the \emph{child} entry rather than the parent.
  Finally, the \textsf{title}-less comments are prime candidates for
  the \texttt{shortextrafield} option, which prints a disambiguating
  field after short notes when they would otherwise be
  indistinguishable.  So endemic is this situation in this context
  that I've set a default means of disambiguating them, which is the
  \textsf{date} and \textsf{time} in \textsf{online} entries or the
  \textsf{eventdate} and \textsf{eventtime} in \textsf{review} and
  \textsf{suppperiodical}, though you can of course override these
  defaults by setting the \texttt{shortextra} options yourself.  See
  their documentation in section~\ref{sec:useropts}, below.
\item[\qquad maintitle:] The 17th edition of the \emph{Manual} has
  deployed, in at least two contexts, a notable syntactic change in
  the presentation of works that form part of other, larger works.
  Generally, the order of presentation, in \textsf{biblatex} terms,
  has always been \textsf{title} -- \textsf{booktitle} --
  \textsf{maintitle}, in increasing order of generality.  In the vast
  majority of cases this order still holds, but in TV episodes, for
  one example, the recommendation now is to present the name of the
  series (\textsf{booktitle}) \emph{before} the name of the episode
  (\textsf{title}). (See the \textbf{video} type in
  section~\ref{sec:entrytypes}, above).  The other context in which
  this reversal occurs is multi-volume works (14.116--22).  Here, the
  preferred format, at least for notes, appears to be
  \textsf{maintitle} -- \textsf{[book]title} or, when all three titles
  are present, \textsf{title} -- \textsf{maintitle} --
  \textsf{booktitle}.  The \emph{Manual} doesn't carry this reordering
  through with absolute consistency, but I think it important at least
  to offer it as a possibility to users of \textsf{biblatex-chicago},
  hence the \texttt{maintitle} \textsf{relatedtype}, which is
  currently the only way to achieve this reversal in this context.

  In its simplest usage, to document one volume of a multi-volume set,
  you would have, e.g., an \textsf{mvcollection} entry with
  \textsf{relatedtype} \texttt{maintitle}, and a \textsf{related}
  field pointing to a \textsf{collection} entry.  When you cite the
  \textsf{mvcollection} entry itself, you'll get a long note like
  \emph{MVCollTitle}, vol.\ 1, \emph{CollTitle}, and a short note like
  \emph{MVCollTitle}, vol.\ 1., or, with a \textsf{postnote} field,
  \emph{MVCollTitle}, 1:12, as the specification requires.  If you
  wanted to cite one essay in the \textsf{collection}, then you would,
  additionally, need an \textsf{incollection} entry with the
  \texttt{maintitle} \textsf{relatedtype} and a \textsf{related} field
  pointing to the \textsf{mvcollection} entry already mentioned, so
  you're creating a chain of three different related entries but
  presenting them in one reference, i.e., (long form)
  \enquote{InCollTitle,} in \emph{MVCollTitle}, vol.\ 1,
  \emph{CollTitle}, and (short form) \enquote{InCollTitle.}  It's
  important to keep in mind here that, in effect, you're \emph{not}
  actually citing the \textsf{mvcollection} entry, but the one volume
  of it represented by the \textsf{collection} entry, or indeed an
  essay in that one volume.

  Now, for the details, which are many.  First, any \textsf{mv*} entry
  without the \texttt{maintitle} \textsf{relatedtype} should behave
  just as it always has, and can still happily be used as the target
  of \textsf{crossref} fields to supply a \textsf{maintitle} to other
  entries.  The abbreviated references created when you have several,
  e.g., \textsf{books crossref'd} to the same \textsf{mvbook} are
  still available, assuming you enable them with the
  \texttt{booklongxref} option.  You can happily mix the new and the
  old methods of presentation in your documents, but please don't mix
  them within individual entries, which means that if you are using a
  \textsf{crossref} field to an \textsf{mvcollection} entry in a
  \textsf{collection} entry, say, and the \textsf{collection} entry is
  itself the target of the \textsf{mvcollection} entry's
  \textsf{related} field, please be careful not to cite that
  \textsf{collection} entry independently, as it can lead to
  unexpected results.  (If things don't look right to you, try
  eliminating the use of \textsf{crossref} entirely from these
  \textsf{related} chains and see if that helps, then send me a bug
  report if it does.)  This restriction also means that, although the
  \emph{Manual} prefers the \textsf{maintitle}-first format in notes
  and allows either syntax in the bibliography, nonetheless with
  \textsf{biblatex-chicago} whichever syntax you choose for the notes
  will also appear in the bibliography.

  As for automatically abbreviating references using the
  \texttt{maintitle} \textsf{relatedtype}, this works differently
  depending on whether the \textsf{related} chain consists of 2 or 3
  works.  In 2-work chains (\emph{MVCollTitle}, \emph{CollTitle}),
  it's actually the first that needs abbreviating, and this didn't
  look right, so these entries will always print in full.  (You can
  still regulate how much information appears in the references to
  individual volumes by regulating how much information appears in the
  .bib entries for those volumes.  In the harleymt:* entries I've used
  as examples below, the individual volumes have a \textsf{crossref}
  field to the multi-volume work, so they inherit the
  \textsf{publisher} and \textsf{location}, for example.  If you were
  to omit the \textsf{crossref} field you would always get an
  abbreviated reference which, were it to appear \emph{after} a
  reference to the whole multi-volume work, would let that reference
  give the complete publication details and itself behave like a
  normal abbreviated cross reference to it.)

  In 3-work chains, when you've cited more than one
  \enquote{InCollTitle} from a given \emph{CollTitle}, you can choose
  for the short note version of the second and third titles (with just
  volume number rather than full \emph{CollTitle}) to appear in the
  bibliography and in long notes after the first one.  This is
  controlled using the same \texttt{booklongxref} option as you would
  use to control the old automatic abbreviation mechanism.  See under
  that option in section~\ref{sec:chicpreset}, below.

  As with the other two \textsf{relatedtypes} I've added to
  \textsf{biblatex's} standard six, the \texttt{maintitle} type is
  somewhat restricted in its relevance.  If you want to use a
  three-work chain to cite one part of one volume, then this is
  possible only by starting with the following entry types:
  \textsf{bookinbook}, \textsf{inbook}, \textsf{incollection},
  \textsf{inproceedings}, and \textsf{letter}.  All two-work chains
  must start with one of the \textsf{mv*} types.  As might be apparent
  from the previous list, \textsf{mvreference} entries are special, in
  that their \textsf{related} field should point to an
  \textsf{inreference} entry if you want to cite an entry in an
  \enquote{alphabetically arranged work}, or to a \textsf{reference}
  entry otherwise.  In other words, \textsf{mvreference} entries
  should only ever be used in 2-work chains.

  It's possible it may have occurred to you that this
  \textsf{relatedtype} could, given the presence of a many-volumed
  collection, require rather a lot of extra entries in your\ .bib
  files, i.e., one extra \textsf{mv*} entry for every volume of the
  collection you wish to cite.  Borrowing an idea from the
  \texttt{multivolume} \textsf{relatedtype}, you can put the entry
  keys of \emph{all} the individual volumes into a single
  \textsf{related} field in a single \textsf{mv*} entry, and
  \textsf{biblatex-chicago} will still allow you to cite each volume
  independently, and for each to appear independently in the
  bibliography, too, unlike the \texttt{multivolume} mechanism.
  Here's how it works.  When \textsf{Biber} detects more than one
  entrykey in an \textsf{mv*} entry with \texttt{maintitle}
  \textsf{relatedtype}, it produces a series of clones of the
  \textsf{mv*} entry, each with the same \textsf{relatedtype} and a
  \textsf{related} field containing exactly \emph{one} of the
  entrykeys from the original \textsf{related} field.  It gives each
  of these clones its own entrykey, of the form
  \texttt{mventrykey-singlevolumeentrykey}, and it is these virtual,
  cloned entries that you should cite.  Such entries don't exist in
  your\ .bib file, but you can see them in your\ .bbl file, assuming
  you've actually cited any of them.  The original
  \texttt{mventrykey}, in this case, refers merely to the original
  \textsf{mv*} entry, as though it had never had a \textsf{related}
  field, so it's available for citing the multi-volume set as a whole,
  should that be necessary.  Indeed, to make the virtual clones
  available to \textsf{Biber} (and \textsf{biblatex}) in the first
  place, you do need to cite (or \cmd{nocite}) the original
  \textsf{mv*} entry somewhere in your document.

  As an example of how this might look, consider the three entries
  from \textsf{notes-test.bib}: harleymt:hoc, harleymt:ancient:cart,
  and harleymt:cartography.  The first, an \textsf{mvcollection}
  entry, has a \textsf{related} field containing both of the others
  (\textsf{collection} entries), and in \textsf{cms-notes-sample.tex}
  you'll see citations of harleymt:hoc,
  harleymt:hoc-harleymt:ancient:cart, and
  harleymt:hoc-harleymt:cartography, which are themselves
  \textsf{mvcollection} entries.  The latter two don't exist in the\
  .bib file, only in the\ .bbl file, where you'd see that each has a
  \textsf{related} field pointing to the entrykey that forms the
  second half of its own hyphenated key.  If I hadn't somewhere cited
  harleymt:hoc then \textsf{Biber} would give up entirely because it
  wouldn't know where to find the two hyphenated keys.

  A similar problem arises when you create a three-work chain in which
  the first, e.g., \textsf{incollection}, entry contains a
  \textsf{related} field pointing to just such a virtual, cloned
  entry.  In this case, if you haven't already cited (or
  \cmd{nocite}'d) the cloned entry, \textsf{Biber} really gets,
  understandably, confused.  As a convenience feature for this
  situation, I have included a very slightly modified version of the
  \texttt{maintitle} \textsf{relatedtype}, called
  \mymarginpar{\texttt{maintitlenc}} \texttt{maintitlenc}, the
  \enquote{nc} standing for \cmd{nocite}.  As you might have guessed,
  every clone produced by an \textsf{mv*} entry with multiple
  entrykeys in its \textsf{related} field and \texttt{maintitlenc} as
  its \textsf{relatedtype} will automatically be \cmd{nocite}'d, and
  will then be available for inclusion in another entry's
  \textsf{related} field.  The \texttt{maintitlenc} type differs in no
  other way whatever from the \texttt{maintitle} type.

  In general, the \textsf{maintitle} \textsf{relatedtype} attempts to
  follow the Chicago specification with as little intervention needed
  from the user as possible.  To that end, \textsf{biblatex-chicago}
  automatically sets the \texttt{related} option to \texttt{true}
  entry-by-entry to ensure that the full data appears both in notes
  and in the bibliography.  It also attempts to spot duplicate
  \textsf{authors} or \textsf{editors} and to print them only when
  needed, and in its short-note version uses the \textsf{volume} and
  \textsf{part} information from the \textsf{related}
  \textsf{collection} entry, say, to specify the \textsf{labeltitle}
  which comes from the \textsf{mvcollection}.  If you want the
  \textsf{mv*} entry's \textsf{volumes} data to appear in notes and
  bibliography, you can do so by setting the \texttt{hidevolumes}
  option to \texttt{false} either in the preamble or in the
  \textsf{options} field of the entry referenced by the \textsf{mv*}
  entry's \textsf{related} field (cf.\ harleymt:cartography).

  Another, trickier intervention involves the problem of sorting
  entries in the bibliography.  The \emph{Manual's} rules are,
  basically, to sort by name, then title, then year, and as every
  \textsf{mv*} entry citing the same multi-volume work will basically
  have identical values for all three, the sorting order in the
  bibliography will fall back on the order in which such works are
  cited, which may not be what you want.  If the pertinent
  \textsf{related} fields in your\ .bib file only contain one
  entrykey, then you can use a series of \textsf{sortkey},
  \textsf{sorttitle}, or \textsf{sortyear} fields to arrange the
  volumes as you wish.  If you are using \textsf{Biber's} cloning
  facilities, however, any such fields in the \textsf{mv*} entry will
  be copied into all of its clones, so you'll be back to square one.
  My current solution to this dilemma \mymarginpar{\textsf{sorttitle}}
  is to treat the \textsf{sorttitle} field as special, so that in
  \textsf{mv*} entries with \textsf{relatedtype} \texttt{maintitle} or
  \texttt{maintitlenc} and a \textsf{related} field containing more
  than one entrykey any \textsf{sorttitle} field is indeed copied to
  all of the clones, but it is also modified by appending the contents
  of the clone's \textsf{related} field to the end of it.  (If there
  is no \textsf{sorttitle} field, then the clones will have none,
  either.)  In effect, the alphabetical order of the entrykeys in the
  \textsf{mv*} entry's \textsf{related} field determines the sorting
  order of the clones each of which contains a \textsf{related} field
  holding one of those keys.  (In the examples from
  \textsf{notes-test.bib}, harleymt:hoc retains its original
  \textsf{sorttitle} field, and sorts first, then
  harleymt:hoc-harleymt:ancient:cart sorts before
  harleymt:hoc-harleymt:cartography --- the second half of the key,
  after the hyphen, is the relevant part, and is what appears appended
  to the original \textsf{sorttitle} field.)  It is rather onerous, I
  know, to be required to choose entrykeys that sort properly; if I
  come up with something better I'll include it in a later release.

  Finally, although I've tested this functionality extensively, it's
  new and rather complicated.  If something doesn't work right please
  let me know at the email address at the head of this documentation.
\item[\qquad reviewof:] Philip Kime's \textsf{biblatex-apa} package
  includes this type, and user Bertold Schweitzer suggested it might
  be a useful addition to \textsf{biblatex-chicago}, so I've added it
  to the standard six detailed above.  It differs from all of them,
  and resembles \texttt{commenton} and \texttt{maintitle}, in that it
  prints the \textsf{relatedstring} (\verb+\bibstring{reviewof}+ by
  default) and the data from the \textsf{related} entry in the middle
  of the parent entry, rather than at the end.  It differs from
  \texttt{commenton} in that it's not possible to create a chain of
  such entries to mimic online thread structures.  Finally, it differs
  from all other \textsf{relatedtypes} in being available only in
  \textsf{article} and \textsf{review} entries (along with the
  latter's clone, \textsf{suppperiodical}).

  In \textsf{article} entries it replaces the \textsf{titleaddon} with
  the \textsf{relatedstring} followed by the \textsf{title} of the
  child entry, and in \textsf{review} entries it replaces the
  \textsf{title} with the same two components.  In both types these
  components will optionally be followed by the \textsf{author},
  \textsf{editor}, \textsf{translator}, etc.,\ of the reviewed item,
  and then any child \textsf{titleaddon} may optionally appear at the
  end, allowing maximum flexibility when presenting, for example,
  reviews of live performances.

  This mechanism automates both the provision of the localized
  \cmd{bibstring} and also the formatting of the \textsf{title} of the
  reviewed work, and it also obviates the need to use any of the
  \cmd{partedit} macros in this context.  Further,
  \textsf{biblatex-chicago} automatically sets the \texttt{related}
  option to \texttt{true} entry-by-entry to ensure that the full data
  appears both in notes and in the bibliography.  Finally, this
  \textsf{relatedtype} has the further peculiarity that, in
  \textsf{review} and \textsf{suppperiodical} entries only, it uses
  the \textsf{related} functionality also in short notes, which means
  that if you want to provide a \textsf{shorttitle} for short notes
  then it goes in the \emph{child} entry rather than the parent.
  Please remember, too, that the standard way of presenting reviewed
  works is still available if the mechanism doesn't work for you in a
  particular context.
\item[\qquad short:] This \textsf{relatedtype} is like the
  \texttt{default} type, only it prints \emph{short} references rather
  than long ones. There is no default \textsf{relatedstring} for this
  type, so if you leave that field blank then the short references
  will simply appear at the end of the long note or bibliography
  entry.
\end{description}

\mybigspace \textbf{NB:} \mymarginpar{\textbf{reprinttitle}}
\textbf{If you have been using this feature, you may want to have a
  look at the} \textsf{relatedtype} \texttt{reprintfrom},
\textbf{documented above, for a better solution to this problem, one
  that also allows you to change the introductory string using the}
\textsf{relatedstring} \textbf{field.  The} \textsf{reprinttitle}
\textbf{field will continue to work as before, however.}  At the
request of Will Small, I have included a means of providing the
original publication details of an essay or a chapter that you are
citing from a subsequent reprint, e.g., a \emph{Collected Essays}
volume.  In such a case, at least according to the \emph{Manual}
(14.181), such details needn't be provided in notes, only in the
bibliography, and then only if these details are \enquote{of specific
  interest.}  The data would follow an introductory phrase like
\enquote{originally published as,} making the problem strictly
parallel to that of including details of a work in the original
language alongside the details of its translation.  I have addressed
the latter problem with the \textsf{userf} field, which provides a
sort of cross-referencing method for this purpose, and
\textsf{reprinttitle} works in \emph{exactly} the same way.  In the
.bib entry for the reprint you include a cross-reference to the cite
key of the original location using the \textsf{reprinttitle} field
(which it may help mnemonically to think of as a \enquote{reprinted
  title} field).  The main difference between the two forms is that
\textsf{userf} prints all but the \textsf{author} of the original
work, whereas \textsf{reprinttitle} suppresses both the
\textsf{author} and the \textsf{title} of the original, giving only
the more general details, beginning with, e.g., the
\textsf{journaltitle} or \textsf{booktitle} and continuing from there.
The string prefacing this information will be \enquote{Reprinted
  from.}  Please see the documentation on \textbf{userf} below for all
the details on how to create .bib entries for presenting your data.

\mybigspace This \mymarginpar{\textbf{userf}} is one of the
supplemental fields which \textsf{biblatex} provides, and is used by
\textsf{biblatex-chicago} for a very specific purpose.  When you cite
both a translation and its original, the \emph{Manual} (14.99)
recommends that, in the bibliography at least, you combine references
to both texts in one entry, though the presentation in notes is pretty
much up to you.  In order to follow this specification, I have
provided a third cross-referencing system (the others being
\textsf{crossref} and \textsf{xref}), and have chosen the name
\textsf{userf} because it might act as a mnemonic for its function.

\mylittlespace In order to use this system, you should start by
entering both the original and its translation into your .bib file,
just as you normally would.  The mechanism works for any entry type,
and the two entries need not be of the same type.  In the entry for
the \emph{translation}, you put the cite key of the original into the
\textsf{userf} field.  In the \emph{original's} entry, you need to
include something that will prevent the entry from being printed
separately in the bibliography --- \texttt{skipbib} in the
\textsf{options} field will work, as would something in the
\textsf{keywords} field in conjunction with a \texttt{notkeyword=}
switch in the \cmd{printbibliography} command.  In this standard case,
the data for the translation will be printed first, followed by the
string \texttt{originally published as}, followed by the original,
author omitted, in what amounts to the same format that the
\emph{Manual} uses for long footnotes (furet:passing:eng,
furet:passing:fr).  As explained above (\textbf{origlanguage}), I have
also included a way to modify the string printed before the original.
In the entry for the \emph{translation}, you put the original's
language in \textsf{origlanguage}, and instead of \texttt{originally
  published as}, you'll get \texttt{French edition:} or \texttt{Latin
  edition:}, etc.\ (aristotle:metaphy:gr, aristotle:metaphy:trans).
\textbf{NB:} You can use the \textsf{relatedtype} \texttt{origpubas}
to replicate the \textsf{userf} functionality, and you can also
customize the \textsf{relatedstring} field to achieve the same result
as with \textsf{origlanguage}.


\subsection{Commands}
\label{sec:commands}

In this section I shall attempt to document all those commands you may
need when using \textsf{biblatex-chicago-notes} that I have either
altered with respect to the standard provided by \textsf{biblatex} or
that I have provided myself.  Some of these, unfortunately, will make
your .bib file incompatible with other \textsf{biblatex} styles, but
I've been unable to avoid this.  Any ideas for more elegant, and more
compatible, solutions will be warmly welcomed.

\subsubsection{Formatting Commands}
\label{sec:formatcommands}

These commands allow you to fine-tune the presentation of your
references in both notes and bibliography.  You can find many examples
of their usage in \textsf{notes-test.bib}, and I shall try to point
you toward a few such entries in what follows.  \textbf{NB:}
\textsf{biblatex's} \cmd{mkbibquote} command is mandatory in some
situations.  See its entry below.

\mybigspace Version \mymarginpar{\textbf{\textbackslash autocap}} 0.8
of \textsf{biblatex} introduced the \cmd{autocap} command, which
capitalizes a word inside a note or bibliography entry if that word
follows sentence-ending punctuation, and leaves it lowercase
otherwise.  As this command is both more powerful and more elegant
than the kludge I designed for a previous version of
\textsf{biblatex-chicago-notes} (see\ \textbf{\textbackslash
  bibstring} below), you should be aware that the use of the
single-letter \cmd{bibstring} commands in your .bib file is obsolete.

\mylittlespace In order somewhat to reduce the burden on users even
further, I have, following \textsf{biblatex's} example, implemented a
system which automatically tracks the capitalization of certain fields
in your .bib file.  I chose these fields after a non-scientific survey
of entries in my own databases, so of course if you have ideas for the
extension of this facility I would be most interested to hear them.
In order to take advantage of this functionality, all you need do is
begin the data in the appropriate field with a lowercase letter,
e.g.,\ \texttt{note = \{with the assistance of X\}}.  If the data
begins with a capital letter --- and this is not infrequent --- that
capital will always be retained.  (cf., e.g., creel:house,
morgenson:market.)  If, on the other hand, you for some reason need
such a field always to start with a lowercase letter, then you can try
putting an empty set of curly braces\ \{\}\ at the start, which turns
off the mechanism without printing anything itself.  Here, then, is
the complete list of fields where this functionality is active:

\begin{enumerate}
\setlength{\parskip}{-4pt}
\item The \textbf{addendum} field in all entry types.
\item The \textbf{booktitleaddon} field in all entry types.
\item The \textbf{edition} field in all entry types.  (Numerals work
  as you expect them to here.)
\item The \textbf{maintitleaddon} field in all entry types.
\item The \textbf{note} field in all entry types.
\item The \textbf{part} field in entry types that use it.
\item The \textbf{prenote} field prefixed to citation commands.
\item The \textbf{relatedstring} field in all entry types.
\item The \textbf{shorttitle} field in the \textsf{review}
  (\textsf{suppperiodical}) entry type and in the \textsf{misc} type,
  in the latter case, however, only when there is an
  \textsf{entrysubtype} defined, indicating that the work cited is
  from an archive.
\item The \textbf{title} field in the \textsf{review}
  (\textsf{suppperiodical}) entry type and in the \textsf{misc} type,
  in the latter case, however, only when there is an
  \textsf{entrysubtype} defined, indicating that the work cited is
  from an archive.
\item The \textbf{titleaddon} field in all entry types.
\item The \textbf{type} field in \textsf{artwork}, \textsf{audio},
  \textsf{image}, \textsf{music}, \textsf{suppbook},
  \textsf{suppcollection}, and \textsf{video} entry types.
\end{enumerate}

In any other cases --- and there are only two examples of this in
\textsf{notes-test.bib} (centinel\hc letters, powell:email) --- you'll
need to provide the \cmd{autocap} command yourself.  Indeed, if you
accidentally do so in one of the above fields, it shouldn't matter at
all, and you'll still get what you want, but taking advantage of the
automatic provisions should at least save some typing.

\mybigspace This \mymarginpar{\textbf{\textbackslash bibstring}} is a
very powerful mechanism to allow \textsf{biblatex} automatically to
provide a localized version of a string, and to determine whether that
string needs capitalization, depending on where it falls in an entry.
\textsf{Biblatex 0.7} introduced functionality which sometimes allows
you simply to input, for example, \texttt{newseries} instead of
\verb+\bibstring{newseries}+, the package auto-detecting when a
bibstring is involved and doing the right thing, though in all such
cases either form will work.  This functionality is available in the
\textsf{series} field of \textsf{article}, \textsf{jurisdiction},
\textsf{legislation}, \textsf{periodical}, and \textsf{review}
entries; in the \textsf{type} field of \textsf{audio},
\textsf{manual}, \textsf{music}, \textsf{patent}, \textsf{report},
\textsf{suppbook}, \textsf{suppcollection}, \textsf{thesis}, and
\textsf{video} entries; in the \textsf{location} field of
\textsf{patent} entries; in the \textsf{language} field in all entry
types; in the \textsf{nameaddon} field in \textsf{customc} entries;
and in the \textsf{editor[abc]type} and \textsf{nameatype} fields in
all entry types.  There may be other places where \textsf{biblatex's}
standard styles support this feature, and I shall add them when they
come to my attention.

\enlargethispage{\baselineskip}

\mybigspace These \mymarginpar{\textbf{\textbackslash
    foottextcite\\\textbackslash foottextcites}} two commands look
like citation commands, but are in fact wrappers for customizing the
behavior of the \cmd{textcite} and \cmd{textcites} citation commands
when they are used inside a foot- or an endnote.  By default, in such
a context these commands print the name of the \textsf{author(s)}
followed by the \emph{short} citation or citations, i.e., usually
\textsf{title} only, enclosed within parentheses.  You can change the
way the citation part is presented by using \cmd{renewcommand} in your
preamble.  The default definitions are: \texttt{\{\textbackslash
  addspace\textbackslash headlessparenshortcite\}} \&
\texttt{\{\textbackslash addspace\textbackslash
  headlessparenshortcites\}}.  If you wanted to return to the default
behavior of previous releases of \textsf{biblatex-chicago} you could
change the first to: \texttt{\{\textbackslash newcunit\textbackslash
  bibstring\{in\}\textbackslash addspace\textbackslash
  headlesscite\}}, and the second similarly, only using
\cmd{headlesscites}.  (There is also, by the way, a
\cmd{headlessparencite(s)} command if you want to retain the long
citations inside the parentheses.)

\mybigspace I \mymarginpar{\textbf{\textbackslash letterdatelong}}
have provided this macro mainly for use in the optional postnote field
of the various citation commands.  When citing a letter (published or
unpublished, \textsf{letter} or \textsf{misc}), it may be useful to
append the date to the usual short note form in order to disambiguate
references.  This macro simply prints the date of a letter, or indeed
of any other sort of correspondence, in day-month-year order, as
recommended by the \emph{Manual} (14.224).  (If your main document
language isn't American, it's better just to use the standard
\textsf{biblatex} command \cmd{printorigdate}.)

\mybigspace This \mymarginpar{\textbf{\textbackslash mkbibquote}} is
the standard \textsf{biblatex} command, which requires attention here
because it is a crucial part of the mechanism for the
\enquote{American} punctuation system.  If you look in
\textsf{chicago-notes.cbx} you'll see that the quoted fields, e.g., an
\textsf{article} or \textsf{incollection title}, have this command in
their formatting, which does most of the work for you.  If, however,
you need to provide additional quotation marks in a field --- a quoted
title within a title, for example --- then you may need to use this
command so that any following period or comma will be brought within
the closing quotation marks.  Its use is \emph{required} when the
quoted material comes at the end of a field, and I recommend always
using it in your .bib database, as it does no harm even when that
condition is not fulfilled.  A few examples from
\textsf{notes-test.bib} should help to clarify this.

\mylittlespace In an \textsf{article} entry, the \textsf{title}
contains a quoted phrase:

\begin{verbatim}
  title = {Diethylstilbestrol and Media Coverage of the
    \mkbibquote{Morning After} Pill}
\end{verbatim}

Here, because the quoted text doesn't come at the end of title, and no
punctuation will ever need to be drawn within the closing quotation
mark, you could instead use \texttt{\cmd{enquote}\{Morning After\}} or
even \texttt{`Morning After'}. (Note the single quotation marks here
--- the other two methods have the virtue of taking care of nesting
for you.)  All of these will produce the formatted
\enquote{Diethylstilbestrol and Media Coverage of the \enquote{Morning
    After} Pill.}  Here, by contrast, is a \textsf{book title}:

\begin{verbatim}
  title = {Annotations to \mkbibquote{Finnegans Wake}}
\end{verbatim}

Because the quoted title within the title comes at the end of the
field, and because this bibliographical unit will be separated from
what follows by a period in the bibliography, then the
\cmd{mkbibquote} command is necessary to bring that period within the
final quotation marks, like so: \emph{Annotations to
  \enquote{Finnegans Wake.}}

\mylittlespace Let me also add that this command interacts well with
Lehman's \textsf{csquotes} package, which I highly recommend, though
the latter isn't strictly necessary in texts using an American style,
to which \textsf{biblatex} defaults when \textsf{csquotes} isn't
loaded.

\mybigspace The \mymarginpar{\textbf{\textbackslash postvolpunct}}
\emph{Manual} (14.116) unequivocally prescribes that when a
\textsf{volume} number appears immediately before a page number,
\enquote{the abbreviation \emph{vol.}\ is omitted and a colon
  separates the volume number from the page number with no intervening
  space.}  The treatment is basically the same whether the citation is
of a book or of a periodical, and it appears to be a surprising and
unwelcome feature for many users, conflicting as it may do with
established typographic traditions in a number of contexts.  Clea~ F.\
Rees requested a way to customize this, so I have provided the
\cmd{postvolpunct} command, which prints the punctuation between a
\textsf{volume} number and a page number.  It is set to \cmd{addcolon}
by default, except when the current language of the entry is French,
in which case it defaults to \verb+\addcolon\addspace+.  You can use
\verb+\renewcommand{\postvolpunct}{...}+ in your preamble to redefine
it, but please note that the command only applies in this limited
context, not more generally to the punctuation that appears between,
e.g., a \textsf{volume} and a \textsf{part} field.

\mybigspace This \mymarginpar{\textbf{\textbackslash reprint}} and the
following 7 macros all help \textsf{biblatex-chicago-notes} cope with
the fact that many bibstrings in the Chicago system differ between
notes and bibliography, the former sometimes using abbreviated forms
when the latter prints them in full.  In the current case, if a book
is a reprint, then the macro \cmd{reprint}, followed by a comma, could
go in the \textsf{location} field before the city of publication.
Simply putting \enquote{\texttt{reprint}} into the \textsf{pubstate}
field is a simpler way to achieve the same result
(aristotle:metaphy:gr, schweitzer:bach).  See \textsf{location} and
\textsf{pubstate} in section~\ref{sec:entryfields}, above.

\mylittlespace \textbf{NB:} The rules for employing abbreviated or
full bibstrings in the \emph{Manual} are remarkably complex, but I
have attempted to make them as transparent for users as possible.  In
\textsf{biblatex-chicago-notes}, if you don't see it mentioned in this
section, then in theory you should always provide an abbreviated
version, using the \cmd{bibstring} mechanism, if necessary
(babb:peru).  The standard \textsf{biblatex} bibstrings should also
work (palmatary:pottery), and any that won't should be covered by the
series of macros beginning here with \cmd{reprint} and ending below
with \cmd{parttransandcomp}.

\mybigspace Since \mymarginpar{\textbf{\textbackslash partcomp}} the
\emph{Manual} specifies that the strings \texttt{editor},
\texttt{translator}, and \texttt{compiler} all require different forms
in notes and bibliography, and since it mentions these three apart
from all the others \textsf{biblatex} provides (\textsf{annotator},
\textsf{commentator}, et al.), and further since it may indeed happen
that the available fields (\textsf{editor}, \textsf{namea},
\textsf{translator}, \textsf{nameb}, and \textsf{namec}) aren't
adequate for presenting some entries, I have provided 7 macros to
allow you to print the correct strings for these functions in both
notes and bibliography.  Their names all begin with \cmd{part}, as
originally I intended them for use when a particular name applied only
to a specific \textsf{title}, rather than to a \textsf{maintitle} or
\textsf{booktitle} (cf.\ \textbf{namea} and \textbf{nameb}, above).

\mylittlespace In the present instance, you can use \cmd{partcomp} to
identify a compiler when \textsf{namec} won't do, e.g., in a
\textsf{note} field or the like.  In such a case,
\textsf{biblatex-chicago-notes} will print the appropriate string in
your references.

\mybigspace Use \mymarginpar{\textbf{\textbackslash partedit}} this
macro when identifying an editor whose name doesn't conveniently fit
into the usual fields (\textsf{editor} or \textsf{namea}).  (N.B.: If
you are writing in French then you no longer need to add either
\texttt{de} or \texttt{d'} after this command in your .bib files.  The
new version of the command should take care of this automatically for
you.)  See chaucer:liferecords.

\mybigspace As \mymarginpar{\textbf{\textbackslash
    partedit-\\andcomp}} before, but for use when an editor is also a
compiler.

\vspace{1.3\baselineskip} As \mymarginpar{\textbf{\textbackslash
    partedit-\\andtrans}} before, but for when when an editor is also a
translator (ratliff:review).

\vspace{1.3\baselineskip} As \mymarginpar{\textbf{\textbackslash
    partedit-\\transandcomp}} before, but for when an editor is also a
translator and a compiler.

\vspace{1.4\baselineskip} As \mymarginpar{\textbf{\textbackslash
    parttrans-\\andcomp}} before, but for when a translator is also a
compiler.

\vspace{1.3\baselineskip} As \mymarginpar{\textbf{\textbackslash
    parttrans}} before, but for use when identifying a translator
whose name doesn't conveniently fit into the usual fields
(\textsf{translator} and \textsf{nameb}).

\mylittlespace Unlike \mymarginpar{\textbf{\textbackslash
    suppress-\\bibfield[]\{\}}} the other commands presented here,
this should be used in your document preamble rather than in your
bibliographical apparatus.  Also unlike them, it has two arguments,
the first of which is optional, the second required.  Jan David Hauck
suggested that, in addition to the field-exclusion package options
provided by \textsf{biblatex-chicago} (see
section~\ref{sec:chicpreset}), I might also provide a general-purpose
macro to clear fields from selected entry types when the package
options aren't quite right for a user's particular needs.  The
\cmd{suppressbibfield} command does this, so that
\verb+\suppressbibfield{note}+ clears the \textsf{note} field from
\emph{all} entries, while \verb+\suppressbibfield[report]{note}+
clears it only from \textsf{report} entries.  Both arguments take
comma-separated lists, so to suppress \textsf{titleaddon} and
\textsf{volumes} fields from \textsf{report} and \textsf{manual}
entries, your preamble could contain
\verb+\suppressbib-field[report,manual]{titleaddon,volumes}+.

\mylittlespace A few usage notes are in order.  First, you can use as
many calls to the command in your preamble as you wish.  Second, the
command is a very basic user interface to \textsf{biblatex's} source
mapping functionality (\textsf{biblatex.pdf} \S~4.5.3), so what it
does is modify what \textsf{biber} takes from your .bib file in order
to produce the .bbl file that \textsf{biblatex} actually reads.  As
far as \textsf{biblatex} is concerned, the fields simply aren't there
in the data source, so they can't appear anywhere in the
bibliographical apparatus, whether in notes, bibliography or shorthand
lists.  Third, because source mapping is involved, you'll need a
complete cycle of \LaTeX-\textsf{biber}-\LaTeX\ runs to make the
commands take effect.  Fourth, source mapping occurs at a very early
stage in \textsf{biber's} operation, so if your field names or entry
types are standard aliases, the command will only work on the names as
they appear in your .bib file, not as they are aliased in the .bbl
file.  If you have a \textsf{techreport} entry, for example, it won't
be affected by a command that alters \textsf{report} entries, and a
\textsf{date} field won't be affected by a command that suppresses the
\textsf{year}.  Fifth, the code for the command resides in
\textsf{biblatex-chicago.sty}, so if you load the styles without
loading that package it won't be available to you.  Sixth and finally,
the \cmd{suppressbibfield} command is new and relatively untested, so
please report any untoward behavior to me.

\subsubsection{Citation Commands}
\label{sec:citecommands}

The \textsf{biblatex} package is particularly rich in citation
commands, some of which (e.g., \cmd{supercite(s)}, \cmd{citeyear})
provide functionality that isn't really needed by the Chicago notes
and bibliography style offered here.  If you are getting unexpected
behavior when using them please have a look in your .log file ---
there may be warnings there that alert you to undefined citation
commands.  Other \textsf{biblatex}-provided commands, though I haven't
tested them extensively, should pretty much work out of the box.  What
remains are the commands I have found most useful and necessary for
following the \emph{Manual}'s specifications, and I document in this
section any alterations I have made to these.  As always, if there are
standard commands that don't work for you, or new commands that would
be useful, please let me know, and it should be possible to fix or add
them.

\mylittlespace A number of users have run into a problem that appears
when they've used a command like \cmd{cite} inside a \cmd{footnote}
macro.  In this situation, the automatic capitalization routines will
not be in operation at the start of the footnote, so instead of
\enquote{Ibid.,} for example, you'll see \enquote{ibid.}  If you need
to use the \cmd{cite} command within a \cmd{footnote} command, the
solution is to use \cmd{Cite} instead.  Alternatively, don't use a
\cmd{footnote} macro at all, rather try \cmd{footcite} or
\cmd{autocite} with the optional prenote and postnote arguments.  Cf.\
\cmd{Citetitle} below, and also section~3.8 of \textsf{biblatex.pdf}.

\mybigspace I \mymarginpar{\textbf{\textbackslash autocite}} haven't
adapted this in the slightest, but I thought it worth pointing out
that \textsf{biblatex-chicago-notes} sets this command to use
\cmd{footcite} as the default option. It is, in my experience, much
the most common citation command you will use, and also works fine in
its multicite form, \textbf{\textbackslash autocites}.

% %\enlargethispage{\baselineskip}

\mybigspace While \mymarginpar{\textbf{\textbackslash cite*}} the
\cmd{cite} command works just as you would expect it to, I have also
provided a starred version for the rare situations when you might need
to turn off the ibidem tracking mechanism.  \textsf{Biblatex} provides
very sophisticated algorithms for this mechanism, so in general you
won't find a need for this command, but in case you'd prefer a longer
citation where you might automatically find the shortest one, I've
provided this.  Of course, you'll need to put it inside a
\cmd{footnote} command manually.  (See also section
\ref{sec:useropts}, below.)

\mybigspace I \mymarginpar{\textbf{\textbackslash citeauthor}} have
adapted this standard \textsf{biblatex} command only very slightly to
bring it into line with \textsf{biblatex-chicago's} needs.  Its main
usage will probably be for references to works from classical
antiquity, when an \textsf{author's} name (abbreviated or not)
sometimes suffices in the absence of a \textsf{title}, e.g., Thuc.\
2.40.2--3 (14.244).  You'll need to put it inside a \cmd{footnote}
command manually.  (Cf.\ also \textsf{entrysubtype} in
section~\ref{sec:entryfields}, above, and please note that the new
\mycolor{\texttt{notitle}} option (section~\ref{sec:useropts})
produces the same effect without the need to worry about citation
commands.)

\mybigspace Although \mymarginpar{\textbf{\textbackslash citeincite}}
the officially-sanctioned, and safest, way to present cross-references
to other works is by using the \textsf{related} mechanism, it would
appear, judging from various bug reports I've received, that there are
plenty of users who need to present such references in ways that
aren't supported by the \textsf{related} code as it stands. This new
citation command is designed to allow you to present short references
to other works inside other references, and to avoid some, I hope
most, of the bugs associated with using cite commands in the fields of
other entries. If you want to present long citations, I'm assuming
that you wouldn't want to do that in fields right in the middle of the
parent entry, and that, therefore, the usual methods detailed in
section~\ref{sec:related} (above) or the \cmd{fullciteincite} command
(below) will serve your needs.  Similarly, if you want to provide
short citations at the end of your parent entry, the new
\textsf{relatedtype} \texttt{short} should work.

\mylittlespace If, however, you must have a short reference in the
midst of another field, then \cmd{citeincite} can help. It is
certified to work properly only in the \textsf{addendum, annotation,
  annote, note, and titleaddon} fields, which I hope is a broad enough
choice to cover most needs. The multicite command \cmd{citeincites} is
also available. In previous releases of \textsf{biblatex} there were
differences in how the case-changing backends behaved with respect to
citation commands inside entry fields, but in my testing the two
backends appear currently to be equivalent on this score. Still,
should you run into problems using \cmd{citeincite}, it might still be
worth setting \texttt{casechanger=latex2e} in your preamble to see if
that helps.

\mybigspace This \mymarginpar{\textbf{\textbackslash citejournal}}
command provides an alternative short form when citing journal
\textsf{articles}, giving the \textsf{journaltitle} and
\textsf{volume} number instead of the article \textsf{title} after the
\textsf{author's} name.  The \emph{Manual} suggests that this format
might be helpful \enquote{in the absence of a full bibliography}
(14.185).  It may also prove useful when you want to provide
parenthetical references to newspaper articles within the text rather
than in the bibliography, a style endorsed by the \emph{Manual}
(14.198).  In such a case, an article's author, if there is one, could
form part of the running text.  As usual with these general citation
commands, if you want the reference to appear in a footnote you need
to put it inside a \cmd{footnote} command manually.

\mybigspace This \mymarginpar{\textbf{\textbackslash Citetitle}}
simply prepends \cmd{bibsentence} to the usual \cmd{citetitle}
command.  Some titles may need this for the automatic contextual
capitalization facility to work correctly.  (Included as standard from
\textsf{biblatex} 0.8d.)

\mybigspace Joseph \mymarginpar{\textbf{\textbackslash citetitles}}
Reagle noticed that, because of the way
\textsf{biblatex-chicago-notes} formats titles in quotation marks,
using the \cmd{citetitle} command will often get you punctuation you
don't want, especially when presenting a list of titles.  I've
included this multicite command to enable you to present such a list,
if the need arises.  Remember that you'll have to put it inside a
\cmd{footnote} command manually.

\mybigspace Another \mymarginpar{\textbf{\textbackslash footfullcite}}
standard \textsf{biblatex} command, modified to work properly with
\textsf{biblatex-chicago-notes}, and provided in case you find
yourself in a situation where you really need the full citation in a
footnote, but where \cmd{autocite} would print a short note or even
engage the ibidem mechanism.  This may be particularly useful if
you've chosen to use all short notes by setting the \texttt{short}
option in the arguments to \cmd{usepackage\{biblatex\}}, yet still
feel the need for the occasional full citation.

\mybigspace This, \mymarginpar{\textbf{\textbackslash fullcite}} too,
is a standard command, and it too provides a full citation, but unlike
the previous command it doesn't automatically place it in a footnote.
It may be useful within long textual notes.

\mybigspace Ordinarily, \mymarginpar{\textbf{\textbackslash
    fullciteincite}} you can use one of the methods discussed in
section~\ref{sec:related} to present, in notes or bibliography, full
references to a work related in some way to the current entry. The new
command \cmd{citeincite} (above) allows you to present short citations
to other works within selected fields of the parent entry, while
attempting to minimize the problems caused by citations within other
citations. The \cmd{fullciteincite} command does the same for full
citations, but really is designed to work only in \textsf{annotation}
fields, the idea being that you might want to refer in an annotated
bibliography to works that you haven't cited in the main body of your
text. The multicite command \cmd{fullciteincites} is also available.

\mybigspace Arne \mymarginpar{\textbf{\textbackslash gentextcite}}
Skj{\ae}rholt requested, for the author-date styles, a variant of the
\cmd{textcite} command that presented the author's name in the
genitive case in running text, thereby simplifying certain syntactic
constructions.  As a convenience for users, I've also ported this to
the notes \& bibliography style.  In most respects it behaves exactly
like \cmd{textcite}, on which see below.  The difference is that I've
added a new optional field to the front of the command to allow you to
choose which declensional ending to add to the name.  If you don't
specify this field, you'll get the standard English \enquote{\,'s\,}.
If you want something different, then you'll need to present a third
option to the command, like so:
\cmd{gentextcite[<ending>][][]\{entry:key\}}.  You must include the
two further sets of square brackets, because with only one set it
will, as with other citation commands, be interpreted as a
\textsf{postnote}, and with two a \textsf{prenote} and a
\textsf{postnote}.  There is a \cmd{gentextcites} command as well, and
for it you may need to specify
\cmd{gentextcites[<ending>]()()[][]\hfill\{entry:key1\}\{entry:key2\}},
though if you don't have a \textsf{pre-} or \textsf{postnote} to the
first citation you can make do with
\cmd{gentextcites[<ending>]()\{entry:key1\}\{entry:key2\}}.

\mylittlespace The syntax of multiple authors' names in running text
is unpredictable.  There is currently no way to add the genitival
ending to all the names attached to a single citation key, so it will
only appear at the end of a group of names in such a case.  (This is
in keeping with the usual syntax when referring to a multi-author
work, at least in English.)  When using \cmd{gentextcites}, however,
you can control whether the ending appears after the name(s) attached
to each citation key, or whether it only appears after the names
attached to the last key.  By default, it only appears after the last,
but the \texttt{genallnames} preamble and/or entry option
set to \texttt{true} will attach the ending to each key's name(s).

\mybigspace Matthew \mymarginpar{\textbf{\textbackslash
    headlesscite}\\\textbf{\textbackslash
    headlesscites}\\\textbf{\textbackslash Headlesscite}}
Lundin requested a more generalized \cmd{headlesscite} macro,
suppressing the author's name in specific contexts while allowing
users not to worry about whether a particular citation needs the long
or short form, a responsibility thereby handed over to
\textsf{biblatex's} tracking mechanisms.  These citation commands
attempt to fulfill this request.  The (new) capitalized command, as
usual, will ensure capitalization of, e.g., \enquote{ibid.}\ at the
beginning of notes, and was made necessary by fixes to a bug
identified by David Purton.  Please note that, in the short form, the
result will be rather like a \cmd{citetitle} command, which may or may
not be what you want.  Note, also, that as I have provided only the
most flexible form of the command, you'll have to wrap it in a
\cmd{footnote} yourself.  Please see the next entry for further
discussion of some of the needs this command might help address.

\mybigspace I \mymarginpar{\textbf{\textbackslash
    headless-\\fullcite}\\\textbf{\textbackslash
    headless-\\fullcites}} have provided these commands in case you
want to print a full citation without the author's name.  The
\emph{Manual} (14.78, 14.104) suggests this for brevity's sake in
cases where that name is already obvious enough from the title, and
where repetition might seem awkward (creel:house, feydeau:farces,
meredith:letters, and sewall:letter).  \textsf{Letter} entries and
\textsf{standard} entries (where the \textsf{author} is usually the
publishing \textsf{organization}) --- and only these entries --- do
this for you automatically, and of course the repetition is tolerated
in bibliographies for the sake of alphabetization, but in notes for
other entry types this command may help achieve greater elegance, even
if it isn't strictly necessary.  As I've provided only the most
flexible form of the command, you'll have to wrap it in a
\cmd{footnote} yourself.

\mybigspace I \mymarginpar{\textbf{\textbackslash shortcite}} have
provided this command in case, for any reason, you specifically
require the short form of a note, and \textsf{biblatex} thinks you
want something else.  Again, I've provided only the most flexible form
of the command, so you'll have to wrap it in a \cmd{footnote}
manually.

\mybigspace The \mymarginpar{\textbf{\textbackslash shortcite*}}
starred version of the command turns off page and citation tracking
for a short note, designed particularly to prevent a \texttt{noteref}
back reference from appearing, should you and the code have a
disagreement over just when such back references might be necessary.

\mybigspace At \mymarginpar{\textbf{\textbackslash shorthandcite}}
the request of Kenneth Pearce, I have included this command which
always prints the \textsf{shorthand}, even at the first citation of a
given work.  Again, I've only provided the most flexible form of the
command, so you'll need to place it inside parentheses or wrap it in a
\cmd{footnote} manually.

\mybigspace The \mymarginpar{\textbf{\textbackslash shorthandcite*}}
starred version of the command turns off page and citation tracking
for a \textsf{shorthand} note, designed particularly to prevent a
\texttt{noteref} back reference from appearing, should you and the
code have a disagreement over just when such back references might be
necessary.

\mybigspace This \mymarginpar{\textbf{\textbackslash
    shorthand-\\refcite}} command, like the next, forces the printing
of a back reference when you are using the new package option
\texttt{noteref}, only it prints a \textsf{shorthand} note rather than
a short note.  It's the opposite of \cmd{shorthandcite*}.

\mybigspace The \mymarginpar{\textbf{\textbackslash shortrefcite}} new
package option \texttt{noteref} provides for the printing of back
references from short notes to their corresponding long notes.
\textsf{Biblatex-chicago} provides several options to help you control
when such back references should appear, but as there may well be
occasions when you and the code disagree, this command forces the
printing of a short note with a back reference.  It's the opposite,
then, of \cmd{shortcite*}, which prevents such a back reference from
appearing.

\mybigspace This \mymarginpar{\textbf{\textbackslash surnamecite}}
command is analogous to \cmd{headlesscite}, but whereas the latter
allows you to omit an \textsf{author's} name when that name is obvious
from the \textsf{title} of a work, \cmd{surnamecite} allows you to
shorten a full note citation in contexts where the full name(s) of the
\textsf{author} have already been provided in the text.  In short
notes this falls back to the standard format, but in long notes it
simply omits the given names of the \textsf{author} and provides only
the surname, along with the full data of the entry (cf.\ 14.57).

\mybigspace Norman \mymarginpar{\textbf{\textbackslash textcite}} Gray
started a discussion on
\href{http://tex.stackexchange.com/questions/67837/citations-as-nouns-in-biblatex-chicago}{Stackexchange}
which established both that \textsf{biblatex} had begun including a
\cmd{textcite} command in its verbose styles and that
\textsf{biblatex-chicago-notes} hadn't kept up.  In that thread Audrey
Boruvka provided some code, adapted from \textsf{verbose.cbx}, to
provide such a command for the Chicago notes \&\ bibliography style.
More recently, Rasmus Pank Rouland pointed out some changes in
\textsf{biblatex} that made the \cmd{textcites} command fit more
elegantly into the flow of text.  I've adapted this solution in this
release.  I'm still not entirely certain how best to accommodate this
request within the package, but there are now at least commands
(\cmd{textcite} and \cmd{textcites}) for users to test.  Their
functionality is a little complicated.  In the main text, they will
provide an \textsf{author's} name(s), followed immediately by a foot-
or endnote which contains the full (or short) reference, following the
usual rules.  If you use \cmd{textcite} inside a foot- or endnote,
then the default behavior, for both \cmd{textcite} and
\cmd{textcites}, specifies that you'll get the \textsf{author's}
name(s) followed by a headless \emph{short} citation (or citations)
placed within parentheses.  Such parentheses are generally discouraged
by the \emph{Manual} (14.38), but are nonetheless somewhat better than
other solutions for smoothing the syntax of sentences that include
such a citation.  I have made the citation short, i.e., \textsf{title}
only, because this again seems likely to be the least awkward solution
syntactically.  If you want to configure this behavior for either
citation command, please see \cmd{foottextcite} and
\cmd{foottextcites} in section~\ref{sec:formatcommands}.

\mylittlespace If you look at \textsf{chicago-notes.cbx}, you'll see a
number of other citation commands, but those are intended for internal
use only, mainly in cross-references of various sorts.  Use at your
own risk.

\subsection{Package Options}
\label{sec:options}

\subsubsection{Pre-Set \textsf{biblatex} Options}
\label{sec:presetopts}

Although a quick glance through \textsf{biblatex-chicago.sty} will
tell you which \textsf{biblatex} options the package sets for you, I
thought I might gather them here also for your perusal.  These
settings are, I believe, consistent with the specification, but you
can alter them in the options to \textsf{biblatex-chicago} in your
preamble or by loading the package via \verb+\usepackage+
\verb+[style=chicago-notes]{biblatex}+, which gives you the
\textsf{biblatex} defaults unless you redefine them yourself inside
the square brackets.

\mylittlespace By \mymarginpar{\texttt{abbreviate=\\false}} default,
\textsf{biblatex-chicago-notes} prints the longer bibstrings, mainly
for use in the bibliography, but since notes require the shorter forms
of many of them, I've had to define many new strings for use there.

\mylittlespace \textsf{Biblatex-chicago-notes}
\mymarginpar{\texttt{autocite=\\footnote}} places references in
footnotes by default.

\mybigspace The \mymarginpar{\texttt{citetracker=\\true}} citetracker
for the \cmd{ifciteseen} test is enabled globally.

\mybigspace The \mymarginpar{\texttt{alldates=comp}} specification calls
for the long format when presenting dates, slightly shortened when
presenting date ranges.

\mylittlespace In \mymarginpar{\texttt{alltimes=12h}} entries which
print time stamps, they will, when the stamp is part of a
\textsf{date}, \textsf{eventdate}, or \textsf{origdate}, appear in
12-hour format, i.e., \enquote{4:45 p.m.}  Stamps that are part of a
\textsf{urldate} are, by default, controlled by the \texttt{urltime}
option, which is set to \texttt{24h}.  See that option below, and
table~\ref{ad:date:extras}.

\mylittlespace The \mymarginpar{\texttt{dateabbrev=\\false}}
\emph{Manual} prefers to use full month names in the notes \&\
bibliography style.

\mybigspace This \mymarginpar{\texttt{datecirca=true}} option enables
\textsf{biblatex's} enhanced \enquote{circa} date specification, which
given a \textsf{date} like \verb+1989~+ will print [ca.\ 1989].  Cf.\
table~\ref{ad:date:extras}.

\mylittlespace This \mymarginpar{\texttt{dateuncertain=\\true}} option
enables \textsf{biblatex's} enhanced \enquote{uncertain} date
specification, which given a \textsf{date} like \verb+1989?+ will
print [1989?].  A field like \verb+1989%+ is both \enquote{circa}
\emph{and} \enquote{uncertain,} like so: [ca.\ 1989?].  Cf.\
table~\ref{ad:date:extras}.

\mylittlespace This \mymarginpar{\texttt{datezeros=false}} ensures
that leading zeros don't appear in date specifications.

\mybigspace This \mymarginpar{\texttt{ibidtracker=\\constrict}}
enables the use of the ibidem mechanism in notes, but only in the most
strictly-defined circumstances.  Whenever there might be any
ambiguity, \textsf{biblatex} should default to printing a more
informative reference.  Remember also that you can use the \cmd{cite*}
command to disable this functionality in any given reference, or
indeed one of the \texttt{fullcite} commands if you need the long note
form for any reason.  Please see section~\ref{sec:useropts} for
options (\texttt{noibid} and \texttt{useibid}) managing how
\textsf{biblatex-chicago} presents ibidem references, as the defaults
have changed in the 17th edition (14.34).

\mylittlespace Roger \mymarginpar{\texttt{idemtracker=\\false}} Hart
suggested that it might be helpful, despite the \emph{Manual's}
objections (14.35), to be able to turn on \textsf{biblatex's}
\texttt{idemtracker}.  This replaces, in notes only, authors' names
with the string \enquote{Idem} when a work by the same author follows
a different work by that author, two consecutive references to the
same work by the same author generating, of course, \enquote{Ibid.}
Indeed, if you are going to use the \texttt{idemtracker}, you should
also set the package option \texttt{useibid} to \texttt{true}, so that
you don't get a mix of \enquote{Idem} and the new 17th-edition ibidem
behavior which doesn't print \enquote{Ibid.}  You can turn this all on
when loading \textsf{biblatex-chicago} by setting
\texttt{idemtracker=constrict,useibid=true}.  It works very much like
the standard \textsf{biblatex} styles which include this option, so
that you never get \enquote{Idem} in long notes, but only in short
ones, and (ideally) never when the repeated name might be somewhat
ambiguous.  Also, if you wish the localized string to be appropriately
gendered, you need to employ the \textsf{biblatex} field
\textsf{gender}, on which see \textsf{biblatex.pdf} \S~2.3.3.

\mylittlespace This \mymarginpar{\texttt{labeltitle=\\true}} option
enables \textsf{biblatex-chicago-notes} to disambiguate shortened
citations to different sources which might otherwise be confusingly
identical.  Though I've set it in \textsf{biblatex-chicago.sty},
you'll need to set the \texttt{shortextrafield} option yourself in
order for it to have any visible effect.  See the documentation of
that option in section~\ref{sec:useropts}, below.

\mylittlespace This \mymarginpar{\texttt{loccittracker\\=constrict}}
allows the package to determine whether two consecutive citations of
the same source also cite the same page of that source.  In such a
case, \texttt{Ibid} alone will be printed, without the page reference,
following the specification (14.29).

\mylittlespace These \mymarginpar{\textsf{\texttt{maxbibnames\\=10\\
      minbibnames\\=7}}} two options control the number of names
printed in the bibliography when that number exceeds 10.  These
numbers follow the recommendations of the \emph{Manual} (14.76), and
they are different from those for use in notes.  With
\textsf{biblatex} 1.6 you can no longer redefine \texttt{maxnames} and
\texttt{minnames} in the \cmd{printbibliography} command at the bottom
of your document, so \textsf{biblatex-chicago} now does this
automatically for you, though of course you can change them in your
document preamble.

\mylittlespace This \mymarginpar{\texttt{pagetracker=\\true}} enables
page tracking for the \cmd{iffirstonpage} and \cmd{ifsamepage}
commands for controlling, among other things, the printing of
\enquote{Ibid.}  It tracks individual pages if \LaTeX\ is in oneside
mode, or whole spreads in twoside mode.

\mylittlespace This \mymarginpar{\texttt{related=bib}} is the standard
\textsf{biblatex} bibliography option, but I have given it some extra
settings and also added entry and type options as well.  By default it
enables the use of \textsf{related} functionality in the bibliography
only, not in long notes.  You can set it either in the preamble or in
individual entries to enable the functionality in long notes only
(\texttt{notes}), in both notes and bibliography (\texttt{true}), or
in neither (\texttt{false}).  When you use the \texttt{commenton},
\texttt{maintitle}, or \texttt{reviewof} \textsf{relatedtypes},
\textsf{biblatex-chicago} automatically sets this option to
\texttt{true} on an entry-by-entry basis, as these
\textsf{relatedtypes} require this for proper functioning.  Cf.\
coolidge:speech, weed:flatiron.

\mylittlespace This \mymarginpar{\texttt{sortcase=\\false}} turns off
the sorting of uppercase and lowercase letters separately, a practice
which the \emph{Manual} doesn't appear to recommend.

\mylittlespace This \mymarginpar{\texttt{sorting=\\}\cmd{cms@choose}}
new setting tests whether you are using \textsf{Biber} as your
backend, and if so enables a custom \textsf{biblatex-chicago} sorting
scheme for the bibliography (\texttt{cms}).  If you are using any
other backend, it reverts to the \textsf{biblatex} default
(\texttt{nty}).  Please see the discussion of
\cmd{DeclareSortingTemplate} just below.

\mylittlespace If \mymarginpar{\texttt{timezones=true}} you provide a
timezone for a time stamp, usually using one of the \textsf{timezone}
fields, this option ensures it will be printed.

\mylittlespace This \colmarginpar{\texttt{uniquename=\\minfull}}
enables the package to distinguish, in short notes, different authors
who share a surname, using initials in the first instance, and whole
names if initials aren't enough (14.32).

\mylittlespace In \mymarginpar{\texttt{urltime=24h}} entries with
\textsf{urldate} fields containing time stamps, that stamp will by
default appear in 24-hour format, i.e., \enquote{16:45.}  Cf.\
\texttt{alltimes}, above, \texttt{urlstamp} in
section~\ref{sec:chicpreset} below, and table~\ref{ad:date:extras}.

\mylittlespace In
\mymarginpar{\texttt{[standard]\\useeditor=false\\usenamec=false}}
\textsf{standard} entries any editors' or compilers' names appear
after the title, according to 14.259, so these entry-type-specific
options encode this.  You can, of course, override these defaults in
your preamble, should you deem it necessary.

\mylittlespace This \mymarginpar{\texttt{usetranslator\\=true}}
enables automatic use of the \textsf{translator} at the head of
entries in the absence of an \textsf{author} or an \textsf{editor}.
In the bibliography, the entry will be alphabetized by the
translator's surname.  You can disable this functionality on a
per-entry basis by setting \texttt{usetranslator=false} in the
\textsf{options} field.  Cf.\ silver:gawain.

\subsubsection*{Other \textsf{biblatex} Formatting Options}
\label{sec:formatopts}

I've chosen defaults for many of the general formatting commands
provided by \textsf{biblatex}, including the vertical space between
bibliography items and between items in the list of shorthands
(\cmd{bibitemsep} and \cmd{lositemsep}).  I define many of these in
\textsf{biblatex-chicago.sty}, and of course you may want to redefine
them to your own needs and tastes.  It may be as well you know that
the \emph{Manual} does state a preference for two of the formatting
options I've implemented by default: the 3-em dash as a replacement
for repeated names in the bibliography (14.67--71, and just below);
and the formatting of note numbers, both in the main text and at the
bottom of the page / end of the essay (superscript in the text,
in-line in the notes; 14.24).  The code for this last formatting is
also in \textsf{biblatex-chicago.sty}, and I've wrapped it in a test
that disables it if you are using the \textsf{memoir} class, which I
believe has its own commands for defining these parameters.  You can
also disable it by using the \texttt{footmarkoff} package option, on
which see below.

\mylittlespace Gildas Hamel pointed out that my default definition, in
\textsf{biblatex-chicago.sty}, of \textsf{biblatex's}
\cmd{bibnamedash} didn't work well with many fonts, leaving a line of
three dashes separated by gaps.  He suggested an alternative, which
I've adopted, with a minor tweak to make the dash thicker, though you
can toy with all the parameters to find what looks right with your
chosen font.  The default definition is:

\mylittlespace\verb+\renewcommand*{bibnamedash}{\rule[.4ex]{3em}{.6pt}}+.

\mylittlespace At \mymarginpar{\texttt{losnotes}
  \&\\\texttt{losendnotes}} the request of Kenneth Pearce, I have
added two \texttt{bibenvironments} to \textsf{chicago-notes\break
  .bbx}, for use with the \texttt{env} option to the
\cmd{printshorthands} command.  The first, \texttt{losnotes}, is
designed to allow a list of shorthands to appear inside footnotes,
while \texttt{losendnotes} does the same for endnotes.  Their main
effect is to change the font size, and in the latter case to clear up
some spurious punctuation and white space that I see on my system when
using endnotes.  (You'll probably also want to use the option
\texttt{heading=none} in order to get rid of the [oversized] default,
providing your own within the \cmd{footnote} command.)  If you use a
command like \verb+\printbiblist{shortjournal}+ to print a list of
journal abbreviations, you can use the \texttt{sjnotes} and
\texttt{sjendnotes} \texttt{bibenvironments} in exactly the same way.
Please see the documentation of \textsf{shorthand} and
\textsf{shortjournal} in section~\ref{sec:entryfields} above for
further options available to you for presenting and formatting these
two types of \texttt{biblist}.

% %\enlargethispage{\baselineskip}

\mylittlespace Formerly
\mymarginpar{\cmd{Declare-}\\\texttt{Labelname}} available only to
those using \textsf{Biber}, but since version 3.0 handled by
\textsf{biblatex} itself, \cmd{DeclareLabelname} allows you to add
name fields for consideration when \textsf{biblatex} is attempting to
find a shortened name for short notes.  This, for example, allows a
compiler (=\textsf{namec}) to appear at the head of short notes
without any other intervention from the user, rather than requiring a
\textsf{shortauthor} field as previous releases of
\textsf{biblatex-chicago} did.  The default search order for the
Chicago styles is \textsf{shortauthor}, \textsf{author},
\textsf{shorteditor}, \textsf{namea}, \textsf{editor}, \textsf{nameb},
\textsf{translator}, \textsf{namec}.  You can set the option
\texttt{use<name>=false} in entries or when loading
\textsf{biblatex-chicago} to exclude individual fields from appearing
in short notes, or indeed at the head of long notes and bibliography
entries.  See the documentation of those name types in
section~\ref{sec:entryfields} for further details.

\mylittlespace I
\mymarginpar{\cmd{Declare-}\\\texttt{Sorting-}\\\texttt{Template}}
have provided, using this declaration, a custom sorting algorithm for
the bibliography.  The idea is that \textsf{biblatex} can use any
field whatsoever for sorting, so that a great many more entries will
be sorted correctly automatically rather than requiring manual
intervention in the form of a \textsf{sortkey} field or the like.
Code in \textsf{biblatex-chicago.sty} loads the custom scheme
\enquote{\texttt{cms},} a Chicago-specific variant of the default
\texttt{nty}.  (You can find its definition in
\textsf{chicago-notes.cbx}.)  The advantages of this scheme are,
specifically, that any entry headed by one of the supplemental name
fields (\textsf{name[a-c]}), a \textsf{manual} or a \textsf{standard}
entry headed by an \textsf{organization}, or an \textsf{article} or
\textsf{review} entry headed by a \textsf{journaltitle} won't need a
\textsf{sortkey} set.  Further, the \texttt{use<name>=false} options
will remove any name field from the sorting order, again reducing the
need for user intervention.

\subsubsection{{Pre-Set \textsf{chicago} Options}}
\label{sec:chicpreset}

At \mymarginpar{\texttt{bookpages=\\true}} the request of Scot Becker,
I have included this rather specialized option, which controls the
printing of the \textsf{pages} field in \textsf{book} entries.  Some
bibliographic managers, apparently, place the total page count in that
field by default, and this option allows you to stop the printing of
this information in notes and bibliography.  It defaults to true,
which means the field is printed, but it can be set to false either in
the preamble, for the whole document or for specific entry types, or
on a per-entry basis in the \textsf{options} field (though rather than
use this latter method it would make sense to eliminate the
\textsf{pages} field from the affected entries).

\mylittlespace This \colmarginpar{\texttt{doi=true}} option controls
whether any \textsf{doi} fields present in the .bib file will be
printed in notes and bibliography.  At the request of Daniel
Possenriede, and keeping in mind the \emph{Manual's} preference for
this field instead of a \textsf{url} (14.6), I have added a third
switch, \texttt{only}, which prints the \textsf{doi} if it is present
and the \textsf{url} only if there is no \textsf{doi}.  Ryo Furue more
recently requested a way to suppress the \textsf{urldate} when using
only the \textsf{doi}, so I've added the \mycolor{\texttt{onlynd}}
switch to do this.  The package default remains the same, however ---
it defaults to true, which will print both \textsf{doi} and
\textsf{url} if both are present.  The option can be set to
\texttt{only}, \mycolor{\texttt{onlynd}}, or to \texttt{false} either
in the preamble, for the whole document or for specific entry types,
or on a per-entry basis in the \textsf{options} field.  In
\textsf{online} entries, the \textsf{doi} field will always be
printed, but the \texttt{only} switch will still eliminate any
\textsf{url}, and \mycolor{\texttt{onlynd}} will still eliminate both
the \textsf{url} and the \textsf{urldate}.

\mylittlespace This \mymarginpar{\texttt{eprint=true}} option controls
whether any \textsf{eprint} fields present in the .bib file will be
printed in notes and bibliography.  It defaults to true, and can be
set to false either in the preamble, for the whole document or for
specific entry types, or on a per-entry basis, in the \textsf{options}
field.  In \textsf{online} entries, the \textsf{eprint} field will
always be printed.

\mylittlespace This \mymarginpar{\texttt{isbn=true}} option controls
whether any \textsf{isan}, \textsf{isbn}, \textsf{ismn},
\textsf{isrn}, \textsf{issn}, and \textsf{iswc} fields present in the
.bib file will be printed in notes and bibliography.  It defaults to
true, and can be set to false either in the preamble, for the whole
document or for specific entry types, or on a per-entry basis, in the
\textsf{options} field.

\mylittlespace Once \mymarginpar{\texttt{numbermonth\\=true}} again at
the request of Scot Becker, I have included this option, which
controls the printing of the \textsf{month} field in all the
periodical-type entries when a \textsf{number} field is also present.
Some bibliographic software, apparently, always includes the month of
publication even when a \textsf{number} is present.  When all this
information is available the \emph{Manual} (14.180, 14.185) prints
everything, so this option defaults to true, which means the field is
printed, but it can be set to false either in the preamble, for the
whole document or for specific entry types, or on a per-entry basis in
the \textsf{options} field.

\mylittlespace This \mymarginpar{\texttt{url=true}} option controls
whether any \textsf{url} fields present in the .bib file will be
printed in notes and bibliography.  It defaults to true, and can be
set to false either in the preamble, for the whole document or for
specific entry types, or on a per-entry basis, in the \textsf{options}
field.  Please note that, as in standard \textsf{biblatex}, the
\textsf{url} field is always printed in \textsf{online} entries,
regardless of the state of this option.

\mylittlespace This \mymarginpar{\texttt{urlstamp=true}} option
controls whether any \textsf{urltime} fields, included as part of the
\textsf{urldate}, will be printed in notes and bibliography.  It
defaults to true, and can be set to false either in the preamble, for
the whole document or for specific entry types, or on a per-entry
basis in the \textsf{options} field.  Please note that, unlike the
\texttt{url} option, this option \emph{does} control what is printed
in \textsf{online} entries.

\mylittlespace This \mymarginpar{\texttt{includeall=\\true}} is the
one option that rules the seven preceding, either printing all the
fields under consideration --- the default --- or excluding all of
them.  It is set to \texttt{true} in \textsf{chicago-notes.cbx}, but
you can change it either in the preamble for the whole document or for
specific entry types, or in the \textsf{options} field of individual
entries.  The seven individual options above are similarly available
in the same places, for finer-grained control.  The rationale for all
of these options is the availability of bibliographic managers that
helpfully present as much data as possible, in every entry, some of
which may not be felt to be entirely necessary.  Setting
\texttt{includeall} to \texttt{true} probably works just fine for
those compiling their .bib databases by hand, but others may find that
some automatic pruning helps clear things up, at least to a first
approximation.  Some per-type or per-entry work afterward may then
polish up the details.  If you find that you need control over fields
that aren't included among these options, I have provided the
\cmd{suppressbibfield} command for your preamble, as suggested by Jan
David Hauck.  It is in fact a user interface to the source mapping
feature of \textsf{biblatex}, and it is something of a nuclear option,
preventing fields from even appearing in the .bbl file generated by
\textsf{biber} from your .bib database.  See the
\cmd{suppressbibfield} command in section~\ref{sec:formatcommands} and
the source mapping docs in \textsf{biblatex.pdf} \S~4.5.3.

\mylittlespace At \mymarginpar{\texttt{addendum=\\true}} the request
of Roger Hart, I have included this option, which controls the
printing of the \textsf{addendum} field, but \emph{only} in long
notes.  It defaults to true, and can be set to false either in the
preamble, for the whole document, or on a per-entry basis, in the
\textsf{options} field.

\mylittlespace According \mymarginpar{\texttt{bookseries=\\true}} to
the \emph{Manual} (14.123), the \textsf{series} field in book-like
entries \enquote{may be omitted to save space (especially in a
  footnote).}  This option allows you to control the printing of that
field in long notes.  It defaults to true, and can be set to false
either in the preamble, for the whole document, or on a per-entry
basis, in the \textsf{options} field.  Several entry types don't use
this field, so the option will have no effect in them, and it is also
ignored in \textsf{article}, \textsf{misc}, \textsf{music},
\textsf{periodical}, and \textsf{review} entries.

\mylittlespace As \mymarginpar{\texttt{notefield=\\true}} with the
previous two options, Roger Hart requested an option to control the
printing of the \textsf{note} field in long notes.  It defaults to
true, and can be set to false either in the preamble, for the whole
document, or on a per-entry basis, in the \textsf{options} field.  The
option will be ignored in \textsf{article}, \textsf{misc},
\textsf{periodical}, and \textsf{review} fields.

\mylittlespace This
\mymarginpar{\vspace{-1\baselineskip}\texttt{completenotes=}%
  \\\texttt{true}} is the one option that rules
the three preceding, either printing all the fields under
consideration --- the default --- or excluding all of them from long
notes.  It is set to \texttt{true} in \textsf{chicago-notes.cbx}, but
you can change it either in the preamble for the whole document or,
for specific fields, in the \textsf{options} field of individual
entries.

\mylittlespace These
\mymarginpar{\texttt{bibannotesep\\=vpar\\citeannotesep\\=period}}
options define the relation of the \textsf{annotation} field to the
main entry, \texttt{bibannotesep} doing so in the bibliography and
\texttt{citeannotesep} in long notes.  (The \texttt{annotation} option
in section~\ref{sec:useropts} determines where, if anywhere, the field
will appear.)  Both options have the same set of keys, though they
have different default settings if you don't define them yourself.
The possible values are:

\begin{description}
\setlength{\parskip}{-4pt}
\item[\qquad none] = no punctuation at all.
\item[\qquad space] = \cmd{addspace}
\item[\qquad comma] = \verb+\addcomma\addspace+
\item[\qquad period] = \verb+\addperiod\addspace+
\item[\qquad colon] = \verb+\addcolon\addspace+
\item[\qquad semicolon] = \verb+\addsemicolon\addspace+
\item[\qquad par:] This starts a new paragraph on the next line.  Page
  breaking is strongly inhibited before the \textsf{annotation}.
\item[\qquad vpar:] This starts a new paragraph, and also inserts
  some vertical space before it.  In long notes that vertical space is
  minimal (1 pt), while in the bibliography it creates a blank line.
  Page breaking is strongly inhibited before the \textsf{annotation}.
\item[\qquad parbreak:] This is the same as \texttt{par}, but it allows
  a page break to occur between the main entry and the start of the
  \textsf{annotation}.
\item[\qquad vparbreak:] This is the same as \texttt{vpar}, but it also
  allows a page break between the main entry and the
  \textsf{annotation}.
\end{description}

Please note that both of these options are available in the preamble
both globally and per-type, and also in the \textsf{options} field of
individual entries.  Each defines a command (\cmd{bibannotesep} and
\cmd{citeannotesep}) which appears in the \textsf{annotation} field's
formatting directive, so it's possible to redefine these commands in
your preamble if you have needs that the available values don't
address.  (You can also try sending an email to encourage me to add
other keys.)  Please also keep in mind that \texttt{bibannotesep}
interacts with the \texttt{entrybreak} and \textsf{formatbib} options
in section~\ref{sec:useropts}, below, to determine the general layout
of the bibliography.  Depending on the settings of those options,
changing the \texttt{bibannotesep} from entry to entry may not work
out well.

\mylittlespace At \mymarginpar{\texttt{booklongxref=\\true}} the
request of Bertold Schweitzer, I have included two options for
controlling whether and where \textsf{biblatex-chicago} will print
abbreviated references when you cite more than one part of a given
collection or series.  This option controls whether multiple
\textsf{book}, \textsf{bookinbook}, \textsf{collection}, and
\textsf{proceedings} entries which are part of the same collection
will appear in this space-saving format.  The parent collection itself
will usually be presented in, e.g., a \textsf{book},
\textsf{bookinbook}, \textsf{mvbook}, \textsf{mvcollection}, or
\textsf{mvproceedings} entry, and using \textsf{crossref} or
\textsf{xref} in the child entries will allow such presentation
depending on the value of the option:

\begin{description}
\setlength{\parskip}{-4pt}
\item[\qquad true:] This is the default.  If you use \textsf{crossref}
  or \textsf{xref} fields in these entry types, by default you will
  \emph{not} get any abbreviated references, either in notes or
  bibliography.
\item[\qquad false:] You'll get abbreviated references in these entry
  types both in notes and in the bibliography.
\item[\qquad notes:] The abbreviated references will not appear in
  notes, but only in the bibliography.
\item[\qquad bib:] The abbreviated references will not appear in the
  bibliography, but only in notes.
\end{description}

This option can be set either in the preamble or in the
\textsf{options} field of individual entries.  For controlling the
behavior of \textsf{inbook}, \textsf{incollection},
\textsf{inproceedings}, and \textsf{letter} entries, please see
\texttt{longcrossref}, below, and also the documentation of
\textsf{crossref} in section~\ref{sec:entryfields}.

\mylittlespace The \mymarginpar{\texttt{compressyears\\=true}}
\emph{Manual} has long recommended (9.64, 14.117, 14.144), as a
space-saving measure, the compression of year ranges when presenting
dates.  I have, finally, implemented this in the current release, and
have made it the default, which you can change in your document
preamble.  Please note that the rules for compressing years are
different from those for compressing other numbers (e.g., page
numbers), and also that the compression code is in
\textsf{biblatex-chicago.sty}, which will have to be loaded for this
option to make any difference.  Cf.\ table~\ref{ad:date:extras}.

\mylittlespace Roger
\mymarginpar{\texttt{ctitleaddon=\\comma\\ptitleaddon=\\period}} Hart
requested a way to control the punctuation printed before the
\textsf{titleaddon}, \textsf{booktitleaddon}, and
\textsf{maintitleaddon} fields.  By default, this is
\verb+\addcomma\addspace+ (\textsf{cti\-tleaddon}) for all occurrences
in notes and for nearly all \textsf{book-} and
\textsf{maintitleaddons} in the bibliography, while
\verb+\addperiod\addspace+ (\textsf{ptitleaddon}) is the default
before most \textsf{titleaddons} in the bibliography.  If the
punctuation printed isn't correct for your needs, you can set the
relevant option either in the preamble or in individual entries.
(Cf.\ coolidge\hc speech and schubert:muellerin.)  The accepted option
keys are:

\begin{description}
\setlength{\parskip}{-4pt}
\item[\qquad none] = no punctuation at all
\item[\qquad space] = \cmd{addspace}
\item[\qquad comma] = \verb+\addcomma\addspace+
\item[\qquad period] = \verb+\addperiod\addspace+
\item[\qquad colon] = \verb+\addcolon\addspace+
\item[\qquad semicolon] = \verb+\addsemicolon\addspace+
\end{description}

If you need something a little more exotic, you can directly
\cmd{renewcommand} either \cmd{ctitleaddonpunct} or
\cmd{ptitleaddonpunct} (or both) in your preamble, but it's worth
remembering that the redefinition will hold for all instances, unless
you use the \textsf{options} field in your other entries with a
\textsf{titleaddon} field.  A simpler solution might be to set the
relevant option to \texttt{none} in your entry and then include the
punctuation in the \textsf{titleaddon} field itself.

\mylittlespace Constanza Cordoni \mymarginpar{\texttt{dashed=true}}
has requested a way to turn off the 3-em dash for replacing repeated
names in the bibliography, and the \emph{Manual} admits that some
publishers prefer this, as the dash can carry with it certain
inconveniences, especially for electronic formats (14.67).  Some of
\textsf{biblatex's} standard styles have a \texttt{dashed} option, so
for compatibility purposes I've provided the same.  By default I have
set it to print the name dash, but you can set \texttt{dashed=false}
globally, per type, or per entry to repeat names as and when required.

\mylittlespace If \mymarginpar{\texttt{hidevolumes=\\true}} both a
\textsf{volume} and a \textsf{volumes} field are present, as may occur
particularly in cross-referenced entries, then
\textsf{biblatex-chicago} will ordinarily suppress the
\textsf{volumes} field.  In some instances, when a \textsf{maintitle}
is present, this may not be the desired result.  In this latter case,
if the \textsf{volume} appears before the \textsf{maintitle}, this new
option, set to \texttt{true} by default, controls whether to print the
\textsf{volumes} field after that title or not.  Set it to
\texttt{false} either in the preamble or in the \textsf{options} field
of your entry to have it appear after the \textsf{maintitle}.

\mylittlespace I \colmarginpar{\texttt{jtitleaddon=\\space}} have
added the standard \textsf{biblatex}
\mycolor{\textsf{journaltitleaddon}} field to the \textsf{article} and
\textsf{review} entry types, and also the \textsf{titleaddon} field to
the \textsf{periodical} type, fields that may, for example, be
particularly useful when you want to provide the original form of a
translated journal title.  The \mycolor{\texttt{jtitleaddon}} option
controls the separator between the main title and the addon, as with
the \texttt{ctitleaddon} and \texttt{ptitleaddon} options, above, and
like them is settable globally, per type, or per entry.  The possible
settings are the same as for those options, but the default is a
\texttt{space}.  You can redefine \cmd{jtitleaddonpunct} directly if
you have more unusual needs.

\mylittlespace This \mymarginpar{\texttt{longcrossref=\\false}} is the
second option, requested by Bertold Schweitzer, for controlling
whether and where \textsf{biblatex-chicago} will print abbreviated
references when you cite more than one part of a given collection or
series.  It controls the settings for the entry types more-or-less
authorized by the \emph{Manual}, i.e., \textsf{inbook},
\textsf{incollection}, \textsf{inproceedings}, \textsf{letter}, and
\textsf{review}.  The mechanism itself is enabled by multiple
\textsf{crossref} or \textsf{xref} references to the same parent,
whether that be, e.g., a \textsf{collection}, an
\textsf{mvcollection}, a \textsf{proceedings}, or an
\textsf{mvproceedings} entry.  Given these multiple cross references,
the presentation in the reference apparatus will be governed by the
following options:

\begin{description}
\setlength{\parskip}{-4pt}
\item[\qquad false:] This is the default.  If you use
  \textsf{crossref} or \textsf{xref} fields in the four mentioned
  entry types, you'll get the abbreviated references in both notes and
  bibliography.
\item[\qquad true:] You'll get no abbreviated references in these
  entry types, either in notes or in the bibliography.
\item[\qquad notes:] The abbreviated references will not appear in
  notes, but only in the bibliography.
\item[\qquad bib:] The abbreviated references will not appear in the
  bibliography, but only in notes.
\item[\qquad none:] This switch is special, allowing you with one
  setting to provide abbreviated references not just to the four entry
  types mentioned but also to \textsf{book}, \textsf{bookinbook},
  \textsf{collection}, and \textsf{proceedings} entries, both in notes
  and in the bibliography.
\end{description}

This option can be set either in the preamble or in the
\textsf{options} field of individual entries.  For controlling the
behavior of \textsf{book}, \textsf{bookinbook}, \textsf{collection},
and \textsf{proceedings} entries, please see \texttt{booklongxref},
above, and also the documentation of \textsf{crossref} in
section~\ref{sec:entryfields}.

\mylittlespace This \mymarginpar{\texttt{nameaddonsep\\=space}} option
sets the punctuation which appears before the \textsf{nameaddon} field
in all entry types except \textsf{customc}.  You can set it globally,
per type or per entry, using one of the six following keys:

\begin{description}
\setlength{\parskip}{-4pt}
\item[\qquad \texttt{space}] = \cmd{addspace}.  This is the default.
\item[\qquad \texttt{none}] = no separator at all.  It presumes that
  you will include one in the \textsf{nameaddon} field itself.
\item[\qquad \texttt{colon}] = \verb+\addcolon\addspace+.
\item[\qquad \texttt{comma}] = \verb+\addcomma\addspace+.
\item[\qquad \texttt{period}] = \verb+\addperiod\addspace+.
\item[\qquad \texttt{semicolon}] = \verb+\addsemicolon\addspace+.
\end{description}

Cf.\ \texttt{nameaddon} and \texttt{nameaddonformat} in
section~\ref{sec:useropts}.

\mylittlespace This \mymarginpar{\texttt{nodates=true}} option means
that \textsf{biblatex-chicago} automatically provides
\verb+\bibstring{nodate}+ in any \enquote{circa} or
\enquote{uncertain} date specification where the user has also turned
off the printing of brackets around the date using the
\texttt{nodatebrackets} or \texttt{noyearbrackets} options
(section~\ref{sec:useropts}; 14.145).  If you set
\texttt{nodates=false} either in your preamble (for global coverage or
for specific entry types) or in individual entries then the package
will simply print the unbracketed date specification in this context.
See table~\ref{ad:date:extras}, below.  (The bibstring expands to
\enquote{\texttt{n.d.}} in English; please note that this option works
quite differently in the author-date styles.)

\mylittlespace This
\mymarginpar{\texttt{shorthand\\punct\\=space}}
option controls the punctuation that appears before the first
introduction of a \textsf{shorthand} field, including the
\textsf{shorthandintro}, in long notes.  The default is
\cmd{addspace}, but if this isn't correct for your needs, especially
if you change the \textsf{shorthandintro} or don't want the whole
phrase inside parentheses, then you can change it in the preamble or
in individual entries.  The accepted option keys are:

\begin{description}
\setlength{\parskip}{-4pt}
\item[\qquad none] = no punctuation at all
\item[\qquad space] = \cmd{addspace}
\item[\qquad comma] = \verb+\addcomma\addspace+
\item[\qquad period] = \verb+\addperiod\addspace+
\item[\qquad colon] = \verb+\addcolon\addspace+
\item[\qquad semicolon] = \verb+\addsemicolon\addspace+
\item[\qquad emdash] = \verb+\addthinspace\textemdash\addthinspace+
\item[\qquad endash] = \verb+\addspace\textendash\addspace+
\end{description}

You can, in emergencies, directly
\verb+\renewcommand{\shorthandpunct}+ in the preamble, but it might
be easier to use the \texttt{none} option to \texttt{shorthandpunct}
and hand-craft solutions inside the \textsf{shorthandintro} fields of
individual entries.

\subsubsection{Style Options -- Preamble}
\label{sec:useropts}

These are parts of the specification that not everyone will wish to
enable.  All except the sixth, seventh, and twelfth can be used even
if you load the package in the old way via a call to
\textsf{biblatex}, but most users can just place the appropriate
string(s) in the options to the \cmd{usepackage}
\texttt{\{biblatex-chicago\}} call in your preamble.

\mylittlespace \textsf{Biblatex-chicago}
\mymarginpar{\texttt{alwaysrange}} now implements \textsf{biblatex's}
enhanced date specifications, one part of which is the presentation of
decades and centuries not as year ranges but as localized strings like
\enquote{19th c.} or \enquote{1970s.}  The \texttt{alwaysrange} option
set to \texttt{true}, either in your preamble or in individual
entries, simply tells the package to present the year range instead.
This allows you to use the efficient enhanced notations in the
\textsf{date} field (\verb+{18XX}+ or \verb+{197X}+) without the
localized strings appearing, should you require it.  The two options
\texttt{centuryrange} and \texttt{decaderange} limit the same effect
to centuries and decades, respectively.  Please see
table~\ref{ad:date:extras}.

\mylittlespace At \mymarginpar{\texttt{annotation}} the request of
Emil Salim, I included in \textsf{biblatex-chicago} the ability to
produce annotated bibliographies.  More recently, Moritz Wemheuer
brought to my attention a
\href{https://tex.stackexchange.com/questions/528374/moving-addendum-field-to-the-end-in-biblatex-chicago/540755#540755}{StackExchange}
question which suggested that the field might be useful in several
other contexts as well, so I've modified the \texttt{annotation}
option to allow the field to appear in the bibliography (\texttt{=bib}
or \texttt{=true}, the default, if no string is given), in long notes
(\texttt{=notes}), in both (\texttt{=all}), or in neither
(\texttt{=false}).  You can now set the option in the preamble both
globally and per-type, and in the \textsf{options} field of individual
entries.  There are two new options (\texttt{bibannotesep} and
\texttt{citeannotesep}) to allow you to choose the separator between
the \textsf{annotation} and the rest of the entry, and also two new
options (\texttt{formatbib} and \texttt{entrybreak}) to give you
fine-grained control over the presentation of the bibliography as a
whole, including an annotated one.  Please have a look at the
documentation for the latter two options just below, for the former
two options in section~\ref{sec:chicpreset}, and for the
\textsf{annotation} field on page~\pageref{sec:annote}.  Please also
see \textsf{biblatex.pdf} \S~3.13.8 for details on how to use external
files to store annotations.

\mylittlespace As \mymarginpar{\texttt{blogurl}} a convenience
feature, this option, available only in the preamble, allows
\textsf{review} entries to inherit a \textsf{url} from
\textsf{article} entries.  The main use, as the name suggests, is when
you want to provide the same \textsf{url} for a blog comment as you
have for a blog post.  You'll need an extra \LaTeX - \textsf{Biber} -
\LaTeX\ run to make sure any changes to this option take effect.

\mylittlespace Like \mymarginpar{\texttt{casechanger}} \texttt{mcite}
and \texttt{natbib}, this is a standard \textsf{biblatex} option which
\textsf{biblatex-chicago} simply passes through to that package.  In
\textsf{biblatex} it defaults to \texttt{auto}, but there were, and
possibly still are, cases when the old \texttt{latex2e} case-changing
code can work around some bugs when using, for example, citation
commands inside fields that have a case-changing element as an
automatic part of their formatting (\textsf{note},
\textsf{titleaddon}, \textsf{type}).  Cf.\ the \cmd{citeincite}
command in section~\ref{sec:citecommands}, and section~3.1.1 of
\textsf{biblatex.pdf}.

\mylittlespace This \mymarginpar{\texttt{centuryrange}} option works
just like \texttt{alwaysrange}, above, but only affects century
presentation, not decade.  Cf.\ table~\ref{ad:date:extras}.

\mylittlespace The \mymarginpar{\texttt{cmsbreakurl}} \emph{Manual}
gives fairly specific instructions about breaking URLs across lines
(14.18), so I have attempted to implement them by tweaking
\textsf{biblatex's} default settings, which are found in
\textsf{biblatex.def}.  In truth, I haven't succeeded in getting
\textsf{biblatex} flawlessly to follow all of the \emph{Manual's}
instructions, nor do the changes I have made work well in all
circumstances, being particularly unsightly if you happen to be using
the \textsf{ragged2e} package.  For these reasons, I have made my
changes dependent on a package option, \texttt{cmsbreakurl}, which you
can set in your preamble.  I have placed all of this code in
\textsf{biblatex-chicago.sty}, so if you load the package with a call
to \textsf{biblatex} instead, then URL line breaking will revert to
the \textsf{biblatex} defaults.  See
\href{file:cms-notes-sample.pdf}{\textsf{cms-notes-sample.pdf}} for a
lot of examples of what URLs look like when the option is set, and
also section~\ref{sec:loading}, below.

\mylittlespace When \mymarginpar{\texttt{compresspages}} set to
\texttt{true}, any page ranges in your .bib file or in the
\textsf{postnote} field of your citation commands will be compressed
in accordance with the \emph{Manual's} specifications (9.61).
Something like 321-{-}328 in your .bib file would become 321--28 in
your document.  See the \textsf{pages} field in
section~\ref{sec:entryfields}, above.  Please note that the code for
this is in \textsf{biblatex-chicago.sty}, so if you load the package
with a call to \textsf{biblatex} instead then you'll get the default
\textsf{biblatex} compression style.

\mylittlespace This \mymarginpar{\texttt{datamodel}} is the standard
\textsf{biblatex} option for loading the named data model file
(excluding its\ .dbx extension).  After a request by Philipp Immel,
you can now set this option when you load the Chicago styles with
\verb+\usepackage{biblatex-chicago}+, and it will be passed through
properly to \textsf{biblatex} itself.  Cf.\ \textsf{biblatex.pdf}
\S~4.5.4.

\mylittlespace This \mymarginpar{\texttt{decaderange}} option works
just like \texttt{alwaysrange}, above, but only affects decade
presentation, not century.  Cf.\ table~\ref{ad:date:extras}.

\mylittlespace The \mymarginpar{\texttt{delayvolume}} presentation of
\textsf{volume} information in the notes \&\ bibliography style is
complicated (\emph{Manual}, 14.116--22).  Depending on entry type and
on the presence or absence of a \textsf{booktitle} or a
\textsf{maintitle}, \textsf{volume} data will be presented, in the
bibliography, either before a \textsf{maintitle} or after a
\textsf{booktitle} or \textsf{maintitle}, that is, just before
publication information.  This, so far, is handled for you
automatically by \textsf{biblatex-chicago-notes}.  In long notes, the
same options apply, but it is also sometimes better to place
\textsf{volume} information \emph{after} the publication information
and just before any page numbers, so I have included this option,
which you can set either for the whole document or on a per-entry
basis, to allow you to move \textsf{volume} data to the end of a long
note.  Please note that this doesn't affect any \textsf{volume} data
printed \emph{before} a \textsf{maintitle}, but only data that would,
without this option, be printed \emph{after} a \textsf{booktitle} or
\textsf{maintitle}.  Cf.\ also \cmd{postvolpunct}, below.

\mylittlespace This \mymarginpar{\texttt{entrybreak}} option only
makes sense when used in conjunction with the \texttt{annotenp} switch
to the \texttt{formatbib} option (below).  The latter allows \LaTeX\
to break a page inside an \textsf{annotation} field printed without
starting a new paragraph, and it does so by allowing such a break in
the entry only after a set number of lines, by default set to 3.  The
idea is that most bibliography entries will fit within 3 lines, so the
break would generally be somewhere inside the \textsf{annotation}.  If
your document needs a value different from 3, provide the integer
using the \texttt{entrybreak} option in your preamble.  Some
experimentation may be needed to find the optimum number for a given
document.

\mylittlespace Although \mymarginpar{\texttt{footmarkoff}} the
\emph{Manual} (14.24) recommends specific formatting for footnote (and
endnote) marks, i.e., superscript in the text and in-line in foot- or
endnotes, Charles Schaum has brought it to my attention that not all
publishers follow this practice, even when requiring Chicago style.  I
have retained this formatting as the default setup, but if you include
the \texttt{footmarkoff} option, \textsf{biblatex-chicago-notes} will
not alter \LaTeX 's (or the \textsf{endnote} package's) defaults in
any way, leaving you free to follow the specifications of your
publisher.  I have placed all of this code in
\textsf{biblatex-chicago.sty}, so if you load the package with a call
to \textsf{biblatex} instead, then once again footnote marks will
revert to the \LaTeX\ default, but of course you also lose a fair
amount of other formatting, as well.  See section~\ref{sec:loading},
below.

\mylittlespace The \mymarginpar{\texttt{formatbib}} \emph{Manual} in
fact says very little about formatting issues in bibliographies, e.g.,
whether to break entries across pages and whether to allow widows and
orphans (single lines at the start or end of a page).  A quick and
non-scientific survey of publications issued by the University of
Chicago Press suggests that actual practices are extremely varied, so
I've tried to provide a number of choices for users of
\textsf{biblatex-chicago}, most of them available as keys to the new
\texttt{formatbib} preamble option, but a few of them also involving
settings of the \texttt{bibannotesep} and the \texttt{entrybreak}
options.  The keys of \texttt{formatbib} are as follows:

\begin{description}
\setlength{\parskip}{-4pt}
\item[\qquad max:] This is the \textsf{biblatex} default, so if you
  don't set \texttt{formatbib} at all it's what you get.  It provides
  maximal intervention, disallowing entries broken across pages,
  including even when an entry includes a lengthy annotation.
\item[\qquad min:] This allows page breaks just about anywhere,
  including inside entries, and it also allows widows and orphans, so
  it will usually provide the most efficient use of available space on
  the pages of your bibliography.
\item[\qquad minwo:] This is like \texttt{min}, but discourages
  widows and orphans.
\item[\qquad annote:] This option treats annotations separately from
  bibliography entries, allowing them to be broken across pages while
  the entry itself won't be.  It is intended for use with, and will
  only properly work with, \texttt{bibannotesep} set to one of the
  modes that start a new paragraph for the annotation, to wit,
  \texttt{par}, \texttt{vpar} (the default), \texttt{parbreak}, or
  \texttt{vparbreak}.  See below for the meaning of the
  \enquote{\texttt{break}} options here.
\item[\qquad annotenp:] This option attempts to treat annotations
  separately from bibliography entries in those settings of
  \texttt{bibannotesep} which don't involve starting a new paragraph.
  It works by setting the number of lines in an entry after which page
  breaking is allowed.  By default entries will only break after 3
  lines, the idea being that most bibliography entries fit into three
  lines, so at that point you're likely to be inside the annotation,
  but you can set the \texttt{entrybreak} option to any integer that
  works for your reference apparatus.
\end{description}

Please note that there is one possible break point that isn't directly
addressed by these options, that is, the one between the main entry
and the annotation when that annotation starts a new paragraph.  If
you set \texttt{bibannotesep} to \texttt{par} or \texttt{vpar}, then
\LaTeX\ will try very hard not to break between entry and annotation,
ensuring that the annotation at least starts on the same page as its
entry.  If you use \texttt{parbreak} or \texttt{vparbreak}, \LaTeX\ is
positively encouraged to break a page there, as is usual between
paragraphs.

\mylittlespace You can of course ignore the \texttt{formatbib} option
and provide your own settings.  \textsf{Biblatex} uses the
\cmd{bibsetup} command which you can renew in your preamble.  You can
find a nice commentary on the default values set by the package in the
file \textsf{biblatex.def}, which you'll find in the main
\textsf{biblatex} directory of your \TeX\ distribution.

\mylittlespace This \mymarginpar{\texttt{genallnames}} option affects
the choice of which names to present in the genitive case when using
the \cmd{gentextcites} command.  Please see the documentation of that
command in section~\ref{sec:citecommands}, above.

\mylittlespace Setting \mymarginpar{\texttt{inheritshorthand}} this
option to \texttt{true} allows child entries to inherit the
\textsf{shorthand} and \textsf{shorthandintro} fields from
cross-referenced parent entries.  This in turn allows abbreviated
references to the parent entry to use the \textsf{shorthand} instead
of the usual and merely short citation, thus allowing for extra space
savings.  There are several other steps required to make this all
function smoothly, so please see the documentation of the
\textbf{shorthand} field in section~\ref{sec:entryfields}, above.

\mylittlespace This \mymarginpar{\texttt{journalabbrev}} option
controls the printing of the \textsf{shortjournal} field in place of
the \textsf{journaltitle} field in notes and bibliography.  It is
\texttt{false} by default, so as shipped
\textsf{biblatex-chicago-notes} will silently ignore such fields, but
you can set it, either in the preamble or in individual entries, to
one of three other values: \texttt{true} prints the abbreviated form
in notes and bibliography, \texttt{notes} in notes only, and
\texttt{bib} in the bibliography only.  Please note that in
\textsf{periodical} entries the \textsf{title} and \textsf{shorttitle}
fields behave in exactly the same manner.  For more details, see the
documentation of \textsf{shortjournal} in
section~\ref{sec:entryfields}, above.

\mylittlespace The \mymarginpar{\texttt{juniorcomma}} \emph{Manual}
(6.43) states that \enquote{commas are not required with \emph{Jr.}\
  and \emph{Sr.},} so by default \textsf{biblatex-chicago} has
followed standard \textsf{biblatex} in using a simple space in names
like \enquote{John Doe Jr.}  Charles Schaum has pointed out that
traditional \textsc{Bib}\TeX\ practice was to include the comma, and
since the \emph{Manual} has no objections to this, I have provided an
option which allows you to turn this behavior back on, either for the
whole document or on a per-entry basis.  Please note, first, that
numerical suffixes (John Doe III) never take the comma.  The code
tests for this situation, and detects cardinal numbers well, but if
you are using ordinals you may need to set this to \texttt{false} in
the \textsf{options} field of some entries.  Second, I have fixed a
bug in older releases which always printed the \enquote{Jr.}\ part of
the name immediately after the surname, even when the surname came
before the given names (as in a bibliography).  The package now
correctly puts the \enquote{Jr.}\ part at the end, after the given
names, and in this position it always takes a comma, the presence of
which is unaffected by this option.

\mylittlespace This \mymarginpar{\texttt{nameaddon}} option determines
where and when the \textsf{nameaddon} field will be printed.  There
are seven possible values, the first four of which are available
globally, per type, and per entry, with the last three only applicable
globally:
\begin{description}
\setlength{\parskip}{-4pt}
\item[\qquad \texttt{all}:] This is the default; if an entry has a
  \textsf{nameaddon}, it will appear in both long notes and in the
  bibliography.
\item[\qquad \texttt{none}:] The field will appear neither in the
  bibliography nor in long notes.
\item[\qquad \texttt{bib}:] The field will appear only in the
  bibliography.
\item[\qquad \texttt{cite}:] The field will appear only in long notes.
\item[\qquad \texttt{first}:] Philipp Immel requested this as a way to
  provide an \textsf{author's} dates in the \textsf{nameaddon} field
  and only have them printed the first time that author appears in the
  bibliography.  A sequence of consecutive long notes citing works by
  the same author will be treated the same way.  The code tests for
  identical \textsf{nameaddon} fields in works by identical
  \textsf{authors}, so other sorts of \textsf{nameaddon} will be
  printed as usual.
\item[\qquad \texttt{bibfirst}:] Like \texttt{first}, but will not
  print the \textsf{nameaddon} field in long notes.
\item[\qquad \texttt{citefirst}:] Like \texttt{first}, but will not
  print the \textsf{nameaddon} field in the bibliography.
\end{description}

Cf.\ \texttt{nameaddonformat} just below, and \texttt{nameaddonsep} in
section~\ref{sec:chicpreset}.

\mylittlespace This \mymarginpar{\texttt{nameaddon-\\format}} option,
available globally, per type, and per entry, allows you to change the
format of the \textsf{nameaddon} field on the fly, so its value should
be a field format that \textsf{biblatex} understands.  This includes
standard formats like \texttt{parens,\,brackets} or \texttt{emph}, and
also custom formats that you provide in your preamble using
\cmd{DeclareFieldFormat}, in case the standard ones aren't adequate.
If you don't define this option, then the usual defaults apply, that
is, no formatting in \textsf{online, review,} and
\textsf{suppperiodical} entries, as well as in \textsf{misc} entries
with an \textsf{entrysubtype}, while square brackets surround the
field in all other entry types with the exception of \textsf{customc},
which has its own rules and ignores this option. Cf.\
\texttt{nameaddonsep} in section~\ref{sec:chicpreset}.

\mylittlespace These \mymarginpar{\texttt{natbib}\\\texttt{mcite}} are
the standard \textsf{biblatex} options, which formerly required
slightly special handling when you loaded the Chicago style with
\verb+\usepackage{biblatex-chicago}+.  Both the forms \texttt{natbib}
and \texttt{natbib=true} (\texttt{mcite} \&\ \texttt{mcite=true})
should now work.

\mylittlespace When \mymarginpar{\texttt{nodatebrackets}\\
  \texttt{noyearbrackets}} you use \textsf{biblatex's} enhanced date
specifications to present a \enquote{circa} date (\verb+{1956~}+), an
uncertain date (\verb+{1956?}+), or one that is both at the same time
(\verb+{1956%}+), the date that by default will appear in your
documents will have square brackets around it.  This accords with the
\emph{Manual's} instructions concerning such dates (14.145), but that
section also includes an alternative form, where the guessed at date
appears, without brackets, after the \verb+\bibstring{nodate}+, e.g.,
\enquote{\texttt{n.d., ca.\ 1750.}}  These two package options, which
may appear in the preamble either for the whole document or for
specific entry types, or in individual entries, allow you to control
when these brackets will appear, while the \texttt{nodates} option,
set to \texttt{true} by default, decides whether to print
\verb+\bibstring{nodate}+ before the date.  In truth, users of the
notes \&\ bibliography style will probably only ever need
\texttt{nodatebrackets}, which controls most of the dates that will
appear in your documents, with the exception of dates in some
\textsf{article}, \textsf{review}, and \textsf{periodical} entries
without an \textsf{entrysubtype}, which are governed by
\texttt{noyearbrackets}.  (The distribution is different in the
author-date styles, so it's impossible to do without both options.)
Cf.\ table~\ref{ad:date:extras}.

\mylittlespace At \colmarginpar{\texttt{noibid}} the request of an
early tester, I have included this option to allow you to turn off the
\emph{ibidem} mechanism that \textsf{biblatex-chicago-notes} uses by
default.  Some publishers, it would appear, require this.  Setting
this option will mean that instead of the \emph{ibidem} mechanism
you'll get the short note form.  Please note that the 17th edition no
longer recommends the use of \enquote{\emph{ibid.}}\ at all (14.34),
so depending on the state of the \texttt{useibid} option, below, what
you'll be turning off may well no longer be the appearance of
\emph{ibid.}\ itself.  The option is settable globally, per type, or
per entry, so that fine-grained control of individual citations is now
possible without the use of specialized citation commands, though
these commands are still available in a pinch.  See section
\ref{sec:citecommands}.

\mylittlespace When \colmarginpar{\texttt{notitle}} citing sources
from antiquity (using the \texttt{classical} \textsf{entrysubtype}),
the \emph{Manual} (14.244--5) recommends using just the
\textsf{author} in short citations if only one \textsf{title} by that
author has come down to us, making the identification of the work
unambiguous.  I previously suggested using a command like
\cmd{citeauthor} to achieve this, but Tobias Becht suggested that a
less clumsy method would be better, so I've provided the
\mycolor{\texttt{notitle}} option, settable by entry type and also in
the \textsf{options} field of specific entries.  The option has no
effect whatever in long notes.  Cf.\ herodotus:wilson.

\mylittlespace As \mymarginpar{\texttt{omitxrefdate}} part of the
abbreviated cross-referencing functionality for \textsf{book},
\textsf{bookinbook}, \textsf{collection}, and \textsf{proceedings}
entries, I have thought it helpful to include, in the abbreviated
references only, a date for any \textsf{title} that's part of a
\textsf{maintitle}, though not for those that are only part of
\textsf{booktitle}.  If these dates annoy you, you can use this option
to turn them off, either in the preamble for the document as a whole
or in the \textsf{options} field of individual entries.  Cf.\
harley:ancient:cart, harley:cartography, and harley:hoc; and
\textsf{crossref} in section~\ref{sec:entryfields}, above.

\mylittlespace This \mymarginpar{\texttt{ordinalgb}} option, which
only affects users of the \texttt{british} language, restores the
previous package defaults, printing the \textsf{day} part of a
\textsf{date} specification as an ordinal number: 26th March 2017.
The new package default prints 26 March 2017, which is more in keeping
both with standard British usage and with the recommendations of the
\emph{Manual} (6.38).  The option is available only in the preamble.

\mylittlespace Several
\mymarginpar{\texttt{postnotepunct}\\(experimental)} users, most
recently David Gohlke, have requested a way to alter the punctuation
that appears just before the \textsf{postnote} argument of citation
commands, usually, but perhaps not always, to allow citations to fit
better into the flow of text.  This punctuation is a complex issue in
the \emph{Manual}, and I've attempted to make
\textsf{biblatex-chicago} follow the specifications closely.  Still,
as a first stab at enabling the greater flexibility in punctuation
that some have requested, I have introduced the \texttt{postnotepunct}
package option.  Set to \texttt{true}, it allows you to start the
\textsf{postnote} field with a punctuation mark (.\,,\,;\,:) and have
it appear as the \cmd{postnotedelim} in place of whatever the package
might otherwise automatically have chosen.  Please note that this
functionality relies on a very nifty macro by Philipp Lehman which I
haven't extensively tested, so I'm labeling this option experimental.
Note also that the option only affects the \textsf{postnote} field of
citation commands, not the \textsf{pages} field in your .bib file.

\mylittlespace This \mymarginpar{\texttt{seriesabbrev}} option
controls the printing of the \textsf{shortseries} field in place of
the \textsf{series} field in book-like entries in notes and
bibliography.  It is \texttt{false} by default, so as shipped
\textsf{biblatex-chicago-notes} will silently ignore such fields, but
you can set it, either in the preamble for the whole document or for
specific entry types, or in individual entries, to one of three other
values: \texttt{true} prints the abbreviated form in notes and
bibliography, \texttt{notes} in notes only, and \texttt{bib} in the
bibliography only.  For more details, see the documentation of
\textsf{shortseries} in section~\ref{sec:entryfields}, above.

\mylittlespace This \colmarginpar{\texttt{short}} option means that
your text will only use the short note form, even in the first
citation of a particular work.  The \emph{Manual} (14.19) recommends
this space-saving format only when you provide a \emph{full}
bibliography, though even with such a bibliography you may feel it
easier for your readers to present long first citations.  Tobias Becht
pointed out that the \emph{Manual} (14.242--54) envisages using short
citations for ancient, medieval, and Renaissance works even in the
first reference, and this without regard to whether you do so with
other works, so now you can set the \mycolor{\texttt{short}} option
for a whole document, for specific entry types, and for specific
entries.  If you do use the \texttt{short} option, remember that there
are several citation commands which allow you to present the full
reference in specific cases (see section \ref{sec:citecommands}).  If
your bibliography is not complete, then you should probably not set
this option globally.

\mylittlespace N.\
Andrew\mymarginpar{\texttt{shortextra-\\field\\shortextra-\\format\\shortextra-\\punct}}
Walsh has remarked that it is quite possible for documents to cite
works that, though perfectly distinguishable in their long form, end
up looking identical in short notes; multiple performances of the same
work by the same artist, for example, might end up producing such a
situation.  (In \textsf{online} and \textsf{review} entries using
\texttt{commenton} \textsf{relatedtype} this situation is so endemic
that I've set a default method of disambiguating short notes there,
though you can still override it with the following options.  See
section~\ref{sec:related}, above.)  While the use of a
\textsf{shorthand} field could provide some sort of remedy, he
requested a way to disambiguate short notes by adding a
user-configurable field to the note, thereby keeping it both short and
unique without the need to consult a list of shorthands.  The
\emph{Manual} (15.29) itself provides just such a mechanism in the
author-date specification, so I've added one to the notes \&\
bibliography style, as well.  It consists of the standard
\textsf{biblatex} option \texttt{labeltitle}, now set to \texttt{true}
by default, along with three package options for the user to
configure.  All three of these options are settable for the whole
document, for individual entries, or for individual entry types.

\mylittlespace The first is \texttt{shortextrafield}, which
\emph{must} be set in order for the mechanism to print anything at
all.  You should set this option to the name of the field you wish to
be printed in addition to the \textsf{author} and \textsf{labeltitle}.
(Possibilities include, but are not limited to, the 4 \textsf{*date}
fields and the 4 \textsf{*time} fields, the latter of which will print
the appropriate \textsf{*date} and the \textsf{*time}.)  By default,
it will be printed after the latter, separated from it by a comma.
You can manually define this punctuation by setting the
\texttt{shortextrapunct} option to one of \texttt{none, space, comma,
  period, colon,} or \texttt{semicolon}.  You can also enclose the
extra field in parentheses or square brackets by setting the
\texttt{shortextraformat} option to \texttt{parens} or
\texttt{brackets}.

\mylittlespace User \mymarginpar{\texttt{shorthand-\\first}} laudecir
requested a simpler way to print the \textsf{shorthand} even in the
first citation of a source, simpler, that is, than remembering to use
the \cmd{shorthandcite} command.  You can set this option to
\texttt{true} either in the preamble or in individual entries.

\mylittlespace Kenneth Pearce \mymarginpar{\texttt{shorthandfull}} has
suggested that, in some fields of study, a list of shorthands
providing full bibliographical information may replace the
bibliography itself.  This option prints this full information in the
list of shorthands, though of course you should remember that any .bib
entry not containing a \textsf{shorthand} field won't appear in such a
list.  Please see the documentation of the \textbf{shorthand} field in
section~\ref{sec:entryfields} above for information on further options
available to you for presenting and formatting the list of shorthands.

% %\enlargethispage{\baselineskip}

\mylittlespace Chris Sparks \mymarginpar{\texttt{shorthandibid}}
pointed out that \textsf{biblatex-chicago-notes} would never use
\emph{ibid.}\ in the case of entries containing a \textsf{shorthand}
field, but rather that consecutive references to such an entry
continued to provide the shorthand, instead.  The \emph{Manual} isn't,
as far as I can tell, completely clear on this question.  In 14.244,
discussing references to works from classical antiquity, it states
that \enquote{when abbreviations are used, these rather than
  \emph{ibid.}\ should be used in succeeding references to the same
  work,} but I can't make out whether this rule is specific to
classical references or has more general scope.  Given this ambiguity,
I don't think it unreasonable to provide an option to allow printing
of \emph{ibid.}\ instead of the shorthand in such circumstances,
though the default behavior remains the same as it always has.

\mylittlespace Fr.\ Norbert \mymarginpar{\texttt{shorthand-\\intro}}
Keliher requested a way to turn off the printing, in the first
citation of a work, of the introduction of a \textsf{shorthand} that
will appear in subsequent citations.  (A possible use case is when
that \textsf{shorthand} is so well known that it needs no
introduction.)  By default, \textsf{biblatex-chicago} prints a notice
--- (hereafter cited as) --- when no \textsf{shorthandintro} field is
present, the latter field allowing you to modify the notice but not
eliminate it.  The new \texttt{shorthandintro} option is available
globally, per-type, and per-entry, and has three possible values:

\begin{description}
\setlength{\parskip}{-4pt}
\item[\qquad none:] This setting turns off the notice entirely.  The
  \textsf{shorthand} field will simply appear in subsequent citations
  with no introduction.
\item[\qquad short:] This setting provides the \textsf{shorthand}
  alone in parentheses at the end of the first citation.  You can get
  a similar effect by providing a \textsf{shorthandintro} field which
  only contains the \textsf{shorthand} wrapped in parentheses.
\item[\qquad full:] This is in effect the default setting, but you
  would only need to use it in a document where one of the other two
  settings was in force either globally or per-type.  In
  \textsf{chicago-notes.cbx} I have set this value for the
  \textsf{jurisdiction}, \textsf{legal}, and \textsf{legislation}
  entry types to fulfil the \emph{Bluebook's} specification even when
  a user changes it for other sorts of entry.  This requires you to
  alter it specifically for these types should you wish to depart from
  the \emph{Bluebook} guidelines. Cf.\ section~\ref{sec:legal}, below.
\end{description}

N.B.\ The list of shorthands will always still be available to you as
an explanatory tool in the absence of any introductions to individual
shorthands.  Please see the documentation of the \textsf{shorthand}
and \textsf{shorthandintro} fields in section~\ref{sec:entryfields}.

\mylittlespace This \mymarginpar{\texttt{strict}} still-experimental
option attempts to follow the \emph{Manual}'s recommendations (14.41)
for formatting footnotes on the page, using no rule between them and
the main text unless there is a run-on note, in which case a short
rule intervenes to emphasize this continuation.  I haven't tested this
code very thoroughly, and it's possible that frequent use of floats
might interfere with it.  Let me know if it causes problems.

\mylittlespace Stefan \mymarginpar{\texttt{urlnotes}} Bj\"{o}rk, for
specialized reasons, requested a way to turn off the printing of
\textsf{url}, \textsf{doi}, and \textsf{eprint} information in notes
but not in the bibliography.  As it's possible this might be of more
general usefulness, I've provided a new option.  You can set it to
\texttt{false} either in the preamble or in individual entries, but
please note that it does not apply to \textsf{online} entries.

\mylittlespace In \colmarginpar{\texttt{useibid}} a change to previous
recommendations, the 17th edition of the \emph{Manual}
\enquote{discourages the use of \emph{ibid.,}}\ preferring instead a
shortened reference with only the author's name (14.34).
\textsf{Biblatex-chicago} now implements these recommendations by
default, including the repetition of page references even when they're
exactly the same as the previous note.  If you prefer to continue
using \emph{ibid.}, then set \texttt{useibid=true} in your document
preamble either globally or per type, or else in the \textsf{options}
field of individual entries, and you'll get the traditional behavior
(To be perfectly accurate, in entries with \texttt{classical}
\textsf{entrysubtype} and which have a \textsf{shortauthor}, the
\emph{ibid.}\ string is suppressed unless you set
\texttt{useibid=true} in the options field of the entry itself.  The
global setting will in this case be ignored.)

\mylittlespace Stefan \mymarginpar{\texttt{xrefurl}} Bj\"{o}rk pointed
out that when, using the \texttt{longcrossref} or
\texttt{booklongxref} options, you turn on the automatic abbreviation
of multiple entries in the same (e.g.) \textsf{collection} or
\textsf{mvcollection}, you could entirely lose a \textsf{url} that
might be helpful for locating a source, as the abbreviated forms in
notes and in the bibliography wouldn't include this information.
Setting this option to \textsf{true} either in the preamble or in
individual entries will allow the \textsf{url}, \textsf{doi}, or
\textsf{eprint} field to appear even in these abbreviated references.

\subsubsection{Back References: The \texttt{noteref} Option}
\label{sec:noteref}

\textsf{Biblatex} has always provided the \texttt{backref} option,
which prints, in the bibliography, those pages on which individual
works have been cited.  The \emph{Manual} (14.31) recommends another,
related system, which involves, at the end of short notes,
cross-references to the note where the reader can find the full, long
citation of the same source, \enquote{especially in the absence of a
  full bibliography.}  The general idea is that, where a short note is
\enquote{far} from the long citation, a back reference to that long
note may prove \enquote{helpful.}  The recommended format is something
like this: \texttt{(see chap.\ 1, n.\ 4)}.  The previous release of
\textsf{biblatex-chicago} provided something similar only for certain
subsets of material from the \emph{Bluebook} guidelines for legal
citations.  That provision is very basic and has a different
rationale, so it remains both unaltered and entirely separate.  With
this release, in the notes \&\ bibliography style only, I have
provided Chicago-style back references for all other entry types,
enabled through the \texttt{noteref} preamble option.  (Cf.\
\href{file:cms-noteref-demo.pdf}{\textsf{cms-noteref-demo.pdf}} for a
brief introduction.)

\mylittlespace Before embarking on a description of this new option,
and its many sub-options, I would like to point out that
\textsf{biblatex} provides a number of mechanisms designed to help
readers navigate long documents.  I have made many improvements to the
\textsf{biblatex-chicago hyperref} interface, so electronic documents
can, at your discretion, contain clickable links from short notes to
long notes or, in a document with all short notes, from such notes to
bibliography entries.  (The \textsf{noteref} mechanism cooperates well
with \textsf{hyperref}, and therefore can add another layer of links
to those already available.)  \textsf{Biblatex} also offers the
\texttt{refsection} and \texttt{citereset} preamble options, which
allow you to choose how its citation trackers behave.  Using these you
can, for example, always have a long note appear for a given source at
its first appearance in a chapter or a section, something which the
\emph{Manual} recommends in any case, and which may mean that your
short notes are never too \enquote{far} from a longer citation.  (See
\textsf{biblatex.pdf, \S~3.1.2.1}.)

\mylittlespace All \mymarginpar{\texttt{noteref}} the same, sometimes
chapters or sections can get rather long, or a too-frequent
reappearance of the long form may not be optimal for your work, so in
such situations the \texttt{noteref} option may well prove useful.
Its general principle is this: if a short note appears on the same
page as its corresponding long note, or on the same double-page spread
for \texttt{twoside} mode, then nothing will appear.  Similarly, if a
short note appears on the same page or double-page spread as a
previous short citation of the same source which \emph{does} have a
\texttt{noteref}, then this subsequent short citation will once again
\emph{not} present any \texttt{noteref}.  (This behavior is
configurable --- please see below.)  If a back reference is to be
printed, then the value of the \texttt{noteref} option determines what
it will look like.  Its six possible values are:

{\renewcommand{\descriptionlabel}[1]{\qquad\texttt{#1}}
\begin{description}
\item[none:] This is the default, and will always produce a back
  reference like this: (see n.~1).  It may well be useful when you are
  using the \texttt{citereset} or \texttt{refsection} options and know
  that any short note will always be in the same chapter or section of
  the text as the full reference to which it points.
\item[page:] This will always produce a back reference like this: (see
  p.\ 1, n.\ 1).  It can be a tidy way of directing the reader, as
  page numbers are usually simpler to track than sections or chapters.
  It's also a good setting if you've set the \LaTeX\ \texttt{secnumdepth}
  counter so that sections aren't numbered.
\item[chapter:] This is the example provided by the \emph{Manual}, and
  as implemented here it produces two different sorts of back
  reference.  If the short note is in the \emph{same} chapter as the
  long note to which it points, then by default it will only mention
  the note number, as with the \texttt{none} switch, above.  If the
  long reference is in a different chapter, then it prints like so:
  (see chap.\ 1, n.\ 4).  All of the options that name an organizing
  division of the text work the same way in footnotes, i.e., only when
  the short note and the long note to which it refers are in different
  \texttt{parts, chapters, sections, or subsections} will the actual
  division type appear in the \texttt{noteref}.  If you want the
  longer form in \emph{all} of your \texttt{noterefs}, you can set
  \mymarginpar{\texttt{fullnoterefs}} the \textsf{biblatex-chicago}
  option \texttt{fullnoterefs} to \texttt{true} when loading the
  package.  In endnotes, depending on which options you've chosen for
  presenting them, you may well never get the short version of the
  back reference.  Please see the details starting on
  page~\pageref{sec:endnoterefs}, below.
\item[section:] This key is particularly intended for documents, like
  the standard \LaTeX\ \textsf{article} class, which don't offer
  chapters, but rather start their divisions at the section level, but
  it's perfectly usable even in a document that also uses chapters.
  Assuming the short note and its long antecedent are in different
  sections, the \texttt{noteref} will look like so: w/o chapter (see
  \S~2, n.\ 6), w/\ chapter (see \S~1.2, n.\ 6).
\item[subsection:] I'm not sure there's any need for this key, but I
  include it for the sake of completeness.  It's usable in documents
  both with and without chapter divisions, and assuming the short note
  and its long antecedent are in different subsections the
  \textsf{noteref} will look like so: w/o chapter (see \S~3.2, n.\ 5),
  w/\ chapter (see \S~2.3.2, n.\ 5).
\item[part:] This is, I suspect, even less likely to be useful than
  \texttt{subsection}, but assuming the short note and its long
  antecedent are in different parts the back reference will look like
  so: (see pt.\ I, n.\ 4).  You'll need to be careful that note
  numbering is continuous across chapters for this to work correctly,
  otherwise the plain note number might well be ambiguous.  Also, if
  you'd like the part number not to be roman, you can try putting this
  in your preamble: \verb+\let\cmsnrpart\relax+.
\end{description}}

Several comments are in order, before moving on to the onerous
details.  In case it's not already clear, the \texttt{noteref} option
is only relevant if your document includes full notes, that is, if you
aren't using the \texttt{short} option.  Even in documents that use
long notes, it can occasionally happen that the \texttt{noteref} code
won't be able to find a full citation of a particular source.  In this
case, no back reference will appear, and you will find a warning in
your\ .log file informing you about it.  (If you combine
\texttt{short} and \texttt{noteref}, you'll see a lot of such
warnings.)  It can happen that even the first citation of a particular
source appears in a somewhat abbreviated form, as when multiple
contributions to the same \textsf{collection} are present in your
reference apparatus.  The \texttt{noteref} will point to this
abbreviated first citation all the same, given that it is at least
somewhat more informative than a short note.  Similarly, a
\texttt{noteref} from a \textsf{collection} may well point to the long
citation of an essay from that collection, as that long citation will
contain all the details of the collection, too.  I hope this doesn't
prove too surprising.  I should also clarify that all of the strings
in the \texttt{noterefs} as printed above are localized, so should
adapt to your document language reasonably smoothly, if not entirely
idiomatically.  Finally, the code assumes that the standard \LaTeX\
counters for parts, chapters, and sections are available, which I
believe is almost universally the case even for classes and styles
that redefine a lot of the relevant functionality, but I confess I
haven't tested \texttt{noteref} at all extensively against the
possibilities offered by CTAN, so please do let me know if something
breaks for you.

\enlargethispage{\baselineskip}

\mylittlespace I \mymarginpar{\texttt{noterefinterval}} mentioned
above that the gap between appearances of a \texttt{noteref} for a
given source was configurable.  What I had in mind was this option,
along with four new citation commands which I discuss below.  The
\texttt{noterefinterval} preamble option allows you to define the
number of references that must have intervened since the last
\texttt{noteref} before another to the same source will appear.  If
you judge that your readers don't need a pointer on every new page but
only after a certain number of other citations have passed, you can
set this to a number higher than zero (the default).  You can't,
currently, use this mechanism to make new pointers appear on the same
page as previous ones, but you can spread them out if they appear too
frequently for your tastes.  Also, the counter that this option uses
is \texttt{instcount}, which will be incremented not only by new
references but also, e.g., by uses of the \textsf{related}
functionality to extract data from other entries.  A value of
\texttt{15} may not delay a reappearance for exactly that many notes,
so you will need to experiment a little to find a value that suits
your document.

\mylittlespace If \mymarginpar{\cmd{shortrefcite}\\
  \cmd{shorthand-}\\\texttt{refcite}\\ \cmd{shortcite*}\\
  \cmd{shorthandcite*}} you require more fine-grained control over the
spacing between \texttt{noterefs}, or indeed if you want them to
appear more frequently than the previous mechanisms allow, then these
new citation commands will allow you to do so, though perhaps without
maximum convenience.  (I suppose that it would be safest to introduce
these commands into your documents at quite a late stage in their
preparation.)  The first two commands present, respectively, short
notes and \textsf{shorthand} notes where the \texttt{noteref} will
absolutely be printed (unless, of course, no full citation can be
found).  The second two commands prevent the printing of the
\texttt{noteref}, no matter where the resulting note appears.  All of
them will need enclosing in a \cmd{footnote} command if you want them
to appear in one, as I've provided only the most general form of each.
I've \colmarginpar{\texttt{suppressnoterefs}} also, after a user
request, provided the \mycolor{\texttt{suppressnoterefs}} option,
settable per entry type and per entry, which prevents any
\texttt{noteref} being printed for the entries concerned.  This may be
more convenient than the specialized citation commands.

\subsubsection*{Zero Sections}
\label{sec:zero}

The \LaTeX\ sectioning counters all start from zero, so if you put a
note into material occurring \emph{before} the first \cmd{part},
\cmd{chapter}, \cmd{section}, or \cmd{subsection} command then any
back reference to this citation will, by default, present that zero
(or zeros).  This may, in fact, be exactly what you want, in which
case you can ignore the following options.  If you don't want a zero
to appear in your \texttt{noterefs}, you can either make sure no
citations occur in contexts that will produce them, or you can use a
combination of the next three options to hide them.

\mylittlespace This \mymarginpar{\texttt{noterefintro}} option is
special in that it handles only zeros that occur in the \emph{first}
position in a sectioning identifier, e.g., \S~0.x.x or chap.~0.  It's
possible that this zero represents some sort of introductory material
before, e.g., the numbered chapters appear.  If you say
\texttt{noterefintro=introduction}, then instead of (see chap.~0,
n.~1), you'll have (see intro., n.~1).  If the value of the option is
a \cmd{bibstring} known to \textsf{biblatex}, then it will appear in
localized, and possibly abbreviated, form, as above.  If the section
title you want isn't a known \cmd{bibstring}, you can either define a
new one for your language in your preamble, or you can just set the
option to whatever it is you want to appear in such \texttt{noterefs}.
Both will work, particularly because you don't need to worry too much
about capitalization because the word always appears after the
\verb+\bibstring{see}+.

% %\enlargethispage{\baselineskip}

\mylittlespace Now, it's perfectly possible for an introduction to
have numbered sections of its own, so a citation there may produce a
back reference like \S~0.1 or \S~0.1.0.  The rules are: 1.\ any back
reference which is all zeros will just print the \texttt{noterefintro}
string alone, assuming you've provided one; 2.\ any back reference that
has the zero \emph{only} in the first place will print something like
(see intro., \S~1, n.~1); 3.\ a back reference of the form 0.x.0 or
0.0.x --- this can only occur if \texttt{noteref=subsection} and the
document class provides a \cmd{chapter} command --- such a reference
will either print the zero after the \texttt{noterefintro} string, or
you can use the \texttt{pagezeros} option, to which I turn.

\mylittlespace This \mymarginpar{\texttt{pagezeros}} boolean option
deals with the problem of zero sections by changing such back
references, and only such back references, to behave as though
\texttt{noteref=page}.  In the absence of a \textsf{noterefintro}
\emph{all} such zero citations will be so treated, but if both options
are set then zero sections with the zero in the first position of the
identifier will fall first under the jurisdiction of the
\texttt{noterefintro} option, only then turning to \texttt{pagezeros}
if there's a zero remaining that hasn't yet been eliminated by the
first option's rules.  Some examples:
\begin{verse}
  \textbf{Without noterefintro:}\\
  Any zero ---> (see p.~3, n.~1)\\
  No zero ---> (see \S~1.2.3, n.~1)

  \textbf{With noterefintro:}\\
  1.2.3 ---> (see \S~1.2.3, n.~1)\\
  0.0.0 ---> (see intro., n.~1)\\
  1.1.0 ---> (see p.~24, n.1)\\
  0.1.0 ---> (see intro., p.~2, n.~1)
\end{verse}

As you see, this produces a mixed system of back references, so you'll
need to decide whether you and your readers might still find it
acceptable.

\mylittlespace This \mymarginpar{\texttt{hidezeros}} boolean option
attempts, with varying degrees of success, to disguise the zeros in
section identifiers without mixing different sorts of back reference
in one document.  Unlike \texttt{pagezeros}, it will never modify
identifiers where the first number is zero.  It leaves all such
identifiers to the \texttt{noterefintro} option, so if your identifier
looks like 0.1.0, the second zero will still appear.  The only thing
you can do about it is to move the citation so that it isn't in a zero
section.  The rationale is that \texttt{hidezeros} places a string ---
by default \enquote{pref.} --- \emph{inside} the section identifier so
that a reader knows that the citation occurred in the prefatory
material to a particular section.  Combining this with another string
for the prefatory material to a whole work is unattractive, and I've
avoided it.  Here are some examples of how it looks:
\begin{verse}
  1.1.1 ---> (see \S~1.1.1, n.~1)\\
  1.0.1 ---> (see \S~1.pref.1, n.~1); should this even be possible?\\
  1.1.0 ---> (see \S~1.1.1 [pref.], n.~1)\\
  1.0.0 ---> (see \S~1.1 [pref.], n.~1)\\
  1.0 ---> (see \S~1.1 [pref.], n.~1)
\end{verse}

The \mymarginpar{\cmd{cmspref}} brackets and placement of the
identifying string are hard coded, but you can change the string
itself with a line something like this in your preamble:\\
\verb+\renewrobustcmd{\cmspref}{\emph{pref.\@}}+ <--- \textbf{NB} the
\verb+\@+
after the dot\\
The \texttt{hidezeros} method at least produces more uniform back
references, though it perhaps sacrifices something in immediate
readability in order to do so.  I would be glad to entertain
suggestions for other solutions.

\mylittlespace If your documentation uses footnotes, then the guide to
the \texttt{noteref} option(s) concludes here.  I have, however, been
determined, at least in this context, to provide for endnote users the
same features as for footnote users, mainly because the \emph{Manual}
caters equally to both.  Unfortunately, because footnotes are a core
part of \LaTeX\ formats and endnotes are provided by means of extra
packages, users of the latter will face some additional complications
if they wish to provide Chicago-style back references.  I document
these forthwith.

\subsubsection*{Endnotes and \texttt{noteref}}
\label{sec:endnoterefs}

The traditional way to provide endnotes instead of footnotes in a
document is to use the \textsf{endnotes} package by John Lavagnino,
and \textsf{biblatex} provides a reliable interface to that package,
making it relatively simple to use and control (cf.\ esp.\ the
\texttt{notetype} option in \textsf{biblatex.pdf, \S~3.1.2.1}).  The
package does have its limitations --- on which see more below --- so I
did have a look at its next-generation replacement, Clemens
Niederberger's \textsf{enotez}.  This adds all the needed
functionality, I think, and is also more future proof, relying as it
does on the work of the \LaTeX 3 project.  The downside is that my
\TeX nical abilities fell short of being able to make it work with the
\texttt{noteref} option, so I have instead created a new package
\mymarginpar{\small\texttt{\textbf{cmsendnotes.sty}}} which combines
functionality from \textsf{endnotes}, from Ulrich Dirr's
\textsf{hyperendnotes}, and from \textsf{biblatex-chicago}.  If you
need some functionality that \textsf{endnotes} doesn't provide, then
you can load \textsf{cmsendnotes} instead.  (I haven't tested any of
this with the \textsf{memoir} class, which has its own endnote
mechanism, so please let me know if it doesn't work and I'll try to
have a look.)  The documentation that follows should clarify when you
might want to load the new package, and also the options available to
get \textsf{cmsendnotes} to do what you want.

\mylittlespace Before we begin, I should just point out that, as usual
with \textsf{biblatex}, you can mix foot- and endnotes in the same
document, but if \texttt{noterefs} are going to appear in both sorts
of note --- surely this situation is highly unlikely --- then you need
to be careful that they refer back \emph{only} to long references in
the \emph{same} sort of note.  A \texttt{noteref} from an endnote to a
long citation in a footnote will be inaccurate, and vice versa, so
careful use of the \cmd{citereset} command (as in
\href{file:cms-noteref-demo.pdf}{\textsf{cms-noteref-demo.pdf}}) or
perhaps of the \textsf{biblatex} \texttt{citereset} option should
allow you to keep the two sorts of note distinct.

\mylittlespace The principle to keep in mind is that back references
to full endnotes point not to the place in the main text where you've
cited a source, but rather to the place where that citation is
actually printed, which may well be in another division of your
document altogether.  If you are providing endnotes at the end of each
chapter, or (less likely) at the end of each section of a long
article, then this means that an endnote to a later chapter or section
will point to the earlier chapter or section after which the full
citation was printed.  This interacts very well with the
\textsf{endnotes} package's \cmd{theendnotes} command, which prints,
and then clears, all the endnotes created up to the point at which you
call it.  Subsequent calls to \cmd{theendnotes} do the same, and short
notes will always have a reasonably accurate sense of where their
antecedent long note has appeared, i.e., in the endnotes to a
particular chapter or section.  (This even has the side effect of
making the zero section problem somewhat more tractable, as the back
reference doesn't mind that the \cmd{endnote} command occurs in
section 1.0, but rather that the citation appears in the notes to
section 1.1.)

\mylittlespace Similarly, if your endnotes appear all together at the
end of an article, then you can just use the \texttt{page} option to
\texttt{noteref}, or no option at all, and the back references will be
both accurate and usable (assuming the notes are all numbered
consecutively, I suppose, which seems a safe assumption).  The upshot
is that, if you are providing endnotes in either of these scenarios,
both of them envisaged by the \emph{Manual} (14.46), and \emph{either}
if you don't need the \textsf{hyperref} functionality, \emph{or} if
the somewhat restricted functionality available through the
\textsf{endnotes} package is good enough, then you can happily ignore
the new \textsf{cmsendnotes} package entirely.

\mylittlespace If back references are, in the scenarios discussed
above, basically working for you, but you want more elaborate
\textsf{hyperref} functionality, then you can load the
\textsf{cmsendnotes} package without any options instead of the
\textsf{endnotes} package.  Please be aware, however, that
\mymarginpar{\small\textbf{NB}} you must load \textsf{cmsendnotes}
\emph{after} \textsf{biblatex-chicago} for it to work properly.  What
you'll then get by default, assuming you've loaded \textsf{hyperref},
are links from endnote numbers in the main text to the corresponding
numbers in the endnotes section itself, and vice versa, along with
more accurate links from the back references to particular endnotes in
preceding sections or chapters.  There are several options available
for changing the default appearance of your endnotes, four of which
are package options to \textsf{cmsendnotes} and two of which are
commands that you can redefine to your liking.

\mylittlespace Two \mymarginpar{\texttt{hyper}} package options
control the \textsf{hyperref} behavior of endnote numbers.  They are
both set to \texttt{true} by default, if the \textsf{hyperref} package
is loaded.  If you set the first, \texttt{hyper}, to \textsf{false}
when loading \textsf{cmsendnotes}, there will be no hyperlinking of
endnote numbers at all.  If you \mymarginpar{\texttt{enotelinks}} set
the second, \texttt{enotelinks}, which I've borrowed from
\textsf{hyperendnotes.sty}, to \texttt{false}, then only endnote
numbers in the main text will function as links, the numbers in the
endnote sections themselves ceasing to act as such.

\mylittlespace This \mymarginpar{\texttt{noheader}} package option, if
set to \texttt{true}, stops the printing of the usual section header
before the endnotes themselves, in case this might help to solve some
formatting problems in your documents.

\mylittlespace This \mymarginpar{\cmd{enoteheading}} is the standard
\textsf{endnotes} package command for defining exactly what is printed
in the heading.  I have slightly redefined it (for reasons I shall
explain later), but you can redefine it in your preamble if you wish,
and that will be respected by \textsf{cmsendnotes}.

\mylittlespace This \mymarginpar{\texttt{blocknotes}} package option,
if set to \texttt{true}, presents the text of each endnote as a
flush-left block, i.e., without the first line being indented.

\mylittlespace This \mymarginpar{\cmd{enoteskip}} command, which was
inspired by a similar provision in \textsf{hyperendnotes.sty}, sets
the vertical space between individual endnotes.  By default it doesn't
change basic interline space, but you can define it in your preamble
to something like \cmd{smallskip} or \cmd{medskip} in case you want a
bit more light inside your endnote sections.

\mylittlespace So far, then, I have discussed contexts where using
\textsf{cmsendnotes.sty} only brings cosmetic changes to functionality
which basically already works using \textsf{endnotes.sty}.  Indeed, if
you are using either of these packages in the ways already outlined,
then the setting of the main \texttt{noteref} option defines how your
back references will look, and \mymarginpar{\texttt{fullnoterefs}} the
\texttt{fullnoterefs} option will still govern the \texttt{chapter},
\texttt{section}, \texttt{subsection}, and \texttt{part} values of the
\texttt{noteref} option, just as in the general discussion above.  In
the methods discussed below, additional steps are required for
defining how your back references will look, and the
\texttt{fullnoterefs} option is irrelevant, as the back references
will always appear in their fullest form.

\mylittlespace The methods of endnote presentation to which I now turn
involve, in the \emph{Manual's} words, when notes to \enquote{each
  chapter of a book are \ldots\ grouped in the end matter} (14.46).
Here, you would want not only a main heading for the endnotes section
but also \enquote{a subhead bearing the chapter number or title or
  both.}  It is perfectly possible to achieve the subdivision and
subheading of a long endnotes section by using \textsf{endnote's}
\cmd{addtoendnotes} command in each chapter of your document, putting
a sectioning command of some sort in its argument, for example.  Using
\texttt{noteref} back references in this context can be a little
complicated, however, mainly because of the principle I explained
above, i.e., that back references point to the place where the long
note was printed, not to the section of the main document where the
source was actually cited.  For our purposes, this means that, unless
you have set \texttt{noteref} to \texttt{page} or \texttt{none},
\textsf{biblatex-chicago} actually has to extract data from the
sectioning command you've included in \cmd{addtoendnotes} and, most
frequently, modify that data to make it work inside a \texttt{noteref}
back reference in way that is both consistent and useful.  The
\textsf{cmsendnotes} package tries to automate this process as much as
possible so that with, in the best-case scenario, only one option
given to the package the whole system can be made to work without
further user intervention.  Further package options can help with
slightly more complicated scenarios, but if your requirements are more
complex than the automatic system can provide, then there are two ways
to handcraft a divided endnote section: one uses traditional
\cmd{addtoendnotes} functionality from the \textsf{endnotes} package,
and the other uses new commands available from \textsf{cmsendnotes}.
I believe the second handcrafted option to be slightly more convenient
than the first, but in any case I'll start by explaining the automatic
provisions, then move on to the two handcrafted options, leaving you
to judge which seems best suited to your needs.

%%\enlargethispage{\baselineskip}

\mylittlespace For \mymarginpar{\texttt{\textbf{split}}} the automatic
subdivision of an endnotes section I have borrowed a concept, if not
its implementation, from \textsf{enotez}, and provided the
\textsf{cmsendnotes} option \texttt{split}, which has 4 possible
values: \texttt{part}, \texttt{chapter}, \texttt{section} and
\texttt{subsection}.  (If you don't provide a key, it defaults to
\texttt{chapter}.)  With this option set, you need to use a new
command for printing the endnotes, the ungainly
\mymarginpar{\cmd{\textbf{theendnotes\-bypart}}} but I hope memorable
\cmd{theendnotesbypart}.  When you do this, \textsf{cmsendnotes} does
something in the background that's worth understanding.  Ordinarily,
in the standard \textsf{endnotes} package, any call to
\cmd{theendnotes} produces an\ .ent file containing all of the endnote
data in the document up to that point, and proceeds to print it.
Another call to \cmd{theendnotes} gathers the endnote data occurring
between it and the first call, overwriting the\ .ent file, and again
printing it, and so on.  Whenever you use any version of
\cmd{theendnotesbypart}, \textsf{cmsendnotes} will write one\ .ent
file per section named by the \texttt{split} option, assuming that
said section actually contains any endnotes.  The plain
\cmd{theendnotesbypart} command, with no further options, proceeds to
print, in sequence, \emph{all} the\ .ent files in numerical order.  In
the first instance, then, the value of the \texttt{split} option
decides how your endnotes are distributed to different\ .ent files.
These files are named after the main document suffixed with the number
of the section, e.g., \textsf{jobname1.ent} for chapters or
\textsf{jobname1.3.ent} for sections.

\mylittlespace In the second instance, the \texttt{split} option
governs what the subheaders of your endnotes section will contain.
The main header is produced by a \cmd{section*} command, and by
default it will read \texttt{\textbf{Notes}}.  The subheaders are
produced by \cmd{subsection*} commands, and will take their name from
the \textsf{split} option and their number from the number of the\
.ent file currently being processed: \texttt{\textbf{Chapter 1}}
\ldots\ \texttt{\textbf{Chapter 2}}.  The headers are localized,
assuming you're using a language supported by
\textsf{biblatex-chicago}.  Even if you're not using \texttt{noteref}
back references in your document, this mechanism can still provide a
convenient means of subdividing an endnotes section.

\mylittlespace If you are using \textsf{noteref}, then the value of
that option leads to two possible outcomes.  If set to \texttt{page}
or \texttt{none}, any back references will point to full notes by page
plus note number or just by note number, as usual, bearing in mind
that the page involved is where the full note was printed, not where
it was cited in the main body of the text.  If set to any of the other
possible sections of your document, then the \texttt{split} option has
a third function, which is to provide the (localized) string for the
back reference itself --- (see chap.~1, n.~2) --- which will refer to
a subsection of the endnotes section named \texttt{\textbf{Chapter~1}}
rather than to the actual first chapter of the main document.  This
setup will usually involve setting \texttt{split} to the same value as
the \texttt{noteref} option itself, but if they differ, and
\texttt{noteref} isn't \texttt{page} or \texttt{none}, then
\texttt{split} takes precedence and governs the appearance of the back
reference.

\mylittlespace Let's say, then, that your document is in English and
you've set the \textsf{biblatex-chicago} option
\texttt{noteref=chapter}, and the \textsf{cmsendnotes} option
\texttt{split}, then what you can expect to see, when you use
\cmd{theendnotesbypart}, is something like this, subject to the usual
rules for the appearance or non-appearance of \texttt{noterefs}, and
remembering that in this context, as noted above, all
\texttt{noterefs} will appear in their long form:
\begin{verse}
  \textbf{Notes}\\

  \textbf{Chapter~1}\\
  1. Book.\\
  2. Article.\\
  3. InCollection.\\

  \textbf{Chapter~2}\\
  1. InCollection (see chap.~1, n.~3).\\
  2. BookInBook.\\
  3. Article (see chap.~1, n.~2).\\

  \textbf{Chapter~3}\\
  1. BookInBook (see chap.~2, n.~2).\\
  2. Book (see chap.~1, n.~1).
\end{verse}

It won't, unfortunately, always be this simple, but it may be a
comfort to know that some of the complications are the same as those
faced by users of \texttt{noteref} with footnotes, in particular
\mymarginpar{\textbf{Zero Sections}} the zero section problem.  To
deal with this issue you use the same options to
\textsf{biblatex-chicago} as you would for footnotes, with one
difference.  In the case of split endnotes, the code has to handle the
zeros both in the \cmd{subsection*} names \emph{and} in the back
references, which means that the \texttt{pagezeros} option is no
longer relevant, as it can't do the right thing in section names.
This leaves the \texttt{noterefintro} option for hiding zeros that
start a section number, and the \texttt{hidezeros} option for zeros
anywhere else.  These \textsf{biblatex-chicago} options work here just
as described in the footnote section above, but with one or two
additional caveats.

\mylittlespace First, I probably shouldn't have bothered trying to
implement the combination of \texttt{noterefintro} and
\texttt{hidezeros}, as any sections of a zero chapter in your document
will appear by default with zeros intact in the text itself, unless
measures are taken.  I did \mymarginpar{\cmd{cmsintrosection}} do this
thing, however, and part of the implementation is a command
\cmd{cmsintrosection}, which provides the identifying string for
subsections of the introduction.  It is set by default to \S, and
though you can redefine it in your preamble, please remember that it
will appear as such in both \cmd{subsection*} names and
\texttt{noterefs}.

\mylittlespace Indeed, it is the need to cater for two quite distinct
contexts that makes the automatic provision of \textsf{noteref} back
references in a divided endnotes section surprisingly tricky.  You
need \mymarginpar{\cmd{introduction\-name}\\\cmd{forewordname}\\
  \cmd{sectionname}\\ \cmd{subsectionname}} one mechanism to take
\texttt{chapter} and turn it into Chapter, and quite another to turn
it into chap., bearing in mind that \cmd{bibstrings} don't work
outside of the reference apparatus, and therefore not in
\cmd{subsection*} names, though obviously they're perfect for the back
references themselves.  My solution is to borrow a principle from
\textsf{babel}, which provides for its languages commands like
\cmd{prefacename} and \cmd{chaptername}, which print the localized
version of the term, usually capitalized.  In the\ .lbx files that
come with \textsf{biblatex-chicago} I have added
\cmd{introductionname} and \cmd{forewordname}, along with
\cmd{sectionname} and \cmd{subsectionname}, so at least the most
common types of prefatory material, when provided to the
\texttt{noterefintro} option, should work properly both in headings
and in \texttt{noterefs}, and across languages.

\mylittlespace So, another caveat.  Should you wish to provide a
\texttt{noterefintro} value that isn't a standard \cmd{bibstring} or
doesn't have a corresponding \cmd{[value]name} command, then it may
not work well for you in one or other of the two contexts in which it
can appear.  The code does test whether the bibstring and command
exist, and it will capitalize anything in section headers, but
otherwise you can just cunningly craft a string that's good in both
contexts or you can define a \cmd{[value]name} command and a new
\cmd{bibstring} for the value in your language, both in your preamble.

\enlargethispage{\baselineskip}

\mylittlespace The \mymarginpar{\texttt{subheadername}} same caveat
applies if you want to supply your own name for the \cmd{subsection*}
commands that divide up the general endnotes \cmd{section*}.  Let's
say for some reason you want subsections called \textbf{Further
  Remarks} instead of \textbf{Chapter}.  Strings of more than one word
are difficult for the code to manage correctly, so instead you could
include in your preamble lines looking approximately like this:
\begin{verbatim}
  \NewBibliographyString{furthrem}
  \DefineBibliographyStrings{american}%
       {furthrem = {furth\adddotspace rem\adddot},}
  \def\furthremname{Further Remarks}
\end{verbatim}

Then you could set \texttt{subheadername=furthrem} in the options to
\textsf{cmsendnotes} and you'll get what you want.  It's not wildly
convenient, but it's slightly less typing than the handcrafted options
I discuss below, though for anything more complicated you'll probably
need those options.

\mylittlespace Before \mymarginpar{\texttt{headername}} I move on to
the handcrafted methods, I should point out two more
\textsf{cmsendnotes} options.  The first, \texttt{headername}, sets
the name of the main \cmd{section*} command at the start of the
endnotes section.  It defaults to the usual \textsf{endnotes} package
command \cmd{notesname}, which gives \textbf{Notes} in English.  I
have kept this separate from the standard \cmd{enoteheading} because
it needs slightly different treatment in a divided endnotes section.
If the definitions I have provided of \cmd{notesname} in the\ .lbx
files that come with \textsf{biblatex-chicago} aren't to your liking,
you can provide a string here instead, which is simpler to do because
it shouldn't be turning up in any \texttt{noterefs}.  If you'd like to
redefine any of the \cmd{*name} commands, the best place to do so is
very near to where you actually print the endnotes, where it can
override the definitions in the\ .lbx files (or in \textsf{babel's}
files).  Remember, too, that you can use the \texttt{noheader} option
to turn off the printing of this header if you just want to provide
your own sectioning command instead.

\mylittlespace This \mymarginpar{\texttt{runningname}}
\textsf{cmsendnotes} option controls the text that appears in running
headers in the endnotes section of your document, should you be
providing them.  I have followed the style of the \textsf{endnotes}
package, so that the default reads something like: \emph{NOTES TO
  CHAPTER 1}.  The section name and number are controlled by the other
options already discussed, but the \enquote{Notes to} part is
controlled by \texttt{runningname}, so if your document isn't in
English, and/or you're unhappy with the default string, you can change
it when loading \textsf{cmsendnotes}.

\mylittlespace Should \mymarginpar{\texttt{\textbf{endnotesplit}}} the
options above not fulfil your needs, you can control more or less all
parts of the subdivision of your endnotes section, of the running
headers there, and of back references from short notes to full ones,
by providing your own sectioning commands in your document.  If you
wish to use \texttt{noteref} back references in this context, you
\mymarginpar{\textbf{NB}} must first set the \textsf{biblatex-chicago}
option \texttt{endnotesplit} to \texttt{true}, no matter which of the
two possible implementation methods you choose.  With the standard
\textsf{endnotes} package you would then use the command
\cmd{addtoendnotes}, while with \textsf{cmsendnotes} it involves
variants of the endnote-printing command \cmd{theendnotesbypart}.
(Please note, first, that \cmd{addtoendnotes} still works with
\texttt{cmsendnotes}, in case that's useful to you; and second, that
with any other endnote implementation, you'll have to consult its
documentation to see if there's a compatible means of dividing the
notes.)

\mylittlespace First, \mymarginpar{Handcrafting
  w/\\\texttt{cmsendnotes}} I introduce the methods provided by
\textsf{cmsendnotes}.  The command \cmd{theendnotesbypart} has three
variants.  The first, \cmd{theendnotesbypart*}, simply suppresses the
printing of the \texttt{headername}, so it works more or less like
setting \texttt{noheader} in the options to \textsf{cmsendnotes}.  The
other two involve an optional argument,
\mymarginpar{\cmd{theendnotes\-bypart[]}} in square brackets,
containing an individual section number, which prints the endnotes
from that section.  This command never prints the general endnote
section header (as controlled by the \texttt{headername} option), but
it will print the individual
\mymarginpar{\cmd{theendnotes\-bypart*[]}} section's subheader, as
controlled by the \texttt{subheadername} option.  To turn that
printing off you can either use the starred version of the command,
i.e., \cmd{theendnotesbypart*[]}, or you can set the
\textsf{cmsendnotes} option \mymarginpar{\texttt{nosubheader}}
\texttt{nosubheader} to \texttt{true}.  A sequence of commands, each
with one section of the document inside square brackets, will give you
a complete endnotes section wherever you decide to place it, while the
starred forms or \texttt{nosubheader} option allow you to create your
own subheaders before each subsection.

\mylittlespace First, \mymarginpar{\textbf{NB}} please note that what
you need to place inside the square brackets is the \emph{number that
  forms part of the name of the\ .ent file in your working directory.}
In other words, it's the number \emph{before} any manipulations by the
\textsf{cmsendnotes} package remove zeros from it.  Depending on the
setting of your \textsf{split} option your commands may look like:

\begin{verbatim}
Chapters                             Sections
\theendnotesbypart*[0]          \theendnotesbypart*[0.0] <-- "introname"
\theendnotesbypart*[1]          \theendnotesbypart*[1.1]
\theendnotesbypart*[2] etc.    \theendnotesbypart*[1.2] etc.
\end{verbatim}

When in doubt, have a look in your working directory for the\ .ent
files produced for your document, and use the numbers from there.  (If,
for some reason, you decide to split by \texttt{part}, you'll probably
have roman numerals there, for example, apart from the zero.)

\mylittlespace The \mymarginpar{\textbf{Sectioning}} next step is to
provide some sort of sectioning command for the subheaders and for the
\texttt{noteref} back references.  This is slightly complicated, but
works the same whether you're using \textsf{cmsendnotes} or
\textsf{endnotes}.  The basic principle is that the main name of the
section appears in the endnotes section, while the optional name
provided for the table of contents [toc] appears in the
\texttt{noteref}:

\begin{verbatim}
\subsubsection[chap. 1]{Chapter 1} --> Chapter 1 ... (see chap. 1, n.1)
\end{verbatim}

You'll notice that the sectioning command isn't starred, as only
unstarred commands provide the optional [toc] argument.  (The
\cmd{addcontentsline} command can also be used with starred forms, but
keeping the [toc] argument out of the actual table of contents remains
an issue, so please read on.)  The unusual form of the [toc] argument
would merely pollute any table of contents you want to provide, and
the actual header in your endnotes section shouldn't have a number in
it provided by the standard \LaTeX\ methods, so you'll have to pick a
section type that falls underneath the thresholds of the \LaTeX\
counters \texttt{tocdepth} and \texttt{secnumdepth}.  By default, in
the standard \texttt{book} and \texttt{report} classes,
\cmd{subsubsection} works for this, while in the \texttt{article}
class you may need \cmd{paragraph}.  (You could, of course, also
change the counters, should you wish.)

\mylittlespace So, let's say you want to subdivide your endnotes
section with subheaders containing both the chapter number and the
chapter title, as envisaged by the \emph{Manual} (14.46).  Your
endnotes section might start like this:

\begin{verbatim}
\section*{Notes}
\subsubsection[intro.]{Introduction: The History of the Problem}
\theendnotesbypart*[0]
\subsubsection[chap. 1]{Chapter 1: Renewing the Question}
\theendnotesbypart*[1]   (etc.)
\end{verbatim}

There remain a couple of formatting issues with this code.  The
\textsf{endnotes} package points out that the first endnote after such
a sectioning command won't be indented properly, so it and
\textsf{cmsendnotes} use \verb+\mbox{}\par\vskip-\baselineskip+ after
sectioning commands to prevent this.  Additionally, both packages
provide code for running headers using \cmd{@mkboth}, so if you use
such headers you can either do the same inside \cmd{makeatletter} and
\cmd{makeatother} commands or just use \cmd{markboth}.  Taking all of
this into account gives code looking something like this, perhaps:

\begin{verbatim}
\section*{Notes}
\subsubsection[intro.]{Introduction: The History of the Problem
\markboth{NOTES TO INTRODUCTION}%
      {NOTES TO INTRODUCTION}}%
\mbox{}\par\vskip-\baselineskip
\theendnotesbypart*[0]
\subsubsection[chap. 1]{Chapter 1: Renewing the Question
\markboth{NOTES TO CHAPTER 1}%
      {NOTES TO CHAPTER 1}}%
\mbox{}\par\vskip-\baselineskip
\theendnotesbypart*[1]   (etc.)
\end{verbatim}

One of the, perhaps minor, advantages of using the
\textsf{cmsendnotes} commands for this is that they will at least all
typically be grouped together in one place in your document, rather
than scattered throughout, as when you use \textsf{endnotes'}
\cmd{addtoendnotes} command, to which we now turn.

\mylittlespace To \mymarginpar{Handcrafting w/ \textsf{endnotes}} use
the \textsf{endnotes} package with its main command \cmd{theendnotes}
to produce a subdivided endnotes section, you must first remember to
set the \textsf{biblatex-chicago} option \texttt{endnotesplit} to
\texttt{true}, that is, assuming you want to provide \texttt{noteref}
back references.  For splitting the endnotes, you
\mymarginpar{\cmd{addtoendnotes}} need the \cmd{addtoendnotes}
command, which you have to place in your document yourself.
Ordinarily, you'll need one such command for each relevant division of
your text, placed just after the sectioning command itself, so that
any endnotes that occur in the section will appear grouped underneath
the heading you provide.  At the next section, another such command
starts a new subsection of endnotes.

\mylittlespace To provide the same endnotes section divided by chapter
that we've already discussed above, your commands will look something
like this:

\begin{verbatim}
\chapter*{Introduction: The History of the Problem}
\addcontentsline{toc}{chapter}{Introduction}
\addtoendnotes{%
  \protect\subsubsection[intro.]{Introduction: The History of the%
   Problem%
    \protect\markboth{NOTES TO INTRODUCTION}%
    {NOTES TO INTRODUCTION}}%
   \mbox{}\par\vskip-\baselineskip} ...
\chapter{Renewing the Question}
\addtoendnotes{%
  \protect\subsubsection[chap. 1]{Chapter 1: Renewing the Question%
    \protect\markboth{NOTES TO CHAPTER 1}%
    {NOTES TO CHAPTER 1}}%
   \mbox{}\par\vskip-\baselineskip} ... (etc.)
\end{verbatim}

The commands you use are the same as with \cmd{theendnotesbypart[]},
but in this context both the sectioning command \emph{and} the command
for running headers need to be \cmd{protected}.  I would also
recommend redefining \cmd{enoteheading} in your preamble, as the
default definition produces too much extra vertical space before the
first subheading.  Something like
\verb+\def\enoteheading{\section*{Notes}}+ will do.  Finally, remember
that you can use \textsf{cmsendnotes} (without a \texttt{split}
option) instead of \textsf{endnotes}, if the extra \textsf{hyperref}
functionality is important to you.  The command sequence above will
continue to work in the same way.

\mylittlespace There are tradeoffs for both systems.  With
\textsf{endnotes}, at least the single \cmd{theendnotes} command keeps
things simple, but you still have to keep track of which sections have
endnotes in them, else spurious subheaders will appear.  Rooting
around in your working directory to make sure you've printed all the\
.ent files is annoying, but at least those represented there will be
those which contain endnotes in the first place.  Both methods are, I
think it's fair to say, a fair amount of labor, but they do give you
complete control over how your endnotes section looks, and over how
\texttt{noteref} back references within it look.  As with all new
functionality, \texttt{noteref} and \textsf{cmsendnotes} may well
contain bugs, so if you find any please let me know, but do please
also send along a minimum working example so I have a chance to
identify what's a bug in the code and what's resulted from inadequate
documentation.

\subsection{General Usage Hints}
\label{sec:hints}

\subsubsection{Loading the Style}
\label{sec:loading}

With the addition of the author-date styles to the package, I have
provided three keys for choosing which style to load, \texttt{notes},
\texttt{authordate}, and \textsf{authordate-trad}, one of which you
put in the options to the \cmd{usepackage} command.  The default way
of loading the notes + bibliography style has therefore slightly
changed.  With early versions of \textsf{bibla\-tex-chicago-notes},
the standard way of loading the package was via a call to
\textsf{biblatex}, e.g.:
\begin{verbatim}
  \usepackage[style=chicago-notes,strict,backend=bibtex8,%
    babel=other,bibencoding=inputenc]{biblatex}
\end{verbatim}
Now, the default way to load the style, and one that will in the
vast majority of standard cases produce the same results as the old
invocation, will look like this:
\begin{verbatim}
  \usepackage[notes,strict,backend=biber,autolang=other,%
    bibencoding=inputenc]{biblatex-chicago}
\end{verbatim}

(In point of fact, the previous \textsf{biblatex-chicago} loading
method without the \texttt{notes} option will still work, but only
because I've made the notes \&\ bibliography style the default if no
style is explicitly requested.)  If you read through
\textsf{biblatex-chicago.sty}, you'll see that it sets a number of
\textsf{biblatex} options aimed at following the Chicago
specification, as well as setting a few formatting variables intended
as reasonable defaults (see section~\ref{sec:presetopts}, above).
Some parts of this specification, however, are plainly more
\enquote{suggested} than \enquote{required,} and indeed many
publishers, while adopting the main skeleton of the Chicago style in
citations, nonetheless maintain their own house styles to which the
defaults I have provided do not conform.

\mylittlespace If you only need to change one or two parameters, this
can easily be done by putting different options in the call to
\textsf{biblatex-chicago} or redefining other formatting variables in
the preamble, thereby overriding the package defaults.  If, however,
you wish more substantially to alter the output of the package,
perhaps to use it as a base for constructing another style altogether,
then you may want to revert to the old style of invocation above.
You'll lose all the definitions in \textsf{biblatex-chicago.sty},
including those to which I've already alluded and also the code that
sets the note number in-line rather than superscript in endnotes or
footnotes, the URL line-breaking code, and the Chicago-specific
number- and date-range compression code.  You'll need to load the
required packages \textsf{xstring} and \textsf{nameref} yourself, as
\textsf{biblatex} doesn't do it for you.  Also, you'll lose the code
that calls \textsf{cms-american.lbx}, which means that you'll lose all
the Chicago-specific bibstrings I've defined unless you provide, in
your preamble, a \cmd{DeclareLanguageMapping} command adapted for your
setup, on which see section~\ref{sec:international} below and also
\S\S~4.9.1 and 4.11.8 in \textsf{biblatex.pdf}.

\mylittlespace What you \emph{will not} lose is the ability to call
the package options \texttt{annotation, strict, short,} and
\texttt{noibid} (section~\ref{sec:useropts}, above), in case these
continue to be useful to you when constructing your own modifications.
There's very little code, therefore, actually in
\textsf{biblatex-chicago.sty}, but I hope that even this minimal
separation will make the package somewhat more adaptable.  Any
suggestions on this score are, of course, welcome.

\subsubsection{Other Hints}
\label{sec:otherhints}

One useful rule, when you are having difficulty creating a .bib entry,
is to ask yourself whether all the information you are providing is
strictly necessary.  The Chicago specification is a very full one, but
the \emph{Manual} is actually, in many circumstances, fairly relaxed
about how much of the data from a work's title page you need to fit
into a reference.  Authors of introductions and afterwords, multiple
publishers in different countries, the real names of authors more
commonly known under pseudonyms, all of these are candidates for
exclusion if you aren't making specific reference to them, and if you
judge that their inclusion won't be of particular interest to your
readers.  Of course, any data that may be of such interest, and
especially any needed to identify and track down a reference, has to
be present, but sometimes it pays to step back and reevaluate how much
information you're providing.  I've tried to make
\textsf{biblatex-chicago-notes} robust enough to handle the most
complex, data-rich citations, but there may be instances where you can
save yourself some typing by keeping it simple.

\mylittlespace Scot Becker pointed out to me that the inverse problem
not only exists but may well become increasingly common, to wit, .bib
database entries generated by bibliographic managers which helpfully
provide as much information as is available, including fields that
users may well wish not to have printed (ISBN, URL, DOI,
\textsf{pagetotal}, inter alia).  The standard \textsf{biblatex}
styles contain a series of options, detailed in \textsf{biblatex.pdf}
\S~3.1.2.2, for controlling the printing of some of these fields, and
I have implemented the ones that are relevant to
\textsf{biblatex-chicago}, along with a couple that Scot requested and
that may be of more general usefulness.  There is also a general
option to excise with one command all the fields under consideration
-- please see section~\ref{sec:chicpreset} above.

\mylittlespace If you are having problems with the interaction of
punctuation and quotation marks in notes or bibliography, first please
check that you've used \cmd{mkbibquote} in the relevant part of your
.bib file.  If you are still getting errors, please let me know, as it
may well be a bug.

\mylittlespace For the \textsf{biblatex-chicago-notes} style, I have
fully adopted \textsf{biblatex's} system for providing punctuation at
the end of entries.  Several users noted insufficiencies in previous
releases of \textsf{biblatex-chicago}, sometimes related to the
semicolon between multiple citations, sometimes to ineradicable
periods after long notes, bugs that were byproducts of my attempt to
fix other end-of-entry errors.  One of the side effects of this older
code was (wrongly) to put a period after a long note produced, e.g.,
by a command like \verb+\footnote{\headlessfullcite}+, whereas only
the \enquote{foot} cite commands (including \cmd{autocite} in the
default \textsf{biblatex-chicago-notes} set up) should do so.  If you
came to rely on this side effect, please note now that you'll have to
put the period in yourself when explicitly calling \cmd{footnote},
like so: \verb+\footnote{\headlessfullcite{key}.}+

\mylittlespace When you use abbreviations at the ends of fields in
your .bib file (e.g., \enquote{\texttt{n.d.}} or
\enquote{\texttt{Inc.},}) \textsf{biblatex-chicago-notes} should deal
automatically with adding (or suppressing) appropriate punctuation
after the final dot.  This includes retaining periods after such dots
when a closing parenthesis intervenes, as in (n.d.).  Merely entering
the abbreviation without informing \textsf{biblatex} that the final
dot is a dot and not a period should always work, though you do have
to provide manual formatting in those rare cases when you need a comma
after the author's initials in a bibliography, usually in a
\textsf{misc} entry (see house:papers).  If you find you need to
provide such formatting elsewhere, please let me know.

\mylittlespace Finally, allow me to re-emphasize that, in its current
form, the notes \&\ bibliography style \emph{requires} the use of
\textsf{biber} as your backend --- variants of \textsc{Bib}\TeX\
simply cannot produce accurate output anymore, given how many features
now depend on the more modern backend.

\section{The Specification: Author-Date}
\label{sec:authdate}

The \textsf{biblatex-chicago} package contains two different
author-date styles.  The first,
\textsf{bibla\-tex-chicago-authordate}, implements the specifications
of the 17th edition of the \emph{Chicago Manual of Style}.  Numbers in
parentheses refer to sections of the \emph{Manual}, though many of
these references will in fact be to the chapter on the notes \&\
bibliography style (chapter 14), which chapter is, by design,
considerably more detailed than that devoted to the author-date style,
and which \textsf{biblatex-chicago-authordate} always modifies
according to the guidelines in chapter 15.  The second author-date
style, \textsf{biblatex-chicago-authordate-trad}, implements the same
specification but with a markedly different style of title
presentation, including sentence-style capitalization and the absence
of any quotation marks around the (plain-text) titles of
\textsf{article} or \textsf{incollection} entries, \emph{inter alia}.
The \textsf{trad} style is so named because older versions of the
\emph{Manual}, up to and including the 15th edition, recommended this
plainer style for author-date titles, and the 17th edition itself
suggests the possibility, when needed, of retaining such title
presentation in combination with its own recommendations for other
parts of the reference apparatus (15.38).  In practice, the
differences between the two styles necessitate separate discussions of
the \textsf{title} field and one extra package option
(\texttt{headline}), and that's about it.

\mylittlespace Generally, then, the following documentation covers
both Chicago author-date styles, and attempts to explain all the parts
of the specification that might be considered somehow \enquote{non
  standard,} at least with respect to the styles included with
\textsf{biblatex} itself.  In the section on entry fields I admit I
have also duplicated a lot of the information in
\textsf{biblatex.pdf}, which I hope won't badly annoy expert users of
the system.  As usual, headings in \mycolor{green}
\colmarginpar{\textsf{New in this release}} indicate either material
new to this release or old material that has undergone significant
revision.  The file \textsf{dates-test.bib} contains many examples
from the \emph{Manual} which, when processed using
\textsf{biblatex-chicago-authordate}, should produce the same output
as you see in the \emph{Manual} itself, or at least compliant output,
where the specifications are vague or open to interpretation, a state
of affairs which does sometimes occur.  If you are using
\textsf{biblatex-chicago-authordate-trad} the same basically holds,
but you'd have to keep one eye on the 15th edition of the
\emph{Manual} (chap.\ 17) for the titles.  I have provided
\href{file:cms-dates-sample.pdf}{\textsf{cms-dates-sample.pdf}} and
\href{file:cms-trad-sample.pdf}{\textsf{cms-trad-sample.pdf}}, which
show how my system processes \textsf{dates-test.bib}, and I have also
included the reference keys from the latter file below in parentheses.

\subsection{Entry Types}
\label{sec:types:authdate}

The complete list of entry types currently available in
\textsf{authordate} and \textsf{authordate-trad}, minus the odd
\textsf{biblatex} alias, is as follows: \textbf{article},
\textbf{artwork}, \textbf{audio}, \textbf{book}, \textbf{bookinbook},
\textbf{booklet}, \textbf{collection}, \textbf{customc},
\textbf{dataset}, \textbf{image}, \textbf{inbook},
\textbf{incollection}, \textbf{inproceedings}, \textbf{inreference},
\textbf{jurisdiction}, \textbf{legal}, \textbf{legislation},
\textbf{letter}, \textbf{manual}, \textbf{misc}, \textbf{music},
\textbf{mvbook}, \textbf{mvcollection}, \textbf{mvproceedings},
\textbf{mvreference}, \textbf{online} (with its alias \textbf{www}),
\textbf{patent}, \textbf{performance}, \textbf{periodical},
\textbf{proceedings}, \textbf{reference}, \textbf{report} (with its
alias \textbf{techreport}), \textbf{review}, \textbf{standard},
\textbf{suppbook}, \textbf{suppcollection}, \textbf{suppperiodical},
\textbf{thesis} (with its aliases \textbf{mastersthesis} and
\textbf{phdthesis}), \textbf{unpublished}, and \textbf{video}.

\mylittlespace What follows is an attempt to specify all the
differences between these types and the standard provided by
\textsf{biblatex}.  If an entry type isn't discussed here, then it is
safe to assume that it works as it does in the standard styles.  In
general, I have attempted not to discuss specific entry fields here,
unless such a field is crucial to the overall operation of a given
entry type.  As a general and important rule, most entry types require
very few fields when you use \textsf{biblatex-chicago-authordate}, so
it seemed to me better to gather information pertaining to fields in
the next section.

\paragraph*{\protect\mymarginpar{\textbf{article}}}
\label{sec:ad:article}
The \emph{Chicago Manual of Style} (14.164) recognizes three different
sorts of periodical publication, \enquote{journals,}
\enquote{magazines,} and \enquote{newspapers.}  The first (14.166) is
\enquote{a scholarly or professional periodical available mainly by
  subscription,} while the second refers to \enquote{weekly or monthly
  (or sometimes daily)} publications that are \enquote{available in
  individual issues at libraries or bookstores or newsstands or
  offered online, with or without a subscription.}
\enquote{Magazines} will tend to be \enquote{more accessible to
  general readers,} and typically won't have a volume number.  The
following paragraphs detail how to construct your .bib entries for all
these sorts of periodical publication.

\mylittlespace For articles in \enquote{journals} you can simply use
the traditional \textsc{Bib}\TeX\ --- and indeed \textsf{biblatex} ---
\textsf{article} entry type, which will work as expected and set off
the page numbers with a colon in the list of references, as required
by the \emph{Manual}.  If, however, you wish to cite a
\enquote{magazine} or a \enquote{newspaper}, then you need to add an
\textsf{entrysubtype} field containing the exact string
\texttt{magazine} or, now, its synonym \texttt{newspaper}.  The main
formatting differences between a \texttt{magazine/newspaper} and a
plain \textsf{article} are that time specifications (month, day,
season) aren't placed within parentheses, and that page numbers are
set off by a comma rather than a colon.  Otherwise, the two sorts of
reference have much in common.  (For \textsf{article}, see
\emph{Manual} 14.168--87, 15.9, 15.46--49; batson, beattie:crime,
chu:panda, connell:chronic, conway:evolution, friedman:learning,
garaud:gatine, garrett, hlatky:hrt, kern, lewis, loften:hamlet,
loomis:structure, rozner:liberation, schneider\hc mittelpleistozaene,
terborgh:preservation, wall:radio, warr:ellison,
white:callimachus. For \textsf{entrysubtype} \texttt{magazine}, cf.\
14.171, 14.188--200, 15.49; assocpress:gun, lakeforester:pushcarts,
morgenson:market, reaves:rosen, stenger:privacy.)

\mylittlespace The \emph{Manual} suggests that, no matter which
citation style you are using, it is \enquote{usually sufficient to
  cite newspaper and magazine articles entirely within the text}
(15.49).  This involves giving the title of the journal and the full
date of publication in a parenthetical reference, including any other
information in the main text (14.198), thereby obviating the need to
present such an entry in the list of references.  To utilize this
method in the author-date styles, in addition to a \texttt{magazine}
\textsf{entrysubtype}, you'll need to place \texttt{cmsdate=full} into
the \textsf{options} field, including \texttt{skipbib} there as well
to stop the entry printing in the list of references.  If the entry
only contains a \textsf{date} and \textsf{journaltitle} that's enough,
but if it's a fuller entry also containing an \textsf{author} then
you'll also need \texttt{useauthor=false} in the \textsf{options}
field.  Other surplus fields will be ignored.  (See osborne:poison.)

\mylittlespace If you are familiar with the notes \&\ bibliography
style, you'll know that the \emph{Manual} treats reviews (of books,
plays, performances, etc.) as a sort of recognizable subset of
\enquote{journals,} \enquote{magazines,} and \enquote{newspapers,}
distinguished mainly by the way one formats the title of the review
itself.  The key rule is this: if a review has a separate, non-generic
title (gibbard; osborne:poison) in addition to something that reads
like \enquote{review of \ldots,} then you need an \textsf{article}
entry, with or without the \texttt{magazine} \textsf{entrysubtype},
depending on the sort of publication containing the review.  If the
only title is the generic \enquote{review of \ldots,} for example,
then you'll need the \textsf{review} entry type, with or without this
same \textsf{entrysubtype} toggle using \texttt{magazine}.  On
\textsf{review} entries, see below.

\mylittlespace In the case of a review with a specific as well as a
generic title, the former goes in the \textsf{title} field, and the
latter in the \textsf{titleaddon} field.  Standard \textsf{biblatex}
intends this field for use with additions to titles that may need to
be formatted differently from the titles themselves, and
\textsf{biblatex-chicago-authordate} uses it in just this way, with
the additional wrinkle that it can, if needed, replace the
\textsf{title} entirely, and this in, effectively, any entry type,
providing a fairly powerful, if somewhat complicated, tool for getting
\textsf{biblatex} to do what you want.  Here, however, if all you need
is a generic title like \enquote{review of \ldots,} then you want to
switch to the \textsf{review} type, where you can simply use the
\textsf{title} field for it.

\mylittlespace \textsf{Biblatex-chicago} also, at the behest of
Bertold Schweitzer, supports the \textsf{relatedtype}
\texttt{reviewof}, which allows you to use the \textsf{related}
mechanism to provide information about the work being reviewed.  In
particular, it relieves you of the need to construct
\textsf{titleaddon} or \textsf{title} fields like: \texttt{review of
  \textbackslash mkbibemph\{Book Title\} by Author Name}, as the
\textsf{related} entry's \textsf{title} automatically provides the
\textsf{titleaddon} in the \textsf{article} type and the
\textsf{title} in the \textsf{review} type, with the \textsf{related}
mechanism providing the connecting string.  This may be particularly
helpful if you need to cite multiple reviews of the same work; please
see section \ref{sec:authrelated} for further details.

\mylittlespace No less than ten more things need explication under
this heading.  First, since the \emph{Manual} specifies that what goes
into the \textsf{titleaddon} field of \textsf{article} entries stays
unformatted --- no italics, no quotation marks --- this plain style is
the default for such text, which means that you'll have to format any
titles within \textsf{titleaddon} yourself, e.g., with
\cmd{mkbibemph\{\}}.  Second, the \emph{Manual} specifies a similar
plain style for the titles of other sorts of material found in
\enquote{magazines} and \enquote{newspapers,} e.g., obituaries,
letters to the editor, interviews, the names of regular columns, and
the like.  References may contain both the title of an individual
article and the name of the regular column, in which case the former
should go, as usual, in a \textsf{title} field, and the latter in
\textsf{titleaddon}.  As with reviews proper, if there is only the
generic title, then you want the \textsf{review} entry type.  (See
14.191, 14.195--96; morgenson:market, reaves:rosen.)

\mylittlespace Third, the \emph{Manual} suggests that
\enquote{unsigned newspaper articles or features are best dealt with
  in text \ldots} (14.199).  As with newspaper or magazine articles in
general, you can place \texttt{cmsdate=full} and \texttt{skipbib} into
the \textsf{options} field to produce an augmented in-text citation
whilst keeping this material out of the reference list.  If you do use
the reference list, then the standard shorter citation will be
sufficient, and in both cases the name of the periodical (in the
\textsf{journaltitle} field) will be used in place of the missing
author.  Just to clarify: in \textsf{article} or \textsf{review}
entries, \textsf{entrysubtype} \texttt{magazine}, a missing
\textsf{author} field results in the name of the periodical (in the
\textsf{journaltitle} field) being used as the missing author.
Without an \textsf{entrysubtype}, and assuming that no name whatsoever
can be found to put at the head of the entry, the \textsf{title} will
be used, not the \textsf{journaltitle}, or so I interpret the
\emph{Manual} (14.168).  The default sorting scheme in
\textsf{biblatex-chicago-authordate} considers the
\textsf{journaltitle} before the \textsf{title}, so if the latter
heads an entry you'll need a \textsf{sortkey}, just as you will if you
retain the definite or indefinite article at the beginning of the
\textsf{journaltitle} in author-less entries with an
\textsf{entrysubtype}.  If you want to abbreviate the
\textsf{journaltitle} for use in citations, but give the full name in
the list of references, then the \textsf{shortjournal} field is the
place for it.  A shortened \textsf{title} should go, as usual, in
\textsf{shorttitle}.  (See section~\ref{sec:authformopts}, below;
lakeforester:pushcarts, nyt:trevorobit, unsigned:ranke.)

\mylittlespace Fourth, Bertold Schweitzer has pointed out, following
the \emph{Manual} (14.183), that while an \textsf{issuetitle} often
has an \textsf{editor}, it is not too unusual for a \textsf{title} to
have, e.g., an \textsf{editor} and/or a \textsf{translator}.  In order
to allow as many permutations as possible on this theme, I have
brought the \textsf{article} entry type into line with most of the
other types in allowing the use of the \textsf{namea} and
\textsf{nameb} fields in order to associate an editor or a translator
specifically with the \textsf{title}.  The \textsf{editor} and
\textsf{translator} fields, in strict homology with other entry types,
are associated with the \textsf{issuetitle} if one is present, and
with the \textsf{title} otherwise.  The usual string concatenation
rules still apply --- cf.\ \textsf{editor} and \textsf{editortype} in
section~\ref{sec:fields:authdate}, below.

\mylittlespace Fifth, in certain fields, just beginning your data with
a lowercase letter activates the mechanism for capitalizing that
letter depending on its context within a reference list entry.  This
is less important in the author-date styles, where this information
only turns up in the reference list and not in citations, but you can
consult\,\textbf{\textbackslash autocap} in
section~\ref{sec:formatting:authdate} below for all the details.  Both
the \textsf{titleaddon} and \textsf{note} fields are among those
treating their data this way, and since both appear regularly in
\textsf{article} entries, I thought the problem merited a preliminary
mention here.

\mylittlespace Sixth, if you need to cite an entire issue of any sort
of periodical, rather than one article in an issue, then the
\textsf{periodical} entry type, once again with or without the
\texttt{magazine} toggle in \textsf{entrysubtype,} is what you'll
need.  (You can also use the \textsf{article} type, placing what would
normally be the \textsf{issuetitle} in the \textsf{title} field and
retaining the usual \textsf{journaltitle} field, but this arrangement
isn't compatible with standard \textsf{biblatex}.)  The \textsf{note}
field is where you place something like \enquote{special issue} (with
the small \enquote{s} enabling the automatic capitalization routines),
whether you are citing one article or the whole issue
(conley:fifthgrade, good:wholeissue).  Indeed, this is a somewhat
specialized use of \textsf{note}, and if you have other sorts of
information you need to include in an \textsf{article} or
\textsf{periodical} entry, then you shouldn't put it in the
\textsf{note} field, but rather in \textsf{titleaddon} or perhaps
\textsf{addendum} (brown:bremer).

\mylittlespace Seventh, I would suggest that if you wish to cite
certain kinds of television or radio broadcast, most notably
interviews but perhaps also news segments or other
\enquote{journalistic} material, then the \textsf{article} type,
\textsf{entrysubtype} \texttt{magazine} is the place for it.  The name
of the program as a whole would go in \textsf{journaltitle}, with the
name of the episode in \textsf{title}.  The network's name goes into
the \textsf{usera} field.  Of course, if the piece you are citing has
only a generic name (an interview, for example), then the
\textsf{review} type would be the best place for it (8.189, 14.213;
see bundy:macneil for an example of how this all might look in a .bib
file.)  Other sorts of broadcast, usually accessible through
commercial recordings, would need one of the audiovisual entry types,
probably \textsf{audio} (danforth:podcast) or \textsf{video}
(friends:leia), while recordings from archives fit best either into
\textsf{online} or into \textsf{misc} entries with an
\textsf{entrysubtype} (coolidge:speech, roosevelt:speech).

\mylittlespace Eighth, the \emph{Manual} (14.208, 15.51) specifies
that blogs and other, similar online material should be presented like
\textsf{articles}, with \texttt{magazine} \textsf{entrysubtype}
(ellis:blog), and needn't appear in a reference list at all, if you'd
prefer to provide relevant details in the text.  I've attempted,
however, to make \textsf{biblatex-chicago-authordate} as useful as
possible when managing references to such sources, so I'll outline
these facilities here.  The title of the specific entry goes in
\textsf{title}, the general title of the blog goes in
\textsf{journaltitle}, and the word \enquote{\texttt{blog}} in the
\textsf{location} field (though you could just use special formatting
in the \textsf{journaltitle} field itself, which may sometimes be
necessary).  The 17th edition specifies that \enquote{blogs that are
  part of a larger publication should include the name of that
  publication.}  This usually involves a newspaper or magazine which
also publishes various blogs on its website, and it means that such
entries need a more general title than the \textsf{journaltitle}.
It's not standard \textsf{biblatex} or anything, but you can now put
such information in \textsf{maintitle} (with \textsf{mainsubtitle} and
\textsf{maintitleaddon}, if needed), but only in \textsf{article} and
\textsf{review} entries with a \texttt{magazine} \textsf{entrysubtype}
(amlen:hoot).

\mylittlespace The \emph{Manual} (15.51) is even more emphatic about
whole blogs (rather than individual posts) and comments on blogs not
appearing in reference lists, but I've kept as many options open as
possible, including fairly simple ways you can provide all the
information needed in text citations alone.  To cite a whole blog,
you'll need the \textsf{periodical} entry type, with a \textsf{title}
instead of a \textsf{journaltitle}, along with a (possible)
\textsf{maintitle} (amlen:wordplay).  Comments on blogs, with generic
titles like \enquote{comment on} or \enquote{reply to,} need a
\textsf{review} entry with the same \textsf{entrysubtype} (viv:amlen).
Such comments make particular use of the \textsf{eventdate} and
\textsf{nameaddon} fields, and also of specialized \textbf{customc}
entries for adding comments to in-text citations.  Please see the
documentation of \textbf{customc}, \textbf{periodical}, and
\textbf{review}, the \textsf{relatedtype} \texttt{commenton} in
section~\ref{sec:authrelated}, and the general discussion of online
sources in the \textbf{online} documentation.

\mylittlespace Ninth, the special \textsf{biblatex} field
\textsf{shortjournal} allows you to present shortened
\textsf{journaltitles} in \textsf{article}, \textsf{review}, and
\textsf{periodical} entries, as well as facilitating the creation of
lists of journal abbreviations in the manner of a \textsf{shorthand}
list.  Please see the documentation of \textbf{shortjournal} in
section~\ref{sec:fields:authdate} for all the details on how this
works.

\mylittlespace Finally, the 17th edition (14.191) specifies that, for
news sites carrying \enquote{stories as they unfold, it may be
  appropriate to include a time stamp for an article that includes
  one.}  You can provide this by using the standard \textsf{biblatex}
time stamp format inside the \textsf{date} field, e.g.,
\texttt{2008-07-01T10:18:00}.  Since the \emph{Manual} prefers the
standard time zone initialisms, a separate \textsf{timezone} field
would be required if you want to provide one.

\mylittlespace If you're still with me, allow me to recommend that you
browse through \textsf{dates-test.bib} to get a feel for just how many
of the \emph{Manual}'s complexities the \textsf{article},
\textsf{periodical}, and \textsf{review} types attempt to address.  It
may be that in future releases of \textsf{biblatex-chicago} I'll be
able to simplify these procedures somewhat, but with any luck the vast
majority of sources won't require knowledge of these onerous details.

\mybigspace Arne \mymarginpar{\textbf{artwork}} Kjell Vikhagen
pointed out to me that none of the standard entry types were
straightforwardly adaptable when referring to visual artworks.  It's
unclear whether the \emph{Manual} (14.235) believes it necessary to
include them in the reference apparatus at all, but it's easy to
conceive of contexts in which a list of artworks studied might be
desirable, and \textsf{biblatex} includes entry types for just this
purpose, though the standard styles leave them undefined.
\textsf{Biblatex-chicago} defines both \textsf{artwork} and
\textsf{image}, which are in fact now clones of each other, so you can
use either of them indifferently, the distinction existing only for
historical reasons.

\mylittlespace As one might expect, the artist goes in \textsf{author}
and the name of the work in \textsf{title}.  The \textsf{type} field
is intended for the medium --- e.g., oil on canvas, charcoal on paper
--- and the \textsf{version} field might contain the state of an
etching.  You can place the dimensions of the work in \textsf{note},
and the current location in \textsf{organization},
\textsf{institution}, and/or \textsf{location}, in ascending order of
generality.  The \textsf{type} field, as in several other entry types,
uses \textsf{biblatex's} automatic capitalization routines, so if the
first word only needs a capital letter at the beginning of a sentence,
use lowercase in the .bib file and let \textsf{biblatex} handle it for
you.  (See \emph{Manual} 3.22, 8.198; leo:madonna, bedford:photo.)

\mylittlespace The 17th edition of the \emph{Manual} has included new
information in some of its examples, so I have added 4 new fields to
the driver.  Alongside the usual \textsf{date} for the creation of a
work, you may also want to include the printing date of a particular
exemplar of a photograph or a print.  The system I have designed uses
the \emph{earlier} of the \textsf{date} and the \textsf{origdate} to
be the date of creation, and the \emph{later} to be the printing date.
The style will automatically prefix the printing date with the
localized \cmd{bibstring} \texttt{printed}, so if that's the wrong
string entirely then you can define \textsf{userd} any way you like to
change it.  If only \emph{one} of those two dates is available, it
will always serve as a creation date.  Any date specification provided
will always appear in full somewhere in the reference list entry,
though sometimes that could be the plain year at the head of the
entry.  This system, which is unlike other entry types, helps to avoid
ambiguities in some situations.

\mylittlespace One of the \emph{Manual's} examples is of a photograph
published in a periodical, and information about this publication
appears late in the entry, after the \textsf{type}.  I have included
the \textsf{howpublished} field so that you can give information about
the periodical (meaning that you'll have to format the title yourself
with \cmd{mkbibemph}), and the \textsf{eventdate} field for you to
provide the date of publication (mccurry:afghangirl).

\mylittlespace Depending on the presence or absence of these three
date fields, and also on how you've set the \texttt{cmsdate} option,
any of the three can appear in citations and at the head of reference
list entries, allowing you to order entries by creation date, printing
date, or publication date.  See the documentation of \textbf{date} on
page~\pageref{sec:ad:date}, below, for all the complicated details.
Please note, when choosing your date presentation, that these new
fields ostensibly replace most of the possible uses of the
\textsf{pubstate} field in \textsf{artwork} entries, though this field
continues to function here more or less as before, should you still
require it.

\mylittlespace As a final complication, the \emph{Manual} (8.198) says
that \enquote{the names of works of antiquity \ldots\,are usually set
  in roman.}  If you should need to include such a work in the
reference apparatus, you can either define an \textsf{entrysubtype}
for an \textsf{artwork} entry --- anything will do --- or you could
try the \textsf{misc} entry type with an \textsf{entrysubtype}.
Assuming the complicated date handling I've just outlined isn't
required for such a work, in this instance the other fields in a
\textsf{misc} entry function pretty much as in \textsf{artwork}.

\mybigspace Following \mymarginpar{\textbf{audio}} the request of
Johan Nordstrom, I have included three entry types, all undefined by
the standard styles, designed to allow users to present audiovisual
sources in accordance with the Chicago specifications.  The
\emph{Manual's} presentation of such sources (14.261--68, 15.57),
though admirably brief, seems to me somewhat inconsistent, though
perhaps I'm merely unable to spot the important regularities.  The
proliferation of online sources has made the task yet more complex,
requiring the inclusion of the \textbf{article}, the \textbf{online},
and even the \textbf{misc} entry types, which see, under the
audiovisual rubric.  I shall attempt to delineate the main differences
here, and though there are likely to be occasions when your choice of
entry type is not obvious, at the very least \textsf{biblatex-chicago}
should help you maintain consistency.

\mylittlespace For users of the author-date styles, the 17th edition
of the \emph{Manual} continues to emphasize the need to provide dating
information for audiovisual materials (14.263), meaning that nearly
all such entries will have some such information and will therefore
fit better, stylistically, with other author-date references.  In
particular, it continues to recommend that \enquote{the date of the
  original recording should be privileged in the citation} (15.57).
Guidance for supplying dates for this class of material will be found
below under the different entry types in use, though it will also be
worthwhile to look at the documentation of \textsf{date},
\textsf{eventdate}, \textsf{origdate}, and \textsf{urldate}, in
section~\ref{sec:fields:authdate}, below.  The \emph{Manual} continues
to suggest, also, that \enquote{it is often more appropriate to list
  such materials in running text and group them in a separate section
  or discography.}

\mylittlespace The \textbf{music} type is intended for all musical
recordings that do not have a video component.  This means, for
example, digital media (whether on CD or hard drive), vinyl records,
and tapes.  The \textbf{video} type includes most visual media,
whether it be films, TV shows, tapes and DVDs of the preceding or of
any sort of performance (including music), or online multimedia.  The
\emph{Manual's} treatment (14.267) of the latter suggests that online
video excerpts, short pieces, and interviews should generally use the
\textbf{online} type (horowitz:youtube, pollan:plant) or the
\textbf{article} type (harwood:biden, kessler:nyt), depending on
whether the pieces originate from an identifiably
\enquote{journalistic} outlet.  The \textbf{audio} type, our current
concern, fills gaps in the others, and presents its sources in a more
\enquote{book-like} manner.  Published musical scores need this type
--- unpublished ones would use \textsf{misc} with an
\textsf{entrysubtype} (shapey:partita) --- as do podcasts and such
favorite educational formats as the slideshow and the filmstrip
(danforth:podcast, greek:filmstrip, schubert:muellerin,
verdi:corsaro).  The \emph{Manual} (14.264) sometimes uses a similar
format for audio books (twain:audio), though, depending on the sorts
of publication facts you wish to present, this sort of material may
fall under \textsf{music} (auden:reading).  Dated audio recordings
that are part of an archive, online or no, may be presented either in
an \textbf{online} or in a \textbf{misc} entry with an
\textsf{entrysubtype}, the difference mainly being in just how closely
associated the \textsf{date} will be with the \textsf{title}
(coolidge:speech, roosevelt:speech).  Actual radio broadcasts (as
opposed to podcasts) pose something of a conundrum.  Interviews and
other sorts of \enquote{journalistic} material fit well into
\textsf{article} or \textsf{review} entries (14.213), but other sorts
of broadcast are not well represented in the \emph{Manual's} examples
(8.189), and what little there is suggests that, counter-intuitively,
the \textsf{video} type is the best fit, as it is well equipped to
present broadcasts of any sort.

\mylittlespace Once you've accepted the analogy of composer to
\textsf{author}, constructing an \textsf{audio} entry should be fairly
straightforward, since many of the fields function just as they do in
\textsf{book} or \textsf{inbook} entries.  Indeed, please note that I
compare it to both these other types as, in common with the other
audiovisual types, \textsf{audio} has to do double duty as an analogue
for both books and collections, so while there will normally be an
\textsf{author}, a \textsf{title}, a \textsf{publisher}, a
\textsf{date}, and a \textsf{location}, there may also be a
\textsf{booktitle} and/or a \textsf{maintitle} --- see
schubert:muellerin for an entry that uses all three in citing one song
from a cycle.  (As with the \textsf{music} and \textsf{video} types,
you can cite an individual piece separate from any large collection by
using the \textsf{title} field and by defining an
\textsf{entrysubtype}, which will stop \textsf{biblatex-chicago}
italicizing your \textsf{title} in the absence of a
\textsf{booktitle}.)  If the medium in question needs specifying, the
\textsf{type} field is the place for it.  Please note, also, that
while the \textsf{titleaddon} field can still specify creative or
editorial functions for which \textsf{biblatex-chicago} provides no
automated, localized handling, you can also now provide the string you
need in an \textsf{editor[abc]type} field, e.g.,
\enquote{\texttt{libretto by}} (verdi:corsaro).

\mylittlespace For podcasts, newly covered by the 17th edition
(14.267), the \textsf{audio} type provides the nearest analogue I
could find, and in general most of the data should fit comfortably
into the fields already discussed above, the episode name in
\textsf{title} and the name of the podcast in \textsf{booktitle}, for
starters.  Two details, however, need mentioning: first, the
\textsf{note} field as the place to specify that it is a podcast, and
the \textsf{eventdate} field for the date of publication of the
specific episode (\textsf{title}) cited, which appears in close
association with that \textsf{title}.  Indeed, the \textsf{eventdate}
field helps \textsf{biblatex-chicago} know that the \textsf{audio}
entry is a podcast episode, and helps it construct the entry
appropriately (danforth:podcast).

% %\enlargethispage{\baselineskip}

\mybigspace This \mymarginpar{\textbf{book}} is the standard
\textsf{biblatex} and \textsc{Bib}\TeX\ entry type, but the package
can automatically provide abbreviated references in the reference list
when you use a \textsf{crossref} or an \textsf{xref} field.  The
functionality is not enabled by default, but you can enable it in the
preamble or in the \textsf{options} field using the
\texttt{booklongxref} option.  Please see \textbf{crossref} in
section~\ref{sec:fields:authdate} and \texttt{booklongxref} in
section~\ref{sec:authpreset}, below.  Also, cf.\ harley:ancient:cart,
harley:cartography, and harley:hoc for how this might look.  The
\textsf{book} type is useful also to present multimedia app content,
the added fields \textsf{version} and \textsf{type} providing
information about the app's version and about the system on which it
runs (14.268, 15.57; angry:birds).

\mybigspace This \mymarginpar{\textbf{bookinbook}} type provides the
means of referring to parts of books that are considered, in other
contexts, themselves to be books, rather than chapters, essays, or
articles.  Such an entry can have a \textsf{title} and a
\textsf{maintitle}, but it can also contain a \textsf{booktitle}, all
three of which will be italicized in the reference matter.  In general
usage it is, therefore, rather like the traditional \textsf{inbook}
type, only with its \textsf{title} in italics rather than in quotation
marks.  As with the \textsf{book} type, you can enable automatically
abbreviated references in the reference list, though this isn't the
default.  Please see \textbf{crossref} in
section~\ref{sec:fields:authdate} and \texttt{booklongxref} in
section~\ref{sec:preset:authdate}, below.  (Cf.\ \emph{Manual} 14.109,
14.122, 14.124; bernhard:boris, bernhard:ritter, and
bernhard:themacher for the abbreviating functionality; also
euripides:orestes [treated differently in 14.122 and 14.124],
plato:republic:gr.)

\mybigspace This \mymarginpar{\textbf{booklet}} is the first of two
entry types --- the other being \textsf{manual}, on which see below
--- which are traditional in \textsc{Bib}\TeX\ styles, but which the
\emph{Manual} (14.220) suggests may well be treated basically as
books.  In the interests of backward compatibility,
\textsf{biblatex-chica\-go-authordate} will so format such an entry,
which uses the \textsf{howpublished} field instead of a standard
\textsf{publisher}, though of course if you do decide just to use a
\textsf{book} entry then any information you might have given in a
\textsf{howpublished} field should instead go in \textsf{publisher}.
(See clark:mesopot.)

\mybigspace This \mymarginpar{\textbf{collection}} is the standard
\textsf{biblatex} entry type, but the package can provide
automatically abbreviated references in the reference list when you
use a \textsf{crossref} or an \textsf{xref} field.  The functionality
is not enabled by default, but you can enable it in the preamble or in
the \textsf{options} field using the \texttt{booklongxref} option.
Please see \textbf{crossref} in section~\ref{sec:fields:authdate} and
\texttt{booklongxref} in section~\ref{sec:authpreset}, below.  See
harley:ancient:cart, harley:cartography, and harley\hc hoc for how
this might look.

\mybigspace This \mymarginpar{\textbf{customc}} entry type allows you
to include alphabetized cross-references to other, separate entries in
the bibliography, particularly to other names or pseudonyms, as
recommended by the \emph{Manual}.  (This is different from the usual
\textsf{crossref}, \textsf{xref}, \textsf{userf}, and \textsf{related}
mechanisms, all primarily designed to include cross-references to
other works.  Cf.\ 14.81--82, 15.35).  I should add immediately that,
as I read the specification, the alphabetized cross-references
provided by \textsf{customc} are particularly encouraged, bordering on
required, when a reference list \enquote{includes two or more works
  published by the same author but under different pseudonyms.}  The
following entries in \textsf{dates-test.bib} show one way of
addressing this: creasey:ashe:blast, creasey:york:death,
creasey:morton:hide, ashe:creasey, york\hc creasey and morton:creasey.
In these latter cases, you would need merely to place the pseudonym in
the \textsf{author} field, and the author's real name, under which his
or her works are presented in the bibliography, in the \textsf{title}
field.  To make sure the cross-reference also appears in the
bibliography, you can either manually include the entry key in a
\cmd{nocite} command, or you can put that entry key in the
\textbf{userc} field in the work's main .bib entry, in which case
\textsf{biblatex-chicago} will print the cross-reference if and only
if you cite the main work.  (Cf.\ \textsf{userc}, below.)

\mylittlespace Under ordinary circumstances, \textsf{biblatex-chicago}
will connect the two parts of the cross-reference with the word
\enquote{\emph{See}} --- or its equivalent in the document's language
--- in italics.  If you wish to present it differently, you can put
the connecting word(s) into the \textsf{nameaddon} field, formatted as
you wish.

\mylittlespace Finally, you may need to use this entry type if you
wish to include a comment inside the parentheses of a citation, as
specified by the \emph{Manual} (15.24).  If you have a
\textsf{postnote}, then you can manually provide the punctuation and
comment there, e.g., \cmd{autocite[4;~the unrevised
  trans.]\{stendhal:parma\}}.  Without a \textsf{postnote}, you have
two solutions.  You can enable the \texttt{postnotepunct} option,
which allows you simply to type \cmd{autocite[;~the unrevised
  trans.]\{stendhal:parma\}}, or you can use a separate
\textsf{customc} entry containing just the text of the comment in the
\textsf{title} field, \textsf{entrysubtype} \texttt{classical}, and
\textsf{options} \texttt{skipbib}.  An \cmd{autocites} command calling
both the main text and the comment then does the trick, e.g.,
\cmd{autocites\{chicago:manual\}\{chicago:comment\}}.  Cf.\
\texttt{postnotepunct} in section~\ref{sec:authuseropts}, below.

\mylittlespace For its 17th edition, the \emph{Manual} has provided a
more detailed treatment of online comments, whether on blogs, mailing
lists, or social media posts (15.50--52).  Such comments \enquote{are
  cited only in the text, in reference to the related post,} an
arrangement most easily created using \textsf{customc} entries
referencing the main post.  The new \texttt{commenton}
\texttt{relatedtype} in \textsf{online} and \textsf{review} entries
attempts to automate this for you, creating a separate
\textsf{customc} entry to be used in an \cmd{autocites} command
together with the comment's own entry.  Please see the details of this
in the \textbf{online} and \textbf{review} entry types, below, and in
the \texttt{commenton} docs in section~\ref{sec:authrelated}.  Cf.\
also ac:comment, diaz:surprise, ellis:blog, licis:diazcomment, and the
associated automatically-created entries ellis:blog-customc and
diaz:surprise-customc.

\mybigspace This \mymarginpar{\textbf{dataset}} entry type, new in
\textsf{biblatex} 3.13, allows you to cite scientific databases, for
which the \emph{Manual} (14.257) presents some rather specific, if
brief, instructions.  To construct your entry, you can put the name of
the database into \textsf{author}, a \enquote{descriptive phrase or
  record locator} in the \textsf{title} field, and if there's a
specific accession number needed beyond the record locator you can put
it into the \textsf{number} field, with the \textsf{type} field
reserved to help explain what sort of \textsf{number} is involved.
The \textsf{howpublished} field can also be used to provide extra
descriptive detail about the \textsf{number}, if needed.  More
generally, a \textsf{url} will locate the database as a whole and a
\textsf{urldate} will specify the date you accessed it.  If, for some
reason, an additional date is relevant, then the \textsf{date} field
is available, while the \textsf{pubstate} field will appear before the
\textsf{date} in case you need to modify the latter.  (See 14.257;
genbank:db, nasa:db.)  Given that usually the only relevant date in
such entries is the access date, which means that they would usually
have \enquote{n.d.}\ at the head of reference-list entries and in
citations, I have thought it sensible to treat \textsf{dataset}
entries by default as author-title instead of author-date in
citations, and to stop the printing of \enquote{n.d.}\ in reference
lists just as in \textsf{misc} entries.  You can shorten the component
parts of the author-title citation using the usual
\textsf{shortauthor}, \textsf{shorthand}, and/or \textsf{shorttitle}
fields, and you can also restore author-date formatting to these
entries by setting the \texttt{authortitle} option to \texttt{false}
either in individual entries or in the preamble for all examples of
the entry type.

%%\enlargethispage{\baselineskip}

\mybigspace This \mymarginpar{\textbf{image}} entry type is now a
clone of the \textsf{artwork} type, which see.  I retain it here for
historical reasons (See 3.22, 8.198; bedford:photo.)

\mybigspace These \mymarginpar{\textbf{inbook}\\\textbf{incollection}}
two standard \textsf{biblatex} types have very nearly identical
formatting requirements as far as the Chicago specification is
concerned, but I have retained both of them for compatibility.
\textsf{Biblatex.pdf} (\S~2.1.1) intends the first for \enquote{a part
  of a book which forms a self-contained unit with its own title,}
while the second would hold \enquote{a contribution to a collection
  which forms a self-contained unit with a distinct author and its own
  title.}  The \textsf{title} of both sorts will be placed within
quotation marks, and in general you can use either type for most
material falling into these categories.  I have, in both types,
implemented the \emph{Manual's} recommendations for space-saving
abbreviations in the reference list when you cite multiple pieces from
the same \textsf{collection}.  These abbreviations are activated by
default when you use the \textsf{crossref} or \textsf{xref} field in
\textsf{incollection} entries \emph{and} in \textsf{inbook} entries,
because although the \emph{Manual} (15.42) here specifies a
\enquote{multiauthor book,} I believe the distinction between the two
is fine enough, and the discussion general enough, to encourage
similar treatments.  (For more on this mechanism see \textbf{crossref}
in section~\ref{sec:fields:authdate}, below, and the option
\texttt{longcrossref} in section~\ref{sec:authpreset}.  Please note
that it is also active by default in \textsf{letter} and
\textsf{inproceedings} entries.)  If the part of a book to which you
are referring has had a separate publishing history as a book in its
own right, then you may wish to use the \textsf{bookinbook} type,
instead, on which see above.  (See \emph{Manual} 14.106--9, 15.42;
\textsf{inbook}: ashbrook:brain, phibbs:diary, will:cohere;
\textsf{incollection}: centinel:letters, contrib:contrib,
sirosh:visualcortex; ellet:galena, keating:dearborn, and
lippincott:chicago [and the \textsf{collection} entry prairie:state]
demonstrate the use of the \textsf{crossref} field with its attendant
abbreviations in the list of references.)

\mylittlespace \textbf{NB}: The \emph{Manual} suggests that, when
referring to a chapter, one use either a chapter number or the
inclusive page numbers, not both.  In-text citations, of course,
require any \textsf{postnote} field to specify if it is a whole
chapter to which you are referring.

\mybigspace This \mymarginpar{\textbf{inproceedings}} entry type works
pretty much as in standard \textsf{biblatex}, even more so now that,
after a request from Patrick Danilevici, I have newly included the
\textsf{eventdate}, \textsf{eventtitle}, \textsf{eventtitleaddon}, and
\textsf{venue} fields for specifying where and when the event occurred
that produced the proceedings.  These four fields are the main
difference between it and \textsf{incollection}, along with the lack
of an \textsf{edition} field and the possibility that an
\textsf{organization} may be cited alongside the \textsf{publisher},
even though the \emph{Manual} doesn't specify the use of any of these
supplementary fields (14.217).  Please note, also, that the
\textsf{crossref} and \textsf{xref} mechanism for shortening citations
of multiple pieces from the same \textsf{proceedings} is operative
here, just as it is in \textsf{incollection} and \textsf{inbook}
entries.  See \textbf{crossref} in section~\ref{sec:fields:authdate}
and the option \texttt{longcrossref} in section~\ref{sec:authpreset},
below, for more details.

\paragraph*{\protect\mymarginpar{\textbf{inreference}}}
\label{sec:ad:inreference}
This entry type is aliased to \textsf{incollection} in the standard
styles, but the \emph{Manual's} requirements prompted a thoroughgoing
revision.  Instructions for the author-date style are not copious, so
what follows is my best guess at following the specification
(14.232--34).

\mylittlespace Please first note that if your reference work can
easily or conveniently be presented like a regular author-date book,
that is, with an author or editor, a year of publication, and a title,
and if you will be citing it by page or section number, then you
should almost certainly simply choose the \textsf{book} entry type for
your .bib entry. (Cf.\ mla:style, schellinger:novel, times:guide.  The
latter was presented as an \textsf{inreference} entry for the notes
\&\ bibliography style, but because the \textsf{book} entry type can
also present references to alphabetized headings [see below], at least
in the list of references, then it seemed better just to choose a
\textsf{book} entry for the author-date styles.)

\mylittlespace If your source simply doesn't fit the standard
author-date template for a \textsf{book}, most especially if it is a
\enquote{well-known} reference work, then you may need to use the
\textsf{inreference} type, the main feature of which is the
\textsf{lista} field, which you use to present citations from
\enquote{alphabetically arranged} works by named article rather than
by page number.  You should present these article names just as they
appear in the work, separated by the keyword \enquote{\texttt{and}} if
there is more than one, and \textsf{biblatex-chicago-authordate} will
provide the appropriate prefatory string (\texttt{s.v.}, plural
\texttt{s.vv.}), and enclose each in its own set of quotation marks
(times:guide).  More relevant to the author-date styles is the fact
that the \textsf{postnote} field works the same way in
\textsf{inreference} entries, the only limitation on this system being
that this field, unlike \textsf{lista}, is not a list, and therefore
for the formatting to work correctly you can only put one article name
in it.  In the case of \enquote{[w]ell-known reference books, such as
  major dictionaries and encyclopedias,} you are encouraged not to
include them in the list of references, so the \textsf{lista} field
actually may be of less use than this special formatting of
\textsf{postnote}.  You may want to look at ency:britannica, where
only a (carefully-formatted) \textsf{shorttitle}, a \textsf{year}, and
an \textsf{options} field are necessary to allow you to produce
in-text citations that look like (\emph{Ency.\ Brit.}, 15th ed. 1985,
s.v. \enquote{Article}).

\mylittlespace If it seems appropriate to include such a work in the
list of references, perhaps because the work is not so well known that
a short citation will be parseable by your readers, or perhaps because
it is an online work, which requires you to provide a \textsf{urldate}
(see below), be aware that the contents of the \textsf{lista} field
will also be presented there, which may not be what you want.  A
separate \textsf{inreference} or \textsf{reference} entry might well
solve this problem.  In a typical \textsf{inreference} entry very few
fields are needed, but \enquote{if a physical edition is cited, not
  only the edition number (if not the first) but also the date the
  volume or set was issued must be specified.}  In practice, this
means a \textsf{title}, \textsf{date}, and possibly an
\textsf{edition} field.  The \textsf{author} field holds the author of
the specific article (in \textsf{lista}), not the author of the
\textsf{title} as a whole.  This name will be printed in parentheses
after the alphabetized entry's title (grove:sibelius).

\mylittlespace Finally, all of these rules apply to online reference
works, along with a few more.  The 17th edition of the \emph{Manual}
now allows, \enquote{subject to editorial discretion,} the alternative
treatment of an online reference work which \enquote{does not have
  (and never had) a printed counterpart} (14.206, 14.233).  In effect
this means that it can be treated more like an \textsf{online} entry
than a \textsf{book}, its \textsf{title} therefore in plain roman
rather than in italics.  You can achieve this in \textsf{inreference}
entries by providing an \textsf{entrysubtype} in the entry.  Online
reference works need not only a \textsf{url} but also, always, a
\textsf{urldate} (instead of a \textsf{date)}, as these sources are in
constant flux.  When that flux is of a particularly high frequency, as
with Wikipedia, then a time stamp may also be needed.  You can provide
this in the \textsf{urldate} field itself, using the standard
\textsf{biblatex} format, e.g., \texttt{2008-07-01T10:18:00}.  It is
\mymarginpar{\texttt{urlstamp=true}} possible to turn off the
printing of the \textsf{urltime} with the new \texttt{urlstamp}
option, which is set to \texttt{true} by default, but which can be
changed in your preamble for the whole document, for specific entry
types, or in the \textsf{options} field of specific entries
(wikiped:bibtex, grove:sibelius).  In keeping with the rules of the
17th edition, changed since the 16th, any \textsf{inreference} entry
with no date, or with only an access date, as opposed to a revision
date or another, more traditional publishing date, will use the
\texttt{nodate} abbreviation at the head of the entry and in citations
(15.44, 15.50).

\mybigspace I
\mymarginpar{\textbf{jurisdiction}\\\textbf{legal}\\\textbf{legislation}}
document these three types in section~\ref{sec:legal} below, both
because they all follow the specifications of the \emph{Bluebook}
instead of the \emph{Manual}, and also because they are the only entry
types treated identically by the notes \&\ bibliography style and the
author-date styles.

\mybigspace This \mymarginpar{\textbf{letter}} entry type was designed
to be used for citing letters, memoranda, or similar texts, but
\emph{only} when they appear in a published collection.  (Unpublished
material of this nature needs a \textsf{misc} entry, for which see
below.)  The author-date specification (15.43), however, recommends
against individual letters appearing in a list of references,
suggesting instead that you put the whole published collection in a
\textsf{book} entry and use a notice in the text to specify the letter
(white:total).

\mylittlespace If you absolutely must include individual letters in
the list of references, for whatever reason, then please consult the
instructions above for the notes \&\ bibliography style in
section~\ref{sec:entrytypes}, s.v.\ \enquote{\textsf{letter,}} which
cover all the details.  There are a few wrinkles, related to date
specifications, that I shall attempt to clarify here.  If you look at
white:ross:memo and white:russ, you'll see two letters from the same
published collection, both written in the same year.  You can simply
use the \textsf{origdate} field in both of them, because in the
absence of a \textsf{date} (or an \textsf{eventdate}) \textsf{Biber}
and \textsf{biblatex} will use the \textsf{origyear} as the
\textsf{labelyear}, putting it at the head of the entry and in the
citation, and also ensuring that the letters \texttt{a,b,c} are
appended to disambiguate the two sources.  In this case, it works
because we are using the \textsf{xref} mechanism to refer to the whole
published collection (white:total), so a separate citation of that
entry provides the \textsf{date} for the shortened cross-reference
included in the list of references, and the \textsf{letter} entry
never sees that \textsf{date} at all.  (Cf.\ also the documentation of
the \textsf{date} field in section~\ref{sec:fields:authdate} below.)

\mybigspace This \mymarginpar{\textbf{manual}} is the second of two
traditional \textsc{Bib}\TeX\ entry types that the \emph{Manual}
suggests formatting as books, the other being \textsf{booklet}. As
with this latter, I have retained it in
\textsf{biblatex-chicago-authordate} for backward compatibility, its
main peculiarity being that, in the absence of a named author, the
\textsf{organization} producing the manual will be provided both as
author and as publisher.  (You can give a shortened form of the
\textsf{organization} in the \textsf{shortauthor} field for text
citations, if needed, or use the \textsf{shorthand} field.)  Of
course, if you were to use a \textsf{book} entry for such a reference,
then you would need to define both \textsf{author} and
\textsf{publisher} using the name you here might have put in
\textsf{organization}.  (See 14.84; chicago:manual, dyna:browser,
natrecoff:camera.  Cp.\ also the new \textbf{standard} entry type.)

\paragraph*{\protect\mymarginpar{\textbf{misc}}}
\label{sec:ad:misc}
As its name suggests, the \textsf{misc} entry type was designed as a
hold-all for citations that didn't quite fit into other categories.
In \textsf{biblatex-chicago}, I have somewhat extended its
applicability, while retaining its traditional use.  Put simply, with
no \textsf{entrysubtype} field, a \textsf{misc} entry will retain
backward compatibility with the standard styles, so the usual
\textsf{howpublished}, \textsf{version}, and \textsf{type} fields are
all available for specifying an otherwise unclassifiable text, and the
\textsf{title} will be italicized.  (The \emph{Manual}, you may wish
to note, doesn't give specific instructions on how such citations
should be formatted, so when using the Chicago style I would recommend
you have recourse to this traditional entry type as sparingly as
possible.)

\mylittlespace If you do provide an \textsf{entrysubtype} field, the
\textsf{misc} type provides a means for citing unpublished letters,
memoranda, private contracts, wills, interviews, and the like, making
it something of an unpublished analogue to the \textsf{letter},
\textsf{article}, and \textsf{review} entry types (which see).  It
also works well for presenting online audio pieces, particularly dated
ones, like speeches.  Typically, such an entry will cite part of an
archive, and equally typically the text cited won't have a specific
title, but only a generic one, whereas an \textsf{unpublished} entry
will ordinarily have a specific author and title, and won't come from
a named archive.  The \textsf{misc} type with an \textsf{entrysubtype}
defined is the least formatted of all those specified by the
\emph{Manual}, so titles are in plain text by default.  (It is quite
possible, though somewhat unusual, for archival material to have a
specific title, rather than a generic one.  In these cases, you will
need to enclose the title inside a \cmd{mkbibquote} command manually.
Cf.\ roosevelt:speech, shapey:partita.)

\mylittlespace If you are presenting part of an unpublished archive,
then it's worth remembering that, as with the \textsf{letter} type,
the \emph{Manual} (15.54) suggests that the list of references will
usually contain only the name of the whole archived collection, with
more specific information about individual items provided in the text,
outside the parentheses.  If, on the other hand, \enquote{only one
  item from a collection has been mentioned in text, the entry may
  begin with the writer's name (if known).}  (See 14.211--12, 14.219,
14.221--31; house:papers cites a whole archive, while creel:house,
dinkel:agassiz, and spock:interview cite individual pieces.)

\mylittlespace As far as constructing your .bib entry goes, you should
first know that the absence of any date will not result in the
\enquote{n.d.}\ abbreviation automatically being provided.  Indeed, no
date at all will be required for entries referring to entire archival
collections.  In such entries, you may wish to use the word
\enquote{\texttt{classical}} as your \textsf{entrysubtype}, which will
have no effect on the list of references but will change the look of
the in-text citations (house:\break papers).  Instead of any date, the
citation will include the \textsf{title}, separated from the
\textsf{author's} name by a space, e.g., (House Papers).  This same
arrangement, happily, allows you easily to cite individual books of
the Bible, and also certain other sacred texts (14.238--41; genesis).
Please see under \textsf{entrysubtype} in
section~\ref{sec:fields:authdate} below for all the details of the
\texttt{classical} toggle.

\mylittlespace As for presenting the date of individual items, the
\emph{Manual} (14.224) allows in these entries, as it does in
documentation generally \enquote{if numerous dates occur} (9.35), for
a more streamlined presentation of dates using the day-month-year
form, different from the standard American month-day-year.  In
\textsf{letter} entries you use the \textsf{origdate} field to give
the date of individual letters, and it is always presented in the more
streamlined form.  Here, the same field will do exactly the same
thing, though with the added wrinkle that if you'd prefer to use the
standard day-month-year form you can, simply by putting the date into
the \textsf{date} field instead.  (Please choose one only in
\textsf{misc} entries with an \textsf{entrysubtype} --- in
\textsf{letter} entries the \textsf{date} refers to the published
collection.)  As with the \textsf{letter} type, if the only date
present is an \textsf{origdate}, you don't need to set the
\texttt{cmsdate} option in your .bib entry to make sure that that year
appears at the head of the entry (and in citations) --- this happens
automatically.  (Cf.\ particularly the documentation in
section~\ref{sec:fields:authdate} below, s.v.\ \enquote{date}, and
also the \textsf{letter} type above for some of the date-related
complications that can arise, and how you can address them with
judicious use of the \textsf{options}, \textsf{date}, and
\textsf{origdate} fields.)

\mylittlespace As in \textsf{letter} entries, the titles of
unpublished letters are of the form \texttt{Author to Recipient},
further information can be given in the \textsf{titleaddon} field,
while the \textsf{origlocation} field can hold the place where the
letter was written.  Interviews or similar pieces will have a
different sort of title, but all types will use the \textsf{note},
\textsf{organization}, \textsf{institution}, and \textsf{location}
fields (in ascending order of generality) to identify the archive,
though the \emph{Manual} specifies (14.227) that well-known
depositories don't usually need a city, state or country specified.
(The traditional \textsf{misc} fields are all still available, also.)

\mylittlespace There are a few more subtleties involved here.  Some
material (roosevelt:speech) includes a venue for the event recorded in
the archive, so I have added the \textsf{venue} field, which prints
\emph{before} the date, with the \textsf{origlocation} appearing after
it.  Somewhat confusingly, in published letters the
\textsf{origlocation} itself prints before the date, rather than
after, so if the inconsistency between published and unpublished
letters bothers you then you could conceivably use \textsf{venue}
instead of \textsf{origlocation} for that purpose here.  Also, the
\emph{Manual} offers several examples of specific location information
\emph{within} an archive, some of which appears \emph{before} the main
archive name, and some of which appears \emph{after} it.  I assume
this may depend on the exact nature of the archive itself, but in any
case you can try the \textsf{type} or \textsf{howpublished} fields for
the first case and the \textsf{number} field for the second.

\mylittlespace In all this class of archived material, the
\emph{Manual} (14.221) quite specifically requires more consistency
within your own work than conformity to some external standard, so it
is the former which you should pursue.  I hope that
\textsf{biblatex-chicago} proves helpful in this regard.

\mybigspace This \mymarginpar{\textbf{music}} is one of three
audiovisual entry types, and is intended primarily to aid in the
presentation of musical recordings that do not have a video component,
though it can also include audio books (auden:reading).  A DVD or VHS
of an opera or other performance, by contrast, should use the
\textbf{video} type instead, while an online music video will probably
need an \textbf{online} entry.  (Cf.\ \textsf{online} and
\textsf{video}; handel:messiah, horowitz:youtube.)  Because
\textsf{biblatex} --- and \textsc{Bib}\TeX\ before it --- were
designed primarily for citing book-like objects, some choices needed
to be made in assigning the various roles found on the back of a CD to
the fields in a typical .bib entry.  I have also implemented several
bibstrings to help in identifying these roles within entries.  The
17th edition of the \emph{Manual} once again revised its
recommendations for this type, but fortunately the changes are
additive, i.e., you can re-use 16th-edition citations but are
encouraged to peruse the following guidelines to see if there's any
information you might think of adding to bring your citations more
into line with the spec.

\mylittlespace These guidelines, in summary form, are:

{\renewcommand{\descriptionlabel}[1]{\qquad\textsf{#1}}
\begin{description}
\item[author =] composer, songwriter, or performer(s),
  depending on whom you wish to emphasize by placing them at the head
  of the entry.
\item[bookauthor:] Somewhat like an \textsf{author}, but it will hold
  the name associated with a whole album rather than an individual
  piece, should both be present, and will therefore appear in close
  association with the \textsf{booktitle} rather than at the head of
  the entry.
\item[editor, editora, editorb =] conductor, director or
  performer(s).  These will ordinarily follow the \textsf{title} of
  the work, though the usual \texttt{useauthor} and \texttt{useeditor}
  options can alter the presentation within an entry.  Because these
  are non-standard roles, you will need to identify them using the
  following:
\item[editortype, editoratype, editorbtype:] The most common roles,
  all associated with specific bibstrings (or their absence), will be
  \texttt{conductor}, \texttt{director}, \texttt{producer}, and,
  oddly, \texttt{none}.  The last is particularly useful when
  identifying the group performing a piece, as it usually doesn't need
  further specifying and this role prevents \textsf{biblatex} from
  falling back on the default \texttt{editor} bibstring.  The 17th
  edition (15.57) also seems to favor, in some circumstances, using
  strings to identify individual performers, e.g., \enquote{vocalist}
  or \enquote{pianist,} so even though there's no \cmd{bibstring}
  associated with these types you can now provide them, or anything
  else you need, in whatever form (\enquote{vocalist} or \enquote{sung
    by}) suits your citation.
\item[note:] This field can also hold contributors, perhaps
  collaborators or featured artists (holiday:fool, rihanna:umbrella).
\item[title, booktitle, maintitle:] As with the other audiovisual
  types, \textsf{music} serves as an analogue both to books and to
  collections, so the title will either be, e.g., the album title or a
  song title, in which latter case the album title would go into
  \textsf{booktitle}.  If you wish to cite a song that, as may happen,
  isn't part of any larger collection, your entry will in such a case
  have only a \textsf{title}, which \textsf{biblatex-chicago} would
  normally interpret as an album title.  You can now define an
  \textsf{entrysubtype} to let it know that the lone \textsf{title} is
  in fact a song (cf.\ naraya).  The \textsf{maintitle} might be
  necessary for something like a box set of \emph{Complete
    Symphonies}.
\item[chapter:] The 17th edition seems more keen on having track
  numbers for individual pieces, whether on a traditional format or on
  a streaming service.  The \textsf{chapter} field is the place for
  this information, and \textsf{biblatex-chicago} will automatically
  prepend the localized string \texttt{track} (cf.\ holiday:fool,
  rihanna:umbrella).
\item[publisher, series, number:] These three closely-associated
  fields are intended for presenting the catalog information provided
  by the music publisher.  The 17th edition generally only requires
  the \textsf{series} and \textsf{number} fields (nytrumpet:art),
  which hold the record label and catalog number, respectively.
  Alternatively, \textsf{publisher} would function as a synonym for
  \textsf{series} (holiday:fool), but there may be cases when you need
  or want to specify a publisher in addition to a label, as perhaps
  when a single publisher oversees more than one label.  You can
  certainly put all of this information into one of the above fields,
  but separating it may help make the .bib entry more readable.
\item[pubstate:] The \textsf{pubstate} field in \textsf{music} entries
  mainly has the usual meaning it has in other entry types, for which
  see the documentation of the field in
  section~\ref{sec:fields:authdate}, below.  If the field contains
  \texttt{reprint}, however, this has a special meaning in
  \textsf{music} entries, where it will transform the
  \textsf{origdate} from a recording date for an entire album into an
  original release date for that album, notice of which will be
  printed towards the end of a reference list entry, always assuming
  that the \textsf{origdate} hasn't already appeared at the head of
  the entry and in citations.  No \texttt{reprint} \cmd{bibstring}
  will be printed, as only the syntax of the reference will have been
  altered.
\item[date, eventdate, origdate:] The 17th edition of the
  \emph{Manual}, like the 16th, considers \textsf{music} citations
  without a date to be \enquote{generally unacceptable} (14.263),
  while if there is more than one date \enquote{the date of the
    original recording should be privileged} (15.57).  Finding these
  dates may take some research, but they will basically fall into two
  types, i.e., the date of the recording or the copyright / publishing
  date.  Recording dates go either in \textsf{origdate} (for complete
  albums) or \textsf{eventdate} (for individual tracks).  The current
  copyright or publishing date goes in the \textsf{date} field, while
  the original release date goes in \textsf{origdate}.  You may have
  noticed that the \textsf{origdate} has two slightly different uses
  --- you can tell \textsf{biblatex-chicago} which sort you intend by
  using the string \texttt{reprint} in the \textsf{pubstate} field,
  which transforms the \textsf{origdate} from a recording date into an
  original release date.  The style will automatically use the
  \textsf{eventdate} or the \textsf{origdate} in citations and at the
  head of the list of references, falling back on a \textsf{date} or
  even a \textsf{urldate} in their absence.  It will also prepend the
  bibstring \texttt{recorded} to any part of the \textsf{eventdate}
  that doesn't appear at the head of the list of references or, in the
  absence of the \textsf{pubstate} mechanism, to the
  \textsf{origdate}, or indeed to both.  You can modify what is
  printed here using the \textsf{userd} field, which acts as a sort of
  date type modifier.  In \textsf{music} entries, \textsf{userd} will
  be prepended to an \textsf{eventdate} if there is one, barring that
  to the \textsf{origdate}, barring that to a \textsf{urldate}, and
  absent those three to a \textsf{date}.  (See floyd:atom,
  holiday:fool, nytrumpet:art.)
\item[type, howpublished:] As in all the audiovisual entry types, the
  \textsf{type} field holds the medium of the recording, e.g., vinyl,
  33 rpm, 8-track tape, cassette, compact disc, mp3, ogg vorbis.  The
  \textsf{howpublished} field, newly included for the 17th edition,
  can hold similar information \enquote{for streaming audio formats
    and downloads} (14.263, 15.57). It can also, alternatively, hold
  the name of the streaming service, e.g., Spotify (cf.\
  rihanna:umbrella).
\end{description}}

The entries in \textsf{notes-test.bib} should at least give you a good
idea of how this all works, and that file also contains an example of
an audio book presented in a \textsf{music} entry.  If you browse the
examples in the \emph{Manual} you will see the sheer variety of
possibilities for presenting these sources, my intention being that
judicious manipulation of\ .bib entries should allow you to make
\textsf{biblatex-chicago} do what you want.  Please let me know if
I've ignored something you need.  (Cf. 14.263--64, 15.57;
\textsf{eventdate}, \textsf{origdate}, \textsf{userd};
\cmd{DeclareLabeldate} in section~\ref{sec:authformopts} and
\texttt{avdate} in section~\ref{sec:authpreset}; auden:reading,
beethoven:sonata29, bernstein:shostakovich, floyd:atom, holiday:fool,
nytrumpet:art, rubinstein:chopin.)

\mybigspace All \mymarginpar{\textbf{mvbook}\\\textbf{mvcollection}%
  \\\textbf{mvproceedings}\\\textbf{mvreference}} four of these entry
types function more or less as in standard \textsf{biblatex}.  I would
like, however, to emphasize a couple of things.  First, each is
aliased to the entry type that results from removing the
\enquote{\textbf{mv}} from their names.  Second, each has an important
role as the target of cross-references from other entries, the
\textsf{title} of the \textbf{mv*} entry \emph{always} providing a
\textsf{maintitle} for the entry referencing it.  If you want to
provide a \textsf{booktitle} for the referencing entry, please use
another entry type, e.g., \textbf{collection} for
\textbf{incollection} or \textbf{book} for \textbf{inbook}.  (These
distinctions are particularly important to the correct functioning of
the abbreviated references that \textsf{biblatex-chicago}, in various
circumstances, provides.  Please see the documentation of the
\textbf{crossref} field in section~\ref{sec:fields:authdate}, below.)

\mylittlespace On the same subject, when multi-volume works are
presented in the reference apparatus, the \emph{Manual} (14.116--22,
15.41) requires that any dates presented should be appropriate to the
specific nature of the citation.  In short, this means that a date
range that is right for the presentation of a multi-volume work in its
entirety isn't right for citing, e.g., a single volume of that work
which appeared in one of the years contained in the date range.
Because child entries will by default inherit all the date fields from
their parent (including the \textsf{endyear} of a date range), I have
turned off the inheritance of \textsf{date} and \textsf{origdate}
fields from all of the \textbf{mv*} entry types to any other entry
type.  When the dates of the parent and of the child in such a
situation are exactly the same, then this unfortunately requires an
extra field in the child's .bib entry.  When they're not the same, as
will, I believe, often be the case, this arrangement saves a lot of
annoying work in the child entry to suppress wrongly-inherited fields.
Other sorts of parent entries aren't affected by this.  See
harley:ancient:cart, harley:cartography, and harley:hoc for how this
might look.

\mylittlespace Finally, in order to cope with parts of the notes \&\
bibliography specification I have fairly thoroughly revised the
\textbf{mv*} types there for the 17th edition.  The author-date
specification is, as I read it, simpler, so I haven't revised these
types here.  If you should happen to leave a stray \texttt{maintitle
  relatedtype} in an entry when switching specifications that entry
will use the \texttt{multivolume relatedtype} instead, which will very
likely not be what you want.

\mybigspace One \mymarginpar{\textbf{online}} of the features of the
17th edition of the \emph{Manual} is the considerably extended, but
still scattered, treatment of online materials (8.189--92, 14.6--18,
14.159--63, 14.175--76, 14.187, 14.189, 14.205--10, 14.233, 15.4,
15.9, 15.49--52, 15.57).  The principles of that treatment have
changed somewhat, as the \emph{Manual} now places greater emphasis on
the \emph{location} of a source, which can in many cases outweigh, as
far as choosing an entry type goes, the \emph{nature} of the source.
Working out the correspondences between online sources and
\textsf{biblatex-chicago} entry types can, therefore, be tricky, so I
have included table~\ref{tab:online:adtypes} summarizing the
increasingly detailed instructions in the \emph{Manual}, along with
some further annotations here that might help to clarify it.

\afterpage{\clearpage\thispagestyle{longtable}

\begin{table}[h!]
  \caption[\hspace{-1em}Online Entry Types - Author-Date]%
  {Online materials and author-date entry types}
  \label{tab:online:adtypes}
  \centering\small\sffamily
  \hspace*{-6em}
  \begin{tabularx}{160mm}{@{}>{\raggedright}p{25mm}>{\raggedright}p{20mm}p{15mm}p{26mm}X@{}}
    \toprule
    Online Material & Entry Type & CMS Ref. & Sample Entry &
    Notes \\
    \cmidrule{1-1}\cmidrule(l){2-2}\cmidrule(l){3-3}\cmidrule(l){4-4}
    \cmidrule(l){5-5}
    Online edition of trad.\ publ.\ matter. &&&& Use the same
    entry type as you would choose were you citing it
    from a printed source.\\\addlinespace[.6mm]
    & @Book & 14.161-62 & james:ambassadors &
    CMS prefers (scanned) original page nos. \\\addlinespace[.6mm]
    & @Article @Review & 14.175 15.47-49 & black:infectious & If
    no \enquote{suitable URL} is available, e.g., if it points
    to a generic portal page rather than to an abstract,
    use the name of the commercial database in an addendum
    field instead.
    \\\addlinespace[.6mm]
    \cmidrule(r{2em}){1-1}\addlinespace[.6mm]
    Blogs&& 15.51 &&\\\addlinespace[1.5mm]
    \hspace{.5em} Single post & @Article & & amlen:hoot & The
    \textsf{maintitle} field holds the larger publication of which the
    blog is a part.\\\addlinespace[.6mm]
    \hspace{.5em} Whole Blog & @Periodical & & amlen:wordplay & This \&\
    the next usually not in the ref.\ list.
    \\\addlinespace[.6mm]
    \hspace{.5em} Comment & @Review & & viv:amlen & The
    \texttt{commenton} \textsf{relatedtype} helps manage these,
    in or out of the reference list. \\\addlinespace[.6mm]
    \cmidrule(r{2em}){1-1}\addlinespace[.6mm]
    Social Media & @Online & 15.52 && This includes anything ---
    posts, photos, videos --- on these and similar sites; the
    \emph{location} of the material defines its treatment. \\\addlinespace[2mm]
    \hspace{.5em} Mailing list or \hspace*{.5em} forum post & & 14.210 &
    powell:email & Posts on private lists are to be treated as
    \enquote{personal communications,} using @Misc w/
    \textsf{entrysubtype}.\\\addlinespace[.6mm]
    \hspace{.5em} Facebook&& 14.209 & diaz:surprise &\\\addlinespace[.6mm]
    \hspace{.5em} Twitter&&& obrien:recycle &\\\addlinespace[.6mm]
    \hspace{.5em} Instagram&&& souza:obama &\\\addlinespace[.6mm]
    \hspace{.5em} Comments / \hspace*{.5em} replies & & 14.210
    & braun:reply &This and the next are usually not included in the
    reference list. The
    \texttt{commenton} \textsf{relatedtype} helps manage them,
    in or out of that list.
    \\
    && 14.209 & licis:diazcomment &\\\addlinespace[.6mm]
    \cmidrule(r{2em}){1-1}\addlinespace[.6mm]
    Online Multimedia && 15.57 &&\\\addlinespace[.6mm]
    \hspace{.5em} Online video & @Online & 14.267 & pollan:plant &
    This category includes TED talks and most informal videos on
    YouTube and similar sites. \\\addlinespace[.6mm]
    \hspace{.5em} Online video, \hspace*{.5em} from a trad.\ \hspace*{.5em}
    journal & @Article &&
    kessler:nyt & You can use @Online, but this requires special
    formatting in the \textsf{note} or \textsf{titleaddon} field.
    \\\addlinespace[.6mm]
    \hspace{.5em} Published films in \hspace*{.5em} an archive & @Video &&
    weed:flatiron &\\\addlinespace[.6mm]
    \hspace{.5em} Podcasts & @Audio && danforth:podcast & Note the
    eventdate of the individual episode.\\\addlinespace[.6mm]
    \hspace{.5em} Archival audio & @Misc w/ \textsf{entrysubtype} & 14.264 &
    roosevelt:speech & Can have both a venue and an origlocation.
    \\\addlinespace[.6mm]
    \cmidrule(r{2em}){1-1}\addlinespace[.6mm]
    Streaming Media&&&& \\\addlinespace[.6mm]
    \hspace*{.5em} TV / Film & @Video & 14.265 & mayberry:brady &
    The streaming service is supplied by the URL.  The
    \texttt{tvepisode} entrysubtype is new in the 17th
    edition. \\\addlinespace[.6mm]
    \hspace*{.5em} Music & @Music & 14.263 & rihanna:umbrella &
    The streaming service is supplied by the howpublished field.
    \\\addlinespace[.6mm]
    \hspace*{.5em} News / Interviews & @Article @Review & 14.213 &
    bundy:macneil & Network information goes in the usera field.
    \\\addlinespace[.6mm]
    \cmidrule(r{2em}){1-1}\addlinespace[.6mm]
    Websites & @Online & 14.206-7 15.50 & evanston:library stenger:privacy &
    An online source \enquote{analogous
    to a traditionally printed
    work but [which] does not have (and never had) a printed counterpart}
    may now use an @Online entry, at your discretion.
    If you only have an access date, \enquote{n.d.}\ will appear as
    the publication date.
    \\\addlinespace[.6mm]
    \hspace*{.5em} Reference works, \hspace*{.5em} cited by alpha-
    \hspace*{.5em} betized entry & @InReference w/ entrysub-\par
    type & 14.233 & wikiped:bibtex & As above, you can choose the
    @Online treatment of the title, but it's best achieved
    using an @InReference entry w/ entrysubtype. \\\addlinespace[.6mm]
    \hspace*{.5em} Scientific data- \hspace*{.5em} bases & @Dataset &
    14.257 & genbank:db & Treated as author-title by
    default.\\\addlinespace[.6 mm]
    \bottomrule
  \end{tabularx}
\end{table}}

\mylittlespace The basic principle, as I've cited in the penultimate
entry of table~\ref{tab:online:adtypes}, is that \enquote{the title of a
  website that is analogous to a traditionally printed work but does
  not have (and never had) a printed counterpart can be treated like
  the titles of other websites, subject to editorial discretion}
(14.206).  This means that an intrinsically online entry like
stenger:privacy (citing CNN.com) need no longer be an \textsf{article}
but can be presented in an \textsf{online} entry.  (The same principle
applies to wikiped:bibtex, but because of the code facilitating
presentation of alphabetized entries in reference works, it's best in
this case to keep the \textsf{inreference} entry but add an
\textsf{entrysubtype} so that the \textsf{title} is presented as it
would be in an \textsf{online} entry.)  The corollary of the
principle, as the first entry in table~\ref{tab:online:adtypes}
suggests, is that an online edition of a printed work will generally
require the same entry type as that printed work itself would.  Blogs
are, therefore, somewhat anomalous in requiring the various periodical
types, though the \emph{Manual} does specify that if you're not sure
whether a website is a blog, then it probably requires the
\textsf{online} type (14.206).  Social media, on the other hand, are
very much subject to the first principle, requiring \textsf{online}
entries no matter whether the citation is of text, a photo, or a
video.  Without pretending that all of the correspondences flow
deductively from the basic principles, I hope that the table might
simplify most of your choices.  If something remains unclear, please
let me know and I'll see if I can improve it.

% %\enlargethispage{\baselineskip}

\mylittlespace A few more notes are in order. I designed the new
\textsf{relatedtype} \texttt{commenton} to facilitate citation of
online comments, and it is available in two entry types,
\textsf{online} and \textsf{review}.  In both types the \emph{Manual}
(15.51--52) recommends that such material appear \emph{only} in the
text and not in the reference list, but I have attempted to simplify
the presentation of such material wherever you want it to appear.
Following the specifications, then, the default when you use
\texttt{commenton} is for \textsf{biblatex-chicago-authordate} to
modify how your .bib entry appears in the .bbl file by setting both
\texttt{skipbib} and \texttt{cmsdate=full} in the \textsf{options}
field, so that nothing appears in the reference list and citations
present the full date and possibly also a time stamp (see below).
Further, the style sets the \textsf{verbc} field so that these entries
don't interfere with the provision of extra date letters --- the full
date and time should be enough to individuate separate comments.
Finally, the style creates a new \textsf{customc} entry in your .bbl
file which you can cite after your initial \texttt{commenton} entry
using \cmd{autocites} and which will, as a comment to your initial
entry, say whether it's a comment or a reply or what have you, and
then giving the short citation of that upon which it is a comment.

\mylittlespace As an example, take the Facebook post diaz:surprise,
which does appear in the reference list.  The entry licis:diazcomment
presents a comment on this post using the \textsf{relatedtype}
\texttt{commenton}, so \textsf{biblatex-chicago-authordate} creates a
new entry, diaz:surprise-customc, which won't appear in your .bib file
(which is never altered) but in the .bbl file that \textsf{biber}
produces to supply \textsf{biblatex} with the data for its citations.
A command like
\verb+\autocites{licis:diazcomment}{diaz:surprise-customc}+ will
produce a citation like (Licis, February 24, 2016; comment on D\'iaz
2016).  You can alter the string connecting the two citations (by
default \verb+\bibstring{commenton}+) by using the
\textsf{relatedstring} field in the first of them (cf.\
powell:comment).  (Note how minimal the .bib entry of a comment using
this system can be --- \textsf{author}, \textsf{related},
\textsf{relatedtype}, and \textsf{date} are pretty much the only
fields required.)

\mylittlespace Those who want \textsf{online} comments to appear in
the reference list can still use the \texttt{commenton}
\textsf{relatedtype}, and the same citation of the commented piece
will appear there, connected by the same string that the
\textsf{customc} entry provides.  Here, though, you can also provide a
separate \textsf{title} for the comment, and/or a separate
\textsf{url} for it, should they exist, which will be printed
before/after the citation of the commented piece, respectively.  (In
\textsf{review} entries, which use the same \textsf{relatedtype}, only
the generic title is available, as is always the case with such
entries.)  If you manually set either (or both) of the
\texttt{cmsdate} or the \texttt{skipbib} options in your entry then
\textsf{biblatex-chicago} will assume you want to hand-craft that
entry without its intervention, though it will still provide the
virtual \textsf{customc} entry in your .bbl file, as that may still
prove convenient.  Note also that any \textsf{verbc} field you provide
will never be altered by the package.

\mylittlespace In general, constructing an \textsf{online}\ .bib file
entry is much the same as in \textsf{biblatex}.  The \textsf{title}
field would contain the title of the page, the \textsf{organization}
field could hold the title or owner of the whole site.  If there is no
specific title for a page, but only a generic one, then such a title
should go in \textsf{titleaddon}, not forgetting to begin that field
with a lowercase letter so that capitalization will work out
correctly.  It is worth remarking here, too, that the \emph{Manual}
(15.50) strongly prefers, if they're available, revision dates to
access dates when documenting online material.  If there is only a
\textsf{urldate} in an entry, and that date is an access date (i.e.,
there's no \textsf{userd} field), then \enquote{n.d.} will appear in
citations and at the head of the entry in the reference list.
Moreover, given how rapidly online sources can change (14.191, 14.209,
14.233), a time stamp may often be necessary further to specify a
revision date (\textsf{urldate}) or the date of a comment or reply
(\textsf{date}).  This time specification should be added to the date
field using \textsf{biblatex's} standard format, i.e.,
\texttt{2008-07-01T10:18:00}.  If a time zone is needed, then a
separate \textsf{timezone} or \textsf{urltimezone} field is the best
way, as it allows you to provide the initialisms that the
\emph{Manual} prefers (10.41, 14.191).  On all of this please see
\textsf{date}, \textsf{timezone}, \textsf{urldate}, \textsf{userd},
and \textsf{verbc} in section~\ref{sec:fields:authdate}, below.  Cf.\
also the documentation of the \texttt{commenton} \textsf{relatedtype}
in section~\ref{sec:authrelated}.

\mybigspace The \mymarginpar{\textbf{patent}} \emph{Manual} is very
brief on the subject of patents (15.55), but very clear about which
information it wants you to present, so such entries may not work well
with other \textsf{biblatex} styles.  Chicago's author-date style
prefers the \emph{later} of the two possible dates to appear in
citations and at the head of the entry in the list of references.  If
a patent has been filed but not yet granted, then you can place the
filing date in either the \textsf{date} field or the \textsf{origdate}
field, and \textsf{biblatex-chicago-authordate} will automatically
prepend the bibstring \texttt{patentfiled} to it.  If the patent has
been granted, then you put the filing date in the \textsf{origdate}
field, and you put the date it was issued in the \textsf{date} field,
to which the bibstring \texttt{patentissued} will automatically be
prepended, and it is this later date that will head the entry and
appear in citations.  The patent number goes in the \textsf{number}
field, and you should use the standard \textsf{biblatex} bibstrings in
the \textsf{type} field.  Though it isn't mentioned by the
\emph{Manual}, \textsf{biblatex-chicago-authordate} will print the
\textsf{holder} after the \textsf{author}, if you provide one.
Finally, the style capitalizes the \textsf{title} sentence-style,
which seems to be the generally-accepted convention across all Chicago
specifications.  If you need to keep a word capitalized then you
should wrap it in curly braces.  See petroff:impurity.

\mybigspace The \mymarginpar{\textbf{performance}} 17th edition of
the \emph{Manual} includes a new section (14.266) on citing live
performances, and even though such references can usually be limited
to the main text (cf.\ 15.57) it may sometimes be useful to include
them in a reference list.  Since \textsf{biblatex} provides the
\textbf{performance} type, albeit without using it in its standard
styles, I though it might be useful to define it for
\textsf{biblatex-chicago}, particularly as the other option for such
material is the \textsf{misc} entry without any \textsf{entrysubtype},
and that entry type is already somewhat overloaded, though you can
still use it if you wish.

\mylittlespace Such entries will generally have a \textsf{title}, a
\textsf{venue}, a \textsf{location} for the venue, and a \textsf{date}
for the performance, along with a possible plethora of authorial
and/or editorial roles depending on which sorts of contributor(s) you
wish to emphasize in the citation.  The \textsf{editor[abc]} and
\textsf{editor[abc]type} fields should be most helpful here.  I have
included strings for \texttt{choreographer} in all localization files,
but for others you may need to provide them in the
\textsf{editor[abc]type} fields as you wish them printed ---
\textsf{biblatex-chicago} will automatically capitalize any that start
with a lowercase letter.  For the author-date styles it will probably
be convenient to allow one of these names to appear at the head of the
entry and in citations, as this will facilitate the appearance of the
\textsf{extradate} to distinguish, e.g., multiple performances of the
same work or performances of different works by the same producer or
choreographer.

\mybigspace This \mymarginpar{\textbf{periodical}} is the standard
\textsf{biblatex} entry type for presenting an entire issue of a
periodical, rather than one article within it.  It has the same
function in \textsf{biblatex-chicago}, and in the main uses the same
fields, though in keeping with the system established in the
\textsf{article} entry type (which see) you'll need to provide
\textsf{entrysubtype} \texttt{magazine} if the periodical you are
citing is a \enquote{newspaper} or \enquote{magazine} instead of a
\enquote{journal.}  Also, remember that the \textsf{note} field is the
place for identifying strings like \enquote{special issue,} with its
initial lowercase letter to activate the automatic capitalization
routines, though this isn't strictly necessary in the author-date
styles.  (See \emph{Manual} 14.187; good:wholeissue.)

\enlargethispage{-\baselineskip}

\mylittlespace It is worth noting a few things.  First, the
\textsf{titleaddon} field is now available in these entries, but as
the \textsf{title} here is analogous to the \textsf{journaltitle} in
\textsf{article} or \textsf{review} entries, the new
\mycolor{\texttt{jtitleaddon}} option (section~\ref{sec:authpreset})
governs the punctuation separating the \textsf{titleaddon} from the
\textsf{title}.  Second, the special \textsf{biblatex} field
\textsf{shortjournal} allows you to present shortened
\textsf{journaltitles} in \textsf{article}, \textsf{review}, and
\textsf{periodical} entries, as well as facilitating the creation of
lists of journal abbreviations in the manner of a \textsf{shorthand}
list.  Because the \textsf{periodical} type uses the \textsf{title}
field instead of \textsf{journaltitle}, \textsf{biblatex-chicago}
automatically copies any \textsf{shorttitle} field, if one is present,
into \textsf{shortjournal}.  Please see the documentation of
\textbf{shortjournal} in section~\ref{sec:fields:authdate} for all the
details on how this works.  Finally, although the 17th edition
recommends that references to whole blogs, as opposed to individual
blog posts, need appear only in the text (15.51), using the
\textsf{periodical} type for such material can help with this, in or
out of the reference list.  The new \texttt{authortitle} entry option
will ensure the presence of the name of the blog (as opposed to the
non-existent date) in citations, and you could also use a
\cmd{citeurl} command to give the URL in the text (or a note).
Alternately, you could let the entry appear in the reference list and
cite it in the usual way.  In that list the \emph{Manual} (14.208)
recommends that you include the name of any larger (usually
periodical) publication of which the blog is a part.  The
\textsf{maintitle} field (with \textsf{mainsubtitle} and
\textsf{maintitleaddon}, if needed) is the place for it. Cf.\
amlen:wordplay.

\mybigspace This \mymarginpar{\textbf{proceedings}} is the standard
\textsf{biblatex} and \textsc{Bib}\TeX\ entry type, now also including
the \textsf{eventdate}, \textsf{eventtitle}, \textsf{eventtitleaddon},
and \textsf{venue} fields for identifying the event that produced the
\textsf{proceedings}.  The package can provide automatically
abbreviated references in the reference list when you use a
\textsf{crossref} or an \textsf{xref} field.  The functionality is not
enabled by default, but you can enable it in the preamble or in the
\textsf{options} field using the \texttt{booklongxref} option.  Please
see \textbf{crossref} in section~\ref{sec:fields:authdate} and
\texttt{booklongxref} in section~\ref{sec:authpreset}, below.

\mybigspace This \mymarginpar{\textbf{reference}} entry type is
aliased to \textsf{collection} by the standard \textsf{biblatex}
styles, but I intend it to be used in cases where you need to cite a
reference work but not an alphabetized article or articles in that
work.  This could be because it doesn't contain such articles, and yet
you still want the entry in the list of references to start with the
\textsf{title}.  Indeed, the only differences between it and
\textsf{inreference} are the lack of a \textsf{lista} field to present
an alphabetized entry, and the fact that any \textsf{postnote} field
will be printed verbatim, rather than formatted as an alphabetized
entry.  (Cf.\ \textsf{inreference}, above.)

\mybigspace This \mymarginpar{\textbf{report}} entry type is a
\textsf{biblatex} generalization of the traditional \textsc{Bib}\TeX\
type \textsf{techreport}.  Instructions for such entries are rather
thin on the ground in the \emph{Manual} (8.186, 14.220), so I have
followed the generic advice about formatting it like a book, and hope
that the results conform to the specification.  At least one user has
indicated a need, now filled, for an \texttt{unpublished}
\textsf{entrysubtype}, which prints the \textsf{title} inside
quotation marks (or, in \textsf{authordate-trad}, in plain roman)
instead of in italics, but affects nothing else.  This detail aside,
the type's main peculiarities are the \textsf{institution} field in
place of a \textsf{publisher}, the \textsf{type} field for identifying
the kind of report in question, and the \textsf{isrn} field containing
the International Standard Technical Report Number of a technical
report.  As in standard \textsf{biblatex}, if you use a
\textsf{techreport} entry, then the \textsf{type} field automatically
defaults to \verb+\bibstring{techreport}+.  As with \textsf{booklet}
and \textsf{manual}, you can also use a \textsf{book} entry, putting
the report type in \textsf{note} and the \textsf{institution} in
\textsf{publisher}.  (See herwign:office.)

\mybigspace As \mymarginpar{\textbf{review}} its name suggests, the
\textsf{review} entry type was designed for reviews published in
periodicals, and if you've already read the \textsf{article}
instructions above --- if you haven't, I recommend doing so now ---
you'll know that \textsf{review} serves as well for citing other sorts
of material with generic titles, like letters to the editor,
obituaries, interviews, online comments and the like.  The primary
rule is that any piece that has only a generic title, like
\enquote{review of \ldots,} \enquote{interview with \ldots,} or
\enquote{obituary of \ldots,} calls for the \textsf{review} type.  Any
piece that also has a specific title, e.g., \enquote{\enquote{Lost in
    \textsc{Bib}\TeX,} an interview with \ldots,} requires an
\textsf{article} entry.  (This assumes the text is found in a
periodical of some sort.  Were it found in a book, then the
\textsf{incollection} type would serve your needs, and you could use
\textsf{title} and \textsf{titleaddon} there.  While we're on the
topic of exceptions, the \emph{Manual} includes an example --- 14.213
--- where the \enquote{Interview} part of the title is considered a
subtitle rather than a titleaddon, said part therefore being included
inside the quotation marks and capitalized accordingly.  Not having
the journal in front of me I'm not sure what prompted that decision,
but \textsf{biblatex-chicago} would obviously have no difficulty
coping with such a situation.)

\mylittlespace Once you've decided to use \textsf{review}, then you
need to determine which sort of periodical you are citing, the rules
for which are the same as for an \textsf{article} entry.  If it is a
\enquote{magazine} or a \enquote{newspaper}, then you need an
\textsf{entrysubtype} \texttt{magazine}, or the synonymous
\textsf{entrysubtype} \texttt{newspaper}.  The generic title goes in
\textsf{title} and the other fields work just as as they do in an
\textsf{article} entry with the same \textsf{entrysubtype}, including
the substitution of the \textsf{journaltitle} for the \textsf{author}
if the latter is missing. (See 14.190--91, 14.195--96, 14.201--4,
14.213, 15.49; barcott:review, bundy:macneil, Clemens:letter,
gourmet:052006, kozinn:review, nyt:trevorobit, unsigned:ranke,
wallraff:word.)  If, on the other hand, the piece comes from a
\enquote{journal,} then you don't need an \textsf{entrysubtype}.  The
generic title goes in \textsf{title}, and the remaining fields work
just as they do in a plain \textsf{article} entry.  (See 14.202;
ratliff:review.)

\mylittlespace The \emph{Manual} now suggests that, no matter which
citation style you are using, it is \enquote{usually sufficient to
  cite newspaper and magazine articles entirely within the text}
(15.47).  This involves giving the title of the journal and the full
date of publication in a parenthetical reference, including any other
information in the main text (14.206), thereby obviating the need to
present such an entry in the list of references.  To utilize this
method in the author-date styles, in addition to a \texttt{magazine}
\textsf{entrysubtype}, you'll need to place \texttt{cmsdate=full} into
the \textsf{options} field, including \texttt{skipbib} there as well
to stop the entry printing in the list of references.  If the entry
only contains a \textsf{date} and \textsf{journaltitle} that's enough,
but if it's a fuller entry also containing an \textsf{author} then
you'll also need \texttt{useauthor=false} in the \textsf{options}
field.  Other surplus fields will be ignored.  (See osborne:poison.)

% %\enlargethispage{\baselineskip}

\mylittlespace \textsf{Biblatex-chicago} also, at the behest of
Bertold Schweitzer, supports the \textsf{relatedtype}
\texttt{reviewof}, which allows you to use the \textsf{related}
mechanism to provide information about the work being reviewed,
thereby simplifying how much information you need to provide in the
reviewing entry.  In particular, it relieves you of the need to
construct \textsf{title} or \textsf{titleaddon} fields like:
\verb+review of \mkbibemph{Book Title} by Author+, as the
\textsf{related} entry's \textsf{title} automatically provides the
\textsf{title} in the \textsf{review} type and the \textsf{titleaddon}
in the \textsf{article} type, with the \textsf{related} mechanism
providing the connecting string.  This may be particularly helpful if
you need to cite multiple reviews of the same work.  Please see
section \ref{sec:authrelated} for further information.

\mylittlespace Most of the onerous details are the same as I described
them in the \textbf{article} section above, but I'll repeat some of
them briefly here.  If anything in the \textsf{title} needs
formatting, you need to provide those instructions yourself, as the
default is completely plain.  \textsf{Author}-less \textsf{reviews}
are treated just like similar \textsf{articles} --- with an
\textsf{entrysubtype}, the \textsf{journaltitle} replaces the author
in citations and heads the entry in the list of references, without an
\textsf{entrysubtype} the \textsf{title} does the same.  In the former
case, \textsf{Biber} handles the sorting for you, but in the latter
you'll need a \textsf{sortkey} because \textsf{journaltitle} comes
before \textsf{title} in the sorting scheme.  (14.204; gourmet:052006,
nyt:trevorobit, unsigned:ranke, and see \cmd{DeclareSortingTemplate}
in section~\ref{sec:authformopts}, below.).  As in \textsf{misc}
entries with an \textsf{entrysubtype}, words like \enquote{interview,}
\enquote{review,} and \enquote{letter} only need capitalization after
a full stop, so you can start the \textsf{title} field with a
lowercase letter and let the automatic field formatting with
\cmd{autocap} do its work, though this isn't strictly necessary with
\textsf{biblatex-chicago-authordate}.

\mylittlespace A few details of the \textsf{review} type are fairly
new, and in particular have changed between the 16th and 17th editions
of the \emph{Manual}.  As I mentioned above, blogs are best treated as
\textsf{articles} with \texttt{magazine} \textsf{entrysubtype},
whereas comments on those blogs --- or replies to those comments,
etc.\ --- need the \textsf{review} type with the same
\textsf{entrysubtype}.  The 17th edition recommends that blog comments
appear only in the text, and not in the reference list (15.51), so
just as with comments in social media threads, for which see the
\textbf{online} type above, I have provided the \texttt{commenton}
\texttt{relatedtype} to simplify the presentation of such material
wherever you want it to appear.  Following the specifications, then,
the default when you use \texttt{commenton} is for
\textsf{biblatex-chicago-authordate} to modify how your .bib entry
appears in the .bbl file by setting both \texttt{skipbib} and
\texttt{cmsdate=full} in the \textsf{options} field, so that nothing
appears in the reference list and citations present the full date and
possibly also a time stamp (see below).  Further, the style sets the
\textsf{verbc} field so that these entries don't interfere with the
provision of extra date letters --- the full date and time should be
enough to individuate separate comments.  Finally, the style creates a
new \textsf{customc} entry in your .bbl file which you can cite after
your initial \texttt{commenton} entry using \cmd{autocites} and which
will, as a comment to your initial entry, say whether it's a comment
or a reply or what have you, and then giving the short citation of
that upon which it is a comment.

\mylittlespace As an example, take the blog ellis:blog, which does
appear in the reference list.  The entry ac:comment presents a comment
on this post using the \textsf{relatedtype} \texttt{commenton}, so
\textsf{biblatex-chicago-authordate} creates a new entry,
ellis:blog-customc, which won't appear in your .bib file (which is
never altered) but in the .bbl file that \textsf{biber} produces to
supply \textsf{biblatex} with the data for its citations.  A command
in your document like
\verb+\autocites{ac:comment}{ellis:blog-customc}+ will produce a
citation like (AC, July 1, 2008, 10:18 a.m.; comment on Ellis 2008).
You can alter the string connecting the two citations (by default
\verb+\bibstring{commenton}+) by using the \textsf{relatedstring}
field in the first of them.  (Note how minimal the .bib entry of a
comment using this system can be --- \textsf{author},
\textsf{entrysubtype}, \textsf{related}, \textsf{relatedtype}, and
\textsf{eventdate} are pretty much the only fields required.)

\mylittlespace Those who want \textsf{online} comments to appear in
the reference list can still use the \texttt{commenton}
\textsf{relatedtype}, and the same citation of the commented piece
will appear there, connected by the same string that the
\textsf{customc} entry provides.  Here, though, you can also provide a
separate \textsf{url} for the comment, should it exist, which will be
printed after the citation of the commented piece.  (In
\textsf{online} entries, which use the same \textsf{relatedtype}, you
can also provide a separate title for the comment.)  If you manually
set either (or both) of the \texttt{cmsdate} or the \texttt{skipbib}
options in your entry then \textsf{biblatex-chicago} will assume you
want to hand-craft that entry without its intervention, though it will
still provide the virtual \textsf{customc} entry in your .bbl file, as
that may still prove convenient.  Note also that any \textsf{verbc}
field you provide will never be altered by the package.  (Please see
the documentation of this \textsf{relatedtype} in
section~\ref{sec:authrelated}, that of \textsf{verbc} in
section~\ref{sec:fields:authdate}, and also the information about
online materials in table~\ref{tab:online:adtypes}.)

\mylittlespace The new edition of the \emph{Manual} retains the
requirement for a date closely associated with the comment (14.208,
15.51), so in such entries you now have a choice.  If you are using
the \texttt{commenton} \textsf{relatedtype}, you can use the
\textsf{date} or \textsf{eventdate} indifferently, as even when you
print the entry in the reference list the reference to the main blog
provides its own date.  If, in 16th-edition style, you print a fuller
entry in the reference list, then you'll need the \textsf{eventdate}
for the comment, as the \textsf{date} applies to the main blog post.
If you need a time stamp in addition, as may frequently be the case
with multiple contributions by the same author to a single thread,
then you should now use the standard \textsf{biblatex} time-stamp
format (e.g., \texttt{2008-07-01T10:18:00}) in whichever of the two
date fields you're using, and not, as previously, in the
\textsf{nameaddon} field.  \textsf{Biblatex-chicago} will format and
print it appropriately.  This change allows the \textsf{nameaddon}
field to revert to its primary use, which is to provide extra
information about the \textsf{author}.  In blog comments, this could
include the commenter's geographical location, which you need to
enclose in parentheses, as I've removed the automatic square brackets
from this field to allow it this more general usefulness.  You can, of
course, still provide your own square brackets in \textsf{review}
entries to indicate pseudonymous authorship, which is the standard
function of \textsf{nameaddon} in most entry types.  The package
options \texttt{nameaddonformat} and \texttt{nameaddonsep} can help
here, as well.  See sections~\ref{sec:authuseropts} and
\ref{sec:authpreset}, below.  Please also see the documentation of
\textsf{date}, \textsf{eventdate}, and \textsf{timezone} in
section~\ref{sec:fields:authdate}, \cmd{DeclareLabeldate} in
section~\ref{sec:authformopts}, and \texttt{avdate} in
section~\ref{sec:authpreset}.

\mylittlespace For the reasons I explained in the \textsf{article}
docs above, I have brought the \textsf{article} and \textsf{review}
entry types into line with most of the other types in allowing the use
of the \textsf{namea} and \textsf{nameb} fields in order to associate
an editor or a translator specifically with the \textsf{title}.  The
\textsf{editor} and \textsf{translator} fields, in strict homology
with other entry types, are associated with the \textsf{issuetitle} if
one is present, and with the \textsf{title} otherwise.  The usual
string concatenation rules still apply --- cf.\ \textsf{editor} and
\textsf{editortype} in section~\ref{sec:fields:authdate}, below.

\mylittlespace Finally, the special \textsf{biblatex} field
\textsf{shortjournal} allows you to present shortened
\textsf{journaltitles} in \textsf{review} entries, as well as in
\textsf{article} and \textsf{periodical} entries, and it facilitates
the creation of lists of journal abbreviations in the manner of a
\textsf{shorthand} list.  Please see the documentation of
\textbf{shortjournal} in section~\ref{sec:fields:authdate} for all the
details on how this works.

\mybigspace In \mymarginpar{\textbf{standard}} older releases it was
fairly straightforward to present published national or international
standards using a \textsf{book} entry, but with some additional
specifications now included in the 17th edition of the \emph{Manual}
(14.259, 15.37) I think it might be helpful to provide a separate
entry type.  The \textbf{standard} type has long existed in
\textsf{biblatex}, though none of its included styles use it.  In
\textsf{biblatex-chicago} constructing such an entry is mostly
straightforward.  The organization responsible for the standard goes
in \textsf{organization}, the title in \textsf{title}, and the
\textsf{series} and \textsf{number} fields provide the ID of the
standard.  The \textsf{date} field generally provides the publication
date, though for some standards there may also be a later
reaffirmation date (or similar), for which you can use the
\textsf{eventdate}.  To choose which year appears in citations, the
\textbf{standard} type follows, by default, the same ordering as
\textsf{review} and \textsf{music} entries, so that the
\textsf{eventdate} will, if present, provide the year.  (Cf.\
\texttt{avdate} in section~\ref{sec:authpreset}, below.)

\mylittlespace Now, for the peculiarities.  In the reference list, the
\textsf{organization} will appear at the head of the entry, and will
be reprinted as the publisher.  If you wish to provide a shortened
version for the second appearance, then the \textsf{publisher} field
is the place for it.  You can also use an \textsf{author} instead of
an \textsf{organization}, but in such a case you'll have to provide a
\textsf{publisher}, and no matter which field you choose to appear at
the head of the entry you'll usually have to think about providing
some sort of abbreviated form for citations.  A \textsf{shortauthor}
will appear only in citations, while a \textsf{shorthand} can also
appear at the head of the entry in the list of references.  (If you
provide the latter, \textsf{biblatex-chicago} will automatically sort
entries by it.)  Any named \textsf{editor} or \textsf{namec} will, as
per the specification, \emph{not} appear at the head of entries.  You
can really only alter this by using a \textsf{book} entry, instead.
(Cf.\ w3c:xml, and the \textsf{shorthand} docs on
page~\pageref{sec:ad:shorthand}, below.)

\mylittlespace Finally, it is distinctly possible that an entry with
two dates will need somehow to specify just what sort of dates are
involved.  The usual \textsf{biblatex-chicago} method is the
\textsf{userd} field, and here that field will act as a date-type for
the \textsf{date} field itself, assuming as usual that there is no
\textsf{urldate}.  For the \textsf{eventdate}, you'll need to use
\textsf{howpublished}, which I have commandeered for this purpose in a
few other entry types, as well.  (Cf.\ niso:bibref and
\textbf{howpublished} in section~\ref{sec:fields:authdate}, below.)

\mybigspace This \mymarginpar{\textbf{suppbook}} is the entry type to
use if the main focus of a reference is supplemental material in a
book or in a collection, e.g., an introduction, afterword, or forward,
either by the same or by a different author.  There are two mechanisms
in \textsf{biblatex-chicago} for producing such a citation.  First,
these three just-mentioned types of material, and only these three
types, can be referenced using the \textsf{introduction},
\textsf{afterword}, or \textsf{foreword} fields, a system that
requires you simply to define one of them in any way and leave the
others undefined.  The macros don't use the text provided by such an
entry, they merely check to see if one of them is defined, in order to
decide which sort of pre- or post-matter is at stake, and to print the
appropriate string before the \textsf{title} in the list of
references, and possibly also in the list of shorthands.  This
mechanism works without modification across multiple languages, but I
have also provided functionality which allows you to cite any sort of
supplemental material whatever, using the \textsf{type} field.  Under
this second system, simply put the nature of the material, including
the relevant preposition, in that field, beginning with a lowercase
letter so \textsf{biblatex} can decide whether it needs capitalization
depending on the context.  Examples might be \enquote{\texttt{preface
    to}} or \enquote{\texttt{colophon of}.}  (Please note, however,
that unless you use a \cmd{bibstring} command in the \textsf{type}
field, the resultant entry will not be portable across languages.)

\mylittlespace The other rules for constructing your .bib entry remain
the same.  The \textsf{author} field refers to the author of the
introduction or afterword, while \textsf{bookauthor} refers to the
author of the main text of the work, if the two differ.  Recent
editions of the \emph{Manual} requires that you include the page range
for the cited part in the list of references.  As ever, if the focus
of the reference is the main text of the book, but you want to mention
the name of the writer of an introduction or afterword for
completeness, then the normal \textsf{biblatex} rules apply, and you
can just put their name in the appropriate field of a \textsf{book}
entry, that is, in the \textsf{foreword}, \textsf{afterword}, or
\textsf{introduction} field.  (See \emph{Manual} 14.110;
friedman:intro, polakow:afterw, prose:intro).

%%\enlargethispage{\baselineskip}

\mybigspace This \mymarginpar{\textbf{suppcollection}} fulfills a
function analogous to \textsf{suppbook}.  Indeed, I believe the
\textbf{suppbook} type can serve to present supplemental material in
both types of work, so this entry type is an alias to
\textsf{suppbook}, which see.

\mybigspace This \mymarginpar{\textbf{suppperiodical}} type is
intended to allow reference to generically-titled works in
periodicals, such as regular columns or letters to the editor.
\textsf{Biblatex} also provides the \textsf{review} type for this
purpose, so in both Chicago styles \textsf{suppperiodical} is an alias
of \textsf{review}.  Please see above under \textbf{review} for the
full instructions on how to construct a .bib entry for such a
reference.

\mybigspace The \mymarginpar{\textbf{unpublished}}
\textsf{unpublished} entry type works largely as it does in standard
\textsf{biblatex}, though it's worth remembering that you should use a
lowercase letter at the start of your \textsf{note} field (or perhaps
an\ \cmd{autocap} command in the somewhat contradictory
\textsf{howpublished}, if you have one) for material that wouldn't
ordinarily be capitalized except at the beginning of a sentence.
Thanks to a bug report by Henry D. Hollithron, such entries will print
information about any \textsf{editor}, \textsf{translator},
\textsf{compiler}, etc., that you include in the .bib file.  Also,
conforming to the indications of the \emph{Manual}, and thanks to the
prompting of Jan David Hauck, you can use the \textsf{venue},
\textsf{eventdate}, \textsf{eventtitle}, and \textsf{eventtitleaddon}
fields further to specify unpublished conference papers and the like
(14.216--18; nass:address).

\mybigspace This \mymarginpar{\textbf{video}} is the last of the
three audiovisual entry types, and as its name suggests it is intended
for citing visual media, be it films of any sort or TV shows,
broadcast, on the Net, on VHS, DVD, or Blu-ray, though it will serve
as well, I think, for radio broadcasts of plays or drama serials.  As
with the \textsf{music} type discussed above, certain choices had to
be made when associating the production roles found, e.g., on a DVD,
to those bookish ones provided by \textsf{biblatex}.  The 17th edition
of the \emph{Manual} once again revised its recommendations for this
type, but fortunately the changes are additive, i.e., you can re-use
16th-edition citations but are encouraged to peruse the following
guidelines to see if there's any information you might think of adding
to bring your citations more into line with the spec.  Here are the
main guidelines:

{\renewcommand{\descriptionlabel}[1]{\qquad\textsf{#1}}
\begin{description}
\item[author:] This will not infrequently be left undefined, as the
  director of a film should be identified as such and therefore placed
  in the \textsf{editor} field with the appropriate
  \textsf{editortype} (see below).  You will need it, however, to
  identify the composer of, e.g., an oratorio on VHS (handel:messiah),
  or perhaps the provider of commentaries or other extras on a film
  DVD (cleese:holygrail).
\item[editor, editora, editorb:] The director or producer, or possibly
  the performer or conductor in recorded musical performances.  These
  will ordinarily follow the \textsf{title} of the work, though the
  usual \texttt{useauthor} and \texttt{useeditor} options can alter
  the presentation within an entry.  Because these are non-standard
  roles, you will need to identify them using the following:
\item[editortype, editoratype, editorbtype:] The most common roles,
  all associated with specific bibstrings (or their absence), will
  likely be \texttt{director}, \texttt{producer}, and, oddly,
  \texttt{none}.  The last is particularly useful if you want to
  identify performers, as they usually don't need further specifying
  and this role prevents \textsf{biblatex} from falling back on the
  default \texttt{editor} bibstring.  Any other roles you want to
  emphasize, even if there is no pre-defined \cmd{bibstring}, can be
  provided here, and will be printed as-is, contextually capitalized.
  (Cf.\ hitchcock:nbynw.)
\item[title, titleaddon, booktitle, booktitleaddon, maintitle:] As
  with the other 2 audiovisual types, \textsf{video} serves as an
  analogue both to books and to collections, so the \textsf{title} may
  be of a whole film DVD or of a TV series, or it may identify one
  episode in a series or one scene in a film.  In the latter cases,
  the title of the whole would go in \textsf{booktitle}.  The
  \textsf{booktitleaddon} field is the place for specifying the season
  and/or episode number of a TV series, while the \textsf{titleaddon}
  is for any information that needs to come between the \textsf{title}
  and the \textsf{booktitle} (american:crime, cleese:holygrail,
  friends:leia, handel:messiah, hitchcock:nbynw, mayberry:brady).  As
  in the \textsf{music} type, \textsf{maintitle} may be necessary for
  a boxed set or something similar.
\item[entrysubtype:] If, for some reason, you want to cite an
  individual episode or scene without reference to any larger unit,
  then your entry will contain only a \textsf{title}, which
  \textsf{biblatex-chicago} would normally interpret as the title of a
  complete film or TV series.  In such a case, you'll need to define
  an \textsf{entrysubtype} to let it know that the lone \textsf{title}
  is such a sub-unit.  In quite a different syntactic transformation,
  the 17th edition (14.265) now recommends that, when presenting
  episodes from a TV series, the name of the series
  (\textsf{booktitle}) comes before the episode name (\textsf{title}).
  The exact string \texttt{tvepisode} in the \textsf{entrysubtype}
  field achieves this reversal, which includes using the
  \textsf{booktitle} as a \textsf{sorttitle} in the reference list and
  also as the \textsf{labeltitle} in short notes.
\item[date, eventdate, origdate, pubstate:] The 17th edition of the
  \emph{Manual} continues to encourage writers to find some way of
  dating audiovisual materials, while if there is more than one date
  \enquote{the date of the original recording should be privileged}
  (15.57).  As with \textsf{music} entries, in order to follow these
  specifications I have had to provide three separate date fields for
  citing \textsf{video} sources, but their uses differ somewhat
  between the two types.  In both, the \textsf{date} will generally
  provide the publishing or copyright date of the medium you are
  referencing.  More specific to this entry type, the
  \textsf{origdate} will generally hold the date of the original
  theatrical release of a film, while the \textsf{eventdate} will most
  commonly present either the broadcast date of a particular TV
  program, or the recording/performance date of, for example, an opera
  on DVD.  The style will automatically prepend the bibstring
  \texttt{broadcast} to such a date, though you can use the
  \textsf{userd} field to change the string printed there.  (Absent an
  \textsf{eventdate}, the \textsf{userd} field in \textsf{video}
  entries will modify the \textsf{urldate}, and absent those two it
  will modify the \textsf{date}.)  Typically, any given \textsf{video}
  entry will only need an \textsf{eventdate} \emph{or} an
  \textsf{origdate}, and it is this date that will appear in citations
  and at the head of the entry in the reference list.  It's
  conceivable that you may need all three dates, in which case you can
  also use the standard \textsf{pubstate} field with \texttt{reprint}
  in it to control the printing of the \textsf{origdate} at the end of
  the entry, though I have altered the string that is printed there.
  Cf.\ friends:leia, handel:messiah, hitchcock:nbynw;
  \textsf{pubstate}, below.
\item[type:] As in all the audiovisual entry types, the \textsf{type}
  field holds the medium of the \textsf{title}, e.g., 8 mm, VHS, DVD,
  Blu-ray, MPEG.
\end{description}}

As with the \textsf{music} type, entries in \textsf{dates-test.bib}
should at least give you a good idea of how all this works.  (Cf.\
14.265, 14.267; \textsf{eventdate}, \textsf{origdate}, \textsf{userd};
\cmd{DeclareLabeldate} in section~\ref{sec:authformopts}, and
\texttt{avdate} in section~\ref{sec:authpreset}; cleese:holygrail,
friends:leia, handel:messiah, hitchcock:nbynw, loc:city,
weed:flatiron.)

\subsection{Entry Fields}
\label{sec:fields:authdate}

The following discussion presents, in alphabetical order, a complete
list of the entry fields you will need to use
\textsf{biblatex-chicago-authordate}.  As in
section~\ref{sec:types:authdate}, I shall include references to the
numbered paragraphs of the \emph{Chicago Manual of Style}, and also to
the entries in \textsf{dates-test.bib}.  Many fields are most easily
understood with reference to other, related fields.  In such cases,
cross references should allow you to find the information you need.

\mybigspace As \mymarginpar{\textbf{addendum}} in standard
\textsf{biblatex}, this field allows you to add miscellaneous
information to the end of an entry, after publication data but before
any \textsf{url} or \textsf{doi} field.  In the \textsf{patent} entry
type (which see), it will be printed in close association with the
filing and issue dates.  In any entry type, if your data begins with a
word that would ordinarily only be capitalized at the beginning of a
sentence, then simply ensure that that word is in lowercase, and the
style will take care of the rest.  Cf.\ \textsf{note}. (See
\emph{Manual} 14.114, 14.159--63; davenport:attention,
natrecoff:camera.)

\mybigspace In most \mymarginpar{\textbf{afterword}} circumstances,
this field will function as it does in standard \textsf{biblatex},
i.e., you should include here the author(s) of an afterword to a given
work.  The \emph{Manual} suggests that, as a general rule, the
afterword would need to be of significant importance in its own right
to require mentioning in the reference apparatus, but this is clearly
a matter for the user's judgment.  As in \textsf{biblatex}, if the
name given here exactly matches that of an editor and/or a translator,
then \textsf{biblatex-chicago} will concatenate these fields in the
formatted references.

\mylittlespace As noted above, however, this field has a special
meaning in the \textsf{suppbook} entry type, used to make an
afterword, foreword, or introduction the main focus of a citation.  If
it's an afterword at issue, simply define \textsf{afterword} any way
you please, leave \textsf{foreword} and \textsf{introduction}
undefined, and \textsf{biblatex-chicago} will do the rest. Cf.\
\textsf{foreword} and \textsf{introduction}. (See \emph{Manual}
14.105, 14.110; polakow:afterw.)

\mybigspace At \mymarginpar{\textbf{annotation}} the request of Emil
Salim, \textsf{biblatex-chicago} provides a package option (see
\texttt{annotation} below, section~\ref{sec:authuseropts}) to allow
you to produce annotated reference lists.  A recent feature request by
Moritz Wemheuer referenced a
\href{https://tex.stackexchange.com/questions/528374/moving-addendum-field-to-the-end-in-biblatex-chicago/540755#540755}{StackExchange}
question which suggested that the possible uses for the
\textsf{annotation} field could well be more extensive, appearing as
it does at the very end of all entry types.  I have therefore modified
the \texttt{annotation} option so that you can print the field in the
reference list (\texttt{=bib} or \texttt{=true}, the default), in long
(probably legal) notes (\texttt{=notes}), in both (\texttt{=all}), or
in neither (\texttt{=false}).  The two options \texttt{bibannotesep}
and \texttt{citeannotesep} allow you to choose the separator between
the rest of the entry and the \textsf{annotation}, and to choose a
different one in reference list and notes.  The default formatting in
the reference list (\texttt{vpar}) is to print the \textsf{annotation}
as a separate block using\ \verb+\par\nobreak\vskip\bibitemsep #1+,
while in long notes the default (\texttt{period}) is to print it
simply as an additional field, separated by a period.  The
\emph{Manual's} guidelines (14.64) allow for both these possibilities,
and I have provided a range of others, for which you should consult
the full documentation in section~\ref{sec:authpreset}.  (Note also
that both options can be set globally or per-type in the preamble, or
per-entry in the \textsf{options} field of individual entries.  For
specialized needs, of course, you can re-declare the format
[\verb+\DeclareFieldFormat{annotation}+] in your preamble, or redefine
the \cmd{bibannotesep} and \cmd{citeannotesep} commands there.)  In
section~\ref{sec:authuseropts} you will find instructions for
employing the new \texttt{formatbib} and \texttt{entrybreak} options
to give you fine-grained control over the formatting of the entire
reference list, particularly with regard to \TeX's page-breaking
algorithms.  The aim is to remove, in most cases, any need for you to
delve into the low-level commands involved in these algorithms.

\mybigspace I \mymarginpar{\textbf{annotator}} have implemented this
\textsf{biblatex} field pretty much as that package's standard styles
do, even though the \emph{Manual} doesn't actually mention it.  It may
be useful for some purposes.  Cf.\ \textsf{commentator}.

\paragraph*{\protect\mymarginpar{\textbf{author}}}
\label{sec:ad:author}
For the most part, I have implemented this field in a completely
standard fashion.  Remember that corporate or organizational authors
need to have an extra set of curly braces around them (e.g.,
\texttt{\{\{Associated Press\}\}}\,) to prevent \textsf{biber} from
treating one part of the name as a surname (14.84, 14.200, 15.37;
assocpress:gun, chicago:manual).  If there is no \textsf{author}, then
\textsf{biblatex-chicago} will look, in order, for a \textsf{namea},
an \textsf{editor}, a \textsf{nameb}, a \textsf{translator}, or a
\textsf{namec} (i.e., a compiler) and use that name (or those names)
instead, followed by the appropriate identifying string (esp.\ 15.36,
also 14.76, 14.103, 14.121, 14.126, 14.180; boxer:china, brown:bremer,
harley:cartography, schellinger:novel, sechzer:women, silver:gawain,
soltes:georgia).  \textsf{Biber} and \textsf{biblatex} take care of
alphabetizing entries no matter which name appears at their head.  In
citations, where the \textsf{labelname} is used, the order searched is
somewhat augmented: \textsf{shortauthor}, \textsf{author},
\textsf{shorteditor}, \textsf{namea}, \textsf{editor}, \textsf{nameb},
\textsf{translator}, and \textsf{namec}.

\mylittlespace If you wish to emphasize the activity of an editor, a
translator, or a compiler (14.104; eliot:pound), you can use the
\textsf{biblatex} options \texttt{useauthor=false},
\texttt{usenamea=false}, \texttt{use\-editor=false},
\texttt{usenameb=false}, \texttt{usetranslator=false}, and
\texttt{usenamec=false} in the \textsf{options} field to choose which
name appears at the head of an entry and in the citation.  You only
need to turn off any fields that are present in the entry, but please
remember to use the new option \texttt{usenamec} instead of the old
\texttt{usecompiler} (which I've deprecated), as the latter doesn't
work as smoothly and completely as \textsf{biblatex's} own name
toggles.  See \cmd{DeclareSortingTemplate} in
section~\ref{sec:authformopts}, and the \textsf{editortype}
documentation, below.

\mylittlespace Of course, in \textsf{collection} and
\textsf{proceedings} entry types, an \textsf{author} isn't expected,
so there the chain of substitutions begins with \textsf{namea} and
\textsf{editor}.  Also, in \textsf{article} and \textsf{review}
entries with \textsf{entrysubtype} \texttt{magazine}, the absence of
an \textsf{author} triggers the use of the \textsf{journaltitle} in
its stead.  Without an \textsf{entrysubtype}, the \textsf{title} will
be used.  See the discussion a few paragraphs down, and those entry
types, for further details.

\mylittlespace Recommendations concerning anonymous authors in other
kinds of references (15.34) emphasize using the \textsf{title} in
citations and at the head of reference list entries, rather than
\enquote{Anonymous.}  The latter may still in some cases be useful
\enquote{in a bibliography in which several anonymous works need to be
  grouped} (14.79), but even with a source like virginia:plantation,
\enquote{the reference list entry should normally begin with the
  title\ldots\ Text citations may refer to a short form of the title
  but must include the first word (other than an initial article)}
(15.34).  The \textsf{shorttitle} field is the place for the short
form, and you'll also need a \textsf{sortkey} of some sort if the full
title begins with an article that is to be ignored when alphabetizing.

\mylittlespace If \enquote{the authorship is known or guessed at but
  was omitted on the title page,} then you need to use the
\textsf{authortype} field to let \textsf{biblatex-chicago} know this
fact (15.34).  If the author is known (horsley:prosodies), then put
\texttt{anon} in the \textsf{authortype} field, if guessed at
(cook:sotweed) put \texttt{anon?}\ there.  (In both cases,
\textsf{biblatex-chicago} tests for these \emph{exact} strings, so
check your typing if it doesn't work.)  This will have the effect of
enclosing the name in square brackets, with or without the question
mark indicating doubt.  As long as you have the right string in the
\textsf{authortype} field, \textsf{biblatex-chicago-authordate} will
also do the right thing automatically in text citations.

\mylittlespace In most entry types (except \textsf{customc}) the
\textsf{nameaddon} field furnishes the means to cope with the case of
pseudonymous authorship.  If the author's real name isn't known,
simply put \texttt{pseud.}\ (or \verb+\bibstring{pseudonym}+) in that
field (centinel:letters).  If you wish to give a pseudonymous author's
real name, simply include it there, formatted as you wish it to
appear, as the contents of this field won't be manipulated as a name
by \textsf{biblatex} (lecarre:quest, stendhal:parma).  If you have
given the author's real name in the \textsf{author} field, then the
pseudonym goes in \textsf{nameaddon}, in the form \texttt{Firstname
  Lastname,\,pseud.}\ (creasey:ashe:blast, creasey:morton:hide,
creasey:york:death).  This latter method will allow you to keep all
references to one author's work under different pseudonyms grouped
together in the list of references, a method recommended by the
\emph{Manual}.  The \emph{Manual} (14.82) recommends using
cross-references from author to pseudonym or vice versa, so in these
latter examples I have included such references from the various
pseudonyms back to the author's name, using the \textsf{customc} entry
type, which see (ashe:creasey, morton:creasey, york:creasey).  Please
see the entry on \textbf{nameaddon}, below, for circumstances where
you may need to provide your own square brackets when presenting a
pseudonym, and also the package options \texttt{nameaddonformat} and
\texttt{nameaddonsep} in sections~\ref{sec:authuseropts} and
\ref{sec:authpreset}, below.

%%\enlargethispage{2\baselineskip}

\mylittlespace As its name suggests, the author-date style very much
wants to have a name of some sort present both for the entries in the
list of references and for the in-text citations.  The \emph{Manual}
is nothing if not flexible, however, so with unsigned articles or
encyclopedia entries the \textsf{journaltitle} or \textsf{title} may
take the place of the \textsf{author} (gourmet:052006,
lakeforester:pushcarts, nyt:trevorobit, unsigned:ranke,
wikipedia:bibtex).  Even in such entries, however, it may be
advantageous to provide either a standard \textsf{shorttitle} or, for
abbreviating a \textsf{journaltitle}, a \textsf{shortjournal} field,
thereby keeping the in-text citations to a reasonable length, though
not at the expense of making it hard to find the relevant entries in
the reference list.  An institutional author's name can also be rather
too long for in-text citations.  In unsigned:ranke I placed an
abbreviated form of the \textsf{journaltitle} into
\textsf{shortjournal}, adapting for a periodical the practice
recommended for books in 15.37.  In iso:electrodoc, I provided a
\textsf{shorthand} field, which by default in
\textsf{biblatex-chicago-authordate} will appear both in text
citations and at the head of the entry in the list of references,
followed, within the entry, by its expansion, this latter placed
within parentheses.  Please see under \textbf{shorthand} below for the
details.  (You can utilize the list of shorthands to clarify the
abbreviation, if you wish, and you can also provide a separate list of
journal abbreviations using the \cmd{printbiblist\{shortjournal\}}
command.  Please cf.\ the \textbf{shortjournal} documentation, below,
and the \texttt{journalabbrev} option in
section~\ref{sec:authpreset}.)

\mybigspace In \mymarginpar{\textbf{authortype}}
\textsf{biblatex-chicago}, this field serves a function very much in
keeping with the spirit of standard \textsf{biblatex}, if not with its
letter.  Instead of allowing you to change the string used to identify
an author, the field allows you to indicate when an author is
anonymous, that is, when their name doesn't appear on the title page
of the work you are citing.  As I've just detailed under
\textsf{author}, the \emph{Manual} generally discourages the use of
\enquote{Anonymous} (or \enquote{Anon.} as an author, though in some
cases it may well be your best option.  If, however, the name of the
author is known or guessed at, then you're supposed to enclose that
name within square brackets, which is exactly what
\textsf{biblatex-chicago} does for you when you put either
\texttt{anon} (author known) or \texttt{anon?} (author guessed at) in
the \textsf{authortype} field.  (Putting the square brackets in
yourself doesn't work right, hence this mechanism.)  The macros test
for these \emph{exact} strings, so check your typing if you don't see
the brackets.  Assuming the strings are correct,
\textsf{biblatex-chicago} will also automatically do the right thing
in citations.  (See the \textsf{author} docs just above.  Also
\emph{Manual} 15.34; cook:sotweed, horsley:prosodies.)

\mylittlespace The \emph{Manual} doesn't clarify how to treat multiple
works by the same \textsf{author}, in one or more of which their name
doesn't appear on the title page.  By default,
\textsf{biblatex-chicago} will, after the first appearance in the
reference list, replace identical \textsf{authors} with the 3-em dash,
regardless of any \textsf{authortype} field that may be present.  If
you want to distinguish between works certainly written by and works
merely ascribed to a given author, then you can use the
\texttt{dashed} option in the \textsf{options} field of individual
entries, and possibly also a \textsf{sortname}, to get the results you
want.


\mybigspace For \mymarginpar{\textbf{bookauthor}} the most part, as in
\textsf{biblatex}, a \textsf{bookauthor} is the author of a
\textsf{booktitle}, so that, for example, if one chapter in a book has
different authorship from the book as a whole, you can include that
fact in a reference (will:cohere).  Keep in mind, however, that the
entry type for introductions, forewords and afterwords
(\textsf{suppbook}) uses \textsf{bookauthor} as the author of
\textsf{title} (polakow:afterw, prose:intro).

\mybigspace This, \mymarginpar{\vspace{-8pt}\textbf{bookpagination}}
a standard \textsf{biblatex} field, allows you automatically to prefix
the appropriate string to information you provide in a \textsf{pages}
field.  If you leave it blank, the default is to print no identifying
string (the equivalent of setting it to \texttt{none}), as this is the
practice the \emph{Manual} recommends for nearly all page numbers.
Even if the numbers you cite aren't pages, but it is otherwise clear
from the context what they represent, you can still leave this blank.
If, however, you specifically need to identify what sort of unit the
\textsf{pages} field represents, then you can either hand-format that
field yourself, or use one of the provided bibstrings in the
\textsf{bookpagination} field.  These bibstrings currently are
\texttt{column,} \texttt{line,} \texttt{paragraph,} \texttt{page,}
\texttt{section,} and \texttt{verse}, all of which are used by
\textsf{biblatex's} standard styles.

\mylittlespace There are two points that may need explaining here.
First, all the bibstrings I have just listed follow the Chicago
specification, which may be confusing if they don't produce the
strings you expect.  Second, remember that \textsf{bookpagination}
applies only to the \textsf{pages} field --- if you need to format a
citation's \textsf{postnote} field, then you must use
\textsf{pagination}, which see (10.42--43, 14.147--56).

\mybigspace The \mymarginpar{\textbf{booksubtitle}} subtitle for a
\textsf{booktitle}.  See the next entry for further information.

\mybigspace In \mymarginpar{\textbf{booktitle}} the
\textsf{bookinbook}, \textsf{inbook}, \textsf{incollection},
\textsf{inproceedings}, and \textsf{letter} entry types, the
\textsf{booktitle} field holds the title of the larger volume in which
the \textsf{title} itself is contained as one part.  It is important
not to confuse this with the \textsf{maintitle}, which holds the more
general title of multiple volumes, e.g., \emph{Collected Works}.  It
is perfectly possible for one .bib file entry to contain all three
sorts of title (euripides:orestes, plato:republic:gr).  You may also
find a \textsf{booktitle} in other sorts of entries (e.g.,
\textsf{book} or \textsf{collection}), but there it will almost
invariably be providing information for the traditional
\textsc{Bib}\TeX\ cross-referencing apparatus (prairie:state), which I
discuss below (\textbf{crossref}).  Such provision is now unnecessary,
assuming you are using \textsf{biber}.  The \textsf{booktitle} no
longer takes sentence-style capitalization in \textsf{authordate},
though it does in \textsf{authordate-trad}.

\mybigspace An \mymarginpar{\textbf{booktitleaddon}} annex to the
\textsf{booktitle}.  It will be printed in the main text font, without
quotation marks.  If your data begins with a word that would
ordinarily only be capitalized at the beginning of a sentence, then
simply ensure that that word is in lowercase, and
\textsf{biblatex-chicago} will automatically do the right thing.  The
package and entry options \texttt{ptitleaddon} and
\texttt{ctitleaddon} (section~\ref{sec:authpreset}) allow you to
customize the punctuation that appears before the
\textsf{booktitleaddon} field.

\mybigspace This \mymarginpar{\textbf{chapter}} field holds the
chapter number, mainly useful only in an \textsf{inbook} or an
\textsf{incollection} entry where you wish to cite a specific chapter
of a book (ashbrook:brain).  It now also holds the track number of
individual pieces of \textsf{music}, whether on a traditional format
or on a streaming service (holiday:fool, rihanna:umbrella).

\mybigspace I \mymarginpar{\textbf{commentator}} have implemented this
\textsf{biblatex} field pretty much as that package's standard styles
do, even though the \emph{Manual} doesn't actually mention it.  It may
be useful for some purposes.  Cf.\ \textsf{annotator}.

\mybigspace \textsf{Biblatex} \mymarginpar{\textbf{crossref}} uses the
standard \textsc{Bib}\TeX\ cross-referencing mechanism, and has also
introduced a modified one of its own (\textsf{xref}).  The latter
works as it always has, attempting to remedy some of the deficiencies
of the traditional mechanism by ensuring that child entries will
inherit no data at all from their parents.  Section~2.4.1 of
\textsf{biblatex.pdf} contains useful notes on managing
cross-referenced entries, and section~3.15 explains some of the
limitations of the traditional backends, which offer only a small
subset of \textsf{Biber's} features.  For the \textsf{crossref} field,
when \textsf{Biber} is the backend, \textsf{biblatex} defines a series
of inheritance rules which make it much more convenient to use.
Appendix B of \textsf{biblatex.pdf} explains the defaults, to which
\textsf{biblatex-chicago} has added several that I should mention
here: \textsf{incollection} entries can inherit from \textsf{book} and
\textsf{mvbook} just as they do from \textsf{mvcollection} entries;
\textsf{letter} entries inherit from \textsf{book},
\textsf{collection}, \textsf{mvbook}, and \textsf{mvcollection}
entries the same way an \textsf{inbook} or an \textsf{incollection}
entry would; the \textsf{namea}, \textsf{nameb}, \textsf{sortname},
\textsf{sorttitle}, and \textsf{sortyear} fields, all highly
single-entry specific, are no longer inheritable; and \textsf{date}
and \textsf{origdate} fields are not inheritable from any of the
\textbf{mv*} entry types.

\mylittlespace Aside from these inheritance questions, the other main
function of the \textsf{crossref} and \textsf{xref} fields in
\textsf{biblatex-chicago} is as a trigger for the provision of
abbreviated entries in the list of references.  The \emph{Manual}
(15.42) specifies that if you cite several contributions to the same
collection, all (including the collection itself) may be listed
separately in the list, which the package does automatically, using
the default inclusion threshold of 2 in the case both of
\textsf{crossref}'ed and \textsf{xref}'ed entries.  (The familiar
\cmd{nocite} command may also help in some circumstances.)  In the
reference list an abbreviated form will be appropriate for all the
child entries.  The \textsf{biblatex-chicago-authordate} package has
always implemented these instructions, but only if you use a
\textsf{crossref} or an \textsf{xref} field, and only in
\textsf{incollection}, \textsf{inproceedings}, or \textsf{letter}
entries (on the last named, see just below).  Recent releases have
considerably extended this functionality.

\mylittlespace First, I added five entry types --- \textbf{book},
\textbf{bookinbook}, \textbf{collection}, \textbf{inbook}, and
\textbf{proceedings} --- to the list of those which use shortened
cross references, and I provided two options --- \texttt{longcrossref}
and \texttt{booklongxref}, on which more below --- which you can use
in the preamble or in the \textsf{options} field of an entry to enable
or disable the automatic provision of abbreviated references.  (The
\textsf{crossref} or \textsf{xref} field are still necessary for this
provision, but they are no longer sufficient on their own.)  The
\textsf{inbook} type works exactly like \textsf{incollection} or
\textsf{inproceedings}; in previous releases, you could use
\textsf{inbook} instead of \textsf{incollection} to avoid the
automatic abbreviation, the two types being otherwise identical.  Now
that you can use an option to turn off abbreviated references even in
the presence of a \textsf{crossref} or \textsf{xref} field, I have
thought it sensible to include this entry type alongside the others.
(Cf.\ ellet:galena, keating:dearborn, lippincott:chicago, and
prairie:state to see this mechanism in action in the reference list.)

\mylittlespace The inclusion of \textbf{book}, \textbf{bookinbook},
\textbf{collection}, and \textbf{proceedings} entries fulfills a
request made by Kenneth L. Pearce, and allows you to obtain shortened
references to, for example, separate volumes within a multi-volume
work, or to different book-length works collected inside a single
volume.  Such references are not part of the \emph{Manual's}
specification, but they are a logical extension of it, so the system
of options for turning on this functionality behaves differently for
these four entry types than for the other 4 (see below).  In
\textsf{dates-test.bib} you can get a feel for how this works by
looking at bernhard:boris, bernhard:ritter, bernhard:themacher,
harley:ancient:cart, harley:cartography, and harley:hoc.

\mylittlespace A published collection of letters requires a somewhat
different treatment (15.40).  In the author-date style, the
\emph{Manual} discourages individual letters from appearing in the
list of references at all, preferring that the \enquote{dates of
  individual correspondence should be woven into the text.}  If you
have special reason to do so, however, you can still present
individual published letters there (using the \textsf{letter} entry
type), and they too can use the system of shortened references just
outlined, even though the \emph{Manual} doesn't explicitly require it.
As with \textsf{book}, \textsf{bookinbook}, \textsf{collection},
\textsf{inbook}, \textsf{incollection}, \textsf{inproceedings}, and
\textsf{proceedings} entries, the use of a \textsf{crossref} or
\textsf{xref} field will activate this mechanism, assuming the 
preamble and entry options are set to enable it.  (See
white:ross:memo, white:russ, and white:total, for examples of the
\textsf{xref} field in action in this way, and please note that the
second of these entries is entirely fictitious, provided merely for
the sake of example.)

\mylittlespace These \mymarginpar{\texttt{longcrossref}} options
function, by default, asymmetrically.  The first,
\texttt{longcrossref}, generally controls the settings for the entry
types more-or-less authorized by the \emph{Manual}: \textsf{inbook},
\textsf{incollection}, \textsf{inproceedings}, and \textsf{letter}.

\begin{description}
\item[\qquad false:] This is the default.  If you use
  \textsf{crossref} or \textsf{xref} fields in the four mentioned
  entry types, you'll get the abbreviated entries in the reference
  list.
\item[\qquad true:] You'll get no abbreviated citations of these entry
  types in the reference list.
\item[\qquad none:] This switch is special, allowing you with one
  setting to provide abbreviated citations not just of the four entry
  types mentioned but also of \textsf{book}, \textsf{bookinbook},
  \textsf{collection}, and \textsf{proceedings} entries.
\item[\qquad notes, bib:] These two options are carried over from the
  notes \&\ bibliography style; here they are synonymous with
  \texttt{false} and \textsf{true}, respectively.
\end{description}

The \mymarginpar{\texttt{booklongxref}} second option,
\texttt{booklongxref}, controls the settings for \textsf{book},
\textsf{bookinbook}, \textsf{collection}, and \textsf{proceedings}
entries:

\begin{description}
\item[\qquad true:] This is the default.  If you use \textsf{crossref}
  or \textsf{xref} fields in these entry types, by default you will
  \emph{not} get any abbreviated citations in the reference list.
\item[\qquad false:] You'll get abbreviated citations in these entry
  types in the reference list.
\item[\qquad notes, bib:] These two options are carried over from the
  notes \&\ bibliography style; here they are synonymous with
  \texttt{false} and \textsf{true}, respectively.
\end{description}

Please note that you can set both of these options either in the
preamble or in the \textsf{options} field of individual entries,
allowing you to change the settings on an entry-by-entry basis.

\mylittlespace Please further note that in previous releases of
\textsf{biblatex-chicago} I recommended against using
\textsf{shorthand}, \textsf{reprinttitle} and/or \textsf{userf} fields
in combination with this abbreviated cross-referencing mechanism.  I
received, however a request from Alexandre Roberts to allow the
shorthand to appear in the place of the abbreviated cross-reference as
an additional space-saving measure, and one from Kenneth Pearce to
permit the combination of the other two fields with \textsf{crossref},
as well.  All three of these fields, in any combination, should just
work in such circumstances in \textsf{biblatex-chicago-authordate},
though if you are using a list of shorthands then you may need to
include \texttt{skipbiblist} in the \textsf{options} field of some
entries to avoid duplicates.  If you come across any problems or
inaccuracies, please report them.

\mylittlespace Finally, there is also an \texttt{xrefurl} option
available to control the printing of \textsf{url}, \textsf{doi}, and
\textsf{eprint} fields in abbreviated references where such
information might otherwise never appear.  See \texttt{xrefurl} in
section~\ref{sec:authuseropts}.

\paragraph*{\protect\mymarginpar{\textbf{date}}}
\label{sec:ad:date}
Predictably, this is one of the key fields for the author-date styles,
and one which, as a general rule, every .bib entry designed for this
system ought to contain.  So important is it, that
\textsf{biblatex-chicago-authordate} will, in most entry types, supply
a missing \cmd{bibstring\{nodate\}} if there is no date otherwise
provided (15.44), or if there is only a \textsf{urldate}, and that
date is an access date, i.e., there's no \textsf{userd} field (15.50).
Citations will look like (Author, n.d.), and entries in the list of
references will begin: Author, Firstname.\ n.d.  This seems simple
enough, but there are a surprising number of complications which
require attention.

\mylittlespace To start, in each entry, \textsf{Biber} attempts to
find something which it can designate a \textsf{labeldate}, which
will, in general and ideally, be the year printed both in citations
and at the head of the entry in the list of references.  The search
for the \textsf{labeldate} is governed by instances of the declaration
\cmd{DeclareLabeldate}, which cannot be set on an entry-by-entry
basis, but rather only in a document preamble (or in files used by
\textsf{biblatex} or its styles, like \textsf{biblatex-chicago}).  The
declaration can set a different search order according to entry type,
but other differentiations are not currently possible.  In all cases,
guided by the instructions given by the \cmd{DeclareLabeldate}
instances, \textsf{Biber} will search each entry in the declared
order, and the first match will provide the \textsf{labeldate}.  Only
when it finds no match at all will it fall back on
\verb+\bibstring{nodate}+.  (In the \textsf{misc} and \textsf{dataset}
types this automatic provision is turned off, as such material may not
be expected in many standard cases to have a usable date provided.)
You can prevent the appearance of \verb+\bibstring{nodate}+ throughout
your document in all entry types with the option
\texttt{nodates=false} when loading \textsf{biblatex-chicago} in your
preamble, or you can set it in the options field of individual
entries.  (See section~\ref{sec:authpreset}, below.)

\mylittlespace The thing to keep in mind is that \emph{only} for a
\textsf{labelyear} will \textsf{biblatex} provide what it calls the
\textsf{extradate} field, which means the alphabetical suffix
(1978\textbf{a}) to differentiate entries with the same author and
year.  A style can print any year it wants in a citation, but only the
\textsf{labelyear} comes equipped with an \textsf{extradate}.  (It is
also, by the way, the field that the sorting algorithm will use for
ordering the list of references.)  So the challenge, in a style
wherein entries can contain more than one date, is to allow different
dates to appear in citations and at the head of reference list
entries, but to ensure that, as often as is possible, that date
\emph{is} the \textsf{labeldate}.  This sounds simple, but in practice
it requires a series of options for date presentation, and multiple
iterations of the \cmd{DeclareLabeldate} command.  There are
\emph{two} standard search orders set up by default: in
\textsf{music}, \textsf{review}, \textsf{standard},
\textsf{suppperiodical}, and \textsf{video} entries, the default order
is \textsf{eventdate, origdate, date, urldate}, while in all other
entry types the default is \textsf{date, eventdate, origdate,
  urldate}.  I believe that these defaults work well for most
reference lists, especially those that contain relatively few entries
with multiple dates, but if they don't work for you then the following
options can help.

\mylittlespace In the case of \textsf{music}, \textsf{review},
\textsf{standard}, \textsf{suppperiodical}, and \textsf{video}
entries, the \texttt{avdate} (i.e., audio-visual date) option, set to
\texttt{true} by default, can be set to \texttt{false} in your
preamble to return these entry types to the general defaults.  Please
see the documentation of the entry types in
section~\ref{sec:types:authdate} above for the details of how multiple
dates will be treated in such entries, and also see \texttt{avdate} in
section~\ref{sec:authpreset}, below.  If you don't alter the
\texttt{avdate} settings, the other settings I am about to describe
won't apply to such entries.  For the entry types not covered by the
\texttt{avdate} option, the \emph{Manual} (15.40) presents a fairly
simple scheme for when a particular entry has more than one date, but
I have been unable to make its implementation quite as
straightforward.  If a reprinted book, say, has both a \textsf{date}
of publication for the reprint edition and an \textsf{origdate} for
the original edition, then by default
\textsf{biblatex-chicago-authordate} will use the \textsf{date} in
citations and at the head of the entry in the reference list.  If you
inform \textsf{biblatex-chicago} that the book is a reprint by putting
the string \texttt{reprint} in the \textsf{pubstate} field, then a
notice will be printed at the end of the entry saying \enquote{First
  published 1898.}  With no \textsf{pubstate} field (and no
\texttt{cmsdate} option), the algorithms will ignore the
\textsf{origdate}.

\mylittlespace If, \mymarginpar{\texttt{cmsdate}\\\emph{in entry}}
for any reason, you wish the \textsf{origdate} to appear at the head
of the entry, then your first option is to use the \texttt{cmsdate}
toggle in the \textsf{options} field of the entry itself.  This has 3
possible states relevant to this context, though there is a fourth
state (\texttt{full}) which I shall discuss below:

\begin{enumerate}
\item \texttt{cmsdate=both} prints both the \textsf{origdate} and the
  \textsf{date}, using the \emph{Manual's}\ standard format: (Author
  [1898] 1952) in parenthetical citations, Author (1898) 1952 outside
  parentheses, e.g., in the reference list.
\item \texttt{cmsdate=off} is the default, discussed above:
  (Author 1952).
\item \texttt{cmsdate=on} prints the \textsf{origdate} at the head of
  the entry in the list of references and in citations: (Author 1898).
  NB: The \emph{Manual} no longer includes this among the approved
  options.  If you want to present the \textsf{origdate} at the head
  of an entry, then generally speaking you should probably use
  \texttt{cmsdate=both}.  I have nevertheless retained this option for
  certain cases where it has proved useful.  The obsolete options
  \texttt{new} and \texttt{old} work like \texttt{both}.
\end{enumerate}

In the first and third cases, if you put the string \texttt{reprint}
in the \textsf{pubstate} field, then the publication data in the list
of references will include a notice, formatted according to the
specifications, that the modern edition is a reprint.  In the third
case, since the \textsf{date} hasn't yet been printed, this
publication data will also include the date of the modern reprint.

\mylittlespace Let us imagine, however, that your list of references
contains another book by the same author, also a reprint edition:
(Author [1896] 1974).  How will these two works be ordered in the list
of references?  By the \textsf{labelyear}, in this case the
\textsf{year} field, which appears first in the default definition
(\textsf{date, eventdate, origdate, urldate}) of
\cmd{DeclareLabeldate}, and which in this case will be wrong, because
the entries should always be ordered by the \emph{first} date to
appear there, in this case the contents of \textsf{origdate}.  In this
example, the solution can be as simple as a \textsf{sortyear} field
set to something earlier than the date of the other work, e.g.,
\texttt{1951}.

\mylittlespace And if the reprint dates --- in the \textsf{date} field
--- of the two works were the same?  Just as when it is ordering
entries, \textsf{biblatex} will always first process the contents of
the \textsf{labelyear} field when it is deciding whether to add the
\textsf{extradate} alphabetical suffix (\texttt{a,b,c} etc.)\ to the
year to distinguish different works by the same author published in
the same year.  Our current hypothetical examples would look like
this: ([1896] 1974a) and ([1898] 1974b), with the suffixes
unnecessary, strictly-speaking, either for ordering or for
disambiguating the entries.  If the original publication dates --- in
the \textsf{origdate} field --- are the same, and the reprint dates
different, you may prefer citations of the two works to read, e.g.,
(Author [1898a] 1952) and (Author [1898b] 1974), when they in fact
read (Author [1898] 1952) and (Author [1898] 1974).  These latter
forms aren't ambiguous, and even if the reprints themselves appeared
in the same year then the alphabetical suffix would appear attached to
the \textsf{date} --- (Author [1898] 1974a) and (Author [1898] 1974b)
--- again avoiding ambiguity.

\mylittlespace The \emph{Manual} doesn't give clear instructions for
how to cope with these situations, but
\textsf{biblatex-chicago-authordate} provides help.  You can't
manually put the alphabetical suffix on an \textsf{origdate} yourself
because that field only accepts numerical data.  Instead, we can
choose between two solutions.  The old way
\mymarginpar{\texttt{cmsdate}\\\emph{in entry} \\+
  \texttt{switchdates}} is an unusual expedient, which amounts to
switching the two date fields, placing the earlier date in
\textsf{date} and the later one in \textsf{origdate}.  The style tests
for this condition using a simple arithmetical comparison between the
two years, then prints the two dates according to the state of the
\texttt{cmsdate} toggle.  The three relevant states of this toggle are
the same as before, but there are only two possible outcomes, as
follows:

\begin{enumerate}
\item \texttt{cmsdate=off} (the default) and \texttt{cmsdate=on}
  \emph{both} print the \textsf{date} at the head of the entry in the
  list of references and in citations: (Author 1898a), (Author 1898b).
  As noted above, this style is no longer recommended by the
  \emph{Manual}, but may still be useful in some cases.
\item \texttt{cmsdate=both} prints both the \textsf{date} and the
  \textsf{origdate}, using the \emph{Manual's}\ preferred format:
  (Author [1898a] 1952), (Author [1898b] 1974).  The obsolete options
  \texttt{old} and \texttt{new} are synonyms for this.
\end{enumerate}

If, for some reason, the automatic switching of the dates cannot be
achieved, perhaps in crossref'd \textsf{letter} entries that you
really want to have in your list of references (white:ross:\break
memo, white:russ), or perhaps in a reprint edition that hasn't yet
appeared in print (preventing the comparison between a year and the
word \enquote{forthcoming}), then you can use the per-entry option
\texttt{switchdates} in the \textsf{options} field to achieve the
required effects.

\mylittlespace The more \mymarginpar{\texttt{cmsdate}\\\emph{in
    preamble}} drastic method of simplifying the creation of databases
with a great many multi-date entries is to use the \texttt{cmsdate}
option \emph{in the preamble}.  Despite warnings in previous releases,
users had plainly already been setting this option in their preambles,
so I thought I might at least attempt to make it work as
\enquote{correctly} as I can.  The switches for this option are the
same as for the entry-only option, that is:

\begin{enumerate}
\item \texttt{cmsdate=off} is the default: (Author 1952).
\item \texttt{cmsdate=both} prints both the \textsf{origdate} and the
  \textsf{date}, using the \emph{Manual's}\ standard format: (Author
  [1898] 1952) in parenthetical citations, Author (1898) 1952 outside
  parentheses, e.g., in the reference list.
\item \texttt{cmsdate=on} prints the \textsf{origdate} at the head of
  the entry in the list of references and in citations: (Author 1898).
  NB: The \emph{Manual} no longer includes this among the approved
  options.  If you want to present the \textsf{origdate} at the head
  of an entry, then generally speaking you should probably use
  \texttt{cmsdate=both}.  I have nevertheless retained this option for
  certain cases where it might proved useful.  The obsolete options
  \texttt{new} and \texttt{old} work like \texttt{both}.
\end{enumerate}

The important change for the user is that, when you set this option in
your preamble to \texttt{on} or \texttt{both} (or to the obsolete
synonyms for the latter, \texttt{new} or \texttt{old}), then
\textsf{biblatex-chicago-authordate} will change the default
\cmd{DeclareLabeldate} definition so that the \textsf{labelyear}
search order will be \textsf{origdate, date, eventdate, urldate}.
This means that for entry types not covered by the \texttt{avdate}
option, and for those types as well if you turn off that option, the
\textsf{labelyear} will, in any entry containing an \textsf{origdate},
be that very date.  If you want \emph{every} such entry to present its
\textsf{origdate} in citations and at the head of reference list
entries, then setting the option this way makes sense, as you should
automatically get the proper \textsf{extradate} and the correct
sorting, without having to switch dates around counter-intuitively in
your .bib file.  A few clarifications may yet be in order.

\mylittlespace Obviously, any entry with only a \textsf{date} should
behave as usual.  Also, since \textsf{patent} entries have fairly
specialized needs, I have exempted them from this change to
\cmd{DeclareLabeldate}.  Third, the per-entry \texttt{cmsdate} options
will still affect which dates are printed in citations and at the head
of reference list entries, but they cannot change the search order for
the \textsf{labeldate}.  This will be fixed by the preamble option.
Fourth, if you have been used to switching the \textsf{date} and the
\textsf{origdate} to get the correct results, then you should be aware
that this mechanism may actually still be useful when using the
\texttt{on} switch to \texttt{cmsdate} in the preamble, but it
produces incorrect results when the \texttt{cmsdate} option is
\texttt{both} in the preamble and the individual entry.  The preamble
option is designed to make the need for this switching as rare as
possible, so some editing of existing databases may be necessary.

\mylittlespace Finally, Bertold Schweitzer has brought to my attention
certain difficult corner cases involving cross-referenced works with
more than one date.  In order to facilitate the accurate presentation
of such sources, I made a slight change to the way the entry-only
\texttt{cmsdate=on} and \texttt{cmsdate=both} work.  If, and only if,
a work has only one date, and there is no \texttt{switchdates} in the
\textsf{options} field, then \texttt{cmsdate=on} and
\texttt{cmsdate=both} will both result in the suppression of the
\textsf{extradate} field in that entry, that is, the year will no
longer be printed with its following lowercase letter used to
distinguish works by the same \textsf{author} published in the same
year.  Obviously, if the same options are set in the preamble, this
behavior is turned off, so that single-date entries will still work
properly without manual intervention.

\mylittlespace Up\mymarginpar{\textsc{iso}8601-2\\Extended\\Format}
to this point, the discussion of the \textsf{date} field has in fact
presented no substantive alterations to the way it behaved in previous
releases of \textsf{biblatex-chicago.}  With this release, however, I
have implemented all of the applicable parts of \textsf{biblatex's}
elegant, and long standing, support for the \textsc{iso}8601-2
Extended Format specification, which means the package now provides
greatly enhanced possibilities for presenting uncertain and
unspecified dates and date ranges, along with date eras, seasons, and
time stamps.  I have also implemented the \emph{Manual's} (9.64)
guidelines for compressing year ranges, as well as providing a few
more extras to help with some of the other tricky corners of the
\emph{Manual's} instructions.  A combination of \textsf{biblatex} and
\textsf{biblatex-chicago} package options allows you to define when,
how, and where any of these extended specifications will appear in
your documents.  I have attempted to provide as compliant a set of
defaults as possible in \textsf{biblatex-chicago.sty}, but you can
alter any of them according to your needs.  All are documented in
section~\ref{sec:opts:authdate}, below, but table~\ref{ad:date:extras}
purports to serve as a convenient reference guide to how this all
works.

\afterpage{\clearpage

\begin{table}[h!]
\hspace*{-2.5em}
\begin{threeparttable}
  \caption[\hspace{-1em}Enhanced Date Specifications]
  {Enhanced Date Specifications in biblatex-chicago}
\begin{tabularx}{140mm}{@{}>{\ttfamily}llX@{}}
\toprule
\multicolumn{1}{@{}H}{Date Specification} &
\multicolumn{2}{H}{Formatted Date (Examples use \texttt{american}
  localization)\tnote{\textmd{a}}} \\
\cmidrule(l){2-3}
&
\multicolumn{1}{H}{Output Format} &
\multicolumn{1}{H}{Output Format Notes} \\
\cmidrule{1-1}\cmidrule(l){2-2}\cmidrule(l){3-3}
1723? & [1723?]\tnote{b} & \texttt{dateuncertain=true} set by default\\
1723\textasciitilde & [ca. 1723]\tnote{b} & \texttt{datecirca=true} set by
default\\
1723\% & [ca. 1723?]\tnote{b} & Both \textsf{biblatex} options \texttt{true} by
default, as above\\\addlinespace[.8mm]
\cmidrule(r{2em}){1-1}\addlinespace[.8mm]
2016-05-24T15:34:00 & May 2, 2016, 3:34 p.m.\tnote{c} &
\texttt{alltimes=12h} set by default \\
2016-05-24T15:34:00 & May 2, 2016, 15:34\tnote{c} &
\texttt{urltime=24h} set by default \\\addlinespace[.8mm]
\cmidrule(r{2em}){1-1}\addlinespace[.8mm]
-0876 & 877 BC & \texttt{dateera=christian} set in your preamble\\
-0876/-0866 & 877--867 BC & \\
0876 & AD 876\tnote{d} & \texttt{dateeraauto=1000} also in preamble\\
-0876/0866 & 877 BC--AD 866 & \\
0866/0876 & AD 866--76 & Cf.\ \texttt{compressyears}, below\\
0343-02-03 & February 3, AD 343 & \\\addlinespace[.8mm]
\cmidrule(r{2em}){1-1}\addlinespace[.8mm]
-0876 & 877 BCE & \texttt{dateera=secular}, \texttt{dateeraauto=1000}\\
-0876/-0866 & 877--867 BCE & \\
0876 & 876 CE &  \\
-0876/0866 & 877 BCE--866 CE & \\
0866/0876 & 866--76 CE & Cf.\ \texttt{compressyears}, below\\
0343-02-03 & February 3, 343 CE & \\\addlinespace[.8mm]
\cmidrule(r{2em}){1-1}\addlinespace[.8mm]
195X & 1950s\tnote{e,f} & Chicago option
\texttt{decaderange=true} gets you 1950--59 \\
19XX & 20th c.\tnote{f} & Chicago option
\texttt{centuryrange=true} gets you 1900--1999;
\texttt{alwaysrange=true} does the same for this and the
previous entry \\\addlinespace[.8mm]
\cmidrule(r{2em}){1-1}\addlinespace[.8mm]
2004-22 & Summer 2004 &  \\
1908/1912 & 1908--12\tnote{g} & Chicago option
\texttt{compressyears=true} set by default \\
\bottomrule
\end{tabularx}
  \def\TPTnoteSettings{%
    \setlength\leftmargin{1.5em}%
    \setlength\labelwidth{.5em}%
    \setlength\labelsep{.2em}%
    \footnotesize
    \rightskip\tabcolsep \leftskip\tabcolsep}%
\begin{tablenotes}
\item[a] In other languages both the strings and their placement with
  respect to the year can and will differ.
  \item[b] The two Chicago options \texttt{nodatebrackets} and
    \texttt{noyearbrackets} can remove the brackets around the year in
    this context, though please note that they work quite differently
    in the notes \&\ bibliography and author-date styles. Please see
    their documentation in sections~\ref{sec:useropts} and
    \ref{sec:authuseropts}, respectively.
  \item[c] Any time stamp that is part of a \textsf{urldate} will
    appear in any entry type, though you can prevent this by setting
    \texttt{urlstamp=false}. Time stamps in \textsf{date} and
    \textsf{origdate} fields will appear only in \textsf{article} and
    \textsf{periodical} entries with a \texttt{magazine}
    \textsf{entrysubtype}, as well as in all \textsf{online},
    \textsf{review}, and \textsf{suppperiodical} entries. Such data in
    \textsf{eventdate} fields will appear only in \textsf{review} and
    \textsf{suppperiodical} entries. For timezones the four
    \textsf{timezone} fields allow you to present Chicago's preferred
    initialisms (\enquote{PST}). The \texttt{timezones} option is
    \texttt{true} by default. By contrast, the \texttt{seconds} option
    is not set by default, though you still need to include the
    seconds, as above, for \textsf{biber} to process the time stamp.
\item[d] The \texttt{annodomini} string appears before the year
  \emph{only} in documents in some variant of English.
\item[e] When the decade string would be ambiguous ---
  \enquote{1900s} --- the style prints \enquote{1900--1909} instead.
\item[f] For decades and centuries, the current state of the
  \textsf{biber} code cannot process dates BCE / BC.
\item[g] The Chicago rules for year-range compression differ from its
  rules for page-range compression (9.61 \&\ 9.64). Dates BCE~/ BC are
  never compressed. You must have loaded \textsf{biblatex-chicago.sty}
  for the compression code to be available.
\end{tablenotes}
\vspace*{2em}
\label{ad:date:extras}
\end{threeparttable}
\end{table}}

\mylittlespace There are several more general remarks about the
\textsf{date} field that may be helpful to users of the author-date
styles.  First, I highly recommend familiarizing yourself with the
extended date specifications, as in many cases they will greatly
simplify the creation of your .bib databases.  A \textsf{date} field
like \verb+{1957?}+ in clark:mesopot not only lets the package provide
the appropriate square brackets for you, it also means that the
\textsf{year} field in the .bbl file sorts just as it should, and can
be tested numerically for its relation to other date years in the
entry.  A \textsf{year} field like \verb+{[1957?]}+ in the .bib
database produces a field in the .bbl that neither sorts nor can be
numerically tested.  The same holds for a compressed year range, as in
tillich:system.  With \texttt{compressyears} set to \texttt{true} by
default, a \textsf{date} field like \verb+{1951/1963}+ lets the
package decide what compression is correct, and provides \textsf{year}
and \textsf{endyear} fields that sort and compare numerically for both
\textsf{switchdates} and \textsf{extradate} tests.  Clearly,
situations may still arise when a specially-crafted \textsf{year} or
\textsf{origyear} field may be necessary, but if you can use the
enhanced specifications then I strongly advocate doing so.

\mylittlespace One\mymarginpar{\textsf{verbc}} possible drawback is
that using these facilities makes a great many more dates available to
the \textsf{extradate} mechanism which, it turns out, is something of
a mixed blessing.  The \emph{Manual} isn't entirely forthcoming about
what to do in the (vanishingly rare) case that two works by the same
\textsf{author} have the same uncertain \textsf{date}.
\textsf{Biblatex-chicago} will print the \textsf{extradate} in such
situations, so that you could have \texttt{[1957?]a} followed by
\texttt{[ca.~1957]b}, which may not be exactly right, nor exactly what
you want.  Here, the new \textsf{verbc} field can help --- giving the
two entries different values of this field will prevent the
\textsf{extradate} from appearing.  Please see the documentation of
the \textsf{verbc} field below for all the details.

\mylittlespace Second, for most entry types, only a year is really
necessary, and in most situations only the year --- or year range ---
will be printed in text citations and at the head of entries in the
list of references.  More specific \textsf{date} fields are often
present, however, in an unpredictably broad range of entries.  In a
change to previous practice, a \textsf{date} with a year, month, and
day will, even if the year appears at the head of the entry, be
repeated in full later in the entry, while if there's only a month (or
a season) and a year the month (or season) alone will follow.  Also
new is the presentation of time stamps, which you can easily provide
in your \textsf{date} fields (see examples and usage notes in
table~\ref{ad:date:extras}).  These finer-grained specifications are
really only necessary for news stories that are frequently updated
\enquote{as they unfold} (14.191), or for online posts, particularly
comments, that may need a time stamp for disambiguation.  If you wish
to specify the time zone, the \emph{Manual} (10.41) prefers
initialisms like \enquote{EST} or \enquote{PDT,} and these are most
easily provided using the \texttt{timezone} field, where you can
include your own parentheses if so desired (cp.\ 14.191).  If you
follow the recommendations of the \emph{Manual} and present newspaper
and magazine articles \enquote{entirely within the text} (15.49), then
the citations need to contain the complete \textsf{date} (and possible
time stamp) along with the \textsf{journaltitle}.  Placing
\mymarginpar{\texttt{cmsdate=full}} \texttt{cmsdate=full} (and
\texttt{skipbib}) in the \textsf{options} field of an \textsf{article}
or a \textsf{review} entry, alongside a possible
\texttt{useauthor=false}, should allow you to achieve this.  For
online comments in \textsf{online} or \textsf{review} entries this
presentation is the default when you use the new \texttt{commenton}
\textsf{relatedtype}.  (See the documentation of those two types in
section~\ref{sec:types:authdate}, as well as \textsf{relatedtype} in
section~\ref{sec:authrelated}.)  While we're on this subject, the
\emph{Manual} is flexible (in both specifications) on abbreviating the
names of months (14.171).  By default,
\textsf{biblatex-chicago-authordate} uses the full names, which you
can change by setting the option \texttt{dateabbrev=true} in the
preamble.  (Cf.\ ac:comment, assocpress:gun, barcott:review, batson,
creel:house, friends:leia, holiday:fool, nass:address,
petroff:impurity, powell:email.)

\mylittlespace Third, in the \textsf{misc} entry type the
\textsf{date} field can help to distinguish between two classes of
archival material, letters and \enquote{letter-like} sources using
\textsf{origdate} while others (interviews, wills, contracts) use
\textsf{date}.  (See \textsf{misc} in section~\ref{sec:types:authdate}
for the details.)  If such an entry, as may well occur, contains only
an \textsf{origdate}, as can also be the case in \textsf{letter}
entries, then \textsf{Biber} and either \cmd{DeclareLabeldate}
definition will make it work without further intervention.  Fourth,
you can in most entry types qualify a \textsf{date} with the
\textsf{userd} field, assuming that the entry contains no
\textsf{urldate}.  For \textsf{music} and \textsf{video} entries,
there are several other requirements --- please see the documentation
of \textsf{userd}, below.  Fifth, and finally, please note that the
\textsf{nameaddon} field, which see, is no longer the place for time
stamps, as it was in the 16th-edition styles.  Any such data there
should be moved into the corresponding date field (either the
\textsf{date} or the \textsf{eventdate}, typically).

\mylittlespace I recommend that you have a look at
\textsf{dates-test.bib} to see how all these complications will affect
the construction of your .bib database, especially at the following
entries: aristotle:metaphy:gr, creel:house, emerson:nature,
james:ambassadors, maitland:canon, maitland:equity, schweitzer:bach,
spock:interview, white:ross:memo, and white:russ.  Cf.\ also
\textsf{origdate}, \textsf{timezone}, \textsf{verbc}, and
\textsf{year}, below; the \texttt{alldates}, \texttt{alltimes},
\texttt{alwaysrange}, \texttt{centuryrange}, \texttt{cmsdate},
\texttt{compressyears}, \texttt{datecirca}, \texttt{dateera},
\texttt{dateeraauto}, \texttt{dateuncertain}, \texttt{decaderange},
\texttt{nodatebrackets}, \texttt{nodates}, \texttt{noyearbrackets},
\texttt{switchdates}, \texttt{timezones}, \texttt{urlstamp}, and
\texttt{urltime} options in sections~\ref{sec:preset:authdate},
\ref{sec:authuseropts}, and \ref{sec:authentryopts}; and
section~4.5.10 in \textsf{biblatex.pdf}, and section
\ref{sec:authformopts}, below, for the \cmd{DeclareLabeldate} command.

\mybigspace This \mymarginpar{\textbf{day}} field, as of
\textsf{biblatex} 0.9, is obsolete, and will be ignored if you use it
in your .bib files.  Use \textsf{date} instead.

\mybigspace Standard \mymarginpar{\textbf{doi}} \textsf{biblatex}
field, providing the Digital Object Identifier of the work.  The
\emph{Manual} specifies that, given their relative permanence compared
to URLs, \enquote{authors should prefer a DOI- or Handle-based URL
  whenever one is available} (14.8).  (14.175; friedman:learning).
Cf.\ \textsf{url}.

\mybigspace Standard \mymarginpar{\textbf{edition}} \textsf{biblatex}
field.  If you enter a plain cardinal number, \textsf{biblatex} will
convert it to an ordinal (chicago:manual), followed by the appropriate
string.  Any other sort of edition information will be printed as is,
though if your data begins with a word (or abbreviation) that would
ordinarily only be capitalized at the beginning of a sentence, then
simply ensure that that word (or abbreviation) is in lowercase, and
\textsf{biblatex-chicago} will automatically do the right thing
(babb:peru, times:guide).  In most situations, the \emph{Manual}
generally recommends the use of abbreviations in the list of
references, but there is room for the user's discretion in specific
citations (emerson:nature).

\mybigspace As \mymarginpar{\textbf{editor}} far as possible, I have
implemented this field as \textsf{biblatex}'s standard styles do, but
the requirements specified by the \emph{Manual} present certain
complications that need explaining.  \textsf{Biblatex.pdf} points out
that the \textsf{editor} field will be associated with a
\textsf{title}, a \textsf{booktitle}, or a \textsf{maintitle},
depending on the sort of entry.  More specifically,
\textsf{biblatex-chicago} associates the \textsf{editor} with the most
comprehensive of those titles, that is, \textsf{maintitle} if there is
one, otherwise \textsf{booktitle}, otherwise \textsf{title}, if the
other two are lacking.  In a large number of cases, this is exactly
the correct behavior (adorno:benj, centinel:letters,
plato:republic:gr, among others).  Predictably, however, there are
numerous cases that require, for example, an additional editor for one
part of a collection or for one volume of a multi-volume work.  For
these cases I have provided the \textsf{namea} field.  You should
format names for this field as you would for \textsf{author} or
\textsf{editor}, and these names will always be associated with the
\textsf{title} (donne:var).

\mylittlespace As you will see below, I have also provided a
\textsf{nameb} field, which holds the translator of a given
\textsf{title} (euripides:orestes).  If \textsf{namea} and
\textsf{nameb} are the same, \textsf{biblatex-chicago} will
concatenate them, just as \textsf{biblatex} already does for
\textsf{editor}, \textsf{translator}, and \textsf{namec} (i.e., the
compiler).  Furthermore, it is conceivable that a given entry will
need separate editors for each of the three sorts of title.  For this,
and for various other tricky situations, there is the \cmd{partedit}
macro (and its siblings), designed to be used in a \textsf{note}
field, in one of the \textsf{titleaddon} fields, or even in a
\textsf{number} field (howell:marriage).  (Because the strings
identifying an editor differ in notes and bibliography, one can't
simply write them out in such a field when using the notes \&\
bibliography style, but you can certainly do so in the author-date
styles, if you wish.  Using the macros will make your .bib file more
portable across both Chicago specifications, and also across multiple
languages, but they are otherwise unnecessary.
Cf. section~\ref{sec:international}, and also \textsf{namea},
\textsf{nameb}, \textsf{namec}, and \textsf{translator}.)

\mybigspace \textsf{Biblatex}
\mymarginpar{\textbf{editora\\editorb\\editorc}} provides these fields
as a means to specify additional contributors to texts in a number of
editorial roles.  In the Chicago styles they seem most relevant for
the audiovisual types, especially \textsf{music} and \textsf{video},
and now also the \textsf{performance} type, in all of which they help
to identify conductors, directors, producers, and performers.  To
specify the role, use the fields \textsf{editoratype},
\textsf{editorbtype}, and \textsf{editorctype}, which see.  (Cf.\
bernstein:shostakovich, hamilton:miranda, handel:messiah.)

\mybigspace Normally, \mymarginpar{\textbf{editortype}} with the
exception of the \textsf{article} and \textsf{review} types with a
\texttt{magazine} \textsf{entrysubtype},
\textsf{biblatex-chicago-authordate} will automatically find a name to
put at the head of an entry, starting with an \textsf{author}, and
proceeding in order through \textsf{namea}, \textsf{editor},
\textsf{nameb}, \textsf{translator}, and \textsf{namec} (the
compiler).  If all six are missing, then the \textsf{title} will be
placed at the head.  (In \textsf{article} and \textsf{review} entries
with a \texttt{magazine} \textsf{entrysubtype}, a missing
\textsf{author} immediately prompts the use of \textsf{journaltitle}
at the head of an entry.  See above under \textsf{article} for
details.)  The \textsf{editortype} field provides even greater
flexibility, allowing you to choose from a variety of editorial roles
while only using the \textsf{editor} field.  You can do this even
though an author is named (eliot:pound shows this mechanism in action
for a standard editor, rather than for some other role).  Two things
are necessary for this to happen.  First, in the \textsf{options}
field you need to set \texttt{useauthor=false} (if there is an
\textsf{author}), then you need to put the name you wish to see at the
head of your entry into the \textsf{editor} or the \textsf{namea}
field.  If the \enquote{\textsf{editor}} is in fact, e.g., a compiler,
then you need to put \texttt{compiler} into the \textsf{editortype}
field, and \textsf{biblatex} will print the correct string after the
name in the list of references.

\mylittlespace In previous releases of \textsf{biblatex-chicago} you
could only use defined \cmd{bibstrings} in this field, at least if you
wanted anything printed.  N.~Andrew Walsh pointed out that the
standard \textsf{biblatex} styles will just print the field as-is in
this case, allowing them to handle a great many unforeseen editorial
roles with comparative ease, so I've implemented this, too, making
sure to capitalize the string if the context demands it.  The string
you choose will differ depending on whether it will be printed after a
name at the head of an entry or before a name later on in the entry,
e.g., \enquote{cartographer} or \enquote{maps created by.}  A bit of
trial and error should see you through.

\mylittlespace There are a few details of which you need to be aware.
Because \textsf{biblatex-chicago} has added the \textsf{namea} field,
which gives you the ability to identify the editor specifically of a
\textsf{title} as opposed to a \textsf{maintitle} or a
\textsf{booktitle}, the \textsf{editortype} mechanism checks first to
see whether a \textsf{namea} is defined.  If it is, that name will be
used at the head of the entry, if it isn't , or if you've set the
option \texttt{usenamea=false}, it will go ahead and look for an
\textsf{editor}.  The \textsf{editortype} field applies only to the
\textsf{editor}, but you can use \textsf{nameatype} to modify
\textsf{namea}.  \textsf{Biblatex}'s sorting algorithms, and also its
\textsf{labelname} mechanism, should both work properly no matter sort
of name you provide, thanks to \textsf{Biber} and the (default)
Chicago-specific definitions of \cmd{DeclareLabelname} and
\cmd{DeclareSortingTemplate}.  (Cf.\ section~\ref{sec:authformopts},
below).  Please be aware that if you want a shortened form to appear
in citations then there's only the \textsf{shorteditor}, which you
should ensure presents whichever of the two editors' names
(\textsf{namea} or \textsf{editor}) appears at the head of the
reference-list entry.

% %\enlargethispage{\baselineskip}

\mylittlespace In \textsf{biblatex} 0.9 Lehman reworked the string
concatenation mechanism, for reasons he outlines in his RELEASE file,
and I have followed his lead.  In short, if you define the
\textsf{editortype} field, then concatenation is turned off, even if
the name of the \textsf{editor} matches, for example, that of the
\textsf{translator}.  In the absence of an \textsf{editortype} (or
\textsf{nameatype}), the usual mechanisms remain in place, that is, if
the \textsf{editor} exactly matches a \textsf{translator} and/or a
\textsf{namec}, or alternatively if \textsf{namea} exactly matches a
\textsf{nameb} and/or a \textsf{namec}, then \textsf{biblatex} will
print the appropriate strings.  The \emph{Manual} specifically (15.7)
recommends not using these identifying strings in citations, and
\textsf{biblatex-chicago-authordate} follows that recommendation.  If
you nevertheless need to provide such a string, you'll have to do it
manually in the \textsf{shorteditor} field, or perhaps, in a different
sort of entry, in a \textsf{shortauthor} field.

\mylittlespace It may also be worth noting that because of certain
requirements in the specification -- absence of an \textsf{author},
for example -- the \texttt{useauthor=false} mechanism is either
unnecessary or won't work properly in the following entry types:
\textsf{collection}, \textsf{letter}, \textsf{patent},
\textsf{periodical}, \textsf{proceedings}, \textsf{review},
\textsf{suppbook}, \textsf{suppcollection}, and
\textsf{suppperiodical}.

\mybigspace These
\mymarginpar{\textbf{editoratype\\editorbtype\\editorctype}} fields
identify the exact role of the person named in the corresponding
\textsf{editor[a-c]} field, just as \textsf{editortype} (q.v.) does
for the \textsf{editor}.  Note that they are not part of the string
concatenation mechanism.  I have implemented them just as the standard
styles do, that is, if the field isn't a pre-defined \cmd{bibstring}
it will be printed as-is, contextually capitalized.  They have found a
use particularly in \textsf{music}, \textsf{performance}, and
\textsf{video} entries.  Cf.\ bernstein:shostakovich,
hamilton:miranda, and handel:messiah.

\mybigspace Standard \mymarginpar{\textbf{eid}} \textsf{biblatex}
field, providing a string or number some journals use uniquely to
identify a particular article.  Only applicable to the
\textsf{article} entry type, and only to those without a
\texttt{magazine} \textsf{entrysubtype}.  The 17th edition of the
\emph{Manual} now specifies where to print this (14.174), and I have
moved it in accordance with its specifications.  It replaces the
\textsf{pages} field in the list of references.

\paragraph*{\protect\mymarginpar{\textbf{entrysubtype}}}
\label{sec:ad:entrysubtype}
Standard and very powerful \textsf{biblatex} field, left undefined by
the standard styles.  In \textsf{bibla\-tex-chicago-authordate} it has
seven very specific uses, the first three of which I have designed in
order to maintain, as much as possible, backward compatibility with
the standard styles.  First, in \textsf{article} and
\textsf{periodical} entries, the field allows you to differentiate
between scholarly \enquote{journals,} on the one hand, and
\enquote{magazines} and \enquote{newspapers} on the other.  Usage is
fairly simple: you need to put the exact string \texttt{magazine} into
the \textsf{entrysubtype} field if you are citing one of the latter
two types of source, whereas if your source is a \enquote{journal,}
then you need do nothing.

\mylittlespace The second use involves references to works from
classical antiquity and, according to the \emph{Manual}, from the
Middle Ages, as well.  When you cite such a work using the traditional
divisions into books, sections, lines, etc., divisions which are
presumed to be the same across all editions, then you need to put the
exact string \texttt{classical} into the \textsf{entrysubtype} field.
This has no effect in the list of references, which will still present
the particular edition you are using, but it does affect the
formatting of in-text citations, in two ways.  First, it suppresses
some of the punctuation.  Second, and more importantly, it suppresses
the \textsf{date} field in favor of the \textsf{title}, so that
citations look like (Aristotle \emph{Metaphysics} 3.2.996b5--8)
instead of (Aristotle 1997, 3.2.996b5--8).  This mechanism may also
prove useful in \textsf{misc} entries for citations from the Bible or
other sacred texts (cf.\ genesis), and for citing archival collections
(house:papers), where it produces citations of the form (House
Papers).  (Cf.\ the next but one paragraph.)

\mylittlespace If you wish to reference a classical or medieval work
by the page numbers of a particular, non-standard edition, then you
shouldn't use the \texttt{classical} \textsf{entrysubtype} toggle.
Also, and the specification is reasonably clear about this, works from
the Renaissance and later, even if cited by the traditional divisions,
seem to have citations formatted normally, and therefore don't need an
\textsf{entrysubtype} field.  (See \emph{Manual} 14.242--54;
aristotle:metaphy:gr, herodotus:wilson, plato:republic:gr;
euripides:orestes is an example of a translation cited by page number
in a modern edition.  Cf.\ also the \mycolor{\texttt{notitle}} option
in section~\ref{sec:authuseropts}.)

\mylittlespace The third use of the \textsf{entrysubtype} field occurs
in \textsf{misc} entries.  If such an entry contains no such field,
then the citation will be treated just as the standard
\textsf{biblatex} styles would, including the use of italics for the
\textsf{title}.  Any string at all in \textsf{entrysubtype} tells
\textsf{biblatex-chicago} to treat the source as part of an
unpublished archive.  Please see section~\ref{sec:types:authdate}
above under \textbf{misc} for all the details on how these citations
work.

\mylittlespace Fourth, the field can be defined in the
\textsf{artwork} entry type in order to refer to a work from antiquity
whose title you do not wish to be italicized.  Please see the
documentation of \textsf{artwork} above for the details.  Fifth, you
can define it in an \textbf{audio}, \textbf{music}, or \textbf{video}
entry if such an entry refers to an individual unit that isn't part of
any larger collection, the entry therefore having only a
\textsf{title} and not a \textsf{booktitle}, a \textsf{title} that
\textsf{biblatex-chicago} would normally interpret as the title of a
larger unit (and therefore italicize).  Sixth, and sticking with the
\textbf{video} type, though enacting quite a different syntactic
transformation, the 17th edition (14.265) now recommends that, when
presenting episodes from a TV series, the name of the series
(\textsf{booktitle}) comes before the episode name (\textsf{title}).
The exact string \texttt{tvepisode} in the \textsf{entrysubtype} field
achieves this reversal, which includes using the \textsf{booktitle} as
a \textsf{sorttitle} in the list of references.

\mylittlespace Seventh, and finally, you can use any
\textsf{entrysubtype} whatever in \textsf{inreference} entries in
order to treat them as inherently online works rather than standard
published works.  See the documentation of \textbf{online} and
\textbf{inreference} entries in section~\ref{sec:types:authdate},
above, and also 14.233 and wikiped:bibtex.

\mybigspace Kazuo
\mymarginpar{\textbf{eprint}\\\textbf{eprintclass}\\\textbf{eprinttype}}
Teramoto suggested adding \textsf{biblatex's} excellent
\textsf{eprint} handling to \textsf{biblatex-chica\-go}, and he sent
me a patch implementing it.  I have applied it, with minor
alterations, so these three fields work more or less as they do in
standard \textsf{biblatex}.  They may prove helpful in providing more
abbreviated references to online content than conventional URLs,
though I can find no specific reference to them in the \emph{Manual}.

\mybigspace This \mymarginpar{\textbf{eventdate}} is a standard
\textsf{biblatex} field which has gradually accumulated functions in
\textsf{biblatex-chicago}.  It can now play a role in
\textsf{artwork}, \textsf{audio}, \textsf{image},
\textsf{inproceedings}, \textsf{music}, \textsf{proceedings},
\textsf{review}, \textsf{standard}, \textsf{suppperiodical},
\textsf{unpublished}, and \textsf{video} entries.  In \textsf{artwork}
and \textsf{image} entries it identifies the publication date of, most
frequently, a photograph, in association with the
\textsf{howpublished} field which identifies the periodical or other
medium in which it was published (mccurry:afghangirl).  In
\textsf{standard} entries it will also usually be associated with a
\textsf{howpublished} field, allowing you to specify a later renewal
or reaffirmation of a standard (niso\hc bibref).  In \textsf{audio}
entries, it specifies the release date of a single episode of a
podcast (danforth:podcast).  In \textsf{music} entries, it identifies
the recording or performance date of a particular song (rather than of
a whole disc, for which you would use \textsf{origdate}), whereas in
\textsf{video} entries it identifies either the original broadcast
date of a particular episode of a TV series or the date of a filmed
musical performance.  In both these cases \textsf{biblatex-chicago}
will automatically prepend a bibstring --- \texttt{recorded} and
\texttt{aired}, respectively --- to the date, but you can change this
string using the \textsf{userd} field, something you'll definitely
want to do for filmed musical performances (friends:leia,
handel:messiah, holiday:fool).

\mylittlespace In \textsf{inproceedings}, \textsf{proceedings}, and
\textsf{unpublished} entries it identifies the date of an event at
which a published or unpublished work was presented, though in truth
the \textsf{date} will do as well in \textsf{unpublished} entries
(nass:address).  The field's use in \textsf{review} and
\textsf{suppperiodical} entries, finally, includes a possible time
stamp.  In this context, an \textsf{eventdate} helps to identify a
particular comment on, or reply to another comment on, a blog post.
Given that many such posts by a single \textsf{author} could appear on
the same day, you can distinguish them by putting a time specification
in the \textsf{eventdate} field itself (ac:comment).  Please see the
\textbf{review} type, above, for the details of how to cite these
materials, possibly with the help of the new \texttt{commenton}
\textsf{relatedtype}.  See also the \textsf{date} field docs above, in
particular table~\ref{ad:date:extras}, for details on how the
\textsc{iso}8601-2 Extended Format specifications offered by
\textsf{biblatex}, including time stamps and much else besides, have
been implemented in \textsf{biblatex-chicago}.

\mylittlespace In the default configuration of \cmd{DeclareLabeldate},
dates for citations and for the head of reference list entries are
searched for in the order \textsf{date, eventdate, origdate, urldate}.
This suits the Chicago author-date styles very well, except for
\textsf{music}, \textsf{review}, \textsf{standard},
\textsf{suppperiodical}, and \textsf{video} entries.  In
\textsf{music} and \textsf{video} entries the general rule is to
emphasize the earliest date, whether that be, for example, the
recording date or original release date (15.57).  The other three
types have special requirements that once again necessitate putting
the \textsf{eventdate} at the head of the queue.  For these five entry
types, then, \cmd{DeclareLabeldate} uses the order \textsf{eventdate,
  origdate, date, urldate}.  (See the \texttt{avdate} option in
section~\ref{sec:authpreset}, below.)

\mybigspace This \mymarginpar{\textbf{eventtimezone}} field can, if
necessary, specify the time zone associated with a time stamp given as
part of an \textsf{eventdate}.  The \emph{Manual} prefers initialisms
like \enquote{EST} for this purpose, and you can provide parentheses
around it at your discretion (cp.\ 10.41 and 14.191).

\mybigspace A \mymarginpar{\textbf{eventtitle}} standard
\textsf{biblatex} field for identifying the name of the event that
produces either a published record (\textsf{inproceedings} and
\textsf{proceedings} entries) or an unpublished one
(\textsf{unpublished}).

\mybigspace Standard \mymarginpar{\textbf{eventtitleaddon}}
\textsf{biblatex} field for adding information about an
\textsf{eventtitle}, and available in the same entry types as that
field.

\mybigspace As \mymarginpar{\textbf{foreword}} with the
\textsf{afterword} field above, \textsf{foreword} will in general
function as it does in standard \textsf{biblatex}.  Like
\textsf{afterword} (and \textsf{introduction}), however, it has a
special meaning in a \textsf{suppbook} entry, where you simply need to
define it somehow (and leave \textsf{afterword} and
\textsf{introduction} undefined) to make a foreword the focus of a
citation.

\mybigspace A \mymarginpar{\textbf{holder}} standard \textsf{biblatex}
field for identifying a \textsf{patent}'s holder(s), if they differ
from the \textsf{author}.  The \emph{Manual} has nothing to say on the
subject, but \textsf{biblatex-chicago} prints it (them), in
parentheses, just after the author(s).

\mybigspace Standard \mymarginpar{\textbf{howpublished}}
\textsf{biblatex} field which, like the \textsf{eventdate} field, is
gradually accumulating functions in \textsf{biblatex-chicago}.  In the
\textsf{booklet} type it retains something of its traditional usage,
replacing the \textsf{publisher}, and has a similar (somewhat
paradoxical) place in \textsf{unpublished} entries.  In the
\textsf{misc} and \textsf{performance} types it works almost as a
second \textsf{note} field, bringing in extra information about a work
in close association with the \textsf{type} and \textsf{version}
fields, while the \textsf{dataset} entry type associates it both with
those two fields and with the \textsf{number} field.  17th-edition
\textsf{music} entries require a field to provide the medium of
downloaded music and/or the name of the streaming service, so
\textsf{howpublished} works there as an online double of \textsf{type}
and of \textsf{publisher}.  Finally, in \textsf{artwork},
\textsf{image}, and \textsf{standard} entries it serves to qualify or
modify an \textsf{eventdate}, almost as a \textsf{userd} field
modifies a \textsf{date} or \textsf{urldate}.  Please see the docs of
those entry types for more information, and also bedford:photo,
clark:mesopot, mccurry:afghangirl, niso:bibref, rihanna:umbrella.

\mybigspace Standard \mymarginpar{\textbf{institution}}
\textsf{biblatex} field.  In the \textsf{thesis} entry type, it will
usually identify the university for which the thesis was written,
while in a \textsf{report} entry it may identify any sort of
institution issuing the report.

\mybigspace As \mymarginpar{\textbf{introduction}} with the
\textsf{afterword} and \textsf{foreword} fields above,
\textsf{introduction} will in general function as it does in standard
\textsf{biblatex}.  Like those fields, however, it has a special
meaning in a \textsf{suppbook} entry, where you simply need to define
it somehow (and leave \textsf{afterword} and \textsf{foreword}
undefined) to make an introduction the focus of a citation.

\mybigspace Standard \mymarginpar{\textbf{isbn}} \textsf{biblatex}
field, for providing the International Standard Book Number of a
publication.  Not typically required by the \emph{Manual}.

%%\enlargethispage{-2\baselineskip}

\mybigspace Standard \mymarginpar{\textbf{isrn}} \textsf{biblatex}
field, for providing the International Standard Technical Report
Number of a report.  Only relevant to the \textsf{report} entry type,
and not typically required by the \emph{Manual}.

\mybigspace Standard \mymarginpar{\textbf{issn}} \textsf{biblatex}
field, for providing the International Standard Serial Number of a
periodical in an \textsf{article} or a \textsf{periodical} entry.  Not
typically required by the \emph{Manual}.

\mybigspace Standard \mymarginpar{\textbf{issue}} \textsf{biblatex}
field, designed for \textsf{article} or \textsf{periodical} entries
identified by something like \enquote{Spring} or \enquote{Summer}
rather than by the usual \textsf{month} or \textsf{number} fields
(brown\hc bremer).  \textsf{Biblatex's} enhanced date handling allows
you to specify a season in the \textsf{date} field, with the
\enquote{months} 21--24 used for Spring, Summer, Autumn, and Winter,
respectively.  Cf.\ table~\ref{ad:date:extras}, above.

\mybigspace The \mymarginpar{\textbf{issuesubtitle}} subtitle for an
\textsf{issuetitle} --- see next entry.

\mybigspace Standard \mymarginpar{\textbf{issuetitle}}
\textsf{biblatex} field, intended to contain the title of a special
issue of any sort of periodical.  If the reference is to one article
within the special issue, then this field should be used in an
\textsf{article} entry (conley:fifthgrade), whereas if you are citing
the entire issue as a whole, then it would go in a \textsf{periodical}
entry, instead (good:wholeissue).  The \textsf{note} field is the
proper place to identify the type of issue, e.g.,\ \texttt{special
  issue}, with the initial letter lower-cased to enable automatic
contextual capitalization.

\mybigspace The \mymarginpar{\textbf{journalsubtitle}} subtitle for a
\textsf{journaltitle} --- see next entry.

\mybigspace Standard \mymarginpar{\textbf{journaltitle}}
\textsf{biblatex} field, replacing the standard \textsc{Bib}\TeX\
field \textsf{journal}, which, however, still works as an alias.  It
contains the name of any sort of periodical publication, and is found
in the \textsf{article} and \textsf{review} entry types.  In the case
where a piece in an \textsf{article} or \textsf{review}
(\textsf{entrysubtype} \texttt{magazine}) doesn't have an author,
\textsf{biblatex-chicago} provides for this field to be used as the
author.  See above (section~\ref{sec:fields:authdate}) under
\textbf{article} for details.  The lakeforester:pushcarts and
nyt:trevorobit entries in \textsf{dates-test.bib} will give you some
idea of how this works.  Please note there is a \textsf{shortjournal}
field which you can use to abbreviate the \textsf{journaltitle} in
citations and/or in the reference list, and you can also use it to
print a list of journal abbreviations.  Cf.\ the \textsf{shortjournal}
documentation below.

\mybigspace An \colmarginpar{\textbf{journaltitleaddon}} annex to the
\textsf{journaltitle}, for which see previous entry.  Such an annex
would be printed in the main text font.  If your data begins with a
word that would ordinarily only be capitalized at the beginning of a
sentence, then simply ensure that that word is in lowercase, and
\textsf{biblatex-chicago-authordate} will automatically do the right
thing.  The package and entry option \mycolor{\texttt{jtitleaddon}}
(section~\ref{sec:authpreset}) allows you to customize the punctuation
that appears before the \textsf{journaltitleaddon} field (hua:cms).
The default is a \texttt{space}.

\mybigspace This \mymarginpar{\textbf{keywords}} field is
\textsf{biblatex}'s powerful and flexible technique for filtering
entries in a list of references, allowing you to subdivide it
according to just about any criteria you care to invent, or indeed to
prevent entries in citations from appearing in reference list, as the
\emph{Manual} sometimes recommends.  See \textsf{biblatex.pdf} (3.7)
for thorough documentation.

\mybigspace A \mymarginpar{\textbf{language}} standard
\textsf{biblatex} field, designed to allow you to specify the
language(s) in which a work is written.  As a general rule, the
Chicago style doesn't require you to provide this information, though
it may well be useful for clarifying the nature of certain works, such
as bilingual editions, for example.  There is at least one situation,
however, when the \emph{Manual} does specify this data, and that is
when the title of a work is given in translation, even though no
translation of the work has been published, something that might
happen when a title is in a language deemed to be unparseable by a
majority of your expected readership (14.99; chu:panda, pirumova,
rozner:liberation).  In such a case, you should provide the
language(s) involved using this field, connecting multiple languages
using the keyword \texttt{and}.  (I have retained \textsf{biblatex's}
\cmd{bibstring} mechanism here, which means that you can use the
standard bibstrings or, if one doesn't exist for the language you
need, just give the name of the language, capitalized as it should
appear in your text.  You can also mix these two modes inside one
entry without apparent harm.)

\mylittlespace An alternative arrangement suggested by the
\emph{Manual} is to retain the original title of a piece but then to
provide its translation, as well.  If you choose this option, you'll
need to make use of the \textbf{usere} field, on which see below.  In
effect, you'll probably only ever need to use one of these two fields
in any given entry, and in fact \textsf{biblatex-chicago} will only
print one of them if both are present, preferring \textsf{usere} over
\textsf{language} for this purpose (see kern, pirumova:russian, and
weresz).  Note also that both of these fields are universally
associated with the \textsf{title} of a work, rather than with a
\textsf{booktitle} or a \textsf{maintitle}.  If you need to attach a
language or a translation to either of the latter two, you could
probably manage it with special formatting inside those fields
themselves.

\mybigspace I \mymarginpar{\textbf{lista}} intend this field
specifically for presenting citations from reference works that are
arranged alphabetically, where the name of the article rather than a
page or volume number should be given.  The field is a
\textsf{biblatex} list, which means you should separate multiple items
with the keyword \texttt{and}.  Each item receives its own set of
quotation marks, and the whole list will be prefixed by the
appropriate string (\enquote{s.v.,} \emph{sub verbo}, pl.\
\enquote{s.vv.}).  \textsf{Biblatex-chicago} will only print such a
field in a \textsf{book} or an \textsf{inreference} entry, and you
should look at the documentation of these entry types for further
details.  (See \emph{Manual} 14.232--33; grove:sibelius, times:guide,
wikiped:bibtex.)

\mybigspace This \mymarginpar{\textbf{location}} is
\textsf{biblatex}'s version of the usual \textsc{Bib}\TeX\ field
\textsf{address}, though the latter is accepted as an alias if that
simplifies the modification of older .bib files.  According to the
\emph{Manual} (14.129), a citation usually need only provide the first
city listed on any title page, though a list of cities separated by
the keyword \enquote{\texttt{and}} will be formatted appropriately.
If the place of publication is unknown, you can use
\verb+\autocap{n}.p.+\ instead (14.132).  For all cities, you should
use the common English version of the name, if such exists (14.131).

\mylittlespace Two other uses need explanation here.  In
\textsf{article}, \textsf{periodical}, and \textsf{review} entries,
there is usually no need for a \textsf{location} field, but
\enquote{if a journal might be confused with another with a similar
  title, or if it might not be known to the users of a bibliography,}
then this field can present the place or institution where it is
published (14.182, 14.191, 14.193--94; garrett, kimluu:diethyl, and
lakeforester:pushcarts).  For blogs cited using \textsf{article}
entries, this is a good place to identify the nature of the source ---
i.e., the word \enquote{blog} --- letting the style automatically
provide the parentheses (15.51; ellis:blog).

\mybigspace The \mymarginpar{\textbf{mainsubtitle}} subtitle for a
\textsf{maintitle} --- see next entry.

\mybigspace The \mymarginpar{\textbf{maintitle}} main title for a
multi-volume work, e.g., \enquote{Opera} or \enquote{Collected Works.}
It no longer takes sentence-style capitalization in
\textsf{authordate}, though it does in \textsf{authordate-trad}.  In
cross references produced using the \textsf{crossref} field, the
\textsf{title} of \textbf{mv*} entry types always becomes a
\textsf{maintitle} in the child entry.  (See donne:var,
euripides:orestes, harley\hc cartography, lach:asia,
pelikan:christian, and plato:republic:gr.)

\mylittlespace Because the 17th edition of the \emph{Manual}
recommends that you present not only the names of blogs but also the
names of their parent (usually periodical) publications, I have added
this field to \textsf{article}, \textsf{periodical}, and
\textsf{review} entries for just this purpose.  See the documentation
of those entry types in section~\ref{sec:types:authdate}, above, and
also table~\ref{tab:online:adtypes} (15.51; amlen:hoot).

\mybigspace An \mymarginpar{\textbf{maintitleaddon}} annex to the
\textsf{maintitle}, for which see previous entry.  Such an annex would
be printed in the main text font.  If your data begins with a word
that would ordinarily only be capitalized at the beginning of a
sentence, then simply ensure that that word is in lowercase, and
\textsf{biblatex-chicago} will automatically do the right thing.  The
package and entry options \texttt{ptitleaddon} and
\texttt{ctitleaddon} (section~\ref{sec:authpreset}) allow you to
customize the punctuation that appears before the
\textsf{maintitleaddon} field (schubert:muellerin).

\mybigspace Standard \mymarginpar{\textbf{month}} \textsf{biblatex}
field, containing the month of publication.  This should be an
integer, i.e., \texttt{month=\{3\}} not \texttt{month=\{March\}}.  See
\textsf{date} for more information.

\mybigspace This \mymarginpar{\textbf{namea}} is one of the fields
\textsf{biblatex} provides for style writers to use, but which it
leaves undefined itself.  In \textsf{biblatex-chicago} it contains the
name(s) of the editor(s) of a \textsf{title}, if the entry has a
\textsf{booktitle} and/or a \textsf{maintitle}, in which situation the
\textsf{editor} would be associated with one of these latter fields
(donne:var).  (In \textsf{article} and \textsf{review} entries,
\textsf{namea} applies to the \textsf{title} instead of the
\textsf{issuetitle}, should the latter be present.)  You should
present names in this field exactly as you would those in an
\textsf{author} or \textsf{editor} field, and the package will
concatenate this field with \textsf{nameb} if they are identical.
When choosing a name for a citation or to head a reference-list entry,
\textsf{biblatex-chicago} gives precedence to \textsf{namea} over
\textsf{editor}.  See under \textbf{editor} and \textbf{editortype}
above for the full details.  Please note that, as the field is highly
single-entry specific, \textsf{namea} isn't inherited from a
\textsf{crossref}'ed parent entry.  Please note, also, that you can
use the \textsf{nameatype} field to redefine this role just as you can
with \textsf{editortype}, which see.  Cf.\ also \textsf{nameb},
\textsf{namec}, \textsf{translator}, and the macros \cmd{partedit},
\cmd{parttrans}, \cmd{parteditandtrans}, \cmd{partcomp},
\cmd{parteditandcomp}, \cmd{parttransandcomp}, and
\cmd{partedittransandcomp}, for which see
section~\ref{sec:formatting:authdate}.

\mybigspace This \mymarginpar{\textbf{nameaddon}} field is provided by
\textsf{biblatex}, though not used by the standard styles.  In
\textsf{biblatex-chicago} its primary use, in most entry types, has
always been to specify that an author's name is a pseudonym, or to
provide either the real name or the pseudonym itself, if the other is
being provided in the \textsf{author} field.  The abbreviation
\enquote{\texttt{pseud.}\hspace{-2pt}}\ (always lowercase in English)
is specified, either on its own or after the pseudonym
(centinel:letters, creasey:ashe:blast, creasey:morton:hide,
creasey:york:death, and lecarre:quest); remember that
\verb+\bibstring{pseudonym}+ does the work for you.  See under
\textbf{author} above for the full details.

\mylittlespace The field has slowly accumulated other functions, so
when Philipp Immel made a feature request, and pointed to a discussion
on Stack Exchange which suggested a few more, I thought I might
generalize the field's functionality, providing three package options
to allow users to mould it to their needs.  Before discussing these,
allow me to emphasize that the package defaults remain exactly the
same as before, so that, absent any of the new options, the style
still provides square brackets around the \textsf{nameaddon} in most
entry types, no brackets of any sort in \textsf{online, review,} and
\textsf{suppperiodical} entries, as well as in \textsf{misc} entries
with an \textsf{entrysubtype}, and rather specialized handling in
\textsf{customc} entries (which ignore all but the first of the new
options --- see below).  If you're happy with the status quo, then no
changes to your documents or\ .bib databases are necessary.

\mylittlespace If you do need or want to put the field to a different
use, the following options may help.  All of them are available
globally, per type, and per entry.  The first
\mymarginpar{\texttt{nameaddon}} new option is simply called
\texttt{nameaddon}, and determines where and when the field will be
printed at all.  There are three possible values:
\begin{description}
\setlength{\parskip}{-4pt}
\item[\qquad \texttt{all}:] This is the default; if an entry has a
  \textsf{nameaddon}, it will appear in the reference list.
\item[\qquad \texttt{none}:] The field will not appear in the
  reference list.
\item[\qquad \texttt{first}:] Philipp Immel requested this as a way to
  provide an \textsf{author's} dates in the \textsf{nameaddon} field
  and only have them printed the first time that author appears in the
  reference list.  The code tests for identical \textsf{nameaddon}
  fields in works by identical \textsf{authors}, so other sorts of
  \textsf{nameaddon} will be printed as usual.
\end{description}

The \mymarginpar{\texttt{nameaddonsep}} \texttt{nameaddonsep} option
controls the punctuation that appears before the
\textsf{nameaddon}. It takes the following six keys:

\begin{description}
\setlength{\parskip}{-4pt}
\item[\qquad \texttt{space}] = \cmd{addspace}.  This is the default.
\item[\qquad \texttt{none}] = no separator at all.  It presumes that
  you will include one in the \textsf{nameaddon} field itself.
\item[\qquad \texttt{colon}] = \verb+\addcolon\addspace+.
\item[\qquad \texttt{comma}] = \verb+\addcomma\addspace+.
\item[\qquad \texttt{period}] = \verb+\addperiod\addspace+.
\item[\qquad \texttt{semicolon}] = \verb+\addsemicolon\addspace+.
\end{description}

The \mymarginpar{\texttt{nameaddon-\\format}} \texttt{nameaddonformat}
option allows you to change the format of the \textsf{nameaddon} field
on the fly, so its value should be a field format that
\textsf{biblatex} understands.  This includes standard formats like
\texttt{parens,\,brackets} or \texttt{emph}, and also custom formats
that you provide in your preamble using \cmd{DeclareFieldFormat}, in
case the standard ones aren't adequate.  If you don't define this
option, then the usual defaults apply, as delineated above, and you
can use your own parentheses in \textsf{online, review,} and
\textsf{suppperiodical} entries, as well as in \textsf{misc} entries
with an \textsf{entrysubtype}, to distinguish screen names or other
authorial information from traditional pseudonyms (in brackets).

\mylittlespace There are two more details.  If you are using the
17th-edition styles for the first time, please note that the 16th
edition of the \emph{Manual} recommended specifying comments to blogs
and other online content using a time stamp in parentheses after the
\textsf{author}, but the 17th edition handles time stamps both
differently and more widely, so in this case you would now put time
data into the \textsf{date} or \textsf{eventdate} field, particularly
when the date itself is too coarse a specification to identify a
comment unambiguously.  Also, the \emph{Manual} (15.52) now specifies
that comments should appear \enquote{only in the text, in reference to
  the related post,} so I've provided some new functionality to enable
this.  Please see the \textbf{online} and \textbf{review} types,
above, especially table~\ref{tab:online:adtypes}, for the details of
how to cite these materials, possibly with the help of the new
\texttt{commenton} \textsf{relatedtype} and a separate
\textbf{customc} entry.  See also the \textsf{date} field docs above,
in particular table~\ref{ad:date:extras}, for details on how the
\textsc{iso}8601-2 Extended Format specifications offered by
\textsf{biblatex}, including time stamps and much else besides, have
been implemented in \textsf{biblatex-chicago}.  (Cf.\ ellis:blog,
obrien:recycle).

\mylittlespace In the \textsf{customc} entry type, finally, which is
used to create alphabetized cross-references to other bibliography
entries, the \textsf{nameaddon} field allows you to change the default
string linking the two parts of the cross-reference.  The code
automatically tests for a known bibstring, which it will italicize.
Otherwise, it prints the string as you've provided in the
\textsf{nameaddon} field itself.  The punctuation is fixed.

\mybigspace You \mymarginpar{\textbf{nameatype}} can use this field
to change the role of a \textsf{namea} just as you can use
\textsf{editortype} to change the role of an \textsf{editor}.  As with
the \textsf{editortype}, using this field prevents string
concatenation with identical \textsf{nameb} or \textsf{namec} fields.
Please see \textbf{editortype}, above, for the details.

\mybigspace Like \mymarginpar{\textbf{nameb}} \textsf{namea}, above,
this is a field left undefined by the standard \textsf{biblatex}
styles.  In \textsf{biblatex-chicago}, it contains the name(s) of the
translator(s) of a \textsf{title}, if the entry has a
\textsf{booktitle} or \textsf{maintitle}, or both, in which situation
the \textsf{translator} would be associated with one of these latter
fields (euripides:orestes).  (In \textsf{article} and \textsf{review}
entries, \textsf{nameb} applies to the \textsf{title} instead of the
\textsf{issuetitle}, should the latter be present.)  You should
present names in this field exactly as you would those in an
\textsf{author} or \textsf{translator} field, and the package will
concatenate this field with \textsf{namea} if they are identical.  See
under the \textbf{translator} field below for the full details.
Please note that, as the field is highly single-entry specific,
\textsf{nameb} isn't inherited from a \textsf{crossref}'ed parent
entry.  Please note, also, that in \textsf{biblatex-chicago's}
name-finding algorithms \textsf{nameb} takes precedence over
\textsf{translator}.  Cf.\ also \textsf{namea}, \textsf{namec},
\textsf{origlanguage}, \textsf{translator}, \textsf{userf} and the
macros \cmd{partedit}, \cmd{parttrans}, \cmd{parteditandtrans},
\cmd{partcomp}, \cmd{parteditandcomp}, \cmd{parttransandcomp}, and
\cmd{partedittransandcomp} in section~\ref{sec:formatting:authdate}.

\mybigspace The \mymarginpar{\textbf{namec}} \emph{Manual} (15.36)
specifies that works without an author may be listed under an editor,
translator, or compiler, assuming that one is available, and it also
specifies the strings to be used with the name(s) of compiler(s).  All
this suggests that the \emph{Manual} considers this to be standard
information that should be made available in a bibliographic
reference, so I have added that possibility to the many that
\textsf{biblatex} already provides, such as the \textsf{editor},
\textsf{translator}, \textsf{commentator}, \textsf{annotator}, and
\textsf{redactor}, along with writers of an \textsf{introduction},
\textsf{foreword}, or \textsf{afterword}.  Since \textsf{biblatex}
doesn't offer a \textsf{compiler} field, I have adopted for this
purpose the otherwise unused field \textsf{namec}.  It is important to
understand that, despite the analogous name, this field does not
function like \textsf{namea} or \textsf{nameb}, but rather like
\textsf{editor} or \textsf{translator}, and therefore if used will be
associated with whichever title field these latter two would be were
they present in the same entry.  Identical fields among these three
will be concatenated by the package, and concatenated too with the
(usually) unnecessary commentator, annotator and the rest.  Also
please note that I've arranged the concatenation algorithms to include
\textsf{namec} in the same test as \textsf{namea} and \textsf{nameb},
so in this particular circumstance you can, if needed, make
\textsf{namec} analogous to these two latter, \textsf{title}-only
fields.  (See above under \textbf{editortype} for details of how you
can use that field, or the \textsf{nameatype} field, to identify a
compiler.)

\mylittlespace It might conceivably be necessary at some point to
identify the compiler(s) of a \textsf{title} separate from the
compiler(s) of a \textsf{booktitle} or \textsf{maintitle}, but for the
moment I've run out of available \textsf{name} fields, so you'll have
to fall back on the \cmd{partcomp} macro or the related
\cmd{parteditandcomp}, \cmd{parttransandcomp}, and
\cmd{partedittransandcomp}, on which see Commands
(section~\ref{sec:formatting:authdate}) below.  (Future releases may
be able to remedy this.)  It may be as well to mention here too that
of the three names that can be substituted for the missing
\textsf{author} at the head of an entry, \textsf{biblatex-chicago}
will choose a \textsf{namea} if present, then an \textsf{editor}, a
\textsf{nameb}, or a \textsf{translator}, with \textsf{namec} coming
last, assuming that the fields aren't identical, and therefore to be
concatenated.  \textsf{Biblatex}'s sorting algorithms, and also its
\textsf{labelname} mechanism, should both work properly no matter what
sort of name you provide, but do please remember that if you want the
package to skip over any names you can employ the
\texttt{use<name>=false} options.  Indeed, \textsf{biblatex's}
\texttt{usenamec} has replaced the old Chicago-specific
\texttt{usecompiler}, which is deprecated.

\mybigspace As \mymarginpar{\textbf{note}} in standard
\textsf{biblatex}, this field allows you to provide bibliographic data
that doesn't easily fit into any other field.  In this sense, it's
very like \textsf{addendum}, but the information provided here will be
printed just before the publication data.  (See chaucer:alt,
cook:sotweed, emerson:nature, and rodman:walk for examples of this
usage in action.)  It also has a specialized use in the periodical
types (\textsf{article}, \textsf{periodical}, and \textsf{review}),
where it holds supplemental information about a \textsf{journaltitle},
such as \enquote{special issue} (conley\hc fifthgrade,
good:wholeissue).  In all uses, if your data begins with a word that
would ordinarily only be capitalized at the beginning of a sentence,
then simply ensure that that word is in lowercase, and
\textsf{biblatex-chicago} will automatically do the right thing.  Cf.\
\textsf{addendum}.

\mybigspace This \mymarginpar{\textbf{number}} is a standard
\textsf{biblatex} field, steadily accumulating uses in
\textsf{biblatex-chicago}.  It may contain the number of a
\textsf{journaltitle} in an \textsf{article} or \textsf{review} entry,
the number of a \textsf{title} in a \textsf{periodical} entry, the
volume/number of a book (or musical recording) in a \textsf{series},
the (generally numerical) specifier of the \textsf{type} in a
\textsf{report} entry, the archive location (or database accession
number) of a \textsf{dataset} entry, and the number of a national or
international standard in a \textsf{standard} entry.  Generally, in an
\textsf{article}, \textsf{periodical}, or \textsf{review} entry, this
will be a plain cardinal number, but in such entries
\textsf{biblatex-chicago} does the right thing if you have a list or
range of numbers (unsigned:ranke).  In any \textsf{book}-like entry it
may well contain considerably more information, including even a
reference to \enquote{2nd ser.,} for example, while the
\textsf{series} field in such an entry will contain the name of the
series, rather than a number.  This field is also the place for the
patent number in a \textsf{patent} entry.  Cf.\ \textsf{issue} and
\textsf{series}.  (See 14.123--25 and boxer:china, palmatary:pottery,
wauchope:ceramics; 14.171 and beattie:crime, conley:fifthgrade,
friedman:learning, garrett, gibbard, hlatky:hrt, mcmillen:antebellum,
rozner:liberation, warr:ellison; 14.257 and genbank:db; 14.259 and
niso:bibref; 14.263 and holiday:fool.)

\mylittlespace \textbf{NB}: This may be an opportune place to point
out that the \emph{Manual} (14.147) prefers arabic to roman numerals
in most circumstances (chapters, volumes, series numbers, etc.), even
when such numbers might be roman in the work cited.  The obvious
exception is page numbers, in which roman numerals indicate that the
citation came from the front matter, and should therefore be retained.

\mybigspace A \mymarginpar{\textbf{options}} standard
\textsf{biblatex} field, for setting certain options on a per-entry
basis rather than globally.  Information about some of the more common
options may be found above under \textsf{author} and \textsf{date},
and below in section~\ref{sec:authuseropts}.  See creel:house,
eliot:pound, emerson:nature, ency:britannica, herwign:office,
lecarre:quest, and maitland:canon for examples of the field in use.

\mybigspace A \mymarginpar{\textbf{organization}} standard
\textsf{biblatex} field, retained mainly for use in the \textsf{misc},
\textsf{online}, and \textsf{manual} entry types, where it may be of
use to specify a publishing body that might not easily fit in other
categories.  In \textsf{biblatex}, it is also used to identify the
organization sponsoring a conference in a \textsf{proceedings} or
\textsf{inproceedings} entry, and I have retained this as a
possibility, though the \emph{Manual} is silent on the matter.

% %\enlargethispage{\baselineskip}

\mybigspace This \mymarginpar{\textbf{origdate}} is a standard
\textsf{biblatex} field which allows more than one full date
specification for those references which need it.  (You can also
provide a time stamp in the field, after an uppercase
\enquote{\texttt{T}}, but I foresee this being very rarely needed in
the author-date styles.  See table~\ref{ad:date:extras} for
\textsf{biblatex-chicago's} implementation of \textsf{biblatex's}
enhanced date specifications.)  As with the analogous \textsf{date}
field, you provide the date (or range of dates) in \textsc{iso}8601
format, i.e., \texttt{yyyy-mm-dd}.  In most entry types, you would use
\textsf{origdate} to provide the date of first publication of a work,
most usually needed only in the case of reprint editions, but also
recommended by the \emph{Manual} for electronic editions of older
works (15.40, 14.114, 14.162; aristotle:metaphy:gr, emerson:nature,
james:ambassadors, schweitzer:bach).  In both the \textsf{letter} and
\textsf{misc} (with \textsf{entrysubtype)} entry types, the
\textsf{origdate} identifies when a letter (or similar) was written.
In such \textsf{misc} entries, some \enquote{non-letter-like}
materials (like interviews) need the \textsf{date} field for this
purpose, while in \textsf{letter} entries the \textsf{date} applies to
the publication of the whole collection.  If such a published
collection were itself a reprint, judicious use of the
\textsf{pubstate} field or perhaps improvisation in the
\textsf{location} field might be able to rescue the situation.  (See
white:ross:memo, white:russ, and white:total for how \textsf{letter}
entries can work; creel:house shows the field in action in a
\textsf{misc} entry, while spock:interview uses \textsf{date}
instead.)

\mylittlespace Because of the importance of date specifications in the
author-date styles, \textsf{biblatex-chica\-go-authordate} and
\textsf{authordate-trad} provide options and automated behaviors that
allow you to emphasize the \textsf{origdate} in citations and at the
head of entries in the list of references.  In entries which have
\emph{only} an \textsf{origdate} --- usually \textsf{misc} with an
\textsf{entrysubtype} --- \textsf{Biber} and the default
\cmd{DeclareLabeldate} configuration make it possible to do without a
\texttt{cmsdate} option, as the \textsf{origdate} will automatically
appear where and as it should.  In \textsf{book}-like entries with
both a \textsf{date} and an \textsf{origdate}, the \emph{Manual}
recommends that you present, in citations and at the head of reference
list entries, only the \textsf{date} or both dates together.  The
latter is accomplished using the \texttt{cmsdate} entry option.  In
some cases it may even be necessary to reverse the two date fields,
putting the earlier year in \textsf{date} and the later in
\textsf{origdate}.  If your reference apparatus contains many such
instances, it may well be convenient for you instead to use the
\mymarginpar{\texttt{cmsdate}\\\emph{in preamble}} \texttt{cmsdate}
preamble option, which I have designed in an attempt to reduce the
amount of manual intervention needed to present lots of entries with
multiple dates.  In short, setting \texttt{cmsdate} to \texttt{both}
or \texttt{on} in the preamble promotes the \textsf{origdate} to the
top of the search for a \textsf{labeldate} to use in citations and at
the head of entries in the reference list.  This can solve many
problems with the \textsf{extradate} field --- 1978\textbf{a} --- and
also with sorting in the reference list.  Please see above under
\textbf{date} for all the details on how these options interact.

\mylittlespace In the default configuration of \cmd{DeclareLabeldate},
dates for citations and for the head of reference list entries are
searched for in the order \textsf{date, eventdate, origdate, urldate}.
If you set the \texttt{cmsdate} preamble options I've just mentioned,
this changes to \textsf{origdate, date, eventdate, urldate}.  These
generally cover the needs of the Chicago author-date styles well,
except for \textsf{music}, \textsf{standard}, and \textsf{video}
entries, and, exceptionally, some \textsf{review} and
\textsf{suppperiodical} entries.  Here the general rule is to
emphasize the earliest date.  For these five entry types, then,
\cmd{DeclareLabeldate} uses the order \textsf{eventdate, origdate,
  date, urldate}.  In \textsf{music} entries, you can use the
\textsf{origdate} in two separate but related ways.  First, it can
identify the recording date of an entire disc, rather than of one
track on that disc, which would go in \textsf{eventdate}.  (Compare
holiday:fool with nytrumpet:art.)  Second, the \textsf{origdate} can
provide the original release date of an album.  For this to happen,
you need to put the string \texttt{reprint} in the \textsf{pubstate}
field, which is the standard mechanism across many other entry types
for identifying a reprinted work.  (See floyd:atom.)  In
\textsf{video} entries, the \textsf{origdate} is intended for the
original release date of a film, whereas the \textsf{eventdate} would
hold the original broadcast date of, e.g., an episode of a TV series.
In both these two entry types, the style will, depending on the
context, automatically prepend appropriate bibstrings to the
\textsf{origdate}.  You can, assuming you've not activated the
\textsf{pubstate} mechanism in a \textsf{music} entry, choose a
different string using the \textsf{userd} field, but please be aware
that if an entry also has an \textsf{eventdate}, then \textsf{userd}
will apply to that, instead, and you'll be forced to accept the
default string.  (Compare friends:leia with hitchcock:nbynw; 15.57,
14.263--65; Cf.\ \texttt{cmsdate} in sections~\ref{sec:authuseropts}
and \ref{sec:authentryopts}, \cmd{DeclareLabeldate} in
section~\ref{sec:authformopts}, and \texttt{avdate} in
section~\ref{sec:authpreset}.)

\mylittlespace A couple of further notes are in order.  First,
\textsf{artwork} and \textsf{image} entries (which see) have their own
scheme, and are not governed by the \texttt{avdate} option.  Here, the
style uses the earlier of two dates as the creation date of the work
while the later is the printing date of, e.g., a particular exemplar
of a photograph or of an etching.  Depending on how you want this
information presented in an entry, you can distribute these dates
between the \textsf{date} and \textsf{origdate} fields as you wish.
Second, because the \textsf{origdate} field only accepts numbers, some
improvisation may be needed if you wish to include \enquote{n.d.}\
(\verb+\bibstring{nodate}+) in an entry.  In \textsf{letter} and
\textsf{misc}, this information can be placed in \textsf{titleaddon},
but in other entry types you may need to use the \textsf{location}
field.  (The \textsf{origyear} field usually works, too.)

\mybigspace See
\vspace{-14.2pt}
\mymarginpar{\textbf{origlanguage}\\
\textbf{origlocation}\\\textbf{origpublisher}}
section~\ref{sec:authrelated}, below.
\vspace{18pt}

\mybigspace This \mymarginpar{\textbf{origtimezone}} field can, if
necessary, specify the time zone associated with a time stamp given as
part of an \textsf{origdate}.  The \emph{Manual} prefers initialisms
like \enquote{EST} for this purpose, and you can provide parentheses
around it at your discretion (cp.\ 10.41 and 14.191).

\mybigspace This \mymarginpar{\textbf{pages}} is the standard
\textsf{biblatex} field for providing page references.  In many
\textsf{article} entries you'll find this contains something other
than a page number, e.g. a section name or edition specification
(14.191; kozinn:review, nyt:trevorobit).  Of course, the same may be
true of almost any sort of entry, though perhaps with less frequency.
Curious readers may wish to look at brown:bremer (14.180) for an
example of a \textsf{pages} field used to facilitate reference to a
two-part journal article.  Cf.\ \textsf{number} for more information
on the \emph{Manual}'s preferences regarding the formatting of
numerals; \textsf{bookpagination} and \textsf{pagination} provide
details about \textsf{biblatex's} mechanisms for specifying what sort
of division a given \textsf{pages} field contains; and \textsf{usera}
discusses a different way to present the section information
pertaining to a newspaper article.

\mylittlespace David Gohlke brought to my attention a discussion that
took place a couple of years ago on
\href{http://tex.stackexchange.com/questions/44492/biblatex-chicago-style-page-ranges}{Stackexchange}
regarding the automatic compression of page ranges, e.g., 101-{-}109
in the .bib file or in the \textsf{postnote} field would become 101--9
in the document.  \textsf{Biblatex} has long had the facilities for
providing this, and though the \emph{Manual's} rules (9.61) are fairly
complicated, Audrey Boruvka fortunately provided in that discussion
code that implements the specifications.  As some users may well be
accustomed to compressing page ranges themselves in their .bib files,
and in their \textsf{postnote} fields, I have made the activation of
this code a package option, so setting \texttt{compresspages=true}
when loading \textsf{biblatex-chicago} should automatically give you
the Chicago-recommended page ranges.  \textbf{NB}: the code now
resides in \textsf{biblatex-chicago.sty}, so if you don't load that
package then you'll need to copy the code into your preamble for the
option to have the desired effect.

\mybigspace This, \mymarginpar{\textbf{pagination}} a standard
\textsf{biblatex} field, allows you automatically to prefix the
appropriate identifying string to information you provide in the
\textsf{postnote} field of a citation command, whereas
\textsf{bookpagination} allows you to prefix a string to the
\textsf{pages} field.  Please see \textbf{bookpagination} above for
all the details on this functionality, as aside from the difference
just mentioned the two fields are equivalent.

%\enlargethispage{\baselineskip}

\mybigspace Standard \mymarginpar{\textbf{part}} \textsf{biblatex}
field, which identifies physical parts of a single logical volume in
\textsf{book}-like entries, not in periodicals.  It has the same
purpose in \textsf{biblatex-chicago}, but because the \emph{Manual}
(14.121) calls such a thing a \enquote{book} and not a \enquote{part,}
the string printed in the list of references will, at least in
English, be \enquote{\texttt{bk.}\hspace{-2pt}}\ instead of the plain
dot between volume number and part number (harley:cartography,
lach:asia).  If the field contains something other than a number,
\textsf{biblatex-chicago} will print it as is, capitalizing it if
necessary, rather than supplying the usual bibstring, so this provides
a mechanism for altering the string to your liking.  The field will be
printed in the same place in any entry as would a \textsf{volume}
number, and although it will most usually be associated with such a
number, it can also function independently, allowing you to identify
parts of works that don't fit into the standard scheme.  If you need
to identify \enquote{parts} or \enquote{books} that are part of a
published \textsf{series}, for example, then you'll need to use a
different field, (which in the case of a series would be
\textsf{number} [palmatary:pottery]).  Cf.\ \textsf{volume};
iso:electrodoc.

\mybigspace Standard \mymarginpar{\textbf{publisher}}
\textsf{biblatex} field.  Remember that \enquote{\texttt{and}} is a
keyword for connecting multiple publishers, so if a publisher's name
contains \enquote{and,} then you should either use the ampersand (\&)
or enclose the whole name in additional braces.  (See \emph{Manual}
14.133--41; aristotle:metaphy:gr, cohen:schiff, creasey:ashe:blast,
dunn:revolutions.)

\mylittlespace There are, as one might expect, a few further
subtleties involved here.  If you give two publishers in the field
they will both be printed, separated by a forward slash in both notes
and bibliography (14.90; sereny:cries).  The 17th edition generally is
rather keener than the 16th on using just one, particularly so in the
case when the parent company of an imprint is also listed on a title
page, in which case only the imprint need be included in your
apparatus (14.138).  If an academic publisher issues \enquote{certain
  books through a special publishing division or under a special
  imprint or as part of a publishing consortium (or joint imprint),}
this arrangement may be specified in the \textsf{publisher} field
(14.139; cohen:schiff).  If a book has two co-publishers \enquote{in
  different countries} (14.140), then the simplest thing to do is to
choose one, probably the nearest one geographically.  If you feel it
necessary to include both, then levistrauss:savage demonstrates one
way of doing so, using a combination of the \textsf{publisher} and
\textsf{location} fields.  If the work is self-published, you can
specify this in the \textsf{pubstate} field (see below), and any
commercial self-publishing platform would go in \textsf{publisher}
(14.137).  Books published before 1900 can, at your discretion,
include only the place (if known) and the date (14.128).  If for some
reason you need to indicate the absence of a publisher, the
abbreviation given by the \emph{Manual} is \texttt{n.p.}, though this
can also stand for \enquote{no place.}  The \emph{Manual} also
mentions {s.n.}\,(= \emph{sine nomine}) to specify the lack of a
publisher (10.42).

\mybigspace In \mymarginpar{\textbf{pubstate}} response to new
specifications in the 17th edition of the \emph{Manual} (esp.\
14.137), I have tried to generalize the functioning of the
\textsf{pubstate} field in all entry types.  Because the author-date
style has fairly complicated rules about presenting reprinted editions
(15.40), the \texttt{reprint} string still has a special status.
Depending on which date(s) you have chosen to appear at the head of
the entry, \textsf{biblatex-chicago-authordate} will either print the
(localized) string \texttt{reprint} in the proper place or otherwise
provide a notice at the end of the entry detailing the original
publication date.  See under \textbf{date} above for the available
permutations.  (Cf.\ aristotle:metaphy:gr, maitland:canon,
maitland:equity, schweitzer:bach.)

\mylittlespace Other strings are divided into two types: those which
\textsf{biblatex-chicago} will print as the \textsf{year}, which
currently means \emph{only} those for which \textsf{biblatex} contains
bibstrings indicating works soon to be published, i.e.,
\texttt{forthcoming}, \texttt{inpreparation}, \texttt{inpress}, and
\texttt{submitted}; and those, i.e., everything else, which will be
printed before, and in close association with, other information about
the publisher of a work.  (This \mymarginpar{\textbf{NB}} is a change
from previous behavior, where non-\texttt{reprint} strings were
printed \emph{after} the publication information, as in the standard
styles.  You can still use the \textsf{addendum} field to present
information here, of course.)  The four strings that replace the
\textsf{year} will always be localized, as will \texttt{reprint} and
\texttt{selfpublished} (and anything else that \textsf{biblatex} finds
to be a \cmd{bibstring}) from the second category.  All other strings
will be printed as-is, capitalized if needed, just before the
publisher (author:forthcoming, contrib:contrib, schweitzer:bach).

\mylittlespace There is one further subtlety of which you ought to be
aware.  In \textsf{music} and \textsf{video} entries, the
\texttt{reprint} string in \textsf{pubstate} will only make a
difference to your entries when the date which it modifies --- the
\textsf{origdate}, typically --- \emph{doesn't} appear in citations
and at the head of reference-list entries.  In this case the date is
treated as an original release date, and it will be printed, preceded
by the appropriate string, near the end of the entry.  Other strings
don't show this special behavior in these entries.

\mybigspace I \mymarginpar{\textbf{redactor}} have implemented this
field just as \textsf{biblatex}'s standard styles do, even though the
\emph{Manual} doesn't actually mention it.  It may be useful for some
purposes.  Cf.\ \textsf{annotator} and \textsf{commentator}.

\mybigspace See \mymarginpar{\textbf{reprinttitle}}
section~\ref{sec:authrelated}, below.

% %\enlargethispage{\baselineskip}

\mybigspace A \mymarginpar{\textbf{series}} standard \textsf{biblatex}
field, usually just a number in an \textsf{article},
\textsf{periodical}, or \textsf{review} entry, almost always the name
of a publication series in \textsf{book}-like entries, and providing
similar identifying information associated with a \textsf{number} in
\textsf{music} and \textsf{standard} entries.  If you need to attach
further information to the \textsf{series} name in a
\textsf{book}-like entry, then the \textsf{number} field is the place
for it, whether it be a volume, a number, or even something like
\enquote{2nd ser.} or \enquote{\cmd{bibstring\{oldseries\}}.}  Of
course, you can also use \cmd{bibstring\{oldseries\}} or
\cmd{bibstring\{newseries\}} in an \textsf{article} entry, but there
you would place it in the \textsf{series} field itself.  (In fact, the
\textsf{series} field in \textsf{article} and \textsf{periodical}
entries is one of the places where \textsf{biblatex} allows you just
to use the plain bibstring \texttt{oldseries}, for example, rather
than making you type \verb+\bibstring{oldseries}+.  The \textsf{type}
field in \textsf{manual}, \textsf{patent}, \textsf{report}, and
\textsf{thesis} entries also has this auto-detection mechanism in
place; see the discussion of \cmd{bibstring} below for details.)  In
whatever entry type, these bibstrings produce the required
abbreviation.  (For books and similar entries, see \emph{Manual}
14.123--26; boxer:china, browning:aurora, palmatary:pottery,
plato:republic:gr, wauchope:ceramics; for periodicals, see 14.184;
garaud:gatine, sewall:letter.)  Cf.\ \textsf{number} for more
information on the \emph{Manual}'s preferences regarding the
formatting of numerals.

\mybigspace This \mymarginpar{\textbf{shortauthor}} is a standard
\textsf{biblatex} field, but \textsf{biblatex-chicago} makes
considerably greater use of it than the standard styles.  For the
purposes of the author-date specification, the field provides the name
to be used in text citations.  In the vast majority of cases, you
don't need to specify it, because the \textsf{biblatex} system selects
the author's last name from the \textsf{author} field and uses it in
such a reference, and if there is no \textsf{author} it will search
\textsf{namea}, \textsf{editor}, \textsf{nameb}, \textsf{translator},
and \textsf{namec}, in that order.  The current versions of
\textsf{biblatex} and \textsf{Biber} will automatically alphabetize by
any of these names if they appear at the head of an entry.  If, in an
author-less \textsf{article} entry (\textsf{entrysubtype}
\texttt{magazine}), you allow \textsf{biblatex-chicago} to use the
\textsf{journaltitle} as the author --- the default behavior --- and
you have been accustomed to using the \textsf{shortauthor} field to
abbreviate it, it may be simpler now to use the \textsf{shortjournal}
field instead, which does all of the formatting for you, and
additionally adds the possibility of printing a list of journal
abbreviations.  See just below for the details.  (Cf.\ gourmet:052006,
lakeforester:pushcarts, nyt:trevorobit, unsigned:ranke).  With long,
institutional authors, a shortened version in \textsf{shortauthor} may
save space in the running text (evanston:library), but see under
\textbf{shorthand} for another method of saving space.

\mylittlespace As mentioned under \textsf{editortype}, the
\emph{Manual} (15.36) recommends against providing the identifying
string (e.g., ed.\ or trans.)\ in text citations, and
\textsf{biblatex-chicago} follows their recommendation.  If you need
to provide these strings in such a citation, then you'll have to do so
by hand in the \textsf{shortauthor} field, or in the
\textsf{shorteditor} field, whichever you are using.

\mybigspace Like \mymarginpar{\textbf{shorteditor}}
\textsf{shortauthor}, a field to provide a name for a text citation,
in this case for, e.g., a \textsf{collection} entry that typically
lacks an author.  The \textsf{shortauthor} field works just as well in
most situations, but if you have set \texttt{useauthor=false} (and not
\texttt{useeditor=false}) in an entry's \textsf{options} field, then
only \textsf{shorteditor} will be recognized.  It may be worth
pointing out that, because \textsf{biblatex-chicago} also provides a
\textsf{namea} field for the editor of a \textsf{title} as opposed to
a \textsf{main-} or \textsf{booktitle}, and because in standard use
the \textsf{namea}, if present, will be chosen to head a reference
list entry before the \textsf{editor}, you should present the
shortened \textsf{namea} here instead of a shortened \textsf{editor}
in such cases.Cf.\ \textsf{editortype}, above.

\paragraph*{\protect\mymarginpar{\textbf{shorthand}}}
\label{sec:ad:shorthand}

This is \textsf{biblatex}'s mechanism for using abbreviations in
citations.  For \textsf{biblatex-chicago-authordate} I have modified
it somewhat to conform to the needs of the specification, though there
is a package option to revert the behavior to something closer to the
\textsf{biblatex} standard --- see below and under \texttt{cmslos} in
section~\ref{sec:authpreset}.  The main problem when presenting
readers with an abbreviation is to ensure that they know how to expand
it.  In the notes \&\ bibliography style this is accomplished with a
notice in the first footnote citing a given work, which explains that
henceforth the abbreviation will be used instead, and also, if needed,
with a list of shorthands that summarizes all the abbreviations used
in a particular text.  The first part of this system isn't available
in the author-date style of citation, and indeed these citations are
in themselves already highly-abbreviated keys to the fuller
information to be found in the list of references.  There are cases,
however, particularly when institutions or \textsf{journaltitles}
appear as authors, when you may feel the need to provide a shortened
version for citations.  I have already discussed two options available
to you just above (cf.\ \textbf{shortauthor} and
\textbf{shortjournal}).  For the former to work the abbreviation must
either be instantly recognizable to your readership or at least easily
parseable by them, while with the latter you can either rely on the
conventions of your field or, alternately, provide a list of journal
abbreviations using \cmd{printbiblist\{shortjournal\}}.

% %\enlargethispage{-\baselineskip}

\mylittlespace For long institutional names the \emph{Manual's}
recommendation (15.37) involves using an abbreviation which will
appear not only in citations but also at the head of the entry in the
list of references.  Such an entry should therefore be alphabetized by
the abbreviation, with its expansion placed (inside parentheses)
between the abbreviation and the date.  This formatting can be
produced in one of two ways: either you can provide a
specially-formatted \textsf{author} field (for the reference list, and
including both the abbreviation and the parenthesized expansion) + a
\textsf{shortauthor} (for the citations), or you can use a normal
\textsf{author} field + a \textsf{shorthand}, in which case
\textsf{biblatex-chicago-authordate} will automatically use the
\textsf{shorthand} in text citations and also place it at the head of
the reference list entry, followed by the \textsf{author} within
parentheses.  This method is simpler and more compatible with other
styles, and will also produce a list that is correctly sorted by the
\textsf{shorthand}.  (Cf.\ niso:bibref, bsi:abbreviation,
iso:electrodoc.)

\mylittlespace I should clarify here that this automatic placement of
the \textsf{shorthand} at the head of the entry will \emph{not} occur
if you set the package option \texttt{cmslos=false} in your preamble.
This allows you to implement other systems of shorthand expansion
using either a list of shorthands (via \cmd{printshorthands}, which is
always available no matter what the state of \texttt{cmslos}) or
cross-references (via \textsf{customc}) within the reference list
itself.  You can place \texttt{skiplos} in the \textsf{options} field
to exclude a particular entry from the list of shorthands if you do
decide to print that list, giving maximum flexibility.

\mylittlespace Indeed, I have provided two options to add to this
flexibility.  First, I have included two \texttt{bibenvironments} for
use with the \texttt{env} option to the \cmd{printshorthands} command:
\texttt{losnotes} is designed to allow a list of shorthands to appear
inside footnotes, while \texttt{losendnotes} does the same for
endnotes.  Their main effect is to change the font size, and in the
latter case to clear up some spurious punctuation and white space that
I see on my system when using endnotes.  (You'll probably also want to
use the option \texttt{heading=none} in order to get rid of the
[oversized] default, providing your own within the \cmd{footnote}
command.)  Second, I have provided a package option,
\texttt{shorthandfull}, which prints entries in the list of shorthands
which contain full bibliographical information, effectively allowing
you to eschew the list of references in favor of a fortified shorthand
list.  This option will only work if used in tandem with
\texttt{cmslos=false}, as otherwise the shorthand will be printed
twice.  (See 15.37, 13.67, 14.59--60, and also \textsf{biblatex.pdf}
for more information.)

\mylittlespace As I mentioned above under \textbf{crossref}, I believe
it is safe to use shorthands in parent entries, as this, in the
standard configuration, gives you the shorthand itself in the child
entry's abbreviated cross-reference, which may well save space in the
list of references.

\mybigspace A \mymarginpar{\textbf{shortjournal}} special
\textsf{biblatex} field, used to provide both an abbreviated form of a
\textsf{journaltitle} in citations and/or the reference list and to
facilitate the creation of a list of journal abbreviations, should
this be needed, rather in the manner of a \textsf{shorthand} list.  As
requested by user BenVB, you can now utilize this functionality in
your documents, but there are a few details worth mentioning here.
First, users in some fields may well already be accustomed to using a
set of standard journal abbreviations (15.46), in which case the
\textsf{journaltitle} field may well already contain the abbreviation,
which will appear wherever that field is printed.  In such cases, it
usually isn't necessary to provide a list of abbreviations in
individual publications, but were you to require such a thing, you'd
have to move the abbreviation from the \textsf{journaltitle} to the
\textsf{shortjournal} field, placing the full title in the former.  In
\textsf{periodical} entries the \textsf{title} field presents what
would be the \textsf{journaltitle} in the \textsf{articles} or
\textsf{reviews}, so in such entries you can provide the standard
\textsf{shorttitle} field to accompany the \textsf{title}, and
\textsf{biblatex-chicago} will automatically copy the
\textsf{shorttitle} into a \textsf{shortjournal}.

\mylittlespace Having done this, you then need to choose where to
print the \textsf{shortjournal}, which is controlled by the
\texttt{journalabbrev} option either in the preamble or in the
\textsf{options} field of individual .bib entries.  By default, and
taking account of the space-saving features of the author-date styles,
this option is set to \texttt{notes}, so your \textsf{shortjournal}
fields will be printed only in those citations where they appear in
place of an \textsf{author}.  There are three other settings:
\texttt{true} prints the shortened fields both in citations and in the
reference list, \texttt{bib} prints them only in the reference list,
and \texttt{false} ignores them.  Should you wish to present a list of
these abbreviations with their expansions, then you need to use the
\verb+\printbiblist{shortjournal}+ command, perhaps with a
\texttt{title} option to differentiate the list from any
\textsf{shorthand} list.  As with \textsf{shorthand} lists, I have
provided two \texttt{bibenvironments} for printing this list in foot-
or endnotes (\texttt{sjnotes} and \texttt{sjendnotes}, respectively),
to be used with the \texttt{env} option to \cmd{printbiblist}.  Again
as with \textsf{shorthands}, you'll probably want to use the option
\texttt{heading=none} when using these environments, just to turn off
the (oversized) default, and perhaps provide your own title within the
\cmd{footnote} command.  Finally, if you don't like the default
formatting of the abbreviations in the list (bold italic), you can
change it with \verb+\DeclareFieldFormat{shortjournal-width}+ --- you
can see its default definition at the top of
\textsf{chicago-authordate.bbx}.

\mybigspace A \mymarginpar{\textbf{shortseries}} special
\textsf{biblatex} field, used both to provide an abbreviated form of a
(book) \textsf{series} in a reference list and to facilitate the
creation of a list of such abbreviations rather in the manner of a
\textsf{shorthand} list.  As with the \textsf{shortjournal} field, its
inclusion in \textsf{biblatex-chicago} was requested by user BenVB,
and it is now available in entry types which have book-like series
titles rather than journal-like numbers in the \textsf{series} field,
to wit: \textsf{audio, book, bookinbook, collection, inbook,
  incollection, inproceedings, inreference, letter, manual, music,
  mvbook, mvcollection, mvproceedings, mvreference, reference, report,
  standard, suppbook,} and \textsf{video}.  There are several steps to
take in order to use the field.  First, you'll need to provide both
\textsf{shortseries} and \textsf{series} fields in the entry, then
you'll need to set the \texttt{seriesabbrev} option either when
loading \textsf{biblatex-chicago}, for the whole document or for
specific entry types, or in the \textsf{options} field of individual
.bib entries.  By default, this option is not set, so your
\textsf{shortseries} fields will be silently ignored.  Setting it to
\texttt{true} prints the shortened fields in the reference list.
Should you wish to present a list of these abbreviations with their
expansions, then you need to use the \verb+\printbiblist{shortseries}+
command, perhaps with a \texttt{title} option to differentiate the
list from any \textsf{shorthand} list.  As with \textsf{shorthand}
lists, I have provided two \texttt{bibenvironments} for printing this
list in foot- or endnotes (\texttt{shsernotes} and
\texttt{shserendnotes}, respectively), to be used with the
\texttt{env} option to \cmd{printbiblist}.  Again as with
\textsf{shorthands}, you'll probably want to use the option
\texttt{heading=none} when using these environments, just to turn off
the (oversized) default, and perhaps provide your own title within the
\cmd{footnote} command.  Finally, if you don't like the default
formatting of the abbreviations in the list (plain roman), you can
roll your own using \verb+\DeclareFieldFormat{shortserieswidth}+ ---
you can see its default definition at the top of
\textsf{chicago-authordate.bbx}.

\mybigspace A \mymarginpar{\textbf{shorttitle}} standard
\textsf{biblatex} field, primarily used to provide an abbreviated
title for citation styles that need one.  (It is also the way to hook
\textsf{periodical} entries into the \textsf{shortjournal} mechanism,
on which see the previous entry.)  In
\textsf{biblatex-chicago-authordate} such a field will be necessary
only very rarely (unlike in the notes \&\ bibliography style), and is
most likely to turn up in \textsf{inreference} or \textsf{reference}
entries (where the \textsf{title} takes the place of the
\textsf{author}), in \textsf{dataset} entries, or in any sort of entry
with a \texttt{classical} \textsf{entrysubtype} or with
\texttt{authortitle} set in its \textsf{options} field.  These latter
three contexts make citations use \textsf{author} and \textsf{title}
instead of \textsf{author} and \textsf{year}, and if an abbreviated
version of that title would save space in your running text this is
the field where you can provide it.  (Cf.\ ency:britannica,
grove:sibelius, aristotle:metaphy:gr.)

\mybigspace Standard \mymarginpar{\textbf{sortkey} \\\textbf{sortname}
  \\\textbf{sorttitle} \\\textbf{sortyear}} \textsf{biblatex} fields,
designed to allow you to specify how you want an entry alphabetized in
a list of references.  The \textsf{sortkey} field trumps all other
sorting information, while the others offer more fine-grained control.
In general, if an entry doesn't turn up where you expect or want it,
one of these fields should provide the solution.  Entries with a
corporate author can omit the definite or indefinite article, which
should help (14.70, 14.84; cotton:manufacture, nytrumpet:art).  The
default settings of \cmd{DeclareSortingTemplate} include the three
supplemental name fields (\textsf{name[a-c]}) and also the
\textsf{journaltitle} in the sorting algorithm, so once again you
should find those algorithms needing less help than before.  Entries
headed by a \textsf{title} beginning with the definite or indefinite
article may well still require such assistance (grove:sibelius).
There may be circumstances --- several reprinted books by the same
author, for example --- when the \textsf{sortyear} field is the best
choice.  Please consult \textsf{biblatex.pdf} for the details.

\mybigspace The \mymarginpar{\textbf{subtitle}} subtitle for a
\textsf{title} --- see next entry.

\mybigspace This \mymarginpar{\textbf{timezone}} field can, if
necessary, specify the time zone associated with a time stamp given as
part of an \textsf{date}.  The \emph{Manual} prefers initialisms like
\enquote{EST} for this purpose, and you can provide parentheses around
it at your discretion (cp.\ 10.41 and 14.191).

\mybigspace \textsf{Biblatex-chicago} \mymarginpar{\textbf{title}}
includes the \textsf{authordate-trad} style, designed as a kind of
hybrid style according to indications contained in the \emph{Manual}
(15.38).  This \textsf{trad} style differs \emph{only} in the way it
treats the \textsf{title} and related fields, which retain the forms
they have traditionally had in the Chicago author-date specifications
prior to the 16th edition.  Where newer editions use headline-style
capitalization, the older editions used sentence-style; where newer
editions place \textsf{article} or \textsf{incollection}
\textsf{titles} within quotation marks, the older editions presented
them in plain text.  I include below, under a separate rubric, full
documentation of \textsf{trad} \textsf{title} fields for those needing
or wishing to use them.  First, though, I document the same field(s)
for the standard author-date style.

\mylittlespace In the vast majority of cases, this field works just as
it always has in \textsc{Bib}\TeX, and just as it does in
\textsf{biblatex}.  The \emph{Manual} recommends that \textsf{titles}
be treated more or less identically across both its systems of
documentation (15.3, 15.6, 15.13).  This means that users of the
author-date style don't need to worry about sentence-style
capitalization when compiling their .bib databases, and so can eschew
the extra curly braces needed to preserve uppercase letters in this
context.  These rules, however, mean that a few complications familiar
to users of the notes \&\ bibliography style do arise.  First,
although nearly every entry will have a \textsf{title}, there are some
exceptions, particularly \textsf{incollection} or \textsf{online}
entries with a merely generic title, instead of a specific one
(centinel:letters, powell:email).  Second, the \emph{Manual}'s rules
for formatting \textsf{titles}, which also hold for
\textsf{booktitles} and \textsf{maintitles}, require additional
attention.  The whole point of using a \textsf{biblatex}-based system
is for it to do the formatting for you, and in most cases
\textsf{biblatex-chicago-authordate} does just that, surrounding
titles with quotation marks, italicizing them, or occasionally just
leaving them alone.  When, however, a title is quoted within a title,
then you need to know some of the rules.  A summary here should serve
to clarify them, and help you to understand when
\textsf{biblatex-chicago-authordate} might need your help in order to
comply with them.

\mylittlespace The internal rules of
\textsf{biblatex-chicago-authordate} are as follows:

\begin{description}
\item[\qquad Italics:] \textsf{booktitle}, \textsf{maintitle}, and
  \textsf{journaltitle} in all entry types; \textsf{title} of
  \textsf{artwork}, \textsf{book}, \textsf{bookinbook},
  \textsf{booklet}, \textsf{collection}, \textsf{image},
  \textsf{manual}, \textsf{misc} (with no \textsf{entrysubtype}),
  \textsf{performance}, \textsf{periodical}, \textsf{proceedings},
  \textsf{report}, \textsf{standard}, \textsf{suppbook}, and
  \textsf{suppcollection} entry types.
\item[\qquad Quotation Marks:] \textsf{title} of \textsf{article},
  \textsf{inbook} \textsf{incollection}, \textsf{inproceedings},
  \textsf{online}, \textsf{periodical}, \textsf{thesis}, and
  \textsf{unpublished} entry types, \textsf{issuetitle} in
  \textsf{article}, \textsf{periodical}, and \textsf{review} entry
  types.
\item[\qquad Unformatted:] \textsf{booktitleaddon},
  \textsf{maintitleaddon}, and \textsf{titleaddon} in all entry types,
  \textsf{title} of \textsf{customc}, \textsf{letter}, \textsf{misc}
  (with an \textsf{entrysubtype}), \textsf{patent}, \textsf{review},
  and \textsf{suppperiodical} entry types.
\item[\qquad Italics or Quotation Marks:] All of the audiovisual entry
  types --- \textsf{audio}, \textsf{music}, and \textsf{video} ---
  have to serve as analogues both to \textsf{book} and to
  \textsf{inbook}.  Therefore, if there is both a \textsf{title} and a
  \textsf{booktitle}, then the \textsf{title} will be in quotation
  marks.  If there is no \textsf{booktitle}, then the \textsf{title}
  will be italicized, unless you provide an \textsf{entrysubtype}.
\end{description}

Now, the rules for which entry type to use for which sort of work tend
to be fairly straightforward, but in cases of doubt you can consult
section~\ref{sec:types:authdate} above, the examples in
\textsf{dates-test.bib}, or go to the \emph{Manual} itself,
8.156--201.  Assuming, then, that you want to present a title within a
title, and you know what sort of formatting each of the two would, on
its own, require, then the following rules apply:

\begin{enumerate}
\item Inside an italicized title, all other titles are enclosed in
  quotation marks and italicized, so in such cases all you need to do
  is provide the quotation marks using \cmd{mkbibquote}, which will
  take care of any following punctuation that needs to be brought
  within the closing quotation mark(s) (14.94; donne:var,
  mchugh:wake).
\item Inside a quoted title, you should present another title as it
  would appear if it were on its own, so in such cases you'll need to
  do the formatting yourself.  Within the double quotes of the title
  another quoted title would take single quotes --- the
  \cmd{mkbibquote} command does this for you automatically, and also,
  I repeat, takes care of any following punctuation that needs to be
  brought within the closing quotation mark(s).  (See 14.94--95;
  garrett, loften:hamlet, murphy:silent, white:callimachus.)
\item Inside a plain title (most likely in a \textsf{review} entry or
  a \textsf{titleaddon} field), you should present another title as it
  would appear on its own, once again formatting it yourself using
  \cmd{mkbibemph} or \cmd{mkbibquote}.  (barcott:review, gibbard,
  osborne\hc poison, ratliff:review, unsigned:ranke).
\end{enumerate}

The \emph{Manual} provides a few more rules, as well.  A word normally
italicized in text should also be italicized in a quoted or plain-text
title, but should be in roman (\enquote{reverse italics}) in an
italicized title.  A quotation used as a (whole) title (with or
without a subtitle) retains, according to the 16th edition, its
quotation marks in an italicized title if it appears that way in the
source, but I can't find similar instructions in the 17th.  Such a
quotation always retains its quotation marks when the surrounding
title is quoted or plain (14.94; lewis).  A word or phrase in
quotation marks, but that isn't a quotation, retains those marks in
all title types (kimluu:diethyl).

\mylittlespace Finally, please note that in all \textsf{review} (and
\textsf{suppperiodical}) entries, and in \textsf{misc} entries with an
\textsf{entrysubtype}, and only in those entries,
\textsf{biblatex-chicago-authordate} will automatically capitalize the
first word of the \textsf{title} after sentence-ending punctuation,
assuming that such a \textsf{title} begins with a lowercase letter in
your .bib database.  See\,\textbf{\textbackslash autocap} in
section~\ref{sec:formatting:authdate} below for more details.

\mybigspace When \mymarginpar{\textbf{title (trad)}} you choose the
\textsf{authordate-trad} style, your \textsf{title} and related fields
will need extra care, familiar to users of the 15th-edition
author-date style.  The whole point of using a \textsf{biblatex}-based
system is for it to do the formatting for you, and in most cases
\textsf{biblatex-chicago-authordate-trad} does just that, capitalizing
titles sentence-style, italicizing them, and sometimes both.  There
are two situations that require user intervention.  First, in titles
that take sentence-style capitalization, you need, as always in
traditional \textsc{Bib}\TeX, to assist the algorithms by placing
anything that needs to remain capitalized within an extra pair of
curly braces.  Second, when a title is quoted within a title, you need
to know some of the rules of the Chicago style.  A summary here should
serve to clarify them, and help you to understand when
\textsf{biblatex-chicago-authordate-trad} might need your help in
order to comply with them.

%%\enlargethispage{2\baselineskip}

\mylittlespace With regard to sentence-style capitalization, the rules
of the Chicago \textsf{authordate-trad} style are fairly simple:

\begin{description}
\item[\qquad Headline Style:] \textsf{journaltitle} in all types,
  \textsf{series} in all \textsf{book}-like entries (i.e., not in
  \textsf{articles}), and \textsf{title} in \textsf{periodical}
  entries.
\item[\qquad Sentence Style:] every other \textsf{title},
  \emph{except} in \textsf{letter} entries, \textsf{review} and
  \textsf{suppperiodical} entries, and in \textsf{misc} entries with
  an \textsf{entrysubtype}.  Also, the \textsf{booktitle},
  \textsf{issuetitle}, and \textsf{maintitle} in all entry types use
  sentence style.
\item[\qquad Contextual Capitalization of First Word:]
  \textsf{titleaddon}, \textsf{booktitleaddon},
  \textsf{maintitleaddon} in all entry types, also the \textsf{title}
  of \textsf{review} entries, of \textsf{suppperiodical} entries, and
  of \textsf{misc} entries with an \textsf{entrysubtype}.
\item[\qquad Plain:] \textsf{title} in \textsf{letter} entries.
\end{description}

What this means in practice is that to get a title like \emph{The
  Chicago manual of style}, your .bib entry needs to have a field that
looks something like this:

\mylittlespace\hspace*{1em}\texttt{title = \{The \{Chicago\} Manual
  of Style\}}

\mylittlespace This is completely straightforward, but remember that
if an \textsf{article} has a title like: Review of \emph{The Chicago
  manual of style}, then the curly braces enclosing material to be
formatted in italics will cause the capitalization algorithm to stop
and leave all of that material as it is, so your .bib entry would need
to have a field something like this:

\mylittlespace\hspace*{1em}\texttt{title = \{review of}
\cmd{mkbibemph\{The Chicago manual of style\}\}}

\mylittlespace (As\colmarginpar{\textsf{biblatex} 3.20} an aside, a
\cmd{bibstring} at the start of a title field no longer works as it
used to in the \textsf{trad} style, due to changes in the handling of
the \cmd{MakeSentenceCase} command.  You can still use them if you set
the \texttt{casechanger} option to \texttt{latex2e}, but otherwise you
can just provide the actual string, as here.  The lowercase letter at
the start isn't strictly necessary unless you're going to cite the
work by title instead of by author-date.  [Compare the notes \&\
bibliography style, and also the \textsf{*titleaddon} fields.])

\mylittlespace With regard to italics, the rules of
\textsf{biblatex-chicago-authordate-trad} are as follows:

\begin{description}
\item[\qquad Italics:] \textsf{booktitle}, \textsf{maintitle}, and
  \textsf{journaltitle} in all entry types; \textsf{title} of
  \textsf{artwork}, \textsf{book}, \textsf{bookinbook},
  \textsf{booklet}, \textsf{collection}, \textsf{manual},
  \textsf{misc} (w/o \textsf{entrysubtype}), \textsf{performance},
  \textsf{periodical}, \textsf{proceedings}, \textsf{report},
  \textsf{standard}, \textsf{suppbook}, and \textsf{suppcollection}
  types.
\item[\qquad Main Text Font (Roman):] \textsf{title} of
  \textsf{article}, \textsf{image}, \textsf{inbook},
  \textsf{incollection}, \textsf{inproceedings}, \textsf{letter},
  \textsf{misc} (with an \textsf{entrysubtype}), \textsf{online},
  \textsf{patent}, \textsf{periodical}, \textsf{review},
  \textsf{suppperiodical}, \textsf{thesis}, and \textsf{unpublished}
  entry types, \textsf{issuetitle} in \textsf{article} and
  \textsf{periodical} entry types.  \textsf{booktitleaddon},
  \textsf{maintitleaddon}, and \textsf{titleaddon} in all entry types.
\item[\qquad Italics or Roman:] All of the audiovisual entry types ---
  \textsf{audio}, \textsf{music}, and \textsf{video} --- have to serve
  as analogues both to \textsf{book} and to \textsf{inbook}.
  Therefore, if there is both a \textsf{title} and a
  \textsf{booktitle}, then the \textsf{title} will be in the main text
  font.  If there is no \textsf{booktitle}, then the \textsf{title}
  will be italicized, unless you provide an \textsf{entrysubtype}.
\end{description}

Now, the rules for which entry type to use for which sort of work tend
to be fairly straightforward, but in cases of doubt you can consult
section~\ref{sec:types:authdate} above, the examples in
\textsf{dates-test.bib}, or go to the \emph{Manual} itself,
8.156--201.  Assuming, then, that you want to present a title within a
title, and you know what sort of formatting each of the two would, on
its own, require, then the following rules apply:

\begin{enumerate}
\item Inside an italicized title, all other titles are enclosed in
  quotation marks and italicized, so in such cases all you need to do
  is provide the quotation marks using \cmd{mkbibquote}, which will
  take care of any following punctuation that needs to be brought
  within the closing quotation mark(s) (14.94; donne:var,
  mchugh:wake).
\item Inside a plain-text title, you should set off other plain-text
  titles with quotation marks, while italicized titles should appear
  as they would if they were on their own.  In such cases you'll need
  to do the formatting yourself, using \cmd{mkbibemph} or
  \cmd{mkbibquote}.  (See barcott:review, garrett, gibbard,
  loften:hamlet, loomis\hc structure, murphy:silent, osborne:poison,
  ratliff:review, unsigned:ranke, white\hc callimachus.)
\end{enumerate}

The \emph{Manual} provides a few more rules, as well.  A word normally
italicized in text should also be italicized in a plain-text title,
but should be in roman (\enquote{reverse italics}) in an italicized
title.  A quotation used as a (whole) title (with or without a
subtitle) retains, according to the 16th edition, its quotation marks
in an italicized title if it appears that way in the source, but I
can't find similar instructions in the 17th.  Such a quotation always
retains its quotation marks when the surrounding title is quoted or
plain (14.94; lewis).  A word or phrase in quotation marks, but that
isn't a quotation, retains those marks in all title types
(kimluu:diethyl).

\mylittlespace Finally, please note that there is also a preamble
option --- \texttt{headline} --- that disables the automatic
sentence-style capitalization routines in \textsf{authordate-trad}.
If you set this option, the word case in your title fields will not be
changed in any way, that is, this doesn't automatically transform your
titles into headline-style, but rather allows the .bib file to
determine capitalization.  It works by redefining the command
\cmd{MakeSentenceCase}, so in the unlikely event you are using the
latter anywhere in your document please be aware that it will also be
turned off there.  See section~\ref{sec:authuseropts}, below.

\mybigspace Standard \colmarginpar{\textbf{titleaddon}}
\textsf{biblatex} intends this field for use with additions to titles
that may need to be formatted differently from the titles themselves,
and \textsf{biblatex-chicago} uses it in just this way, with the
additional wrinkle that it can, if needed, replace the \textsf{title}
entirely, and this in, effectively, any entry type, providing a fairly
powerful, if somewhat complicated, tool for getting \textsf{biblatex}
to do what you want (cf.\ centinel:letters).  This field will always
be unformatted, that is, neither italicized nor placed within
quotation marks, so any formatting you may need within it you'll need
to provide manually yourself.  The single exception to this rule is
when your data begins with a word that would ordinarily only be
capitalized at the beginning of a sentence, in which case you need
then simply ensure that that word is in lowercase, and
\textsf{biblatex-chicago} will automatically do the right thing.
See\,\textbf{\textbackslash autocap} in
section~\ref{sec:formatting:authdate} below.  The package and entry
options \texttt{ptitleaddon} and \texttt{ctitleaddon}
(section~\ref{sec:authpreset}) can help you customize the punctuation
that appears before the \textsf{titleaddon} field.  Please note,
however, that I have added this field to the \textsf{periodical} entry
type, and that the punctuation there is governed by the
\mycolor{\texttt{jtitleaddon}} option, which defaults to a
\texttt{space}.  (Cf.\ brown:bremer, osborne:poison, reaves:rosen, and
white:ross:memo for examples where the field starts with a lowercase
letter; morgenson:market provides an example where the
\textsf{titleaddon} field, holding the name of a regular column in a
newspaper, is capitalized, a situation that is handled as you would
expect; coolidge:speech shows an entry option for controlling the
punctuation.)

\mybigspace As \mymarginpar{\textbf{translator}} far as possible, I
have implemented this field as \textsf{biblatex}'s standard styles do,
but the requirements specified by the \emph{Manual} present certain
complications that need explaining.  \textsf{Biblatex.pdf} points out
that the \textsf{translator} field will be associated with a
\textsf{title}, a \textsf{booktitle}, or a \textsf{maintitle},
depending on the sort of entry.  More specifically,
\textsf{biblatex-chicago} associates the \textsf{translator} with the
most comprehensive of those titles, that is, \textsf{maintitle} if
there is one, otherwise \textsf{booktitle}, otherwise \textsf{title},
if the other two are lacking.  In a large number of cases, this is
exactly the correct behavior (adorno:benj, centinel:letters,
plato:republic:gr, among others).  Predictably, however, there are
numerous cases that require, for example, an additional translator for
one part of a collection or for one volume of a multi-volume work.
For these cases I have provided the \textsf{nameb} field.  You should
format names for this field as you would for \textsf{author} or
\textsf{editor}, and these names will always be associated with the
\textsf{title} (euripides:orestes).  In the algorithm for finding a
name for the head of a reference list entry or for a citation,
\textsf{nameb} takes precedence over \textsf{translator}.

\mylittlespace I have also provided a \textsf{namea} field, which
holds the editor of a given \textsf{title} (euripides\hc orestes).  If
\textsf{namea} and \textsf{nameb} are the same,
\textsf{biblatex-chicago} will concatenate them, just as
\textsf{biblatex} already does for \textsf{editor},
\textsf{translator}, and \textsf{namec} (i.e., the compiler).
Furthermore, it is conceivable that a given entry will need separate
translators for each of the three sorts of title.  For this, and for
various other tricky situations, there is the \cmd{parttrans} macro
(and its siblings), designed to be used in a \textsf{note} field or in
one of the \textsf{titleaddon} fields (ratliff:review).  (Because the
strings identifying a translator differ in notes and bibliography, one
can't simply write them out in such a field when using the notes \&\
bibliography style, but you can certainly do so in the author-date
styles, if you wish.  Using the macros will make your .bib file more
portable across both Chicago specifications, and also across multiple
languages, but they are otherwise unnecessary.  [See
section~\ref{sec:international}].)

\mylittlespace Finally, as I detailed above under \textbf{author}, in
the absence of an \textsf{author}, a \textsf{namea}, an
\textsf{editor}, and a \textsf{nameb}, the \textsf{translator} will be
used at the head of an entry (silver:gawain), and the reference list
entry alphabetized by the translator's name, behavior that can be
controlled with the \texttt{use<name>} switches in the
\textsf{options} field.  Cf.\ \textsf{author}, \textsf{editor},
\textsf{namea}, \textsf{nameb}, and \textsf{namec}.

\mybigspace This \mymarginpar{\textbf{type}} is a standard
\textsf{biblatex} field, and in its normal usage serves to identify
the type of a \textsf{manual}, \textsf{patent}, \textsf{report}, or
\textsf{thesis} entry.  \textsf{Biblatex} implements the possibility,
in some circumstances, to use a bibstring without inserting it in a
\cmd{bibstring} command, and in some entry types (\textsf{audio,
  manual, music, patent, report, suppbook, suppcollection, thesis,}
and \textsf{video}) the \textsf{type} field works this way, allowing
you simply to input, e.g., \texttt{patentus} rather than
\verb+\bibstring{patentus}+, though both will work.  (See
petroff:impurity; herwign:office, murphy:silent, and ross:thesis all
demonstrate how the \textsf{type} field may sometimes be automatically
set in such entries by using one of the standard entry-type aliases).
In other entry types (\textsf{artwork, image, book, online, article,
  review,} and \textsf{suppperiodical}) \textsf{biblatex-chicago} will
merely capitalize the contents according to context.

\mylittlespace Another use for the field is to generalize the
functioning of the \textsf{suppbook} entry type, and of its alias
\textsf{suppcollection}.  In such entries, the \textsf{type} field can
specify what sort of supplemental material you are citing, e.g.,
\enquote{\texttt{preface to}} or \enquote{\texttt{postscript to}.}
Cf.\ \textsf{suppbook} above for the details.  (See \emph{Manual}
14.110; polakow:afterw, prose:intro).

\mylittlespace You can use the \textsf{type} field in
\textsf{artwork}, \textsf{audio}, \textsf{image}, \textsf{music}, and
\textsf{video} entries to identify the medium of the work, e.g.,
\texttt{oil on canvas}, \texttt{albumen print}, \texttt{compact disc}
or \texttt{MPEG}.  In \textsf{book} entries it will normally hold
system information about multimedia app content, while in
\textsf{online, article,} and \textsf{review} entries it will hold the
medium of online multimedia (15.57, 14.267--68).  Cf.\ under these
entry types in section~\ref{sec:types:authdate}, above, for more
details.  (See auden:reading, bedford:photo, cleese:holygrail,
leo:madonna, nytrumpet:art.)

\mybigspace Standard \mymarginpar{\textbf{url}} \textsf{biblatex}
field, it holds the url of an online publication, though you can
provide one for all entry types.  The \emph{Manual} expresses a strong
preference for DOIs over URLs if the former is available --- cf.\
\textsf{doi} above, and also \textsf{urldate} just below.  The
required \LaTeX\ package \textsf{url} will ensure that your documents
format such references properly, in the text and in the reference
apparatus.  It may be worth noting that child entries no longer
inherit \textsf{url} fields from their parents --- the information
seems entry-specific enough to warrant a little bit of extra typing if
you need to present the same locator in several entries.

\mybigspace Standard \mymarginpar{\textbf{urldate}} \textsf{biblatex}
field, it identifies exactly when you accessed a given url.  The
\emph{Manual} prefers DOIs to URLs; in the latter case it allows the
use of access dates, particularly in contexts that require it, but
prefers that you use revision dates, if these are available.  To
enable you to specify which date is at stake, I have provided the
\textbf{userd} field, documented below.  If an entry doesn't have a
\textsf{userd}, then the \textsf{urldate} will be treated as an access
date (14.8, 14.12--13, 15.50; evanston:library, grove:sibelius,
hlatky:hrt, osborne:poison, sirosh:visualcortex, wikiped:bibtex).  In
the default setting of \cmd{DeclareLabeldate}, any entry without a
\textsf{date}, \textsf{eventdate}, or \textsf{origdate} will use the
\textsf{urldate} to find a year for citations and the list of
references (grove:sibelius, wikiped:bibtex), but \emph{only} if the
\textsf{urldate} isn't an access date, that is, only if a
\textsf{userd} field is present.  If the only date available is an
online access date, then the entry is considered to have no date, and
\enquote{n.d.} will appear instead, though of course the access date
will still be printed later in the reference list entry.  (If you were
to put the string \texttt{accessed} into the \textsf{userd} field, you
could work around this prohibition.)

\mylittlespace You can also use the \textsf{urldate} field to specify
a time stamp, should the date alone not be specific enough.  The time
stamp follows the date, separated by an uppercase \enquote{T}, like
so: \texttt{yyyy-mm-dd\textbf{T}hh:mm:ss}.  If you wish to specify the
time zone, the \emph{Manual} (10.41) prefers initialisms like
\enquote{EST} or \enquote{PDT,} and these are most easily provided
using the \texttt{urltimezone} field, where you can provide your own
parentheses if so desired (cp.\ 14.191).  Following the examples in
the \emph{Manual}, any \textsf{urldate} will by default be printed in
24-hour format, though other time stamps use 12-hour format.  The
\textsf{biblatex} option \texttt{urltime}, discussed in
section~\ref{sec:preset:authdate}, allows you to change this in your
preamble.

\mylittlespace A \textsf{urldate} time stamp (and
\texttt{urltimezone}) can appear in any entry whatsoever, if you judge
the online source to be the sort that changes rapidly enough for a
time stamp to be necessary (14.207, 14.233; wikiped:bibtex).  You can
stop it printing by setting the new \mymarginpar{\texttt{urlstamp}}
\texttt{urlstamp} option to \texttt{false} in your preamble for the
whole document or for specified entry types, or in the
\textsf{options} field of individual entries.  Please see the
documentation of \textbf{date} and also table~\ref{ad:date:extras},
above, for more details about time stamps and other parts of
\textsf{biblatex's} enhanced date specifications.
Table~\ref{tab:online:adtypes} contains a summary of the current state
of \textsf{biblatex-chicago's} handling of online materials.

\mybigspace This \mymarginpar{\textbf{urltimezone}} field can, if
necessary, specify the time zone associated with a time stamp given as
part of an \textsf{urldate}.  The \emph{Manual} prefers initialisms
like \enquote{EST} for this purpose, and you can provide parentheses
around it at your discretion (cp.\ 10.41 and 14.191).

%%\enlargethispage{\baselineskip}

\mybigspace A \mymarginpar{\textbf{usera}} supplemental
\textsf{biblatex} field which in certain contexts in
\textsf{biblatex-chicago} will identify the broadcast network when you
cite a radio or television program.  In \textsf{article},
\textsf{periodical}, and \textsf{review} entries with
\textsf{entrysubtype} \texttt{magazine}, it acts almost as a
\enquote{\textsf{journaltitleaddon}} field, and its contents will be
placed, unformatted and between commas, after the
\textsf{journaltitle} and before the \textsf{date}.  In \textsf{video}
entries it comes after the \textsf{eventdate}, i.e., the date of first
broadcast, and is separated from that date by the \cmd{bibstring}
\enquote{\texttt{on}} (14.213, 14.265; american:crime, bundy:macneil,
friends:leia, mayberry:brady).

\mybigspace I \mymarginpar{\textbf{userc}} have implemented this
supplemental \textsf{biblatex} field as part of the Chicago
author-date style's handling of cross-references within the list of
references.  (The \enquote{c} part is meant as a sort of mnemonic for
this latter function.)  In recent editions of the \emph{Manual} you no
longer need to use the \textbf{customc} entry type to include
alphabetized expansions of \textsf{shorthands} in the reference list,
but you may still need to provide cross-references of some sort to
separate entries in that list, perhaps when a single author uses
multiple pseudonyms.  In such a case it is unlikely that you will cite
the \textsf{customc} entry itself in the body of your text.
Therefore, in order for it to appear in the reference list, you have
two choices.  You can either include the entry key of the
\textsf{customc} entry in a \cmd{nocite} command inside your document,
or you can place that entry key in the \textsf{userc} field of the
.bib entry that actually contains one of the full citations.  In the
latter case, \textsf{biblatex-chicago} will call \cmd{nocite} for you
when you cite the main entry.  (See 14.81--82; creasey:ashe:blast,
creasey:morton:hide, creasey:york:death, lecarre:quest.)

\mybigspace The \mymarginpar{\textbf{userd}} \textsf{userd} field acts
as a sort of \enquote{\textsf{datetype}} field, allowing you in most
entry types to identify whether a \textsf{urldate} is an access date
or a revision date.  The general usage is fairly simple.  If this
field is absent, then a \textsf{urldate} will be treated as an access
date, as has long been the default in \textsf{biblatex} and in
\textsf{biblatex-chicago}.  If you need to identify it in any other
way, what you include in \textsf{userd} will be printed \emph{before}
the \textsf{urldate}, so phrases like \enquote{\texttt{last modified}}
or \enquote{\texttt{last revised}} are what the field will typically
contain (14.12--13; wikiped:bibtex).  In the absence of a
\textsf{urldate}, you can in most entry types include a \textsf{userd}
field to qualify a \textsf{date} in the same way it would have
modified a \textsf{urldate}.  If an entry contains \emph{only} a
\textsf{urldate} and no other sort of date, and has no \textsf{userd}
field, that entry will now be treated as though it had no date, and
\enquote{n.d.} will appear in citations and at the head of entries in
the reference list (15.50).

\mylittlespace Because of the rather specialized needs of some
audio-visual references, this basic schema changes for \textsf{music}
and \textsf{video} entries.  In \textsf{music} entries where an
\textsf{eventdate} is present, \textsf{userd} will modify that date
instead of any \textsf{urldate} that may also be present, and it will
modify an \textsf{origdate} if it is present and there is no
\textsf{eventdate}.  It will modify a \textsf{date} only in the
absence of the other three.  In \textsf{video} entries it will modify
an \textsf{eventdate} if it is present, and in its absence the
\textsf{urldate}.  Given the absence of those two, it can modify a
\textsf{date}.  In all these cases, \textsf{userd} will modify what
remains of any date, i.e., the month and the day, if that date's year
has been printed at the head of the entry.  Please see the
documentation of the \textsf{music} and \textsf{video} entry types,
and especially of the \textsf{eventdate}, \textsf{origdate}, and
\textsf{urldate} fields, above (14.276--279, 15.53; nytrumpet:art).

\mylittlespace In all cases, you can start the \textsf{userd} field
with a lowercase letter, and \textsf{biblatex} will take care of
automatic contextual capitalization for you.

\mybigspace Another \mymarginpar{\textbf{usere}} supplemental
\textsf{biblatex} field, which \textsf{biblatex-chicago} uses
specifically to provide a translated \textsf{title} of a work,
something that may be needed if you deem the original language
unparseable by a significant portion of your likely readership.  The
\emph{Manual} offers two alternatives in such a situation: either you
can translate the title and use that translation in your
\textsf{title} field, providing the original language in
\textsf{language}, or you can give the original title in
\textsf{title} and the translation in \textsf{usere}.  Cf.\
\textbf{language}, above.  (See 14.99; kern, pirumova:russian,
weresz.)

\mybigspace See \mymarginpar{\textbf{userf}}
section~\ref{sec:authrelated}, below.

\mybigspace Standard \mymarginpar{\textbf{venue}} \textsf{biblatex}
offers this field for use in \textsf{proceedings} and
\textsf{inproceedings} entries, and after a request from Patrick
Danilevici I have followed suit.  I have also implemented the field in
the \textsf{misc} entry type, both with and without an
\textsf{entrysubtype}, in the \textsf{performance} type, and in the
\textsf{unpublished} type.  In all uses it will normally present the
actual venue of an event, as opposed, e.g., to the
\textsf{origlocation}, which might present where a letter was written
or where an earlier edition was printed.

\mybigspace Author-date \mymarginpar{\textbf{verbc}} styles in
\textsf{biblatex} use the \textsf{extradate} field, automatically
provided by \textsf{biber}, to distinguish citations of different
works by the same \textsf{author} that were published in the same
year, e.g., (Surname 1978\textbf{a}).  The Chicago author-date styles
recommend that some sorts of material --- online comments, newspaper
articles, and live performances, \emph{inter alia} --- needn't appear
in reference lists, but only in the text, often accompanied by a full
date reference (\texttt{cmsdate=full}) rather than by the rather less
informative year on its own.  In most circumstances a simple
\texttt{skipbib} in the \textsf{options} field will suffice, but,
especially with online materials, it is possible, even probable, that
users will have .bib databases containing different works by the same
\textsf{author} from the same \textsf{year}, only \emph{some} of which
need to appear in the reference list.  \textsf{Biber} will provide
\textsf{extradate} fields for all these entries, however, so it is
easy to get an extradate letter in a reference even when only one work
by that \textsf{author} appears in the list, or perhaps a series of
letters with some missing from the sequence.

\mylittlespace The \textsf{verbc} field allows you manually to
intervene to control these side effects.  (It's a standard
\textsf{biblatex} field, but isn't used in the standard styles.)
Putting anything in the field prevents that entry from interfering in
the \textsf{extradate} provision of entries that don't contain such a
field, and in more complicated scenarios you could group entries by
identical \textsf{verbc} field to prevent them from interfering both
with entries not having any \textsf{verbc} field \emph{and} with
entries having a different value for that field.  By default, the
\texttt{commenton} \textsf{relatedtype} for \textsf{online} and
\textsf{review} entries adds a \textsf{verbc} field to its entry, but
you can in all cases control this and provide your own in any
circumstances and in any entry type you wish.  Please see the
documentation of those two entry types in
section~\ref{sec:types:authdate}, and of the \textsf{commenton}
\textsf{relatedtype} in section~\ref{sec:authrelated}.

\mybigspace Standard \mymarginpar{\textbf{version}} \textsf{biblatex}
field, formerly only available in \textsf{artwork}, \textsf{image},
\textsf{misc}, \textsf{music}, and \textsf{patent} entries in
\textsf{biblatex-chicago-authordate}, but now also in \textsf{book}
and \textsf{performance} entries.  In most entry types it prints a
localized \enquote{\texttt{version}} string, but there may be
specialist needs in \textsf{artwork} and \textsf{image} entries, so
there you'll need to specify the type of data inside the field itself.
In the \textsf{book} type it is particularly needed for presenting
multimedia app content (15.57, 14.268).

\mybigspace Standard \mymarginpar{\textbf{volume}} \textsf{biblatex}
field.  It holds the volume of a \textsf{journaltitle} in
\textsf{article} entries, and also the volume of a multi-volume work
in many other sorts of entry.  The treatment and placement of
\textsf{volume} information in \textsf{book}-like entries is rather
complicated in the \emph{Manual} (14.116--22, 15.41).  In the
reference list, the \textsf{volume} appears either before the
\textsf{maintitle} or before the publication information, while in
citations you may need to provide it in the \textsf{postnote} field
--- see the \textsf{volumes} field, just below.  In a number of these
contexts, and in both books and periodicals, \textsf{volume}
information can appear \emph{immediately before} the page number(s).
In such a case, the \emph{Manual} (14.116) prescribes the same
treatment for both sorts of sources, that is, that \enquote{a colon
  separates the volume number from the page number with no intervening
  space.}  I have implemented this, but at the request of Clea~F.\
Rees \mymarginpar{\cmd{postvolpunct}} I have made this punctuation
customizable, using the command \cmd{postvolpunct}.  By default it
prints \cmd{addcolon}, so use \verb+\renewcommand{\postvolpunct}{...}+
in your preamble to redefine it.  Cf.\ \textsf{part}, and the command
documentation in section~\ref{sec:formatting:authdate};
conway:evolution shows how sometimes this field may hold series
information, as well.

\mybigspace Standard \mymarginpar{\textbf{volumes}} \textsf{biblatex}
field.  It holds the total number of volumes of a multi-volume work,
and in such references you should provide the volume and page numbers
in the \textsf{postnote} field of the relevant \cmd{cite} command,
e.g.:

\begin{verbatim}
\autocite[3:25]{bibfile:key}.
\end{verbatim}

Cf.\ 15.22; meredith:letters, tillich:system, weber:saugetiere,
wright:evolution.  The entry wright:theory presents one volume of such
a multi-volume work, so you would no longer need to give the volume in
any \textsf{postnote} field when citing it.  If both a \textsf{volume}
and a \textsf{volumes} field are present, as may occur particularly in
cross-referenced entries, then \textsf{biblatex-chicago} will
ordinarily suppress the \textsf{volumes} field, except in some cases
when a \textsf{maintitle} is present.  In this latter case, if the
\textsf{volume} appears before the \textsf{maintitle}, the option
\texttt{hidevolumes}, \mymarginpar{\texttt{hidevolumes}} set to
\texttt{true} by default, controls whether to print the
\textsf{volumes} field after that title or not.  Set it to
\texttt{false} either in the preamble or in the \textsf{options} field
of your entry to have it appear after the \textsf{maintitle}.  See the
option's documentation in section~\ref{sec:authpreset}, below.

\mybigspace A \mymarginpar{\textbf{xref}} modified \textsf{crossref}
field provided by \textsf{biblatex}, which prevents inheritance of any
data from the parent entry.  See \textbf{crossref}, above.

\mybigspace Standard \mymarginpar{\textbf{year}} \textsf{biblatex}
field, especially important for the author-date specification.  Please
see all the details under \textbf{date} above.  Unlike the
\textsf{date} field \textsf{year} allows non-numeric input, so you can
put \verb+\bibstring{nodate}+ here if required, or indeed any other
sort of non-numerical date information.  For many kinds of uncertain
and unspecified dates it is now much simpler to make use of
\textsf{biblatex's} enhanced date specifications in the \textsf{date}
field, instead.  Please see table~\ref{ad:date:extras} for a summary
of how \textsf{biblatex-chicago} implements these enhancements.  Cf.\
bedford:photo, clark:mesopot, leo:madonna, ross:thesis.

\subsubsection{Fields for Related Entries}
\label{sec:authrelated}

As \textsf{biblatex.pdf} puts it (\S~3.4), \enquote{Almost all
  bibliography styles require authors to specify certain types of
  relationship between entries such as \enquote{Reprint of},
  \enquote{Reprinted in,} etc. It is impossible to provide data fields
  to cover all of these relationships and so \textsf{biblatex}
  provides a general mechanism for this using the entry fields
  \textsf{related}, \textsf{relatedtype} and \textsf{relatedstring}.}
Before this mechanism was available \textsf{biblatex-chicago}
attempted to provide a similar but much more limited set of
inter-entry relationships using the \textsf{biblatex} fields
\textsf{origlanguage}, \textsf{origlocation}, \textsf{origpublisher},
\textsf{pubstate}, \textsf{reprinttitle}, and \textsf{userf}.  All of
these still work just as they always have or, I hope, somewhat better
than they always have after many recent bug fixes, but the more
general and more powerful \textsf{biblatex} \texttt{related} mechanism
is also available.  It can provide much of what the older system
provided and a great deal that it couldn't.  What follows is a
field-by-field discussion of the options now available.

\mybigspace In \mymarginpar{\textbf{origlanguage}} keeping with the
\emph{Manual}'s specifications, I have fairly thoroughly redefined
\textsf{biblatex}'s facilities for treating translations.  The
\textsf{origtitle} field isn't used, while the \textsf{language} and
\textsf{origdate} fields have been press-ganged for other duties.  The
\textsf{origlanguage} field, for its part, retains a dual role in
presenting translations in a list of references.  The details of the
\emph{Manual}'s suggested treatment when both a translation and an
original are cited may be found below under \textbf{userf}.  Here,
however, I simply note that the introductory string used to connect
the translation's citation with the original's is \enquote{Originally
  published as,} which I suggest may well be inaccurate in a great
many cases, as for instance when citing a work from classical
antiquity, which will most certainly not \enquote{originally} have
been published in the Loeb Classical Library.  Although not, strictly
speaking, authorized by the \emph{Manual}, I have provided another way
to introduce the original text, using the \textsf{origlanguage} field,
which must be provided \emph{in the entry for the translation, not the
  original text} (aristotle:metaphy:trans).  If you put one of the
standard \textsf{biblatex} bibstrings there (enumerated below), then
the entry will work properly across multiple languages.  Otherwise,
just put the name of the language there, localized as necessary, and
\textsf{biblatex-chicago} will eschew \enquote{Originally published
  as} in favor of, e.g., \enquote{Greek edition:} or \enquote{French
  edition:}.  This has no effect in citations, where only the work
cited --- original or translation --- will be printed, but it may help
to make the \emph{Manual}'s suggestions for the list of references
more palatable.  \textbf{NB:} You can use the \textsf{relatedtype}
\texttt{origpubas} with a customized \textsf{relatedstring} field to
achieve the same ends.

\mylittlespace That was the first usage, in keeping at least with the
spirit of the \emph{Manual}.  I have also, perhaps less in keeping
with that specification, retained some of \textsf{biblatex}'s
functionality for this field.  If an entry doesn't have a
\textsf{userf} field, and therefore won't be combining a text and its
translation in the list of references, you can also use
\textsf{origlanguage} as \textsf{biblatex} intended it, so that
instead of saying, e.g., \enquote{translated by X,} the entry will
read \enquote{translated from the German by X.}  The \emph{Manual}
doesn't mention this, but it may conceivably help avoid certain
ambiguities in some citations.  As in \textsf{biblatex}, if you wish
to use this functionality, you have to provide \emph{not} the name of
the language, but rather a bibstring, which may, at the time of
writing, be one of \texttt{american}, \texttt{brazilian},
\texttt{danish}, \texttt{dutch}, \texttt{english}, \texttt{french},
\texttt{german}, \texttt{greek}, \texttt{italian}, \texttt{latin},
\texttt{norwegian}, \texttt{portuguese}, \texttt{spanish}, or
\texttt{swedish}, to which I've added \texttt{russian}.

\mybigspace This \mymarginpar{\textbf{origlocation}} field mainly
serves to help document reprint editions and their corresponding
originals (14.114, 15.40).  In \textsf{biblatex-chicago} you can
provide both an \textsf{origlocation} and an \textsf{origpublisher} to
go along with the \textsf{origdate}, should you so wish, and all of
this information will be printed in the reference list.  You can also
use this field in a \textsf{letter} or \textsf{misc} (with
\textsf{entrysubtype}) entry to give the place where a published or
unpublished letter was written (14.111, 14.229).  (Jonathan Robinson has
suggested that the \textsf{origlocation} may in some circumstances
actually be helpful for disambiguation, his example being early
printed editions of the same material printed in the same year but in
different cities.  The new functionality should make this simple to
achieve.  Cf.\ \textsf{origdate} [section~\ref{sec:fields:authdate}],
\textsf{origpublisher} and \textsf{pubstate}; schweitzer:bach.)
\textbf{NB:} It is impossible to present this same information, as
here, \emph{inside} a single entry using a \texttt{related} field,
though the \textsf{relatedtype} \texttt{origpubin} presents much the
same information \emph{after} the entry, using data extracted from a
separate entry.

\mybigspace As \mymarginpar{\textbf{origpublisher}} with the
\textsf{origlocation} field just above, this field mainly serves to
help document reprint editions and their corresponding originals
(14.114, 15.40).  You can provide an \textsf{origpublisher} and/or an
\textsf{origlocation} in addition to the \textsf{origdate}, and all
will be presented in the reference list.  (Cf.\ \textsf{origdate}
[section~\ref{sec:fields:authdate}], \textsf{origlocation}, and
\textsf{pubstate}; schweitzer:bach.)  \textbf{NB:} It is impossible to
present this same information, as here, \emph{inside} a single entry
using a \texttt{related} field, though the \textsf{relatedtype}
\texttt{origpubin} presents much the same information \emph{after} the
entry, using data extracted from a separate entry.

\mybigspace In \mymarginpar{\textbf{pubstate}} response to new
specifications in the 17th edition of the \emph{Manual} (esp.\
14.137), I have tried to generalize the functioning of the
\textsf{pubstate} field in all entry types.  Because the author-date
style has fairly complicated rules about presenting reprinted editions
(15.40), the \texttt{reprint} string still has a special status.
Depending on which date(s) you have chosen to appear at the head of
the entry, \textsf{biblatex-chicago-authordate} will either print the
(localized) string \texttt{reprint} in the proper place or otherwise
provide a notice at the end of the entry detailing the original
publication date.  See under \textbf{date} above for the available
permutations.  (Cf.\ aristotle:metaphy:gr, maitland:canon,
maitland:equity, schweitzer:bach.)

\mylittlespace Other strings are divided into two types: those which
\textsf{biblatex-chicago} will print as the \textsf{year}, which
currently means \emph{only} those for which \textsf{biblatex} contains
bibstrings indicating works soon to be published, i.e.,
\texttt{forthcoming}, \texttt{inpreparation}, \texttt{inpress}, and
\texttt{submitted}; and those, i.e., everything else, which will be
printed before, and in close association with, other information about
the publisher of a work.  (This \mymarginpar{\textbf{NB}} is a change
from previous behavior, where non-\texttt{reprint} strings were
printed \emph{after} the publication information, as in the standard
styles.  You can still use the \textsf{addendum} field to present
information here, of course.)  The four strings that replace the
\textsf{year} will always be localized, as will \texttt{reprint} and
\texttt{selfpublished} (and anything else that \textsf{biblatex} finds
to be a \cmd{bibstring}) from the second category.  All other strings
will be printed as-is, capitalized if needed, just before the
publisher (author:forthcoming, contrib:contrib, schweitzer:bach).

\mylittlespace There is one further subtlety of which you ought to be
aware.  In \textsf{music} and \textsf{video} entries, the
\texttt{reprint} string in \textsf{pubstate} will only make a
difference to your entries when the date which it modifies --- the
\textsf{origdate}, typically --- \emph{doesn't} appear in citations
and at the head of reference-list entries.  In this case the date is
treated as an original release date, and it will be printed, preceded
by the appropriate string, near the end of the entry.  Other strings
don't show this special behavior in these entries.  \textbf{NB:} For
those uses of the \textsf{pubstate} field that print a notice at the
end of the entry, the \textsf{relatedtype} \texttt{origpubin} provides
much the same information, using data extracted from a different
entry.  If the information appears inside the entry then there is no
equivalent \textsf{related} functionality.

\enlargethispage{1.5\baselineskip}

\mybigspace This \mymarginpar{\textbf{related}} field is required to
use \textsf{biblatex's} \textsf{related} functionality, and it should
contain the entry key or keys from which \textsf{biblatex} should
extract data for presentation not on its own, but rather in the
reference list entry which contains the \textsf{related} field itself.
Indeed, unless you change the defaults using the
\textsf{relatedoptions} field this data will only appear in such
entries, never on its own and never in citations.  Without a
\textsf{relatedtype} field, this will print the default type,
equivalent to a full reference list entry \emph{immediately after} the
entry containing the \textsf{related} field, with no intervening
string.  You can specify a string using the \textsf{relatedstring}
field, so in effect this presents a powerful mechanism for presenting
full references to related material of any sort whatsoever.

\mylittlespace By \mymarginpar{\texttt{related=true}} default, the
package option \texttt{related} is set to print \textsf{related}
entries in the list of references.  If you would like to turn this off
you can set this option, either in your preamble or in the
\textsf{options} or \textsf{relatedoptions} field of the relevant
entry, to \texttt{false}.  For the two \textsf{relatedtypes} that
construct a single entry using data extracted from related entries ---
\texttt{commenton} and \texttt{reviewof} --- you'll need to make sure
this is \texttt{true} to get properly-formatted citations in the
reference list.  See below for the details.

\mybigspace This \mymarginpar{\textbf{relatedoptions}} field will, I
should expect, only be needed very rarely.  If you want to set
entry-level options for a \textsf{related} entry this is where you can
do it, though please remember one important detail.  By default,
\textsf{Biber} sets this option to \texttt{dataonly}, which among
other things prevents the \textsf{related} entry from appearing
separately in the list of references, assuming you don't specifically
cite it elsewhere.  If you use the field yourself, then you'll need to
include \texttt{dataonly} as one of the options therein to maintain
this effect.  Of course, it may be you don't want all the effects of
\texttt{dataonly}, so you can tailor it however you wish.  See
\textsf{biblatex.pdf} \S~3.4.

\mybigspace The \mymarginpar{\textbf{relatedstring}} procedure for
choosing a string to connect the main entry with its related entry/ies
is straightforward, the default being a \texttt{bibstring}, if any,
with the same name as the \textsf{relatedtype}, or alternately a
string or strings defined within the driver for that
\textsf{relatedtype}, as happens with the types \texttt{origpubin} and
\texttt{bytranslator}.  Failing these, you can supply your own in the
\textsf{relatedstring} field, either in the form of the name of a
pre-defined \texttt{bibstring} or as any text you choose, and anything
in this field always takes precedence over the automatic choices.  If
your non-\texttt{bibstring} starts with a lowercase letter then
\textsf{biblatex-chicago} will capitalize it automatically for you
depending on context (coolidge:speech, weed:flatiron).  I have not
altered the standard \textsf{relatedtype} strings, and have in fact
modified the \textsf{reprinttitle} mechanism to use the
\texttt{reprintfrom} string, which works better syntactically in this
context, and modified the \textsf{pubstate} mechanism to use the
\texttt{origpubin} string, which brings it into line with the notes
\&\ bibliography style.

\mybigspace The \mymarginpar{\textbf{relatedtype}} standard
\textsf{biblatex} styles define six \textsf{relatedtypes}, and I have
either simply adopted them wholesale or adapted them to the needs of
the Chicago style, retaining the basic syntax as much as possible.  I
have also added three to these six (see below):

\begin{description}
\item[\qquad bytranslator:] This prints a full reference to a
  translation, starting with the (localized) string
  \enquote{Translated by \textsf{translator} as \textsf{Title},
    \ldots} The reference is fuller in \textsf{biblatex-chicago} than
  in the standard styles, and for the first time allows users to
  choose the \emph{Manual's} alternate method for presenting original
  + translation (14.99; furet\hc passing:fr).  The old \textsf{userf}
  mechanism provides the other, as does the \texttt{origpubas}
  \textsf{relatedtype} (see below).
\item[\qquad default:] This is the macro used when no
  \textsf{relatedtype} is defined.  It prints, as in the standard
  styles, and with no intervening string, full references to the
  \textsf{related} entries.
\item[\qquad multivolume:] This briefly lists the individual volumes
  in a multi-volume work, and works much as in the standard styles.
  The \emph{Manual}, as far as I can see, has little to say on the
  matter.
\item[\qquad origpubas:] This type can, if you want, replace the old
  \textsf{userf} mechanism, described below, for presenting an
  original with its translation.  It's quite similar to the
  \texttt{default} type, but with a \texttt{bibstring} automatically
  connecting the entry with its \textsf{related} entries.  You can
  identify other sorts of relationships if you change the introductory
  string using \textsf{relatedstring}.
\item[\qquad origpubin:] I have barely altered this from the
  \textsf{biblatex} default, and it will present reprint information
  \emph{after} the main entry rather than within it.  The
  \emph{Manual} seems to prefer the latter for the notes \&
  bibliography style and, in some circumstances, the former for
  author-date.
\item[\qquad reprintfrom:] This type provides a replacement for the
  old \textsf{reprinttitle} mechanism described below.  As in the
  standard styles, it presents a fuller reference to the reprinted
  material than does \texttt{origpubin}, and is designed particularly
  for presenting pieces formerly printed in other collections or
  perhaps essays collected from various periodicals.  (In
  \textsf{biblatex-chicago} it contains some kludges to cope with
  possible \textsf{babel} language environments, so if you find it
  behaving oddly please let me know, including whether you are using
  \textsf{babel} [which I've tested] or \textsf{polyglossia} [which
  I've tested somewhat less].)
\end{description}
Now, the Chicago-specific types:
\begin{description}
\item[\qquad commenton:] I designed the new \textsf{relatedtype}
  \texttt{commenton} to facilitate citation of online comments, and it
  is available in two entry types, \textsf{online} and \textsf{review}
  (with its clone \textsf{suppperiodical}).  In both types the
  \emph{Manual} (15.51--52) recommends that such material appear
  \emph{only} in the text and not in the reference list, but I have
  attempted to simplify the presentation of such material wherever you
  want it to appear.  Following the specifications, then, the default
  when you use \texttt{commenton} is for
  \textsf{biblatex-chicago-authordate} to modify how your .bib entry
  appears in the .bbl file by setting both \texttt{skipbib} and
  \texttt{cmsdate=full} in the \textsf{options} field, so that nothing
  appears in the reference list and citations present the full date
  and possibly also a time stamp (see below).  Further, the style sets
  the \textsf{verbc} field so that these entries don't interfere with
  the provision of extra date letters --- the full date and time
  should be enough to individuate separate comments.  Finally, the
  style creates a new \textsf{customc} entry in your .bbl file which
  you can cite after your initial \texttt{commenton} entry using
  \cmd{autocites} and which will, as a comment to your initial entry,
  say whether it's a comment or a reply or what have you, and then
  giving the short citation of that upon which it is a comment.  (Just
  to be clear: your .bib file will itself never be altered, only the
  .bbl file, which is produced by \textsf{biber} and which provides
  the data from which \textsf{biblatex} actually typesets citations.)
  Because of how this works, you can currently only use \emph{one}
  entry key in the \textsf{related} field for this
  \textsf{relatedtype}, though I may lift this restriction in the
  future.

  As an example, take the Facebook post diaz:surprise, which does
  appear in the reference list.  The entry licis:diazcomment presents
  a comment on this post using the \textsf{relatedtype}
  \texttt{commenton}, so \textsf{biblatex-chicago-authordate} creates
  a new entry, diaz:surprise-customc.  When you cite the comment in
  your document a command like
  \verb+\autocites{licis:diazcomment}{diaz:surprise-customc}+ will
  produce a citation like (Licis, February 24, 2016; comment on D\'iaz
  2016).  You can alter the string connecting the two citations (by
  default \verb+\bibstring{commenton}+) by using the
  \textsf{relatedstring} field in the first of them (cf.\
  powell:comment).  (Note how minimal the .bib entry of a comment
  using this system can be --- \textsf{author}, \textsf{related},
  \textsf{relatedtype}, and \textsf{date} are pretty much the only
  fields required.)

  Those who want \textsf{online} comments to appear in the reference
  list can still use the \texttt{commenton} \textsf{relatedtype}, and
  the same citation of the commented piece will appear there,
  connected by the same string that the \textsf{customc} entry
  provides.  Here, though, you can also provide a separate
  \textsf{title} for the comment, and/or a separate \textsf{url} for
  it, should they exist, which will be printed before/after the
  citation of the commented piece, respectively.  (In \textsf{review}
  entries, which use the same \textsf{relatedtype}, only the generic
  title is available, as is always the case with such entries.)  If
  you manually set either (or both) of the \texttt{cmsdate} or the
  \texttt{skipbib} options in your entry then
  \textsf{biblatex-chicago} will assume you want to hand-craft that
  entry without its intervention, so it won't alter the
  \textsf{options} field or indeed provide any \textsf{verbc} field,
  though it will still provide the virtual \textsf{customc} entry in
  your .bbl file, as that may still prove convenient.  Note also that
  any \textsf{verbc} field you provide will never be altered by the
  package.
\item[\qquad reviewof:] Philip Kime's \textsf{biblatex-apa} package
  includes this type, and user Bertold Schweitzer suggested it might
  be a useful addition to \textsf{biblatex-chicago}, so I've added it
  to the standard six detailed above.  It differs from all of them in
  that it prints the \textsf{relatedstring} (by default
  \cmd{bibstring\{reviewof\}}) and the data from the \textsf{related}
  entry in the middle of the parent entry, rather than at the end.  It
  also differs from them in being available only in \textsf{article}
  and \textsf{review} entries (along with the latter's clone,
  \textsf{suppperiodical}).

%  %\enlargethispage{\baselineskip}

  In \textsf{article} entries it replaces the \textsf{titleaddon} with
  the \textsf{relatedstring} followed by the \textsf{title} of the
  child entry, and in \textsf{review} entries it replaces the
  \textsf{title} with the same two components.  In both types these
  components will optionally be followed by the \textsf{author},
  \textsf{editor}, \textsf{translator}, etc.,\ of the reviewed item,
  and then any child \textsf{titleaddon} may optionally appear at the
  end, allowing maximum flexibility when presenting, for example,
  reviews of live performances.

  This mechanism automates both the provision of the localized
  \cmd{bibstring} and also the formatting of the \textsf{title} of the
  reviewed work, and it also obviates the need to use any of the
  \cmd{partedit} macros in this context.  If you've changed the
  default setting of the \texttt{related} option in the preamble, then
  you'll need to ensure that it is set to \texttt{true} in the
  individual entries where you use this \textsf{relatedtype} to ensure
  that the entry's full data appears in the list of references.  Also,
  if the mechanism doesn't work for you in a particular context,
  remember that the standard way of presenting reviewed works is still
  available.
\item[\qquad short:] This \textsf{relatedtype} is like the
  \texttt{default} type, only it prints author-date citations rather
  than full reference-list entries. There is no default
  \textsf{relatedstring} for this type, so if you leave that field
  blank then the short references will simply appear at the end of the
  parent entry in the reference list.  Any entries you reference in
  this way will by default appear separately in the reference list,
  even if you haven't cited them anywhere else in your document.
  Otherwise the short reference won't actually be decipherable.
\end{description}

\mybigspace \textbf{NB:} \mymarginpar{\textbf{reprinttitle}}
\textbf{If you have been using this feature, you may want to have a
  look at the} \textsf{relatedtype} \texttt{reprintfrom},
\textbf{documented above, for a better solution to this problem, one
  that also allows you to change the introductory string using the}
\textsf{relatedstring} \textbf{field.  The} \textsf{reprinttitle}
\textbf{field will continue to work as before, however.}  At the
request of Will Small, I have included a means of providing the
original publication details of an essay or a chapter that you are
citing from a subsequent reprint, e.g., a \emph{Collected Essays}
volume.  In such a case, at least according to the \emph{Manual}
(14.181), these details would only appear in the reference list, and
then only if these details are \enquote{of particular interest.}  The
data would follow an introductory phrase like \enquote{originally
  published as,} making the problem strictly parallel to that of
including details of a work in the original language alongside the
details of its translation.  I have addressed the latter problem with
the \textsf{userf} field, which provides a sort of cross-referencing
method for this purpose, and \textsf{reprinttitle} works in
\emph{exactly} the same way.  In the .bib entry for the reprint you
include a cross-reference to the cite key of the original location
using the \textsf{reprinttitle} field (which it may help mnemonically
to think of as a \enquote{reprinted title} field).  The main
difference between the two forms is that \textsf{userf} prints all but
the \textsf{author} of the original work, whereas
\textsf{reprinttitle} suppresses both the \textsf{author} and the
\textsf{title} of the original, giving only the more general details,
beginning with, e.g., the \textsf{journaltitle} or \textsf{booktitle}
and continuing from there.  The string prefacing this information will
be \enquote{Originally published in.}  Please see the documentation on
\textsf{userf} below for all the details on how to create .bib entries
for presenting your data.

\mybigspace This \mymarginpar{\textbf{userf}} is one of the
supplemental fields which \textsf{biblatex} provides, and is used by
\textsf{biblatex-chicago} for a very specific purpose.  When you cite
both a translation and its original, the \emph{Manual} (14.99)
recommends that, in a reference list at least, you combine references
to both texts in one entry.  Lacking specific instructions about the
author-date style, I have nonetheless chosen to implement this
possibility also for a list of references, though in-text citations
will still only refer to individual works.  In order to follow this
specification, I have provided a third cross-referencing system (the
others being \textsf{crossref} and \textsf{xref}), and have chosen the
name \textsf{userf} because it might act as a mnemonic for its
function.

\mylittlespace In order to use this system, you should start by
entering both the original and its translation into your .bib file,
just as you normally would.  The mechanism works for any entry type,
and the two entries need not be of the same type.  In the entry for
the \emph{translation}, you put the cite key of the original into the
\textsf{userf} field.  In the \emph{original's} entry, you need to
include some means of preventing it appearing separately in the list
of references, either a toggle in the \textsf{keywords} field or
perhaps \texttt{skipbib} in the \textsf{options} field.  In this
standard case, the data for the translation will be printed first,
followed by the string \texttt{orig. pub. as}, followed by the
original, author omitted.  As explained above (\textbf{origlanguage}),
I have also included a way to modify the string printed before the
original.  In the entry for the \emph{translation}, you put the
original's language in \textsf{origlanguage}, and instead of
\texttt{originally published as}, you'll get \texttt{French edition:}
or \texttt{Latin edition:}, etc.\ (aristotle:metaphy:gr,
aristotle:metaphy:trans).  \textbf{NB:} You can use the
\textsf{relatedtype} \texttt{origpubas} to replicate the
\textsf{userf} functionality, and you can also customize the
\textsf{relatedstring} field to achieve the same result as with
\textsf{origlanguage}.

\subsection{Commands}
\label{sec:commands:authdate}

In this section I shall attempt to document all those commands you may
need when using \textsf{biblatex-chicago-authordate} that I have either
altered with respect to the standard provided by \textsf{biblatex} or
that I have provided myself.  Some of these, unfortunately, will make
your .bib file incompatible with other \textsf{biblatex} styles, but
I've been unable to avoid this.  Any ideas for more elegant, and more
compatible, solutions will be warmly welcomed.

\subsubsection{Formatting Commands}
\label{sec:formatting:authdate}

These commands allow you to fine-tune the presentation of your
references in both citations and list of references.  You can find
many examples of their usage in \textsf{dates-test.bib}, and I shall
try to point you toward a few such entries in what follows.
\textbf{NB:} \textsf{biblatex's} \cmd{mkbibquote} command is mandatory
in some situations.  See its entry below.

\mybigspace Version \mymarginpar{\textbf{\textbackslash autocap}} 0.8
of \textsf{biblatex} introduced the \cmd{autocap} command, which
capitalizes a word inside a citation or list of references entry if
that word follows sentence-ending punctuation, and leaves it lowercase
otherwise.  The whole question of capitalization is considerably more
complicated in the notes \&\ bibliography style, where the former uses
commas and the latter (often) periods to separate blocks of
information, whereas the more streamlined author-date specification
has few such issues.  In \textsf{dates-test.bib} there are only two
places where the \cmd{autocap} macro is necessary, and they both
involve the string \texttt{forthcoming} in the \textsf{year} field
(author:forthcoming, contrib:contrib), though you can now avoid even
this necessity by placing \texttt{forthcoming} in the
\textsf{pubstate} field.

\mylittlespace I have nonetheless retained the system developed,
following \textsf{biblatex's} example, for the notes \&\ bibliography
style, which automatically tracks the capitalization of certain fields
in your .bib file.  I chose these fields after a non-scientific survey
of entries in my own databases, so of course if you have ideas for the
extension of this facility I would be most interested to hear them.
In order to take advantage of this functionality, all you need do is
begin the data in the appropriate field with a lowercase letter,
e.g.,\ \texttt{note = \{with the assistance of X\}}.  If the data
begins with a capital letter --- and this is not infrequent --- that
capital will always be retained.  (cf., e.g., creel:house,
morgenson:market.)  If, on the other hand, you for some reason need
such a field always to start with a lowercase letter, then you can try
putting an empty set of curly braces\ \{\}\ at the start, which turns
off the mechanism without printing anything itself.  Here, then, for
reference purposes, is the complete list of fields where this
functionality is active:

\begin{enumerate}
\setlength{\parskip}{-4pt}
\item The \textbf{addendum} field in all entry types.
\item The \textbf{booktitleaddon} field in all entry types.
\item The \textbf{edition} field in all entry types.  (Numerals work
  as you expect them to here.)
\item The \textbf{maintitleaddon} field in all entry types.
\item The \textbf{note} field in all entry types.
\item The \textbf{part} field in entry types that use it.
\item The \textbf{prenote} field prefixed to citation commands.
\item The \textbf{relatedstring} field in all entry types.
\item The \textbf{shorttitle} field in the \textsf{review}
  (\textsf{suppperiodical}) entry type and in the \textsf{misc} type,
  in the latter case, however, only when there is an
  \textsf{entrysubtype} defined, indicating that the work cited is
  from an archive.
\item The \textbf{title} field in the \textsf{review}
  (\textsf{suppperiodical}) entry type and in the \textsf{misc} type,
  in the latter case, however, only when there is an
  \textsf{entrysubtype} defined, indicating that the work cited is
  from an archive.
\item The \textbf{titleaddon} field in all entry types.
\item The \textbf{type} field in \textsf{artwork}, \textsf{audio},
  \textsf{image}, \textsf{music}, \textsf{suppbook},
  \textsf{suppcollection}, and \textsf{video} entry types.
\end{enumerate}

If you accidentally use the \cmd{autocap} macro in one of the above
fields, it really shouldn't matter at all, and you'll still get what
you want, but taking advantage of the automatic provisions should at
least save some typing.

\mybigspace This \mymarginpar{\textbf{\textbackslash bibstring}} is a
very powerful mechanism to allow \textsf{biblatex} automatically to
provide a localized version of a string, and to determine whether that
string needs capitalization, depending on where it falls in an entry.
\textsf{Biblatex} also provides functionality which allows you
sometimes simply to input, for example, \texttt{newseries} instead of
\verb+\bibstring{newseries}+, the package auto-detecting when a
bibstring is involved and doing the right thing, though in all such
cases either form will work.  This functionality is available in the
\textsf{series} field of \textsf{article}, \textsf{periodical}, and
\textsf{review} entries; in the \textsf{type} field of
\textsf{manual}, \textsf{patent}, \textsf{report}, and \textsf{thesis}
entries; in the \textsf{location} field of \textsf{patent} entries; in
the \textsf{language} field in all entry types; and in the
\textsf{nameaddon} field in \textsf{customc} entries.  These are the
places, as far as I can make out, where \textsf{biblatex's} standard
styles support this feature, though I have added the last,
style-specific, one.  If the \textsf{biblatex} authors generalize it
still further in a future release, I shall do the same, if possible.

\mybigspace I \mymarginpar{\textbf{\textbackslash letterdatelong}}
have provided this macro mainly for use in the optional postnote field
of the various citation commands.  When citing a letter (published or
unpublished, \textsf{letter} or \textsf{misc}), it may be useful to
include the date in the citation in order to disambiguate references.
This macro simply prints the date of a letter, or indeed of any other
sort of correspondence, in day-month-year order, as recommended by the
\emph{Manual} (14.224).  (If your main document language isn't
American, it's better just to use the standard \textsf{biblatex}
command \cmd{printorigdate}.)

\mybigspace This \mymarginpar{\textbf{\textbackslash mkbibquote}} is
the standard \textsf{biblatex} command, which requires attention here
because it is a crucial part of the mechanism of that package's
\enquote{American} punctuation system.  Quotation marks around the
\textsf{title} field in various entry types are automatically provided
by \textsf{biblatex-chicago}, but titles-within-titles frequently also
require them, so it is best to get accustomed to using this command to
make sure any periods or commas appearing in the neighborhood of the
closing quotes will appear inside them automatically.  A few examples
from \textsf{dates-test.bib} should help to clarify this.

\mylittlespace In an \textsf{article} entry, the \textsf{title}
contains a quoted phrase:

\begin{quotation}
  \noindent\texttt{title = \{Diethylstilbestrol and Media Coverage of the \\
    \indent\cmd{mkbibquote}\{Morning After\} Pill\}}
\end{quotation}

Here, because the quoted text doesn't come at the end of title, and no
punctuation will ever need to be drawn within the closing quotation
mark, you could instead use \texttt{\cmd{enquote}\{Morning After\}} or
even \texttt{`Morning After'}. (Note the single quotation marks here
--- the other two methods have the virtue of taking care of nesting
for you.)  All of these will produce the formatted:
\enquote{Diethylstilbestrol and Media Coverage of the \enquote{Morning
    After} Pill.}

\mylittlespace Here, by contrast, is a \textsf{book title}:

\begin{quotation}
  \noindent \texttt{title = \{Annotations to
    \cmd{mkbibquote}\{Finnegans Wake\}\}}
\end{quotation}

Because the quoted title within the title comes at the end of the
field, and because this reference unit will be separated from
what follows by a period in the list of references, then the
\cmd{mkbibquote} command is necessary to bring that period within the
final quotation marks, like so: \emph{Annotations to
  \enquote{Finnegans Wake.}}

\mylittlespace Note in both cases that you only need to be careful
with the capitalization inside the curly brackets if you are using
\textsf{authordate-trad}, as recent editions of the \emph{Manual} have
unified the title formatting for the two remaining styles, which means
that, for them, all lower- and uppercase letters remain as they are
typed in your .bib file.

\enlargethispage{\baselineskip}

\mylittlespace Let me also add that this command interacts well with
Lehman's \textsf{csquotes} package, which I highly recommend, though
the latter isn't strictly necessary in texts using an American style,
to which \textsf{biblatex} defaults when \textsf{csquotes} isn't
loaded.

\mybigspace The \mymarginpar{\textbf{\textbackslash postvolpunct}}
\emph{Manual} (14.116) unequivocally prescribes that when a
\textsf{volume} number appears immediately before a page number,
\enquote{the abbreviation \emph{vol.}\ is omitted and a colon
  separates the volume number from the page number with no intervening
  space.}  The treatment is basically the same whether the citation is
of a book or of a periodical, and it appears to be a surprising and
unwelcome feature for many users, conflicting as it may do with
established typographic traditions in a number of contexts.  Clea~F.\
Rees has requested a way to customize this, so I have provided the
\cmd{postvolpunct} command, which prints the punctuation between a
\textsf{volume} number and a page number.  It is set to \cmd{addcolon}
by default, except when the current language of the entry is French,
in which case it defaults to \cmd{addcolon\textbackslash addspace}.
You can use \cmd{renewcommand\{\textbackslash
  postvolpunct\}\{\ldots\}} in your preamble to redefine it, but
please note that the command only applies in this limited context, not
more generally to the punctuation that appears between, e.g., a
\textsf{volume} and a \textsf{part} field.

\mybigspace This \mymarginpar{\textbf{\textbackslash partcomp}} and
the following 6 macros were all designed to help
\textsf{biblatex-chicago} cope with the fact that many bibstrings in
the notes \&\ bibliography style differ between notes and
bibliography, the former sometimes using abbreviated forms when the
latter prints them in full.  These problems do not arise in the
author-date styles, but using these macros will make your .bib database
more portable across languages and across both Chicago styles, and may
be slightly easier to remember than the strings themselves.  On the
other hand, of course, they will make your .bib file less portable
across multiple \textsf{biblatex} styles.

\mylittlespace These macros allow you to provide an \texttt{editor}, a
\texttt{translator}, and/or a \texttt{compiler} in situations where
the available fields (\textsf{editor}, \textsf{namea},
\textsf{translator}, \textsf{nameb}, and \textsf{namec}) aren't
adequate.  Their names all begin with \cmd{part}, as originally I
intended them for use when a particular name applied only to a
specific \textsf{title}, rather than to a \textsf{maintitle} or
\textsf{booktitle} (cf.\ \textbf{namea} and \textbf{nameb}, above).

\mylittlespace In the present instance, you can use \cmd{partcomp} to
identify a compiler when \textsf{namec} (or \textsf{editortype}) won't
do, e.g., in a \textsf{note} field or the like.  In such a case,
\textsf{biblatex-chicago} will print the appropriate string in your
references.

\mybigspace Use \mymarginpar{\textbf{\textbackslash partedit}} this
macro when identifying an editor whose name doesn't conveniently fit
into the usual fields (\textsf{editor} or \textsf{namea}).  (N.B.: If
you are writing in French then you no longer need to add either
\texttt{de} or \texttt{d'} after this command in your .bib files.  The
new version of the command should take care of this automatically for
you.)  See howell:marriage.

\mybigspace As \mymarginpar{\textbf{\textbackslash
    partedit-\\andcomp}} before, but for use when an editor is also a
compiler.

\vspace{1.3\baselineskip} As \mymarginpar{\textbf{\textbackslash
    partedit-\\andtrans}} before, but for when when an editor is also a
translator (ratliff:review).

\vspace{1.3\baselineskip} As \mymarginpar{\textbf{\textbackslash
    partedit-\\transandcomp}} before, but for when an editor is also a
translator and a compiler.

\vspace{1.4\baselineskip} As \mymarginpar{\textbf{\textbackslash
    parttrans-\\andcomp}} before, but for when a translator is also a
compiler.

\vspace{1.3\baselineskip} As \mymarginpar{\textbf{\textbackslash
    parttrans}} before, but for use when identifying a translator
whose name doesn't conveniently fit into the usual fields
(\textsf{translator} and \textsf{nameb}).

\mylittlespace This \mymarginpar{\textbf{\textbackslash reprint}} is
equivalent to \verb+\bibstring{reprint}+.  It is useful in the notes
\&\ bibliography style, and I include it in the author-date styles for
compatibility.

\mylittlespace Unlike \mymarginpar{\textbf{\textbackslash
    suppress-\\bibfield[]\{\}}} the other commands presented here,
this should be used in your document preamble rather than in your
bibliographical apparatus.  Also unlike them, it has two arguments,
the first of which is optional, the second required.  Jan David Hauck
suggested that, in addition to the field-exclusion package options
provided by \textsf{biblatex-chicago} (see
section~\ref{sec:authpreset}), I might also provide a general-purpose
macro to clear fields from selected entry types when the package
options aren't quite right for a user's particular needs.  The
\cmd{suppressbibfield} command does this, so that
\verb+\suppressbibfield{note}+ clears the \textsf{note} field from
\emph{all} entries, while \verb+\suppressbibfield[report]{note}+
clears it only from \textsf{report} entries.  Both arguments take
comma-separated lists, so to suppress \textsf{titleaddon} and
\textsf{volumes} fields from \textsf{report} and \textsf{manual}
entries, your preamble could contain
\verb+\suppressbib-field[report,manual]{titleaddon,volumes}+.

\mylittlespace A few usage notes are in order.  First, you can use as
many calls to the command in your preamble as you wish.  Second, the
command is a very basic user interface to \textsf{biblatex's} source
mapping functionality (\textsf{biblatex.pdf} \S~4.5.3), so what it
does is modify what \textsf{biber} takes from your .bib file in order
to produce the .bbl file that \textsf{biblatex} actually reads.  As
far as \textsf{biblatex} is concerned, the fields simply aren't there
in the data source, so they can't appear anywhere in the
bibliographical apparatus, whether in citations, reference lists, or
shorthand lists.  Third, because source mapping is involved, you'll
need a complete cycle of \LaTeX-\textsf{biber}-\LaTeX\ runs to make
the commands take effect.  Fourth, source mapping occurs at a very
early stage in \textsf{biber's} operation, so if your field names or
entry types are standard aliases, the command will only work on the
names as they appear in your .bib file, not as they are aliased in the
.bbl file.  If you have a \textsf{techreport} entry, for example, it
won't be affected by a command that alters \textsf{report} entries,
and a \textsf{date} field won't be affected by a command that
suppresses the \textsf{year}.  Fifth, the code for the command resides
in \textsf{biblatex-chicago.sty}, so if you load the styles without
loading that package it won't be available to you.  Sixth and finally,
the \cmd{suppressbibfield} command is new and relatively untested, so
please report any untoward behavior to me.

\subsubsection{Citation Commands}
\label{sec:cite:authordate}

The \textsf{biblatex} package is particularly rich in citation
commands, most of which, in \textsf{biblatex-chicago-authordate} and
\textsf{authordate-trad}, function as they do in the standard
author-date styles.  If you are getting unexpected behavior when using
them please have a look in your .log file.  A command like
\cmd{supercite}, listed in \S~3.6.2 of the \textsf{biblatex} manual
but not defined by \textsf{biblatex-chicago-authordate} or by core
\textsf{biblatex}, defaults to \cmd{cite}, and leaves a warning in the
.log.  The following commands may require some minimal explanation,
but if there are standard commands that don't work for you, or new
commands that would be useful, please let me know, and it should be
possible to fix or add them.

\mybigspace These \mymarginpar{\textbf{\textbackslash atcite
    \\\textbackslash atpcite}} two new citation commands allow you
quickly and easily to provide an author-title citation of any entry,
instead of an author-date citation.  The \texttt{classical}
\textsf{entrysubtype} field does this, but it also changes punctuation
in the citation, so I've provided other means to achieve the same end.
The first of these new citation commands presents the plain citation,
the second includes it in parentheses for inclusion in running text.
The new \texttt{authortitle} type and entry option
(section~\ref{sec:authuseropts}) has the same effect when using the
standard citation commands, but it's possible that using these new
commands instead may give added flexibility.

\mybigspace I \mymarginpar{\textbf{\textbackslash autocite}} haven't
adapted this in the slightest, but I thought it worth pointing out
that \textsf{biblatex-chicago-authordate} sets this command to use
\cmd{parencite} as the default option.  It is, in my experience, much
the most common citation command you will use, and also works fine in
its multicite form, \textbf{\textbackslash autocites}.

\mybigspace Although \mymarginpar{\textbf{\textbackslash citeincite}}
the officially-sanctioned, and safest, way to present cross-references
to other works is by using the \textsf{related} mechanism, it would
appear, judging from various bug reports I've received, that there are
plenty of users who need to present such references in ways that
aren't supported by the \textsf{related} code as it stands. This new
citation command is designed to allow you to present short references
to other works inside other references, and to avoid some, I hope
most, of the bugs associated with using cite commands in the fields of
other entries. If you want to present long citations, I'm assuming
that you wouldn't want to do that in fields right in the middle of the
parent entry, and that, therefore, the usual methods detailed in
section~\ref{sec:authrelated} (above) or the \cmd{fullciteincite}
command (below) will serve your needs.  Similarly, if you want to
provide short citations at the end of your parent entry, the new
\textsf{relatedtype} \texttt{short} should work.

\mylittlespace If, however, you must have a short reference in the
midst of another field, then \cmd{citeincite} can help. It is
certified to work properly only in the \textsf{addendum, annotation,
  annote, note, and titleaddon} fields, and in the author-date styles
also in the \textsf{title} fields of \textsf{customc} entries, which
allows you (via a multicite command) to include such citations as
comments to other in-text citations. The multicite command
\cmd{citeincites} is also available in the same contexts. In previous
releases of \textsf{biblatex} there were differences in how the
case-changing backends behaved with respect to citation commands
inside entry fields, but in my testing the two backends appear
currently to be equivalent on this score. Still, should you run into
problems using \cmd{citeincite}, it might still be worth setting
\texttt{casechanger=latex2e} in your preamble to see if that helps.

\mybigspace Ordinarily, \mymarginpar{\textbf{\textbackslash
    fullciteincite}} you can use one of the methods discussed in
section~\ref{sec:authrelated} to present, in a reference list, a full
citation of a work related in some way to the current entry.  The new
command \cmd{citeincite} (above) allows you to present short citations
to other works within selected fields of the parent entry, or indeed
as part of comments inside other in-text citations, while attempting
to minimize the problems caused by citations within other
citations. The \cmd{fullciteincite} command works similarly for full
citations, but really is designed for use only in \textsf{annotation}
fields, the idea being that you might want to refer in an annotated
reference list to works that you haven't cited in the main body of
your text. The multicite command \cmd{fullciteincites} is also
available. It may be worth noting here that, because the full
reference list drivers don't contain code for printing
\textsf{postnote} fields, any postnotes that you provide will be
silently ignored. Inconsistently enough, \textsf{prenote},
\textsf{multiprenote}, and \textsf{multipostnote} fields will appear.

\mybigspace Arne \mymarginpar{\textbf{\textbackslash gentextcite}}
Skj{\ae}rholt requested, for the author-date styles, a variant of the
\cmd{textcite} command that presented the author's name in the
genitive case in running text, thereby simplifying certain syntactic
constructions (15.25).  The \cmd{gentextcite} command, in effect,
provides a way to include almost anything in between the name and the
parenthesized date in a \cmd{textcite}, so its use may well not be
limited to the possessive.  In most respects it behaves exactly like
\cmd{textcite}, on which see below.  The difference is that I've added
a new optional field to the front of the command to allow you to
choose which declensional ending to add to the name.  If you don't
specify this field, you'll get the standard English \enquote{\,'s\,}.
If you want something different, you'll need to present a third option
to the command, like so:\
\cmd{gentextcite[<ending>][][]\{entry:key\}}.  You must include the
two further sets of square brackets, because with only one set it
will, as with other citation commands, be interpreted as a
\textsf{postnote}, and with two a \textsf{prenote} and a
\textsf{postnote}.  There is a \cmd{gentextcites} command as well,
though currently you can only specify one genitival ending for all
keys, like so:\
\cmd{gentextcites[<ending>]()()[][]\{entry:key1\}\{entry:key2\}},
though if you don't have a \textsf{pre-} or \textsf{postnote} to the
first citation you can make do with
\cmd{gentextcites[<ending>]()\{entry:key1\}\{entry:key2\}}.

\mylittlespace The syntax of multiple authors' names in running text
is unpredictable.  There is currently no way to add the genitival
ending to all the names attached to a single citation key, so it will
only appear at the end of a group of names in such a case.  (This is
in keeping with the usual syntax when referring to a multi-author
work, at least in English.)  When using \cmd{gentextcites}, however,
you can control whether the ending appears after the name(s) attached
to each citation key, or whether it only appears after the names
attached to the last key.  By default, it only appears after the last,
but the \texttt{genallnames} preamble and/or entry option set to
\texttt{true} will attach the ending to each key's name(s).  When
using one citation command to cite more than one work by the same
author, it is the \emph{first} occurrence of the name which
\textsf{biblatex} prints, eliding subsequent ones.  In order to get
the possessive ending on that name you'll need to set
\texttt{genallnames} to \texttt{true}.

\mybigspace In \mymarginpar{\textbf{\textbackslash textcite}} standard
\textsf{biblatex} this command searches first for a
\textsf{labelname}, usually taken from the \textsf{author} or
\textsf{shortauthor} field, then uses the \textsf{shorthand} field if
the former doesn't exist.  Because of the way the Chicago author-date
specification recommends handling abbreviations, I have switched this
around, and the command searches for a \textsf{shorthand} first.
This holds also for the multicite form \textbf{\textbackslash
  textcites}, though both commands revert to their standard
\textsf{biblatex} behavior when you give the \texttt{cmslos=false}
option in the preamble.

\subsection{Package Options}
\label{sec:opts:authdate}

\subsubsection{Pre-set \textsf{biblatex} Options}
\label{sec:preset:authdate}

Although a quick glance through \textsf{biblatex-chicago.sty} will
tell you which \textsf{biblatex} options the package sets for you, I
thought I might gather them here also for your perusal.  These
settings are, I believe, consistent with the specification, but you
can alter them in the options to \textsf{biblatex-chicago} in your
preamble or by loading the package using \verb+\usepackage+
\verb+[style=chicago-authordate]{biblatex}+, which gives you the
\textsf{biblatex} defaults unless you redefine them yourself inside
the square brackets.

\mylittlespace \textsf{Biblatex-chicago-authordate}
\mymarginpar{\texttt{autocite=\\inline}} and \textsf{authordate-trad}
place references in parentheses by default.

\mybigspace The \mymarginpar{\texttt{citetracker=\\true}} citetracker
for the \cmd{ifciteseen} test is enabled globally.

\mybigspace The \mymarginpar{\texttt{alldates=comp}} specification
calls for the long format when presenting dates, slightly shortened
when presenting date ranges.  The new style option
\texttt{compressyears}, enabled by default, means that year ranges are
now compressed automatically according to the \emph{Manual's}
instructions (9.64; section~\ref{sec:authpreset}).

\enlargethispage{-\baselineskip}

\mylittlespace In \mymarginpar{\texttt{alltimes=12h}} entries which
print time stamps, they will, when the stamp is part of a
\textsf{date}, \textsf{eventdate}, or \textsf{origdate}, appear in
12-hour format, i.e., \enquote{4:45 p.m.}  Stamps that are part of a
\textsf{urldate} are, by default, controlled by the \texttt{urltime}
option, which is set to \texttt{24h}.  See that option below, and
table~\ref{ad:date:extras}.

\mylittlespace This \mymarginpar{\texttt{datecirca=true}} option
enables \textsf{biblatex's} enhanced \enquote{circa} date
specification, which given a \textsf{date} like \verb+1989~+ will
print [ca.\ 1989].  Cf.\ table~\ref{ad:date:extras}.

\mylittlespace This \mymarginpar{\texttt{dateuncertain=\\true}}
option enables \textsf{biblatex's} enhanced \enquote{uncertain} date
specification, which given a \textsf{date} like \verb+1989?+ will
print [1989?].  A field like \verb+1989%+ is both \enquote{circa}
\emph{and} \enquote{uncertain,} like so: [ca.\ 1989?].  Cf.\
table~\ref{ad:date:extras}.

\mylittlespace This \mymarginpar{\texttt{datezeros=false}} ensures
that leading zeros don't appear in date specifications.

\mylittlespace This \mymarginpar{\texttt{ibidtracker=\\constrict}}
enables an \emph{ibidem} mechanism in citations, but only in the most
strictly-defined circumstances.  The Chicago author-date style doesn't
print \enquote{Ibid} in citations, but in general a repeated citation
on the same page will print only the page reference.  Technically,
this should only occur when a source is cited \enquote{more than once
  in one paragraph} (15.27), so you can use the \cmd{citereset}
command from \textsf{biblatex} to achieve the greatest compliance, as
the package only offers automatic resetting on part, chapter, section,
and subsection boundaries, while \textsf{biblatex-chicago}
automatically resets the tracker at page breaks.  (Cf.\
\textsf{biblatex.pdf} \S~3.1.2.1.)  Whenever there might be any
ambiguity, \textsf{biblatex} should default to printing a more
informative reference.

\mylittlespace If you are going to repeat a source, make sure that the
cite command provides a postnote --- you'll no longer get any annoying
empty parentheses, but you will get another standard citation, which
may add too much clutter.

\mylittlespace This \mymarginpar{\texttt{labeldateparts\\=true}}
option tells \textsf{biblatex} to provide the special
\textsf{labelyear} and \textsf{extradate} fields for author-date
styles.  (This is the option formerly known as \texttt{labelyear} and
then \texttt{labeldate}, both of which are obsolete.)

\mylittlespace These \mymarginpar{\textsf{\texttt{maxbibnames\\=10\\
      minbibnames\\=7}}} two options control the number of names
printed in the list of references when that number exceeds 10.  These
numbers follow the recommendations of the \emph{Manual} (14.76, 15.9),
and they are different from those for use in citations.  Please see
section~\ref{sec:otherhints:auth} below (and the file
\href{file:cms-dates-intro.pdf}{\textsf{cms-dates-intro.pdf}}) for
hints on dealing with entries with more than three authors.

\mylittlespace This \mymarginpar{\texttt{pagetracker=\\true}} enables
page tracking for the \cmd{iffirstonpage} and \cmd{ifsamepage}
commands for controlling, among other things, the \emph{ibidem}
mechanism.  It tracks individual pages if \LaTeX\ is in oneside mode,
or whole spreads in twoside mode.

\mylittlespace This \mymarginpar{\texttt{punctfont=\\true}} fixes a
minor problem with punctuation in titles, ensuring that the colon
between a title and a subtitle appears in the correct, matching font.

%\enlargethispage{-\baselineskip}

\mylittlespace This \mymarginpar{\texttt{related=true}} is the
standard \textsf{biblatex} bibliography option, and it enables the use
of \textsf{related} functionality in the list of references.  I have
added an entry option, as well, so if you set this to \texttt{false}
in your preamble, in the \textsf{options} field, or in the
\textsf{relatedoptions} field, you can make the package ignore the
\textsf{related} fields.

\mylittlespace This \mymarginpar{\texttt{sortcase=\\false}} turns off
the sorting of uppercase and lowercase letters separately, a practice
which the \emph{Manual} doesn't appear to recommend.

\mylittlespace This \mymarginpar{\texttt{sorting=cms}} setting takes
advantage of the \cmd{DeclareSortingTemplate} command provided by
\textsf{biblatex} and \textsf{Biber}, in effect implementing a default
sorting order in the list of references tailored to comply with the
Chicago author-date specification.  Please see the documentation of
\cmd{DeclareSortingTemplate} in section~\ref{sec:authformopts}, below.

\mylittlespace If \mymarginpar{\texttt{timezones=true}} you provide a
timezone for a time stamp, usually using one of the \textsf{timezone}
fields, this option ensures it will be printed.

\mylittlespace This \mymarginpar{\texttt{uniquelist=\\minyear}} option
enables \textsf{biblatex-chicago-authordate} to disambiguate entries
which have more than three \textsf{authors}, but which differ
\emph{after} the first name in the list.  This will only occur when
two such entries have the same \textsf{year} (15.29).  The option is
\textsf{Biber}-only, like the following, which means that this
next-generation \textsc{Bib}\TeX\ replacement is required for the
author-date styles.  Please see
\href{file:cms-dates-intro.pdf}{\textsf{cms-dates-intro.pdf}} and
section~\ref{sec:otherhints:auth}, below, for further details.

\mylittlespace This \mymarginpar{\texttt{uniquename=\\minfull}}
enables the package to distinguish different authors who share a
surname, using initials in the first instance, and whole names if
initials aren't enough (15.22).  The option is \textsf{Biber}-only,
like the previous one.

\mylittlespace In \mymarginpar{\texttt{urltime=24h}} entries with
\textsf{urldate} fields containing time stamps, that stamp will by
default appear in 24-hour format, i.e., \enquote{16:45.}  Cf.\
\texttt{alltimes}, above, \texttt{urlstamp} in
section~\ref{sec:authpreset} below, and table~\ref{ad:date:extras}.

\mylittlespace In
\mymarginpar{\texttt{[standard]\\useeditor=false\\usenamec=false}}
\textsf{standard} entries any editors' or compilers' names appear
after the title, according to 14.259, so these entry-type-specific
options encode this.  You can, of course, override these defaults in
your preamble, should you deem it necessary.

\mylittlespace This \mymarginpar{\texttt{usetranslator\\=true}}
enables automatic use of the \textsf{translator} at the head of
entries in the absence of an \textsf{author} or an \textsf{editor}.
In the list of references, the entry will be alphabetized by the
translator's surname.  You can disable this functionality on a
per-entry basis by setting \texttt{usetranslator=false} in the
\textsf{options} field.  Cf.\ silver:gawain.

\subsubsection*{Other \textsf{biblatex} Formatting Options}
\label{sec:authformopts}

I've chosen defaults for many of the general formatting commands
provided by \textsf{biblatex}, including the vertical space between
items in the list of references and between items in the list of
shorthands (\cmd{bibitemsep} and \cmd{lositemsep}).  I define many of
these in \textsf{biblatex-chicago.sty}, and of course you may want to
redefine them to your own needs and tastes.  It may be as well you
know that the \emph{Manual} does state a preference for two of the
formatting options I've implemented by default: the 3-em dash as a
replacement for repeated names in the list of references (15.17--19,
and just below); and the formatting of note numbers, both in the main
text and at the bottom of the page / end of the essay (superscript in
the text, in-line in the notes; 14.24).  The code for this last
formatting is also in \textsf{biblatex-chicago.sty}, and I've wrapped
it in a test that disables it if you are using the \textsf{memoir}
class, which I believe has its own commands for defining these
parameters.  You can also disable it by using the \texttt{footmarkoff}
package option, on which see below.

\enlargethispage{\baselineskip}

\mylittlespace Gildas Hamel pointed out that my default definition, in
\textsf{biblatex-chicago.sty}, of \textsf{biblatex's}
\cmd{bibnamedash} didn't work well with many fonts, leaving a line of
three dashes separated by gaps.  He suggested an alternative, which
I've adopted, with a minor tweak to make the dash thicker, though you
can toy with all the parameters to find what looks right with your
chosen font.  The default definition is:
\cmd{renewcommand*\{\textbackslash bibname\-dash\}\{\textbackslash
  rule[.4ex]\{3em\}\{.6pt\}\}}.

\mylittlespace At \mymarginpar{\texttt{losnotes}
  \&\\\texttt{losendnotes}} the request of Kenneth Pearce, I have
added two \texttt{bibenvironments} to
\textsf{chicago-author\-date.bbx}, for use with the \texttt{env}
option to the \cmd{printshorthands} command.  The first,
\texttt{losnotes}, is designed to allow a list of shorthands to appear
inside footnotes, while \texttt{losendnotes} does the same for
endnotes.  Their main effect is to change the font size, and in the
latter case to clear up some spurious punctuation and white space that
I see on my system when using endnotes.  (You'll probably also want to
use the option \texttt{heading=none} in order to get rid of the
[oversized] default, providing your own within the \cmd{footnote}
command.)  If you use a command like
\verb+\printbiblist{shortjournal}+ to print a list of journal
abbreviations, you can use the \texttt{sjnotes} and
\texttt{sjendnotes} \texttt{bibenvironments} in exactly the same way.
Please see the documentation of \textsf{shorthand} and
\textsf{shortjournal} in section~\ref{sec:fields:authdate} above for
further options available to you for presenting and formatting these
two types of \texttt{biblist}.

\mylittlespace The next-generation backend \textsf{Biber} and
\textsf{biblatex} offer enhanced functionality in many areas,
including the next three declarations.  If the default definitions
don't work well for you, you can redefine all of them in your document
preamble --- see \textsf{biblatex.pdf} \S\S~4.5.6 and 4.5.10.

\mylittlespace This \mymarginpar{\cmd{Declare-}\\\texttt{Labelname}}
option allows you to add name fields for consideration when
\textsf{biblatex} is attempting to find a shortened name for in-text
citations.  This, for example, allows a compiler (=\textsf{namec}) to
appear in citations without any other intervention from the user,
rather than requiring a \textsf{shortauthor} field as previous
releases of \textsf{biblatex-chicago} did.  The default definition
currently is
\texttt{\{shortauthor,author,shorteditor,namea,editor,nameb,%
  translator,namec\}}.

\mylittlespace This \mymarginpar{\cmd{Declare-}\\\texttt{Labeldate}}
option allows you to alter the order in which \textsf{Biber} and
\textsf{biblatex} search for the year to use both in citations and at
the head of entries in the list of references.  This will also be the
year to which an alphabetical suffix will be appended when an author
has published more than one work in the same year, and the year by
which works will be sorted in the list of references.  In the default
configuration, a year will be searched for in the order \textsf{date,
  eventdate, origdate, urldate}.  This generally suits the Chicago
author-date styles well, except for two situations.  First, when a
reference apparatus contains many entries with multiple dates, it may
be simplest to promote the \textsf{origdate} to the head of the list,
which you can do using the \texttt{cmsdate} preamble option.  This
changes the order to \textsf{origdate, date, eventdate, urldate}.
Second, in \textsf{music} and \textsf{video} entries, and,
exceptionally, some \textsf{review} entries, the general rule is to
emphasize the earliest date.  For these three entry types, then,
\cmd{DeclareLabeldate} uses the order \textsf{eventdate, origdate,
  date, urldate}.  See \texttt{avdate} in
section~\ref{sec:authpreset}, \texttt{cmsdate} in
section~\ref{sec:authuseropts}, and the \textbf{date} docs in
section~\ref{sec:fields:authdate}.

\mylittlespace The
\mymarginpar{\cmd{Declare-}\\\texttt{Sorting-}\\\texttt{Template}}
third \textsf{Biber} enhancement I have implemented allows you to
include almost any field whatsoever in \textsf{biblatex's} sorting
algorithms for the list of references, so that a great many more
entries will be sorted correctly automatically rather than requiring
manual intervention in the form of a \textsf{sortkey} field or the
like.  Code in \textsf{biblatex-chicago.sty} sets the
\textsf{biblatex} option \texttt{sorting=cms}, which is a custom
scheme, basically a Chicago-specific variant of the default
\texttt{nyt}.  You can find its definition in
\textsf{chicago-authordate.cbx}.  (Please note that it uses the
\textsf{labelyear} as its main year component, which should help
improve the automatic sorting of entries by the same \textsf{author}.)

\mylittlespace The advantages of this scheme are, specifically, that
any entry headed by one of the supplemental name fields
(\textsf{name[a-c]}), a \textsf{manual} or a \textsf{standard} entry
headed by an \textsf{organization}, or an \textsf{article} or
\textsf{review} entry with an \textsf{entrysubtype} and headed by a
\textsf{journaltitle} will no longer need a \textsf{sortkey} set.
Further, the \textsf{biblatex} \texttt{use<name>=false} options will
remove any name field from the sorting order, again reducing the need
for user intervention.  The main disadvantage should only occur very
rarely.  In \textsf{author}-less \textsf{article} and \textsf{review}
entries without an \textsf{entrysubtype}, the \textsf{title} will
appear instead of the \textsf{journaltitle}, and since the latter
appears before the former in the sorting scheme, you'll need a
\textsf{sortkey} for proper alphabetization.

\subsubsection{{Pre-set \textsf{chicago} Options}}
\label{sec:authpreset}

At \mymarginpar{\texttt{bookpages=\\true}} the request of Scot Becker,
I have included this rather specialized option, which controls the
printing of the \textsf{pages} field in \textsf{book} entries.  Some
bibliographic managers, apparently, place the total page count in that
field by default, and this option allows you to stop the printing of
this information in the reference list.  It defaults to true, which
means the field is printed, but it can be set to false either in the
preamble, for the whole document or for specific entry types, or on a
per-entry basis in the \textsf{options} field (though rather than use
this latter method it would make sense to eliminate the \textsf{pages}
field from the affected entries).

%\enlargethispage{\baselineskip}

\mylittlespace This \colmarginpar{\texttt{doi=true}} option controls
whether any \textsf{doi} fields present in the .bib file will be
printed in the reference list.  At the request of Daniel Possenriede,
and keeping in mind the \emph{Manual's} preference for this field
instead of a \textsf{url} (14.7--8), I have added a third switch,
\texttt{only}, which prints the \textsf{doi} if it is present and the
\textsf{url} only if there is no \textsf{doi}.  Ryo Furue more
recently requested a way to suppress the \textsf{urldate} when using
only the \textsf{doi}, so I've added the \mycolor{\texttt{onlynd}}
switch to do this.  The package default remains the same, however ---
it defaults to true, which will print both \textsf{doi} and
\textsf{url} if both are present.  The option can be set to
\texttt{only}, \mycolor{\texttt{onlynd}}, or to \texttt{false} either
in the preamble, for the whole document or for specific entry types,
or on a per-entry basis in the \textsf{options} field.  In
\textsf{online} entries, the \textsf{doi} field will always be
printed, but the \texttt{only} switch will still eliminate any
\textsf{url}, and \mycolor{\texttt{onlynd}} will still eliminate both
the \textsf{url} and the \textsf{urldate}.

\mylittlespace This \mymarginpar{\texttt{eprint=true}} option controls
whether any \textsf{eprint} fields present in the .bib file will be
printed in the list of references.  It defaults to true, and can be
set to false either in the preamble, for the whole document or for
specific entry types, or on a per-entry basis, in the \textsf{options}
field.  In \textsf{online} entries, the \textsf{eprint} field will
always be printed.

\mylittlespace This \mymarginpar{\texttt{isbn=true}} option controls
whether any \textsf{isan}, \textsf{isbn}, \textsf{ismn},
\textsf{isrn}, \textsf{issn}, and \textsf{iswc} fields present in the
.bib file will be printed in the list of references.  It defaults to
true, and can be set to false either in the preamble, for the whole
document or for specific entry types, or on a per-entry basis, in the
\textsf{options} field.

\mylittlespace Once \mymarginpar{\texttt{numbermonth\\=true}} again at
the request of Scot Becker, I have included this option, which
controls the printing of the \textsf{month} field in all the
periodical-type entries when a \textsf{number} field is also present.
Some bibliographic software, apparently, always includes the month of
publication even when a \textsf{number} is present.  When all this
information is available the \emph{Manual} (14.171) prints everything,
so this option defaults to true, which means the field is printed, but
it can be set to false either in the preamble, for the whole document
or for specific entry types, or on a per-entry basis in the
\textsf{options} field.

\mylittlespace This \mymarginpar{\texttt{url=true}} option controls
whether any \textsf{url} fields present in the .bib file will be
printed in the reference list.  It defaults to true, and can be set to
false either in the preamble, for the whole document or for specific
entry types, or on a per-entry basis, in the \textsf{options} field.
Please note that, as in standard \textsf{biblatex}, the \textsf{url}
field is always printed in \textsf{online} entries, regardless of the
state of this option.

\mylittlespace This \mymarginpar{\texttt{urlstamp=true}} option
controls whether any \textsf{urltime} fields, included as part of the
\textsf{urldate}, will be printed in citations and reference list.  It
defaults to true, and can be set to false either in the preamble, for
the whole document or for specific entry types, or on a per-entry
basis in the \textsf{options} field.  Please note that, unlike the
\texttt{url} option, this option \emph{does} control what is printed
in \textsf{online} entries.

\mylittlespace This \mymarginpar{\texttt{includeall=\\true}} is the
one option that rules the seven preceding, either printing all the
fields under consideration --- the default --- or excluding all of
them.  It is set to \texttt{true} in \textsf{chicago-authordate.cbx},
but you can change it either in the preamble for the whole document or
for specific entry types, or in the \textsf{options} field of
individual entries.  The seven individual options above are similarly
available in the same places, for finer-grained control.  The
rationale for all of these options is the availability of
bibliographic managers that helpfully present as much data as
possible, in every entry, some of which may not be felt to be entirely
necessary.  Setting \texttt{includeall} to \texttt{true} probably
works just fine for those compiling their .bib databases by hand, but
others may find that some automatic pruning helps clear things up, at
least to a first approximation.  Some per-type or per-entry work
afterward may then polish up the details.  If you find that you need
control over fields that aren't included among these options, I have
provided the \cmd{suppressbibfield} command for your preamble, as
suggested by Jan David Hauck.  It is in fact a user interface to the
source mapping feature of \textsf{biblatex}, and it is something of a
nuclear option, preventing fields from even appearing in the .bbl file
generated by \textsf{biber} from your .bib database.  See the
\cmd{suppressbibfield} command in
section~\ref{sec:formatting:authdate} and the source mapping docs in
\textsf{biblatex.pdf} \S~4.5.3.

\paragraph*{\protect\mymarginpar{\texttt{avdate=true}}}
\label{sec:ad:avdate}\vspace{-.5\baselineskip}
For \textsf{music} and \textsf{video} entries, the \emph{Manual}
(14.263, 15.57) strongly recommends both that you provide a recording,
release, or broadcast date for your references and also that this
earlier date should appear in citations and at the head of reference
list entries.  In the default setting of \cmd{DeclareLabeldate},
\textsf{biblatex} searches for dates in the following order:
\textsf{year, eventyear, origyear, urlyear}.  This option changes the
default ordering in \textsf{music} and \textsf{video} entries to the
following: \textsf{eventyear, origyear, year, urlyear}.
\textsf{Review} entries presenting on-line comments have similar
needs, as do \textsf{standard} entries, so the same reordering applies
to those entry types, too.  If you simply want to apply the defaults
to these four entry types, you can use \texttt{avdate=false} in the
options when loading \textsf{biblatex-chicago}.  If, however, you want
to tailor the algorithm to your own needs, then you can use
\cmd{DeclareLabeldate} commands in your preamble.  Please be aware,
however, that some parts of the style hard-code the search syntax, and
although they take account of the \texttt{avdate} setting, if you use
your own definitions of \cmd{DeclareLabeldate} the results may, in
some corner cases, surprise.  Please see \textsf{music},
\textsf{review}, \textsf{standard}, and \textsf{video} in
section~\ref{sec:types:authdate}; \textsf{date}, \textsf{eventdate},
\textsf{origdate}, and \textsf{urldate} in
section~\ref{sec:fields:authdate}; and \cmd{DeclareLabeldate} in
section~\ref{sec:authformopts}.

\mylittlespace These
\mymarginpar{\texttt{bibannotesep\\=vpar\\citeannotesep\\=period}}
options define the relation of the \textsf{annotation} field to the
main entry, \texttt{bibannotesep} doing so in the reference list and
\texttt{citeannotesep} in long (legal) notes.  (The
\texttt{annotation} option in section~\ref{sec:authuseropts}
determines where, if anywhere, the field will appear.)  Both options
have the same set of keys, though they have different default settings
if you don't define them yourself.  The possible values are:

\begin{description}
\setlength{\parskip}{-4pt}
\item[\qquad none] = no punctuation at all.
\item[\qquad space] = \cmd{addspace}
\item[\qquad comma] = \verb+\addcomma\addspace+
\item[\qquad period] = \verb+\addperiod\addspace+
\item[\qquad colon] = \verb+\addcolon\addspace+
\item[\qquad semicolon] = \verb+\addsemicolon\addspace+
\item[\qquad par:] This starts a new paragraph on the next line.  Page
  breaking is strongly inhibited before the \textsf{annotation}.
\item[\qquad vpar:] This starts a new paragraph, and also inserts some
  vertical space before it.  In long notes that vertical space is
  minimal (1 pt), while in the reference list it creates a blank line.
  Page breaking is strongly inhibited before the \textsf{annotation}.
\item[\qquad parbreak:] This is the same as \texttt{par}, but it allows
  a page break to occur between the main entry and the start of the
  \textsf{annotation}.
\item[\qquad vparbreak:] This is the same as \texttt{vpar}, but it also
  allows a page break between the main entry and the
  \textsf{annotation}.
\end{description}

Please note that both of these options are available in the preamble
both globally and per-type, and also in the \textsf{options} field of
individual entries.  Each defines a command (\cmd{bibannotesep} and
\cmd{citeannotesep}) which appears in the \textsf{annotation} field's
formatting directive, so it's possible to redefine these commands in
your preamble if you have needs that the available values don't
address.  (You can also try sending an email to encourage me to add
other keys.)  Please also keep in mind that \texttt{bibannotesep}
interacts with the \texttt{entrybreak} and \textsf{formatbib} options
in section~\ref{sec:authuseropts}, below, to determine the general
layout of the reference list.  Depending on the settings of those
options, changing the \texttt{bibannotesep} from entry to entry may
not work out well.

\mylittlespace At \mymarginpar{\texttt{booklongxref=\\true}} the
request of Bertold Schweitzer, I have included two options for
controlling whether and where \textsf{biblatex-chicago} will print
abbreviated references when you cite more than one part of a given
collection or series.  This option controls whether multiple
\textsf{book}, \textsf{bookinbook}, \textsf{collection}, and
\textsf{proceedings} entries which are part of the same collection
will appear in this space-saving format.  The parent collection itself
will usually be presented in, e.g., a \textsf{book},
\textsf{bookinbook}, \textsf{mvbook}, \textsf{mvcollection}, or
\textsf{mvproceedings} entry, and using \textsf{crossref} or
\textsf{xref} in the child entries will allow such presentation
depending on the value of the option:

\begin{description}
\item[\qquad true:] This is the default.  If you use \textsf{crossref}
  or \textsf{xref} fields in these entry types, by default you will
  \emph{not} get any abbreviated citations in the reference list.
\item[\qquad false:] You'll get abbreviated citations in these entry
  types in the reference list.
\item[\qquad notes, bib:] These two options are carried over from the
  notes \&\ bibliography style; here they are synonymous with
  \texttt{false} and \textsf{true}, respectively.
\end{description}

This option can be set either in the preamble or in the
\textsf{options} field of individual entries.  For controlling the
behavior of \textsf{inbook}, \textsf{incollection},
\textsf{inproceedings}, and \textsf{letter} entries, please see
\texttt{longcrossref}, below, and also the documentation of
\textsf{crossref} in section~\ref{sec:fields:authdate}.

\mylittlespace This \mymarginpar{\texttt{cmslos=true}} option alters
\textsf{biblatex's} standard behavior when processing the
\textsf{shorthand} field.  Chicago's author-date style only seems to
recommend the use of shorthands as abbreviations for long authors'
names, particularly institutional names, which means the shorthand
will replace only the name part in citations rather than the whole
citation (15.37; bsi\hc abbreviation, iso:electrodoc).  Recent editions
suggest placing the abbreviation at the head of the entry, followed by
its expansion inside parentheses, an arrangement automatically
provided by \textsf{biblatex-chicago-authordate} when you use the
\textsf{shorthand} field, assuming you retain the default setting of
this option.  Please note that you can still print a list of
shorthands if you wish, and you can also get back something
approaching the \enquote{standard} behavior of shorthands if you give
the \texttt{cmslos=false} option to \textsf{biblatex-chicago} in your
document preamble.  Cf.\ section~\ref{sec:fields:authdate},
s.v. \enquote{\textbf{shorthand}} above, and also
\href{file:cms-dates-intro.pdf}{\textsf{cms-dates-intro.pdf}}.

\mylittlespace The \mymarginpar{\texttt{compressyears\\=true}}
\emph{Manual} has long recommended (9.64, 15.41), as a space-saving
measure, the compression of year ranges when presenting dates.  I
have, finally, implemented this in the current release, and have made
it the default, which you can change in your document preamble.
Please note that the rules for compressing years are different from
those for compressing other numbers (e.g., page numbers), and also
that the compression code is in \textsf{biblatex-chicago.sty}, which
will have to be loaded for this option to make any difference.  Cf.\
table~\ref{ad:date:extras}.

\mylittlespace Roger
\mymarginpar{\texttt{ctitleaddon=\\comma\\ptitleaddon=\\period}} Hart
requested a way to control the punctuation printed before the
\textsf{titleaddon}, \textsf{booktitleaddon}, and
\textsf{maintitleaddon} fields.  By default, this is
\cmd{addcomma\cmd{addspace}} (\textsf{cti\-tleaddon}) for nearly all
\textsf{book-} and \textsf{maintitleaddons} in the list of references,
while \cmd{addperiod\cmd{addspace}} (\textsf{ptitleaddon}) is the
default before most \textsf{titleaddons} there.  If the punctuation
printed isn't correct for your needs, you can set the relevant option
either in the preamble or in individual entries.  (Cf.\
coolidge:speech and schubert:muellerin.) The accepted option keys are:

\begin{description}
\setlength{\parskip}{-4pt}
\item[\qquad none] = no punctuation at all
\item[\qquad space] = \cmd{addspace}
\item[\qquad comma] = \cmd{addcomma\cmd{addspace}}
\item[\qquad period] = \cmd{addperiod\cmd{addspace}}
\item[\qquad colon] = \cmd{addcolon\cmd{addspace}}
\item[\qquad semicolon] = \cmd{addsemicolon\cmd{addspace}}
\end{description}

If you need something a little more exotic, you can directly
\cmd{renewcommand} either \cmd{ctitleaddonpunct} or
\cmd{ptitleaddonpunct} (or both) in your preamble, but it's worth
remembering that the redefinition will hold for all instances, unless
you use the \textsf{options} field in your other entries with a
\textsf{titleaddon} field.  A simpler solution might be to set the
relevant option to \texttt{none} in your entry and then include the
punctuation in the \textsf{titleaddon} field itself.

\mylittlespace Constanza Cordoni \mymarginpar{\texttt{dashed=true}}
has requested a way to turn off the 3-em dash for replacing repeated
names in the reference list, and the \emph{Manual} admits that some
publishers prefer this, as the dash can carry with it certain
inconveniences, especially for electronic formats (15.17).  Some of
\textsf{biblatex's} standard styles have a \texttt{dashed} option, so
for compatibility purposes I've provided the same.  By default I have
set it to print the name dash, but you can set \texttt{dashed=false}
globally, per type, or per entry to repeat names as and when required.

\mylittlespace If \mymarginpar{\texttt{hidevolumes=\\true}} both a
\textsf{volume} and a \textsf{volumes} field are present, as may occur
particularly in cross-referenced entries, then
\textsf{biblatex-chicago} will ordinarily suppress the
\textsf{volumes} field.  In some instances, when a \textsf{maintitle}
is present, this may not be the desired result.  In this latter case,
if the \textsf{volume} appears before the \textsf{maintitle}, this
option, set to \texttt{true} by default, controls whether to print the
\textsf{volumes} field after that title or not.  Set it to
\texttt{false} either in the preamble or in the \textsf{options} field
of your entry to have it appear after the \textsf{maintitle}.

\mylittlespace This \mymarginpar{\texttt{journalabbrev\\=notes}}
option controls the printing of the \textsf{shortjournal} field in
place of the \textsf{journaltitle} field in citations and reference
lists.  It is set to \texttt{notes} by default, so as shipped
\textsf{biblatex-chicago-authordate} will print such fields only in
citations, but you can set it, either in the preamble or in individual
entries, to one of three other values: \texttt{true} prints the
abbreviated form both in citations and reference lists, \texttt{bib}
in the reference list only, and \texttt{false} in neither.  Please
note that in \textsf{periodical} entries the \textsf{title} and
\textsf{shorttitle} fields behave in exactly the same manner.  For
more details, see the documentation of \textbf{shortjournal} in
section~\ref{sec:fields:authdate}, above.

\mylittlespace I \colmarginpar{\texttt{jtitleaddon=\\space}} have
added the standard \textsf{biblatex}
\mycolor{\textsf{journaltitleaddon}} field to the \textsf{article} and
\textsf{review} entry types, and also the \textsf{titleaddon} field to
the \textsf{periodical} type, fields that may, for example, be
particularly useful when you want to provide the original form of a
translated journal title.  The \mycolor{\texttt{jtitleaddon}} option
controls the separator between the main title and the addon, as with
the \texttt{ctitleaddon} and \texttt{ptitleaddon} options, above, and
like them is settable globally, per type, or per entry.  The possible
settings are the same as for those options, but the default is a
\texttt{space}.  You can redefine \cmd{jtitleaddonpunct} directly if
you have more unusual needs.

\mylittlespace This \mymarginpar{\texttt{longcrossref=\\false}} is
the second option, requested by Bertold Schweitzer, for controlling
whether and where \textsf{biblatex-chicago} will print abbreviated
references when you cite more than one part of a given collection or
series.  It controls the settings for the entry types more-or-less
authorized by the \emph{Manual}, i.e., \textsf{inbook},
\textsf{incollection}, \textsf{inproceedings}, and \textsf{letter}.
The mechanism itself is enabled by multiple \textsf{crossref} or
\textsf{xref} references to the same parent, whether that be, e.g., a
\textsf{collection}, an \textsf{mvcollection}, a \textsf{proceedings},
or an \textsf{mvproceedings} entry.  Given these multiple cross
references, the presentation in the reference apparatus will be
governed by the following options:

\begin{description}
\item[\qquad false:] This is the default.  If you use
  \textsf{crossref} or \textsf{xref} fields in the four mentioned
  entry types, you'll get the abbreviated entries in the reference
  list.
\item[\qquad true:] You'll get no abbreviated citations of these entry
  types in the reference list.
\item[\qquad none:] This switch is special, allowing you with one
  setting to provide abbreviated citations not just of the four entry
  types mentioned but also of \textsf{book}, \textsf{bookinbook},
  \textsf{collection}, and \textsf{proceedings} entries.
\item[\qquad notes, bib:] These two options are carried over from the
  notes \&\ bibliography style; here they are synonymous with
  \texttt{false} and \textsf{true}, respectively.
\end{description}

This option can be set either in the preamble or in the
\textsf{options} field of individual entries.  For controlling the
behavior of \textsf{book}, \textsf{bookinbook}, \textsf{collection},
and \textsf{proceedings} entries, please see \texttt{booklongxref},
above, and also the documentation of \textsf{crossref} in
section~\ref{sec:fields:authdate}.

\mylittlespace This \mymarginpar{\texttt{nameaddonsep\\=space}} option
sets the punctuation which appears before the \textsf{nameaddon} field
in all entry types except \textsf{customc}.  You can set it globally,
per type, or per entry, using one of the six following keys:

\begin{description}
\setlength{\parskip}{-4pt}
\item[\qquad \texttt{space}] = \cmd{addspace}.  This is the default.
\item[\qquad \texttt{none}] = no separator at all.  It presumes that
  you will include one in the \textsf{nameaddon} field itself.
\item[\qquad \texttt{colon}] = \verb+\addcolon\addspace+.
\item[\qquad \texttt{comma}] = \verb+\addcomma\addspace+.
\item[\qquad \texttt{period}] = \verb+\addperiod\addspace+.
\item[\qquad \texttt{semicolon}] = \verb+\addsemicolon\addspace+.
\end{description}

Cf.\ \texttt{nameaddon} and \texttt{nameaddonformat} in
section~\ref{sec:authuseropts}.

\mylittlespace This \mymarginpar{\texttt{nodates=true}} option means
that for all entry types except \textsf{misc} and \textsf{dataset}
\textsf{biblatex-chicago} will automatically provide
\verb+\bibstring{nodate}+ for any entry that doesn't otherwise provide
a date for citations and for the heads of entries in the list of
references.  If you set \texttt{nodates=false} either in your preamble
(for global or for per-type coverage) or in individual entries then
the package won't perform this substitution.  (The bibstring expands
to \enquote{\texttt{n.d.}} in English.)

\subsubsection{Style Options -- Preamble}
\label{sec:authuseropts}

These are parts of the specification that not everyone will wish to
enable.  All except the sixth, ninth, and thirteenth can be used even
if you load the package in the old way via a call to
\textsf{biblatex}, but most users can just place the appropriate
string(s) in the options to the \cmd{usepackage\{biblatex-chicago\}}
call in your preamble.

\mylittlespace \textsf{Biblatex-chicago}
\mymarginpar{\texttt{alwaysrange}} now implements \textsf{biblatex's}
enhanced date specifications, one part of which is the presentation of
decades and centuries not as year ranges but as localized strings like
\enquote{19th c.} or \enquote{1970s.}  The \texttt{alwaysrange} option
set to \texttt{true}, either in your preamble or in individual
entries, simply tells the package to present the year range instead.
This allows you to use the efficient enhanced notations in the
\textsf{date} field (\verb+{18XX}+ or \verb+{197X}+) without the
localized strings appearing, should you require it.  The two options
\texttt{centuryrange} and \texttt{decaderange} limit the same effect
to centuries and decades, respectively.  Please see
table~\ref{ad:date:extras}.

\mylittlespace At \mymarginpar{\texttt{annotation}} the request of
Emil Salim, I included in \textsf{biblatex-chicago} the ability to
produce annotated reference lists.  More recently, Moritz Wemheuer
brought to my attention a
\href{https://tex.stackexchange.com/questions/528374/moving-addendum-field-to-the-end-in-biblatex-chicago/540755#540755}{StackExchange}
question which suggested that the field might be useful in several
other contexts as well, so I've modified the \texttt{annotation}
option to allow the field to appear in the reference list
(\texttt{=bib} or \texttt{=true}, the default, if no string is given),
in long (legal) notes (\texttt{=notes}), in both (\texttt{=all}), or
in neither (\texttt{=false}).  You can now set the option in the
preamble both globally and per-type, and in the \textsf{options} field
of individual entries.  There are two new options
(\texttt{bibannotesep} and \texttt{citeannotesep}) to allow you to
choose the separator between the \textsf{annotation} and the rest of
the entry, and also two new options (\texttt{formatbib} and
\texttt{entrybreak}) to give you fine-grained control over the
presentation of the reference list as a whole, including an annotated
one.  Please have a look at the documentation for the latter two
options just below, for the former two options in
section~\ref{sec:authpreset}, and for the \textsf{annotation} field in
section~\ref{sec:fields:authdate}.  Please also see
\textsf{biblatex.pdf} \S~3.13.8 for details on how to use external
files to store annotations.

\mylittlespace In \mymarginpar{\texttt{authortitle}} a few contexts
--- classical references, some archival material, perhaps scientific
databases --- the provision of a date for citations may well be
impossible, irrelevant, or both.  The \textsf{entrysubtype} value
\texttt{classical} results in author-title citations for the entry
containing it, but it modifies punctuation in those citations in a way
that might be wrong for some sources, and it's also possible that you
may need the \textsf{entrysubtype} field for some other purpose yet
still wish to present author-title citations.  You can set
\texttt{authortitle} to \texttt{true} either for a specific entry type
in the preamble or in the \textsf{options} field of individual entries
to achieve this.  You can also use the citation commands \cmd{atcite}
and \cmd{atpcite}, instead, if that's more convenient.  The
\textsf{shorttitle} field provides a way to abbreviate long titles in
this context.  Please note that \textsf{biblatex-chicago} by default
sets this to \texttt{true} for the \textsf{dataset} entry type, so
you can set it to \texttt{false} if you want to present such entries
differently.

\mylittlespace Like \mymarginpar{\texttt{casechanger}} \texttt{mcite}
and \texttt{natbib}, this is a standard \textsf{biblatex} option which
\textsf{biblatex-chicago} simply passes through to that package.  In
\textsf{biblatex} it defaults to \texttt{auto}, but there were, and
possibly still are, cases when the old \texttt{latex2e} case-changing
code can work around some bugs when using, for example, citation
commands inside fields that have a case-changing element as an
automatic part of their formatting (\textsf{note},
\textsf{titleaddon}, \textsf{type}).  Cf.\ section~3.1.1 of
\textsf{biblatex.pdf}.

\mylittlespace This \mymarginpar{\texttt{centuryrange}} option works
just like \texttt{alwaysrange}, above, but only affects century
presentation, not decade.  Cf.\ table~\ref{ad:date:extras}.

\mylittlespace The \mymarginpar{\texttt{cmsbreakurl}} \emph{Manual}
gives fairly specific instructions about breaking URLs across lines
(14.18), so I have attempted to implement them by tweaking
\textsf{biblatex's} default settings, which are found in
\textsf{biblatex.def}.  In truth, I haven't succeeded in getting
\textsf{biblatex} flawlessly to follow all of the \emph{Manual's}
instructions, nor do the changes I have made work well in all
circumstances, being particularly unsightly if you happen to be using
the \textsf{ragged2e} package.  For these reasons, I have made my
changes dependent on a package option, \texttt{cmsbreakurl}, which you
can set in your preamble.  I have placed all of this code in
\textsf{biblatex-chicago.sty}, so if you load the package with a call
to \textsf{biblatex} instead, then URL line breaking will revert to
the \textsf{biblatex} defaults.  See
\href{file:cms-dates-sample.pdf}{\textsf{cms-dates-sample.pdf}} for a
lot of examples of what URLs look like when the option is set, and
also section~\ref{sec:loading:auth}, below.

\mylittlespace This \mymarginpar{\texttt{cmsdate}} option used
\emph{in the preamble} provides a method for simplifying the creation
of databases with a great many multi-date entries.  Despite warnings
in previous releases, users have plainly already been setting this
option in their preambles, so I thought I might at least attempt to
make it work as \enquote{correctly} as I can.  The switches for it are
basically the same as for the entry-only option, that is, assuming an
entry presents a reprinted edition of a work by Smith, first published
in 1926 (the \textsf{origdate}) and reprinted in 1985 (the
\textsf{date}):

\begin{enumerate}
\item \texttt{cmsdate=off} is the default: (Smith 1985).
\item \texttt{cmsdate=both} prints both the \textsf{origdate} and the
  \textsf{date}, using the \emph{Manual's}\ standard format: (Author
  [1926] 1985) in parenthetical citations, Author (1926) 1985 outside
  parentheses, e.g., in the reference list.
\item \texttt{cmsdate=on} prints the \textsf{origdate} at the head of
  the entry in the list of references and in citations: (Author 1926).
  NB: The \emph{Manual} no longer includes this among the approved
  options.  If you want to present the \textsf{origdate} at the head
  of an entry, then generally speaking you should probably use
  \texttt{cmsdate=both}.  I have nevertheless retained this option for
  certain cases where it has proved useful.  The old options
  \texttt{new} and \texttt{old} work like \texttt{both}.
\end{enumerate}

The important information for the user is that, when you set this
option in your preamble to \texttt{on} or \texttt{both} (or to the old
synonyms for the latter, \texttt{new} or \texttt{old}), then
\textsf{biblatex-chicago-authordate} (and \textsf{authordate-trad})
will change the default \cmd{DeclareLabeldate} definition so that the
\textsf{labelyear} search order will be \textsf{origdate, date,
  eventdate, urldate}.  This means that for entry types not covered by
the \texttt{avdate} option, and for those types as well if you turn
off that option, the \textsf{labelyear} will, in any entry containing
an \textsf{origdate}, be that very date.  If you want \emph{every}
such entry to present its \textsf{origdate} in citations and at the
head of reference list entries, then setting the option this way makes
sense, as you should automatically get the proper \textsf{extradate}
letter (1926\textbf{a}) and the correct sorting, without having to use
the counter-intuitive .bib file date switching that sometimes
accompanied the entry-only \texttt{cmsdate} option.  A few
clarifications may yet be in order.

\mylittlespace Obviously, any entry with only a \textsf{date} should
behave as usual.  Also, since \textsf{patent} entries have fairly
specialized needs, I have exempted them from this change to
\cmd{DeclareLabeldate}.  Third, the per-entry \texttt{cmsdate} options
will still affect which dates are printed in citations and at the head
of reference list entries, but they cannot change the search order for
the \textsf{labeldate}.  This will be fixed by the preamble option.
Fourth, if you have been used to switching the \textsf{date} and the
\textsf{origdate} to get the correct results, then you should be aware
that this mechanism may actually still be useful when using the
\texttt{on} switch to \texttt{cmsdate} in the preamble, but it
produces incorrect results when the \texttt{cmsdate} option is
\texttt{both} in the preamble and the individual entry.  The preamble
option is designed to make the need for this switching as rare as
possible, so some editing of existing databases may be necessary.
Fifth, the entry-only option \texttt{full} has no effect at all when
used in the preamble; you must set it in individual entries.  Finally,
please see the documentation of the \textbf{date} field in
section~\ref{sec:fields:authdate} for the fullest discussion of date
presentation in the \textsf{authordate} styles.

\mylittlespace Although \mymarginpar{\texttt{cmsorigdate}} I can't
currently think of any reason why anyone would want to use it on its
own, I should nonetheless mention that the \texttt{cmsorigdate} option
in your preamble will change the default \cmd{DeclareLabeldate}
settings to \textsf{origdate, date, eventdate, urldate}.  Setting
\texttt{cmsdate} to \texttt{on} or \texttt{both} in the preamble ---
see the previous option --- sets this to true, but if for some reason
you want to set it to true without any of the other effects of the
\texttt{cmsdate} option, then you can.  The effects may surprise.

\mylittlespace When \mymarginpar{\texttt{compresspages}} set to
\texttt{true}, any page ranges in your .bib file or in the
\textsf{postnote} field of your citation commands will be compressed
in accordance with the \emph{Manual's} specifications (9.61).
Something like 321-{-}328 in your .bib file would become 321--28 in
your document.  See the \textsf{pages} field in
section~\ref{sec:fields:authdate}, above.  Please note that the code
for this is in \textsf{biblatex-chicago.sty}, so if you load the
package with a call to \textsf{biblatex} instead then you'll get the
default \textsf{biblatex} compression style.

\mylittlespace This \mymarginpar{\texttt{datamodel}} is the standard
\textsf{biblatex} option for loading the named data model file
(excluding its\ .dbx extension).  After a request by Philipp Immel,
you can now set this option when you load the Chicago styles with
\verb+\usepackage{biblatex-chicago}+, and it will be passed through
properly to \textsf{biblatex} itself.  Cf.\ \textsf{biblatex.pdf}
\S~4.5.4.

\mylittlespace This \mymarginpar{\texttt{decaderange}} option works
just like \texttt{alwaysrange}, above, but only affects decade
presentation, not century.  Cf.\ table~\ref{ad:date:extras}.

\mylittlespace This \mymarginpar{\texttt{entrybreak}} option only
makes sense when used in conjunction with the \texttt{annotenp} switch
to the \texttt{formatbib} option (below).  The latter allows \LaTeX\
to break a page inside an \textsf{annotation} field printed without
starting a new paragraph, and it does so by allowing such a break in
the entry only after a set number of lines, by default set to 3.  The
idea is that most reference list entries will fit within 3 lines, so
the break would generally be somewhere inside the \textsf{annotation}.
If your document needs a value different from 3, provide the integer
using the \texttt{entrybreak} option in your preamble.  Some
experimentation may be needed to find the optimum number for a given
document.

\mylittlespace Although \mymarginpar{\texttt{footmarkoff}} the
\emph{Manual} (14.19) recommends specific formatting for footnote (and
endnote) marks, i.e., superscript in the text and in-line in foot- or
endnotes, Charles Schaum has brought it to my attention that not all
publishers follow this practice, even when requiring Chicago style.  I
have retained this formatting as the default setup, but if you include
the \texttt{footmarkoff} option, \textsf{biblatex-chicago} will not
alter \LaTeX 's (or the \textsf{endnote} package's) defaults in any
way, leaving you free to follow the specifications of your publisher.
I have placed all of this code in \textsf{biblatex-chicago.sty}, so if
you load the package with a call to \textsf{biblatex} instead, then
once again footnote marks will revert to the \LaTeX\ default, but of
course you also lose a fair amount of other formatting, as well.  See
section~\ref{sec:loading:auth}, below.

\mylittlespace The \mymarginpar{\texttt{formatbib}} \emph{Manual} in
fact says very little about formatting issues in reference lists,
e.g., whether to break entries across pages and whether to allow
widows and orphans (single lines at the start or end of a page).  A
quick and non-scientific survey of publications issued by the
University of Chicago Press suggests that actual practices are
extremely varied, so I've tried to provide a number of choices for
users of \textsf{biblatex-chicago}, most of them available as keys to
the \texttt{formatbib} preamble option, but a few of them also
involving settings of the \texttt{bibannotesep} and the
\texttt{entrybreak} options.  The keys of \texttt{formatbib} are as
follows:

\begin{description}
\setlength{\parskip}{-4pt}
\item[\qquad max:] This is the \textsf{biblatex} default, so if you
  don't set \texttt{formatbib} at all it's what you get.  It provides
  maximal intervention, disallowing entries broken across pages,
  including even when an entry includes a lengthy annotation.
\item[\qquad min:] This allows page breaks just about anywhere,
  including inside entries, and it also allows widows and orphans, so
  it will usually provide the most efficient use of available space on
  the pages of your bibliography.
\item[\qquad minwo:] This is like \texttt{min}, but discourages
  widows and orphans.
\item[\qquad annote:] This option treats annotations separately from
  reference list entries, allowing them to be broken across pages
  while the entry itself won't be.  It is intended for use with, and
  will only properly work with, \texttt{bibannotesep} set to one of
  the modes that start a new paragraph for the annotation, to wit,
  \texttt{par}, \texttt{vpar} (the default), \texttt{parbreak}, or
  \texttt{vparbreak}.  See below for the meaning of the
  \enquote{\texttt{break}} options here.
\item[\qquad annotenp:] This option attempts to treat annotations
  separately from reference list entries in those settings of
  \texttt{bibannotesep} which don't involve starting a new paragraph.
  It works by setting the number of lines in an entry after which page
  breaking is allowed.  By default entries will only break after 3
  lines, the idea being that most reference list entries fit into
  three lines, so at that point you're likely to be inside the
  annotation, but you can set the \texttt{entrybreak} option to any
  integer that works for your reference apparatus.
\end{description}

Please note that there is one possible break point that isn't directly
addressed by these options, that is, the one between the main entry
and the annotation when that annotation starts a new paragraph.  If
you set \texttt{bibannotesep} to \texttt{par} or \texttt{vpar}, then
\LaTeX\ will try very hard not to break between entry and annotation,
ensuring that the annotation at least starts on the same page as its
entry.  If you use \texttt{parbreak} or \texttt{vparbreak}, \LaTeX\ is
positively encouraged to break a page there, as is usual between
paragraphs.

\mylittlespace You can of course ignore the \texttt{formatbib} option
and provide your own settings.  \textsf{Biblatex} uses the
\cmd{bibsetup} command which you can renew in your preamble.  You can
find a nice commentary on the default values set by the package in the
file \textsf{biblatex.def}, which you'll find in the main
\textsf{biblatex} directory of your \TeX\ distribution.

\mylittlespace This \mymarginpar{\texttt{genallnames}} option affects
the choice of which names to present in the genitive case when using
the \cmd{gentextcites} command.  Please see the documentation of that
command in section~\ref{sec:cite:authordate}, above.

\mylittlespace Several \mymarginpar{\texttt{headline}\\\texttt{(trad
    only)}} users requested an option that turned off the automatic
transformations that produce sentence-style capitalization in the
title fields of the 15th-edition author-date style.  I have,
therefore, also included it in \textsf{authordate-trad}.  If you set
this option, the word case in your title fields will not be changed in
any way, that is, this doesn't automatically transform your titles
into headline-style, but rather allows the .bib file to determine
capitalization.  It works by redefining the command
\cmd{MakeSentenceCase}, so in the unlikely event you are using the
latter anywhere in your document please be aware that it will also be
turned off there.

\mylittlespace I \colmarginpar{\texttt{hyperall}} have received
requests to allow more than one part of a citation to act as a
hyperlink to the entry in the reference list (assuming you're using
the \textsf{hyperref} package, that is).  Setting this option for the
whole document, for specific entry types, or for individual entries
will make any names, titles, or dates present in the citation act as
such a link.  It acts, in other words, like a combination of
\mycolor{\texttt{hypername}} and \texttt{hypertitle}, below.

\mylittlespace This \colmarginpar{\texttt{hypername}} option ensures
that the name component of a citation links to the entry in the
reference list, again assuming that you're using \textsf{hyperref}.
You can set it for the whole document, for specific entry types, or
for individual entries.  Should a citation appear that contains
\emph{only} a name, then it will always be such a link, no matter what
the state of this option.

\mylittlespace When \mymarginpar{\texttt{hypertitle}} you use the
\textsf{hyperref} package with the author-date styles, the in-text
citations will provide a hyperlink to the full information in the list
of references.  Timo Thoms rightly pointed out that, generally, one
only wants one piece of the citation to provide the hyperlink, usually
the \textsf{date} part (though see the two previous options).  The
author-date styles will instead link the \textsf{title} or the
\textsf{shorthand} if there isn't a \textsf{date}, but if you set this
option to \texttt{true} globally in your preamble then all
\textsf{titles} and \textsf{shorthands} will link, regardless of
whether a \textsf{date} is also present.  You can also set
\texttt{hypertitle} by entry type or in the \textsf{options} field of
individual entries, allowing you to provide a hyperlink in cases where
the automatic mechanism gets it wrong (ency:britannica).

\mylittlespace The \mymarginpar{\texttt{juniorcomma}} \emph{Manual}
(6.43) states that \enquote{commas are not required around \emph{Jr.}\
  and \emph{Sr.},} so by default \textsf{biblatex-chicago} has
followed standard \textsf{biblatex} in using a simple space in names
like \enquote{John Doe Jr.}  Charles Schaum has pointed out that
traditional \textsc{Bib}\TeX\ practice was to include the comma, and
since the \emph{Manual} has no objections to this, I have provided an
option which allows you to turn this behavior back on, either for the
whole document or on a per-entry basis.  Please note, first, that
numerical suffixes (John Doe III) never take the comma.  The code
tests for this situation, and detects cardinal numbers well, but if
you are using ordinals you may need to set this to \texttt{false} in
the \textsf{options} field of some entries.  Second, I have fixed a
bug in older releases which always printed the \enquote{Jr.}\ part of
the name immediately after the surname, even when the surname came
before the given names (as in a reference list).  The package now
correctly puts the \enquote{Jr.}\ part at the end, after the given
names, and in this position it always takes a comma, the presence of
which is unaffected by this option.

\mylittlespace This \mymarginpar{\texttt{nameaddon}} option determines
where and when the \textsf{nameaddon} field will be printed.  There
are three possible values, available globally, per type, and per
entry:
\begin{description}
\setlength{\parskip}{-4pt}
\item[\qquad \texttt{all}:] This is the default; if an entry has a
  \textsf{nameaddon}, it will appear in the reference list.
\item[\qquad \texttt{none}:] The field will not appear in the
  reference list.
\item[\qquad \texttt{first}:] Philipp Immel requested this as a way to
  provide an \textsf{author's} dates in the \textsf{nameaddon} field
  and only have them printed the first time that author appears in the
  reference list.  The code tests for identical \textsf{nameaddon}
  fields in works by identical \textsf{authors}, so other sorts of
  \textsf{nameaddon} will be printed as usual.
\end{description}

Cf.\ \texttt{nameaddonformat} just below, and \texttt{nameaddonsep} in
section~\ref{sec:authpreset}.

\mylittlespace This \mymarginpar{\texttt{nameaddon-\\format}} option,
available globally, per type, and per entry, allows you to change the
format of the \textsf{nameaddon} field on the fly, so its value should
be a field format that \textsf{biblatex} understands.  This includes
standard formats like \texttt{parens,\,brackets} or \texttt{emph}, and
also custom formats that you provide in your preamble using
\cmd{DeclareFieldFormat}, in case the standard ones aren't adequate.
If you don't define this option, then the usual defaults apply, that
is, no formatting in \textsf{online, review,} and
\textsf{suppperiodical} entries, as well as in \textsf{misc} entries
with an \textsf{entrysubtype}, while square brackets surround the
field in all other entry types with the exception of \textsf{customc},
which has its own rules and ignores this option. Cf.\
\texttt{nameaddonsep} in section~\ref{sec:authpreset}.

\mylittlespace These \mymarginpar{\texttt{natbib}\\\texttt{mcite}} are
the standard \textsf{biblatex} options, which formerly required
slightly special handling when you loaded the Chicago style with
\verb+\usepackage{biblatex-chicago}+.  Both the forms \texttt{natbib}
and \texttt{natbib=true} (\texttt{mcite} \&\ \texttt{mcite=true})
should now work.

\mylittlespace When \mymarginpar{\texttt{nodatebrackets}\\
  \texttt{noyearbrackets}} you use \textsf{biblatex's} enhanced date
specifications to present a \enquote{circa} date (\verb+{1956~}+), an
uncertain date (\verb+{1956?}+), or one that is both at the same time
(\verb+{1956%}+), the date that by default will appear in your
documents will have square brackets around it.  This accords with the
\emph{Manual's} instructions concerning such dates (15.44), but may in
some circumstances prove syntactically awkward, or may perhaps be out
of step with a specific house style.  These two options, which may
appear in the preamble either for the whole document or for specific
entry types, or in individual entries, allow you to control when these
brackets will appear.  The first controls, mainly, dates that appear
in the body of an entry in the reference list, while the second
controls dates in citations and at the head of entries in the
reference list.  Cf.\ table~\ref{ad:date:extras}.

\mylittlespace At \colmarginpar{\texttt{noibid}} the request of an
early tester, I have included this option to allow you to turn off the
\texttt{ibidem} mechanism that \textsf{biblatex-chicago-authordate}
uses by default.  This mechanism doesn't actually print
\enquote{Ibid,} but rather includes only the \textsf{postnote}
information in a citation, i.e., it will print (224) instead of
(Author 2000, 224).  Setting this option will mean that these
shortened citations no longer appear automatically.  The option is
settable globally, per type, or per entry, so that fine-grained
control of individual citations is now possible without the use of the
\cmd{citereset} command and \textsf{biblatex}\ \texttt{citereset}
option, though these are of course still available.  See
\textsf{biblatex.pdf} \S~3.1.2.1.

\mylittlespace When \colmarginpar{\texttt{notitle}} citing sources
from antiquity (using the \texttt{classical} \textsf{entrysubtype}),
the \emph{Manual} (14.244--5) recommends using just the
\textsf{author} in short citations if only one \textsf{title} by that
author has come down to us, making the identification of the work
unambiguous.  I previously suggested using a command like
\cmd{citeauthor} to achieve this, but Tobias Becht suggested that a
less clumsy method would be better, so I've provided the
\mycolor{\texttt{notitle}} option, settable by entry type and also in
the \textsf{options} field of specific entries.  In effect, this turns
an author-title citation into an author-only one, but it won't affect
author-date citations.  Cf.\ herodotus:wilson.

\mylittlespace This \mymarginpar{\texttt{ordinalgb}} option, which
only affects users of the \texttt{british} language, restores the
previous package defaults, printing the \textsf{day} part of a
\textsf{date} specification as an ordinal number: 26th March 2017.
The new package default prints 26 March 2017, which is more in keeping
both with standard British usage and with the recommendations of the
\emph{Manual} (9.35).  The option is available only in the preamble.

\mylittlespace Originally
\mymarginpar{\texttt{postnotepunct}\\(experimental)} designed for the
notes \&\ bibliography style, this option may in fact be more useful
in the \textsf{authordate} styles.  If set to \texttt{true}, it allows
you to alter the punctuation that appears just before the
\textsf{postnote} argument of citation commands, simplifying in
particular the provision of comments within parenthetical citations.
In previous releases, you either needed to include the comment after a
page number, e.g., \cmd{autocite[16; some comment]\{citekey\}}, or
provide a separate .bib entry using the \textsf{customc} entry type,
e.g., \cmd{autocites\{chicago:manual\}\{chicago:comment\}}.  With this
option enabled, \cmd{autocite[; some comment]\{citekey\}} will do.
More generally, the \texttt{postnotepunct} option allows you to start
the \textsf{postnote} field with a punctuation mark (.\,,\,;\,:) and
have it appear as the \cmd{postnotedelim} in place of whatever the
package might otherwise automatically have chosen.  Please note that
this functionality relies on a very nifty macro by Philipp Lehman
which I haven't extensively tested, so I'm labeling this option
\enquote{experimental.}  Note also that the option only affects the
\textsf{postnote} field of citation commands, not the \textsf{pages}
field in your .bib file.

% %\enlargethispage{\baselineskip}

\mylittlespace This \mymarginpar{\texttt{seriesabbrev}} option
controls the printing, in the reference list, of the
\textsf{shortseries} field in place of the \textsf{series} field in
book-like entries.  It is \texttt{false} by default, so as shipped
\textsf{biblatex-chicago-authordate} will silently ignore such fields,
but you can set it to \texttt{true} either in the preamble for the
whole document or for specific entry types, or in individual entries,
and it will appear in the reference list.  For more details, see the
documentation of \textsf{shortseries} in
section~\ref{sec:fields:authdate}, above.

\mylittlespace Kenneth Pearce \mymarginpar{\texttt{shorthandfull}} has
suggested that, in some fields of study, a list of shorthands
providing full bibliographical information may replace the list of
references itself.  This option, which must be used in tandem with
\texttt{cmslos=false}, prints this full information in the list of
shorthands, though of course you should remember that any .bib entry
not containing a \textsf{shorthand} field won't appear in such a list.
Please see the documentation of the \textsf{shorthand} field in
section~\ref{sec:fields:authdate} above for information on further
options available to you for presenting and formatting the list of
shorthands.

\mylittlespace This \mymarginpar{\texttt{strict}} still-experimental
option attempts to follow the \emph{Manual}'s recommendations (14.36)
for formatting footnotes on the page, using no rule between them and
the main text unless there is a run-on note, in which case a short
rule intervenes to emphasize this continuation.  I haven't tested this
code very thoroughly, and it's possible that frequent use of floats
might interfere with it.  Let me know if it causes problems.

\mylittlespace Stefan \mymarginpar{\texttt{xrefurl}} Bj\"{o}rk pointed
out that when, using the \texttt{longcrossref} or
\texttt{booklongxref} options, you turn on the automatic abbreviation
of multiple entries in the same (e.g.) \textsf{collection} or
\textsf{mvcollection}, you could entirely lose a \textsf{url} that
might be helpful for locating a source, as the abbreviated forms in
the reference list wouldn't include this information.  Setting this
option to \textsf{true} either in the preamble or in individual
entries will allow the \textsf{url}, \textsf{doi}, or \textsf{eprint}
field to appear even in these abbreviated references.

\subsubsection{Style Options -- Entry}
\label{sec:authentryopts}

These options are settable on a per-entry basis in the
\textsf{options} field.

\mylittlespace Recent \mymarginpar{\texttt{cmsdate}} editions of the
\emph{Manual} have simplified the options for entries with more than
one date (15.40).  You can choose among them using the
\texttt{cmsdate} entry option.  It has 3 possible states relevant to
this problem, alongside a fourth which I discuss below.  An example
should make this clearer.  Let us assume that an entry presents a
reprinted edition of a work by Smith, first published in 1926 (the
\textsf{origdate}) and reprinted in 1985 (the \textsf{date}):

\begin{description}
\item[\qquad off:] This is the default.  The citation will look like
  (Smith 1985).
\item[\qquad both:] The citation will look like (Smith [1926] 1985).
\item[\qquad on:] The citation will look like (Smith 1926).  NB: The
  \emph{Manual} no longer includes this among the approved options.
  If you want to present the \textsf{origdate} at the head of an
  entry, then generally speaking you should probably use
  \texttt{cmsdate=both}.  I have retained the option because in some
  cases it is still useful.  The old options \texttt{new} and
  \texttt{old} work like \texttt{both}.
\end{description}

As I explained in detail above in section~\ref{sec:fields:authdate},
s.v.\ \enquote{\textbf{date},}\ because \textsf{biblatex's} sorting
algorithms and automatic creation of the \textsf{extradate} field
refer by default to the \textsf{date} before the \textsf{origdate}
when both are present, there may be situations when you need to have
the \emph{earlier} year in the \textsf{date} field, and the later one
in \textsf{origdate}, e.g., if you have another reprinted work by the
same author originally printed in the same year.
\textsf{Biblatex-chicago-authordate} will automatically detect this
switch, and given the same reprinted work as above, the results will
be as follows:

\begin{description}
\item[\qquad off:] This is the default.  The citation will look like
  (Smith 1926a).  This style is no longer recommended by the
  \emph{Manual}.
\item[\qquad both:] The citation will look like (Smith [1926a] 1985).
  The old options \texttt{old} and \texttt{new} are synonyms for this.
\item[\qquad on:] The citation will look like (Smith 1926a).  As noted
  above, this style is no longer recommended by the \emph{Manual}.
\end{description}

If, \mymarginpar{\texttt{switchdates}} for any reason, simply
switching the \textsf{date} and the \textsf{origdate} isn't possible
in a given entry, then you can put \texttt{switchdates} in the
\textsf{options} field to achieve the same result.  Also, you can use
the preamble version \mymarginpar{\texttt{cmsdate}\\\emph{in
    preamble}} of \texttt{cmsdate} to change the default order of
\cmd{DeclareLabeldate}, generally making this date-switching in your
.bib file unnecessary.  Please take a look at the full documentation
of the \textbf{date} field to which I referred just above, at the
preamble \texttt{cmsdate} documentation in
section~\ref{sec:authuseropts}, and also at
\href{file:cms-dates-sample.pdf}{\textsf{cms-dates-sample.pdf}} and
\textsf{dates-test.bib} for examples of how all this works.

\mylittlespace Bertold Schweitzer has brought to my attention certain
difficult corner cases involving cross-referenced works with more than
one date.  In order to facilitate the accurate presentation of such
sources, I made a slight change to the way \texttt{cmsdate=on}
and \texttt{cmsdate=both} work.  If, and only if, a work has only one
date, and there is no \texttt{switchdates} in the \textsf{options}
field, then \texttt{cmsdate=on} and \texttt{cmsdate=both} will both
result in the suppression of the \textsf{extradate} field in that
entry.  Obviously, if the same options are set in the preamble, this
behavior is turned off, so that single-date entries will still work
properly without manual intervention.

\mylittlespace Recent editions of the \emph{Manual} specify that it is
\enquote{usually sufficient to cite newspaper and magazine articles
  entirely within the text} (15.49).  This will apply mainly to
\textsf{article} and \textsf{review} entries with
\textsf{entrysubtype} \texttt{magazine}, and involves a parenthetical
citation giving the \textsf{journaltitle} and then the full
\textsf{date}, not just the year, with any other relevant identifying
information incorporated into running text.  (Cf.\ 14.198.)\ In order
to facilitate this, I have added a further switch to the
\texttt{cmsdate} option \mymarginpar{\texttt{cmsdate=full}} ---
\texttt{full} --- which \emph{only} affects the presentation of
citations, and causes the printing of the full date specification
there.  You can use the standard \textsf{biblatex} \texttt{skipbib}
option to keep such entries from appearing in the list of references,
and you may, if your .bib entry is a complete one, also need
\texttt{useauthor=false} in order to ensure that the
\textsf{journaltitle} appears in the citations rather than the
\textsf{author}.

\subsection{General Usage Hints}
\label{sec:hints:auth}

\subsubsection{Loading the Styles}
\label{sec:loading:auth}

With the addition of the \textsf{authordate-trad} style to the
package, there are now three keys for choosing which style to load,
\texttt{notes}, \texttt{authordate}, and \textsf{authordate-trad}, one
of which you put in the options to the \cmd{usepackage} command.  With
early versions of \textsf{biblatex-chicago}, the standard way of
loading the package was via a call to \textsf{biblatex}, e.g.:
\begin{verbatim}
  \usepackage[style=chicago-authordate,strict,backend=biber,%
    babel=other,bibencoding=inputenc]{biblatex}
\end{verbatim}
Now, the default way to load the style, and one that will in the
vast majority of standard cases produce the same results as the old
invocation, will look like this:
\begin{verbatim}
  \usepackage[authordate,strict,backend=biber,autolang=other,%
    bibencoding=inputenc]{biblatex-chicago}
\end{verbatim}

If you read through \textsf{biblatex-chicago.sty}, you'll see that it
sets a number of \textsf{biblatex} options aimed at following the
Chicago specification, as well as setting a few formatting variables
intended as reasonable defaults (see section~\ref{sec:preset:authdate},
above).  Some parts of this specification, however, are plainly more
\enquote{suggested} than \enquote{required,} and indeed many
publishers, while adopting the main skeleton of the Chicago style in
citations, nonetheless maintain their own house styles to which the
defaults I have provided do not conform.

\mylittlespace If you only need to change one or two parameters, this
can easily be done by putting different options in the call to
\textsf{biblatex-chicago} or redefining other formatting variables in
the preamble, thereby overriding the package defaults.  If, however,
you wish more substantially to alter the output of the package,
perhaps to use it as a base for constructing another style altogether,
then you may want to revert to the old style of invocation above.
You'll lose all the definitions in \textsf{biblatex-chicago.sty},
including those to which I've already alluded and also the code that
sets the note number in-line rather than superscript in endnotes or
footnotes, the URL line-breaking code, and the Chicago-specific
number- and date-range compression code.  You'll need to load the
required packages \textsf{xstring} and \textsf{nameref} yourself, as
\textsf{biblatex} doesn't do it for you.  Also in this file is the
code that calls all of the package's localization files, which means
that you'll lose all the Chicago-specific bibstrings I've defined
unless you provide, in your preamble, a \cmd{DeclareLanguageMapping}
command, or several, adapted for your setup, on which see
section~\ref{sec:international} below and also \S\S~4.9.1 and 4.11.8
in \textsf{biblatex.pdf}.

\mylittlespace What you \emph{will not} lose is the ability to call
the package options \texttt{annotation, strict, cmslos=false} and
\texttt{noibid} (section~\ref{sec:authuseropts}, above), in case these
continue to be useful to you when constructing your own modifications.
There's very little code, therefore, actually in
\textsf{biblatex-chicago.sty}, but I hope that even this minimal
separation will make the package somewhat more adaptable.  Any
suggestions on this score are, of course, welcome.

\subsubsection{Other Hints}
\label{sec:otherhints:auth}

Starting with \textsf{biblatex} version 1.5, in order to adhere to the
author-date specification you will need to use \textsf{Biber} to
process your .bib files, as \textsc{Bib}\TeX\ (and its more recent
variants) will no longer provide all the required features.  This
document assumes that you are using \textsf{Biber}; if you wish to
continue using \textsc{Bib}\TeX\ then you need \textsf{biblatex}
version 1.4c and \textsf{biblatex-chicago} 0.9.7a.

\mylittlespace If your .bib file contains a large number of entries
with more than three authors, then you may run into some limitations
of the \textsf{biblatex-chicago} code.  The default settings are
\texttt{maxnames=3,minnames=1} in citations and
\texttt{maxbibnames=10,minbibnames=7} in the list of references.  In
practice, this means that an entry like hlatky:hrt, with 5 authors,
will present all of them in the list of references but will truncate
to one in citations, like so: (Hlatky et al. 2002).  For the vast
majority of circumstances, these settings are exactly right for the
Chicago author-date specification.  However, if \enquote{a reference
  list includes another work of the same date that would also be
  abbreviated as [\enquote{Hlatky et al.}] but whose coauthors are
  different persons or listed in a different order, the text citations
  must distinguish between them} (15.29).  The (\textsf{Biber}-only)
\textsf{biblatex} option \texttt{uniquelist}, set for you in
\textsf{biblatex-chicago.sty}, will automatically handle many of these
situations for you, but it is as well to understand that it does so by
temporarily suspending the limits, listed above, on how many names to
print in a citation.  Without \texttt{uniquelist}, \textsf{biblatex}
would present such a work as, e.g., (Hlatky et al. 2002b), while
hlatky:hrt would be (Hlatky et al. 2002a).  This does distinguish
between them, but inaccurately, as it suggests that the two different
author lists are exactly the same.  With \texttt{uniquelist}, the two
citations might look like (Hlatky, Boothroyd et al.\ 2002) and
(Hlatky, Smith et al.\ 2002), which is what the specification
requires.

\mylittlespace If, however, the distinguishing name occurs further
down the author list --- in fourth or fifth position in our examples
--- then the default settings would produce citations with all 4 or 5
names printed, which can become awkwardly long.  In such a situation,
you can provide \textsf{shortauthor} fields that look like this:
\{\{Hlatky et al., \textbackslash mkbibquote\{Quality of Life,\}\}\}
and \{\{Hlatky et al., \textbackslash mkbibquote\{Depressive
Symptoms,\}\}\}, using a shortened title to distinguish the
references.  This would produce (Hlatky et al., \enquote{Quality of
  Life,} 2002) and (Hlatky et al., \enquote{Depressive Symptoms,}
2002), again as the spec requires.  There is, unfortunately, no
simpler way that I know of to deal with this situation.

\mylittlespace One useful rule, when you are having difficulty
creating a .bib entry, is to ask yourself whether all the information
you are providing is strictly necessary.  The Chicago specification is
a very full one, but the \emph{Manual} is actually, in many
circumstances, fairly relaxed about how much of the data from a work's
title page you need to fit into a reference.  Authors of introductions
and afterwords, multiple publishers in different countries, the real
names of authors more commonly known under pseudonyms, all of these
are candidates for exclusion if you aren't making specific reference
to them, and if you judge that their inclusion won't be of particular
interest to your readers.  Of course, any data that may be of such
interest, and especially any needed to identify and track down a
reference, has to be present, but sometimes it pays to step back and
reevaluate how much information you're providing.  I've tried to make
\textsf{biblatex-chicago} robust enough to handle the most complex,
data-rich citations, but there may be instances where you can save
yourself some typing by keeping it simple.

\mylittlespace Scot Becker has pointed out to me that the inverse
problem not only exists but may well become increasingly common, to
wit, .bib database entries generated by bibliographic managers which
helpfully provide as much information as is available, including
fields that users may well wish not to have printed (ISBN, URL, DOI,
\textsf{pagetotal}, inter alia).  The standard \textsf{biblatex}
styles contain a series of options, detailed in \textsf{biblatex.pdf}
\S~3.1.2.2, for controlling the printing of some of these fields, and
I have implemented others that are relevant to
\textsf{biblatex-chicago}, along with a couple that Scot requested and
that may be of more general usefulness.  There is also a general
option to excise with one command all the fields under consideration
-- please see section~\ref{sec:authpreset} above.

\section{The \textsf{Jurisdiction},
  \textsf{Legislation}, and \textsf{Legal} Entry
  Types}
\label{sec:legal}

I have received numerous requests over the years to include some means
of referring to legal and public documents which, broadly speaking,
don't fit easily into any of the standard \textsf{biblatex} entry
types.  The \emph{Manual} (14.269--305) recommends using the
\emph{Bluebook} as a guide for formatting such references, while also
suggesting certain modifications to this formatting to bring it more
into line with Chicago's usual practices.  \textsf{Biblatex-chicago}
now offers three entry types --- \textbf{jurisdiction},
\textbf{legal}, and \textbf{legislation} --- which allow you to
present at least a substantial subset of what the \emph{Bluebook}
offers.  As the rules for your .bib entries are the same in the notes
\&\ bibliography style and in the author-date styles, and as these
rules mainly come from a source outside the \emph{Manual,} and
additionally as these rules apparently require even the author-date
styles to use a system of foot- or endnotes (15.58), I have documented
these types in a section of their own, applicable to all the Chicago
styles.  (Some few changes needed when using the author-date styles,
mainly to do with citation commands, will be outlined at the end.)
You can also consult the example files \textsf{legal-test.bib} and
\href{file:cms-legal-sample.pdf}{\textsf{cms-legal-sample.pdf}} to see
how you might construct your database entries.

\subsection{Types, Subtypes, and Fields}
\label{sec:legal:types}

Anyone who has used the \emph{Bluebook} will realize that it is
hopeless to attempt to fit its labyrinthine complexities into three
entry types, but with the addition of numerous \textsf{entrysubtypes}
and some parsing by \textsf{Biber} under the hood, I hope to have
covered the main sorts of material discussed by the \emph{Manual}.  As
a first approximation, all three types begin from a structure
analogous to the standard \textsf{biblatex article} type, with a
number of subtle differences that I have attempted to make consistent
across the three.  Standard practice is to present the references
\emph{only} in notes, and not in a bibliography, so by default
\textsf{biblatex-chicago} excludes these types from the latter, though
you can control this using an option (see below).

\mybigspace This \mymarginpar{\textbf{jurisdiction}} type is for
presenting legal cases and court decisions.  A typical entry will
contain the following fields:

{\renewcommand{\descriptionlabel}[1]{\quad\textsf{#1}:}
\begin{description}
\setlength{\parskip}{-2pt}
\item[title] The case name as seen in the first, long note.
\item[shorttitle] The case name for subsequent, short
  notes, ordinarily either the plaintiff or the nongovernmental party.
\item[journaltitle \textrm{and/or} shortjournal] The reporter for the
  case, \emph{always} presented in a standard abbreviated form
  available in the \emph{Bluebook}.  You can place the abbreviation
  either in the \textsf{journaltitle} or in the \textsf{shortjournal}
  field.  If you wish to present your readers with a list of
  abbreviations with their expansions, then the expansion goes in
  \textsf{journaltitle} and the abbreviation in \textsf{shortjournal}.
  (Cf.\
  \href{file:cms-legal-sample.pdf}{\textsf{cms-legal-sample.pdf}} to
  see how this might look.)
\item[pages \emph{or} issue] When using a standard official reporter,
  this will contain the opening page of the decision in that reporter,
  while any \textsf{postnote} field will contain the specific page on
  which a particular citation appears (a \enquote{pincite}).  When
  citing a commercial electronic database, on the other hand, you
  should give, instead of a \textsf{pages} field, the identifying
  number of the case using the \textsf{issue} field.
  \textsf{Biblatex-chicago} uses the presence of the \textsf{issue}
  field to provide the slightly different formatting required for
  citations from databases like Westlaw or LexisNexis.  When an
  \textsf{issue} field is present, then both the \textsf{pages} field
  and the \textsf{postnote} field can provide a pincite.  (Cf.\
  federal:case and database:case.)
\item[series] If you are citing an official reporter, then it may have
  a \textsf{series} number, which will be printed immediately after
  the name of the reporter.
\item[volume] The volume number of the reporter.  It will often be the
  same as the year when using a commercial electronic database, but
  you still need to provide it separately.
\item[number] The docket number of the case, generally required when
  the reporter is a commercial database.
\item[date] The date of the decision.
\item[location] The abbreviated name of the court, if it isn't clear
  from the reporter cited.  It will be associated with the
  \textsf{date} in American cases, but not in Canadian or UK cases
  (see below).  Being a list field, it can contain more than one item,
  in case you need a separate set of parentheses to identify a
  jurisdiction as well as a court name in Canadian or UK cases (cf.\
  uk:case:square).
\end{description}}

These are, so to speak, the basic elements of a \textsf{jurisdiction}
citation, which may, depending on specific circumstances, require
supplementation by the following:

{\renewcommand{\descriptionlabel}[1]{\quad\textsf{#1}:}
\begin{description}
\setlength{\parskip}{-2pt}    
\item[entrysubtype\{square\} \textrm{or} \{round\}] The \emph{Manual}
  includes examples for citing cases in the United Kingdom and in
  Canada, and the \texttt{square} \textsf{entrysubtype} identifies the
  reporter either as Canadian or as a UK reporter for which the year
  is essential to locating the case, that is, when \enquote{there is
    either no volume number or the volumes for each year are numbered
    anew, not cumulatively} (14.298).  The \texttt{round} subtype, by
  contrast, identifies a UK reporter where the volumes are numbered
  cumulatively, making the year inessential.  (The names refer to the
  shapes of the brackets placed around the year in each case.  Cf.\
  canada:case, uk:case:round, and uk:case:square.)
\item[origlocation \textrm{or} origpublisher] If you need to cite more
  than one reporter for a given case, then there are two
  possibilities.  Either the second (and subsequent) reporter(s)
  use(s) the same pagination as the first (\textsf{origpublisher}) or
  the reporters use different pagination (\textsf{origlocation}).
  Since both are list fields, you can in theory provide several
  reporters, but please note that these fields are currently only
  provided for citations of American cases.  (Cf.\ state:case:2reps.)
\item[related] It may be necessary sometimes to indicate further
  action by another, higher court, such as the US Supreme Court's
  grant or denial of \emph{certiorari}.  The usual \textsf{related}
  mechanism is useful in such situations, particularly with a tailored
  \textsf{relatedstring} field.  (Cf.\ federal:lower:related.)
\end{description}}

This \mymarginpar{\textbf{legislation}} is the most complicated of
the new entry types, with several \textsf{entrysubtypes} and a number
of tricky corners, particularly with regard to the provision of
subsequent short notes after the first full citation.  It is intended
to cope with constitutions and with legislative and executive
documents of all kinds, with the exception of treaties, for which you
can use the \textsf{legal} type, below.  In effect, the type tries to
cover federal, state, and municipal laws and ordinances, statutes,
bills, resolutions, reports, debates, hearings, presidential and
congressional documents, and constitutions, none of which it does with
particular elegance, so consider it a work in progress.  Many of the
fields have close analogues in the \textsf{jurisdiction} type, so at
least there is some bare minimum of consistency when dealing with
public and legal material.

{\renewcommand{\descriptionlabel}[1]{\quad\textsf{#1}:}
\begin{description}
\setlength{\parskip}{-2pt}
\item[author] Some kinds of material, usually reports, may have an
  \textsf{author}, often an organizational one.  (Cf.\ congress:report and
  uk:command.)
\item[title] Reports, bills, hearings and the like frequently have a
  \textsf{title} which, please note, quite frequently will not turn up
  in short notes, depending on which other fields are present.
\item[titleaddon] This field is considerably more important in
  \textsf{legislation} entries than the \textsf{shorttitle} field,
  mainly because it will turn up in many short notes where the
  \textsf{title} will not.  It will frequently contain specifying
  information on legislative material, and will therefore often allow
  short notes to differentiate citations of sources that might have
  the same \textsf{title} but differ in other respects.  (Cf.\
  congress:publiclaw, congress:bill, and congress:report.)
\item[number] Gives an identifying number to a \textsf{title} or a
  \textsf{titleaddon}, with \cmd{bibstring\{number\}} prefixed to it.
  It too can appear in short notes.
\item[note] This gives a section or other specifying information
  related to a \textsf{titleaddon} and \textsf{number}.  (Cf.\
  congress:publiclaw.)
\item[journaltitle \textrm{and/or} shortjournal] There is usually a
  standard place for publishing various sorts of legislative material,
  and as in \textsf{jurisdiction} entries it is \emph{always}
  presented in a standard abbreviated form available in the
  \emph{Bluebook}.  You can place the abbreviation either in the
  \textsf{journaltitle} or in the \textsf{shortjournal} field.  If you
  wish to present your readers with a list of abbreviations with their
  expansions, then the expansion goes in \textsf{journaltitle} and the
  abbreviation in \textsf{shortjournal}.  Like \textsf{titleaddon},
  this field will often appear in short notes.  (Cf.\
  \href{file:cms-legal-sample.pdf}{\textsf{cms-legal-sample.pdf}} to
  see how this might look.)
\item[volume] The volume number of the \textsf{shortjournal}.  It can
  be a cardinal or an ordinal, depending on the \textsf{shortjournal}.
  (Cf.\ congress:publiclaw and congress:bill.)
\item[series] In citations of American material, this will usually
  contain session information pertaining to a legislative publication.
  Elsewhere it will often contain just be a plain number, not unlike
  in \textsf{jurisdiction} entries.  (Cf.\ congress:debate:globe,
  state\hc statute:okla, and uk:hansard.)
\item[issue] This field can provide an identifying number in some
  circumstances, particularly when you don't want it prefixed by any
  bibstring --- cf.\ uk:command.
\item[pages] Somewhat similar to its use in \textsf{jurisdiction}
  entries, this will usually contain the opening page, or sometimes
  the section number, of the material in the \textsf{shortjournal},
  while any \textsf{postnote} field will contain the specific page on
  which a particular citation appears (a \enquote{pincite}).
\item[part \textrm{or} chapter] Some sources use \textsf{part} or
  \textsf{chapter} numbers instead of \textsf{pages} or sections.
  (Cf.\ canada:statute and uk:statute.)
\item[date] The date of publication of the material, usually just a
  \textsf{year}, though sometimes a full date, e.g., see
  executive:proclamation.
\item[location] If it is not clear from the \textsf{title} or the
  \textsf{shortjournal}, this field can specify, in abbreviated form,
  the US state where the legislative material originates.  It will be
  associated with the \textsf{date} in long notes, but will appear
  elsewhere in short notes.  (Cf.\ state:statute:okla.)
\item[usera] This specifies a particular edition, possibly from a
  commercial electronic database, of a legislative publication.  It
  will be associated with the \textsf{date} in long notes but won't
  appear in short ones.  (Cf.\ congress:debate:new and
  state:statute:ky.)%%\enlargethispage{\baselineskip}
\item[addendum] You can use this field to specify the speaker at
  hearings or in debates, the Canadian or British jurisdiction of some
  laws if not otherwise clear from the citation, or possibly simply
  additional information about a source.  (Cf.\ canada:statute,
  congress:debate:new, congress:hearing, state:statute:ky, and
  uk:statute.)
\item[entrysubtype] The sheer variety of sources included under the
  \textsf{legislation} type, and the specialized rules for presenting
  them, have necessitated the introduction of a substantial network of
  \textsf{entrysubtypes}:
  {\renewcommand{\descriptionlabel}[1]{\quad\textsf{#1}:}
  \begin{description}
    \setlength{\parskip}{-2pt}
  \item[canada] Identifies Canadian statutes (canada:statute).
  \item[constitution] For constitutions, be they federal, state, or
    local (constitution\hc arkansas and constitution:federal).
    \textsf{Biber} will automatically detect if the \textsf{title}
    contains the string \texttt{Const} and provide the
    \textsf{entrysubtype} for you, but in other cases you'll have to
    provide it yourself.
  \item[hansard] Identifies UK parliamentary debates as published in
    \emph{Hansard}.  (Cf.\ uk\hc hansard.)
  \item[hearing] For congressional hearings (congress:hearing).
  \item[uk] Identifies UK statutes and command papers.  (Cf.\
    uk:command, uk:statute, and uk:statute:regnal.)
  \item[un] For UN documents (un:resolution).
  \end{description}}
\end{description}}

A glance through the \textsf{legal-test.bib} file should help
enormously when you're trying to work out how to present a particular
source, and all suggestions for pruning the foliage will be welcome.

\mybigspace This \mymarginpar{\textbf{legal}} type is intended as a
catch-all for miscellaneous public documents not included in the
previous two types, but for the moment the only sort of material for
which it is required is international treaties (14.290; treaty).  The
usual fields for such material include:

{\renewcommand{\descriptionlabel}[1]{\quad\textsf{#1}:}
\begin{description}
\setlength{\parskip}{-2pt}
\item[title] The treaty name as seen in the first, long note.
\item[shorttitle] The treaty name for subsequent, short notes.  You
  can also use the \textsf{shorthand} field in such entries.
\item[titleaddon] This contains the names of the countries involved
  in the treaty, in abbreviated form.
\item[journaltitle \textrm{and/or} shortjournal] The standard
  publication containing the treaty, \emph{always} presented in an
  abbreviated form available in the \emph{Bluebook}.  You can place
  the abbreviation either in the \textsf{journaltitle} or in the
  \textsf{shortjournal} field.  If you wish to present your readers
  with a list of abbreviations with their expansions, then the
  expansion goes in \textsf{journaltitle} and the abbreviation in
  \textsf{shortjournal}.  (Cf.\
  \href{file:cms-legal-sample.pdf}{\textsf{cms-legal-sample.pdf}} to
  see how this might look.)
\item[issue] This will contain the identifying number provided by the
  \textsf{shortjournal}.
\item[volume] The volume number of the \textsf{shortjournal}.
\item[date] The exact date of signing, as the year of publication can
  differ from it.
\item[pages] This, or a \textsf{postnote} field, can contain a
  specific page reference (\enquote{pincite}).
\end{description}}

\subsection{Citation Commands}
\label{sec:legal:citcommands}

The \emph{Bluebook} style mandates footnotes without a bibliography,
so it should be simple to include such references in the Chicago notes
\&\ bibliography style, which uses foot- or endnotes as standard.  The
usual citation commands should work as you expect, though I wouldn't
recommend the \cmd{textcite} commands, as they will produce surprising
and unsatisfactory results.  For users of the author-date styles,
however, the \emph{Bluebook} more or less requires you to adopt a
separate set of notes in addition to the standard author-year
citations, which means that for these three legal entry types you'll
have to remember to use new citation commands that I've provided:
\cmd{fullcite}, \cmd{footfullcite}, and \cmd{parenfullcite}.  The
first prints the reference, the second does so in a footnote, and the
third does so inside parentheses.  For both styles,
\mymarginpar{\cmd{runcite}} there is a new command that you should use
if you are citing a \textsf{jurisdiction} entry in the running flow of
text rather than as a stand-alone citation, whether that text is in a
note or in the main body (14.276).  This requires a different
presentation of the \textsf{title} field, and using \cmd{runcite}
alerts \textsf{biblatex-chicago} to this circumstance.

\subsection{Options}
\label{sec:legal:options}

Several new options allow you to control the presentation of legal
notes in your document.  The default settings are indicated in the
margins.

\mylittlespace This \mymarginpar{\texttt{legalnotes=true}} option
prevents the printing of legal citations in a bibliography or
reference list, as the \emph{Bluebook} recommends.  You can change
this to \texttt{false} in the preamble of your document, but you
should be aware that the reference printed in the bibliography will be
a clone of the long-note form, as the \emph{Bluebook} doesn't provide
an alternative version.

\mylittlespace This \mymarginpar{\texttt{noneshort=\\false}} option
controls the availability of the short form of the note, intended for
use in subsequent citations of entries already presented in full
notes.  By default, \textsf{biblatex-chicago} attempts to provide
\emph{Bluebook}-authorized short versions of citations in this
situation, and some of the many \textsf{entrysubtypes} are at least
partially designed to cope with the complexities of the specification.
The \emph{Manual}, for its part, suggests that \enquote{works that
  cite only a few legal documents may be better off using the full
  form for each citation} (14.275).  In the author-date styles, you
can set this option to \texttt{true} globally, per type, or in
individual .bib entries to accomplish this.  Assuming you've only used
the \cmd{fullcite} commands for the \emph{Bluebook} entry types, the
option will only apply to such entries.  In the notes \&\ bibliography
style the global option would apply to all entry types, but you can
also set this to \texttt{true} on a type-by-type basis or in
individual entries.  It is designed mainly for use with
\emph{Bluebook} entries, but it might perhaps be useful elsewhere.
Please be aware that, even with this option on, the \emph{ibidem}
mechanism remains in operation for repeated citations, and also that
the option may give surprising results in the presence of
\textsf{shorthand} fields and/or the \texttt{shorthandfirst} and
\texttt{short} options.

\mylittlespace I \mymarginpar{\texttt{short=false}} have ported this
option, already present in the notes \&\ bibliography style, to the
author-date styles to allow users to present short notes from the very
first citation.  I'm not certain what the use case might be for this,
as it's intended for saving space in documents where short notes can
point to references in a full bibliography.  Still, if for any reason
you need this you can set the option to \texttt{true} globally, by
entry type, or for specific entries.

\mylittlespace This \mymarginpar{\texttt{supranotes=\\true}} is a
\emph{Bluebook}-specific option, and it produces, for some entry types
and subtypes, a back reference to the first, long note at the end of
subsequent, short citations.  It takes the form \enquote{\emph{supra}
  note \#,} and is available in all Chicago styles, though you'll only
see it in certain sorts of citation, automatically controlled by
\textsf{biblatex-chicago} in accordance with the \emph{Bluebook}
specification.  If you prefer not to see such back references, you can
set the option to \texttt{false} either in the preamble or in
individual .bib entries.

\section{Internationalization}
\label{sec:international}

Several users have requested that, in line with analogous provisions
in other \enquote{American} \textsf{biblatex} styles (e.g.,
\textsf{biblatex-apa} and \textsf{biblatex-mla}), I include facilities
for producing a Chicago-like style in other languages.  I have
supplied three lbx files, \textsf{cms-german.lbx}, its clone
\textsf{cms-ngerman.lbx}, and \textsf{cms-french.lbx}, in at least
partial fulfillment of this request.  For this release Mar\c{c}al
Orteu Punsola has provided \textsf{cms-spanish.lbx}, thereby adding to
the generous contributions of Patrick Danilevici
(\textsf{cms-romanian.lbx}), Wouter Lancee (\textsf{cms-dutch.lbx}),
Gustavo Barros (\textsf{cms-brazilian.lbx}), Stefan Bj\"{o}rk
(\textsf{cms-swedish.lbx}), Antti-Juhani Kaijahano
(\textsf{cms-finnish.lbx}), Baldur Kristinsson
(\textsf{cms-icelandic.lbx}), and H{\aa}kon Malmedal
(\textsf{cms-norsk.lbx}, \textsf{cms-nynorsk.lbx}, and
\textsf{cms-norwegian.lbx}).  I include \textsf{cms-british.lbx} in
order to simplify and to improve the package's handling of
non-American typographical conventions in English.  This means that
all --- or at least most --- of the Chicago-specific bibstrings are
now available for documents and reference apparatuses written in these
languages, with, as I intend, more languages to follow, limited mainly
by my finite time and even-more-finite competence.  (If you would like
to provide bibstrings for a language in which you want to work, or
indeed correct deficiencies in the lbx files contained in the package,
please contact me.)

\mylittlespace Using \mymarginpar{\textbf{babel}} these facilities is
fairly simple.  By default, and this functionality remains the same as
it was in the previous release of \textsf{biblatex-chicago}, calls to
\cmd{DeclareLanguageMapping} in \textsf{biblatex-chicago.sty} will
automatically load the American strings, and also \textsf{biblatex's}
American-style punctuation tracking, when you:
\begin{enumerate}
\item Load \textsf{babel} with \texttt{american} as the main text
  language.
\item Load \textsf{babel} with \texttt{english} as the main text
  language.
\item[] \qquad \emph{or}
\item Do not load \textsf{babel} at all.
\end{enumerate}
(This last is a change from the \textsf{biblatex} defaults --- cp.\
\S~3.12.2 in \textsf{biblatex.pdf} --- but it seems to me reasonable,
in an American citation style, to expect this arrangement to work well
for the majority of users.)

\mylittlespace If, for whatever reason, you wanted to use
\textsf{biblatex-chicago} but retain British typographical conventions
--- punctuation outside of quotation marks, outer quotes single rather
than double, etc.\ --- then you no longer need to follow the
complicated rules outlined in previous releases of
\textsf{biblatex-chicago}.  Instead, simply load \textsf{babel} with
the \texttt{british} option.

\mylittlespace If you want to use Brazilian Portuguese, Dutch,
Finnish, French, German, Icelandic, Norwegian, Romanian, Spanish, or
Swedish strings in the reference apparatus, then you can load
\textsf{babel} with \texttt{brazilian}, \texttt{dutch},
\texttt{finnish}, \texttt{french}, \texttt{german},
\texttt{icelandic}, \texttt{ngerman}, \texttt{norsk},
\texttt{nynorsk}, \texttt{romanian}, \texttt{spanish}, or
\texttt{swedish} as the main document language.  You no longer need
any calls to \cmd{DeclareLanguageMapping} in your document preamble,
since \textsf{biblatex-chicago.sty} automatically provides these if
you load the package in the standard way.

\mylittlespace You can also define which bibstrings to use on an
entry-by-entry basis by using the \textsf{hyphenation} field in your
bib file, but you will have to make sure that all the languages you
want to use in this way are included in the call to load
\textsf{babel} in your preamble, even if not as the main text
language.  The \cmd{DeclareLanguageMapping} calls in
\textsf{biblatex-chicago.sty} should do the rest, assuming you've
loaded the package that way.

\mylittlespace Three other hints may be in order here.  Please note,
first, that I haven't altered the standard punctuation procedures used
in any of the other available languages, so commas and full stops will
appear outside of quotation marks, and those quotation marks
themselves will be language-specific.  If, for whatever reason, you
wish to follow the Chicago specification and move punctuation inside
quotation marks, then you'll need a declaration of this sort in your
preamble:

\begin{verbatim}
  \DefineBibliographyExtras{german}{%
    \DeclareQuotePunctuation{.,}}
\end{verbatim}

Second, depending on the nature of your bibliography database, it will
only rarely be possible to process the same bib file in different
languages and obtain completely satisfactory results.  Fields like
\textsf{note} and \textsf{addendum} will often contain
language-specific information that won't be translated when you switch
languages, so manual intervention will be necessary.  If you suspect
you may have a need to use the same bib file in different languages,
you can minimize the amount of manual intervention required by using
the bibstrings defined either by \textsf{biblatex} or by
\textsf{biblatex-chicago}.  Here, a quick read through
\textsf{notes-test.bib} and/or \textsf{dates-test.bib} should give you
an idea of what is available for this purpose --- see esp.\ the
strings \texttt{by}, \texttt{nodate}, \texttt{newseries},
\texttt{number}, \texttt{numbers}, \texttt{oldseries},
\texttt{pseudonym}, \texttt{reviewof}, \texttt{revisededition}, and
\texttt{volume}, and also section \ref{sec:formatcommands} above,
esp.\ s.v.\ \enquote{\cmd{partedit}.}

\mylittlespace Finally, the French and German bibstrings I have
provided may well break with established bibliographical traditions in
those languages, but my main concern when choosing them was to remain
as close as possible to the quirks of the Chicago specification.  I
have relied on the judgment of the creators of the Brazilian, Dutch,
Finnish, Icelandic, Norwegian, Romanian, Spanish, and Swedish
localizations in those instances, as well as attempting to maintain
harmony where necessary with the standard \textsf{biblatex}
translations.  If you have strong objections to any of the strings, or
indeed to any of my formatting decisions, please let me know.

\subsection{The \mycolor{\texttt{cmsnameparts}} Option}
\label{sec:nameparts}

\textsf{Biblatex} provides a mechanism for localizing name
presentation in citations and in bibliographies, so for example making
it possible to present Russian patronymics or Chinese names (both in
the original script and in transliteration) in ways that the standard
macros can't provide, focusing as those macros do on the usual Western
European approaches.  The files \textsf{93-nameparts.tex} and
\textsf{93-nameparts.dbx} outline and demonstrate the changes
required, providing code for Russian, Ethiopian, and CJK names.  The
\emph{Manual} (11.90) too provides some examples, admitting the
necessity of such treatments in some contexts.  At the request of
Philipp Immel, I have provided the \mycolor{\texttt{cmsnameparts}}
option in all styles to facilitate access to this mechanism for
\textsf{biblatex-chicago} users.  In practice what I have done is the
following:

\begin{enumerate}
\item Include \textsf{93-nameparts.dbx} in the package under the name
  \mycolor{\textsf{cms.dbx}}.  This extends the default data model,
  adding new kinds of name part to it.
\item Copy the code from \textsf{93-nameparts.tex} into the .cbx
  files, along with some shims that allow the addition of new name
  templates and also allow any of the templates to serve as the
  default template, that is, the template to be used when a .bib entry
  doesn't specify one.
\item Define the \mycolor{\texttt{cmsnameparts}} option in such a way
  that none of this code will be loaded unless you call the option,
  which meant that I had to define it in \textsf{biblatex-chicago.sty}
  instead of in the individual styles themselves.  If you don't load the
  style in the standard way with a call to \textsf{biblatex-chicago}
  then a little more work will be required. 
\end{enumerate}

The \colmarginpar{\texttt{cmsnameparts}} basic usage is
straightforward.  If you just put \mycolor{\texttt{cmsnameparts}}
among the options when loading \textsf{biblatex-chicago} and add no
additional options to your .bib entries, nothing visible will change.
Any entry that contains one of \textsf{biblatex's}
\texttt{nametemplate} options will then be able to format its names
according to the Russian, Ethiopian, or CJK specifications.  You can
see this in action in the entry hua:cms in the
\textsf{biblatex-chicago} documentation
(\href{file:cms-notes-sample.pdf}{\textsf{cms-notes-sample.pdf}} and
\href{file:cms-dates-sample.pdf}{\textsf{cms-dates-sample.pdf}}).  It
uses \texttt{nametemplates=cjk} in its \textsf{options} field, which
is the simplest approach, though more fine-grained approaches are
possible with the \textsf{biblatex} options
\texttt{sortingnamekeytemplate} and \texttt{uniquenametemplatename}
(\textsf{biblatex.pdf} \S~3.8.10 and appendix E; cf.\ also
\textsf{93-nameparts.tex}).  Please note that when you use a
non-standard name part in a name list, you'll need to identify even
those standard parts that are present, e.g., \texttt{family=Hua,
  given=Linfu}.

\mylittlespace It's
\colmarginpar{\texttt{cmsnameparts=\\<nametemplate>}} possible that
you might have a .bib database for a specific project that contains a
majority of entries requiring a change of name template.  In this case
you can change the default template by giving its name to the
\mycolor{\texttt{cmsnameparts}} option.  Any entry not having a
\texttt{nametemplate} option will then use the default, which may save
you some typing.  If an entry requires the \textsf{biblatex} standard
default template, its name, following the example of
\textsf{93-nameparts.tex}, is \texttt{western}.

\mylittlespace It's also perfectly possible that none of the available
templates will suit your needs.  The simplest course, at least for
you, is to send me a feature request and a fairly thorough description
of the specification, and I'll try to add it to the list of available
templates.  If you prefer to code it yourself while still using the
\mycolor{\texttt{cmsnameparts}} mechanism, you'll need to look over
the code in \textsf{chicago-notes.cbx} or
\textsf{chicago-dates-common.cbx} and provide several macros according
to a specific naming scheme.  You'll need at least:

{\small
\begin{verbatim}
\DeclareSortingNamekeyTemplate[<templatename>]{<specification>}
\DeclareUniquenameTemplate[<templatename>]{<specification>}
\DeclareNamehashTemplate[<templatename>]{<specification>}
\newbibmacro*{name:<templatename>}[<no. of nameparts as digit>]{<code>}
\csdef{cmssort:<templatename>}{\usebibmacro{name:<templatename>}...}
\csdef{cmslabel:<templatename>}{...\usebibmacro{name:<templatename>}...}
\end{verbatim}}

If your name specification reverses the order of the first name in the
bibliography or reference list (as in the \texttt{western} style),
then you'll also need:

{\small
\begin{verbatim}
\newbibmacro*{name:<templatename>-rev}[<no. of nameparts as digit>]{<code>}
\csdef{cmssort:<templatename>-rev}{\usebibmacro{name:<templatename>-rev}...}
\end{verbatim}}

If you want to be able to use the your new template name as the
default template, you'll need to wrap the first three declarations like
so:

{\small
\begin{verbatim}
\cms@template@wrapper{\DeclareSortingNamekeyTemplate}[<templatename>]{<spec>}
\cms@template@wrapper{\DeclareUniquenameTemplate}[<templatename>]{<spec>}
\cms@template@wrapper{\DeclareNamehashTemplate}[<templatename>]{<spec>}
\end{verbatim}}

If your specification requires new name parts, you can copy
\mycolor{\texttt{cms.dbx}} to your working directory and add them to
\cmd{DeclareDatamodelConstant} there.  Similarly, if you already load
your own datamodel (.dbx) file using the \texttt{datamodel} option,
then if you want to use \mycolor{\texttt{cmsnameparts}}, which calls
that option internally, you should add the code from your own .dbx
file to \texttt{cms.dbx} and clear the \texttt{datamodel} option from
your preamble.  (That option can only load one file.)

\mylittlespace Finally, a few notes if you use
\textsf{biblatex-chicago} without loading
\textsf{biblatex-chicago.sty}.  You can easily load
\mycolor{\texttt{cms.dbx}} using \textsf{biblatex's}
\texttt{datamodel} option, or else just include the declaration there
in your own .dbx file.  In order to load the code in the .cbx files,
however, you'll need to define \cmd{cms@ldt@cmsnameparts} in your
preamble before loading \textsf{biblatex}.  A line like:

\begin{verbatim}
\gdef\cms@ldt@cmsnameparts{western}
\end{verbatim}

should do, or use another template name if you want something else as
the default.

\mylittlespace No matter how you load the code, I recommend at least a
quick read of \textsf{93-nameparts.tex} just to get an idea of how to
construct your .bib entries.  If your entries are entirely devoid of
western script, you may need a \cmd{DeclareSortTranslit} in your
preamble to tell \textsf{biblatex} how to sort such entries when mixed
with other scripts.  An example of its use for Russian entries can be
found in that file.  More generally, the nameparts feature is new to
\textsf{biblatex-chicago}, so please let me know if something doesn't
work as you expect.

\section{One\,.bib Database, Two Chicago Styles}
\label{sec:twostyles}

I have, when designing this package, attempted to keep at least half
an eye on the possibility that users might want to re-use a .bib
database in documents using the two different Chicago styles.  The
extensive unification of the two styles in recent editions of the
\emph{Manual} has simplified things, and though I have no idea whether
this will even be a common concern, I still thought I might gather in
this section the issues that a hypothetical user might face.  The two
possible conversion vectors are by no means symmetrical, so I provide
two lists, items within the lists appearing in no particular order.
These may well be incomplete, so any additions are welcome.

\subsection{Notes -> Author-Date}
\label{sec:conv:notesauth}

This is, I believe, the simpler conversion, as most well-constructed
.bib entries for the notes \&\ bibliography style will nearly
\enquote{just work} in author-date, but here are a few caveats
nonetheless:

\begin{enumerate}
\item \textbf{NB:} Unless you are using \textsf{authordate-trad}, the
  formatting of titles in the two styles is now the same, which means
  you would no longer need to worry about extra curly brackets and
  their effects on capitalization.  If you are using
  \textsf{authordate-trad}, please see the caveats in the
  documentation of the \textsf{title} field in
  section~\ref{sec:fields:authdate}, above.
\item You may need to reevaluate your use of shorthands, given that by
  default the author-date styles use them in place of authors rather
  than in place of the whole citation.  The preamble option
  \textsf{cmslos=false} may help, but this may leave your document
  out-of-spec.
\item The potential problem with multiple author lists containing more
  than three names doesn't arise in the notes \&\ bibliography style,
  so the \textsf{shortauthor} fields in such entries may need
  alteration according to the instructions in
  section~\ref{sec:otherhints:auth} above.
\item Date presentation is relatively simple in notes \&\
  bibliography, so you'll need to contemplate the \texttt{cmsdate}
  options from sections~\ref{sec:authuseropts} and
  \ref{sec:authentryopts} when doing the conversion to author-date.
\end{enumerate}

\subsection{Author-Date -> Notes}
\label{sec:conv:authnotes}

It is my impression that an author-date .bib database is somewhat
easier to construct in the first instance, but subsequently converting
it to notes \&\ bibliography is a little more onerous.  Here are some
of the things you may need to address:

\begin{enumerate}
\item If you've decided against using the \cmd{partedit} macro and
  friends from section~\ref{sec:formatting:authdate} above, commands
  not strictly necessary for author-date, you'll need to insert them
  now.
\item In general, you need to be more careful in notes \&\ bibliography
  about capitalization issues.  Fields which only appear once in
  author-date --- in the list of references --- may appear in both
  long notes and in the bibliography, in different syntactic contexts,
  so a quick perusal of the documentation of the \cmd{autocap} macro
  in section~\ref{sec:formatting:authdate} above may help.
\item You also need to be more careful about the use of abbreviations,
  e.g., in journal names, where the author-date style is more liberal
  in their use than the notes \&\ bibliography style.  (Cf.\ 14.170.)
  The bibstrings mechanism and package options sort much of this out
  automatically, but not all.
\item The \textsf{shorttitle} field is used extensively in notes \&\
  bibliography to keep short notes short, so you may find that you
  need to add a fair number of these to an author-date database.  In
  general this field is ignored by the latter style, so this, too,
  will be a one-time conversion.
\item You may need to add \textsf{letter} entries if you are citing
  just one letter from a published collection.  See
  section~\ref{sec:entrytypes}, s.v. \enquote{letter,} above.
\item The default shorthand presentation differs from one style to the
  other.  You may need to reconsider how you use this field when
  making the conversion.
\item As I explained above in section~\ref{sec:entryfields}, s.v.\
  \enquote{date,} I have included compatibility code in
  \textsf{biblatex-chicago-notes} for the \texttt{cmsdate} (silently
  ignored) and \texttt{switchdates} options, along with the automatic
  mechanism for reversing \textsf{date} and \textsf{origdate}.  This
  means that you can, in theory, leave all of this alone in your .bib
  file when making the conversion, though I'm retaining the right to
  revise this if the code in question demonstrably interferes with the
  functioning of the notes \&\ bibliography style.
\end{enumerate}

\section{Interaction with Other Packages}
\label{sec:otherpacks}

For \mymarginpar{\textbf{endnotes}} users of the \textsf{endnotes}
package --- or of \textsf{pagenote} --- \textsf{biblatex} offers
extensive compatibility options.  Please read the documentation of the
\texttt{notetype} option in \textsf{biblatex.pdf} \S~3.1.2.1.  If you
are using the \texttt{noteref} option with the notes \&\ bibliography
style and \textsf{endnotes}, please read the documentation in
section~\ref{sec:endnoterefs} for your options, which include the
\textsf{cmsendnotes.sty} package.

\mylittlespace Another \mymarginpar{\textbf{memoir}} problem I have
found occurs because the \textsf{memoir} class provides its own
commands for the formatting of foot- and end-note marks.  By default,
\textsf{biblatex-chicago} uses superscript numbers in the text, and
in-line numbers in foot- or end-notes, but I have turned this off when
the \textsf{memoir} class is loaded, reasoning that users of that
package may well have their own ideas about such formatting.

\mylittlespace The \mymarginpar{\textbf{ragged2e}} footnote mark code
I've just mentioned also causes problems for the \textsf{ragged2e}
package, but in this case a simple workaround is to load
\textsf{biblatex} \emph{after} you've loaded \textsf{ragged2e} in your
document preamble.  The URL line-breaking code activated by
\texttt{cmsbreakurl} doesn't play well with \textsf{ragged2e}, and so
far I've not found a workaround.

\section{TODO \&\ Known Bugs}
\label{sec:bugs}

This release implements the 17th edition of the \emph{Chicago Manual
  of Style}.  It contains a version of the author-date style
(\textsf{authordate-trad}) with traditional title formatting,
alongside the \textsf{authordate} code which unifies the treatment of
titles between itself and the notes \&\ bibliography style.  I
strongly encourage users to migrate to one of the styles implementing
the most recent specification, as I am focusing all of my development
and testing time there.

\mylittlespace Regardless of which edition you are considering, there
remain things I haven't implemented.  The solution in brown:bremer to
multi-part journal articles obviously isn't optimal, and I should
investigate a way of making it simpler.  If the kludge presented there
doesn't appeal, you can always, for the time being, refer separately
to the various parts.  If you have other issues with particular sorts
of citation, I'm of course happy to take them on board.  The
\emph{Manual} covers an enormous range of materials, but it seems to
me that the available entry types could be pressed into service to
address the vast majority of them.  If this optimism proves misguided,
please let me know.

\mylittlespace Kenneth L. Pearce has reported a bug that appears when
using multiple citation commands inside the \textsf{annotation} field
of annotated bibliographies.  I have tried to remedy this with the
\cmd{citeincite} and \cmd{fullciteincite} commands, which print short
and long citations, respectively. The first is intended for use inside
a small selection of fields, while the long version should only be
used in the \textsf{annotation} field.  See
sections~\ref{sec:citecommands} and \ref{sec:cite:authordate}, above,
for details.  If you still encounter problems with these commands then
it may be worth setting the \textsf{biblatex} option
\texttt{casechanger=latex2e} in the preamble to see if that helps,
though in my testing the two different case-changing backends are now
equivalent in this regard.

\mylittlespace Roger Hart, Pierric Sans, and a number of other users
have reported a bug in the formatting of title fields.  This, as far
as I can tell, has to do with the interaction between
\cmd{MakeSentenceCase} and certain characters at the start of the
title, particularly Unicode ones.  If you are using
\textsf{authordate-trad}, it may help for the moment to put an empty
set of curly braces \{\}\ at the start of the field, but I shall look
into this further.

\mylittlespace Please\colmarginpar{\textsf{biblatex} 3.20} note, also,
that in the \textsf{authordate-trad} style you can no longer by
default use a \cmd{bibstring} at the start of a \textsf{title} field,
due to the case-changing code in \cmd{MakeSentenceCase}.  You can
either set \texttt{casechanger=latex2e} in the preamble or just write
the string you need in the\ .bib file.

\mylittlespace Patrick Danilevici's Romanian localization
(\textsf{cms-romanian.lbx}) required some fairly extensive changes to
the basic package code, so it's possible that I could improve the
package's handling of other languages in a similar way.

\mylittlespace This release fixes the formatting errors of which I am
aware.  There remain the larger issues I've discussed throughout this
documentation, which mainly represent my inability to make all of
\textsf{biblatex-chicago's} formatting functions transparent for the
user, but thankfully \textsf{biblatex's} superb punctuation-tracking
code preemptively fixed a great many small errors, some of which I
hadn't even noticed before I began testing that functionality.  That
there are other micro-bugs seems certain --- if you report them I'll
do my best to fix them.

\section{Revision History}
\label{sec:history}

\textbf{2.3b: Released \today}

\mylittlespace This is an interim bug-fix release, mainly to address
the following breaking change:

\begin{itemize}
\item The most recent \textsf{biblatex} (3.20) has added a new name
  template (\mycolor{\cmd{DeclareNamehashTemplate}}) for users
  requiring non-Western name handling in their reference apparatus, so
  if you use \textsf{biblatex-chicago's} \texttt{cmsnameparts} option
  then running \textsf{biber} will fail and \LaTeX\ will not process
  your documents.  This update should rectify that, while still being
  compatible with earlier versions of \textsf{biblatex}.  Thanks to
  Philipp Immel for alerting me.
\item I have also added a package-specific definition of the standard
  citation command \cmd{citetitle} to the author-date styles, as the
  standard definition sometimes fails there.  Thanks to Geoffery Zheng
  for reporting this.
\item Users of the \textsf{authordate-trad} style may need to remove
  some \cmd{bibstring} commands from the start of \textsf{title}
  fields in their\ .bib files, or else set the \textsf{biblatex}
  option \texttt{casechanger=latex2e}.  Please see
  section~\ref{sec:bugs} or section~\ref{sec:fields:authdate}, s.v.\
  \textsf{title}.
\end{itemize}

\textbf{2.3a: Released November 11, 2022}

\mylittlespace This release backports a bug-fix from the most recent
\LaTeX 3 programming layer.  If you are using \textsf{biblatex} 3.18b
and an older L3 layer then the processing of your documents could fail
and fall into an infinite loop, leaving behind an inscrutable message
in your log file.  This update should prevent that, and also fixes a
few other minor glitches.

\mybigspace\textbf{2.3: Released July 2, 2022}

\mylittlespace All styles require the current \textsf{biblatex} (3.18)
and \textsf{biber} (2.18).

\begin{itemize}
\item Since the release of \textsf{biblatex} 3.17, the full array of
  \textsc{iso}8601-2 year divisions (seasons, quarters,
  quadrimestrals, semestrals, etc.)\ is available to its users, and
  \textsf{biblatex-chicago} has followed suit in all styles.  The
  \emph{Manual} itself, as far as I know, mentions only the seasons,
  but if you use a month specification higher than 24 in a
  \textsf{date} field you can access these new divisions.  Cf.\
  Table~\ref{ad:date:extras}.
\item Since \textsf{biblatex} provides a mechanism for localizing name
  presentation in the reference apparatus, Philipp Immel suggested I
  add an option to facilitate access to it for users of
  \textsf{biblatex-chicago}.  The \mycolor{\texttt{cmsnameparts}}
  option is the result, and is available in all styles.  Please
  consult section~\ref{sec:nameparts} above for all the details.
\item Philipp also pointed out that the name parts code requires the
  standard \textsf{biblatex} option \texttt{uniquename} to be turned
  on in the notes \&\ bibliography style.  (It's now set to
  \texttt{minfull}.)  This fixes an ancient bug, as Chicago has always
  required authors with the same surname to be distinguished by
  initials or full given names in short notes (14.32).
\item Patrick Danilevici pointed out that, in languages that don't use
  the Oxford comma, it was generally wrong to have a comma after the
  first, reversed name in an author list consisting of only two names.
  I have removed it in such languages, but if you still want the old
  behavior you can use \cmd{DeclareDelimFormat\{revsdnamedelim\}} to
  set it to \cmd{addcomma}.
\item After a request by Ryo Furue, I've added the
  \mycolor{\texttt{onlynd}} switch to the \texttt{doi} package option
  in all styles in order to suppress both the \textsf{url} and the
  \textsf{urldate}, rather than just the former.  See
  sections~\ref{sec:chicpreset} and \ref{sec:authpreset}.
\item After a bug report from Jon Arnold, I have tried to fix the
  field-inheritance schemes for \textsf{audio}, \textsf{music}, and
  \textsf{video} entries, which should now work more intuitively in
  all styles if you use \textsf{crossref} fields in them.  Concretely,
  this means that a \textsf{title} in one such entry (identifying a
  complete work or collection) can become a \textsf{booktitle} in the
  child entry, identifying the work of which the \textsf{title} is
  only one part.
\item Tobias Becht suggests that using specialized citation commands
  for printing author-only citations of ancient works (cf.\
  \emph{Manual} 14.244--5) is unnecessarily awkward.  I have therefore
  provided the \mycolor{\texttt{notitle}} option
  (sections~\ref{sec:useropts} and \ref{sec:authuseropts}) to turn a
  short author-title citation (whether in the notes \&\ bibliography
  style or the author-date styles) into an author-only citation.  You
  can set this option by entry type or in the \textsf{options} field
  of individual entries.
\item Tobias also pointed out that the \emph{Manual} permits using
  short citations of ancient, medieval, and Renaissance works right
  from the first citation in the notes \&\ bibliography style, even if
  you don't want to do the same for the entire document.  You can
  therefore now set the package option \mycolor{\texttt{short}} for
  the whole document, for specific entry types, and for specific
  entries.
\item Again at Tobias' prompting, I've modified how the ibidem
  mechanism works for entries with \texttt{classical}
  \textsf{entrysubtype}.  This includes following the \emph{Manual's}
  guidelines about suppressing the \emph{ibid.}\ string when using
  abbreviated forms of authors' names (unless you've set
  \texttt{useibid=true} in the \textsf{options} field of the entry
  itself).  I've also made the \mycolor{\texttt{noibid}} and
  \mycolor{\texttt{useibid}} options more granular, the former in all
  styles and the latter in notes \&\ bibliography.  You can now set
  them globally, per type, and per entry.
\item Having noticed that the \mycolor{\texttt{journaltitleaddon}}
  field is standard in \textsf{biblatex}, I have added it to the
  \textsf{article} and \textsf{review} types, as well as adding the
  \textsf{titleaddon} field to the \textsf{periodical} type, and this
  in all styles.  The \mycolor{\texttt{jtitleaddon}} option controls
  the separator between the main title and the annex in all these
  cases.  It is set to \texttt{space} by default.  See
  sections~\ref{sec:chicpreset} and \ref{sec:authpreset}.
\item Alanna Warner-Smith reported that the author-date styles weren't
  distinguishing citations with \textsf{extradate} letters when a
  \textsf{journaltitle} replaced the \textsf{author}.  This is, I
  hope, fixed.
\item After input from Claudius Ellsel, I've revised the
  \textsf{hyperref} functionality in the author-date styles, allowing
  you to choose which parts of a citation will act as links to the
  entry in the reference list.  Two new user options,
  \mycolor{\texttt{hyperall}} and \mycolor{\texttt{hypername}},
  complement the pre-existing \texttt{hypertitle} to give more
  fine-grained control.  All of these options are available globally,
  per type, and per entry.  See section~\ref{sec:authuseropts}.
\item After a request from Samuel Foster, I've provided the
  \mycolor{\texttt{suppressnoterefs}} option for the notes \&\
  bibliography style, settable per entry type and per entry, which
  prevents the printing of the \texttt{noteref} back reference after
  such entries.  This may be more convenient than using the
  specialized citation commands to achieve the same end.  Cf.\
  section~\ref{sec:noteref}.
\item There were many other bugfixes, including another attempt to
  make the \textsf{postnotepunct} option work more reliably,
  improvements to \textsf{shorthand} behavior in the author-date
  styles, and fixes, suggested by David Thiel, to the
  \textsf{hyperref} functionality when using \textsf{cmsendnotes.sty}.
\end{itemize}

\textbf{2.2: Released June 30, 2021}

\mylittlespace All styles still require the current \textsf{biblatex}
(3.16) and \textsf{biber} (2.16).  The next release of
\textsf{biblatex} (3.17) will necessitate some changes to Chicago's
date handling, so a new release of \textsf{biblatex-chicago} will, I
hope, follow swiftly on that of \textsf{biblatex}.

\begin{itemize}
\item Mar\c{c}al Orteu Punsola has very generously provided a Spanish
  localization for \textsf{bibla\-tex-chicago}
  (\mycolor{\textsf{cms-spanish.lbx}}).
\item I have modified several other localization files, fixing typos
  and ensuring minimal harmony when Chicago provides the short string
  (in notes) and standard \textsf{biblatex} the same string in long
  form (in the bibliography).
\end{itemize}

\textbf{2.1: Released March 27, 2021}

\mylittlespace This release contains several new features, some
modifications for compatibility with the latest \textsf{biblatex}, and
many bug fixes.  All styles require the latest \textsf{biblatex}
(3.16) and \textsf{biber} (2.16).

\begin{itemize}
\item Patrick Danilevici has very generously provided a Romanian
  localization for \textsf{bibla\-tex-chicago}
  (\mycolor{\textsf{cms-romanian.lbx}}) to go with the
  \textsf{romanian.lbx} he wrote for \textsf{biblatex} itself.  The
  only hitch is that the latter file won't be in the standard
  \textsf{biblatex} package until the next release, so you'll have to
  grab it from
  \href{https://github.com/plk/biblatex/tree/dev/tex/latex/biblatex/lbx}{GitHub}
  until then.
\item Moritz Wemheuer recommended that I include the standard
  \textsf{biblatex} macro \texttt{begentry} at the start of all
  drivers in both Chicago styles, with an empty definition by
  default.  This allows users to redefine it in their preambles and
  thereby inject their code into the driver before anything is printed.
\item Moritz also drew my attention to a question on
  \href{https://tex.stackexchange.com/questions/528374/moving-addendum-field-to-the-end-in-biblatex-chicago/540755#540755}{StackExchange}
  which suggested that I should expand the applicability of the
  \textsf{annotation} field, so I've modified the \texttt{annotation}
  option in all styles to allow the field to appear both in
  bibliography (reference list) and in long notes, and also added two
  new options in all styles, \mycolor{\texttt{bibannotesep}} and
  \mycolor{\texttt{citeannotesep}}, which allow you to define how the
  \textsf{annotation} field will relate to the main entry.  The
  options have different defaults, and although they have the same
  pre-defined keys those keys may have slightly different meanings
  depending on whether the \textsf{annotation} is to appear in notes
  or bibliography.  Please see the docs of field and options in
  sections~\ref{sec:annote}, \ref{sec:chicpreset}, \ref{sec:useropts},
  \ref{sec:fields:authdate}, \ref{sec:authpreset}, and
  \ref{sec:authuseropts}.
\item On a related subject, I have included two new options in all
  styles --- \mycolor{\texttt{formatbib}} and
  \mycolor{\texttt{entrybreak}} --- which provide fine-grained control
  of page breaking in bibliographies and reference lists.  Some of the
  values of these options are designed to be used with specific values
  of the \mycolor{\texttt{bibannotesep}} option, so please see the
  documentation in sections~\ref{sec:useropts} and
  \ref{sec:authuseropts}.
\item In the notes \&\ bibliography style Fr.\ Norbert Keliher
  requested a way to turn off the printing, in the first citation of a
  work, of the introduction --- (hereafter cited as \ldots)\ --- of a
  \textsf{shorthand} that will appear in subsequent citations.  The
  \mycolor{\texttt{shorthandintro}} option is now available globally,
  per type, and per entry to allow you to control the (non)appearance
  of this clarifying notification.  By default, the first citation of
  a work will continue to present the \textsf{shorthand} as it always
  has.  Please see section~\ref{sec:useropts} for the details.
\item Philipp Immel requested a feature involving the
  \textsf{nameaddon} field, and pointed to a discussion on
  StackExchange which suggested that others were also unsatisfied with
  that field's functionality, so in this release I've provided three
  new package options to allow users to mould it to their needs.  The
  options are called \mycolor{\texttt{nameaddon}},
  \mycolor{\texttt{nameaddonformat}}, and
  \mycolor{\texttt{nameaddonsep}}, and you can find all the details in
  the \textsf{nameaddon} docs in sections~\ref{sec:entryfields} and
  \ref{sec:fields:authdate}, above.  Cf.\ also the discussions of the
  options in sections~\ref{sec:chicpreset}, \ref{sec:useropts},
  \ref{sec:authpreset}, and \ref{sec:authuseropts}.  Allow me to
  emphasize that the package defaults remain exactly the same as
  before, so that if you don't set any of the new options no changes
  to your documents or\ .bib databases are necessary.
\item Philipp Immel also requested a way to use the standard
  \textsf{biblatex} option \texttt{datamodel} even when loading the
  Chicago styles with
  \verb+\usepackage{biblatex-chicago}+. Implementing this allowed me
  to regularize how the standard \texttt{mcite} and \texttt{natbib}
  options were passed through to \textsf{biblatex}, and also to
  support the new \textsf{biblatex} \mycolor{\texttt{casechanger}}
  option, which can help solve a bug reported by Peter Mukunda
  Pasedach.  See sections~\ref{sec:useropts} and
  \ref{sec:authuseropts}, above.
\item The new \textsf{relatedtype} \mycolor{\texttt{short}} is
  available in all styles to allow the presentation of short
  references rather than full ones when using the \textsf{related}
  mechanism. See sections~\ref{sec:related} and \ref{sec:authrelated}.
\item In a similar vein, I have provided new
  \mycolor{\cmd{citeincite}} and \mycolor{\cmd{fullciteincite}}
  commands to allow you, in some contexts, to cite a related work
  inside the fields of another entry, something which, although not
  recommended, is common practice, and which has frequently not worked
  quite right. The \mycolor{\cmd{citeincite}} and
  \mycolor{\cmd{citeincites}} commands always present short citations,
  and are available in several fields, while
  \mycolor{\cmd{fullciteincite}} and \mycolor{\cmd{fullciteincites}}
  always present long references, and are designed only to be used in
  the \textsf{annotation} field for cases where an annotated
  bibliography or reference list includes references to works not
  cited in the main text. See sections~\ref{sec:citecommands} and
  \ref{sec:cite:authordate} for the details.  If I've missed yet
  further tricky corners in this functionality please let me know.
\item After a request by Patrick Danilevici, I have added the four
  fields \mycolor{\textsf{eventdate}}, \mycolor{\textsf{eventtitle}},
  \mycolor{\textsf{eventtitleaddon}}, and \mycolor{\textsf{venue}} to
  the \textsf{inproceedings}, \textsf{mvproceedings}, and
  \textsf{proceedings} drivers in all styles, for compatibility with
  standard \textsf{biblatex} usage.  See the docs of those entry types
  in sections~\ref{sec:entrytypes} and \ref{sec:types:authdate}.
\item The \texttt{postnotepunct} option should generally work now,
  both in more contexts and with arbitrary \textsf{postnote} field
  contents.
\item The \mycolor{\texttt{dashed}} option governing the use of the
  3-em dash in bibliographies and reference lists is now available
  globally, per type, and per entry.
\end{itemize}

Note on the 16th-edition files:
\begin{itemize}
\item These have been updated for compatibility with the latest
  \textsf{biblatex} and \textsf{biber}, and there are also a number of
  bug fixes included.  I haven't added any new functionality to the
  files, so for instance if you want to use the new Romanian
  localization you'll need the current 17th-edition styles.  With this
  release I am marking the 16th-edition files as deprecated, and will
  remove them in a future release.  Please consider moving your
  documents to the current specification, where all of my development
  time is now spent.\label{deprec:obsol}
\end{itemize}

\textbf{2.0: Released April 20, 2020}

\mylittlespace Converting from the 16th to the 17th edition in your
.bib files and \LaTeX\ documents:
\begin{itemize}
\item The 17th edition of the \emph{Manual} no longer encourages use
  of \emph{ibid.}\ to replace repeated citations of the same work in
  the notes \&\ bibliography style, preferring instead to use the
  author's name alone, along with any page number(s).  If you wish to
  continue using \emph{ibid.}\ in that style, you need to set the new
  option \mycolor{\texttt{useibid=true}} when loading
  \textsf{biblatex-chicago} in your preamble.
\item If you are loading the package the old-fashioned way, with
  \textsf{biblatex} and the \texttt{style=} option instead of with
  \textsf{biblatex-chicago}, please be aware that there are two
  standard packages required by \textsf{biblatex-chicago} that aren't
  automatically loaded by \textsf{biblatex}:
  \mycolor{\textsf{xstring}} and \mycolor{\textsf{refname}}.  You'll
  need to load them in your preamble yourself.
\item If you've been using the \textsf{year} field to present decades
  like \texttt{1950s}, this will no longer work accurately in
  author-date citations.  The correct way to do so now is to use one
  of \textsf{biblatex's} \textsc{iso}8601-2 date specifications in the
  \textsf{date} field instead, to wit, \texttt{195X}.  Generally, I've
  tried to make \textsf{year} fields like \texttt{[1957?]} or
  \texttt{[ca.\ 1850]} continue to work properly, but here too the
  best thing to do is to use the new \textsf{date} features and
  present them like \texttt{1957?} or \texttt{1850\textasciitilde},
  respectively.  This will ensure that both sorting and punctuation
  work out properly.  See table~\ref{ad:date:extras}, and the
  \textsf{date} docs in sections~\ref{sec:entryfields} and
  \ref{sec:ad:date}, above.
\item If you have been using the \textsf{nameaddon} field to hold time
  stamps for online comments, then you should put the time stamp into
  the \textsf{date} or possibly \textsf{eventdate} field, instead,
  using the \textsc{iso}8601-2 format implemented by
  \textsf{biblatex}.  See the \textsf{date} and \textsf{nameaddon}
  field docs in sections~\ref{sec:entryfields} and
  \ref{sec:fields:authdate}, along with tables~\ref{tab:online:types},
  \ref{tab:online:adtypes}, and \ref{ad:date:extras}.
\item Following on from the previous item, the 17th edition of the
  \emph{Manual} includes more plentiful and more detailed instructions
  for presenting online materials than were available in previous
  editions.  For users of \textsf{biblatex-chicago} this means that
  there is now guidance for many more sorts of material than before,
  so if you have been improvising citations of this sort of material
  in previous releases it will be worth checking to see whether there
  are now clearer instructions available.
  Tables~\ref{tab:online:types} and \ref{tab:online:adtypes} summarize
  the new specifications for the notes \&\ bibliography and
  author-date styles, respectively.  Cf.\ in particular the new
  \mycolor{\texttt{commenton}} \textsf{relatedtype} in
  sections~\ref{sec:related} and \ref{sec:authrelated}.  Also, the
  \textbf{online} entry type now prints both \textsf{author} and
  \textsf{editor} (or other editorial role) if they exist, and I've
  moved the \textsf{addendum} field \emph{before} the \textsf{url},
  which fits better with indications in the \emph{Manual}.  If you've
  been using the \textsf{addendum} field to present citations of other
  entries (as in older versions of \textsf{biblatex-chicago}), please
  switch to the \textsf{related} mechanism, which works better anyway.
\item On the same subject, if you are using the notes \&\ bibliography
  style and are retaining the \textsf{crossref} field (instead of
  using the \mycolor{\texttt{commenton}} \textsf{relatedtype}) in
  \textsf{review} entries as a means of presenting comments on blogs,
  such entries are now subject to the settings of the
  \texttt{longcrossref} option and will appear in abbreviated form in
  some full notes and in the bibliography, as has always occurred in
  \textsf{incollection} entries, for example.  You can set
  \texttt{longcrossref} to \texttt{true} to get back the old behavior.
\item The 17th edition generally encourages more strongly than the
  16th the use of only one publisher in the \textsf{publisher} field.
  If you decide to retain more than one, and one of them is a part of
  an academic publishing consortium, it encourages you to specify this
  relationship rather than merely listing the two using the keyword
  \enquote{and} in the field.  Please see the documentation of the
  field in sections~\ref{sec:entryfields} and
  \ref{sec:fields:authdate}, above, for the rather minor (and rare)
  changes this might mandate for your\ .bib files.
\item The 17th-edition presentation of \textbf{music} entries has
  added a few pieces of information it seems to find desirable ---
  track number in \textsf{chapter} and specification of a lead
  performer's role in, e.g., \textsf{editortype} --- though the basic
  structure of a 16th-edition\ .bib entry remains unchanged.  Please
  see the documentation of that entry type in
  sections~\ref{sec:entrytypes} and \ref{sec:types:authdate}, above.
\item The 17th edition has added a couple of wrinkles to the
  \textbf{video} type specifications.  You can now put the broadcast
  network of a TV show in the \textsf{usera} field, and you can also
  supply the new \textsf{entrysubtype} \mycolor{\texttt{tvepisode}} to
  print the series title (\textsf{booktitle}) \emph{before} the
  episode title (\textsf{title}).  Please see the documentation of the
  entry type in sections~\ref{sec:entrytypes} and
  \ref{sec:types:authdate}, above.
\item Both Chicago styles now sentence case the \textsf{title} field
  in \textbf{patent} entries, so you may need to put curly braces
  around words that shouldn't appear in lowercase.
\item The \textsf{pubstate} field now has a more generalized
  functionality, while maintaining the specialized uses present in
  earlier releases.  In particular, please note now that almost
  anything you put in the field will be printed somewhere in the
  entry, and in the case of the author-date styles may appear in a
  somewhat different part of the entry than that to which you may have
  become accustomed.  If you wish to move this data back to the end of
  the entry in the author-date styles, then the \textsf{addendum}
  field might be of service.  The documentation in
  sections~\ref{sec:entryfields} and \ref{sec:fields:authdate} should
  help.
\item \textsf{Biber} is now the \emph{required} backend for all
  Chicago styles, including the 16th-edition files still included in
  the package.  If you have somehow been using some variant of
  \textsc{Bib}\TeX\ in the notes \&\ bibliography style up to now, I'm
  fairly confident it will no longer work.  Please switch to
  \textsf{biber}.
\end{itemize}

Other new features common to the notes \&\ bibliography and
author-date styles:
\begin{itemize}
\item Wouter Lancee has very generously provided a Dutch localization
  for \textsf{biblatex-chi\-cago}, called
  \mycolor{\textsf{cms-dutch.lbx}}.  You can use it by including
  \enquote{\texttt{dutch}} when loading \textsf{babel}.  Gustavo
  Barros has also very kindly provided a much-revised version of his
  \textsf{cms-brazilian.lbx}.
\item As mentioned above, this release for the first time implements
  \textsf{biblatex's} elegant and long-standing support for the
  \textsc{iso}8601-2 Extended Format date specification, which means
  there are now greatly enhanced possibilities for presenting
  uncertain and unspecified dates and date ranges, along with date
  eras, seasons, time stamps, and time zones.
  Table~\ref{ad:date:extras} summarizes the implementation for all
  Chicago styles, but see also the \textsf{date} field in
  sections~\ref{sec:entryfields} and \ref{sec:fields:authdate}, along
  with the new package options \mycolor{\texttt{alwaysrange}},
  \mycolor{\texttt{centuryrange}}, \mycolor{\texttt{decaderange}},
  \mycolor{\texttt{nodatebrackets}}, and
  \mycolor{\texttt{noyearbrackets}}.
\item I have also implemented year-range compression in all styles,
  governed by the new \mycolor{\texttt{compressyears}} option, set to
  \texttt{true} by default.
\item Constanza Cordoni requested a way to turn off the printing of
  the 3-em dash for repeated names in the bibliography or reference
  list, and the \emph{Manual} concedes that some publishers prefer
  this style.  Some of \textsf{biblatex's} standard styles have a
  \mycolor{\texttt{dashed}} option, so for compatibility purposes I've
  provided the same.  By default, I have set it to \texttt{true} to
  print the name dash, but you can set \mycolor{\texttt{dashed=false}}
  in your preamble to repeat names instead throughout your document.
\item Jan David Hauck suggested I extend the field-exclusion
  functionality beyond the package options already provided
  (sections~\ref{sec:chicpreset} and \ref{sec:authpreset}) by
  \textsf{biblatex-chicago}.  First, I made sure that all of those
  options could be set globally, per type, and per entry.  Second, I
  added the command \mycolor{\cmd{suppressbibfield}}, designed to
  appear in the preamble, and which will look something like:
\begin{verbatim}
  \suppressbibfield[entrytype,entrytype,...]{field,field,field,...}
\end{verbatim}
  Please see sections~\ref{sec:formatcommands} and
  \ref{sec:formatting:authdate} for the details.
\item After a request by user BenVB, I have added support for the
  \textsf{biblatex} \mycolor{\textbf{shortseries}} field, which allows
  you to present abbreviated \textsf{series} in book-like entries in
  all the styles.  You can use the \mycolor{\texttt{seriesabbrev}}
  option to control where in your document these abbreviated forms
  will appear.  By default, the field is ignored in all styles.  You
  can also print a list of series abbreviations, rather in the manner
  of a list of shorthands, using a command like:
  \verb+\printbiblist{shortseries}+.  Please see
  \mycolor{\textsf{shortseries}} in sections~\ref{sec:entryfields} and
  \ref{sec:fields:authdate}.
\item I have added a new preamble option,
  \mycolor{\texttt{cmsbreakurl}}, which attempts to follow the
  \emph{Manual's} instructions for line-breaking inside URLs.  It
  doesn't work 100\%\ accurately, and it doesn't play well with the
  \textsf{ragged2e} package, but in most circumstances it is at least
  closer to the Chicago ideal than the \textsf{biblatex} defaults.
  See sections~\ref{sec:useropts} and \ref{sec:authuseropts}.
\item The \emph{Manual} now specifies how to present \textsf{articles}
  with a unique numeric or alphanumeric ID, which you can place in the
  \textsf{eid} field.  If you've been using this field in previous
  releases you'll notice that the ID has moved to a different place in
  long notes, bibliography, and list of references.
\item In \textsf{jurisdiction} entries, the presentation of the
  \textsf{title} changes depending on whether it appears in a
  stand-alone citation or as part of the flow of running text, no
  matter whether the citation is in the main body or in a note.  I
  have provided the \mycolor{\cmd{runcite}} command, in both Chicago
  styles, for \textsf{jurisdiction} citations that appear in running
  text.
\item N.\ Andrew Walsh suggested that I allow editorial roles that
  aren't pre-defined \cmd{bibstrings} to appear as-is in entries, just
  as the standard \textsf{biblatex} styles do.  I have followed this
  advice for the \textsf{editortype}, \textsf{editoratype},
  \textsf{editorbtype}, \textsf{editorctype}, and \textsf{nameatype}
  fields, making sure to capitalize the string according to its
  context.
\item I have added the \textsf{venue} field to \textbf{misc} entries,
  both with and without an \textsf{entrysubtype}. It also appears in
  the new \mycolor{\textbf{performance}} type.
\item I have added the \textsf{version} and \textsf{type} fields to
  \textbf{book} entries to help with multimedia app content (14.268).
  This material fits quite well in such entries but needs extra fields
  to present information about the version of the app and also the
  system type on which it runs.  I added the \textsf{type} field to
  \textbf{article}, \textbf{review}, and \textbf{online} entries for
  presenting medium information for online multimedia (14.267).
\item I have added a new entry type, \mycolor{\textbf{dataset}}, to
  allow the citation of scientific databases.  Cf.\
  sections~\ref{sec:entrytypes} and \ref{sec:types:authdate}.
\item I have added the \textsf{number} field to \textbf{misc} entries
  with an \textsf{entrysubtype} to help cope with the varieties of
  location information in different archives.
\item The new entry type \mycolor{\textbf{standard}} is now available
  to cite standards published by national or international standards
  organizations.  If you have been using the \textsf{book} type for
  such material it might be worth switching to make sure your entries
  are more in line with the \emph{Manual's} specifications.  See the
  docs in sections~\ref{sec:entrytypes} and \ref{sec:types:authdate}
  for the details.
\item The new entry type \mycolor{\textbf{performance}} is now
  available for citing live performances.  You can sometimes also use
  a \textbf{misc} entry without an \textsf{entrysubtype}.
\item I have added the \textsf{eventdate} field to the \textbf{audio}
  entry type to help it cope with the presentation of podcasts, which
  are new to the 17th edition.  Please see the documentation of the
  entry type in sections~\ref{sec:entrytypes} and
  \ref{sec:types:authdate}, above.
\item I have added the \textsf{origdate}, \textsf{eventdate},
  \textsf{userd}, and \textsf{howpublished} fields to the
  \textbf{artwork} and \textbf{image} entry types, in response to
  additional information given in some of the \emph{Manual's}
  examples.  Please see the documentation of \textbf{artwork} in
  sections~\ref{sec:entrytypes} and \ref{sec:types:authdate}, above.
\item I have added the \textsf{maintitle}, \textsf{mainsubtitle}, and
  \textsf{maintitleaddon} fields to the \textbf{article},
  \textbf{periodical}, and \textbf{review} entry types, where it
  (they) will hold the the name of any larger (usually periodical)
  publication of which a blog is a part.  This departs from standard
  \textsf{biblatex} usage, but the need for two italicized titles
  demanded something like it.
\item I have added a new field-exclusion option,
  \mycolor{\texttt{urlstamp}}, set to \texttt{true} by default, which
  means that any time stamp associated with the \textsf{urldate} will
  always be printed.  You can set it to \texttt{false} in the preamble
  either for the whole document or for specific entry types, or in the
  \textsf{options} field of individual entries.  See the docs in
  sections~\ref{sec:chicpreset} and \ref{sec:authpreset}, above.
\item The \textsf{howpublished} field has accumulated a series of new
  functions in various entry types, bringing it far from its origins
  in \textsf{booklet}, \textsf{misc}, and \textsf{unpublished}
  entries.  Please see its documentation in
  sections~\ref{sec:entryfields} and \ref{sec:fields:authdate}.
\item In \textsf{inreference}, \textsf{mvreference}, and
  \textsf{reference} entries \textsf{biblatex-chicago} no longer
  considers any of the name fields (\textsf{author}, \textsf{editor},
  etc.) for sorting purposes in the bibliography or reference list,
  thus leaving the \textsf{title} as the first field to be considered.
  This may simplify the creation of .bib database entries.
\end{itemize}

New notes \&\ bibliography features:
\begin{itemize}
\item In keeping with indications in the 17th edition of the
  \emph{Manual}, I have provided a means for altering the syntax when
  presenting multi-volume works, i.e., for presenting the title of the
  whole series (\textsf{maintitle}) \emph{before} the title of
  individual volumes of that series (\textsf{title} or
  \textsf{booktitle}).  This involves the use of the new
  \textsf{relatedtypes} \mycolor{\texttt{maintitle}} and
  \mycolor{\texttt{maintitlenc}}, which may be used in
  \textsf{bookinbook}, \textsf{inbook}, \textsf{incollection},
  \textsf{inproceedings}, \textsf{letter}, \textsf{mvbook},
  \textsf{mvcollection}, \textsf{mvproceedings}, and
  \textsf{mvreference} entries.  Please see the detailed documentation
  of this feature in section~\ref{sec:related}, s.v.\
  \textsf{relatedtype} \mycolor{\texttt{maintitle}}.
\item I have implemented a new system of back references from short
  notes to long notes to help readers find fuller information about a
  source more quickly and conveniently, as envisaged by the
  \textsf{Manual}.  The feature is enabled with the
  \mycolor{\texttt{noteref}} option, and there are several sub-options
  to control where and what is printed:
  \mycolor{\texttt{fullnoterefs}}, \mycolor{\texttt{noterefinterval}},
  \mycolor{\texttt{noterefintro}}, \mycolor{\texttt{pagezeros}},
  \mycolor{\texttt{hidezeros}}, and \mycolor{\texttt{endnotesplit}}.
  The dependent package \mycolor{\textsf{cmsendnotes.sty}} can assist
  if you use endnotes instead of footnotes in this context.  It too
  has numerous options: \mycolor{\texttt{hyper}},
  \mycolor{\texttt{enotelinks}}, \mycolor{\texttt{noheader}},
  \mycolor{\texttt{blocknotes}}, \mycolor{\texttt{split}},
  \mycolor{\texttt{subheadername}}, \mycolor{\texttt{headername}},
  \mycolor{\texttt{runningname}}, and \mycolor{\texttt{nosubheader}},
  alongside the new commands \mycolor{\cmd{theendnotesbypart}} and
  \mycolor{\cmd{cmsintrosection}}.  Four new citation commands
  complete the provisions: \mycolor{\cmd{shortrefcite}},
  \mycolor{\cmd{shorthandrefcite}}, \mycolor{\cmd{shortcite*}}, and
  \mycolor{\cmd{shorthandcite*}}.  Please see
  section~\ref{sec:noteref} for all the details, and also
  \href{file:cms-noteref-demo.pdf}{\textsf{cms-noteref-demo.pdf}} for
  a brief example and explanation of some of the functionality.
\item I have ported, with modifications, the author-date package
  option \mycolor{\texttt{nodates}} to the notes \&\ bibliography
  style.  It is set to \texttt{true} by default.  In conjunction with
  the \mycolor{\texttt{nodatebrackets}} and
  \mycolor{\texttt{noyearbrackets}} options it provides an alternative
  presentation of uncertain dates.  See section~\ref{sec:chicpreset}.
\item Pursuant to a bug report by David Purton, I have recoded the
  various \cmd{headlesscite} commands and included a new one,
  \mycolor{\cmd{Headlesscite}}, which is the actually functional way
  to enforce capitalization at the start of such a citation, should
  you need to do so.
\end{itemize}

New author-date features:
\begin{itemize}
\item The \mycolor{\textbf{verbc}} field, which is standard but unused
  in the styles included in \textsf{biblatex}, allows the user
  fine-grained control over if and when an \textsf{extradate} letter
  (1976\textbf{a}) will appear after the year in citations and the
  list of references.  See its documentation in
  section~\ref{sec:fields:authdate}.
\item The new \mycolor{\texttt{authortitle}} type and entry option
  allows you to provide author-title citations in the text instead of
  author-date.  The \textsf{entrysubtype} value \texttt{classical}
  does the same, but there may be cases where using such an
  \textsf{entrysubtype} is impossible.  This is set to \texttt{true}
  by default for \mycolor{\textbf{dataset}} entries.
\item On the same subject, you can also use the new citation commands
  \mycolor{\cmd{atcite}} and \mycolor{\cmd{atpcite}} to achieve the
  same end.  The former prints a plain citation, the latter places it
  in parentheses.
\item In the default configuration, when you use a \textsf{shorthand}
  field the style will now sort properly by that field, which is the
  first thing to appear in reference list entries.  If you set
  \texttt{cmslos=false} in your preamble then this no longer applies,
  as the \textsf{shorthand} no longer appears in the reference list.
\end{itemize}

Note on the 16th-edition files:
\begin{itemize}
\item These have been updated for compatibility with the latest
  \textsf{biblatex} and \textsf{biber}, and there are also a number of
  bug fixes included, many of them already mentioned in changelog
  items above.  The \mycolor{\texttt{compressyears}} option is
  available and turned on by default, and so is the
  \mycolor{\texttt{dashed}} option.  Most of the new
  \textsc{iso}8601-2 Extended Format date specifications are
  available, also, though time stamps won't be printed, as that
  edition of the \emph{Manual} is mostly silent about them.
\end{itemize}

\textbf{1.0rc5: Released January 16, 2018}

\begin{itemize}
\item As Nikola Le\v{c}i\'c spotted, recent releases of
  \textsf{biblatex} have introduced some compatibility problems for
  \textsf{biblatex-chicago}, particularly with regard to the handling
  of the \textsf{origlanguage} field (now a list), but also through
  the renaming of several other fields and declarations (e.g.\
  \mycolor{\cmd{DeclareSortingTemplate}}).  I have improved the
  handling of the \textsf{origlanguage} list by including many new
  bibstrings in the package's localization files, but other changes to
  formatting macros have made backward compatibility with older
  releases of \textsf{biblatex} difficult or impossible.  Please
  upgrade to version 3.10 --- which has received the most testing ---
  to use these styles.
\item As I mentioned in the Notice (section~\ref{sec:Notice}), the
  17th edition of the \emph{Manual} has now appeared, and my
  development energies from this point will be devoted to upgrading
  all styles to conform to it.  You can still file bug reports against
  the 16th edition, but the next major feature release will be based
  on the 17th.  In preparation for these changes, I have removed all
  the 15th-edition files from the package.
\end{itemize}

Other New Features:

\begin{itemize}
\item After fielding multiple requests over the years, I have added
  three new entry types --- \mycolor{\textbf{jurisdiction}},
  \mycolor{\textbf{legal}}, and \mycolor{\textbf{legislation}} --- to
  allow the presentation of court cases, laws, treaties, congressional
  (parliamentary) debates and hearings, constitutions, and executive
  documents.  The first (\mycolor{\texttt{round}} and
  \mycolor{\texttt{square}}) and last (\mycolor{\texttt{canada}},
  \mycolor{\texttt{constitution}}, \mycolor{\texttt{hansard}},
  \mycolor{\texttt{hearing}}, \mycolor{\texttt{uk}}, and
  \mycolor{\texttt{un}}) introduce a number of new
  \textsf{entrysubtypes} to help with formatting quirks, including the
  presentation of Canadian and UK materials for inclusion in an
  otherwise US context.  There are also several new options
  (\mycolor{\texttt{legalnotes}}, \mycolor{\texttt{noneshort}},
  \mycolor{\texttt{short}}, and \mycolor{\texttt{supranotes}}) for
  controlling the output.  I have documented all of this in
  section~\ref{sec:legal} above, a separate section both because the
  specification really comes from the \emph{Bluebook} rather than the
  \emph{Manual}, and also because they are the only entry types
  treated identically by the notes \&\ bibliography style and the
  author-date styles (itself a formatting quirk).  You can also look
  at the sample files \mycolor{\textsf{legal-test.bib}} and
  \mycolor{\textsf{cms-legal-sample.pdf}} to see how you might
  construct your database entries.  Support for \emph{Bluebook}
  citations is in its infancy, so if you have ideas for sorting out
  its complexities more elegantly or spot any inaccuracies then I
  would be happy to hear about it.  The implementation is intended
  mainly for American documents, but there is some rudimentary
  localization for the other languages supported by
  \textsf{biblatex-chicago}.  The actual citations in such contexts
  would, let it be noted, fall outside of the \emph{Bluebook} spec.
\item I am grateful to Gustavo Barros for providing a Brazilian
  Portuguese localization for \textsf{biblatex-chicago}, contained in
  the \mycolor{\textsf{cms-brazilian.lbx}} file.
\item Gustavo also pointed out some instances where the package's
  \textsf{bibstrings} couldn't accommodate the needs of his
  localization, so with his help I've split the \texttt{recorded}
  string into \texttt{discrecorded} and \texttt{songrecorded}, then
  added it to all the .lbx files.  I've also added two new
  \textsf{bibstrings} for the \textsf{lista} field format:
  \mycolor{\texttt{subverbo}} and \mycolor{\texttt{subverbis}}.  I've
  added them to all the .lbx files, but only
  \mycolor{\textsf{cms-brazilian.lbx}} differs from the default.  If
  other languages need this change please let me know.
\item The same user also suggested a fix to \textbf{patent} entries:
  removing the comma from between the dates when the language doesn't
  use a comma in lists.
\item Timo Thoms pointed out some annoying inconsistencies when using
  the \textsf{hyperref} package with the author-date styles, and I
  have attempted to rectify them.  In citations, only the
  \textsf{date} portion should act as a link, if there is a
  \textsf{date}, otherwise a \textsf{title} or perhaps a
  \textsf{shorthand} will link to the entry in the list of references.
  If you have entries that you believe should present hyperlinks but
  don't, you can try setting the new \mycolor{\texttt{hypertitle}}
  option in their \textsf{options} fields.  Alternately, you can set
  the option to \texttt{true} globally in the preamble and then
  \textsf{titles} and \textsf{shorthands} will serve as links whether
  there's a \textsf{date} or not.  Cf.\
  section~\ref{sec:authuseropts}, above.
\item Bertold Schweitzer requested that the styles allow using the
  string \texttt{forthcoming} in the \textsf{pubstate} field to
  present sources that are yet to be published.  This is now supported
  in all styles, and has the additional benefit of rendering recourse
  to the \cmd{autocap} command unnecessary, as the styles print
  \cmd{bibstring\{forthcoming\}} where the \textsf{year} would
  normally appear.  Using the \textsf{year} field itself is, of
  course, still supported too.
\item The same user requested that I allow
  \mycolor{\texttt{newspaper}} as an exact synonym of
  \texttt{magazine} in the \textsf{entrysubtype} field of
  \textsf{article}, \textsf{review}, \textsf{periodical}, and
  \textsf{suppperiodical} entries.  I have provided this in all
  styles, and wherever you see \texttt{magazine} in this
  documentation then \texttt{newspaper} will work in exactly the same
  way.
\item Bertold also suggested that, following the example of Philip
  Kime's \textsf{biblatex-apa} package, I support the use of
  \textsf{related} functionality when presenting reviews, so that you
  can, for example, easily present multiple reviews of the same item.
  I have provided this functionality in all styles.  To enable it
  you'll need to set the \textsf{relatedtype} field to
  \mycolor{\texttt{reviewof}} in \textsf{article}, \textsf{review}, or
  \textsf{suppperiodical} types.  You should also read the
  documentation in section \ref{sec:related} or \ref{sec:authrelated},
  above, as this \textsf{relatedtype} works somewhat differently from
  the others.  The standard, manual way of citing such works remains,
  of course, available.
\item Jan David Hauck suggested that there was a need for an
  \mycolor{\texttt{unpublished}} \textsf{entrysubtype} to the
  \textsf{report} type, which would present the \textsf{title} in
  quotation marks (or plain roman in \textsf{authordate-trad}) instead
  of italics.  I can't quite tell if the \emph{Manual} agrees, but I
  have fulfilled this request in all styles.
\item The same user pointed out that standard \textsf{biblatex} and
  the discussion in the \emph{Manual} both suggest providing
  \textsf{venue}, \textsf{eventdate}, \textsf{eventtitle}, and
  \textsf{eventtitleaddon} fields for the \textsf{unpublished} type,
  thereby allowing for the further specification of unpublished
  conference papers and the like.  I have added these fields in all
  styles.
\item At the request of N.\ Andrew Walsh, the notes \&\ bibliography
  style now offers a way to disambiguate references to different
  sources which would ordinarily produce identical short notes, that
  is, where the \textsf{author} and \textsf{labeltitle} are the same.
  \textsf{Biblatex's} \texttt{uniquework} option is now active by
  default, and \textsf{biblatex-chicago} provides three new user
  options, one for choosing a disambiguating field, one for setting
  the punctuation between that field and the rest of the short note,
  and one for formatting the field using parentheses or square
  brackets --- \mycolor{\texttt{shortextrafield}},
  \mycolor{\texttt{shortextrapunct}}, and
  \mycolor{\texttt{shortextraformat}}, respectively.  Please see
  section~\ref{sec:useropts}, above, for the details, and note that
  \mycolor{\texttt{shortextrafield}} has to be set for the mechanism
  to print anything at all.
\item User P\'{e}t\`{u}r spotted two long-standing bugs: first, that
  the \texttt{url=false} option didn't stop the printing of the
  \textsf{urldate}, and second that empty parentheses would appear in
  some circumstances around non-existent dates in the author-date
  styles.  I have fixed both.
\item Philipp Immel wondered whether I could address a long-standing
  bug when presenting a \textsf{subtitle} after a \textsf{title} that
  ends in an exclamation point or question mark.  This bug has existed
  since the first release of the 16th-edition styles, and I think I've
  finally solved it now after the release of the \emph{Manual's} 17th
  edition.  (Cf.\ batson.)
\end{itemize}

\textbf{1.0rc4: Released May 2, 2017}

\mylittlespace Another bug-fix release.

\begin{itemize}
\item Marko Wenzel reported, and helped to fix, a fairly major problem
  with the date handling in the author-date styles, an issue I hadn't
  spotted when doing the date-related updates for 1.0rc2.
\item I've also fixed a long-standing inaccuracy in the date-handling
  code of \textbf{patent} entries in the author-date styles.  Such
  entries now behave as the documentation claims they do.
\end{itemize}

\textbf{1.0rc3: Released April 20, 2017}

\mylittlespace This is a minor bug-fix release.

\begin{itemize}
\item Charles Schaum reported a whitespace bug that appeared when
  using multiple languages with \textsf{Babel}.  This was introduced
  in the last release by some careless editing by me, and should be
  fixed now.
\item Charles also pointed me to a discussion about a problem using
  \textsc{Bib}\TeX\ with \textsf{bibla\-tex-chicago}.  Ulrike Fischer
  very kindly suggested an elegant solution, and I have integrated it
  into this release.
\end{itemize}

\textbf{1.0rc2: Released March 26, 2017}

\mylittlespace This is an interim release designed mainly to fix a
number of subtle issues, pointed out by several users, that appear
when you use the newest version of \textsf{biblatex} (3.7).  These
were mostly concentrated in the date-handling code, which I believe
now behaves correctly, and should do both with the newest
\textsf{biblatex} and with somewhat older releases.  A much larger set
of new features is still pending, but I have fixed some other bugs and
added a few new options:

\begin{itemize}
\item J.~P.~E.~Harper-Scott pointed out that, in ordinary British
  usage, day numbers are presented as plain cardinals rather than
  ordinals.  The \emph{Manual} itself also prefers this format, not
  only for American-style dates but also for British ones, so I think
  the previous behavior of the package was a bug.  I have in both
  styles set the default presentation of British day numerals to be
  plain cardinals, providing a new preamble option
  \mycolor{\texttt{ordinalgb}} restoring the previous default and
  printing ordinal dates when using the \texttt{british} language with
  \textsf{Babel}.
\item I have, in both styles, attempted to provide an improved
  \cmd{partedit} macro, the old one being inconvenient for users
  writing in French.  The new macro should work now without manual
  intervention to provide the correct form of the preposition (de or
  d').  If you are using the \texttt{french} option to \textsf{Babel},
  please take care to remove any hand-formatting you might have
  provided in these contexts.
\item Jan David Hauck has both reported a bug in the \cmd{gentextcite}
  code in the author-date styles and also pointed me to its solution,
  as provided by moewe on
  \href{http://tex.stackexchange.com/questions/326472/gentextcites-multiple-books-same-author-biblatex-chicago/326628#326628}{Stackexchange}.
  It turned out there were other bugs in that code, now also fixed.
\item User laudecir requested a way to present a \textsf{shorthand}
  even in the first citation of a given work.  The new
  \mycolor{\texttt{shorthandfirst}} option in the notes \&\
  bibliography style can be set to \texttt{true} either in the
  preamble or in individual entries, and should make this
  functionality simpler to activate than the \cmd{shorthandcite}
  command.
\item Also in the notes \&\ bibliography style, Stefan Bj\"{o}rk
  requested a way to turn off the printing of \textsf{url},
  \textsf{doi}, and \textsf{eprint} information in notes but not in
  the bibliography.  The new \mycolor{\texttt{urlnotes}} option, which
  you can set to \texttt{false} in the preamble or in individual
  entries, provides this.  Please note that it does not apply to
  \textsf{online} entries.
\item Several users pointed out the presence of warnings in .log files
  caused by deprecated grammar in the default Sorting Schemes of both
  styles.  These should now be fixed.
\end{itemize}

\textbf{1.0rc1: Released June 7, 2016}

\mylittlespace Obsolete and Deprecated Features:
\begin{itemize}
\item The 15th-edition styles are now obsolete, and have been moved to
  a new \mycolor{\texttt{obsolete}} subdirectory.  You can still use
  them as they stand, but they won't compile against the newest
  \textsf{biblatex}, so you'll have to make sure that you have an
  older version (2.9a, perhaps).  If you are still using them, I
  strongly urge you to consider switching to the the 16th-edition
  styles, which contain many new features and bug-fixes.
\item The old Chicago-specific option \texttt{usecompiler} is
  deprecated, and has been replaced by the standard \textsf{biblatex}
  \mycolor{\texttt{usenamec}}.  If you have been using the former in
  your preamble or in your .bib entries, please replace it with the
  latter, which works better across the board.  \texttt{Usecompiler}
  still \enquote{works,} just not very well.
\end{itemize}

Other New Features:

\begin{itemize}
\item Stefan Bj\"{o}rk has very generously provided a Swedish
  localization file for the package ---
  \mycolor{\textsf{cms-swedish.lbx}} --- which can be loaded and used
  with \textsf{babel} just like the other localizations.
\item I have added support for \mycolor{\textbf{related}}
  functionality to all the Chicago styles, including all the standard
  \textsf{biblatex} \mycolor{\textbf{relatedtypes}}.  It is turned on
  by default in all styles, but you can turn it off, or alter where
  the information is printed, using the \mycolor{\texttt{related}}
  option in the preamble or in individual entries.  In the notes \&\
  bibliography style, \textsf{related} information is printed by
  default only in the bibliography, but you can change that by setting
  the option.  In the author-date styles, it will only ever print in
  the list of references, depending on the option's setting.  Please
  see sections~\ref{sec:related} and \ref{sec:authrelated} for the
  details.
\item I have improved the name-handling code in all styles,
  regularizing the functioning of the \textsf{namea}, \textsf{nameb},
  and \textsf{namec} fields with respect to the other, standard
  \textsf{biblatex} names.  The former two in particular are newly
  available in the \textsf{collection} and \textsf{periodical} entry
  types, and \textsf{biblatex-chicago} now recognizes the standard
  \mycolor{\texttt{usenamea}}, \mycolor{\texttt{usenameb}}, and
  \mycolor{\texttt{usenamec}} toggles, the last replacing the
  deprecated \texttt{usecompiler} (as above).  You can also now use
  the \mycolor{\textbf{nameatype}} field just as you would an
  \textsf{editortype}, extending the possibilities for identifying
  certain roles attached specifically to \textsf{titles} as opposed to
  \textsf{booktitles} or \textsf{maintitles}.
\item After a request by user BenVB, I have added support for the
  \textsf{biblatex} \mycolor{\textbf{shortjournal}} field, which
  allows you to present abbreviated \textsf{journaltitles} in all the
  styles.  You can use the \mycolor{\texttt{journalabbrev}} option to
  control where in your document these abbreviated forms will appear.
  By default, the field is ignored in the notes \&\ bibliography
  style, and appears only in citations in the author-date styles.  You
  can also print a list of journal abbreviations, rather in the manner
  of a list of shorthands, using a command like:
  \cmd{printbiblist\{shortjournal\}}.  Even though the
  \textsf{periodical} entry type uses the \textsf{title} and
  \textsf{shorttitle} fields in place of \textsf{journaltitle} and
  \textsf{shortjournal}, these entries are included in this
  functionality, and controlled by the same \texttt{journalabbrev}
  option.  Please see s.v.\ \enquote{shortjournal} in
  sections~\ref{sec:entryfields} and \ref{sec:fields:authdate}.
\item Following a request by Arne Skj{\ae}rholt, and his generous
  provision of some code to get me started, I have implemented a new
  \mycolor{\cmd{gentextcite}} citation command in all styles.  The
  \enquote{gen} part of the name refers to the genitive case, and it
  adds a possessive ending --- \textbf{'s} by default --- to the
  author's name in what is otherwise an ordinary \cmd{textcite}.  You
  can change the added ending however you want, using a third optional
  field to the citation command, and you can control to which names
  the ending is added in a \mycolor{\cmd{gentextcites}} multicite
  command by using the \mycolor{\texttt{genallnames}} preamble and
  entry option.  Please see sections~\ref{sec:citecommands} and
  \ref{sec:cite:authordate} for the details.
\item Stefan Bj\"{o}rk pointed out that \textsf{url}, \textsf{doi},
  and \textsf{eprint} information could be totally ignored in some
  entries when you used the abbreviated cross-referencing
  functionality accessed through the \textsf{crossref} and
  \textsf{xref} fields.  At his request, I have provided a new
  \mycolor{\texttt{xrefurl}} entry and preamble option for all the
  styles to control the printing of this information in abbreviated
  notes or bibliography (reference list) entries.  Please see
  sections~\ref{sec:useropts} and \ref{sec:authuseropts} for the
  details.
\item In a related change, I have stopped child entries inheriting
  \textsf{url}, \textsf{doi}, and \textsf{eprint} fields from their
  cross-ref'd parents, so if your documents rely on this behavior
  please note that you'll have to provide such fields manually in the
  child entries.
\item Roger Hart long ago requested a way to control the punctuation
  before \textsf{book-}, \textsf{main-}, or plain \textsf{titleaddon}
  fields, and I have finally added it in this release in the form of
  two entry and preamble options, \mycolor{\texttt{ptitleaddon}} and
  \mycolor{\texttt{ctitleaddon}}, available in all styles.  By
  default, the former prints \cmd{addperiod}\cmd{addspace}, hence its
  name, and the latter\,\cmd{addcomma}\cmd{addspace}, but you can
  change either or both depending on which field you are using and
  which sort of entry it appears in --- the default output can be your
  guide to which option(s) to change.  Please see the available valid
  option keys in sections~\ref{sec:chicpreset} and
  \ref{sec:authpreset}.
\item The same user also long ago requested that the notes \&\
  bibliography style make it possible to use \emph{Idem} when two
  consecutive notes cite different works by the same author.  You can
  now use the standard \textsf{biblatex} option
  \texttt{idemtracker=constrict} in your preamble to activate this in
  your documents, but please be aware, first, that the \emph{Manual}
  doesn't exactly approve of this and, second, that you'll only see
  \emph{Idem} in short notes, never in full ones, which seems to be
  the standard (\textsf{biblatex}) way of implementing this.
\item Also only in the notes \& bibliography style, I have added a
  \mycolor{\texttt{shorthandpunct}} option to control the punctuation
  that appears before the first appearance of a \textsf{shorthand}
  and/or a \textsf{shorthandintro} in a long note.  The default is
  \cmd{addspace}, but you can change it in your preamble or in
  individual entries.  Please see the available valid option keys in
  section~\ref{sec:chicpreset}.
\item After reading a discussion started by Ryo Furue at
  \href{https://github.com/plk/biblatex/issues/363}{github}, I have
  added, in the author-date styles only, a test to some spacing
  commands to prevent line breaks immediately after abbreviation dots.
  These tests apply only in running text, never in the list of
  references, where good line breaks are already hard enough to find.
\item In addition to moving all the 15th-edition styles into the
  \texttt{obsolete} subdirectory, I have also reorganized the
  author-date style files, adding
  \mycolor{\textsf{chicago-dates-common.cbx}} which contains the code
  that is common to the \texttt{trad} and the standard
  \texttt{authordate} styles.  Nothing has changed in terms of loading
  the styles, the changes being designed primarily to ease
  maintenance.
\item I have created two new documentation files (and an appendix) to
  provide short introductions to the Chicago styles, introductions
  which attempt to fill the gap between the Quickstart section
  (\ref{quickstart}) and the fuller documentation contained in
  sections~\ref{sec:Spec} and \ref{sec:authdate} of this file.  Both
  \mycolor{\textsf{cms-notes-intro.pdf}} and
  \mycolor{\textsf{cms-dates-intro.pdf}} are fully hyperlinked so you
  can move easily from formatted citations and (annotated) references
  to .bib entries and back, with marginal references to the fuller
  discussions here.  There is also a short
  \mycolor{\textsf{cms-trad-appendix.pdf}} file to discuss a few
  entries that would need special treatment for the \texttt{trad}
  style.  The \textsf{sample} files for each style still exist, but I
  intend them mainly for testing purposes, while many more (annotated)
  entries are still available for consultation in
  \textsf{notes-test.bib} and \textsf{dates-test.bib}.
\item I have made a number of other small enhancements to and fixed
  numerous bugs in all the styles, including some subtle inaccuracies
  in author-date citations spotted by Arne Skj{\ae}rholt and some
  macros in \textsf{inproceedings} entries that had been missing for
  years.  I have provided some default values for counters in
  \textsf{biblatex-chicago.sty} that aid in breaking long
  \textsf{urls} across lines, but I make no pretense that these fully
  adhere to the \emph{Manual's} specifications.  I have added a few
  \texttt{bibstrings}, currently missing in standard
  \textsf{biblatex}, to \textsf{cms-german.lbx} for use with the
  \mycolor{\textsf{related}} functionality.  Recommendations for
  better ones would be gratefully received.
\end{itemize}

\textbf{0.9.9i: Released May 16, 2016}
\begin{itemize}
\item This is another interim release, allowing the use of
  \textsf{biblatex} 3.4 for those who want to try it.  I have also
  fixed one old formatting error when \enquote{n.d.} appears in
  author-date citations.  A full feature-release based on 3.4 is
  imminent.
\end{itemize}

\textbf{0.9.9h: Released March 22, 2016}
\begin{itemize}
\item This is an interim bug-fix release, updating the styles so that
  they will work with \textsf{biblatex} 3.3.  The notes \&\
  bibliography style, as pointed out by several users, wouldn't
  compile at all with the newest \textsf{biblatex} version, and all
  styles had inaccuracies in the presentation of names due to changes
  in the name-handling code in \textsf{biblatex}.  I've done some
  testing against \textsf{biblatex} 3.3, and fixed all the errors I've
  spotted, but there may still be parts of my code that need updating
  to work well with the current version, so you can still downgrade to
  an earlier \textsf{biblatex} --- I recommend 2.9a --- if 3.3 doesn't
  work for you.  The next release will be a feature release, so if
  you've made a request, it should be fulfilled then.

\item I've also fixed a couple of long-standing bugs, one in the entry
  options controlling abbreviated cross-references and another in the
  formatting of the \textsf{prenote} field, the latter identified
  (ages ago) by Bernd Rellermeyer.

\end{itemize}

\textbf{0.9.9g: Released August 21, 2014}
\begin{itemize}
\item Alexandre Roberts found a showstopper in the functionality
  related to the new \mycolor{\texttt{inheritshorthand}} option in the
  notes \&\ bibliography style, and I found an unpleasant bug in the
  formatting of abbreviated cross-references in the same style.  This
  release, I hope, fixes both, but is in all other respects identical
  to 0.9.9f.
\end{itemize}

\textbf{0.9.9f: Released August 15, 2014}
\begin{itemize}
\item I've made the alterations needed to bring the styles into line
  with the latest version of \textsf{biblatex} (2.9a).  This is the
  version that has been tested most thoroughly with
  \textsf{biblatex-chicago}, so I strongly recommend using it.
\item I fixed several inaccuracies in the presentation of abbreviated
  cross-references in all the Chicago styles, and while I was working
  on that portion of the code it seemed an opportune moment to fulfill
  some feature requests bearing on the same area of functionality.
\item First, following a request from Alexandre Roberts, I have added
  the \mycolor{\texttt{inheritshort\-hand}} option to the notes \&\
  bibliography style, which allows child entries to inherit the
  \textsf{shorthand} field from their parents.  This in turn allows
  the \textsf{shorthand} itself to appear in place of the usual
  abbreviated citation of parent entries cross-referenced by several
  different child entries, thereby saving some space.  (This behavior
  was already available in the author-date styles, so the option is
  unnecessary there.)  You'll need to use \texttt{skipbiblist} in the
  \textsf{options} field of child entries to make the list of
  shorthands work correctly.  Please see the documentation of the
  \textsf{shorthand} field for the full explanation.
\item Second, following a request from Kenneth Pearce, I have added to
  all Chicago styles the capacity to combine abbreviated
  cross-references with the presentation of the original text of
  translations (via the \textsf{userf} field) or of the original
  publication details of an essay or chapter you are citing from a
  subsequent reprint (via the \textsf{reprinttitle} field).  See the
  documentation of those fields, and also of \textsf{crossref}, and
  note that you can now, taking certain precautions as outlined in the
  \textsf{shorthand} docs, combine the \textsf{userf},
  \textsf{crossref}, and \textsf{shorthand} fields.  This mechanism
  contains a great many moving parts, so please report any problems
  you might have with it.
\item Third, and finally, following a bug report by Mark van Atten I
  have fixed all Chicago styles so that the \textsf{biblatex}
  \texttt{backref} mechanism works properly in
  \textsf{biblatex-chicago}, including in those entries that use
  abbreviated cross-references, and in those that use the
  \textsf{userf} or \textsf{reprinttitle} fields.  I can't see any
  instructions concerning this in the \emph{Manual}, so I've left the
  formatting of \texttt{backref} lists in the hands of
  \textsf{biblatex} itself.  If the default behavior doesn't match
  your needs, let me know, as it's possible I could add some further
  options for modifying it.
\item I have added a new \mycolor{\texttt{compresspages}} option to
  all the Chicago styles.  If set to \texttt{true} it automatically
  compresses page ranges in the \textsf{pages} and \textsf{postnote}
  fields, allowing you to type ranges naturally, e.g., 101-{-}109, and
  letting the package follow the \emph{Manual's} rules for you.  (In
  this case, it would yield 101--9 in the document.)  Thanks are due
  to David Gohlke who brought to my attention a discussion that took
  place a couple of years ago on
  \href{http://tex.stackexchange.com/questions/44492/biblatex-chicago-style-page-ranges}{Stackexchange}
  regarding the automatic compression of page ranges.
  \textsf{Biblatex} has long had the facilities for providing this,
  and though the \emph{Manual's} rules (9.60) are fairly complicated,
  Audrey Boruvka fortunately provided in that discussion code that
  implements the specifications.  As some users may well be accustomed
  to compressing page ranges themselves in their .bib files, and in
  their \textsf{postnote} fields, I have made the activation of this
  code a package option.
\item Several users, most recently David Gohlke, have requested a way
  to alter the punctuation that appears just before the
  \textsf{postnote} argument of citation commands.  This allows, in
  the notes \&\ bibliography style, citations to fit better into the
  flow of text, while in the \textsf{authordate} styles it allows you
  very easily to insert comments, which follow a semi-colon, inside
  parenthetical text citations.  This punctuation is a complex issue
  in the \emph{Manual}, but as a first stab at enabling this greater
  flexibility, I have introduced the \mycolor{\texttt{postnotepunct}}
  package option.  Set to \texttt{true}, it allows you to start the
  \textsf{postnote} field with a punctuation mark (.\,,\,;\,:) and
  have it appear as the \cmd{postnotedelim} in place of whatever the
  package might otherwise automatically have chosen.  Please note that
  this functionality relies on a very nifty macro by Philipp Lehman
  which I haven't extensively tested, so I'm labeling this option
  \mycolor{experimental}.  Note also that the option only affects the
  \textsf{postnote} field of citation commands, not the \textsf{pages}
  field in your .bib file.  Note, finally, that if you are using the
  new \mycolor{\texttt{compresspages}} option then any
  \textsf{postnote} field starting with a punctuation mark will
  require you to do the compression of page ranges yourself.
\item I've added a new inheritance declaration so that
  \textbf{incollection} entries can inherit from \textbf{book} entries
  the same way they inherit from \textbf{mvbook}.
\item I've fixed a fair number of other bugs, including two in the
  \emph{Ibidem} mechanism identified by Bernd Rellermeyer, one in the
  printing of dates, and one in the \cmd{textcite} command in the
  notes \&\ bibliography style, these last two pointed out by Kenneth
  Beesley.  The presentation of all the periodical entry types
  (without an \textsf{entrysubtype}) has also been made more accurate.
\end{itemize}

\textbf{0.9.9e: Released January 29, 2014}
\begin{itemize}
\item This minor release fixes a regression in the \emph{Ibidem}
  mechanism in the notes \&\ bibliography style, spotted by Harold
  Bellemare, and present in the package since version 0.9.9c.  In all
  other respects this release is identical to 0.9.9d.
\end{itemize}

\textbf{0.9.9d: Released October 30, 2013}
\begin{itemize}
\item Following requests by Kenneth~L.\ Pearce and Bertold Schweitzer,
  I have modified and extended the mechanism for creating abbreviated
  citations when several parts of the same collection are included in
  a reference apparatus.  To the \textbf{incollection},
  \textbf{inproceedings}, and \textbf{letter} entries of previous
  releases, I have added \textbf{inbook}, \textbf{book},
  \textbf{bookinbook}, \textbf{collection}, and \textbf{proceedings}
  entries.  Only \textbf{inbook} entries join the former three in
  having this functionality turned on by default --- if you don't want
  this, it will require intervention either in the preamble or in the
  \texttt{options} field of individual entries.  This intervention
  will be via the new \mycolor{\texttt{longcrossref}} option, which
  controls the behavior of the four essay-like entry types and
  defaults to \texttt{false}, while the new
  \mycolor{\texttt{booklongxref}} option controls the four book-like
  types and defaults to \texttt{true}.  The useful settings for the
  options differ slightly between the author-date and the notes \&\
  bibliography specifications, so please see all the details in the
  docs of the \textsf{crossref} field in
  sections~\ref{sec:entryfields} and \ref{sec:fields:authdate}, above.
\item On the same subject, in the notes \&\ bibliography style, I
  should mention that in the first, full citation of one part of a
  collection in a note, the code no longer uses a separate citation of
  the parent entry to supply parts of what you see printed.  (This led
  to numerous inaccuracies.)  If your setup uses a side-effect of the
  old code to print data that hasn't even been inherited by the child,
  you may find that you need to change some \textsf{xref} fields to
  \textsf{crossref} fields to make it work correctly now.  This
  situation will, I imagine, be very rare, but you can look at
  white:ross:memo in \textsf{notes-test.bib} to see an example.
\item In the author-date styles, several users have been frustrated by
  the lack of an approved way of setting the \texttt{cmsdate} option
  in the preambles of their documents, and Kenneth~L.\ Pearce
  requested that I attempt to ease the burden on users by looking at
  this again.  With this release, you can now set \texttt{cmsdate}
  either to \texttt{both} or \texttt{on} in the preamble, and it will
  affect all entries (except \textbf{music}, \textbf{review}, and
  \textbf{video}) with multiple dates.  You can still change this
  setting in the \texttt{options} field of individual entries, but
  what you won't be able to change there is the new call to
  \cmd{DeclareLabeldate} which puts the \textsf{origdate} first in the
  list of dates when \textsf{Biber} searches for a \textsf{labelyear}
  to use in citations and in the list of references.  If you have been
  using the \texttt{switchdates} mechanism to get the
  \textsf{origdate} as the \textsf{labeldate}, your .bib files may
  need some editing in order to use the new preamble options.  Please
  see the documentation of the \textsf{date} field in
  section~\ref{sec:fields:authdate} above for all the (voluminous)
  details.
\item Following a request by Rasmus Pank Rouland, I adapted new
  \textsf{biblatex} code in the \cmd{textcite(s)} commands in all
  styles to make them fit more elegantly in the flow of text.  Upon
  reconsideration of the commands in the notes \&\ bibliography style,
  I slightly modified them, but \emph{only} when used inside a foot-
  or endnote.  In this context, by default, for both \cmd{textcite}
  and \cmd{textcites}, you'll now get the \textsf{author's} name(s)
  followed by a headless \emph{short} citation (or citations) placed
  within parentheses.  You can use \cmd{renewcommand} in the preamble
  of your document to redefine the new \mycolor{\cmd{foottextcite}}
  and \mycolor{\cmd{foottextcites}} commands to change this
  formatting.  See section~\ref{sec:formatcommands}, above.
\item This release includes support, in all styles, for
  \textsf{biblatex's} multi-volume entry types: \textbf{mvbook},
  \textbf{mvcollection}, \textbf{mvproceedings}, and
  \textbf{mvreference}.  See sections~\ref{sec:entrytypes} and
  \ref{sec:types:authdate}.
\item If you use \textsf{Biber}, I have added several new inheritance
  schemes to all styles to make cross-referenced entries work more
  smoothly: \textbf{incollection} entries can now inherit from
  \textbf{mvbook} just as they do from \textbf{mvcollection} entries;
  \textbf{letter} entries now inherit from \textbf{book},
  \textbf{collection}, \textbf{mvbook}, and \textbf{mvcollection}
  entries the same way an \textbf{inbook} or an \textbf{incollection}
  entry would; the \textsf{namea}, \textsf{nameb}, \textsf{sortname},
  \textsf{sorttitle}, and \textsf{sortyear} fields, all highly
  single-entry specific, are no longer inheritable; and the
  \textsf{date} and \textsf{origdate} fields of any \textbf{mv*} entry
  will \emph{not} be inherited by any other entry type.
\item Following a bug report by Henry~D.\ Hollithron, I've added to
  \textbf{unpublished} entries in all styles the possibility of
  including an \textsf{editor}, \textsf{translator}, etc.
\item Thanks to bug reports from Denis Maier and Bertold Schweitzer, I
  corrected inaccuracies and outright bugs in many entry types in all
  Chicago styles that appeared when there was a \textsf{booktitle} and
  not a \textsf{maintitle} or vice versa.  This also involved another
  rewrite of the code handling the \textsf{volume} field and other
  related fields in all non-periodical entry types that use them.
\item On the subject of the \textsf{volume} field, I added a new
  preamble and entry option, \mycolor{\texttt{delayvolume}}, to the
  notes \&\ bibliography style.  In long notes where this data isn't
  printed before a \textsf{maintitle}, this option allows you to print
  it \emph{after} the publication information rather than
  \emph{before} it, as may sometimes help clarify things, according to
  the \emph{Manual}.  This applies to the non-periodical entry types
  only.  See section~\ref{sec:useropts}.
\item On the same subject, in all styles, I have added a new preamble
  and entry option, \mycolor{\texttt{hidevolumes}}.  This controls
  whether, in entries where a \textsf{volume} has been printed before
  a \textsf{maintitle}, any \textsf{volumes} field present will also
  be printed, in this case \emph{after} the \textsf{maintitle}.  By
  default, this is set to \texttt{true}, so that the \textsf{volumes}
  field won't appear in such circumstances.  See
  sections~\ref{sec:chicpreset} and \ref{sec:authpreset}.
\item On the same subject, I have modified, in all styles, the field
  format for the \textbf{part} field, so that if the field contains
  something other than a number, \textsf{biblatex-chicago} will print
  it as is, capitalizing it if necessary, rather than supplying the
  usual bibstring, thus providing a mechanism for altering the string
  to your liking.  I have also decoupled the \textsf{part} field from
  the \textsf{volume} field, allowing it to be printed even in the
  absence of the latter, thereby providing a means to refer to
  segments of a larger work that don't easily fit the established
  schemes.  The iso:electrodoc entry in \textsf{dates-test.bib} shows
  an example of how this might work.
\item There is a new \mycolor{\texttt{omitxrefdate}} preamble and
  entry option in the notes \&\ bibliography style.  It turns off the
  printing of the child's \textsf{date} next to its \textsf{title} in
  abbreviated book-like entries \emph{only}, in both notes and
  bibliography.  See section~\ref{sec:useropts}.
\item Clea~F.\ Rees requested a way to customize the punctuation when
  a \textsf{volume} and a \textsf{page} number appear together like
  so: \enquote{2:204.}  You can use \cmd{renewcommand} in your
  preamble to redefine the new \mycolor{\cmd{postvolpunct}} command to
  achieve this, in all styles.  If your document language is French,
  \textsf{cms-french.lbx} redefines this already and prints something
  like \enquote{2 : 204.}  See sections~\ref{sec:formatcommands} and
  \ref{sec:formatting:authdate}.
\item I extended, in all styles, the functions of the \textsf{userd}
  field, allowing it to modify a \textsf{date} field if it hasn't
  already been captured by another date specification in the entry.
  See the documentation of the field in sections~\ref{sec:entryfields}
  and \ref{sec:fields:authdate}.
\item A bug report from Mathias Legrand helped clear up inaccuracies
  in the presentation of ordinal numbers in all styles.
\item For the author-date styles, another bug report by Kenneth Pearce
  resulted in the addition of the \textsf{labelyear} to the default
  \texttt{cms} sorting scheme so that more entries in the reference
  list are sorted properly without further user intervention.
\item George Pigman found an odd punctuation-tracking bug in the
  author-date styles.  This has been fixed.
\item Marc Sommer found a bug in the presentation of the
  \textsf{prenote} field in the author-date styles.  This has been
  fixed.
\item In the notes \&\ bibliography style, I improved the behavior of
  abbreviated foot- and endnotes when using the \textsf{hyperref}
  package.
\item I modified the date-presentation code in all the language files
  (\textsf{cms-*.lbx}) provided by the package.  Now, if an entry
  contains a \textsf{(*)year} and an \textsf{(*)endyear} that are
  exactly the same, and there aren't any further month or day
  specifications, then the \textsf{year} alone will be printed.  This
  allows for the clearing of spurious \textsf{endyears} inherited from
  parent entries.
\item I discovered some unpleasant side effects of my arrangement of
  the \textsf{.lbx} files devoted to Norwegian, and reverted to the
  arrangement as originally provided by H{\aa}kon Malmedal.
\end{itemize}

\textbf{0.9.9c: Released March 15, 2013}

\begin{itemize}
\item Antti-Juhani Kaijahano has very kindly provided a new Finnish
  localization for \textsf{biblatex-chicago}, called
  \mycolor{\textsf{cms-finnish.lbx}}.  As you will see if you look
  through it, it is still something of a work in progress.  If you
  would like to fill some of its lacunae, please do let me know.
\item Following a report by Bertold Schweitzer, I have added the
  \textbf{namea} and \textbf{nameb} fields to \textbf{article} and
  \textbf{review} entries in all three Chicago styles.  As in all the
  book-like entry types, they allow you to associate an editor or a
  translator specifically with a \textsf{title}, rather than, in these
  cases, with an \textsf{issuetitle}.  See the docs on these entry
  types in sections~\ref{sec:entrytypes} and \ref{sec:types:authdate},
  above.
\item Thanks to another report by Bertold I have, in all three Chicago
  styles, corrected inaccuracies in the presentation of the
  \textbf{report} entry type.  The \textsf{number} now appears
  immediately after the \textsf{type}, and the \textsf{type} itself is
  now capitalized properly depending on its context in an entry.
\item A third report by Bertold, detailing inaccuracies in the
  treatment of the \textbf{volume} and \textbf{volumes} fields in
  certain contexts, has resulted in a complete rewrite of the
  presentation of these (and several related) fields in all
  non-periodical entry types in all three Chicago styles.  This won't
  require any changes to your .bib files, but the output you see may,
  in some reasonably unusual situations, change.  Please let me know
  if something doesn't look right to you.
\item A fourth report by Bertold revealed some inadequacies with
  multiple \textsf{date} presentation in the two Chicago author-date
  styles, issues that particularly involved cross-referenced entries.
  In addition to some general fixes in the code, I have also slightly
  changed the functioning of the \texttt{cmsdate=both} and
  \texttt{cmsdate=on} switches.  If, and only if, a work has only one
  date, and there is no \texttt{switchdates} in the \textsf{options}
  field, then \texttt{cmsdate=on} and \texttt{cmsdate=both} will both
  result in the suppression of the \textsf{extrayear} field in that
  entry.  See the \textbf{date} field docs in
  section~\ref{sec:fields:authdate}, above.
\item Following a report by Antti-Juhani Kaijahano, I have modified
  the presentation of author-less \textbf{article} and \textbf{review}
  entries in the reference list of both Chicago author-date styles.
  If such a source had a \texttt{magazine} \textsf{entrysubtype}, the
  styles would already use the \textsf{journaltitle} at the head of
  the entry in the list of references, but if there was no
  \textsf{entrysubtype} the entry would appear in the list
  \textsf{date} first.  Now, in keeping with the \emph{Manual}
  (14.175), the \textsf{title} will appear first, in both reference
  lists and in-text citations.  See especially under \textbf{article}
  in section~\ref{sec:types:authdate}, above.
\item Several users have pointed out annoying formatting errors in the
  styles.  Evan Cortens spotted two bugs in the notes \&\ bibliography
  style, one of which, under various circumstances, introduced extra
  spaces into long notes and the other of which affected the
  formatting of the \textsf{type} field in \textsf{thesis} entries.  I
  have fixed both, also applying the latter fix to several other entry
  types that use the \textsf{type} field.  Bertold Schweitzer pointed
  out a formatting bug with the \textsf{issuesubtitle} field in the
  author-date style, now fixed.  Mark Sprevak reported some spurious
  spaces appearing in headers and footers when using the
  \textsf{titleps} package; the culprits were errors in the
  \textsf{cms-*.lbx} files, now cleaned up.
\item I have rectified a number of other errors, in particular making
  the automatic provision of abbreviated cross-references more robust
  in \textsf{incollection}, \textsf{inproceedings}, and
  \textsf{letter} entries, improving the behavior of the
  \textsf{postnote} field in certain corner cases, fixing bugs in the
  handling of \textsf{pagination} and \textsf{bookpagination} fields,
  and slightly altering the placement of the \textsf{addendum} field
  in book-like entries to bring it closer to the \emph{Manual's}
  specification.  A number of other, smaller improvements should also
  bring the styles into closer conformity with the specification.
\end{itemize}

\textbf{0.9.9b: Released December 6, 2012}
\begin{itemize}
\item This release contains a new variant of the author-date style,
  available as the\break \mycolor{\texttt{authordate-trad}} option
  when loading \textsf{biblatex-chicago}.  This provides the
  traditional, plain, pre-16th-edition Chicago title handling ---
  sentence-style capitalization, absence of quotation marks in
  \textsf{article} titles and the like --- but in all other respects
  follows the 16th-edition specification, as suggested by the
  \emph{Manual} (15.45).  Remember that the \texttt{headline} package
  option can be used to turn off the automatic sentence-style
  capitalization, meaning that titles will appear as presented in the
  .bib file, at least as far as capitalization is concerned.  Please
  see especially the documentation of \textsf{\textbf{title}} in
  section~\ref{sec:fields:authdate}, above, for the details.
\item I have updated calls to \cmd{DeclareLabelname} and
  \cmd{DeclareLabelyear} in several .cbx files so that the package
  works correctly with the most recent version (2.4) of
  \textsf{biblatex}.
\item Following a request by Norman Gray, I have included a
  \cmd{textcite} (and a \cmd{textcites}) command in the notes \&\
  bibliography style for the first time.  Please see
  section~\ref{sec:citecommands}, above, for the details.
\item Following a request by Daniel Possenriede, I have added in all
  three 16th-edition styles a new switch, \mycolor{\texttt{only}}, to
  the \texttt{doi} option, which prints the \textsf{doi} when present
  and the \textsf{url} only when there is no \textsf{doi}.  The
  package default remains, however, \texttt{true}.
\item I am grateful to Baldur Kristinsson for providing an Icelandic
  localization file for \textsf{biblatex-chicago}, called
  \mycolor{\textsf{cms-icelandic.lbx}}.  You'll see if you look
  through it that it is still something of a work in progress, but it
  should cover most needs in that language very well.  If you would
  like to fill in some of the gaps please let me know.
\item I am also grateful to H{\aa}kon Malmedal for providing Norwegian
  localizations for \textsf{biblatex-chicago}, contained in the files
  \mycolor{\textsf{cms-norsk.lbx}},
  \mycolor{\textsf{cms-norwegian.lbx}}, and
  \mycolor{\textsf{cms-nynorsk.lbx}}.
\item I have added a new British localization
  (\mycolor{\textsf{cms-british.lbx}}) that should make it much
  simpler for users to produce documents adhering to that tradition.
  For further details on the usage of all these localizations please
  see section~\ref{sec:international}, above.
\item Several users have reported a bug that resulted in doubled
  bibstrings in certain contexts.  This happened only when using
  localizations for which \textsf{biblatex-chicago} didn't have
  explicit support, and it should now be fixed.
\item I have changed the way the 16th-edition author-date styles
  handle the \emph{Ibidem} mechanism.  In the absence of a
  \textsf{postnote} field you no longer get empty parentheses, but
  rather a standard in-text citation.  If you do have a
  \textsf{postnote} field, then only that will appear.
\end{itemize}

\textbf{0.9.9a: Released July 30, 2012}
\begin{itemize}
\item I have made a few changes to \textsf{biblatex-chicago.sty} to
  allow the package to work with the latest version (2.0) of
  \textsf{biblatex}.  In all other respects this release is identical
  to 0.9.9.  If you do use the package with \textsf{biblatex} 2.0,
  please let me know if there are issues I need to address.  Thanks to
  Charles Schaum for alerting me to some of them.
\end{itemize}

\textbf{0.9.9: Released July 5, 2012}

\mylittlespace Converting 15th-Edition .bib Files to Use the 16th
Edition:

\mylittlespace \textbf{Notes and Bibliography Style}

\begin{itemize}
\item The specification for \textbf{music} entries has been
  significantly altered for the new edition.  You no longer need to
  worry about the \texttt{\textcircledP} and \texttt{\textcopyright}
  signs in the \textsf{howpublished} field, which will be silently
  ignored, and the \textsf{pubstate} field now reverts to its usual
  function of identifying reprints or, in this case, reissues.  The
  spec really only requires a record label (\textsf{series}) and
  catalog number (\textsf{number}), though \textsf{publisher} is still
  available if you need it.  There is a new emphasis, finally, on the
  dating of musical recordings, so that the \textsf{eventdate} gives
  the recording date of a particular song or other portion of a
  recording, the \textsf{origdate} the recording date of an entire
  album, and the \textsf{date} the publishing date of that album.
  Please see the full documentation in section~\ref{sec:entrytypes},
  above.
\item The specification for \textbf{video} entries has also been
  clarified.  For television series, the episode and series numbers go
  in \textsf{booktitleaddon} instead of \textsf{titleaddon} and, as
  with \textsf{music} entries, the \textsf{eventdate} will hold the
  original broadcast date of such an episode, or perhaps the
  recording/performance date of, e.g., an opera on DVD.  The
  \textsf{origdate} will still hold the original release date of a
  film, and the \textsf{date} the publishing or copyright date of the
  medium you are referencing.  Please see the full documentation in
  section~\ref{sec:entrytypes}, above.
\item You should add \textbf{customc} entries to provide
  bibliographical cross-references from multiple pseudonyms back to
  the author's name.
\item In \textbf{suppbook} entries, the \emph{Manual} now requires you
  to provide the page range (in the \textsf{pages} field) for the
  specific part you are citing, e.g., an introduction, foreword, or
  afterword.%%%\enlargethispage{-2\baselineskip}
\item In \textbf{patent} entries, the \emph{Manual} now prefers
  sentence-style capitalization for titles, which you'll need to
  provide yourself by hand.
\item When a descriptive phrase is used as an \textsf{author}, you can
  now omit an initial definite or indefinite article, which will help
  with alphabetization in the bibliography.
\item A DOI is now preferred to a URL, if both are available.
\item On the same subject, a revision date (or similar) is preferred
  to an access date for online material.  You can use the new
  \textbf{userd} field to change the string introducing the
  \textsf{urldate}, which defaults to being an access date.
\item Special imprints are now separated from their parent press by a
  forward slash rather than a comma, so can just be added to the
  \textsf{publisher} field with the usual keyword \texttt{and}.
\item I have implemented a reasonable, less-flexible facsimile of the
  \textsf{Biber}-only command \mycolor{\cmd{DeclareLabelname}} which
  should work for those using any backend.  It allows
  \textsf{biblatex} to find a name for short notes outside the
  standard name fields, including, notably, in the \textsf{name[a-c]}
  fields.  This should reduce the instances where you need a
  \textsf{shortauthor} field to provide such a name.
\item The Chicago-specific setting of another \textsf{Biber}-only
  command, \mycolor{\cmd{DeclareSortingScheme=cms}}, allows
  non-standard fields to be considered by \textsf{biblatex's} sorting
  algorithms, which should reduce the instances where you need a
  \textsf{sortkey} or the like in your entries.  If you aren't using
  \textsf{Biber}, the package reverts to the standard \texttt{nty}
  sorting scheme.
\end{itemize}

\textbf{Author-Date Style}

\begin{itemize}
\item All title fields now follow the rules for the notes \&\
  bibliography style as far as punctuation, formatting, and
  capitalization are concerned.  \textsf{Biblatex-chicago-authordate}
  will deal with most of this automatically, but if you have any hand
  formatting of lowercase letters within curly braces in your .bib
  file, you will need to restore the headline-style capitalization
  there.  Also, you'll need to be more careful when you provide
  quotation marks inside titles, remembering to use \cmd{mkbibquote}
  so that punctuation can be brought inside nested quotation marks.
  These revisions will apply particularly to \textbf{title},
  \textbf{booktitle}, and \textbf{maintitle} fields.
\item The one exception to these rules is in \textbf{patent} entries,
  where sentence-style capitalization of the \textsf{title} is now
  specified.  You'll have to provide this by hand yourself, as in the
  notes \&\ bibliography style.
\item Because of these changes to title formatting, you'll need to
  observe the difference between \textbf{article} and \textbf{review}
  entries, where the latter contain generic, \enquote{Review of
    \ldots} titles and the former standard, specific titles.
\item The presentation of \textbf{shorthand} fields has changed.  You
  no longer need to use the \textbf{customc} entry type to include
  cross-references from shorthands to expansions in the list of
  references.  Now, simply using a \textsf{shorthand} field in an
  entry places that \textsf{shorthand} in citations and at the head of
  the entry in the list of references, where it will be followed by
  its expansion within parentheses.  The new system will require help
  with sorting in the reference list --- placing the
  \textsf{shorthand} also in a \textsf{sortkey} should do the trick.
\item On the subject of \textbf{customc} entries, the \emph{Manual}
  now recommends using cross-refer\-ences in several contexts,
  particularly when a single author uses more than one pseudonym.
  Adding \textsf{customc} entries makes this happen.
\item There have been significant changes when presenting book-like
  entries with more than one date.  If you are using the
  \texttt{cmsdate=on} option, or indeed simply placing the earlier
  date in the \textbf{date} field and the later one in
  \textbf{origdate}, the presentation will be the same as before, but
  you should understand that the \emph{Manual} no longer recommends
  this \textsf{origdate}-only style.  It prefers, instead, to present
  either the \textsf{date} alone or both dates in citations and at the
  head of reference list entries.  When presenting both dates, there
  is now no longer a choice between the \texttt{old} and \texttt{new}
  options for \texttt{cmsdate}, but only the \texttt{both} option.  If
  you have \texttt{old} or \texttt{new} in your .bib files, they will
  be treated as synonyms of \texttt{both}.
\item The specification for \textbf{music} entries has been
  significantly altered for the new edition.  You no longer need to
  worry about the \texttt{\textcircledP} and \texttt{\textcopyright}
  signs in the \textsf{howpublished} field, which will be silently
  ignored, and the \textsf{pubstate} field reverts to its more usual
  function of identifying reprints or, in this case, reissues.  The
  spec really only requires a record label (\textsf{series}) and
  catalog number (\textsf{number}), though \textsf{publisher} is still
  available if you need it.  There is a new emphasis, finally, on the
  dating of musical recordings, which means that such entries will fit
  better with the author-date style.  It also means that I have had to
  redefine the various date fields.  The \textsf{eventdate} gives the
  recording date of a particular song or other portion of a recording,
  the \textsf{origdate} the recording date of an entire album, and the
  \textsf{date} the publishing date of that album.  The earlier date
  is the one that will appear in citations and at the head of
  reference list entries.  Please see the full documentation in
  section~\ref{sec:types:authdate}, above.
\item The specification for \textbf{video} entries has also been
  clarified.  For television series, the episode and series numbers go
  in \textsf{booktitleaddon} instead of \textsf{titleaddon} and, as
  with \textsf{music} entries, the \textsf{eventdate} will hold the
  original broadcast date of such an episode, or perhaps the
  recording/performance date of, e.g., an opera on DVD.  The
  \textsf{origdate} will still hold the original release date of a
  film, and the \textsf{date} the publishing or copyright date of the
  medium you are referencing.  The earlier date, once again, is the
  one that will appear in citations and at the head of reference list
  entries.  Please see the full documentation in
  section~\ref{sec:types:authdate}, above.
\item In \textbf{suppbook} entries, the \emph{Manual} now requires you
  to provide the page range (in the \textsf{pages} field) for the
  specific part you are citing, e.g., an introduction, foreword, or
  afterword.
\item The author-date style now prefers longer bibstrings in the list
  of references, bringing it into line with the notes \&\ bibliography
  style.  Generally, the package will take care of this for you, but
  if you've been using abbreviated strings in \textsf{note} fields,
  for example, you may want to change them so that they conform with
  the strings the package provides.  In some circumstances the
  \cmd{partedit} macro, and its relatives, may help.  See
  section~\ref{sec:formatting:authdate}.
\item When a descriptive phrase is used as an \textsf{author}, you can
  now omit an initial definite or indefinite article, which will help
  with alphabetization in the bibliography.
\item A DOI is now preferred to a URL, if both are available.
\item On the same subject, a revision date (or similar) is preferred
  to an access date for online material.  You can use the new
  \textbf{userd} field to change the string introducing the
  \textsf{urldate}, which defaults to being an access date.
\item Special imprints are now separated from their parent press by a
  forward slash rather than a comma, so can just be added to the
  \textsf{publisher} field with the usual keyword \texttt{and}.
\item The 16th edition of the \emph{Manual} is less than enthusiastic
  about the use of \enquote{Anon.}\ as the \textsf{author}, preferring
  instead that the \textsf{title} or the \textsf{journaltitle} take
  its place.  If you do decide to get rid of \enquote{Anon.,} new
  facilities provided by \textsf{Biber} --- see next entry --- should
  mean that \textsf{biblatex} no longer requires assistance when
  alphabetizing such author-less entries.
\item The Chicago-specific setting of the \textsf{Biber}-only command,
  \mycolor{\cmd{DeclareSortingScheme\break =cms}}, allows non-standard
  fields to be considered by \textsf{biblatex's} sorting algorithms,
  which should reduce the instances where you need a \textsf{sortkey}
  or the like in your entries.
\item The Chicago-specific setting of the \textsf{Biber}-only command
  \mycolor{\cmd{DeclareLabelname}} allows \textsf{biblatex} to find a
  name (\enquote{\textsf{label}}) for citations outside the standard
  name fields, including, notably, in the \textsf{name[a-c]} fields.
  This should reduce the instances where you need a
  \textsf{shortauthor} field to provide such a name.
\end{itemize}

Other New Features:

\begin{itemize}
\item For reprinted books, you can now present more detailed
  publishing information about the original edition using the new
  \mycolor{\textbf{origlocation}} and \mycolor{\textbf{origpublisher}}
  fields.  You can also use the \textsf{origlocation} in
  \textsf{letter} or \textsf{misc} (with \textsf{entrysubtype})
  entries to identify where a published or unpublished letter was
  written.  These uses apply to both Chicago styles.
\item Thanks to a patch sent by Kazuo Teramoto, you can now take
  advantage of \textsf{biblatex's} facilities for citing
  \mycolor{\textbf{eprint}} resources.  There is also a new
  \mycolor{\texttt{eprint}} option, set to \texttt{true} by default,
  which controls the printing of this field in both Chicago styles.
  You can set the option both in the preamble and in the
  \textsf{options} field of individual entries.  The field will always
  print in \textbf{online} entries.
\item I have added a new citation command,
  \mycolor{\cmd{citejournal}}, to the notes \&\ bibliography style to
  allow you to present journal articles using an alternative short
  note form, which may be a clearer form of reference in certain
  circumstances.  Such short notes will present the name of the
  \textsf{author}, the \textsf{journaltitle}, and the \textsf{volume}
  number.
\item I have included a very slightly modified version of the standard
  \textsf{biblatex} \cmd{citeauthor} command, which may be useful for
  references to works from classical antiquity.
\item I have added a new \texttt{cmsdate=\mycolor{full}} switch to the
  author-date style, which only affects citations in the text, and
  means that a full date specification will appear there, rather than
  just the year.  If you follow the \emph{Manual's} recommendations
  concerning newspaper and magazine articles only appearing in running
  text and not in the reference list, this option will help.
\item I have added a new \mycolor{\texttt{avdate}} option to the
  author-date style, set to \texttt{true} by default in
  \textsf{biblatex-chicago.sty}.  This alters the default setting of
  \cmd{DeclareLabelyear} in \textbf{music}, \textbf{review}, and
  \textbf{video} entries to take account of specialized instructions
  in the \emph{Manual} for finding dates to appear in citations and at
  the head of reference list entries.  Setting \texttt{avdate=false}
  in the options when you load \textsf{biblatex-chicago} restores the
  default settings for all entry types.  See \texttt{avdate} in
  section~\ref{sec:authpreset}.
\item The \emph{Manual} has added recommendations for citing blogs,
  which generally will need an \textbf{article} entry with
  \texttt{magazine} \textsf{entrysubtype}.  You can identify a blog as
  such by placing \enquote{blog} in the \textsf{location} field.  If
  you want to cite a comment to a blog or to other online material,
  the \textbf{review} entry type, \textsf{entrysubtype}
  \texttt{magazine} will serve.  The \textbf{eventdate} dates the
  comment, and any time stamp that is required can go in
  \textsf{nameaddon}.  These instructions work in both specifications.
\item Photographs are no longer presented differently from other sorts
  of artworks so, in effect, in both styles, the \textbf{image} type
  is now a clone of \textbf{artwork}, though retained for backward
  compatibility.
\item Following a request by Kenneth Pearce, I have added new
  facilities for presenting \textbf{shorthands} in both Chicago
  styles.  In both, there are two new \texttt{bibenvironments} which
  you can set using the \texttt{env} option to the
  \cmd{printshorthands} command: \mycolor{\texttt{losnotes}} formats
  the list of shorthands so that it can be presented in a footnote,
  while \mycolor{\texttt{losendnotes}} does the same for endnotes.  In
  both styles, there is a new preamble option,
  \mycolor{\texttt{shorthandfull}}, which prints the full
  bibliographical information of each entry inside the list of
  shorthands, allowing such a list effectively to replace a
  bibliography or list of references.  In the author-date style, you
  need to set the \texttt{cmslos=false} option as well, in order for
  this to work.  In the notes \&\ bibliography style, I have added a
  new citation command, \mycolor{\cmd{shorthandcite}}, which prints
  the \textsf{shorthand} even in the first citation of a given work.
\item Following suggestions by Roger Hart, I have implemented three
  new field-exclusion options in the notes \&\ bibliography style.
  In all three cases, the field in question will always appear in the
  bibliography, but not in long notes, which may help to save space.
  The fields at stake are \textsf{addendum}, \textsf{note}, and
  \textsf{series}, controlled respectively by the new
  \mycolor{\texttt{addendum}}, \mycolor{\texttt{notefield}}, and
  \mycolor{\texttt{bookseries}} options.  All of these are set to
  \texttt{true} using the new \mycolor{\texttt{completenotes}} option
  in \textsf{chicago-notes.cbx}, but you can change the settings
  either in the preamble or in the \textsf{options} field of
  individual entries.  Please see the documentation of these options
  in section~\ref{sec:chicpreset}, above, for details on which entry
  types are excluded from their scope.
\item Thanks to a coding suggestion from Gildas Hamel, I have
  redefined the \cmd{bibnamedash} in \textsf{biblatex-chicago.sty},
  which should now by default look a little better in a wider variety
  of fonts.
\item At the request of Baldur Kristinsson, I have added
  \cmd{DeclareLanguageMapping} commands to
  \textsf{biblatex-chicago.sty} for all the languages
  \textsf{biblatex-chicago} currently provides.  If you load the style
  in the standard way, you no longer need to provide these mappings
  manually yourself.
\item I have improved the date handling in both styles, particularly
  with regard to date ranges.
\end{itemize}

\textbf{0.9.8d: Released November 15, 2011}
\begin{itemize}
\item Some minor fixes to both styles for compatibility with
  \textsf{biblatex} 1.7.
\item Kenneth Pearce found an error in the formatting of
  \textsf{bookinbook} titles in the author-date style's list of
  shorthands.  This should work properly now.
\item Jonathan Robinson spotted some inconsistencies in the way the
  notes \&\ bibliography style interacts with the \textsf{hyperref}
  package.  Following his suggestion, short notes now point to long
  notes when the latter are available, but to bibliography entries
  instead when you have set the \texttt{short} option.
\end{itemize}

\textbf{0.9.8c: Released October 12, 2011}
\begin{itemize}
\item Emil Salim pointed out some rather basic errors in the
  presentation of \textsf{inproceedings} and \textsf{proceedings}
  entries, errors that have been present from the first release of the
  style(s).  These should now, belatedly, have been put right.
\item Minor improvements to coding and documentation.
\end{itemize}

\textbf{0.9.8b: Released September 29, 2011}
\begin{itemize}
\item Bad Dates: Christian Boesch alerted me to some date-formatting
  errors produced when using the styles with the \texttt{german}
  option to \textsf{babel}.  A little further investigation revealed
  similar problems with \texttt{french}, and before long it became
  clear that date handling in \textsf{biblatex-chicago} was generally,
  and significantly, sub-optimal.  The whole system should now be more
  robust and more accurate.
\item The new date-handling code shouldn't require any changes to your
  .bib files, but users of the author-date style may want to have a
  look at the documentation of the \textsf{letter} and \textsf{misc}
  entry types, and of the four date fields, for some information about
  how the changes could simplify the creation of their databases.
\item Various other minor improvements.
\end{itemize}

\textbf{0.9.8a: Released September 21, 2011}
\begin{itemize}
\item Fixed a series of unsightly errors in the author-date style,
  discovered while working on the pending update to the 16th edition.
\item Fixed bugs uncovered in both the author-date and the notes \&\
  bibliography styles thanks to Charles Schaum's adventurous use of
  the \textsf{origyear} field.
\item Added two new bibstrings to the cms-*.lbx files to fix potential
  bugs in some of the audiovisual entry types.
\end{itemize}

\textbf{0.9.8: Released August 31, 2011}

\begin{itemize}
\item Starting with \textsf{biblatex} version 1.5, in order to adhere
  to the author-date specification you will need to use \textsf{Biber}
  to process your .bib files, as \textsc{Bib}\TeX\ (and its more
  recent variants) will no longer provide all the required features.
  Unfortunately, however, the current release of \textsf{Biber}
  (0.9.5) contains bugs that make it tricky to use with
  \textsf{biblatex-chicago}.  These bugs have been addressed in 0.9.6
  beta, which is available for various operating systems in the
  \texttt{development} subdirectory of your SourceForge mirror, e.g.,
  \href{http://www.mirrorservice.org/sites/download.sourceforge.net/pub/sourceforge/b/project/bi/biblatex-biber/biblatex-biber/development/binaries/}{UK
    mirror}.  (If, by the time you read this, \textsf{Biber} 0.9.6 has
  already been released, then so much the better.)  Please see the
  start of \textsf{cms-dates-sample.pdf} for more details.
\item The switch to \textsf{Biber} for the author-date specification
  means that \textsf{biblatex} now provides considerably enhanced
  handling of the various date fields.  I have attempted to document
  the relevant changes in \textsf{cms-dates-sample.pdf} and in the
  \textbf{date} discussion in section~\ref{sec:fields:authdate},
  above, but in my testing the only alterations I've so far had to
  make to my .bib files involve adhering more closely to the
  instructions for specifying date ranges.  \textsf{Biber} doesn't
  like \{\texttt{1968/75}\}, and will ignore it.  Either use
  \{\texttt{1968/1975}\} or use \{\texttt{1968-{}-75}\} in the
  \textsf{year} field.
\item In the notes \&\ bibliography style, and mainly in
  \textsf{article}, \textsf{letter}, \textsf{misc}, and
  \textsf{review} entries, previous releases of
  \textsf{biblatex-chicago} recommended using the \cmd{isdot} macro
  when you needed both to define a field and not have it appear in the
  printed output.  This mechanism no longer works in \textsf{biblatex}
  1.6, and while addressing the problem I realized that relying on it
  covered over some inconsistencies and bugs in my code, so from this
  release forward you will need to modify your .bib and .tex files to
  use other, more standard mechanisms to achieve the same ends, in
  particular the \cmd{headlesscite} commands and declaring
  \texttt{useauthor=false} in the \textsf{options} field.  Please
  consult the documentation in section~\ref{sec:formatcommands}, s.v.\
  \enquote{\cmd{isdot},} for a list of example entries where you can
  see these changes at work.
\end{itemize}

Other New Features:

\begin{itemize}
\item Fixed the \cmd{smartcite} citation command in, and added a
  \cmd{smartcites} command to, \textsf{chicago-notes.cbx}, so that the
  notes \&\ bibliography style no longer prints parentheses around
  citations produced using \cmd{autocite(s)} commands inside
  \cmd{footnote} commands.  Many thanks to Louis-Dominique Dubeau for
  pointing out this error.
\item Rembrandt Wolpert and Aaron Lambert pointed out an issue with a
  command (\cmd{lbx\\@fromlang}) that \textsf{biblatex} no longer
  defines, and Charles Schaum very kindly suggested a temporary
  workaround in a newsgroup post, a workaround that should no longer
  be necessary.
\item Version 1.6 of \textsf{biblatex} no longer allows you to
  redefine the \texttt{minnames} and \texttt{maxnames} options in the
  \cmd{printbibliography} command, so I've defined
  \texttt{minbibnames} and \texttt{maxbibnames} in
  \textsf{biblatex-chicago.sty}, instead.  These parameters have been
  available since version 1.1, so this is now the earliest version of
  \textsf{biblatex} that will work with the Chicago styles.  Of
  course, if the (Chicago-recommended) values of these options don't
  suit your needs, you can redefine them in your document preamble.
\end{itemize}

\textbf{0.9.7a: Released March 17, 2011}
\begin{itemize}
\item Added \cmd{smartcite} command to \textsf{chicago-notes.cbx} so
  that the notes \&\ bibliography style will work with
  \textsf{biblatex} 1.3.
\item Added bibstrings \texttt{byconductor} and \texttt{cbyconductor}
  to the .lbx files, mistakenly omitted in version 0.9.7.
\item Minor fixes to the docs.
\end{itemize}

\textbf{0.9.7: Released February 15, 2011}

\mylittlespace Obsolete and Deprecated Features:
\begin{itemize}
\item The \textbf{customa} and \textbf{customb} entry types are now
  obsolete.  Any such entries will be ignored.  Please change any that
  remain to \textbf{letter} and \textbf{bookinbook}, respectively.
\item If you still have any \textbf{customc} entries containing
  introductions, prefaces, or the like, please change them to
  \textbf{suppbook}.  I have recycled \textsf{customc} for another
  purpose, on which see below.
\end{itemize}

Other New Features:

\begin{itemize}
\item At the request of Johan Nordstrom, I have added three new
  audiovisual entry types to both styles, \textbf{audio},
  \textbf{music}, and \textbf{video}.  The documentation of
  \textsf{audio} in sections~ \ref{sec:entrytypes} and
  \ref{sec:types:authdate} above contains an overview of the three,
  and the details for each type are to be found under their individual
  headings.
\item I have transformed the \textbf{customc} entry type to enable
  alphabetized cross-referen\-ces --- the \enquote{c} is meant to be
  mnemonic --- to other, separate entries in a reference list or
  bibliography.  In particular, this facilitates cross-references to
  other names in a list, rather than to other works.  In author-date,
  in a procedure recommended by the \emph{Manual}, this now allows you
  to expand shorthands inside the reference list rather than in a list
  of shorthands.  In both styles, you can now provide a pointer to the
  main entry if a reader is looking an author up under, e.g., a
  pseudonym or other alternative name.
\item I have introduced the \textbf{userc} field, intended to simplify
  the printing of the cross-references provided by \textsf{customc}
  entries.  The standard \cmd{nocite} command works as well, but the
  additional mechanism may be more convenient in some circumstances.
\item You can now provide an \textbf{eventdate} in \textsf{music}
  entries to identify, e.g., a particular recording session.  It will
  be printed just after the \textsf{title}.
\item In the notes \&\ bibliography style, I have now implemented the
  \textbf{shorthandintro} field, which allows you to change the string
  introducing a shorthand in the first, long note.  It works just as
  it does in the standard \textsf{biblatex} styles.
\item At the request of Scot Becker, I have added six new
  field-exclusion options to both styles, all of which can be set both
  in the document preamble and/or in the \textsf{options} field of
  individual .bib entries.  Three of these --- \texttt{doi},
  \texttt{isbn}, and \texttt{url} --- are standard \textsf{biblatex}
  options, the others --- \texttt{bookpages}, \texttt{includeall}, and
  \texttt{numbermonth} --- are \textsf{chicago}-specific.  See the
  docs in sections~\ref{sec:chicpreset} and \ref{sec:authpreset},
  above.
\item At the request of Charles Schaum, I've added the
  \texttt{juniorcomma} option to both styles, which can be set in the
  document preamble and/or in the \textsf{options} field of individual
  entries.  It allows you to get the traditional comma between a
  surname and \enquote{Jr.} or \enquote{Sr.}
\item Fixed an old inaccuracy in the presentation of \enquote{Jr.} and
  \enquote{Sr.,} so that they now appear at the end of names printed
  surname first in bibliographies and reference lists.
\item Thanks to Andrew Goldstone, I fixed some old inaccuracies in the
  syntax of shortened notes and bibliography entries presenting
  multiple contributions to one multi-author (or single-author)
  volume.
\item I've altered the directory structure of the archive containing
  this release.  Files were multiplying, and look set to multiply
  still further, so I've copied the structure used by Lehman for
  \textsf{biblatex} itself.
\item Fixed an old bug, which I'd guess was triggered quite rarely, in
  the formatting of publication information in long notes.
\item Fixed another bug in author-date where the colon separating
  titles and subtitles was in the wrong font.  The \textsf{biblatex}
  \texttt{punctfont} option solved this.
\item Fixed a punctuation bug in \textsf{InReference} entries in the
  notes \&\ bibliography style.  Also fixed \textsf{title}
  presentation in \textsf{Reference} entries in author-date.
\item Fixed some inaccuracies in the tests establishing priority
  between \textsf{date} and \textsf{origdate} fields.  These arose
  when date ranges were involved, and it's possible I haven't yet
  addressed all possible permutations of the problem.
\item Added several new bibstrings to the \textsf{cms-*.lbx} files for
  the new audiovisual entry types.  This means that the
  \textsf{editortype} fields can now be set to \texttt{director},
  \texttt{producer}, or \texttt{conductor}, depending on your needs.
  You can also set the fields to \texttt{none}, which eliminates all
  identifying strings, and which is useful for identifying performers
  of various sorts.
\item Minor improvements to documentation.
\end{itemize}

\textbf{0.9.5a: Released September 7, 2010}
\begin{itemize}
\item Quick fix for an elementary and show-stopping mistake in
  \textsf{biblatex-chicago.sty}, a mistake disguised if you load
  \textsf{csquotes}, which I do in all my test files.  Mea culpa.
  Many thanks indeed to Israel Jacques and Emil Salim for pointing
  this out to me.
\end{itemize}

\textbf{0.9.5: Released September 3, 2010}

\mylittlespace Obsolete and Deprecated Features:
\begin{itemize}
\item All the custom entry types --- \textbf{customa},
  \textbf{customb}, and \textbf{customc} --- are now deprecated.  They
  will still work for the time being, but please be aware that in the
  next major release they will no longer function, at least not as you
  might be expecting.  Please change your .bib files to use
  \textbf{letter} (=\textbf{customa}), \textbf{bookinbook}
  (=\textbf{customb}), and \textbf{suppbook} (=\textbf{customc})
  instead.
\item If by some chance anyone is still using the old \cmd{custpunctc}
  macro, it is now obsolete.  It really shouldn't be needed, but let
  me know if I'm wrong.
\end{itemize}

Other New Features:
\begin{itemize}
\item The Chicago author-date style is now implemented in the
  package, and is fully documented in section~\ref{sec:authdate},
  above.
\item The default way of loading the style(s) has slightly changed.
  You should put either \texttt{notes} or \texttt{authordate} in the
  options to \textsf{biblatex-chicago}, e.g.:
  \begin{quote}
    \cmd{usepackage[authordate,more options%
     \,\ldots]\{biblatex-chicago\}}
  \end{quote}
\item With the addition of the second Chicago style, I have thought it
  appropriate to alter both the name of the package and the names of
  the files it contains.  The package is now \textsf{biblatex-chicago}
  instead of \textsf{biblatex-chicago-notes-df}, and the following
  files have been renamed:
  \begin{itemize}
  \item \textsf{chicago-notes-df.cbx} is now \textsf{chicago-notes.cbx}
  \item \textsf{chicago-notes-df.bbx} is now \textsf{chicago-notes.bbx}
  \item \textsf{sample.tex} is now \textsf{cms-notes-sample.tex}
  \item \textsf{sample.pdf} is now \textsf{cms-notes-sample.pdf}
  \item \textsf{chicago-test.bib} is now \textsf{notes-test.bib}
  \item \textsf{biblatex-chicago-notes-df.pdf} (this file) is now
    \textsf{biblatex-chicago.pdf}
  \end{itemize}
  The following files have been added:
  \begin{itemize}
  \item \textsf{chicago-authordate.cbx}
  \item \textsf{chicago-authordate.bbx}
  \item \textsf{cms-dates-sample.tex}
  \item \textsf{cms-dates-sample.pdf}
  \item \textsf{dates-test.bib}
  \end{itemize}
  The following files have retained their old names:
  \begin{itemize}
  \item \textsf{cms-american.lbx}
  \item \textsf{cms-french.lbx}
  \item \textsf{cms-german.lbx}
  \item \textsf{cms-ngerman.lbx}
  \item \textsf{biblatex-chicago.sty}
  \end{itemize}
\item I have implemented the \textsf{pubstate} field, slightly
  differently yet compatibly in the two styles, to provide a simpler
  mechanism for identifying a reprinted book.  In the author-date
  style, it is highly recommended you use it, as it sorts out some
  complicated formatting questions automatically.  In the notes \&\
  bibliography style it isn't strictly necessary, but may be useful
  anyway and easier to remember than the old system.  See the
  documentation under \textsf{pubstate} in
  sections~\ref{sec:entryfields} and \ref{sec:fields:authdate}, above.
\item Users of \textsf{biblatex-chicago-notes} no longer need a
  \textsf{shortauthor} field in author-less \textsf{manual} entries,
  or in author-less \textsf{article} or \textsf{review} entries with a
  \texttt{magazine} \textsf{entrysubtype}.  The package will now
  automatically take an author for short notes from the
  \textsf{organization} field for \textsf{manual} entries and from the
  \textsf{journaltitle} field for the others.  You can still use a
  \textsf{shortauthor} field if you want, but it's no longer
  necessary.  (This also holds for \textsf{chicago-authordate}.)
\item Date presentation in the \textsf{misc} entry type (with
  \textsf{entrysubtype}) has changed to fix an inaccuracy.  You can
  now use the \textsf{date} and \textsf{origdate} fields to
  distinguish between two sorts of archival source: letters and
  \enquote{letter-like} sources use \textsf{origdate}, interviews and
  other non-letters use \textsf{date}.  The only difference is in how
  the date is printed, so current .bib entries will continue to work
  fine, albeit with minor inaccuracies in the case of non-letter-like
  sources.  See the docs on \textbf{misc} in
  sections~\ref{sec:entrytypes} and \ref{sec:types:authdate}, above.
\item When only one date is presented in a \textsf{patent} entry ---
  either in the \textsf{date} or \textsf{origdate} field --- this will
  now always be used as the filing date.  In
  \textsf{biblatex-chicago-notes}, this makes a change from the
  previous (incorrect) behavior.
\item I have included the option \texttt{dateabbrev=false} in the
  default settings for \textsf{biblatex-chicago-notes}.  This ensures
  that the long month names are printed, as otherwise recent releases
  of \textsf{biblatex} print the abbreviated ones by default.
\item The provision of punctuation in \textsf{entrysubtype}
  \texttt{classical} entries has been improved, allowing the comma to
  appear before certain kinds of location specifiers even when citing
  works by their traditional divisions.  See \emph{Manual} 17.253.
  (This applies to both Chicago styles.)
\item The \textsf{number} field in \textsf{article},
  \textsf{periodical}, and \textsf{review} entries now allows you to
  include a series or range of numbers in the field, with the style
  automatically providing the correct bibstring (singular or plural).
\item I have removed and altered bibstrings in the .lbx files to take
  advantage of the new \cmd{bibsstring} and \cmd{biblstring} commands
  in \textsf{biblatex}, and added one new string
  (\texttt{origpubyear}) needed by
  \textsf{biblatex-chicago-authordate}.
\end{itemize}

\textbf{0.9a: Released March 20, 2010}
\begin{itemize}
\item Quick fixes for compatibility with \textsf{biblatex} 0.9a.
\end{itemize}

\textbf{0.9: Released March 18, 2010}

\mylittlespace Obsolete and Deprecated Features:
\begin{itemize}
\item The \textbf{userd} field is now obsolete.  All information it
  used to hold should be placed in the \textsf{edition} field.
\item The \textbf{origyear} field is now obsolete in
  \textsf{biblatex}.  It has been replaced by \textbf{origdate}, and
  because the latter allows a full date specification, I have been
  able to make the operation of \textsf{customa} (=\,\textsf{letter}),
  \textsf{misc} (with an \textsf{entrysubtype}), and \textsf{patent}
  entries more intuitive.  The RELEASE file contained in this package
  gives the short instructions on how to update your .bib files, and
  you can also consult the documentation of those entry types above.
\item The modified \textsf{csquotes.cfg} file I provided in earlier
  releases is now obsolete, and has been removed from the package.
  Please upgrade to the latest version of \textsf{csquotes} and, if
  you are still using my modified .cfg file, remove it from your \TeX\
  search path, or at the very least excise the code I provided.
\end{itemize}

Other New Features:
\begin{itemize}
\item Added the files \textsf{cms-german.lbx} (with its clone
  \textsf{cms-ngerman.lbx}) and \textsf{cms-french\break .lbx}, which
  allow the creation of Chicago-like references in those languages.
  See section \ref{sec:international} above for details on usage.
\item Added the \texttt{annotation} package option to allow the
  creation of annotated bibliographies.  This code is still not
  entirely polished yet, but it is usable.  Please see page
  \pageref{sec:annote} above for instructions and hints.
\item Added \textsf{biblatex's} new \textbf{bookinbook} entry type,
  which currently functions as an alias of the \textsf{customb} type.
  As \textsf{biblatex} now provides standard equivalents for all of
  the custom types I initially found it necessary to provide ---
  \textsf{letter}~= \textsf{customa}, \textsf{bookinbook}~=
  \textsf{customb}, and \textsf{suppbook} \& \textsf{suppcollection}~=
  \textsf{customc} --- it may soon be time to prune out the custom
  types to enhance compatibility with other \textsf{biblatex} styles.
  I shall give plenty of warning before I do so.
\item In line with the new system adopted in \textsf{biblatex} 0.9,
  using the \textsf{editortype} field turns off the usual string
  concatenation mechanisms of the Chicago style.  See Lehman's RELEASE
  file for a discussion of this.
\item I have added support for the new \textsf{editor[a--c]} and
  \textsf{editor[a--c]type} fields, and they work just as in standard
  \textsf{biblatex}, though I'm uncertain how much use they'll get
  from users of the Chicago style.
\item I have added many bibstrings to the .lbx files to help with
  internationalization.  The new ones that you might want to use in
  your .bib files include: \texttt{pseudonym}, \texttt{nodate},
  \texttt{revisededition}, \texttt{numbers}, and \texttt{reviewof}.
  Please see section~\ref{sec:international} for a fuller list.
\end{itemize}

\textbf{0.8.9d: Released February 17, 2010}
\begin{itemize}
\item Chris Sparks and Aaron Lambert both found formatting bugs in the
  0.8.9c code.  I've fixed these bugs, and am releasing this version
  now, the last in the 0.8.9 series.  The next release of
  \textsf{biblatex-chicago-notes-df}, due as soon as possible, will
  contain many more significant changes, including those necessary for
  it to function properly with the recently-released \textsf{biblatex}
  version 0.9.  In the meantime, at least version 0.8.9d should produce
  more accurate output.
\end{itemize}

\textbf{0.8.9c: Released November 4, 2009}
\begin{itemize}
\item Emil Salim noticed that the \emph{ibidem} mechanism wasn't
  working properly, printing the page number after \enquote{Ibid} even
  when the page reference of the preceding citation was identical.
  The fix for this involved setting \texttt{loccittracker=constrict}
  in \textsf{biblatex-chicago.sty}, something you'll have to do
  manually yourself if you're loading the package via a call to
  \textsf{biblatex} rather than to \textsf{biblatex-chicago}.
\item Several users have reported unwanted behavior when repeated
  names in bibliographies are replaced with the \texttt{bibnamedash}.
  This release should fix both when the \texttt{bibnamedash} appears
  and what punctuation follows it.
\end{itemize}

\textbf{0.8.9b: Released September 9, 2009}
\begin{itemize}
\item Fixed a long-standing bug in formatting names in the
  bibliography.  The package now correctly places a comma after the
  reversed name that begins the entry, using \textsf{biblatex's}
  \cmd{revsdnamedelim} command.  Many thanks to Johanna Pink for
  catching my rather egregious error.
\item While fixing some formatting errors that cropped up when using
  the newest version of \textsf{biblatex} (0.8h at time of writing), I
  also spotted some more venerable bugs in the code for using
  shortened cross-references for citing multiple entries in a
  collection of essays or letters.  I believe this now works
  correctly, but please let me know if you discover differently.
\item Joseph Reagle noticed that endnote marks (produced using the
  \textsf{endnotes} package) did not receive the
  same treatment as footnote marks.  I have rectified this, placing
  the code in \textsf{biblatex-chicago.sty} so that you can turn it
  off either by using the old package-loading system or by setting the
  \texttt{footmarkoff} package option when loading
  \textsf{biblatex-chicago}.
\item Updates to Lehman's \textsf{csquotes} package have rendered my
  modifications in \textsf{csquotes.\break cfg} obsolete.  Please use
  the latest version of \textsf{csquotes} (4.4a at time of writing)
  and ignore my file, which will disappear in a later release.
\item At the request of Will Small, I have included some code, still
  in an alpha state, to allow you to specify, in the bibliography, the
  original publication details of essays which you are citing from
  later reprints (a \emph{Collected Essays} volume, for example).  See
  the documentation above under the \textsf{\mycolor{reprinttitle}}
  field if you would like to test this functionality.
\end{itemize}

\textbf{0.8.9a: Released July 5, 2009}
\begin{itemize}
\item Slight changes for compatibility with \textsf{biblatex} 0.8e.
  The package still works with 0.8c and 0.8d, as well.
\end{itemize}

\textbf{0.8.9: Released July 2, 2009}

\mylittlespace Obsolete and Deprecated Features:
\begin{itemize}
\item The \textbf{single-letter bibstrings} (\cmd{bibstring\{a\}},
  \cmd{bibstring\{b\}}, etc.) are now obsolete.  You should replace
  any still present in your .bib file with \cmd{autocap} commands ---
  see \S~3.8.4 of \textsf{biblatex.pdf}.
\end{itemize}

Other New Features:
\begin{itemize}
\item The default way of loading the package is now with

  \cmd{usepackage[further-options]\{biblatex-chicago\}}

  rather than

  \cmd{usepackage[style=chicago-notes-df,further-options]\{biblatex\}}.

  Please see section~\ref{sec:loading} above for details and hints.
\item Package-specific bibstrings have been removed from the .cbx and
  .bbx files and are now gathered in a new file,
  \textbf{cms-american.lbx}, which changes the way the package
  interacts with \textbf{babel}.  It is now somewhat simpler if you
  want the defaults, but somewhat more complex if you require
  non-standard features.  Please see section~\ref{sec:otherpacks}
  above for more details.
\item Two new entry types have been added: \textbf{artwork} for works
  of visual art excluding photographs, and \textbf{image} for
  photographs.  See the documentation of \textsf{artwork} for how to
  create .bib entries for both types.
\item Added the new bibliography and entry option
  \textbf{usecompiler}, set to \texttt{true} by default.  This
  streamlines the code that finds a name to head an entry
  (\textbf{author -> editor [or namea] -> translator [or nameb] ->
    compiler [namec] -> title}).  The whole system should work more
  consistently now, but do see the \textsf{author} and \textsf{namec}
  documentation for improved notes on how to use it.
\item Added the new bibliography option \textbf{footmarkoff}, to turn
  off the optional in-line (as opposed to superscript) formatting of
  the marks in foot- or endnotes.  You only need this if you load the
  package with the new default \cmd{usepackage\{biblatex-chicago\}};
  users loading it the old way get default \LaTeX\ formatting.
\item At Matthew Lundin's request, I have added the citation command
  \textbf{\textbackslash headlesscite}, which works like
  \cmd{headlessfullcite} but allows \textsf{biblatex} to decide
  whether to print the full or the short note.
\item Fully adopted \textsf{biblatex's} system for providing
  end-of-entry punctuation, which will solve some of the bugs users
  have been finding.  See section~\ref{sec:otherhints}, above, and do
  please let me know if inconsistencies remain.
\item Added a modified \textbf{csquotes.cfg} file to address issues
  users were having when using the \textbf{Xe\LaTeX} engine in
  combination with \textsf{biblatex-chicago}.  See
  section~\ref{sec:otherpacks}, above.
\item Added \texttt{natbib} option to allow users of the default setup
  to continue to benefit from \textsf{biblatex's} \textsf{natbib}
  compatibility code.  Thanks to Bennett Helm for pointing out this
  issue.
\item Added a \textbf{shorthandibid} option to allow the printing of
  \emph{ibid.}\ in consecutive references to an entry that contains a
  \textsf{shorthand} field.  Thanks to Chris Sparks for calling my
  attention to this problem.
\item While investigating the preceding, I noticed failures when
  combining the \texttt{short} option with a \textsf{shorthand} field.
  The package now actually does what it has always claimed to do under
  \textbf{shorthand}.
\item Many small bug fixes and improvements to the documentation.
\end{itemize}

To Do:

\begin{itemize}
\item The shorthand vs \emph{ibid.}\ question may need more careful
  addressing in some cross references, and also in relation to the
  \texttt{noibid} package option.
\item Charles Schaum has quite rightly pointed out the inconsistency
  in my naming conventions --- \textsf{biblatex-chicago.sty} as
  opposed to \textsf{chicago-notes-df.cbx}, for example.  I'm going to
  delay a decision on which way to go with this until a later release.
\end{itemize}

\textbf{0.8.5a: Released June 14, 2009}

\begin{itemize}
\item Quick and dirty fixes to bibliography strings to allow
  compatibility with \textsf{biblatex} version 0.8d.  If you are still
  using 0.8c, then I would wait for the next version of
  \textsf{biblatex-chicago-notes-df}, which is due soon.  See README.
\end{itemize}

\textbf{0.8.5: Released January 10, 2009}

\mylittlespace Obsolete and Deprecated Features:

  \begin{itemize}
  \item The \textbf{\textbackslash custpunct} commands are now
    deprecated --- Lehman's \enquote{American} punctuation tracking
    facilities should handle quoted text automatically, assuming you
    remember always to use \textbf{\textbackslash mkbibquote} in your
    database.  If you still need \cmd{custpunct}, please let me know,
    as it may be an error in the style.
  \item With \cmd{custpunct} no longer needed, the toggles activated
    by placing \enquote{\texttt{plain}} in the \textbf{type} or
    \textbf{userb} fields are also deprecated.
  \end{itemize}

Other New Features:

\begin{itemize}
\item At least \textbf{biblatex 0.8b} is now required --- 0.8c works
  fine, as well.
\item I now \emph{strongly recommend} that you use \textbf{babel} with
  \enquote{\texttt{american}} as the main text language.  See
  section~\ref{sec:otherpacks} above for further details.
\item The \textbf{customc} entry type has been revised, allowing you
  to cite any sort of supplementary material using the \textbf{type}
  field instead of relying on toggles in the \textsf{introduction},
  \textsf{afterword}, and \textsf{foreword} fields, though these
  latter still work.  The two new entry types \textbf{suppbook} and
  \textbf{suppcollection} are both aliased to \textsf{customc}, and
  therefore work in exactly the same way.
\item The new entry type \textbf{suppperiodical} is aliased to
  \textbf{review}.
\item The new entry type \textbf{letter} is aliased to
  \textbf{customa}.
\item In \textbf{inreference} entries the \textsf{postnote} field of
  all \cmd{cite} commands is now treated like data in \textsf{lista},
  that is, it will be placed within quotation marks and prefaced with
  the appropriate string.  The only difference is that you can only
  put one such article name in \textsf{postnote}, as it isn't a list
  field.
\item I've set the new \textsf{biblatex} option \texttt{usetranslator}
  to \texttt{true} by default, which means entries will automatically
  be alphabetized by their \textsf{translator} in the absence of an
  \textsf{author} or an \textsf{editor}.
\item A host of small formatting errors were eliminated, nearly all of
  them through adopting Lehman's punctuation tracker.
\item In the main body of this documentation, I've added some
  \mycolor{\textbf{color coding}} to help you more quickly to identify
  entry types and fields that are either new or that have undergone
  significant revision.
\end{itemize}

To Do:

\begin{itemize}
\item Separate out \enquote{options} from the basic citation
  \enquote{style,} using a \LaTeX\ style file.  This is an
  architectural change recommended by Lehman.
\end{itemize}

\textbf{0.8.2.2: Released November 24, 2008}

\begin{itemize}
\item Fixed spurious commas appearing in some bibliography entries,
  spotted by Nick Andrewes.  While investigating this I noticed a more
  general problem with punctuation after italicized titles ending with
  question marks or exclamation points.  This will be addressed in
  forthcoming revisions both of \textsf{biblatex} and of this package.
\item Nick also reported some problems with spurious punctuation in
  the bibliography when using XeLaTeX.  I haven't yet been able to pin
  down the exact cause of these, but if you are using XeLaTeX and are
  having (or have solved) similar problems I'd be interested to hear
  from you.
\end{itemize}

\textbf{0.8.2: Released November 3, 2008}

\begin{itemize}
\item Fixed several formatting glitches between citations in multicite
  commands (spotted by Joseph Reagle) and also after some prenotes. 
\end{itemize}

\textbf{0.8.1: Released October 22, 2008}

\mylittlespace Obsolete and Deprecated Features:

\begin{itemize}
\item The \textbf{origlocation} field is now obsolete, and has been
  replaced by \textbf{lista}.  Please update your .bib files
  accordingly.
\item The single-letter \textbf{\textbackslash bibstring} commands I
  provided in version 0.7 are now deprecated.  In most cases, you'll
  be able to take advantage of the automatic contextual capitalization
  facilities introduced in this release, but if you still need the
  single-letter \cmd{bibstring} functionality then you should switch
  to \cmd{autocap}, as I shall be removing the single-letter
  \texttt{bibstrings} in a future release.  See above under
  \textbf{\textbackslash autocap} for all the details.
\item The \textbf{userd} field is now deprecated, as \textsf{biblatex}
  0.8 allows all forms of data to be included in the \textsf{edition}
  field.  I shall be removing \textsf{userd} in a future release, so
  please update your .bib files as soon as is convenient.
\end{itemize}

Other New Features:

\begin{itemize}
\item Updated the .bbx and .cbx files to work with \textsf{biblatex}
  0.8.  This most recent version of \textsf{biblatex} is now required
  for \textsf{biblatex-chicago-notes-df} to work.
\item Added the \textbf{usera} field, which holds supplemental
  information about a \textsf{journaltitle} in \textsf{article} and
  \textsf{review} entries.  See the documentation of the field for
  details.
\item Added the \textbf{\textbackslash citetitles} multicite command
  to fix a problem with spurious punctuation when multiple titles were
  listed.
\item Added the \textbf{\textbackslash Citetitle} command to help with
  automatic capitalization of titles when they occur at the beginning
  of a note.
\item Minor punctuation fixes in \textsf{biblatex-chicago-notes-df.bbx}.
\end{itemize}

To Do:

\begin{itemize}
\item Integrate \textsf{biblatex's} American punctuation facilities.
\item Separate out \enquote{options} from the basic citation
  \enquote{style,} using a \LaTeX\ style file.  This is an
  architectural change recommended by Lehman.
\item Investigate and possibly integrate the new entry types provided
  in \textsf{biblatex} 0.8.
\end{itemize}

\textbf{0.7: First public release, September 18, 2008}

\end{document}

%%% Local Variables:
%%% mode: latex
%%% TeX-master: t
%%% TeX-engine: xetex
%%% End:
